% Auto-generated from data/segments.json + layout.json (+ transforms) — do not edit by hand.
% volume: 11Ma03
% mode: reading
% segments: 1661
% transform_rules: 0
% toc_marks: 126

% Volume layout defaults (layout.json ``layout``).
\csromanlayoutapply{1}{1.5}{21.6pt}{5pt}{6.3pt}{65pt}{2.5em}%

% Reading physical-page layout (layout.json ``page_layout_reading_mode``).
\csromanreadingpagelayoutvolume{1}{1.5}{21.6pt}{5pt}{6.3pt}{65pt}{2.5em}%
\csromanreadingpagelayoutenable

\csromanheader{2}{\textbf{มชฺฌิมนิกาย}}
\csromanmatikaanchor{mtk.0}
\csromantocmarknopage{chapter}{อุปริปณฺณาสปาฬิ}{mtk.0}
\bookmark[dest={mtk.0},level=0]{อุปริปณฺณาสปาฬิ}
\gambhira{\textbf{อุปริปณฺณาสปาฬิ}}
\namakkaram{นโม ตสฺส ภควโต อรหโต สมฺมาสมฺพุทฺธสฺส.}
\csromanmatikaanchor{mtk.1}
\csromantocmarknopage{chapter}{1. เทวทหวคฺค}{mtk.1}
\bookmark[dest={mtk.1},level=0]{1. เทวทหวคฺค}
\chapterhead{1. เทวทหวคฺค}
\csromanmatikaanchor{mtk.2}
\csromantocmarkref{subsection}{1. เทวทหสุตฺต}{mtk.2}
\bookmark[dest={mtk.2},level=2]{1. เทวทหสุตฺต}
\csromanheader{2}{1. เทวทหสุตฺต}
\proseitem{1}{\csromanfolio{1}เอวํ เม สุตํ\csromandash{}เอกํ สมยํ ภควา สกฺเกสุ วิหรติ เทวทหํ นาม สกฺยานํ นิคโม.\csromanspacer{} ตตฺร โข ภควา ภิกฺขู อามนฺเตสิ “ภิกฺขโว”ติ. “ภทนฺเต”ติ เต ภิกฺขู ภควโต ปจฺจสฺโสสุํ.\csromanspacer{} ภควา เอตทโวจ\csromandash{}สนฺติ ภิกฺขเว เอเก สมณพฺราหฺมณา เอวํวาทิโน เอวํทิฏฺฐิโน “ยํ กิญฺจายํ ปุริสปุคฺคโล ปฏิสํเวเทติ สุขํ วา ทุกฺขํ วา อทุกฺขมสุขํ วา, สพฺพํ ตํ ปุพฺเพกตเหตุ.\csromanspacer{} อิติ ปุราณานํ กมฺมานํ ตปสา พฺยนฺตีภาวา นวานํ กมฺมานํ อกรณา อายติํ อนวสฺสโว, อายติํ อนวสฺสวา กมฺมกฺขโย, กมฺมกฺขยา ทุกฺขกฺขโย, ทุกฺขกฺขยา เวทนากฺขโย, เวทนากฺขยา สพฺพํ ทุกฺขํ นิชฺชิณฺณํ ภวิสฺสตี”ติ.\csromanspacer{} เอวํวาทิโน ภิกฺขเว นิคณฺฐา.}

\prose{เอวํวาทาหํ ภิกฺขเว นิคณฺเฐ อุปสงฺกมิตฺวา เอวํ วทามิ “สจฺจํ กิร ตุเมฺห อาวุโส นิคณฺฐา เอวํวาทิโน เอวํทิฏฺฐิโน ‘ยํ กิญฺจายํ ปุริสปุคฺคโล ปฏิสํเวเทติ สุขํ วา ทุกฺขํ วา อทุกฺขมสุขํ วา, สพฺพํ ตํ ปุพฺเพกตเหตุ.\csromanspacer{} อิติ ปุราณานํ กมฺมานํ ตปสา พฺยนฺตีภาวา นวานํ กมฺมานํ อกรณา อายติํ อนวสฺสโว, อายติํ อนวสฺสวา กมฺมกฺขโย, กมฺมกฺขยา ทุกฺขกฺขโย, ทุกฺขกฺขยา เวทนากฺขโย, เวทนากฺขยา สพฺพํ ทุกฺขํ นิชฺชิณฺณํ ภวิสฺสตี’ติ”.\csromanspacer{} เต จ เม ภิกฺขเว นิคณฺฐา เอวํ ปุฏฺฐา “อามา”ติ ปฏิชานนฺติ.}

\prose{\csromanfolio{2}ตฺยาหํ เอวํ วทามิ “กิํ ปน ตุเมฺห อาวุโส นิคณฺฐา ชานาถ ‘อหุวเมฺหว มยํ ปุพฺเพ, น นาหุวมฺหา’ติ”.\csromanspacer{} โน หิทํ อาวุโส.}

\prose{“กํ ปน ตุเมฺห อาวุโส นิคณฺฐา ชานาถ ‘อกรเมฺหว มยํ ปุพฺเพ ปาปกมฺมํ, น นากรมฺหา’ติ”.\csromanspacer{} โน หิทํ อาวุโส.}

\prose{“กิํ ปน ตุเมฺห อาวุโส นิคณฺฐา ชานาถ ‘เอวรูปํ วา เอวรูปํ วา ปาปกมฺมํ อกรมฺหา’ติ”.\csromanspacer{} โน หิทํ อาวุโส.}

\prose{“กิํ ปน ตุเมฺห อาวุโส นิคณฺฐา ชานาถ ‘เอตฺตกํ วา ทุกฺขํ นิชฺชิณฺณํ, เอตฺตกํ วา ทุกฺขํ นิชฺชีเรตพฺพํ, เอตฺตกมฺหิ วา ทุกฺเข นิชฺชิณฺเณ สพฺพํ ทุกฺขํ นิชฺชิณฺณํ ภวิสฺสตี’ติ”.\csromanspacer{} โน หิทํ อาวุโส.}

\prose{“กิํ ปน ตุเมฺห อาวุโส นิคณฺฐา ชานาถ ‘ทิฏฺเฐว ธมฺเม อกุสลานํ ธมฺมานํ ปหานํ กุสลานํ ธมฺมานํ อุปสมฺปทนฺ’ติ”.\csromanspacer{} โน หิทํ อาวุโส.}

\proseitem{2}{อิติ กิร ตุเมฺห อาวุโส นิคณฺฐา น ชานาถ “อหุวเมฺหว มยํ ปุพฺเพ, น นาหุวมฺหา”ติ.\csromanspacer{} น ชานาถ “อกรเมฺหว มยํ ปุพฺเพ ปาปกมฺมํ, น นากรมฺหา”ติ.\csromanspacer{} น ชานาถ “เอวรูปํ วา เอวรูปํ วา ปาปกมฺมํ อกรมฺหา”ติ.\csromanspacer{} น ชานาถ “เอตฺตกํ วา ทุกฺขํ นิชฺชิณฺณํ, เอตฺตกํ วา ทุกฺขํ นิชฺชีเรตพฺพํ, เอตฺตกมฺหิ วา ทุกฺเข นิชฺชิณฺเณ สพฺพํ ทุกฺขํ นิชฺชิณฺณํ ภวิสฺสตี”ติ.\csromanspacer{} น ชานาถ\csromandash{}ทิฏฺเฐว ธมฺเม อกุสลานํ ธมฺมานํ ปหานํ กุสลานํ ธมฺมานํ อุปสมฺปทํ.\csromanspacer{} เอวํ สนฺเต อายสฺมนฺตานํ นิคณฺฐานํ น กลฺลมสฺส เวยฺยากรณาย “ยํ กิญฺจายํ ปุริสปุคฺคโล ปฏิสํเวเทติ สุขํ วา ทุกฺขํ วา อทุกฺขมสุขํ วา, สพฺพํ ตํ ปุพฺเพกตเหตุ.\csromanspacer{} อิติ ปุราณานํ กมฺมานํ ตปสา พฺยนฺตีภาวา นวานํ กมฺมานํ อกรณา อายติํ อนวสฺสโว, อายติํ อนวสฺสวา กมฺมกฺขโย, กมฺมกฺขยา ทุกฺขกฺขโย, ทุกฺขกฺขยา เวทนากฺขโย, เวทนากฺขยา สพฺพํ ทุกฺขํ นิชฺชิณฺณํ ภวิสฺสตี”ติ.}

\prose{สเจ ปน ตุเมฺห อาวุโส นิคณฺฐา ชาเนยฺยาถ “อหุวเมฺหว มยํ ปุพฺเพ, น นาหุวมฺหา”ติ.\csromanspacer{} ชาเนยฺยาถ “อกรเมฺหว มยํ ปุพฺเพ ปาปกมฺมํ, น นากรมฺหา”ติ.\csromanspacer{} ชาเนยฺยาถ “เอวรูปํ วา เอวรูปํ วา ปาปกมฺมํ อกรมฺหา”ติ.\csromanspacer{} ชาเนยฺยาถ “เอตฺตกํ วา ทุกฺขํ นิชฺชิณฺณํ, เอตฺตกํ วา ทุกฺขํ นิชฺชีเรตพฺพํ, เอตฺตกมฺหิ วา ทุกฺเข นิชฺชิณฺเณ สพฺพํ ทุกฺขํ นิชฺชิณฺณํ ภวิสฺสตี”ติ.\csromanspacer{} ชาเนยฺยาถ\csromandash{}ทิฏฺเฐว ธมฺเม อกุสลานํ ธมฺมานํ ปหานํ กุสลานํ \csromanfolio{3}ธมฺมานํ อุปสมฺปทํ.\csromanspacer{} เอวํ สนฺเต อายสฺมนฺตานํ นิคณฺฐานํ กลฺลมสฺส เวยฺยากรณาย “ยํ กิญฺจายํ ปุริสปุคฺคโล ปฏิสํเวเทติ สุขํ วา ทุกฺขํ วา อทุกฺขมสุขํ วา, สพฺพํ ตํ ปุพฺเพกตเหตุ.\csromanspacer{} อิติ ปุราณานํ กมฺมานํ ตปสา พฺยนฺตีภาวา นวานํ กมฺมานํ อกรณา อายติํ อนวสฺสโว, อายติํ อนวสฺสวา กมฺมกฺขโย, กมฺมกฺขยา ทุกฺขกฺขโย, ทุกฺขกฺขยา เวทนากฺขโย, เวทนากฺขยา สพฺพํ ทุกฺขํ นิชฺชิณฺณํ ภวิสฺสตี”ติ.}

\proseitem{3}{เสยฺยถาปิ อาวุโส นิคณฺฐา ปุริโส สลฺเลน วิทฺโธ อสฺส สวิเสน คาฬฺหูปเลปเนน\footnote{คาฬฺหปเลปเนน (ก)}, โส สลฺลสฺสปิ เวธนเหตุ\footnote{เวทนาเหตุ (สี, อิ, ก)} ทุกฺขา ติพฺพา\footnote{ติปฺปา (สี, สฺยา, กํ, อิ)} กฏุกา เวทนา เวทิเยยฺย.\csromanspacer{} ตสฺส มิตฺตามจฺจา ญาติสาโลหิตา ภิสกฺกํ สลฺลกตฺตํ อุปฏฺฐาเปยฺยุํ.\csromanspacer{} ตสฺส โส ภิสกฺโก สลฺลกตฺโต สตฺเถน วณมุขํ ปริกนฺเตยฺย, โส สตฺเถนปิ วณมุขสฺส ปริกนฺตนเหตุ ทุกฺขา ติพฺพา กฏุกา เวทนา เวทิเยยฺย.\csromanspacer{} ตสฺส โส ภิสกฺโก สลฺลกตฺโต เอสนิยา สลฺลํ เอเสยฺย, โส เอสนิยาปิ สลฺลสฺส เอสนาเหตุ ทุกฺขา ติพฺพา กฏุกา เวทนา เวทิเยยฺย.\csromanspacer{} ตสฺส โส ภิสกฺโก สลฺลกตฺโต สลฺลํ อพฺพุเหยฺย\footnote{อพฺพุเยฺหยฺย (สี), อพฺภูเณฺหยฺย (สฺยา, กํ)}, โส สลฺลสฺสปิ อพฺพุหนเหตุ\footnote{อพฺพุยฺหนเหตุ (สี), อพฺภูณฺหนเหตุ (สฺยา, กํ)} ทุกฺขา ติพฺพา กฏุกา เวทนา เวทิเยยฺย.\csromanspacer{} ตสฺส โส ภิสกฺโก สลฺลกตฺโต อคทงฺคารํ วณมุเข โอทเหยฺย, โส อคทงฺคารสฺสปิ วณมุเข โอทหนเหตุ ทุกฺขา ติพฺพา กฏุกา เวทนา เวทิเยยฺย.\csromanspacer{} โส อปเรน สมเยน รูเฬฺหน วเณน สญฺฉวินา อโรโค อสฺส สุขี เสรี สยํวสี เยน กามงฺคโม.\csromanspacer{} ตสฺส เอวมสฺส “อหํ โข ปุพฺเพ สลฺเลน วิทฺโธ อโหสิํ สวิเสน คาฬฺหูปเลปเนน, โสหํ สลฺลสฺสปิ เวธนเหตุ ทุกฺขา ติพฺพา กฏุกา เวทนา เวทิยิํ.\csromanspacer{} ตสฺส เม มิตฺตามจฺจา ญาติสาโลหิตา ภิสกฺกํ สลฺลกตฺตํ อุปฏฺฐเปสุํ.\csromanspacer{} ตสฺส เม โส ภิสกฺโก สลฺลกตฺโต สตฺเถน วณมุขํ ปริกนฺติ, โสหํ สตฺเถนปิ วณมุขสฺส ปริกนฺตนเหตุ ทุกฺขา ติพฺพา กฏุกา เวทนา เวทิยิํ.\csromanspacer{} ตสฺส เม โส ภิสกฺโก สลฺลกตฺโต เอสนิยา สลฺลํ เอสิ, โส อหํ เอสนิยาปิ สลฺลสฺส เอสนาเหตุ ทุกฺขา ติพฺพา กฏุกา เวทนา เวทิยิํ.\csromanspacer{} ตสฺส เม โส \csromanfolio{4}ภิสกฺโก สลฺลกตฺโต สลฺลํ อพฺพุหิ\footnote{อพฺพุยฺหิ (สี), อพฺภูณฺหิ (สฺยา, กํ)}, โสหํ สลฺลสฺสปิ อพฺพุหนเหตุ ทุกฺขา ติพฺพา กฏุกา เวทนา เวทิยิํ.\csromanspacer{} ตสฺส เม โส ภิสกฺโก สลฺลกตฺโต อคทงฺคารํ วณมุเข โอทหิ, โสหํ อคทงฺคารสฺสปิ วณมุเข โอทหนเหตุ ทุกฺขา ติพฺพา กฏุกา เวทนา เวทิยิํ.\csromanspacer{} โสมฺหิ เอตรหิ รูเฬฺหน วเณน สญฺฉวินา อโรโค สุขี เสรี สยํวสี เยน กามงฺคโม”ติ.}

\prose{เอวเมว โข อาวุโส นิคณฺฐา สเจ ตุเมฺห ชาเนยฺยาถ “อหุวเมฺหว มยํ ปุพฺเพ, น นาหุวมฺหา”ติ.\csromanspacer{} ชาเนยฺยาถ “อกรเมฺหว มยํ ปุพฺเพ ปาปกมฺมํ, น นากรมฺหา”ติ.\csromanspacer{} ชาเนยฺยาถ “เอวรูปํ วา เอวรูปํ วา ปาปกมฺมํ อกรมฺหา”ติ.\csromanspacer{} ชาเนยฺยาถ “เอตฺตกํ วา ทุกฺขํ นิชฺชิณฺณํ, เอตฺตกํ วา ทุกฺขํ นิชฺชีเรตพฺพํ, เอตฺตกมฺหิ วา ทุกฺเข นิชฺชิณฺเณ สพฺพํ ทุกฺขํ นิชฺชิณฺณํ ภวิสฺสตี”ติ.\csromanspacer{} ชาเนยฺยาถ\csromandash{}ทิฏฺเฐว ธมฺเม อกุสลานํ ธมฺมานํ ปหานํ กุสลานํ ธมฺมานํ อุปสมฺปทํ.\csromanspacer{} เอวํ สนฺเต อายสฺมนฺตานํ นิคณฺฐานํ กลฺลมสฺส เวยฺยากรณาย “ยํ กิญฺจายํ ปุริสปุคฺคโล ปฏิสํเวเทติ สุขํ วา ทุกฺขํ วา อทุกฺขมสุขํ วา, สพฺพํ ตํ ปุพฺเพกตเหตุ.\csromanspacer{} อิติ ปุราณานํ กมฺมานํ ตปสา พฺยนฺตีภาวา นวานํ กมฺมานํ อกรณา อายติํ อนวสฺสโว, อายติํ อนวสฺสวา กมฺมกฺขโย, กมฺมกฺขยา ทุกฺขกฺขโย, ทุกฺขกฺขยา เวทนากฺขโย, เวทนากฺขยา สพฺพํ ทุกฺขํ นิชฺชิณฺณํ ภวิสฺสตี”ติ.}

\prose{ยสฺมา จ โข ตุเมฺห อาวุโส นิคณฺฐา น ชานาถ “อหุวเมฺหว มยํ ปุพฺเพ, น นาหุวมฺหา”ติ.\csromanspacer{} น ชานาถ “อกรเมฺหว มยํ ปุพฺเพ ปาปกมฺมํ, น นากรมฺหา”ติ.\csromanspacer{} น ชานาถ “เอวรูปํ วา เอวรูปํ วา ปาปกมฺมํ อกรมฺหา”ติ.\csromanspacer{} น ชานาถ “เอตฺตกํ วา ทุกฺขํ นิชฺชิณฺณํ, เอตฺตกํ วา ทุกฺขํ นิชฺชีเรตพฺพํ, เอตฺตกมฺหิ วา ทุกฺเข นิชฺชิณฺเณ สพฺพํ ทุกฺขํ นิชฺชิณฺณํ ภวิสฺสตี”ติ.\csromanspacer{} น ชานาถ\csromandash{}ทิฏฺเฐว ธมฺเม อกุสลานํ ธมฺมานํ ปหานํ กุสลานํ ธมฺมานํ อุปสมฺปทํ.\csromanspacer{} ตสฺมา อายสฺมนฺตานํ นิคณฺฐานํ น กลฺลมสฺส เวยฺยากรณาย “ยํ กิญฺจายํ ปุริปุคฺคโล ปฏิสํเวเทติ สุขํ วา ทุกฺขํ วา อทุกฺขมสุขํ วา, สพฺพํ ตํ ปุพฺเพกตเหตุ.\csromanspacer{} อิติ ปุราณานํ กมฺมานํ ตปสา พฺยนฺตีภาวา นวานํ กมฺมานํ อกรณา อายติํ อนวสฺสโว, อายติํ อนวสฺสวา กมฺมกฺขโย, กมฺมกฺขยา ทุกฺขกฺขโย, ทุกฺขกฺขยา เวทนากฺขโย, เวทนากฺขยา สพฺพํ ทุกฺขํ นิชฺชิณฺณํ ภวิสฺสตี”ติ.}

\proseitem{4}{\csromanfolio{5}เอวํ วุตฺเต ภิกฺขเว เต นิคณฺฐา มํ เอตทโวจุํ “นิคณฺโฐ อาวุโส นาฏปุตฺโต\footnote{นาถปุตฺโต (สี)} สพฺพญฺญู สพฺพทสฺสาวี อปริเสสํ ญาณทสฺสนํ ปฏิชานาติ ‘จรโต จ เม ติฏฺฐโต จ สุตฺตสฺส จ ชาครสฺส จ สตตํ สมิตํ ญาณทสฺสนํ ปจฺจุปฏฺฐิตนฺ’ติ.\csromanspacer{} โส เอวมาห ‘อตฺถิ โข โว อาวุโส นิคณฺฐา ปุพฺเพว ปาปกมฺมํ กตํ, ตํ อิมาย กฏุกาย ทุกฺกรการิกาย นิชฺชีเรถ.\csromanspacer{} ยํ ปเนตฺถ เอตรหิ กาเยน สํวุตา วาจาย สํวุตา มนสา สํวุตา, ตํ อายติํ ปาปกมฺมสฺส อกรณํ.\csromanspacer{} อิติ ปุราณานํ กมฺมานํ ตปสา พฺยนฺตีภาวา นวานํ กมฺมานํ อกรณา อายติํ อนวสฺสโว, อายติํ อนวสฺสวา กมฺมกฺขโย, กมฺมกฺขยา ทุกฺขกฺขโย, ทุกฺขกฺขยา เวทนากฺขโย, เวทนากฺขยา สพฺพํ ทุกฺขํ นิชฺชิณฺณํ ภวิสฺสตีติ, ตญฺจ ปนมฺหากํ รุจฺจติ เจว ขมติ จ, เตน จมฺหา อตฺตมนา”ติ.}

\proseitem{5}{เอวํ วุตฺเต อหํ ภิกฺขเว เต นิคณฺเฐ เอตทโวจํ “ปญฺจ โข อิเม อาวุโส นิคณฺฐา ธมฺมา ทิฏฺเฐว ธมฺเม ทฺวิธาวิปากา.\csromanspacer{} กตเม ปญฺจ, สทฺธา รุจิ อนุสฺสโว อาการปริวิตกฺโก ทิฏฺฐินิชฺฌานกฺขนฺติ.\csromanspacer{} อิเม โข อาวุโส นิคณฺฐา ปญฺจ ธมฺมา ทิฏฺเฐว ธมฺเม ทฺวิธาวิปากา.\csromanspacer{} ตตฺรายสฺมนฺตานํ นิคณฺฐานํ กา อตีตํเส สตฺถริ สทฺธา, กา รุจิ, โก อนุสฺสโว, โก อาการปริวิตกฺโก, กา ทิฏฺฐินิชฺฌานกฺขนฺตี”ติ.\csromanspacer{} เอวํวาที\footnote{เอวํวาทีสุ (ก)} โข อหํ ภิกฺขเว นิคณฺเฐสุ น กญฺจิ\footnote{กิญฺจิ (สี, อิ, ก)} สหธมฺมิกํ วาทปฏิหารํ สมนุปสฺสามิ.}

\prose{ปุน จปราหํ\footnote{ปุน จ ปนาหํ (สี, อิ, ก)} ภิกฺขเว เต นิคณฺเฐ เอวํ วทามิ “ตํ กิํ มญฺญถ อาวุโส นิคณฺฐา, ยสฺมิํ โว สมเย ติพฺโพ\footnote{ติปฺโป (อิ)} อุปกฺกโม โหติ ติพฺพํ ปธานํ, ติพฺพา ตสฺมิํ สมเย โอปกฺกมิกา ทุกฺขา ติพฺพา กฏุกา เวทนา เวทิเยถ.\csromanspacer{} ยสฺมิํ ปน โว สมเย น ติพฺโพ อุปกฺกโม โหติ น ติพฺพํ ปธานํ, น ติพฺพา ตสฺมิํ สมเย โอปกฺกมิกา ทุกฺขา ติพฺพา กฏุกา เวทนา เวทิเยถา”ติ? ยสฺมิํ โน อาวุโส โคตม สมเย ติพฺโพ อุปกฺกโม โหติ ติพฺพํ ปธานํ, ติพฺพา ตสฺมิํ สมเย โอปกฺกมิกา ทุกฺขา ติพฺพา กฏุกา เวทนา เวทิยาม.\csromanspacer{} ยสฺมิํ ปน โน สมเย น ติพฺโพ อุปกฺกโม โหติ น ติพฺพํ ปธานํ, น ติพฺพา ตสฺมิํ สมเย โอปกฺกมิกา ทุกฺขา ติพฺพา กฏุกา เวทนา เวทิยามาติ.}

\proseitem{6}{\csromanfolio{6}อิติ กิร อาวุโส นิคณฺฐา ยสฺมิํ โว สมเย ติพฺโพ อุปกฺกโม โหติ ติพฺพํ ปธานํ, ติพฺพา ตสฺมิํ สมเย โอปกฺกมิกา ทุกฺขา ติพฺพา กฏุกา เวทนา เวทิเยถ.\csromanspacer{} ยสฺมิํ ปน โว สมเย น ติพฺโพ อุปกฺกโม โหติ น ติพฺพํ ปธานํ, น ติพฺพา ตสฺมิํ สมเย โอปกฺกมิกา ทุกฺขา ติพฺพา กฏุกา เวทนา เวทิเยถ.\csromanspacer{} เอวํ สนฺเต อายสฺมนฺตานํ นิคณฺฐานํ น กลฺลมสฺส เวยฺยากรณาย “ยํ กิญฺจายํ ปุริสปุคฺคโล ปฏิสํเวเทติ สุขํ วา ทุกฺขํ วา อทุกฺขมสุขํ วา, สพฺพํ ตํ ปุพฺเพกตเหตุ.\csromanspacer{} อิติ ปุราณานํ กมฺมานํ ตปสา พฺยนฺตีภาวา นวานํ กมฺมานํ อกรณา อายติํ อนวสฺสโว, อายติํ อนวสฺสวา กมฺมกฺขโย, กมฺมกฺขยา ทุกฺขกฺขโย, ทุกฺขกฺขยา เวทนากฺขโย, เวทนากฺขยา สพฺพํ ทุกฺขํ นิชฺชิณฺณํ ภวิสฺสตี”ติ.\csromanspacer{} สเจ อาวุโส นิคณฺฐา ยสฺมิํ โว สมเย ติพฺโพ อุปกฺกโม โหติ ติพฺพํ ปธานํ, น ติพฺพา ตสฺมิํ สมเย โอปกฺกมิกา ทุกฺขา ติพฺพา กฏุกา เวทนา เวทิเยถ\footnote{ปธานํ, ติฏฺเฐยฺเยว ตสฺมิํ สมเย ... เวทนา (สี, สฺยา, กํ, อิ)}.\csromanspacer{} ยสฺมิํ ปน โว สมเย น ติพฺโพ อุปกฺกโม โหติ น ติพฺพํ ปธานํ, ติพฺพา ตสฺมิํ สมเย โอปกฺกมิกา ทุกฺขา ติพฺพา กฏุกา เวทนา เวทิเยถ1.\csromanspacer{} เอวํ สนฺเต อายสฺมนฺตานํ นิคณฺฐานํ กลฺลมสฺส เวยฺยากรณาย “ยํ กิญฺจายํ ปุริสปุคฺคโล ปฏิสํเวเทติ สุขํ วา ทุกฺขํ วา อทุกฺขมสุขํ วา, สพฺพํ ตํ ปุพฺเพกตเหตุ.\csromanspacer{} อิติ ปุราณานํ กมฺมานํ ตปสา พฺยนฺตีภาวา นวานํ กมฺมานํ อกรณา อายติํ อนวสฺสโว, อายติํ อนวสฺสวา กมฺมกฺขโย, กมฺมกฺขยา ทุกฺขกฺขโย, ทุกฺขกฺขยา เวทนากฺขโย, เวทนากฺขยา สพฺพํ ทุกฺขํ นิชฺชิณฺณํ ภวิสฺสตี”ติ.}

\prose{ยสฺมา จ โข อาวุโส นิคณฺฐา ยสฺมิํ โว สมเย ติพฺโพ อุปกฺกโม โหติ ติพฺพํ ปธานํ, ติพฺพา ตสฺมิํ สมเย โอปกฺกมิกา ทุกฺขา ติพฺพา กฏุกา เวทนา เวทิเยถ.\csromanspacer{} ยสฺมิํ ปน โว สมเย น ติพฺโพ อุปกฺกโม โหติ น ติพฺพํ ปธานํ, น ติพฺพา ตสฺมิํ สมเย โอปกฺกมิกา ทุกฺขา ติพฺพา กฏุกา เวทนา เวทิเยถ.\csromanspacer{} เต ตุเมฺห สามํเยว โอปกฺกมิกา ทุกฺขา ติพฺพา กฏุกา เวทนา เวทยมานา อวิชฺชา อญฺญาณา สมฺโมหา วิปจฺเจถ “ยํ กิญฺจายํ ปุริสปุคฺคโล ปฏิเวเทติ สุขํ วา ทุกฺขํ วา อทุกฺขมสุขํ วา, สพฺพํ ตํ ปุพฺเพกตเหตุ.\csromanspacer{} อิติ ปุราณานํ กมฺมานํ ตปสา พฺยนฺตีภาวา นวานํ กมฺมานํ อกรณา อายติํ อนวสฺสโว, อายติํ อนวสฺสวา กมฺมกฺขโย, กมฺมกฺขยา ทุกฺขกฺขโย, ทุกฺขกฺขยา \csromanfolio{7}เวทนากฺขโย, เวทนากฺขยา สพฺพํ ทุกฺขํ นิชฺชิณฺณํ ภวิสฺสตี”ติ.\csromanspacer{} เอวํวาทีปิ\footnote{เอวํวาทีสุปิ (ก)} โข อหํ ภิกฺขเว นิคณฺเฐสุ น กญฺจิ สหธมฺมิกํ วาทปฏิหารํ สมนุปสฺสามิ.}

\proseitem{7}{ปุน จปราหํ ภิกฺขเว เต นิคณฺเฐ เอวํ วทามิ “ตํ กิํ มญฺญถาวุโส นิคณฺฐา, ยมิทํ กมฺมํ ทิฏฺฐธมฺมเวทนียํ, ตํ อุปกฺกเมน วา ปธาเนน วา ‘สมฺปรายเวทนียํ โหตู’ติ ลพฺภเมตนฺ”ติ.\csromanspacer{} โน หิทํ อาวุโส. “ยํ ปนิทํ กมฺมํ สมฺปรายเวทนียํ, ตํ อุปกฺกเมน วา ปธาเนน วา ‘ทิฏฺฐธมฺมเวทนียํ โหตู’ติ ลพฺภเมตนฺ”ติ.\csromanspacer{} โน หิทํ อาวุโส. “ตํ กิํ มญฺญถาวุโส นิคณฺฐา, ยมิทํ กมฺมํ สุขเวทนียํ, ตํ อุปกฺกเมน วา ปธาเนน วา ‘ทุกฺขเวทนียํ โหตู’ติ ลพฺภเมตนฺ”ติ.\csromanspacer{} โน หิทํ อาวุโส. “ยํ ปนิทํ กมฺมํ ทุกฺขเวทนียํ, ตํ อุปกฺกเมน วา ปธาเนน วา ‘สุขเวทนียํ โหตู’ติ ลพฺภเมตนฺ”ติ.\csromanspacer{} โน หิทํ อาวุโส. “ตํ กิํ มญฺญถาวุโส นิคณฺฐา, ยมิทํ กมฺมํ ปริปกฺกเวทนียํ, ตํ อุปกฺกเมน วา ปธาเนน วา ‘อปริปกฺกเวทนียํ โหตู’ติ ลพฺภเมตนฺ”ติ.\csromanspacer{} โน หิทํ อาวุโส. “ยํ ปนิทํ กมฺมํ อปริปกฺกเวทนียํ, ตํ อุปกฺกเมน วา ปธาเนน วา ‘ปริปกฺกเวทนียํ โหตู’ติ ลพฺภเมตนฺ”ติ.\csromanspacer{} โน หิทํ อาวุโส. “ตํ กิํ มญฺญถาวุโส นิคณฺฐา, ยมิทํ กมฺมํ พหุเวทนียํ, ตํ อุปกฺกเมน วา ปธาเนน วา ‘อปฺปเวทนียํ โหตู’ติ ลพฺภเมตนฺ”ติ.\csromanspacer{} โน หิทํ อาวุโส. “ยํ ปนิทํ กมฺมํ อปฺปเวทนียํ, ตํ อุปกฺกเมน วา ปธาเนน วา ‘พหุเวทนียํ โหตู’ติ ลพฺภเมตนฺ”ติ.\csromanspacer{} โน หิทํ อาวุโส. “ตํ กิํ มญฺญถาวุโส นิคณฺฐา, ยมิทํ กมฺมํ สเวทนียํ, ตํ อุปกฺกเมน วา ปธาเนน วา ‘อเวทนียํ โหตู’ติ ลพฺภเมตนฺ”ติ.\csromanspacer{} โน หิทํ อาวุโส. “ยํ ปนิทํ กมฺมํ อเวทนียํ, ตํ อุปกฺกเมน วา ปธาเนน วา ‘สเวทนียํ โหตู’ติ ลพฺภเมตนฺ”ติ.\csromanspacer{} โน หิทํ อาวุโส.}

\proseitem{8}{อิติ กิร อาวุโส นิคณฺฐา ยมิทํ กมฺมํ ทิฏฺฐธมฺมเวทนียํ, ตํ อุปกฺกเมน วา ปธาเนน วา “สมฺปรายเวทนียํ โหตู”ติ อลพฺภเมตํ.\csromanspacer{} ยํ ปนิทํ กมฺมํ สมฺปรายเวทนียํ, ตํ อุปกฺกเมน วา ปธาเนน วา “ทิฏฺฐธมฺมเวทนียํ โหตู”ติ อลพฺภเมตํ.\csromanspacer{} ยมิทํ กมฺมํ สุขเวทนียํ, ตํ อุปกฺกเมน วา ปธาเนน วา “ทุกฺขเวทนียํ โหตู”ติ อลพฺภเมตํ.\csromanspacer{} ยมิทํ กมฺมํ \csromanfolio{8}ทุกฺขเวทนียํ, ตํ อุปกฺกเมน วา ปธาเนน วา “สุขเวทนียํ โหตู”ติ อลพฺภเมตํ.\csromanspacer{} ยมิทํ กมฺมํ ปริปกฺกเวทนียํ, ตํ อุปกฺกเมน วา ปธาเนน วา “อปริปกฺกเวทนียํ โหตู”ติ อลพฺภเมตํ.\csromanspacer{} ยมิทํ กมฺมํ อปริปกฺกเวทนียํ, ตํ อุปกฺกเมน วา ปธาเนน วา “ปริปกฺกเวทนียํ โหตู”ติ อลพฺภเมตํ.\csromanspacer{} ยมิทํ กมฺมํ พหุเวทนียํ, ตํ อุปกฺกเมน วา ปธาเนน วา “อปฺปเวทนียํ โหตู”ติ อลพฺภเมตํ.\csromanspacer{} ยมิทํ กมฺมํ อปฺปเวทนียํ, ตํ อุปกฺกเมน วา ปธาเนน วา “พหุเวทนียํ โหตู”ติ อลพฺภเมตํ.\csromanspacer{} ยมิทํ กมฺมํ สเวทนียํ, ตํ อุปกฺกเมน วา ปธาเนน วา “อเวทนียํ โหตู”ติ อลพฺภเมตํ.\csromanspacer{} ยมิทํ กมฺมํ อเวทนียํ, ตํ อุปกฺกเมน วา ปธาเนน วา “สเวทนียํ โหตู”ติ อลพฺภเมตํ.\csromanspacer{} เอวํ สนฺเต อายสฺมนฺตานํ นิคณฺฐานํ อผโล อุปกฺกโม โหติ อผลํ ปธานํ.\csromanspacer{} เอวํวาที ภิกฺขเว นิคณฺฐา, เอวํวาทีนํ ภิกฺขเว นิคณฺฐานํ ทส สหธมฺมิกา วาทานุวาทา คารยฺหํ ฐานํ อาคจฺฉนฺติ.}

\proseitem{9}{สเจ ภิกฺขเว สตฺตา ปุพฺเพกตเหตุ สุขทุกฺขํ ปฏิสํเวเทนฺติ, อทฺธา ภิกฺขเว นิคณฺฐา ปุพฺเพ ทุกฺกฏกมฺมการิโน, ยํ เอตรหิ เอวรูปา ทุกฺขา ติพฺพา กฏุกา เวทนา เวทิยนฺติ.\csromanspacer{} สเจ ภิกฺขเว สตฺตา อิสฺสรนิมฺมานเหตุ สุขทุกฺขํ ปฏิสํเวเทนฺติ, อทฺธา ภิกฺขเว นิคณฺฐา ปาปเกน อิสฺสเรน นิมฺมิตา, ยํ เอตรหิ เอวรูปา ทุกฺขา ติพฺพา กฏุกา เวทนา เวทิยนฺติ.\csromanspacer{} สเจ ภิกฺขเว สตฺตา สงฺคติภาวเหตุ สุขทุกฺขํ ปฏิสํเวเทนฺติ, อทฺธา ภิกฺขเว นิคณฺฐา ปาปสงฺคติกา, ยํ เอตรหิ เอวรูปา ทุกฺขา ติพฺพา กฏุกา เวทนา เวทิยนฺติ.\csromanspacer{} สเจ ภิกฺขเว สตฺต อภิชาติเหตุ สุขทุกฺขํ ปฏิสํเวเทนฺติ, อทฺธา ภิกฺขเว นิคณฺฐา ปาปาภิชาติกา, ยํ เอตรหิ เอวรูปา ทุกฺขา ติพฺพา กฏุกา เวทนา เวทิยนฺติ.\csromanspacer{} สเจ ภิกฺขเว สตฺตา ทิฏฺฐธมฺมูปกฺกมเหตุ สุขทุกฺขํ ปฏิสํเวเทนฺติ, อทฺธา ภิกฺขเว นิคณฺฐา เอวรูปา ทิฏฺฐธมฺมูปกฺกมา, ยํ เอตรหิ เอวรูปา ทุกฺขา ติพฺพา กฏุกา เวทนา เวทิยนฺติ.}

\prose{สเจ ภิกฺขเว สตฺตา ปุพฺเพกตเหตุ สุขทุกฺขํ ปฏิสํเวเทนฺติ, คารยฺหา นิคณฺฐา.\csromanspacer{} โน เจ สตฺตา ปุพฺเพกตเหตุ สุขทุกฺขํ ปฏิสํเวเทนฺติ, คารยฺหา นิคณฺฐา.\csromanspacer{} สเจ ภิกฺขเว สตฺตา อิสฺสรนิมฺมานเหตุ สุขทุกฺขํ ปฏิสํเวเทนฺติ, คารยฺหา นิคณฺฐา.\csromanspacer{} โน เจ สตฺตา อิสฺสรนิมฺมานเหตุ สุขทุกฺขํ ปฏิสํเวเทนฺติ, คารยฺหา นิคณฺฐา.\csromanspacer{} สเจ ภิกฺขเว สตฺตา \csromanfolio{9}สงฺคติภาวเหตุ สุขทุกฺขํ ปฏิสํเวเทนฺติ, คารยฺหา นิคณฺฐา.\csromanspacer{} โน เจ สตฺตา สงฺคติภาวเหตุ สุขทุกฺขํ ปฏิสํเวเทนฺติ, คารยฺหา นิคณฺฐา.\csromanspacer{} สเจ ภิกฺขเว สตฺตา อภิชาติเหตุ สุขทุกฺขํ ปฏิสํเวเทนฺติ, คารยฺหา นิคณฺฐา.\csromanspacer{} โน เจ สตฺตา อภิชาติเหตุ สุขทุกฺขํ ปฏิสํเวเทนฺติ, คารยฺหา นิคณฺฐา.\csromanspacer{} สเจ ภิกฺขเว สตฺตา ทิฏฺฐธมฺมูปกฺกมเหตุ สุขทุกฺขํ ปฏิสํเวเทนฺติ, คารยฺหา นิคณฺฐา.\csromanspacer{} โน เจ สตฺตา ทิฏฺฐธมฺมูปกฺกมเหตุ สุขทุกฺขํ ปฏิสํเวเทนฺติ, คารยฺหา นิคณฺฐา.\csromanspacer{} เอวํวาที ภิกฺขเว นิคณฺฐา, เอวํวาทีนํ ภิกฺขเว นิคณฺฐานํ อิเม ทส สหธมฺมิกา วาทานุวาทา คารยฺหํ ฐานํ อาคจฺฉนฺติ.\csromanspacer{} เอวํ โข ภิกฺขเว อผโล อุปกฺกโม โหติ อผลํ ปธานํ.}

\proseitem{10}{กถญฺจ ภิกฺขเว สผโล อุปกฺกโม โหติ สผลํ ปธานํ.\csromanspacer{} อิธ ภิกฺขเว ภิกฺขุ น เหว อนทฺธภูตํ อตฺตานํ ทุกฺเขน อทฺธภาเวติ, ธมฺมิกญฺจ สุขํ น ปริจฺจชติ, ตสฺมิํ จ สุเข อนธิมุจฺฉิโต โหติ.\csromanspacer{} โส เอวํ ปชานาติ “อิมสฺส โข เม ทุกฺขนิทานสฺส สงฺขารํ ปทหโต สงฺขารปฺปธานา วิราโค โหติ, อิมสฺส ปน เม ทุกฺขนิทานสฺส อชฺฌุเปกฺขโต อุเปกฺขํ ภาวยโต วิราโค โหตี”ติ.\csromanspacer{} โส ยสฺส หิ ขฺวาสฺส\footnote{ยสฺส โข ปนสฺส (สี), ยสฺส ขฺวาสฺส (อิ)} ทุกฺขนิทานสฺส สงฺขารํ ปทหโต สงฺขารปฺปธานา วิราโค โหติ, สงฺขารํ ตตฺถ ปทหติ.\csromanspacer{} ยสฺส ปนสฺส ทุกฺขนิทานสฺส อชฺฌุเปกฺขโต อุเปกฺขํ ภาวยโต วิราโค โหติ, อุเปกฺขํ ตตฺถ ภาเวติ.\csromanspacer{} ตสฺส ตสฺส ทุกฺขนิทานสฺส สงฺขารํ ปทหโต สงฺขารปฺปธานา วิราโค โหติ.\csromanspacer{} เอวมฺปิสฺส ตํ ทุกฺขํ นิชฺชิณฺณํ โหติ.\csromanspacer{} ตสฺส ตสฺส ทุกฺขนิทานสฺส อชฺฌุเปกฺขโต อุเปกฺขํ ภาวยโต วิราโค โหติ.\csromanspacer{} เอวมฺปิสฺส ตํ ทุกฺขํ นิชฺชิณฺณํ โหติ.}

\proseitem{11}{เสยฺยถาปิ ภิกฺขเว ปุริโส อิตฺถิยา สารตฺโต ปฏิพทฺธจิตฺโต ติพฺพจฺฉนฺโท ติพฺพาเปกฺโข, โส ตํ อิตฺถิํ ปสฺเสยฺย อญฺเญน ปุริเสน สทฺธิํ สนฺติฏฺฐนฺติํ สลฺลปนฺติํ สญฺชคฺฆนฺติํ สํหสนฺติํ.\csromanspacer{} ตํ กิํ มญฺญถ ภิกฺขเว, อปิ นุ ตสฺสปุริสสฺส อมุํ อิตฺถิํ ทิสฺวา อญฺเญน ปุริเสน สทฺธิํ สนฺติฏฺฐนฺติํ สลฺลปนฺติํ สญฺชคฺฆนฺติํ สํหสนฺติํ อุปฺปชฺเชยฺยุํ โสกปริเทวทุกฺขโทมนสฺสูปายาสาติ.\csromanspacer{} เอวํ ภนฺเต.\csromanspacer{} ตํ กิสฺส เหตุ, อมุ หิ ภนฺเต ปุริโส อมุสฺสา อิตฺถิยา สารตฺโต ปฏิพทฺธจิตฺโต ติพฺพจฺฉนฺโท ติพฺพาเปกฺโข, ตสฺมา ตํ อิตฺถิํ ทิสฺวา อญฺเญน ปุริเสน สทฺธิํ สนฺติฏฺฐนฺติํ สลฺลปนฺติํ \csromanfolio{10}สญฺชคฺฆนฺติํ สํหสนฺติํ อุปฺปชฺเชยฺยุํ โสกปริเทวทุกฺขโทมนสฺสูปายาสาติ.\csromanspacer{} อถ โข ภิกฺขเว ตสฺส ปุริสสฺส เอวมสฺส “อหํ โข อมุสฺสา อิตฺถิยา สารตฺโต ปฏิพทฺธจิตฺโต ติพฺพจฺฉนฺโท ติพฺพาเปกฺโข, ตสฺส เม อมุํ อิตฺถิํ ทิสฺวา อญฺเญน ปุริเสน สทฺธิํ สนฺติฏฺฐนฺติํ สลฺลปนฺติํ สญฺชคฺฆนฺติํ สํหสนฺติํ อุปฺปชฺชนฺติ โสกปริเทวทุกฺขโทมนสฺสูปายาสา.\csromanspacer{} ยํนูนาหํ โย เม อมุสฺสา อิตฺถิยา ฉนฺทราโค, ตํ ปชเหยฺยนฺ”ติ.\csromanspacer{} โส โย อมุสฺสา อิตฺถิยา ฉนฺทราโค, ตํ ปชเหยฺย.\csromanspacer{} โส ตํ อิตฺถิํ ปสฺเสยฺย อปเรน สมเยน อญฺเญน ปุริเสน สทฺธิํ สนฺติฏฺฐนฺติํ สลฺลปนฺติํ สญฺชคฺฆนฺติํ สํหสนฺติํ.\csromanspacer{} ตํ กิํ มญฺญถ ภิกฺขเว, อปิ นุ ตสฺส ปุริสสฺส อมุํ อิตฺถิํ ทิสฺวา อญฺเญน ปุริเสน สทฺธิํ สนฺติฏฺฐนฺติํ สลฺลปนฺติํ สญฺชคฺฆนฺติํ สํหสนฺติํ อุปฺปชฺเชยฺยุํ โสกปริเทวทุกฺขโทมนสฺสูปายาสาติ.\csromanspacer{} โน เหตํ ภนฺเต.\csromanspacer{} ตํ กิสฺส เหตุ, อมุ หิ ภนฺเต ปุริโส อมุสฺสา อิตฺถิยา วิราโค, ตสฺมา ตํ อิตฺถิํ ทิสฺวา อญฺเญน ปุริเสน สทฺธิํ สนฺติฏฺฐนฺติํ สลฺลปนฺติํ สญฺชคฺฆนฺติํ สํหสนฺติํ น อุปชฺเชยฺยุํ โสกปริเทวทุกฺขโทมนสฺสูปายาสาติ.}

\prose{เอวเมว โข ภิกฺขเว ภิกฺขุ น เหว อนทฺธภูตํ อตฺตานํ ทุกฺเขน อทฺธภาเวติ, ธมฺมิกญฺจ สุขํ น ปริจฺจชติ, ตสฺมิํ จ สุเข อนธิมุจฺฉิโต โหติ.\csromanspacer{} โส เอวํ ปชานาติ\csromandash{}“อิมสฺส โข เม ทุกฺขนิทานสฺส สงฺขารํ ปทหโต สงฺขารปฺปธานา วิราโค โหติ.\csromanspacer{} อิมสฺส ปน เม ทุกฺขนิทานสฺส อชฺฌุเปกฺขโต อุเปกฺขํ ภาวยโต วิราโค โหตี”ติ.\csromanspacer{} โส ยสฺส หิ ขฺวาสฺส ทุกฺขนิทานสฺส สงฺขารํ ปทหโต สงฺขารปฺปธานา วิราโค โหติ, สงฺขารํ ตตฺถ ปทหติ.\csromanspacer{} ยสฺส ปนสฺส ทุกฺขนิทานสฺส อชฺฌุเปกฺขโต อุเปกฺขํ ภาวยโต วิราโค โหติ, อุเปกฺขํ ตตฺถ ภาเวติ.\csromanspacer{} ตสฺส ตสฺส ทุกฺขนิทานสฺส สงฺขารํ ปทหโต สงฺขารปฺปธานา วิราโค โหติ.\csromanspacer{} เอวมฺปิสฺส ตํ ทุกฺขํ นิชฺชิณฺณํ โหติ.\csromanspacer{} ตสฺส ตสฺส ทุกฺขนิทานสฺส อชฺฌุเปกฺขโต อุเปกฺขํ ภาวยโต วิราโค โหติ.\csromanspacer{} เอวมฺปิสฺส ตํ ทุกฺขํ นิชฺชิณฺณํ โหติ.\csromanspacer{} เอวมฺปิ ภิกฺขเว สผโล อุปกฺกโม โหติ สผลํ ปธานํ.}

\proseitem{12}{ปุน จปรํ ภิกฺขเว ภิกฺขุ อิติ ปฏิสญฺจิกฺขติ “ยถาสุขํ โข เม วิหรโต อกุสลา ธมฺมา อภิวฑฺฒนฺติ, กุสลา ธมฺมา ปริหายนฺติ.\csromanspacer{} ทุกฺขาย ปน เม อตฺตานํ ปทหโต อกุสลา ธมฺมา ปริหายนฺติ, กุสลา ธมฺมา อภิวฑฺฒนฺติ.\csromanspacer{} ยํนูนาหํ ทุกฺขาย อตฺตานํ ปทเหยฺยนฺ”ติ.\csromanspacer{} โส ทุกฺขาย อตฺตานํ ปทหติ, ตสฺส ทุกฺขาย อตฺตานํ ปทหโต \csromanfolio{11}อกุสลา ธมฺมา ปริหายนฺติ, กุสลา ธมฺมา อภิวฑฺฒนฺติ.\csromanspacer{} โส น อปเรน สมเยน ทุกฺขาย อตฺตานํ ปทหติ.\csromanspacer{} ตํ กิสฺส เหตุ, ยสฺส หิ โส ภิกฺขเว ภิกฺขุ อตฺถาย ทุกฺขาย อตฺตานํ ปทเหยฺย, สฺวาสฺส อตฺโถ อภินิปฺผนฺโน โหติ.\csromanspacer{} ตสฺมา น อปเรน สมเยน ทุกฺขาย อตฺตานํ ปทหติ.\csromanspacer{} เสยฺยถาปิ ภิกฺขเว อุสุกาโร เตชนํ ทฺวีสุ อลาเตสุ อาตาเปติ ปริตาเปติ อุชุํ กโรติ กมฺมนิยํ.\csromanspacer{} ยโต โข ภิกฺขเว อุสุการสฺส เตชนํ ทฺวีสุ อลาเตสุ อาตาปิตํ โหติ ปริตาปิตํ อุชุํ กตํ\footnote{อุชุํ กตํ โหติ (สี)} กมฺมนิยํ, น โส ตํ อปเรน สมเยน อุสุกาโร เตชนํ ทฺวีสุ อลาเตสุ อาตาเปติ ปริตาเปติ อุชุํ กโรติ กมฺมนิยํ.\csromanspacer{} ตํ กิสฺส เหตุ, ยสฺส หิ โส ภิกฺขเว อตฺถาย อุสุกาโร เตชนํ ทฺวีสุ อลาเตสุ อาตาเปยฺย ปริตาเปยฺย อุชุํ กเรยฺย กมฺมนิยํ, สฺวาสฺส อตฺโถ อภินิปฺผนฺโน โหติ.\csromanspacer{} ตสฺมา น อปเรน สมเยน อุสุกาโร เตชนํ ทฺวีสุ อลาเตสุ อาตาเปติ ปริตาเปติ อุชุํ กโรติ กมฺมนิยํ.\csromanspacer{} เอวเมว โข ภิกฺขเว ภิกฺขุ อิติ ปฏิสญฺจิกฺขติ “ยถาสุขํ โข เม วิหรโต อกุสลา ธมฺมา อภิวฑฺฒนฺติ, กุสลา ธมฺมา ปริหายนฺติ.\csromanspacer{} ทุกฺขาย ปน เม อตฺตานํ ปทหโต อกุสลา ธมฺมา ปริหายนฺติ, กุสลา ธมฺมา อภิวฑฺฒนฺติ.\csromanspacer{} ยํนูนาหํ ทุกฺขาย อตฺตานํ ปทเหยฺยนฺ”ติ.\csromanspacer{} โส ทุกฺขาย อตฺตานํ ปทหติ, ตสฺส ทุกฺขาย อตฺตานํ ปทหโต อกุสลา ธมฺมา ปริหายนฺติ, กุสลา ธมฺมา อภิวฑฺฒนฺติ.\csromanspacer{} โส น อปเรน สมเยน ทุกฺขาย อตฺตานํ ปทหติ.\csromanspacer{} ตํ กิสฺส เหตุ, ยสฺส หิ โส ภิกฺขเว ภิกฺขุ อตฺถาย ทุกฺขาย อตฺตานํ ปทเหยฺย, สฺวาสฺส อตฺโถ อภินิปฺผนฺโน โหติ.\csromanspacer{} ตสฺมา น อปเรน สมเยน ทุกฺขาย อตฺตานํ ปทหติ.\csromanspacer{} เอวมฺปิ ภิกฺขเว สผโล อุปกฺกโม โหติ สผลํ ปธานํ.}

\proseitem{13}{ปุน จปรํ ภิกฺขเว อิธ ตถาคโต โลเก อุปฺปชฺชติ อรหํ สมฺมาสมฺพุทฺโธ วิชฺชาจรณสมฺปนฺโน สุคโต โลกวิทู อนุตฺตโร ปุริสทมฺมสารถิ สตฺถา เทวมนุสฺสานํ พุทฺโธ ภควา.\csromanspacer{} โส อิมํ โลกํ สเทวกํ สมารกํ สพฺรหฺมกํ สสฺสมณพฺราหฺมณิํ ปชํ สเทวมนุสฺสํ สยํ อภิญฺญา สจฺฉิกตฺวา ปเวเทติ, โส ธมฺมํ เทเสติ อาทิกลฺยาณํ มชฺเฌกลฺยาณํ ปริโยสานกลฺยาณํ สาตฺถํ สพฺยญฺชนํ เกวลปริปุณฺณํ \csromanfolio{12}ปริสุทฺธํ พฺรหฺมจริยํ ปกาเสติ, ตํ ธมฺมํ สุณาติ คหปติ วา คหปหิปุตฺโต วา อญฺญตรสฺมิํ วา กุเล ปจฺจาชาโต.\csromanspacer{} โส ตํ ธมฺมํ สุตฺวา ตถาคเต สทฺธํ ปฏิลภติ, โส เตน สทฺธาปฏิลาเภน สมนฺนาคโต อิติ ปฏิสญฺจิกฺขติ “สมฺพาโธ ฆราวาโส รชาปโถ, อพฺโภกาโส ปพฺพชฺชา, นยิทํ สุกรํ อคารํ อชฺฌาวสตา เอกนฺตปริปุณฺณํ เอกนฺตปริสุทฺธํ สงฺขลิขิตํ พฺรหฺมจริยํ จริตุํ, ยํนูนาหํ เกสมสฺสุํ โอหาเรตฺวา กาสายานิ วตฺถานิ อจฺฉาเทตฺวา อคารสฺมา อนคาริยํ ปพฺพเชยฺยนฺ”ติ.\csromanspacer{} โส อปเรน สมเยน อปฺปํ วา โภคกฺขนฺธํ ปหาย มหนฺตํ วา โภคกฺขนฺธํ ปหาย อปฺปํ วา ญาติปริวฏฺฏํ ปหาย มหนฺตํ วา ญาติปริวฏฺฏํ ปหาย เกสมสฺสุํ โอหาเรตฺวา กาสายานิ วตฺถานิ อจฺฉาเทตฺวา อคารสฺมา อนคาริยํ ปพฺพชติ.}

\proseitem{14}{โส เอวํ ปพฺพชิโต สมาโน ภิกฺขูนํ สิกฺขาสาชีวสมาปนฺโน ปาณาติปาตํ ปหาย ปาณาติปาตา ปฏิวิรโต โหติ นิหิตทณฺโฑ นิหิตสตฺโถ ลชฺชี ทยาปนฺโน, สพฺพปาณภูตหิตานุกมฺปี วิหรติ.\csromanspacer{} อทินฺนาทานํ ปหาย อทินฺนาทานา ปฏิวิรโต โหติ ทินฺนาทายี ทินฺนปาฏิกงฺขี, อเถเนน สุจิภูเตน อตฺตานา วิหรติ.\csromanspacer{} อพฺรหฺมจริยํ ปหาย พฺรหฺมจารี โหติ อาราจารี วิรโต เมถุนา คามธมฺมา.\csromanspacer{} มุสาวาทํ ปหาย มุสาวาทา ปฏิวิรโต โหติ สจฺจวาที สจฺจสนฺโธ เถโต ปจฺจยิโก อวิสํวาทโก โลกสฺส.\csromanspacer{} ปิสุณํ วาจํ ปหาย ปิสุณาย วาจาย ปฏิวิรโต โหติ, อิโต สุตฺวา น อมุตฺร อกฺขาตา อิเมสํ เภทาย, อมุตฺร วา สุตฺวา น อิเมสํ อกฺขาตา อมูสํ เภทาย, อิติ ภินฺนานํ วา สนฺธาตา สหิตานํ วา อนุปฺปทาตา สมคฺคาราโม สมคฺครโต สมคฺคนนฺที สมคฺคกรณิํ วาจํ ภาสิตา โหติ.\csromanspacer{} ผรุสํ วาจํ ปหาย ผรุสาย วาจาย ปฏิวิรโต โหติ, ยา สา วาจา เนลา กณฺณสุขา เปมนียา หทยงฺคมา โปรี พหุชนากนฺตา พหุชนมนาปา, ตถารูปิํ วาจํ ภาสิตา โหติ.\csromanspacer{} สมฺผปฺปลาปํ ปหาย สมฺผปฺปลาปา ปฏิวิรโต โหติ กาลวาที ภูตวาที อตฺถวาที ธมฺมวาที วินยวาที, นิธานวติํ วาจํ ภาสิตา กาเลน สาปเทสํ ปริยนฺตวติํ อตฺถสํหิตํ.\csromanspacer{} โส พีชคามภูตคามสมารมฺภา ปฏิวิรโต โหติ.\csromanspacer{} เอกภตฺติโก โหติ รตฺตูปรโต, วิรโต วิกาลโภชนา.\csromanspacer{} นจฺจคีตวาทิตวิสูกทสฺสนา ปฏิวิรโต โหติ. \csromanfolio{13}มาลาคนฺธวิเลปนธารณมณฺฑนวิภูสนฏฺฐานา ปฏิวิรโต โหติ.\csromanspacer{} อุจฺจาสยนมหาสยนา ปฏิวิรโต โหติ.\csromanspacer{} ชาตรูปรชตปฏิคฺคหณา ปฏิวิรโต โหติ.\csromanspacer{} อามกธญฺญปฏิคฺคหณา ปฏิวิรโต โหติ.\csromanspacer{} อามกมํสปฏิคฺคหณา ปฏิวิรโต โหติ.\csromanspacer{} อิตฺถิกุมาริกปฏิคฺคหณา ปฏิวิรโต โหติ.\csromanspacer{} ทาสิทาสปฏิคฺคหณา ปฏิวิรโต โหติ.\csromanspacer{} อเชฬกปฏิคฺคหณา ปฏิวิรโต โหติ.\csromanspacer{} กุกฺกุฏสูกรปฏิคฺคหณา ปฏิวิรโต โหติ.\csromanspacer{} หตฺถิควสฺสวฬวปฏิคฺคหณา ปฏิวิรโต โหติ.\csromanspacer{} เขตฺตวตฺถุปฏิคฺคหณา ปฏิวิรโต โหติ.\csromanspacer{} ทูเตยฺยปหิณคมนานุโยคา ปฏิวิรโต โหติ.\csromanspacer{} กยวิกฺกยา ปฏิวิรโต โหติ.\csromanspacer{} ตุลากูฏกํสกูฏมานกูฏา ปฏิวิรโต โหติ.\csromanspacer{} อุกฺโกฏนวญฺจนนิกติสาจิโยคา\footnote{สาวิโยคา (สฺยา, กํ, ก) เอตฺถ สาจิสทฺโท กุฏิลปริยาโย.} ปฏิวิรโต โหติ.\csromanspacer{} เฉทนวธพนฺธนวิปราโมส-อาโลปสหสาการา ปฏิวิรโต โหติ\footnote{ปสฺส ม 1 จูฬหตฺถิปโทปเม (238) ปิฏฺเฐ.}.}

\prose{โส สนฺตุฏฺโฐ โหติ กายปริหาริเกน จีวเรน กุจฺฉิปริหาริเกน ปิณฺฑปาเตน, โส เยน เยเนว ปกฺกมติ, สมาทาเยว ปกฺกมติ.\csromanspacer{} เสยฺยถาปิ นาม ปกฺขี สกุโณ เยน เยเนว เฑติ, สปตฺตภาโรว เฑติ.\csromanspacer{} เอวเมว ภิกฺขุ สนฺตุฏฺโฐ โหติ กายปริหาริเกน จีวเรน กุจฺฉิปริหาริเกน ปิณฺฑปาเตน, โส เยน เยเนว ปกฺกมติ, สมาทาเยว ปกฺกมติ.\csromanspacer{} โส อิมินา อริเยน สีลกฺขนฺเธน สมนฺนาคโต อชฺฌตฺตํ อนวชฺชสุขํ ปฏิสํเวเทติ.}

\proseitem{15}{โส จกฺขุนา รูปํ ทิสฺวา น นิมิตฺตคฺคาหี โหติ นานุพฺยญฺชนคฺคาหี, ยตฺวาธิกรณเมนํ จกฺขุนฺทฺริยํ อสํวุตํ วิหรนฺตํ อภิชฺฌาโทมนสฺสา ปาปกา อกุสลา ธมฺมา อนฺวาสฺสเวยฺยุํ, ตสฺส สํวราย ปฏิปชฺชติ, รกฺขติ จกฺขุนฺทฺริยํ, จกฺขุนฺทฺริเย สํวรํ อาปชฺชติ.\csromanspacer{} โสเตน สทฺทํ สุตฺวา ฯเปฯ.\csromanspacer{} ฆาเนน คนฺธํ ฆายิตฺวา ฯเปฯ.\csromanspacer{} ชิวฺหาย รสํ สายิตฺวา ฯเปฯ.\csromanspacer{} กาเยน โผฏฺฐพฺพํ ผุสิตฺวา ฯเปฯ.\csromanspacer{} มนสา ธมฺมํ วิญฺญาย น นิมิตฺตคฺคาหี โหติ นานุพฺยญฺชนคฺคาหี, ยตฺวาธิกรณเมนํ มนินฺทฺริยํ อสํวุตํ วิหรนฺตํ อภิชฺฌาโทมนสฺสา ปาปกา อกุสลา ธมฺมา อนฺวาสฺสเวยฺยุํ, ตสฺส สํวราย ปฏิปชฺชติ, รกฺขติ มนินฺทฺริยํ, มนินฺทฺริเย สํวรํ อาปชฺชติ.\csromanspacer{} โส อิมินา อริเยน อินฺทฺริยสํวเรน สมนฺนาคโต อชฺฌตฺตํ อพฺยาเสกสุขํ ปฏิสํเวเทติ.}

\prose{\csromanfolio{14}โส อภิกฺกนฺเต ปฏิกฺกนฺเต สมฺปชานการี โหติ, อาโลกิเต วิโลกิเต สมฺปชานการี โหติ, สมิญฺชิเต\footnote{สมฺมิญฺชิเต (สี, สฺยา, กํ, อิ)} ปสาริเต สมฺปชานการี โหติ, สํฆาฏิปตฺตจีวรธารเณ สมฺปชานการี โหติ, อสิเต ปีเต ขายิเต สายิเต สมฺปชานการี โหติ, อุจฺจารปสฺสาวกมฺเม สมฺปชานการี โหติ, คเต ฐิเต นิสินฺเน สุตฺเต ชาคริเต ภาสิเต ตุณฺหีภาเว สมฺปชานการี โหติ.}

\proseitem{16}{โส อิมินา จ อริเยน สีลกฺขนฺเธน สมนฺนาคโต (อิมาย จ อริยาย สนฺตุฏฺฐิยา สมนฺนาคโต)\footnote{ปสฺส ม 1 จูฬหตฺถิปโทปเม (239) ปิฏฺเฐ.} อิมินา จ อริเยน อินฺทฺริยสํวเรน สมนฺนาคโต อิมินา จ อริเยน สติสมฺปชญฺเญน สมนฺนาคโต วิวิตฺตํ เสนาสนํ ภชติ อรญฺญํ รุกฺขมูลํ ปพฺพตํ กนฺทรํ คิริคุหํ สุสานํ วนปตฺถํ อพฺโภกาสํ ปลาลปุญฺชํ.\csromanspacer{} โส ปจฺฉาภตฺตํ ปิณฺฑปาตปฏิกฺกนฺโต นิสีทติ ปลฺลงฺกํ อาภุชิตฺวา อุชุํ กายํ ปณิธาย ปริมุขํ สติํ อุปฏฺฐเปตฺวา, โส อภิชฺฌํ โลเก ปหาย วิคตาภิชฺเฌน เจตสา วิหรติ, อภิชฺฌาย จิตฺตํ ปริโสเธติ.\csromanspacer{} พฺยาปาทปโทสํ ปหาย อพฺยาปนฺนจิตฺโต วิหรติ สพฺพปาณภูตหิตานุกมฺปี, พฺยาปาทปโทสา จิตฺตํ ปริโสเธติ.\csromanspacer{} ถินมิทฺธํ ปหายวิคตถินมิทฺโธ วิหรติ อาโลกสญฺญี สโต สมฺปชาโน, ถินมิทฺธา จิตฺตํ ปริโสเธติ.\csromanspacer{} อุทฺธจฺจกุกฺกุจฺจํ ปหาย อนุทฺธโต วิหรติ อชฺฌตฺตํ วูปสนฺตจิตฺโต, อุทฺธจฺจกุกฺกุจฺจา จิตฺตํ ปริโสเธติ.\csromanspacer{} วิจิกิจฺฉํ ปหาย ติณฺณวิจิกิจฺโฉ วิหรติ อกถํกถี กุสเลสุ ธมฺเมสุ, วิจิกิจฺฉาย จิตฺตํ ปริโสเธติ.\csromanspacer{} โส อิเม ปญฺจ นีวรเณ ปหาย เจตโส อุปกฺกิเลเส ปญฺญาย ทุพฺพลีกรเณ วิวิจฺเจว กาเมหิ วิวิจฺจ อกุสเลหิ ธมฺเมหิ สวิตกฺกํ สวิจารํ วิเวกชํ ปีติสุขํ ปฐมํ ฌานํ อุปสมฺปชฺช วิหรติ.\csromanspacer{} เอวมฺปิ ภิกฺขเว สผโล อุปกฺกโม โหติ สผลํ ปธานํ.}

\prose{ปุน จปรํ ภิกฺขเว ภิกฺขุ วิตกฺกวิจารานํ วูปสมา อชฺฌตฺตํ สมฺปสาทนํ เจตโส เอโกทิภาวํ อวิตกฺกํ อวิจารํ สมาธิชํ ปีติสุขํ ทุติยํ ฌานํ อุปสมฺปชฺช วิหรติ.\csromanspacer{} เอวมฺปิ ภิกฺขเว สผโล อุปกฺกโม โหติ สผลํ ปธานํ.}

\prose{\csromanfolio{15}ปุน จปรํ ภิกฺขเว ภิกฺขุ ปีติยา จ วิราคา อุเปกฺขโก จ วิหรติ สโต จ สมฺปชาโน สุขญฺจ กาเยน ปฏิสํเวเทติ, ยํ ตํ อริยา อาจิกฺขนฺติ “อุเปกฺขโก สติมา สุขวิหารี”ติ, ตติยํ ฌานํ อุปสมฺปชฺช วิหรติ.\csromanspacer{} เอวมฺปิ ภิกฺขเว สผโล อุปกฺกโม โหติ สผลํ ปธานํ.}

\prose{ปุน จปรํ ภิกฺขเว ภิกฺขุ สุขสฺส จ ปหานา ทุกฺขสฺส จ ปหานา ปุพฺเพว โสมนสฺสโทมนสฺสานํ อตฺถงฺคมา อทุกฺขมสุขํ อุเปกฺขาสติปาริสุทฺธิํ จตุตฺถํ ฌานํ อุปสมฺปชฺช วิหรติ.\csromanspacer{} เอวมฺปิ ภิกฺขเว สผโล อุปกฺกโม โหติ สผลํ ปธานํ.}

\proseitem{17}{โส เอวํ สมาหิเต จิตฺเต ปริสุทฺเธ ปริโยทาเต อนงฺคเณ วิคตูปกฺกิเลเส มุทุภูเต กมฺมนิเย ฐิเต อาเนญฺชปฺปตฺเต ปุพฺเพนิวาสานุสฺสติญาณาย จิตฺตํ อภินินฺนาเมติ.\csromanspacer{} โส อเนกวิหิตํ ปุพฺเพนิวาสํ อนุสฺสรติ, เสยฺยถิทํ\footnote{เสยฺยถีทํ (สี, สฺยา, กํ, อิ)}, เอกมฺปิ ชาติํ เทฺวปิ ชาติโย ติสฺโสปิ ชาติโย จตสฺโสปิ ชาติโย ปญฺจปิ ชาติโย ทสปิ ชาติโย วีสมฺปิ ชาติโย ติํสมฺปิ ชาติโย จตฺตาลีสมฺปิ ชาติโย ปญฺญาสมฺปิ ชาติโย ชาติสตมฺปิ ชาติสหสมฺปิ ชาติสตสหสฺสมฺปิ อเนเกปิ สํวฏฺฏกปฺเป อเนเกปิ วิวฏฺฏกปฺเป อเนเกปิ สํวฏฺฏวิวฏฺฏกปฺเป, “อมุตฺราสิํ เอวํนาโม เอวํโคตฺโต เอวํวณฺโณ เอวมาหาโร เอวํสุขทุกฺขปฺปฏิสํเวที เอวมายุปริยนฺโต, โส ตโต จุโต อมุตฺร อุทปาทิํ, ตตฺราปาสิํ เอวํนาโม เอวํโคตฺโต เอวํวณฺโณ เอวมาหาโร เอวํสุขทุกฺขปฺปฏิสํเวที เอวมายุปริยนฺโต, โส ตโต จุโต อิธูปปนฺโน”ติ, อิติ สาการํ ส-อุทฺเทสํ อเนกวิหิตํ ปุพฺเพนิวาสํ อนุสฺสรติ.\csromanspacer{} เอวมฺปิ ภิกฺขเว สผโล อุปกฺกโม โหติ สผลํ ปธานํ.}

\proseitem{18}{โส เอวํ สมาหิเต จิตฺเต ปริสุทฺเธ ปริโยทาเต อนงฺคเณ วิคตูปกฺกิเลเส มุทุภูเต กมฺมนิเย ฐิเต อาเนญฺชปฺปตฺเต สตฺตานํ จุตูปปาตญาณาย จิตฺตํ อภินินฺนาเมติ.\csromanspacer{} โส ทิพฺเพน จกฺขุนา วิสุทฺเธน อติกฺกนฺตมานุสเกน สตฺเต ปสฺสติ จวมาเน อุปปชฺชมาเน หีเน ปณีเต สุวณฺเณ ทุพฺพณฺเณ สุคเต ทุคฺคเต ยถากมฺมูปเค สตฺเต \csromanfolio{16}ปชานาติ. “อิเม วต โภนฺโต สตฺตา กายทุจฺจริเตน สมนฺนาคตา วจีทุจฺจริเตน สมนฺนาคตา มโนทุจฺจริเตน สมนฺนาคตา, อริยานํ อุปวาทกา มิจฺฉาทิฏฺฐิกา มิจฺฉาทิฏฺฐิกมฺมสมาทานา, เต กายสฺส เภทา ปรํ มรณา อปายํ ทุคฺคติํ วินิปาตํ นิรยํ อุปปนฺนา.\csromanspacer{} อิเม วา ปน โภนฺโต สตฺตา กายสุจริเตน สมนฺนาคตา วจีสุจริเตน สมนฺนาคตา มโนสุจริเตน สมนฺนาคตา อริยานํ อนุปวาทกา สมฺมาทิฏฺฐิกา สมฺมาทิฏฺฐิกมฺมสมาทานา, เต กายสฺส เภทา ปรํ มรณา สุคติํ สคฺคํ โลกํ อุปปนฺนา”ติ.\csromanspacer{} อิติ ทิพฺเพน จกฺขุนา วิสุทฺเธน อติกฺกนฺตมานุสเกน สตฺเต ปสฺสติ จวมาเน อุปปชฺชมาเน หีเน ปณีเต สุวณฺเณ ทุพฺพณฺเณ สุคเต ทุคฺคเต ยถากมฺมูปเค สตฺเต ปชานาติ.\csromanspacer{} เอวมฺปิ ภิกฺขเว สผโล อุปกฺกโม โหติ สผลํ ปธานํ.}

\proseitem{19}{โส เอวํ สมาหิเต จิตฺเต ปริสุทฺเธ ปริโยทาเต อนงฺคเณ วิคตูปกฺกิเลเส มุทุภูเต กมฺมนิเย ฐิเต อาเนญฺชปฺปตฺเต อาสวานํ ขยญาณาย จิตฺตํ อภินินฺนาเมติ.\csromanspacer{} โส อิทํ ทุกฺขนฺติ ยถาภูตํ ปชานาติ, อยํ ทุกฺขสมุทโยติ ยถาภูตํ ปชานาติ, อยํ ทุกฺขนิโรโธติ ยถาภูตํ ปชานาติ, อยํ ทุกฺขนิโรธคามินี ปฏิปทาติ ยถาภูตํ ปชานาติ.\csromanspacer{} อิเม อาสวาติ ยถาภูตํ ปชานาติ, อยํ อาสวสมุทโยติ ยถาภูตํ ปชานาติ, อยํ อาสวนิโรโธติ ยถาภูตํ ปชานาติ, อยํ อาสวนิโรธคามินี ปฏิปทาติ ยถาภูตํ ปชานาติ.\csromanspacer{} ตสฺส เอวํ ชานโต เอวํ ปสฺสโต กามาสวาปิ จิตฺตํ วิมุจฺจติ, ภวาสวาปิ จิตฺตํ วิมุจฺจติ, อวิชฺชาสวาปิ จิตฺตํ วิมุจฺจติ, วิมุตฺตสฺมิํ “วิมุตฺตมฺ”อิติ ญาณํ โหติ, “ขีณา ชาติ, วุสิตํ พฺรหฺมจริยํ, กตํ กรณียํ, นาปรํ อิตฺถตฺตายา”ติ ปชานาติ.\csromanspacer{} เอวมฺปิ โข ภิกฺขเว สผโล อุปกฺกโม โหติ สผลํ ปธานํ.\csromanspacer{} เอวํวาที ภิกฺขเว ตถาคตา, เอวํวาทีนํ ภิกฺขเว ตถาคตานํ\footnote{ตถาคโต, เอวํวาทิํ ภิกฺขเว ตถาคตํ (สี, สฺยา, กํ, อิ)} ทส สหธมฺมิกา ปาสํสฏฺฐานา อาคจฺฉนฺติ.}

\proseitem{20}{สเจ ภิกฺขเว สตฺตา ปุพฺเพกตเหตุ สุขทุกฺขํ ปฏิสํเวเทนฺติ, อทฺธา ภิกฺขเว ตถาคโต ปุพฺเพ สุกตกมฺมการี, ยํ เอตรหิ เอวรูปา อนาสวา สุขา เวทนา เวเทติ.\csromanspacer{} สเจ ภิกฺขเว สตฺตา อิสฺสรนิมฺมานเหตุ สุขทุกฺขํ ปฏิสํเวเทนฺติ, อทฺธา ภิกฺขเว ตถาคโต ภทฺทเกน \csromanfolio{17}อิสฺสเรน นิมฺมิโต, ยํ เอตรหิ เอวรูปา อนาสวา สุขา เวเทนา เวเทติ.\csromanspacer{} สเจ ภิกฺขเว สตฺตา สงฺคติภาวเหตุ สุขทุกฺขํ ปฏิสํเวเทนฺติ, อทฺธา ภิกฺขเว ตถาคโต กลฺยาณสงฺคติโก, ยํ เอตรหิ เอวรูปา อนาสวา สุขา เวทนา เวเทติ.\csromanspacer{} สเจ ภิกฺขเว สตฺตา อภิชาติเหตุ สุขทุกฺขํ ปฏิสํเวเทนฺติ, อทฺธา ภิกฺขเว ตถาคโต กลฺยาณาภิชาติโก, ยํ เอตรหิ เอวรูปา อนาสวา สุขา เวทนา เวเทติ.\csromanspacer{} สเจ ภิกฺขเว สตฺตา ทิฏฺฐธมฺมูปกฺกมเหตุ สุขทุกฺขํ ปฏิสํเวเทนฺติ, อทฺธา ภิกฺขเว ตถาคโต กลฺยาณทิฏฺฐธมฺมูปกฺกโม, ยํ เอตรหิ เอวรูปา อนาสวา สุขา เวทนา เวเทติ.}

\prose{สเจ ภิกฺขเว สตฺตา ปุพฺเพกตเหตุ สุขทุกฺขํ ปฏิสํเวเทนฺติ, ปาสํโส ตถาคโต.\csromanspacer{} โน เจ สตฺตา ปุพฺเพกตเหตุ สุขทุกฺขํ ปฏิสํเวเทนฺติ, ปาสํโส ตถาคโต.\csromanspacer{} สเจ ภิกฺขเว สตฺตา อิสฺสรนิมฺมานเหตุ สุขทุกฺขํ ปฏิสํเวเทนฺติ, ปาสํโส ตถาคโต.\csromanspacer{} โน เจ สตฺตา อิสฺสรนิมฺมานเหตุ สุขทุกฺขํ ปฏิสํเวเทนฺติ, ปาสํโส ตถาคโต.\csromanspacer{} สเจ ภิกฺขเว สตฺตา สงฺคติภาวเหตุ สุขทุกฺขํ ปฏิสํเวเทนฺติ, ปาสํโส ตถาคโต.\csromanspacer{} โน เจ สตฺตา สงฺคติภาวเหตุ สุขทุกฺขํ ปฏิสํเวเทนฺติ, ปาสํโส ตถาคโต.\csromanspacer{} สเจ ภิกฺขเว สตฺตา อภิชาติเหตุ สุขทุกฺขํ ปฏิสํเวเทนฺติ, ปาสํโส ตถาคโต.\csromanspacer{} โน เจ สตฺตา อภิชาติเหตุ สุขทุกฺขํ ปฏิสํเวเทนฺติ, ปาสํโส ตถาคโต.\csromanspacer{} สเจ ภิกฺขเว สตฺตา ทิฏฺฐธมฺมูปกฺกมเหตุ สุขทุกฺขํ ปฏิสํเวเทนฺติ, ปาสํโส ตถาคโต.\csromanspacer{} โน เจ สตฺตา ทิฏฺฐธมฺมูปกฺกมเหตุ สุขทุกฺขํ ปฏิสํเวเทนฺติ, ปาสํโส ตถาคโต.\csromanspacer{} เอวํวาที ภิกฺขเว ตถาคตา, เอวํวาทีนํ ภิกฺขเว ตถาคตานํ อิเม ทส สหธมฺมิกา ปาสํสฏฺฐานา อาคจฺฉนฺตีติ.}

\prose{อิทมโวจ ภควา.\csromanspacer{} อตฺตมนา เต ภิกฺขู ภควโต ภาสิตํ อภินนฺทุนฺติ.}

\nitthitam{เทวทหสุตฺตํ นิฏฺฐิตํ ปฐมํ.}
\csromanmatikaanchor{mtk.3}
\csromantocmarkref{subsection}{2. ปญฺจตฺตยสุตฺต}{mtk.3}
\bookmark[dest={mtk.3},level=2]{2. ปญฺจตฺตยสุตฺต}
\csromanheader{2}{2. ปญฺจตฺตยสุตฺต\footnote{ปญฺจายตนสุตฺต (ก)}}
\proseitem{21}{\csromanfolio{18}เอวํ เม สุตํ\csromandash{}เอกํ สมยํ ภควา สาวตฺถิยํ วิหรติ เชตวเน อนาถปิณฺฑิกสฺส อาราเม.\csromanspacer{} ตตฺร โข ภควา ภิกฺขู อามนฺเตสิ “ภิกฺขโว”ติ. “ภทนฺเต”ติ เต ภิกฺขู ภควโต ปจฺจสฺโสสุํ.\csromanspacer{} ภควา เอตทโวจ\csromandash{}สนฺติ ภิกฺขเว เอเก สมณพฺราหฺมณา อปรนฺตกปฺปิกา อปรนฺตานุทิฏฺฐิโน อปรนฺตํ อารพฺภ อเนกวิหิตานิ อธิวุตฺติปทานิ\footnote{อธิมุตฺติปทานิ (สฺยา, กํ, ก)} อภิวทนฺติ, “สญฺญี อตฺตา โหติ อโรโค ปรํ มรณา”ติ อิตฺเถเก อภิวทนฺติ, “อสญฺญี อตฺตา โหติ อโรโค ปรํ มรณา”ติ อิตฺเถเก อภิวทนฺติ, “เนวสญฺญีนาสญฺญี อตฺตา โหติ อโรโค ปรํ มรณา”ติ อิตฺเถเก อภิวทนฺติ, สโต วา ปน สตฺตสฺส อุจฺเฉทํ วินาสํ วิภวํ ปญฺญเปนฺติ\footnote{ปญฺญาเปนฺติ (สี, สฺยา, กํ, อิ) 4-4. ปรํ มรณา. อิติ อิมานิ (ก)}, ทิฏฺฐธมฺมนิพฺพานํ วา ปเนเก อภิวทนฺติ.\csromanspacer{} อิติ สนฺตํ วา อตฺตานํ ปญฺญเปนฺติ อโรคํ4 ปรํ มรณา, สโต วา ปน สตฺตสฺส อุจฺเฉทํ วินาสํ วิภวํ ปญฺญเปนฺติ, ทิฏฺฐธมฺมนิพฺพานํ วา ปเนเก อภิวทนฺติ.\csromanspacer{} อิติ อิมานิ4 ปญฺจ หุตฺวา ตีณิ โหนฺติ, ตีณิ หุตฺวา ปญฺจ โหนฺติ.\csromanspacer{} อยมุทฺเทโส ปญฺจตฺตยสฺส.}

\proseitem{22}{ตตฺร ภิกฺขเว เย เต สมณพฺราหฺมณา สญฺญิํ อตฺตานํ ปญฺญเปนฺติ อโรคํ ปรํ มรณา.\csromanspacer{} รูปิํ วา เต โภนฺโต สมณพฺราหฺมณา สญฺญิํ อตฺตานํ ปญฺญเปนฺติ อโรคํ ปรํ มรณา.\csromanspacer{} อรูปิํ วา เต โภนฺโต สมณพฺราหฺมณา สญฺญิํ อตฺตานํ ปญฺญเปนฺติ อโรคํ ปรํ มรณา.\csromanspacer{} รูปิํ จ อรูปิํ จ วา เต โภนฺโต สมณพฺราหฺมณา สญฺญิํ อตฺตานํ ปญฺญเปนฺติ อโรคํ ปรํ มรณา.\csromanspacer{} เนวรูปิํ นารูปิํ วา เต โภนฺโต สมณพฺราหฺมณา สญฺญิํ อตฺตานํ ปญฺญเปนฺติ อโรคํ ปรํ มรณา.\csromanspacer{} เอกตฺตสญฺญิํ วา เต โภนฺโต สมณพฺราหฺมณา สญฺญิํ อตฺตานํ ปญฺญเปนฺติ อโรคํ ปรํ มรณา.\csromanspacer{} นานตฺตสญฺญิํ วา เต โภนฺโต สมณพฺราหฺมณา สญฺญิํ อตฺตานํ ปญฺญเปนฺติ อโรคํ ปรํ มรณา.\csromanspacer{} ปริตฺตสญฺญิํ วา เต โภนฺโต สมณพฺราหฺมณา สญฺญิํ อตฺตานํ ปญฺญเปนฺติ อโรคํ ปรํ มรณา.\csromanspacer{} อปฺปมาณสญฺญิํ วา เต โภนฺโต สมณพฺราหฺมณา สญฺญิํ อตฺตานํ ปญฺญเปนฺติ อโรคํ ปรํ มรณา.\csromanspacer{} เอตํ\footnote{เอวํ (ก)} วา ปเนเกสํ\footnote{ปเนเตสํ (สฺยา, กํ)} อุปาติวตฺตตํ วิญฺญาณกสิณเมเก อภิวทนฺติ อปฺปมาณํ}

\vspace{\tipitakaparskip}
\csromancenter{ปญฺจตฺตยสุตฺต (102) อาเนญฺชํ.\csromanspacer{} ตยิทํ ภิกฺขเว ตถาคโต อภิชานาติ\footnote{ปชานาติ (สี, สฺยา, กํ, อิ) อฏฺฐกถา โอโลเกตพฺพา.}.\csromanspacer{} เย โข เต โภนฺโต สมณพฺราหฺมณา สญฺญิํ อตฺตานํ ปญฺญเปนฺติ อโรคํ ปรํ มรณา.\csromanspacer{} รูปิํ วา เต โภนฺโต สมณพฺราหฺมณา สญฺญิํ อตฺตานํ ปญฺญเปนฺติ อโรคํ ปรํ มรณา.\csromanspacer{} อรูปิํ วา เต โภนฺโต สมณพฺราหฺมณา สญฺญิํ อตฺตานํ ปญฺญเปนฺติ อโรคํ ปรํ มรณา.\csromanspacer{} รูปิํ จ อรูปิํ จ วา เต โภนฺโต สมณพฺราหฺมณา สญฺญิํ อตฺตานํ ปญฺญเปนฺติ อโรคํ ปรํ มรณา.\csromanspacer{} เนวรูปิํ นารูปิํ วา เต โภนฺโต สมณพฺราหฺมณา สญฺญิํ อตฺตานํ ปญฺญเปนฺติ อโรคํ ปรํ มรณา.\csromanspacer{} เอกตฺตสญฺญิํ วา เต โภนฺโต สมณพฺราหฺมณา สญฺญิํ อตฺตานํ ปญฺญเปนฺติ อโรคํ ปรํ มรณา.\csromanspacer{} นานตฺตสญฺญิํ วา เต โภนฺโต สมณพฺราหฺมณา สญฺญิํ อตฺตานํ ปญฺญเปนฺติ อโรคํ ปรํ มรณา.\csromanspacer{} ปริตฺตสญฺญิํ วา เต โภนฺโต สมณพฺราหฺมณา สญฺญิํ อตฺตานํ ปญฺญเปนฺติ อโรคํ ปรํ มรณา.\csromanspacer{} อปฺปมาณสญฺญิํ วา เต โภนฺโต สมณพฺราหฺมณา สญฺญิํ อตฺตานํ ปญฺญเปนฺติ อโรคํ ปรํ มรณา\footnote{มรณาติ (ก)}.\csromanspacer{} ยา วา ปเนตาสํ สญฺญานํ ปริสุทฺธา ปรมา อคฺคา อนุตฺตริยา อกฺขายติ, ยทิ รูปสญฺญานํ, ยทิ อรูปสญฺญานํ, ยทิ เอกตฺตสญฺญานํ, ยทิ นานตฺตสญฺญานํ.\csromanspacer{} นตฺถิ กิญฺจีติ อากิญฺจญฺญายตนเมเก อภิวทนฺติ อปฺปมาณํ อาเนญฺชํ.\csromanspacer{} ตยิทํ สงฺขตํ โอฬาริกํ, อตฺถิ โข ปน สงฺขารานํ นิโรโธ, อตฺเถตนฺติ อิติ วิทิตฺวา ตสฺส นิสฺสรณทสฺสาวี ตถาคโต ตทุปาติวตฺโต.}
\proseitem{23}{\csromanfolio{19}ตตฺร ภิกฺขเว เย เต สมณพฺราหฺมณา อสญฺญิํ อตฺตานํ ปญฺญเปนฺติ อโรคํ ปรํ มรณา.\csromanspacer{} รูปิํ วา เต โภนฺโต สมณพฺราหฺมณา อสญฺญิํ อตฺตานํ ปญฺญเปนฺติ อโรคํ ปรํ มรณา.\csromanspacer{} อรูปิํ วา เต โภนฺโต สมณพฺราหฺมณา อสญฺญิํ อตฺตานํ ปญฺญเปนฺติ อโรคํ ปรํ มรณา.\csromanspacer{} รูปิํ จ อรูปิํ จ วา เต โภนฺโต สมณพฺราหฺมณา อสญฺญิํ อตฺตานํ ปญฺญเปนฺติ อโรคํ ปรํ มรณา.\csromanspacer{} เนวรูปิํ นารูปิํ วา เต โภนฺโต สมณพฺราหฺมณา อสญฺญิํ อตฺตานํ ปญฺญเปนฺติ อโรคํ ปรํ มรณา.\csromanspacer{} ตตฺร ภิกฺขเว เย เต สมณพฺราหฺมณา สญฺญิํ อตฺตานํ ปญฺญเปนฺติ อโรคํ ปรํ มรณา.\csromanspacer{} เตสเมเต ปฏิกฺโกสนฺติ.\csromanspacer{} ตํ กิสฺส เหตุ, สญฺญา โรโค สญฺญา คณฺโฑ สญฺญา สลฺลํ.\csromanspacer{} เอตํ สนฺตํ เอตํ ปณีตํ, ยทิทํ อสญฺญนฺติ.\csromanspacer{} ตยิทํ ภิกฺขเว ตถาคโต อภิชานาติ.\csromanspacer{} เย โข เต โภนฺโต สมณพฺราหฺมณา อสญฺญิํ อตฺตานํ ปญฺญเปนฺติ อโรคํ ปรํ มรณา.\csromanspacer{} รูปิํ วา เต \csromanfolio{20}โภนฺโต สมณพฺราหฺมณา อสญฺญิํ อตฺตานํ ปญฺญเปนฺติ อโรคํ ปรํ มรณา.\csromanspacer{} อรูปิํ วา เต โภนฺโต สมณพฺราหฺมณา อสญฺญิํ อตฺตานํ ปญฺญเปนฺติ อโรคํ ปรํ มรณา.\csromanspacer{} รูปิํ จ อรูปิํ จ วา เต โภนฺโต สมณพฺราหฺมณา อสญฺญิํ อตฺตานํ ปญฺญเปนฺติ อโรคํ ปรํ มรณา.\csromanspacer{} เนวรูปิํ นารูปิํ วา เต โภนฺโต สมณพฺราหฺมณา อสญฺญิํ อตฺตานํ ปญฺญเปนฺติ อโรคํ ปรํ มรณา.\csromanspacer{} โย หิ โกจิ ภิกฺขเว สมโณ วา พฺราหฺมโณ วา เอวํ วเทยฺย “อหมญฺญตฺร รูปา อญฺญตฺร เวทนาย อญฺญตฺร สญฺญาย อญฺญตฺร สงฺขาเรหิ วิญฺญาณสฺส\footnote{อญฺญตฺร วิญฺญาณา (สฺยา, กํ), อญฺญตฺร วิญฺญาเณน (ก)} อาคติํ วา คติํ วา จุติํ วา อุปปตฺติํ วา วุทฺธิํ วา วิรูฬฺหิํ วา เวปุลฺลํ วา ปญฺญเปสฺสามี”ติ, เนตํ ฐานํ วิชฺชติ.\csromanspacer{} ตยิทํ สงฺขตํ โอฬาริกํ, อตฺถิ โข ปน สงฺขารานํ นิโรโธ, อตฺเถตนฺติ อิติ วิทิตฺวา ตสฺส นิสฺสรณทสฺสาวี ตถาคโต ตทุปาติวตฺโต.}

\proseitem{24}{ตตฺร ภิกฺขเว เย เต สมณพฺราหฺมณา เนวสญฺญีนาสญฺญิํ อตฺตานํ ปญฺญเปนฺติ อโรคํ ปรํ มรณา.\csromanspacer{} รูปิํ วา เต โภนฺโต สมณพฺราหฺมณา เนวสญฺญีนาสญฺญิํ อตฺตานํ ปญฺญเปนฺติ อโรคํ ปรํ มรณา.\csromanspacer{} อรุปิํ วา เต โภนฺโต สมณพฺราหฺมณา เนวสญฺญีนาสญฺญิํ อตฺตานํ ปญฺญเปนฺติ อโรคํ ปรํ มรณา.\csromanspacer{} รูปิํ จ อรูปิํ จ วา เต โภนฺโต สมณพฺราหฺมณา เนวสญฺญีนาสญฺญิํ อตฺตานํ ปญฺญเปนฺติ อโรคํ ปรํ มรณา.\csromanspacer{} เนวรูปิํ นารูปิํ วา เต โภนฺโต สมณพฺราหฺมณา เนวสญฺญีนาสญฺญิํ อตฺตานํ ปญฺญเปนฺติ อโรคํ ปรํ มรณา.\csromanspacer{} ตตฺร ภิกฺขเว เย เต สมณพฺราหฺมณา สญฺญิํ อตฺตานํ ปญฺญเปนฺติ อโรคํ ปรํ มรณา.\csromanspacer{} เตสเมเต ปฏิกฺโกสนฺติ.\csromanspacer{} เยปิ เต โภนฺโต สมณพฺราหฺมณา อสญฺญิํ อตฺตานํ ปญฺญเปนฺติ อโรคํ ปรํ มรณา.\csromanspacer{} เตสเมเต ปฏิกฺโกสนฺติ.\csromanspacer{} ตํ กิสฺส เหตุ, สญฺญา โรโค สญฺญา คณฺโฑ สญฺญา สลฺลํ อสญฺญา สมฺโมโห.\csromanspacer{} เอตํ สนฺตํ เอตํ ปณีตํ, ยทิทํ เนวสญฺญานาสญฺญนฺติ\footnote{เนวสญฺญานาสญฺญาติ (สฺยา, กํ, อิ, ก) เอตนฺติปทํ มนสิกาตพฺพํ.}.\csromanspacer{} ตยิทํ ภิกฺขเว ตถาคโต อภิชานาติ.\csromanspacer{} เย โข เต โภนฺโต สมณพฺราหฺมณา เนวสญฺญีนาสญฺญิํ อตฺตานํ ปญฺญเปนฺติ อโรคํ ปรํ มรณา.\csromanspacer{} รูปิํ วา เต โภนฺโต สมณพฺราหฺมณา เนวสญฺญีนาสญฺญิํ อตฺตานํ ปญฺญเปนฺติ อโรคํ ปรํ มรณา.\csromanspacer{} อรูปิํ วา เต โภนฺโต สมณพฺราหฺมณา เนวสญฺญีนาสญฺญิํ}

\vspace{\tipitakaparskip}
\csromancenter{ปญฺจตฺตยสุตฺต (102) อตฺตานํ ปญฺญเปนฺติ อโรคํ ปรํ มรณา.\csromanspacer{} รูปิํ จ อรูปิํ จ วา เต โภนฺโต สมณพฺราหฺมณา เนวสญฺญีนาสญฺญิํ อตฺตานํ ปญฺญเปนฺติ อโรคํ ปรํ มรณา.\csromanspacer{} เนวรูปิํ นารูปิํ วา เต โภนฺโต สมณพฺราหฺมณา เนวสญฺญีนาสญฺญิํ อตฺตานํ ปญฺญเปนฺติ อโรคํ ปรํ มรณา.\csromanspacer{} เย หิ เกจิ ภิกฺขเว สมณา วา พฺราหฺมณา วา\footnote{สมณพฺราหฺมณา (สี, อิ)} ทิฏฺฐสุตมุตวิญฺญาตพฺพ- สงฺขารมตฺเตน เอตสฺส อายตนสฺส อุปสมฺปทํ ปญฺญเปนฺติ.\csromanspacer{} พฺยสนํ เหตํ ภิกฺขเว อกฺขายติ\footnote{อายตนมกฺขายติ (ก)} เอตสฺส อายตนสฺส อุปสมฺปทาย, น เหตํ ภิกฺขเว อายตนํ สงฺขารสมาปตฺติปตฺตพฺพมกฺขายติ.\csromanspacer{} สงฺขาราวเสสสมาปตฺติปตฺตพฺพเมตํ ภิกฺขเว อายตนมกฺขายติ.\csromanspacer{} ตยิทํ สงฺขตํ โอฬาริกํ, อตฺถิ โข ปน สงฺขารานํ นิโรโธ, อตฺเถตนฺติ อิติ วิทิตฺวา ตสฺส นิสฺสรณทสฺสาวี ตถาคโต ตทุปาติวตฺโต.}
\proseitem{25}{\csromanfolio{21}ตตฺร ภิกฺขเว เย เต สมณพฺราหฺมณา สโต สตฺตสฺส อุจฺเฉทํ วินาสํ วิภวํ ปญฺญเปนฺติ.\csromanspacer{} ตตฺร ภิกฺขเว เย เต สมณพฺราหฺมณา สญฺญิํ อตฺตานํ ปญฺญเปนฺติ อโรคํ ปรํ มรณา, เตสเมเต ปฏิกฺโกสนฺติ.\csromanspacer{} เยปิ เต โภนฺโต สมณพฺราหฺมณา อสญฺญิํ อตฺตานํ ปญฺญเปนฺติ อโรคํ ปรํ มรณา, เตสเมเต ปฏิกฺโกสนฺติ.\csromanspacer{} เยปิ เต โภนฺโต สมณพฺราหฺมณา เนวสญฺญีนาสญฺญิํ อตฺตานํ ปญฺญเปนฺติ อโรคํ ปรํ มรณา, เตสเมเต ปฏิกฺโกสนฺติ.\csromanspacer{} ตํ กิสฺส เหตุ, สพฺเพปิเม โภนฺโต สมณพฺราหฺมณา อุทฺธํ สรํ\footnote{อุทฺธํสรา (สี, อิ), อุทฺธํปรามสนฺติ (สฺยา, กํ) 4. ขีเล (สี, สฺยา, กํ, อิ)} อาสตฺติํเยว อภิวทนฺติ “อิติ เปจฺจ ภวิสฺสาม, อิติ เปจฺจ ภวิสฺสามา”ติ.\csromanspacer{} เสยฺยถาปิ นาม วาณิชสฺส วาณิชฺชาย คจฺฉโต เอวํ โหติ “อิโต เม อิทํ ภวิสฺสติ, อิมินา อิทํ ลจฺฉามี”ติ.\csromanspacer{} เอวเมวิเม โภนฺโต สมณพฺราหฺมณา วาณิชูปมา มญฺเญ ปฏิภนฺติ “อิติ เปจฺจ ภวิสฺสาม, อิติ เปจฺจ ภวิสฺสามา”ติ.\csromanspacer{} ตยิทํ ภิกฺขเว ตถาคโต อภิชานาติ.\csromanspacer{} เย โข เต โภนฺโต สมณพฺราหฺมณา สโต สตฺตสฺส อุจฺเฉทํ วินาสํ วิภวํ ปญฺญเปนฺติ, เต สกฺกายภยา สกฺกายปริเชคุจฺฉา สกฺกายญฺเญว อนุปริธาวนฺติ อนุปริวตฺตนฺติ.\csromanspacer{} เสยฺยถาปิ นาม สา คทฺทุลพทฺโธ ทเฬฺห ถมฺเภ วา ขิเล4 วา อุปนิพทฺโธ ตเมว ถมฺภํ วา ขิลํ วา อนุปริธาวติ \csromanfolio{22}อนุปริวตฺตติ.\csromanspacer{} เอวเมวิเม โภนฺโต สมณพฺราหฺมณา สกฺกายภยา สกฺกายปริเชคุจฺฉา สกฺกายญฺเญว อนุปริธาวนฺติ อนุปริวตฺตนฺติ.\csromanspacer{} ตยิทํ สงฺขตํ โอฬาริกํ, อตฺถิ โข ปน สงฺขารานํ นิโรโธ, อตฺเถตนฺติ อิติ วิทิตฺวา ตสฺส นิสฺสรณทสฺสาวี ตถาคโต ตทุปาติวตฺโต.}

\proseitem{26}{เย หิ เกจิ ภิกฺขเว สมณา วา พฺราหฺมณา วา อปรนฺตกปฺปิกา อปรนฺตานุทิฏฺฐิโน อปรนฺตํ อารพฺภ อเนกวิหิตานิ อธิวุตฺติปทานิ อภิวทนฺติ, สพฺเพ เต อิมาเนว ปญฺจายตนานิ อภิวทนฺติ, เอเตสํ วา อญฺญตรํ.}

\proseitem{27}{สนฺติ ภิกฺขเว เอเก สมณพฺราหฺมณา ปุพฺพนฺตกปฺปิกา ปุพฺพนฺตานุทิฏฺฐิโน ปุพฺพนฺตํ อารพฺภ อเนกวิหิตานิ อธิวุตฺติปทานิ อภิวทนฺติ, “สสฺสโต อตฺตา จ โลโก จ, อิทเมว สจฺจํ โมฆมญฺญนฺ”ติ อิตฺเถเก อภิวทนฺติ, “อสสฺสโต อตฺตา จ โลโก จ, อิทเมว สจฺจํ โมฆมญฺญนฺ”ติ อิตฺเถเก อภิวทนฺติ, “สสฺสโต จ อสสฺสโต จ อตฺตา จ โลโก จ, อิทเมว สจฺจํ โมฆมญฺญนฺ”ติ อิตฺเถเก อภิวทนฺติ, “เนวสสฺสโต นาสสฺสโต อตฺตา จ โลโก จ, อิทเมว สจฺจํ โมฆมญฺญนฺ”ติ อิตฺเถเก อภิวทนฺติ. “อนฺตวา อตฺตา จ โลโก จ, อิทเมว สจฺจํ โมฆมญฺญนฺ”ติ อิตฺเถเก อภิวทนฺติ, “อนนฺตวา อตฺตา จ โลโก จ, อิทเมว สจฺจํ โมฆมญฺญนฺ”ติ อิตฺเถเก อภิวทนฺติ, “อนฺตวา จ อนนฺตวา จ อตฺตา จ โลโก จ, อิทเมว สจฺจํ โมฆมญฺญนฺ”ติ อิตฺเถเก อภิวทนฺติ, “เนวนฺตวา นานนฺตวา อตฺตา จ โลโก จ, อิทเมว สจฺจํ โมฆมญฺญนฺ”ติ อิตฺเถเก อภิวทนฺติ. “เอกตฺตสญฺญี อตฺตา จ โลโก จ, อิทเมว สจฺจํ โมฆมญฺญนฺ”ติ อิตฺเถเก อภิวทนฺติ, “นานตฺตสญฺญี อตฺตา จ โลโก จ, อิทเมว สจฺจํ โมฆมญฺญนฺ”ติ อิตฺเถเก อภิวทนฺติ, “ปริตฺตสญฺญี อตฺตา จ โลโก จ, อิทเมว สจฺจํ โมฆมญฺญนฺ”ติ อิตฺเถเก อภิวทนฺติ, “อปฺปมาณสญฺญี อตฺตา จ โลโก จ, อิทเมว สจฺจํ โมฆมญฺญนฺ”ติ อิตฺเถเก อภิวทนฺติ. “เอกนฺตสุขี อตฺตา จ โลโก จ, อิทเมว สจฺจํ โมฆมญฺญนฺ”ติ อิตฺเถเก อภิวทนฺติ, “เอกนฺตทุกฺขี อตฺตา จ โลโก จ, อิทเมว สจฺจํ โมฆมญฺญนฺ”ติ อิตฺเถเก อภิวทนฺติ, “สุขทุกฺขี อตฺตา จ โลโก จ, อิทเมว สจฺจํ โมฆมญฺญนฺ”ติ อิตฺเถเก อภิวทนฺติ, “อทุกฺขมสุขี อตฺตา จ โลโก จ, อิทเมว สจฺจํ โมฆมญฺญนฺ”ติ อิตฺเถเก อภิวทนฺติ.}

\csromanheader{2}{2. ปญฺจตฺตยสุตฺต (102)}
\proseitem{28}{\csromanfolio{23}ตตฺร ภิกฺขเว เย เต สมณพฺราหฺมณา เอวํวาทิโน เอวํทิฏฺฐิโน “สสฺสโต อตฺตา จ โลโก จ, อิทเมว สจฺจํ โมฆมญฺญนฺ”ติ.\csromanspacer{} เตสํ วต อญฺญเตฺรว สทฺธาย อญฺญตฺร รุจิยา อญฺญตฺร อนุสฺสวา อญฺญตฺร อาการปริวิตกฺกา อญฺญตฺร ทิฏฺฐินิชฺฌานกฺขนฺติยา ปจฺจตฺตํเยว ญาณํ ภวิสฺสติ ปริสุทฺธํ ปริโยทาตนฺติ เนตํ ฐานํ วิชฺชติ.\csromanspacer{} ปจฺจตฺตํ โข ปน ภิกฺขเว ญาเณ อสติ ปริสุทฺเธ ปริโยทาเต ยทปิ\footnote{ยทิปิ (ก)} เต โภนฺโต สมณพฺราหฺมณา ตตฺถ ญาณภาคมตฺตเมว ปริโยทเปนฺติ.\csromanspacer{} ตทปิ เตสํ ภวตํ สมณพฺราหฺมณานํ อุปาทานมกฺขายติ.\csromanspacer{} ตยิทํ สงฺขตํ โอฬาริกํ, อตฺถิ โข ปน สงฺขารานํ นิโรโธ, อตฺเถตนฺติ อิติ วิทิตฺวา ตสฺส นิสฺสรณทสฺสาวี ตถาคโต ตทุปาติวตฺโต.}

\proseitem{29}{ตตฺร ภิกฺขเว เย เต สมณพฺราหฺมณา เอวํวาทิโน เอวํทิฏฺฐิโน “อสสฺสโต อตฺตา จ โลโก จ, อิทเมว สจฺจํ โมฆมญฺญนฺ”ติ ฯเปฯ\footnote{ยถา สสฺสตวาเร, ตถา วิตฺถาเรตพฺพํ.} สสฺสโต จ อสสฺสโต จ อตฺตา จ โลโก จ.\csromanspacer{} เนวสสฺสโต นาสสฺสโต อตฺตา จ โลโก จ.\csromanspacer{} อนฺตวา อตฺตา จ โลโก จ.\csromanspacer{} อนนฺตวา อตฺตา จ โลโก จ.\csromanspacer{} อนฺตวา จ อนนฺตวา จ อตฺตา จ โลโก จ.\csromanspacer{} เนวนฺตวา นานนฺตวา อตฺตา จ โลโก จ.\csromanspacer{} เอกตฺตสญฺญี อตฺตา จ โลโก จ.\csromanspacer{} นานตฺตสญฺญี อตฺตา จ โลโก จ.\csromanspacer{} ปริตฺตสญฺญี อตฺตา จ โลโก จ.\csromanspacer{} อปฺปมาณสญฺญี อตฺตา จ โลโก จ.\csromanspacer{} เอกนฺตสุขี อตฺตา จ โลโก จ.\csromanspacer{} เอกนฺตทุกฺขี อตฺตา จ โลโก จ.\csromanspacer{} สุขทุกฺขี อตฺตา จ โลโก จ. “อทุกฺขมสุขี อตฺตา จ โลโก จ, อิทเมว สจฺจํ โมฆมญฺญนฺ”ติ.\csromanspacer{} เตสํ วต อญฺญเตฺรว สทฺธาย อญฺญตฺร รุจิยา อญฺญตฺร อนุสฺสวา อญฺญตฺร อาการปริวิตกฺกา อญฺญตฺร ทิฏฺฐินิชฺฌานกฺขนฺติยา ปจฺจตฺตํเยว ญาณํ ภวิสฺสติ ปริสุทฺธํ ปริโยทาตนฺติ เนตํ ฐานํ วิชฺชติ.\csromanspacer{} ปจฺจตฺตํ โข ปน ภิกฺขเว ญาเณ อสติ ปริสุทฺเธ ปริโยทาเต ยทปิ เต โภนฺโต สมณพฺราหฺมณา ตตฺถ ญาณภาคมตฺตเมว ปริโยทเปนฺติ.\csromanspacer{} ตทปิ เตสํ ภวตํ สมณพฺราหฺมณานํ อุปาทานมกฺขายติ.\csromanspacer{} ตยิทํ สงฺขตํ โอฬาริกํ, อตฺถิ โข ปน สงฺขารานํ นิโรโธ, อตฺเถตนฺติ อิติ วิทิตฺวา ตสฺส นิสฺสรณทสฺสาวี ตถาคโต ตทุปาติวตฺโต.}

\proseitem{30}{\csromanfolio{24}อิธ ภิกฺขเว เอกจฺโจ สมโณ วา พฺราหฺมโณ วา ปุพฺพนฺตานุทิฏฺฐีนญฺจ ปฏินิสฺสคฺคา อปรนฺตานุทิฏฺฐีนญฺจ ปฏินิสฺสคฺคา สพฺพโส กามสํโยชนานํ อนธิฏฺฐานา ปวิเวกํ ปีติํ อุปสมฺปชฺช วิหรติ “เอตํ สนฺตํ เอตํ ปณีตํ, ยทิทํ ปวิเวกํ ปีติํ อุปสมฺปชฺช วิหรามี”ติ.\csromanspacer{} ตสฺส สา ปวิเวกา ปีติ นิรุชฺฌติ, ปวิเวกาย ปีติยา นิโรธา อุปฺปชฺชติ โทมนสฺสํ, โทมนสฺสสฺส นิโรธา อุปฺปชฺชติ ปวิเวกา ปีติ.\csromanspacer{} เสยฺยถาปิ ภิกฺขเว ยํ ฉายา ชหติ, ตํ อาตโป ผรติ.\csromanspacer{} ยํ อาตโป ชหติ, ตํ ฉายา ผรติ.\csromanspacer{} เอวเมว โข ภิกฺขเว ปวิเวกาย ปีติยา นิโรธา อุปฺปชฺชติ โทมนสฺสํ, โทมนสฺสสฺส นิโรธา อุปฺปชฺชติ ปวิเวกา ปีติ.\csromanspacer{} ตยิทํ ภิกฺขเว ตถาคโต อภิชานาติ.\csromanspacer{} อยํ โข ภวํ สมโณ วา พฺราหฺมโณ วา ปุพฺพนฺตานุทิฏฺฐีนญฺจ ปฏินิสฺสคฺคา อปรนฺตานุทิฏฺฐีนญฺจ ปฏินิสฺสคฺคา สพฺพโส กามสํโยชนานํ อนธิฏฺฐานา ปวิเวกํ ปีติํ อุปสมฺปชฺช วิหรติ “เอตํ สนฺตํ เอตํ ปณีตํ, ยทิทํ ปวิเวกํ ปีติํ อุปสมฺปชฺช วิหรามี”ติ.\csromanspacer{} ตสฺส สา ปวิเวกา ปีติ นิรุชฺฌติ, ปวิเวกาย ปีติยา นิโรธา อุปฺปชฺชติ โทมนสฺสํ, โทมนสฺสสฺส นิโรธา อุปฺปชฺชติ ปวิเวกา ปีติ.\csromanspacer{} ตยิทํ สงฺขตํ โอฬาริกํ, อตฺถิ โข ปน สงฺขารานํ นิโรโธ, อตฺเถตนฺติ อิติ วิทิตฺวา ตสฺส นิสฺสรณทสฺสาวี ตถาคโต ตทุปาติวตฺโต.}

\proseitem{31}{อิธ ปน ภิกฺขเว เอกจฺโจ สมโณ วา พฺราหฺมโณ วา ปุพฺพนฺตานุทิฏฺฐีนญฺจ ปฏินิสฺสคฺคา อปรนฺตานุทิฏฺฐีนญฺจ ปฏินิสฺสคฺคา สพฺพโส กามสํโยชนานํ อนธิฏฺฐานา ปวิเวกาย ปีติยา สมติกฺกมา นิรามิสํ สุขํ อุปสมฺปชฺช วิหรติ “เอตํ สนฺตํ เอตํ ปณีตํ, ยทิทํ นิรามิสํ สุขํ อุปสมฺปชฺช วิหรามี”ติ.\csromanspacer{} ตสฺส ตํ นิรามิสํ สุขํ นิรุชฺฌติ, นิรามิสสฺส สุขสฺส นิโรธา อุปฺปชฺชติ ปวิเวกา ปีติ, ปวิเวกาย ปีติยา นิโรธา อุปฺปชฺชติ นิรามิสํ สุขํ.\csromanspacer{} เสยฺยถาปิ ภิกฺขเว ยํ ฉายา ชหติ, ตํ อาตโป ผรติ.\csromanspacer{} ยํ อาตโป ชหติ, ตํ ฉายา ผรติ.\csromanspacer{} เอวเมว โข ภิกฺขเว นิรามิสสฺส สุขสฺส นิโรธา อุปฺปชฺชติ ปวิเวกา ปีติ, ปวิเวกาย ปีติยา นิโรธา อุปฺปชฺชติ นิรามิสํ สุขํ.\csromanspacer{} ตยิทํ ภิกฺขเว ตถาคโต อภิชานาติ.\csromanspacer{} อยํ โข ภวํ สมโณ วา พฺราหฺมโณ วา ปุพฺพนฺตานุทิฏฺฐีนญฺจ ปฏินิสฺสคฺคา อปรนฺตานุทิฏฺฐีนญฺจ ปฏินิสฺสคฺคา สพฺพโส กามสํโยชนานํ อนธิฏฺฐานา ปวิเวกาย ปีติยา สมติกฺกมา นิรามิสํ สุขํ อุปสมฺปชฺช วิหรติ “เอตํ สนฺตํ เอตํ ปณีตํ, ยทิทํ นิรามิสํ สุขํ อุปสมฺปชฺช วิหรามี”ติ.}

\vspace{\tipitakaparskip}
\csromancenter{ปญฺจตฺตยสุตฺต (102) ตสฺส ตํ นิรามิสํ สุขํ นิรุชฺฌติ, นิรามิสสฺส สุขสฺส นิโรธา อุปฺปชฺชติ ปวิเวกา ปีติ, ปวิเวกาย ปีติยา นิโรธา อุปฺปชฺชติ นิรามิสํ สุขํ.\csromanspacer{} ตยิทํ สงฺขตํ โอฬาริกํ, อตฺถิ โข ปน สงฺขารานํ นิโรโธ, อตฺเถตนฺติ อิติ วิทิตฺวา ตสฺส นิสฺสรณทสฺสาวี ตถาคโต ตทุปาติวตฺโต.}
\proseitem{32}{\csromanfolio{25}อิธ ปน ภิกฺขเว เอกจฺโจ สมโณ วา พฺราหฺมโณ วา ปุพฺพนฺตานุทิฏฺฐีนญฺจ ปฏินิสฺสคฺคา อปรนฺตานุทิฏฺฐีนญฺจ ปฏินิสฺสคฺคา สพฺพโส กามสํโยชนานํ อนธิฏฺฐานา ปวิเวกาย ปีติยา สมติกฺกมา นิรามิสสฺส สุขสฺส สมติกฺกมา อทุกฺขมสุขํ เวทนํ อุปสมฺปชฺช วิหรติ “เอตํ สนฺตํ เอตํ ปณีตํ, ยทิทํ อทุกฺขมสุขํ เวทนํ อุปสมฺปชฺช วิหรามี”ติ.\csromanspacer{} ตสฺส สา อทุกฺขมสุขา เวทนา นิรุชฺฌติ, อทุกฺขมสุขาย เวทนาย นิโรธา อุปฺปชฺชติ นิรามิสํ สุขํ, นิรามิสสฺส สุขสฺส นิโรธา อุปฺปชฺชติ อทุกฺขมสุขา เวทนา.\csromanspacer{} เสยฺยถาปิ ภิกฺขเว ยํ ฉายา ชหติ, ตํ อาตโป ผรติ.\csromanspacer{} ยํ อาตโป ชหติ, ตํ ฉายา ผรติ.\csromanspacer{} เอวเมว โข ภิกฺขเว อทุกฺขมสุขาย เวทนาย นิโรธา อุปฺปชฺชติ นิรามิสํ สุขํ, นิรามิสสฺส สุขสฺส นิโรธา อุปฺปชฺชติ อทุกฺขมสุขา เวทนา.\csromanspacer{} ตยิทํ ภิกฺขเว ตถาคโต อภิชานาติ.\csromanspacer{} อยํ โข ภวํ สมโณ วา พฺราหฺมโณ วา ปุพฺพนฺตนุทิฏฺฐีนญฺจ ปฏินิสฺสคฺคา อปรนฺตานุทิฏฺฐีนญฺจ ปฏินิสฺสคฺคา สพฺพโส กามสํโยชนานํ อนธิฏฺฐานา ปวิเวกาย ปีติยา สมติกฺกมา นิรามิสสฺส สุขสฺส สมติกฺกมา อทุกฺขมสุขํ เวทนํ อุปสมฺปชฺช วิหรติ “เอตํ สนฺตํ เอตํ ปณีตํ, ยทิทํ อทุกฺขมสุขํ เวทนํ อุปสมฺปชฺช วิหรามี”ติ.\csromanspacer{} ตสฺส สา อทุกฺขมสุขา เวทนา นิรุชฺฌติ, อทุกฺขมสุขาย เวทนาย นิโรธา อุปฺปชฺชติ นิรามิสํ สุขํ, นิรามิสสฺส สุขสฺส นิโรธา อุปฺปชฺชติ อทุกฺขมสุขา เวทนา.\csromanspacer{} ตยิทํ สงฺขตํ โอฬาริกํ, อตฺถิ โข ปน สงฺขารานํ นิโรโธ, อตฺเถตนฺติ อิติ วิทิตฺวา ตสฺส นิสฺสรณทสฺสาวี ตถาคโต ตทุปาติวตฺโต.}

\proseitem{33}{อิธ ปน ภิกฺขเว เอกจฺโจ สมโณ วา พฺราหฺมโณ วา ปุพฺพนฺตานุทิฏฺฐีนญฺจ ปฏินิสฺสคฺคา อปรนฺตานุทิฏฺฐีนญฺจ ปฏินิสฺสคฺคา สพฺพโส กามสํโยชนานํ อนธิฏฺฐานา ปวิเวกาย ปีติยา สมติกฺกมา นิรามิสสฺส สุขสฺส สมติกฺกมา อทุกฺขมสุขาย เวทนาย สมติกฺกมา “สนฺโตหมสฺมิ นิพฺพุโตหมสฺมิ อนุปาทาโนหมสฺมี”ติ สมนุปสฺสติ.\csromanspacer{} ตยิทํ ภิกฺขเว ตถาคโต อภิชานาติ.\csromanspacer{} อยํ โข ภวํ สมโณ วา \csromanfolio{26}พฺราหฺมโณ วา ปุพฺพนฺตานุทิฏฺฐีนญฺจ ปฏินิสฺสคฺคา อปรนฺตานุทิฏฺฐีนญฺจ ปฏินิสฺสคฺคา สพฺพโส กามสํโยชนานํ อนธิฏฺฐานา ปวิเวกาย ปีติยา สมติกฺกมา นิรามิสสฺส สุขสฺส สมติกฺกมา อทุกฺขมสุขาย เวทนาย สมติกฺกมา “สนฺโตหมสฺมิ นิพฺพุโตหมสฺมิ อนุปาทาโนหมสฺมี”ติ สมนุปสฺสติ.\csromanspacer{} อทฺธา อยมายสฺมา นิพฺพานสปฺปายํเยว ปฏิปทํ อภิวทติ.\csromanspacer{} อถ จ ปนายํ ภวํ สมโณ วา พฺราหฺมโณ วา ปุพฺพนฺตานุทิฏฺฐิํ วา อุปาทิยมาโน อุปาทิยติ, อปรนฺตานุทิฏฺฐิํ วา อุปาทิยมาโน อุปาทิยติ, กามสํโยชนํ วา อุปาทิยมาโน อุปาทิยติ, ปวิเวกํ วา ปีติํ อุปาทิยมาโน อุปาทิยติ, นิรามิสํ วา สุขํ อุปาทิยมาโน อุปาทิยติ, อทุกฺขมสุขํ วา เวทนํ อุปาทิยมาโน อุปาทิยติ.\csromanspacer{} ยญฺจ โข อยมายสฺมา “สนฺโตหมสฺมิ นิพฺพุโตหมสฺมิ อนุปาทาโนหมสฺมี”ติ สมนุปสฺสติ.\csromanspacer{} ตทปิ อิมสฺส โภโต สมณสฺส พฺราหฺมณสฺส อุปาทานมกฺขายติ.\csromanspacer{} ตยิทํ สงฺขตํ โอฬาริกํ, อตฺถิ โข ปน สงฺขารานํ นิโรโธ, อตฺเถตนฺติ อิติ วิทิตฺวา ตสฺส นิสฺสรณทสฺสาวี ตถาคโต ตทุปาติวตฺโต.}

\prose{อิทํ โข ปน ภิกฺขเว ตถาคเตน อนุตฺตรํ สนฺติวรปทํ อภิสมฺพุทฺธํ, ยทิทํ ฉนฺนํ ผสฺสายตนานํ สมุทยญฺจ อตฺถงฺคมญฺจ อสฺสาทญฺจ อาทีนวญฺจ นิสฺสรณญฺจ ยถาภูตํ วิทิตฺวา1 อนุปาทาวิโมกฺโขติ1.}

\prose{อิทมโวจ ภควา.\csromanspacer{} อตฺตมนา เต ภิกฺขู ภควโต ภาสิตํ อภินนฺทุนฺติ.}

\nitthitamruled{ปญฺจตฺตยสุตฺตํ นิฏฺฐิตํ ทุติยํ.}
\csromanmatikaanchor{mtk.4}
\csromantocmarkref{subsection}{3. กินฺติสุตฺต}{mtk.4}
\bookmark[dest={mtk.4},level=2]{3. กินฺติสุตฺต}
\csromanheader{2}{3. กินฺติสุตฺต}
\proseitem{34}{เอวํ เม สุตํ\csromandash{}เอกํ สมยํ ภควา ปิสินารายํ\footnote{กุสินารายํ (สี)} วิหรติ พลิหรเณ วนสณฺเฑ.\csromanspacer{} ตตฺร โข ภควา ภิกฺขู อามนฺเตสิ “ภิกฺขโว”ติ. “ภทนฺเต”ติ เต ภิกฺขู ภควโต ปจฺจสฺโสสุํ.\csromanspacer{} ภควา เอตทโวจ “กินฺติ}

\vspace{\tipitakaparskip}
\csromancenter{กินฺติสุตฺต (103) โว ภิกฺขเว มยิ โหติ, จีวรเหตุ วา สมโณ โคตโม ธมฺมํ เทเสติ, ปิณฺฑปาตเหตุ วา สมโณ โคตโม ธมฺมํ เทเสติ, เสนาสนเหตุ วา สมโณ โคตโม ธมฺมํ เทเสติ, อิติภวาภวเหตุ วา สมโณ โคตโม ธมฺมํ เทเสตี”ติ.\csromanspacer{} น โข โน ภนฺเต ภควติ เอวํ โหติ “จีวรเหตุ วา สมโณ โคตโม ธมฺมํ เทเสติ, ปิณฺฑปาตเหตุ วา สมโณ โคตโม ธมฺมํ เทเสติ, เสนาสนเหตุ วา สมโณ โคตโม ธมฺมํ เทเสติ, อิติภวาภวเหตุ วา สมโณ โคตโม ธมฺมํ เทเสตี”ติ.}
\prose{\csromanfolio{27}น จ กิร โว ภิกฺขเว มยิ เอวํ โหติ “จีวรเหตุ วา สมโณ โคตโม ธมฺมํ เทเสติ ฯเปฯ อิติภวาภวเหตุ วา สมโณ โคตโม ธมฺมํ เทเสตี”ติ.\csromanspacer{} อถ กินฺติ จรหิ โว\footnote{อถ กินฺติ โว (สี, อิ), อถ กิญฺจรหิ โว (ก)} ภิกฺขเว มยิ โหตีติ.\csromanspacer{} เอวํ โข โน ภนฺเต ภควติ โหติ “อนุกมฺปโก ภควา หิเตสี อนุกมฺปํ อุปาทาย ธมฺมํ เทเสตี”ติ.\csromanspacer{} เอวญฺจ\footnote{เอวํ (สี, อิ)} กิร โว ภิกฺขเว มยิ โหติ “อนุกมฺปโก ภควา หิเตสี อนุกมฺปํ อุปาทาย ธมฺมํ เทเสตี”ติ.}

\proseitem{35}{ตสฺมาติห ภิกฺขเว เย โว\footnote{เย เต (ก)} มยา ธมฺมา อภิญฺญา เทสิตา.\csromanspacer{} เสยฺยถิทํ, จตฺตาโร สติปฏฺฐานา จตฺตาโร สมฺมปฺปธานา จตฺตาโร อิทฺธิปาทา ปญฺจินฺทฺริยานิ ปญฺจ พลานิ สตฺต โพชฺฌงฺคา อริโย อฏฺฐงฺคิโก มคฺโค.\csromanspacer{} ตตฺถ สพฺเพเหว สมคฺเคหิ สมฺโมทมาเนหิ อวิวทมาเนหิ สิกฺขิตพฺพํ.\csromanspacer{} เตสญฺจ โว ภิกฺขเว สมคฺคานํ สมฺโมทมานานํ อวิวทมานานํ สิกฺขตํ สิยํสุ\footnote{สิยุํ (สี, สฺยา, กํ) สทฺทนีติ โอโลเกตพฺพา.} เทฺว ภิกฺขู อภิธมฺเม นานาวาทา, ตตฺร เจ ตุมฺหากํ เอวมสฺส “อิเมสํ โข อายสฺมนฺตานํ อตฺถโต เจว นานํ พฺยญฺชนโต จ นานนฺ”ติ.\csromanspacer{} ตตฺถ ยํ ภิกฺขุํ สุวจตรํ\footnote{สุพฺพจตรํ (ก)} มญฺเญยฺยาถ, โส อุปสงฺกมิตฺวา เอวมสฺส วจนีโย “อายสฺมนฺตานํ โข อตฺถโต เจว นานํ พฺยญฺชนโต จ นานํ, ตทมินาเปตํ\footnote{ตทิมินาเปตํ (สฺยา, กํ)} อายสฺมนฺโต ชานาถ ‘ยถา อตฺถโต เจว นานํ พฺยญฺชนโต จ นานํ, มายสฺมนฺโต วิวาทํ อาปชฺชิตฺถา’ติ”.\csromanspacer{} อถาปเรสํ เอกโตปกฺขิกานํ ภิกฺขูนํ ยํ ภิกฺขุํ สุวจตรํ มญฺเญยฺยาถ, โส อุปสงฺกมิตฺวา เอวมสฺส วจนีโย “อายสฺมนฺตานํ โข อตฺถโต เจว นานํ พฺยญฺชนโต จ นานํ, ตทมินาเปตํ อายสฺมนฺโต ชานาถ ‘ยถา \csromanfolio{28}อตฺถโต เจว นานํ พฺยญฺชนโต จ นานํ, มายสฺมนฺโต วิวาทํ อาปชฺชิตฺถา’ติ”.\csromanspacer{} อิติ ทุคฺคหิตํ ทุคฺคหิตโต ธาเรตพฺพํ, สุคฺคหิตํ สุคฺคหิตโต ธาเรตพฺพํ.\csromanspacer{} ทุคฺคหิตํ ทุคฺคหิตโต ธาเรตฺวา สุคฺคหิตํ สุคฺคหิตโต ธาเรตฺวา\footnote{อิติ ทุคฺคหิตํ ทุคฺคหิตโต ธาเรตพฺพํ, ทุคฺคหิตํ ทุคฺคหิตโต ธาเรตฺวา (สี, สฺยา, กํ, อิ) อนนฺตรวารตฺตเย ปน อิทํ ปาฐนานตฺตํ นตฺติ.} โย ธมฺโม โย วินโย โส ภาสิตพฺโพ.}

\proseitem{36}{ตตฺร เจ ตุมฺหากํ เอวมสฺส “อิเมสํ โข อายสฺมนฺตานํ อตฺถโต หิ โข นานํ พฺยญฺชนโต สเมตี”ติ.\csromanspacer{} ตตฺถ ยํ ภิกฺขุํ สุวจตรํ มญฺเญยฺยาถ, โส อุปสงฺกมิตฺวา เอวมสฺส วจนีโย “อายสฺมนฺตานํ โข อตฺถโต หิ นานํ พฺยญฺชนโต สเมติ, ตทมินาเปตํ อายสฺมนฺโต ชานาถ ‘ยถา อตฺถโต หิ โข นานํ พฺยญฺชนโต สเมติ, มายสฺมนฺโต วิวาทํ อาปชฺชิตฺถา’ติ”.\csromanspacer{} อถาปเรสํ เอกโตปกฺขิกานํ ภิกฺขูนํ ยํ ภิกฺขุํ สุวจตรํ มญฺเญยฺยาถ, โส อุปสงฺกมิตฺวา เอวมสฺส วจนีโย “อายสฺมนฺตานํ โข อตฺถโต หิ โข นานํ พฺยญฺชนโต สเมติ, ตทมินาเปตํ อายสฺมนฺโต ชานาถ ‘ยถา อตฺถโต หิ โข นานํ พฺยญฺชนโต สเมติ, มายสฺมนฺโต วิวาทํ อาปชฺชิตฺถา’ติ”.\csromanspacer{} อิติ ทุคฺคหิตํ ทุคฺคหิตโต ธาเรตพฺพํ, สุคฺคหิตํ สุคฺคหิตโต ธาเรตพฺพํ.\csromanspacer{} ทุคฺคหิตํ ทุคฺคหิตโต ธาเรตฺวา สุคฺคหิตํ สุคฺคหิตโต ธาเรตฺวา โย ธมฺโม โย วินโย โส ภาสิตพฺโพ.}

\proseitem{37}{ตตฺร เจ ตุมฺหากํ เอวมสฺส “อิเมสํ โข อายสฺมนฺตานํ อตฺถโต หิ โข สเมติ พฺยญฺชนโต นานนฺ”ติ.\csromanspacer{} ตตฺถ ยํ ภิกฺขุํ สุวจตรํ มญฺเญยฺยาถ, โส อุปสงฺกมิตฺวา เอวมสฺส วจนีโย “อายสฺมนฺตานํ โข อตฺถโต หิ สเมติ พฺยญฺชนโต นานํ, ตทมินาเปตํ อายสฺมนฺโต ชานาถ ‘ยถา อตฺถโต หิ โข สเมติ พฺยญฺชนโต นานํ, อปฺปมตฺตกํ โข ปเนตํ ยทิทํ พฺยญฺชนํ, มายสฺมนฺโต อปฺปมตฺตเก วิวาทํ อาปชฺชิตฺถา’ติ”.\csromanspacer{} อถาปเรสํ เอกโตปกฺขิกานํ ภิกฺขูนํ ยํ ภิกฺขุํ สุวจตรํ มญฺเญยฺยาถ, โส อุปสงฺกมิตฺวา เอวมสฺส วจนีโย “อายสฺมนฺตานํ โข อตฺถโต หิ สเมติ พฺยญฺชนโต นานํ, ตทมินาเปตํ อายสฺมนฺโต ชานาถ ‘ยถา อตฺถโต หิ โข สเมติ พฺยญฺชนโต นานํ, อปฺปมตฺตกํ โข ปเนตํ}

\vspace{\tipitakaparskip}
\csromancenter{กินฺติสุตฺต (103) ยทิทํ พฺยญฺชนํ, มายสฺมนฺโต อปฺปมตฺตเก\footnote{อปฺปมตฺตเกหิ (สี, อิ)} วิวาทํ อาปชฺชิตฺถา’ติ”.\csromanspacer{} อิติ สุคฺคหิตํ สุคฺคหิตโต ธาเรตพฺพํ, ทุคฺคหิตํ ทุคฺคหิตโต ธาเรตพฺพํ.\csromanspacer{} สุคฺคหิตํ สุคฺคหิตโต ธาเรตฺวา ทุคฺคหิตํ ทุคฺคหิตโต ธาเรตฺวา โย ธมฺโม โย วินโย โส ภาสิตพฺโพ.}
\proseitem{38}{\csromanfolio{29}ตตฺร เจ ตุมฺหากํ เอวมสฺส “อิเมสํ โข อายสฺมนฺตานํ อตฺถโต เจว สเมติ พฺยญฺชนโต จ สเมตี”ติ.\csromanspacer{} ตตฺถ ยํ ภิกฺขุํ สุวจตรํ มญฺเญยฺยาถ, โส อุปสงฺกมิตฺวา เอวมสฺส วจนีโย “อายสฺมนฺตานํ โข อตฺถโต เจว สเมติ พฺยญฺชนโต จ สเมติ ตทมินาเปตํ อายสฺมนฺโต ชานาถ ‘ยถา อตฺถโต เจว สเมติ พฺยญฺชนโต จ สเมติ, มายสฺมนฺโต วิวาทํ อาปชฺชิตฺถา’ติ”.\csromanspacer{} อถาปเรสํ เอกโตปกฺขิกานํ ภิกฺขูนํ ยํ ภิกฺขุํ สุวจตรํ มญฺเญยฺยาถ, โส อุปสงฺกมิตฺวา เอวมสฺส วจนีโย “อายสฺมนฺตานํ โข อตฺถโต เจว สเมติ พฺยญฺชนโต จ สเมติ, ตทมินาเปตํ อายสฺมนฺโต ชานาถ ‘ยถา อตฺถโต เจว สเมติ พฺยญฺชนโต จ สเมติ, มายสฺมนฺโต วิวาทํ อาปชฺชิตฺถา’ติ”.\csromanspacer{} อิติ สุคฺคหิตํ สุคฺคหิตโต ธาเรตพฺพํ, สุคฺคหิตํ สุคฺคหิตโต ธาเรตฺวา โย ธมฺโม โย วินโย โส ภาสิตพฺโพ.}

\proseitem{39}{เตสญฺจ โว ภิกฺขเว สมคฺคานํ สมฺโมทมานานํ อวิวทมานานํ สิกฺขตํ สิยา อญฺญตรสฺส ภิกฺขุโน อาปตฺติ, สิยา วีติกฺกโม.\csromanspacer{} ตตฺร ภิกฺขเว น โจทนาย ตริตพฺพํ\footnote{โจทิตพฺพํ (สฺยา, กํ, ก) ตุริตพฺพํ (?)} ปุคฺคโล อุปปริกฺขิตพฺโพ “อิติ มยฺหญฺจ อวิเหสา ภวิสฺสติ ปรสฺส จ ปุคฺคลสฺส อนุปฆาโต, ปโร หิ ปุคฺคโล อกฺโกธโน อนุปนาหี อทฬฺหทิฏฺฐี สุปฺปฏินิสฺสคฺคี, สกฺโกมิ จาหํ เอตํ ปุคฺคลํ อกุสลา วุฏฺฐาเปตฺวา กุสเล ปติฏฺฐาเปตุนฺ”ติ.\csromanspacer{} สเจ ภิกฺขเว เอวมสฺส, กลฺลํ วจนาย.}

\prose{สเจ ปน ภิกฺขเว เอวมสฺส “มยฺหํ โข อวิเหสา ภวิสฺสติ ปรสฺส จ ปุคฺคลสฺส อุปฆาโต, ปโร หิ ปุคฺคโล โกธโน อุปนาหี อทฬฺหทิฏฺฐี สุปฺปฏินิสฺสคฺคี, สกฺโกมิ จาหํ เอตํ ปุคฺคลํ อกุสลา วุฏฺฐาเปตฺวา กุสเล ปติฏฺฐาเปตุํ.\csromanspacer{} อปฺปมตฺตกํ โข ปเนตํ ยทิทํ ปรสฺส\footnote{ยทิทํ มยฺหญฺจ วิเหสา ภวิสฺสติ ปรสฺส จ (ก)} ปุคฺคลสฺส อุปฆาโต.\csromanspacer{} อถ โข เอตเทว พหุตรํ, สฺวาหํ สกฺโกมิ เอตํ ปุคฺคลํ อกุสลา \csromanfolio{30}วุฏฺฐาเปตฺวา กุสเล ปติฏฺฐาเปตุนฺ”ติ.\csromanspacer{} สเจ ภิกฺขเว เอวมสฺส กลฺลํ วจนาย.}

\prose{สเจ ปน ภิกฺขเว เอวมสฺส “มยฺหํ โข วิเหสา ภวิสฺสติ ปรสฺส จ ปุคฺคลสฺส อนุปฆาโต.\csromanspacer{} ปโร หิ ปุคฺคโล อกฺโกธโน อนุปนาหี ทฬฺหทิฏฺฐี ทุปฺปฏินิสฺสคฺคี, สกฺโกมิ จาหํ เอตํ ปุคฺคลํ อกุสลา วุฏฺฐาเปตฺวา กุสเล ปติฏฺฐาเปตุํ.\csromanspacer{} อปฺปมตฺตกํ โข ปเนตํ ยทิทํ มยฺหํ วิเหสา\footnote{มยฺหญฺจ วิเหสา ภวิสฺสติ ปรสฺส จ ปุคฺคลสฺส อุปฆาโต (ก)}.\csromanspacer{} อถ โข เอตเทว พหุตรํ, สฺวาหํ สกฺโกมิ เอตํ ปุคฺคลํ อกุสลา วุฏฺฐาเปตฺวา กุสเล ปติฏฺฐาเปตุนฺ”ติ.\csromanspacer{} สเจ ภิกฺขเว เอวมสฺส กลฺลํ วจนาย.}

\prose{สเจ ปน ภิกฺขเว เอวมสฺส “มยฺหญฺจ โข วิเหสา ภวิสฺสติ ปรสฺส จ ปุคฺคลสฺส อุปฆาโต.\csromanspacer{} ปโร หิ ปุคฺคโล โกธโน อุปนาหี ทฬฺหทิฏฺฐี ทุปฺปฏินิสฺสคฺคี, สกฺโกมิ จาหํ เอตํ ปุคฺคลํ อกุสลา วุฏฺฐาเปตฺวา กุสเล ปติฏฺฐาเปตุํ.\csromanspacer{} อปฺปมตฺตกํ โข ปเนตํ ยทิทํ มยฺหญฺจ วิเหสา ภวิสฺสติ ปรสฺส จ ปุคฺคลสฺส อุปฆาโต.\csromanspacer{} อถ โข เอตเทว พหุตรํ, สฺวาหํ สกฺโกมิ เอตํ ปุคฺคลํ อกุสลา วุฏฺฐาเปตฺวา กุสเล ปติฏฺฐาเปตุนฺ”ติ.\csromanspacer{} สเจ ภิกฺขเว เอวมสฺส กลฺลํ วจนาย.}

\prose{สเจ ปน ภิกฺขเว เอวมสฺส “มยฺหญฺจ โข วิเหสา ภวิสฺสติ ปรสฺส จ ปุคฺคลสฺส อุปฆาโต.\csromanspacer{} ปโร หิ ปุคฺคโล โกธโน อุปนาหี ทฬฺหทิฏฺฐี ทุปฺปฏินิสฺสคฺคี, น จาหํ สกฺโกมิ เอตํ ปุคฺคลํ อกุสลา วุฏฺฐาเปตฺวา กุสเล ปติฏฺฐาเปตุนฺ”ติ.\csromanspacer{} เอวรูเป ภิกฺขเว ปุคฺคเล อุเปกฺขา นาติมญฺญิตพฺพา.}

\proseitem{40}{เตสญฺจ โว ภิกฺขเว สมคฺคานํ สมฺโมทมานํ อวิวทมานานํ สิกฺขตํ อญฺญมญฺญสฺส วจีสํหาโร\footnote{วจีสงฺขาโร (สี, อิ)} อุปฺปชฺเชยฺย ทิฏฺฐิปฬาโส\footnote{ทิฏฺฐิปลาโส (สี, ก)} เจตโส อาฆาโต อปฺปจฺจโย อนภิรทฺธิ.\csromanspacer{} ตตฺถ เอกโตปกฺขิกานํ ภิกฺขูนํ ยํ ภิกฺขุํ สุวจตรํ มญฺเญยฺยาถ, โส อุปสงฺกมิตฺวา เอวมสฺส วจนีโย “ยํ โน อาวุโส อมฺหากํ สมคฺคานํ สมฺโมทมานานํ อวิวทมานานํ สิกฺขตํ อญฺญมญฺญสฺส วจีสํหาโร อุปฺปนฺโน ทิฏฺฐิปฬาโส เจตโส อาฆาโต อปฺปจฺจโย อนภิรทฺธิ, ตํ ชานมาโน สมโณ\footnote{สมาโน (สี, ก)} ครเหยฺยา”ติ? สมฺมาพฺยากรมาโน ภิกฺขเว เอวํ พฺยากเรยฺย “ยํ โน อาวุโส}

\vspace{\tipitakaparskip}
\csromancenter{กินฺติสุตฺต (103) อมฺหากํ สมคฺคานํ สมฺโมทมานานํ อวิวทมานานํ สิกฺขตํ อญฺญมญฺญสฺส วจีสํหาโร อุปฺปนฺโน ทิฏฺฐิปฬาโส เจตโส อาฆาโต อปฺปจฺจโย อนภิรทฺธิ, ตํ ชานมาโน สมโณ ครเหยฺยา”ติ.\csromanspacer{} เอตํ ปนาวุโส ธมฺมํ อปฺปหาย นิพฺพานํ สจฺฉิกเรยฺยาติ? สมฺมา พฺยากรมาโน ภิกฺขเว ภิกฺขุ เอวํ พฺยากเรยฺย “เอตํ อาวุโส ธมฺมํ อปฺปหาย น นิพฺพานํ สจฺฉิกเรยฺยา”ติ.}
\prose{\csromanfolio{31}อถาปเรสํ เอกโตปกฺขิกานํ ภิกฺขูนํ ยํ ภิกฺขุํ สุวจตรํ มญฺเญยฺยาถ, โส อุปสงฺกมิตฺวา เอวมสฺส วจนีโย “ยํ โน อาวุโส อมฺหากํ สมคฺคานํ สมฺโมทมานานํ อวิวทมานานํ สิกฺขตํ อญฺญมญฺญสฺส วจีสํหาโร อุปฺปนฺโน ทิฏฺฐิปฬาโส เจตโส อาฆาโต อปฺปจฺจโย อนภิรทฺธิ, ตํ ชานมาโน สมโณ ครเหยฺยา”ติ? สมฺมา พฺยากรมาโน ภิกฺขเว ภิกฺขุ เอวํ พฺยากเรยฺย “ยํ โน อาวุโส อมฺหากํ สมคฺคานํ สมฺโมทมานานํ อวิวทมานานํ สิกฺขตํ อญฺญมญฺญสฺส วจีสํหาโร อุปฺปนฺโน ทิฏฺฐิปฬาโส เจตโส อาฆาโต อปฺปจฺจโย อนภิรทฺธิ, ตํ ชานมาโน สมโณ ครเหยฺยา”ติ.\csromanspacer{} เอตํ ปนาวุโส ธมฺมํ อปฺปหาย นิพฺพานํ สจฺฉิกเรยฺยาติ? สมฺมา พฺยากรมาโน ภิกฺขเว ภิกฺขุ เอวํ พฺยากเรยฺย “เอตํ โข อาวุโส ธมฺมํ อปฺปหาย น นิพฺพานํ สจฺฉิกเรยฺยา”ติ.}

\prose{ตํ เจ ภิกฺขเว ภิกฺขุํ ปเร เอวํ ปุจฺเฉยฺยุํ “อายสฺมตา โน เอเต ภิกฺขู อกุสลา วุฏฺฐาเปตฺวา กุสเล ปติฏฺฐาปิตา”ติ.\csromanspacer{} สมฺมา พฺยากรมาโน ภิกฺขเว ภิกฺขุ เอวํ พฺยากเรยฺย “อิธาหํ อาวุโส เยน ภควา เตนุปสงฺกมิํ, ตสฺส เม ภควา ธมฺมํ เทเสสิ, ตาหํ ธมฺมํ สุตฺวา เตสํ ภิกฺขูนํ อภาสิํ, ตํ เต ภิกฺขู ธมฺมํ สุตฺวา อกุสลา วุฏฺฐหิํสุ กุสเล ปติฏฺฐหิํสู”ติ เอวํ พฺยากรมาโน โข ภิกฺขเว ภิกฺขุ น เจว อตฺตานํ อุกฺกํเสติ, น ปรํ วมฺเภติ, ธมฺมสฺส จานุธมฺมํ พฺยากโรติ, น จ โกจิ สหธมฺมิโก วาทานุวาโท คารยฺหํ ฐานํ อาคจฺฉตีติ.}

\prose{อิทมโวจ ภควา.\csromanspacer{} อตฺตมนา เต ภิกฺขู ภควโต ภาสิตํ อภินนฺทุนฺติ.}

\nitthitam{กินฺติ สุตฺตํ นิฏฺฐิตํ ตติยํ.}
\csromanmatikaanchor{mtk.5}
\csromantocmarkref{subsection}{4. สามคามสุตฺต}{mtk.5}
\bookmark[dest={mtk.5},level=2]{4. สามคามสุตฺต}
\csromanheader{2}{4. สามคามสุตฺต}
\proseitem{41}{\csromanfolio{32}เอวํ เม สุตํ\csromandash{}เอกํ สมยํ ภควา สกฺเกสุ วิหรติ สามคาเม.\csromanspacer{} เตน โข ปน สมเยน นิคณฺโฐ นาฏปุตฺโต\footnote{นาถปุตฺโต (สี, อิ)} ปาวายํ อธุนากาลงฺกโต\footnote{กาลกโต (สี, สฺยา, กํ, อิ)} โหติ, ตสฺส กาลงฺกิริยาย ภินฺนา นิคณฺฐา เทฺวธิกชาตา\footnote{เทฺวฬฺหกชาตา (สฺยา, กํ, ก)} ภณฺฑนชาตา กลหชาตา วิวาทาปนฺนา อญฺญมญฺญํ มุขสตฺตีหิ วิตุทนฺตา วิหรนฺติ “น ตฺวํ อิมํ ธมฺมวินยํ อาชานาสิ, อหํ อิมํ ธมฺมวินยํ อาชานามิ.\csromanspacer{} กิํ ตฺวํ อิมํ ธมฺมวินยํ อาชานิสฺสสิ, มิจฺฉาปฏิปนฺโน ตฺวมสิ, อหมสฺมิ สมฺมาปฏิปนฺโน.\csromanspacer{} สหิตํ เม, อสหิตํ เต.\csromanspacer{} ปุเรวจนียํ ปจฺฉา อวจ, ปจฺฉาวจนียํ ปุเร อวจ.\csromanspacer{} อธิจิณฺณํ\footnote{อวิจิณฺณํ (สี, อิ)} เต วิปราวตฺตํ, อาโรปิโต เต วาโท, นิคฺคหิโตสิ, จร วาทปฺปโมกฺขาย.\csromanspacer{} นิพฺเพเฐหิ วา สเจ ปโหสี”ติ.\csromanspacer{} วโธเยว โข\footnote{วโธเยเวโก (สฺยา, กํ, ก)} มญฺเญ นิคณฺเฐสุ นาฏปุตฺติเยสุ นิคณฺเฐสุ วตฺตติ.\csromanspacer{} เยปิ นิคณฺฐสฺส นาฏปุตฺตสฺส สาวกา คิหี โอทาตวสนา, เตปิ นิคณฺเฐสุ นาฏปุตฺติเยสุ นิพฺพินฺนรูปา\footnote{นิพฺพินฺทรูปา (สฺยา, กํ, ก)} วิรตฺตรูปา ปฏิวานรูปา, ยถา ตํ ทุรกฺขาเต ธมฺมวินเย ทุปฺปเวทิเต อนิยฺยานิเก อนุปสมสํวตฺตนิเก อสมฺมาสมฺพุทฺธปฺปเวทิเต ภินฺนถูเป อปฺปฏิสรเณ.}

\proseitem{42}{อถ โข จุนฺโท สมณุทฺเทโส ปาวายํ วสฺสํวุฏฺโฐ\footnote{วสฺสํวุตฺโถ (สี, สฺยา, กํ, อิ)} เยน สามคาโม, เยนายสฺมา อานนฺโท เตนุปสงฺกมิ, อุปสงฺกมิตฺวา อายสฺมนฺตํ อานนฺทํ อภิวาเทตฺวา เอกมนฺตํ นิสีทิ, เอกมนฺตํ นิสินฺโน โข จุนฺโท สมณุทฺเทโส อายสฺมนฺตํ อานนฺทํ เอตทโวจ “นิคณฺโฐ ภนฺเต นาฏปุตฺโต ปาวายํ อธุนากาลงฺกโต, ตสฺส กาลงฺกิริยาย ภินฺนา นิคณฺฐา เทฺวธิกชาตา ฯเปฯ ภินฺนถูเป อปฺปฏิสรเณ”ติ.\csromanspacer{} เอวํ วุตฺเต อายสฺมา อานนฺโท จุนฺทํ สมณุทฺเทสํ เอตทโวจ “อตฺถิ โข อิทํ อาวุโส จุนฺท กถาปาภตํ ภควนฺตํ ทสฺสนาย, อายาม อาวุโส จุนฺท, เยน ภควา เตนุปสงฺกมิสฺสาม, อุปสงฺกมิตฺวา เอตมตฺถํ ภควโต อาโรเจสฺสามา”ติ. “เอวํ ภนฺเต”ติ โข จุนฺโท สมณุทฺเทโส อายสฺมโต อานนฺทสฺส ปจฺจสฺโสสิ.}

\csromanheader{2}{4. สามคามสุตฺต (104)}
\prose{\csromanfolio{33}อถ โข อายสฺมา จ อานนฺโท จุนฺโท จ สมณุทฺเทโส เยน ภควา เตนุปสงฺกมิํสุ, อุปสงฺกมิตฺวา ภควนฺตํ อภิวาเทตฺวา เอกมนฺตํ นิสีทิํสุ, เอกมนฺตํ นิสินฺโน โข อายสฺมา อานนฺโท ภควนฺตํ เอตทโวจ “อยํ ภนฺเต จุนฺโท สมณุทฺเทโส เอวมาห ‘นิคณฺโฐ ภนฺเต นาฏปุตฺโตปาวายํ อธุนากาลงฺกโต, ตสฺส กาลงฺกิริยาย ภินฺนา นิคณฺฐาเทฺวธิกชาตา ฯเปฯ ภินฺนถูเป อปฺปฏิสรเณ’ติ.\csromanspacer{} ตสฺส มยฺหํ ภนฺเต เอวํ โหติ ‘มาเหว ภควโต อจฺจเยน สํเฆ วิวาโท อุปฺปชฺชิ, สฺวาสฺส\footnote{โส (สี, อิ), สฺวายํ (ก)} วิวาโท พหุชนาหิตาย พหุชนาสุขาย พหุโน ชนสฺส อนตฺถาย อหิตาย ทุกฺขาย เทวมนุสฺสานนฺ’ติ”.}

\proseitem{43}{ตํ กิํ มญฺญสิ อานนฺท, เย โว มยา ธมฺมา อภิญฺญา เทสิตา.\csromanspacer{} เสยฺยถิทํ, จตฺตาโร สติปฏฺฐานา จตฺตาโร สมฺมปฺปธานา จตฺตาโร อิทฺธิปาทา ปญฺจินฺทฺริยานิ ปญฺจ พลานิ สตฺต โพชฺฌงฺคา อริโย อฏฺฐงฺคิโก มคฺโค.\csromanspacer{} ปสฺสสิ โน ตฺวํ อานนฺท อิเมสุ ธมฺเมสุ เทฺวปิ ภิกฺขู นานาวาเทติ.\csromanspacer{} เย เม ภนฺเต ธมฺมา ภควตา อภิญฺญา เทสิตา.\csromanspacer{} เสยฺยถิทํ, จตฺตาโร สติปฏฺฐานา จตฺตาโร สมฺมปฺปธานา จตฺตาโร อิทฺธิปาทา ปญฺจินฺทฺริยานิ ปญฺจ พลานิ สตฺต โพชฺฌงฺคา อริโย อฏฺฐงฺคิโก มคฺโค.\csromanspacer{} นาหํ ปสฺสามิ อิเมสุ ธมฺเมสุ เทฺวปิ ภิกฺขู นานาวาเท.\csromanspacer{} เย จ โข\footnote{สนฺติ จ โข (สฺยา,กํ), สนฺติ จ (ก)} ภนฺเต ปุคฺคลา ภควนฺตํ ปติสฺสยมานรูปา วิหรนฺติ, เตปิ ภควโต อจฺจเยน สํเฆ วิวาทํ ชเนยฺยุํ อชฺฌาชีเว วา อธิปาติโมกฺเข วา.\csromanspacer{} สฺวาสฺส\footnote{โสสฺส (สี, อิ), สฺวายํ (ก)} วิวาโท พหุชนาหิตาย พหุชนาสุขาย พหุโน ชนสฺส อนตฺถาย อหิตาย ทุกฺขาย เทวมนุสฺสานนฺติ.\csromanspacer{} อปฺปมตฺตโก โส อานนฺท วิวาโท, ยทิทํ อชฺฌาชีเว วา อธิปาติโมกฺเข วา.\csromanspacer{} มคฺเค วา หิ อานนฺท ปฏิปทาย วา สํเฆ วิวาโท อุปฺปชฺชมาโน อุปฺปชฺเชยฺย, สฺวาสฺส วิวาโท พหุชนาหิตาย พหุชนาสุขาย พหุโน ชนสฺส อนตฺถาย อหิตาย ทุกฺขาย เทวมนุสฺสานํ.}

\proseitem{44}{ฉยิมานิ อานนฺท วิวาทมูลานิ.\csromanspacer{} กตมานิ ฉ, อิธานนฺท ภิกฺขุ โกธโน โหติ อุปนาหี.\csromanspacer{} โย โส อานนฺท ภิกฺขุ โกธโน โหติ อุปนาหี, โส สตฺถริปิ อคารโว วิหรติ อปฺปติสฺโส, \csromanfolio{34}ธมฺเมปิ อคารโว วิหรติ อปฺปติสฺโส, สํเฆปิ อคารโว วิหรติ อปฺปติสฺโส, สิกฺขายปิ น ปริปูรการี โหติ, โย โส อานนฺท ภิกฺขุ สตฺถริ อคารโว วิหรติ อปฺปติสฺโส, ธมฺเม.\csromanspacer{} สํเฆ อคารโว วิหรติ อปฺปติสฺโส, สิกฺขาย น ปริปูรการี โหติ, โส สํเฆ วิวาทํ ชเนติ, โย โหติ วิวาโท พหุชนาหิตาย พหุชนาสุขาย พหุโน ชนสฺส อนตฺถาย อหิตาย ทุกฺขาย เทวมนุสฺสานํ, เอวรูปญฺเจ ตุเมฺห อานนฺท วิวาทมูลํ อชฺฌตฺตํ วา พหิทฺธา วา สมนุปสฺเสยฺยาถ, ตตฺร ตุเมฺห อานนฺท ตสฺเสว ปาปกสฺส วิวาทมูลสฺส ปหานาย วายเมยฺยาถ, เอวรูปญฺเจ ตุเมฺห อานนฺท วิวาทมูลํ อชฺฌตฺตํ วา พหิทฺธา วา น สมนุปสฺเสยฺยาถ, ตตฺร ตุเมฺห อานนฺท ตสฺเสว ปาปกสฺส วิวาทมูลสฺส อายติํ อนวสฺสวาย ปฏิปชฺเชยฺยาถ, เอวเมตสฺส ปาปกสฺส วิวาทมูลสฺส ปหานํ โหติ, เอวเมตสฺส ปาปกสฺส วิวาทมูลสฺส อายติํ อนวสฺสโว โหติ.}

\proseitem{45}{ปุน จปรํ อานนฺท ภิกฺขุ มกฺขี โหติ ปฬาสี ฯเปฯ อิสฺสุกี โหติ มจฺฉรี ฯเปฯ สโฐ โหติ มายาวี ฯเปฯ ปาปิจฺโฉ โหติ มิจฺฉาทิฏฺฐิ\footnote{มิจฺฉาทิฏฺฐี (สฺยา, กํ, อิ, ก)} ฯเปฯ สนฺทิฏฺฐิปรามาสี โหติ อาธานคฺคาหี ทุปฺปฏินิสฺสคฺคี, โย โส อานนฺท ภิกฺขุ สนฺทิฏฺฐิปรามาสี โหติ อาธานคฺคาหี ทุปฺปฏินิสฺสคฺคี, โส สตฺถริปิ อคารโว วิหรติ อปฺปติสฺโส, ธมฺเมปิ อคารโว วิหรติ อปฺปติสฺโส, สํเฆปิ อคารโว วิหรติ อปฺปติสฺโส, สิกฺขายปิ น ปริปูรการี โหติ, โย โส อานนฺท ภิกฺขุ สตฺถริ อคารโว วิหรติ อปฺปติสฺโส, ธมฺเม.\csromanspacer{} สํเฆ.\csromanspacer{} สิกฺขาย น ปริปูรการี โหติ, โส สํเฆ วิวาทํ ชเนติ, โย โหติ วิวาโท พหุชนาหิตาย พหุชนาสุขาย พหุโน ชนสฺส อนตฺถาย อหิตาย ทุกฺขาย เทวมนุสฺสานํ, เอวรูปญฺเจ ตุเมฺห อานนฺท วิวาทมูลํ อชฺฌตฺตํ วา พหิทฺธา สมนุปสฺเสยฺยาถ, ตตฺร ตุเมฺห อานนฺท ตสฺเสว ปาปกสฺส วิวาทมูลสฺส ปหานาย วายเมยฺยาถ, เอวรูปญฺเจ ตุเมฺห อานนฺท วิวาทมูลํ อชฺฌตฺตํ วา พหิทฺธา วา น สมนุปสฺเสยฺยาถ, ตตฺร ตุเมฺห อานนฺท ตสฺเสว ปาปกสฺส วิวาทมูลสฺส อายติํ อนวสฺสวาย ปฏิปชฺเชยฺยาถ, เอวเมตสฺส ปาปกสฺส วิวาทมูลสฺส ปหานํ โหติ, เอวเมตสฺส ปาปกสฺส วิวาทมูลสฺส อายติํ อนวสฺสโว โหติ.\csromanspacer{} อิมานิ โข อานนฺท ฉ วิวาทมูลานิ.}

\csromanheader{2}{4. สามคามสุตฺต (104)}
\proseitem{46}{\csromanfolio{35}จตฺตาริมานิ อานนฺท อธิกรณานิ.\csromanspacer{} กตมานิ จตฺตาริ, วิวาทาธิกรณํ อนุวาทาธิกรณํ อาปตฺตาธิกรณํ กิจฺจาธิกรณํ.\csromanspacer{} อิมานิ โข อานนฺท จตฺตาริ อธิกรณานิ.\csromanspacer{} สตฺต โข ปนิเม อานนฺท อธิกรณสมถา.\csromanspacer{} อุปฺปนฺนุปฺปนฺนานํ อธิกรณานํ สมถาย วูปสมาย สมฺมุขาวินโย ทาตพฺโพ, สติวินโย ทาตพฺโพ, อมูฬฺหวินโย ทาตพฺโพ, ปฏิญฺญาย กาเรตพฺพํ, เยภุยฺยสิกา, ตสฺสปาปิยสิกา, ติณวตฺถารโก.}

\proseitem{47}{กถญฺจานนฺท สมฺมุขาวินโย โหติ, อิธานนฺท ภิกฺขู วิวทนฺติ ธมฺโมติ วา อธมฺโมติ วา วินโยติ วา อวินโยติ วา, เตหานนฺท ภิกฺขูหิ สพฺเพเหว สมคฺเคหิ สนฺนิปติตพฺพํ, สนฺนิปติตฺวา ธมฺมเนตฺติ สมนุมชฺชิตพฺพา, ธมฺมเนตฺตํ สมนุชฺชิตฺวา ยถา ตตฺถ สเมติ, ตถา ตํ อธิกรณํ วูปสเมตพฺพํ.\csromanspacer{} เอวํ โข อานนฺท สมฺมุขาวินโย โหติ.\csromanspacer{} เอวญฺจ ปนิเธกจฺจานํ อธิกรณานํ วูปสโม โหติ, ยทิทํ สมฺมุขาวินเยน.}

\proseitem{48}{กถญฺจานนฺท เยภุยฺยสิกา โหติ, เต เจ อานนฺท ภิกฺขู น สกฺโกนฺติ ตํ อธิกรณํ ตสฺมิํ อาวาเส วูปสเมตุํ, เตหานนฺท ภิกฺขูหิ ยสฺมิํ อาวาเส พหุตรา ภิกฺขู, โส อาวาโส คนฺตพฺโพ.\csromanspacer{} ตตฺถ สพฺเพเหว สมคฺเคหิ สนฺนิปติตพฺพํ, สนฺนิปติตฺวา ธมฺมเนตฺติ สมนุมชฺชิตพฺพา, ธมฺมเนตฺติ สมนุมชฺชิตฺวา ยถา ตตฺถ สเมติ, ตถา ตํ อธิกรณํ วูปสเมตพฺพํ.\csromanspacer{} เอวํ โข อานนฺท เยภุยฺยสิกา โหติ.\csromanspacer{} เอวญฺจ ปนิเธกจฺจานํ อธิกรณานํ วูปสโม โหติ, ยทิทํ เยภุยฺยสิกาย.}

\proseitem{49}{กถญฺจานนฺท สติวินโย โหติ, อิธานนฺท ภิกฺขู ภิกฺขุํ เอวรูปาย ครุกาย อาปตฺติยา โจเทนฺติ ปาราชิเกน วา ปาราชิกสามนฺเตน วา “สรตายสฺมา เอวรูปิํ\footnote{เอวรูปํ (สี, สฺยา, กํ, อิ) เอวรูปาย-อิติ วุจฺจมานวจเนน สเมติ. วินเยนปิ สํสนฺเทตพฺพํ.} ครุกํ อาปตฺติํ อาปชฺชิตา ปาราชิกํ วา ปาราชิกสามนฺตํ วา”ติ.\csromanspacer{} โส เอวมาห “น โข อหํ อาวุโส สรามิ เอวรูปิํ ครุกํ อาปตฺติํ อาปชฺชิตา ปาราชิกํ วา ปาราชิกสามนฺตํ วา”ติ.\csromanspacer{} ตสฺส โข\footnote{ตสฺส โข เอวํ (สพฺพตฺถ)} อานนฺท ภิกฺขุโน สติวินโย ทาตพฺโพ.\csromanspacer{} เอวํ โข อานนฺท สติวินโย โหติ.\csromanspacer{} เอวญฺจ ปนิเธกจฺจานํ อธิกรณานํ วูปสโม โหติ, ยทิทํ สติวินเยน.}

\proseitem{50}{\csromanfolio{36}กถญฺจานนฺท อมูฬฺหวินโย โหติ, อิธานนฺท ภิกฺขู ภิกฺขุํ เอวรูปาย ครุกาย อาปตฺติยา โจเทนฺติ ปาราชิเกน วา ปาราชิกสามนฺเตน วา “สรตายสฺมา เอวรูปิํ ครุกํ อาปตฺติํ อาปชฺชิตา ปาราชิกํ วา ปาราชิกสามนฺตํ วา”ติ. (โส เอวมาห “น โข อหํ อาวุโส สรามิ เอวรูปิํ ครุกํ อาปตฺติํ อาปชฺชิตา ปาราชิกํ วา ปาราชิกสามนฺตํ วา”ติ.\csromanspacer{} ตเมนํ โส นิพฺเพเฐนฺตํ อติเวเฐติ “อิงฺฆายสฺมา สาธุกเมว ชานาหิ, ยทิ สรสิ เอวรูปิํ ครุกํ อาปตฺติํ อาปชฺชิตา ปาราชิกํ วา ปาราชิกสามนฺตํ วา”ติ.)\footnote{( ) เอตฺถนฺตเร ปาโฐ วิ 4 จูฬวคฺเค (230) ปิฏฺเฐ นตฺถิ, ตสฺสปาปิยสิกาวาเรเยเวเตน ภวิตพฺพํ.} โส เอวมาห “อหํ โข อาวุโส อุมฺมาทํ ปาปุณิํ, เจตโส วิปริยาสํ, เตน เม อุมฺมตฺตเกน พหุํ อสฺสามณกํ อชฺฌาจิณฺณํ ภาสิตปริกฺกนฺตํ\footnote{ภาสิตปริกนฺตํ (สี, สฺยา, กํ, อิ) 3. ตสฺส โข เอวํ (สฺยา, กํ, ก)} นาหํ ตํ สรามิ, มูเฬฺหน เม เอตํ กตนฺ”ติ.\csromanspacer{} ตสฺส โข3 อานนฺท ภิกฺขุโน อมูฬฺหวินโย ทาตพฺโพ.\csromanspacer{} เอวํ โข อานนฺท อมูฬฺหวินโย โหติ.\csromanspacer{} เอวญฺจ ปนิเธกจฺจานํ อธิกรณานํ วูปสโม โหติ, ยทิทํ อมูฬฺหวินเยน.}

\proseitem{51}{กถญฺจานนฺท ปฏิญฺญาตกรณํ โหติ, อิธานนฺท ภิกฺขุ โจทิโต วา อโจทิโต วา อาปตฺติํ สรติ วิวรติ อุตฺตานีกโรติ\footnote{อุตฺตานิํ กโรติ (ก)}.\csromanspacer{} เตน อานนฺท ภิกฺขุนา วุฑฺฒตรํ ภิกฺขุํ\footnote{วุฑฺฒตโร ภิกฺขุ (สี, สฺยา, กํ, อิ)} อุปสงฺกมิตฺวา เอกํสํ จีวรํ กตฺวา ปาเท วนฺทิตฺวา อุกฺกุฏิกํ นิสีทิตฺวา อญฺชลิํ ปคฺคเหตฺวา เอวมสฺส วจนีโย “อหํ ภนฺเต อิตฺถนฺนามํ อาปตฺติํ อาปนฺโน, ตํ ปฏิเทเสมี”ติ.\csromanspacer{} โส เอวมาห “ปสฺสสี”ติ. “อาม ปสฺสามี”ติ. “อายติํ สํวเรยฺยาสี”ติ. (สํวริสฺสามีติ.)\footnote{( ) วินเย นตฺถิ.} เอวํ โข อานนฺท ปฏิญฺญาตกรณํ โหติ.\csromanspacer{} เอวญฺจ ปนิเธกจฺจานํ อธิกรณานํ วูปสโม โหติ, ยทิทํ ปฏิญฺญาตกรเณน.}

\proseitem{52}{กถญฺจานนฺท ตสฺสปาปิยสิกา โหติ, อิธานนฺท ภิกฺขุ ภิกฺขุํ เอวรูปาย ครุกาย อาปตฺติยา โจเทติ ปาราชิเกน วา ปาราชิกสามนฺเตน วา “สรตายสฺมา เอวรูปิํ ครุกํ อาปตฺติํ อาปชฺชิตา ปาราชิกํ วา ปาราชิกสามนฺตํ วา”ติ.\csromanspacer{} โส เอวมาห “น โข อหํ อาวุโส สรามิ เอวรูปิํ ครุกํ อาปตฺติํ อาปชฺชิตา ปาราชิกํ วา}

\vspace{\tipitakaparskip}
\csromancenter{สามคามสุตฺต (104) ปาราชิกสามนฺตํ วา”ติ.\csromanspacer{} ตเมนํ โส นิพฺเพเฐนฺตํ อติเวเฐติ “อิงฺฆายสฺมา สาธุกเมว ชานาหิ, ยทิ สรสิ เอวรูปิํ ครุกํ อาปตฺติํ อาปชฺชิตา ปาราชิกํ วา ปาราชิกสามนฺตํ วา”ติ.\csromanspacer{} โส เอวมาห “น โข อหํ อาวุโส สรามิ เอวรูปิํ ครุกํ อาปตฺติํ อาปชฺชิตา ปาราชิกํ วา ปาราชิกสามนฺตํ วา, สรามิ จ โข อหํ อาวุโส เอวรูปิํ อปฺปมตฺติกํ อาปตฺติํ อาปชฺชิตา”ติ.\csromanspacer{} ตเมนํ โส นิพฺเพเฐนฺตํ อติเวเฐติ “อิงฺฆายสฺมา สาธุกเมว ชานาหิ, ยทิ สรสิ เอวรูปิํ ครุกํ อาปตฺติํ อาปชฺชิตา ปาราชิกํ วา ปาราชิกสามนฺตํ วา”ติ.\csromanspacer{} โส เอวมาห “อิมํ หิ นามาหํ อาวุโส อปฺปมตฺติกํ อาปตฺติํ อาปชฺชิตฺวา อปุฏฺโฐ ปฏิชานิสฺสามิ, กิํ ปนาหํ เอวรูปิํ ครุกํ อาปตฺติํ อาปชฺชิตฺวา ปาราชิกํ วา ปาราชิกสามนฺตํ วา ปุฏฺโฐ นปฏิชานิสฺสามี”ติ.\csromanspacer{} โส เอวมาห “อิมํ หิ นาม ตฺวํ อาวุโส อปฺปมตฺติกํ อาปตฺติํ อาปชฺชิตฺวา อปุฏฺโฐ นปฏิชานิสฺสสิ, กิํ ปน ตฺวํ เอวรูปิํ ครุกํ อาปตฺติํ อาปชฺชิตฺวา ปาราชิกํ วา ปาราชิกสามนฺตํ วา ปุฏฺโฐ\footnote{อปุฏฺโฐ (สฺยา, กํ, ก)} ปฏิชานิสฺสสิ, อิงฺฆายสฺมา สาธุกเมว ชานาหิ, ยทิ สรสิ เอวรูปิํ ครุกํ อาปตฺติํ อาปชฺชิตา ปาราชิกํ วา ปาราชิกสามนฺตํ วา”ติ.\csromanspacer{} โส เอวมาห “สรามิ โข อหํ อาวุโส เอวรูปิํ ครุกํ อาปตฺติํ อาปชฺชิตา ปาราชิกํ วา ปาราชิกสามนฺตํ วา, ทวา เม เอตํ วุตฺตํ, รวา เม เอตํ วุตฺตํ, นาหํ ตํ สรามิ เอวรูปิํ ครุกํ อาปตฺติํ อาปชฺชิตา ปาราชิกํ วา ปาราชิกสามนฺตํ วา”ติ.\csromanspacer{} เอวํ โข อานนฺท ตสฺสปาปิยสิกา โหติ.\csromanspacer{} เอวญฺจ ปนิเธกจฺจานํ อธิกรณานํ วูปสโม โหติ, ยทิทํ ตสฺสปาปิยสิกาย.}
\proseitem{53}{\csromanfolio{37}กถญฺจานนฺท ติณวตฺถารโก โหติ, อิธานนฺท ภิกฺขูนํ ภณฺฑนชาตานํ กลหชาตานํ วิวาทาปนฺนานํ วิหรตํ พหุํ อสฺสามณกํ อชฺฌาจิณฺณํ โหติ ภาสิตปริกฺกนฺตํ, เตหานนฺท ภิกฺขูหิ สพฺเพเหว สมคฺเคหิ สนฺนิปติตพฺพํ, สนฺนิปติตฺวา เอกโตปกฺขิกานํ ภิกฺขูนํ พฺยตฺเตน\footnote{พฺยตฺตตเรน (สี, อิ, ก)} ภิกฺขุนา อุฏฺฐายาสนา เอกํสํ จีวรํ กตฺวา อญฺชลิํ ปณาเมตฺวา สํโฆ ญาเปตพฺโพ\csromandash{}}

\prose{“สุณาตุ เม ภนฺเต สํโฆ, อิทํ อมฺหากํ ภณฺฑนชาตานํ กลหชาตานํ วิวาทาปนฺนานํ วิหรตํ พหุํ อสฺสามณกํ อชฺฌาจิณฺณํ \csromanfolio{38}ภาสิตปริกฺกนฺตํ, ยทิ สํฆสฺส ปตฺตกลฺลํ, อหํ ยา เจว อิเมสํ อายสฺมนฺตานํ อาปตฺติ, ยา จ อตฺตโน อาปตฺติ, อิเมสํ เจว อายสฺมนฺตานํ อตฺถาย อตฺตโน จ อตฺถาย สํฆมชฺเฌ ติณวตฺถารเกน เทเสยฺยํ ฐเปตฺวา ถุลฺลวชฺชํ ฐเปตฺวา คิหิปฏิสํยุตฺตนฺ”ติ.}

\prose{อถาปเรสํ เอกโตปกฺขิกานํ ภิกฺขูนํ พฺยตฺเตน ภิกฺขุนา อุฏฺฐายา สนา เอกํสํ จีวรํ กตฺวา อญฺชลิํ ปณาเมตฺวา สํโฆ ญาเปตพฺโพ\csromandash{}}

\prose{“สุณาตุ เม ภนฺเต สํโฆ, อิทํ อมฺหากํ ภณฺฑนชาตานํ กลหชาตานํ วิวาทาปนฺนานํ วิหรตํ พหุํ อสฺสามณกํ อชฺฌาจิณฺณํ ภาสิตปริกฺกนฺตํ, ยทิ สํฆสฺส ปตฺตกลฺลํ, อหํ ยา เจว อิเมสํ อายสฺมนฺตานํ อาปตฺติ, ยา จ อตฺตโน อาปตฺติ, อิเมสญฺเจว อายสฺมนฺตานํ อตฺถาย อตฺตโน จ อตฺถาย สํฆมชฺเฌ ติณวตฺถารเกน เทเสยฺยํ ฐเปตฺวา ถุลฺลวชฺชํ ฐเปตฺวา คิหิปฏิสํยุตฺตนฺ”ติ.}

\prose{เอวํ โข อานนฺท ติณวตฺถารโก โหติ.\csromanspacer{} เอวญฺจ ปนิเธกจฺจานํ อธิกรณานํ วูปสโม โหติ, ยทิทํ ติณวตฺถารเกน.}

\proseitem{54}{ฉยิเม อานนฺท ธมฺมา สารณียา ปิยกรณา ครุกรณา สงฺคหาย อวิวาทาย สามคฺคิยา เอกีภาวาย สํวตฺตนฺติ.\csromanspacer{} กตเม ฉ, อิธานนฺท ภิกฺขุโน เมตฺตํ กายกมฺมํ ปจฺจุปฏฺฐิตํ โหติ สพฺรหฺมจารีสุ อาวิ เจว รโห จ, อยมฺปิ ธมฺโม สารณีโย ปิยกรโณ ครุกรโณ สงฺคหาย อวิวาทาย สามคฺคิยา เอกีภาวาย สํวตฺตติ.}

\prose{ปุน จปรํ อานนฺท ภิกฺขุโน เมตฺตํ วจีกมฺมํ ปจฺจุปฏฺฐิตํ โหติ สพฺรหฺมจารีสุ อาวิ เจว รโห จ, อยมฺปิ ธมฺโม สารณีโย ปิยกรโณ ครุกรโณ สงฺคหาย อวิวาทาย สามคฺคิยา เอกีภาวาย สํวตฺตติ.}

\prose{ปุน จปรํ อานนฺท ภิกฺขุโน เมตฺตํ มโนกมฺมํ ปจฺจุปฏฺฐิตํ โหติ สพฺรหฺมจารีสุ อาวิ เจว รโห จ, อยมฺปิ ธมฺโม สารณีโย ปิยกรโณ ครุกรโณ สงฺคหาย อวิวาทาย สามคฺคิยา เอกีภาวาย สํวตฺตติ.}

\prose{ปุน จปรํ อานนฺท ภิกฺขุ เย เต ลาภา ธมฺมิกา ธมฺมลทฺธา อนฺตมโส ปตฺตปริยาปนฺนมตฺตมฺปิ, ตถารูเปหิ ลาเภหิ อปฏิวิภตฺตโภคี โหติ สีลวนฺเตหิ สพฺรหฺมจารีหิ สาธารณโภคี, อยมฺปิ}

\vspace{\tipitakaparskip}
\csromancenter{สุนกฺขตฺตสุตฺต (105) ธมฺโม สารณีโย ปิยกรโณ ครุกรโณ สงฺคหาย อวิวาทาย สามคฺคิยา เอกีภาวาย สํวตฺตติ.}
\prose{\csromanfolio{39}ปุน จปรํ อานนฺท ภิกฺขุ ยานิ ตานิ สีลานิ อขณฺฑานิ อจฺฉิทฺทานิ อสพลานิ อกมฺมาสานิ ภุชิสฺสานิ วิญฺญุปฺปสตฺถานิ อปรามฏฺฐานิ สมาธิสํวตฺตนิกานิ, ตถารูเปสุ สีเลสุ สีลสามญฺญคโต วิหรติ สพฺรหฺมจารีหิ อาวิ เจว รโห จ, อยมฺปิ ธมฺโม สารณีโย ปิยกรโณ ครุกรโณ สงฺคหาย อวิวาทาย สามคฺคิยา เอกีภาวาย สํวตฺตติ.}

\prose{ปุน จปรํ อานนฺท ภิกฺขุ ยายํ ทิฏฺฐิ อริยา นิยฺยานิกา นิยฺยาติ ตกฺกรสฺส สมฺมา ทุกฺขกฺขยาย, ตถารูปาย ทิฏฺฐิยา ทิฏฺฐิสามญฺญคโต วิหรติ สพฺรหฺมจารีหิ อาวิ เจว รโห จ, อยมฺปิ ธมฺโม สารณีโย ปิยกรโณ ครุกรโณ สงฺคหาย อวิวาทาย สามคฺคิยา เอกีภาวาย สํวตฺตติ.\csromanspacer{} อิเม โข อานนฺท ฉ สารณียา ธมฺมา ปิยกรณา ครุกรณา สงฺคหาย อวิวาทาย สามคฺคิยา เอกีภาวาย สํวตฺตนฺติ.}

\prose{อิเม เจ ตุเมฺห อานนฺท ฉ สารณีเย ธมฺเม สมาทาย วตฺเตยฺยาถ, ปสฺสถ โน ตุเมฺห อานนฺท ตํ วจนปถํ อณุํ วา ถูลํ วา, ยํ ตุเมฺห นาธิวาเสยฺยาถาติ.\csromanspacer{} โน เหตํ ภนฺเต.\csromanspacer{} ตสฺมาติหานนฺท อิเม ฉ สารณีเย ธมฺเม สมาทาย วตฺตถ, ตํ โว ภวิสฺสติ ทีฆรตฺตํ หิตาย สุขายาติ.}

\prose{อิทมโวจ ภควา.\csromanspacer{} อตฺตมโน อายสฺมา อานนฺโท ภควโต ภาสิตํ อภินนฺทีติ.}

\nitthitamruled{สามคามสุตฺตํ นิฏฺฐิตํ จตุตฺถํ.}
\csromanmatikaanchor{mtk.6}
\csromantocmarkref{subsection}{5. สุนกฺขตฺตสุตฺต}{mtk.6}
\bookmark[dest={mtk.6},level=2]{5. สุนกฺขตฺตสุตฺต}
\csromanheader{2}{5. สุนกฺขตฺตสุตฺต}
\proseitem{55}{เอวํ เม สุตํ\csromandash{}เอกํ สมยํ ภควา เวสาลิยํ วิหรติ มหาวเน กูฏาคารสาลายํ.\csromanspacer{} เตน โข ปน สมเยน สมฺพหุเลหิ \csromanfolio{40}ภิกฺขูหิ ภควโต สนฺติเก อญฺญา พฺยากตา โหติ “ขีณา ชาติ, วุสิตํ พฺรหฺมจริยํ, กตํ กรณียํ, นาปรํ อิตฺถตฺตายาติ ปชานามา”ติ.\csromanspacer{} อสฺโสสิ โข สุนกฺขตฺโต ลิจฺฉวิปุตฺโต “สมฺพหุเลหิ กิร ภิกฺขูหิ ภควโต สนฺติเก อญฺญา พฺยากตา โหติ ‘ขีณา ชาติ, วุสิตํ พฺรหฺมจริยํ, กตํ กรณียํ, นาปรํ อิตฺถตฺตายาติ ปชานามา’ติ”.\csromanspacer{} อถ โข สุนกฺขตฺโต ลิจฺฉวิปุตฺโต เยน ภควา เตนุปสงฺกมิ, อุปสงฺกมิตฺวา ภควนฺตํ อภิวาเทตฺวา เอกมนฺตํ นิสีทิ, เอกมนฺตํ นิสินฺโน โข สุนกฺขตฺโต ลิจฺฉวิปุตฺโต ภควนฺตํ เอตทโวจ\csromandash{}สุตํ เมตํ ภนฺเต “สมฺพหุเลหิ กิร ภิกฺขูหิ ภควโต สนฺติเก อญฺญา พฺยากตา ‘ขีณา ชาติ, วุสิตํ พฺรหฺมจริยํ, กตํ กรณียํ, นาปรํ อิตฺถตฺตายาติ ปชานามา’ติ”.\csromanspacer{} เย เต ภนฺเต ภิกฺขู ภควโต สนฺติเก อญฺญํ พฺยากํสุ “ขีณา ชาติ, วุสิตํ พฺรหฺมจริยํ, กตํ กรณียํ, นาปรํ อิตฺถตฺตายาติ ปชานามา”ติ, กจฺจิ เต ภนฺเต ภิกฺขู สมฺมเทว อญฺญํ พฺยากํสุ, อุทาหุ สนฺเตตฺเถกจฺเจ ภิกฺขู อธิมาเนน อญฺญํ พฺยากํสูติ.}

\proseitem{56}{เย เต สุนกฺขตฺต ภิกฺขู มม สนฺติเก อญฺญํ พฺยากํสุ “ขีณา ชาติ, วุสิตํ พฺรหฺมจริยํ, กตํ กรณียํ, นาปรํ อิตฺถตฺตายาติ ปชานามา”ติ, สนฺเตตฺเถกจฺเจ ภิกฺขู สมฺมเทว อญฺญํ พฺยากํสุ, สนฺติ ปนิเธกจฺเจ ภิกฺขู อธิมาเนนปิ\footnote{อธิมาเนน (?)} อญฺญํ พฺยากํสุ.\csromanspacer{} ตตฺร สุนกฺขตฺต เย เต ภิกฺขู สมฺมเทว อญฺญํ พฺยากํสุ.\csromanspacer{} เตสํ ตํ ตเถว โหติ.\csromanspacer{} เย ปน เต ภิกฺขู อธิมาเนน อญฺญํ พฺยากํสุ.\csromanspacer{} ตตฺร สุนกฺขตฺต ตถาคตสฺส เอวํ โหติ “ธมฺมํ เนสํ เทเสสฺสนฺ”ติ\footnote{เทเสยฺยนฺติ (อิ, ก)}.\csromanspacer{} เอวญฺเจตฺถ สุนกฺขตฺต ตถาคตสฺส โหติ “ธมฺมํ เนสํ เทเสสฺสนฺ”ติ.\csromanspacer{} อถ จ ปนิเธกจฺเจ โมฆปุริสา ปญฺหํ อภิสงฺขริตฺวา อภิสงฺขริตฺวา ตถาคตํ อุปสงฺกมิตฺวา ปุจฺฉนฺติ.\csromanspacer{} ตตฺร สุนกฺขตฺต ยมฺปิ ตถาคตสฺส เอวํ โหติ “ธมฺมํ เนสํ เทเสสฺสนฺ”ติ, ตสฺสปิ โหติ อญฺญถตฺตนฺติ.\csromanspacer{} เอตสฺส ภควา กาโล, เอตสฺส สุคต กาโล, ยํ ภควา ธมฺมํ เทเสยฺย, ภควโต สุตฺวา ภิกฺขู ธาเรสฺสนฺตีติ.\csromanspacer{} เตน หิ สุนกฺขตฺต สุณาหิ สาธุกํ มนสิ กโรหิ ภาสิสฺสามีติ. “เอวํ ภนฺเต”ติ โข สุนกฺขตฺโต ลิจฺฉวิปุตฺโต ภควโต ปจฺจสฺโสสิ.\csromanspacer{} ภควา เอตทโวจ\csromandash{}}

\csromanheader{2}{5. สุนกฺขตฺตสุตฺต (105)}
\proseitem{57}{\csromanfolio{41}ปญฺจ โข อิเม สุนกฺขตฺต กามคุณา.\csromanspacer{} กตเม ปญฺจ, จกฺขุวิญฺเญยฺยา รูปา อิฏฺฐา กนฺตา มนาปา ปิยรูปา กามูปสํหิตา รชนียา.\csromanspacer{} โสตวิญฺเญยฺยา สทฺทา ฯเปฯ.\csromanspacer{} ฆานวิญฺเญยฺยา คนฺธา.\csromanspacer{} ชิวฺหาวิญฺเญยฺยา รสา.\csromanspacer{} กายวิญฺเญยฺยา โผฏฺฐพฺพา อิฏฺฐา กนฺตา มนาปา ปิยรูปา กามูปสํหิตา รชนียา.\csromanspacer{} อิเม โข สุนกฺขตฺต ปญฺจ กามคุณา.}

\proseitem{58}{ฐานํ โข ปเนตํ สุนกฺขตฺต วิชฺชติ, ยํ อิเธกจฺโจ ปุริสปุคฺคโล โลกามิสาธิมุตฺโต อสฺส, โลกามิสาธิมุตฺตสฺส โข สุนกฺขตฺต ปุริสปุคฺคลสฺส ตปฺปติรูปี เจว กถา สณฺฐาติ, ตทนุธมฺมญฺจ อนุวิตกฺเกติ อนุวิจาเรติ, ตญฺจ ปุริสํ ภชติ, เตน จ วิตฺติํ อาปชฺชติ.\csromanspacer{} อาเนญฺชปฏิสํยุตฺตาย จ ปน กถาย กจฺฉมานาย น สุสฺสูสติ, น โสตํ โอทหติ, น อญฺญา จิตฺตํ อุปฏฺฐาเปติ\footnote{อุปฏฺฐเปติ (สี, สฺยา, กํ, อิ)}, น จ ตํ ปุริสํ ภชติ, น จ เตน วิตฺติํ อาปชฺชติ.\csromanspacer{} เสยฺยถาปิ สุนกฺขตฺต ปุริโส สกมฺหา คามา วา นิคมา วา จิรวิปฺปวุตฺโถ อสฺส, โส อญฺญตรํ ปุริสํ ปสฺเสยฺย ตมฺหา คามา วา นิคมา วา อจิรปกฺกนฺตํ, โส ตํ ปุริสํ ตสฺส คามสฺส วา นิคมสฺส วา เขมตญฺจ สุภิกฺขตญฺจ อปฺปาพาเตญฺจ ปุจฺเฉยฺย, ตสฺส โส ปุริโส ตสฺส คามสฺส วา นิคมสฺส วา เขมตญฺจ สุภิกฺขตญฺจ อปฺปาพาธตญฺจ สํเสยฺย.\csromanspacer{} ตํ กิํ มญฺญสิ สุนกฺขตฺต, อปิ นุ โส ปุริโส ตสฺส ปุริสสฺส สุสฺสูเสยฺย, โสตํ โอทเหยฺย, อญฺญา จิตฺตํ อุปฏฺฐาเปยฺย, ตญฺจ ปุริสํ ภเชยฺย, เตน จ วิตฺติํ อาปชฺเชยฺยาติ.\csromanspacer{} เอวํ ภนฺเต.\csromanspacer{} เอวเมว โข สุนกฺขตฺต ฐานเมตํ วิชฺชติ, ยํ อิเธกจฺโจ ปุริสปุคฺคโล โลกามิสาธิมุตฺโต อสฺส, โลกามิสาธิมุตฺตสฺส โข สุนกฺขตฺต ปุริสปุคฺคลสฺส ตปฺปติรูปี เจว กถา สณฺฐาติ, ตทนุธมฺมญฺจ อนุวิตกฺเกติ อนุวิจาเรติ, ตญฺจ ปุริสํ ภชติ, เตน จ วิตฺติํ อาปชฺชติ.\csromanspacer{} อาเนญฺชปฏิสํยุตฺตาย จ ปน กถาย กจฺฉมานาย น สุสฺสูสติ, น โสตํ โอทหติ, น อญฺญา จิตฺตํ อุปฏฺฐาเปติ, น จ ตํ ปุริสํ ภชติ, น จ เตน วิตฺติํ อาปชฺชติ.\csromanspacer{} โส เอวมสฺส เวทิตพฺโพ “อาเนญฺชสํโยชเนน หิ โข วิสํยุตฺโต\footnote{อาเนญฺชสํโยชเนน หิ โข วิสํยุตฺโต-อิติ ปาโฐ สี-สฺยา-กํ-อิ-โปตฺถเกสุ นตฺถิ, อฏฺฐกถาสุ ปน ตพฺพณฺณนา ทิสฺสติเยว.} โลกามิสาธิมุตฺโต ปุริสปุคฺคโล”ติ.}

\proseitem{59}{\csromanfolio{42}ฐานํ โข ปเนตํ สุนกฺขตฺต วิชฺชติ, ยํ อิเธกจฺโจ ปุริสปุคฺคโล อาเนญฺชาธิมุตฺโต อสฺส, อาเนญฺชาธิมุตฺตสฺส โข สุนกฺขตฺต ปุริสปุคฺคลสฺส ตปฺปติรูปี เจว กถา สณฺฐาติ, ตทนุธมฺมญฺจ อนุวิตกฺเกติ อนุวิจาเรติ, ตญฺจ ปุริสํ ภชติ, เตน จ วิตฺติํ อาปชฺชติ.\csromanspacer{} โลกามิสปฏิสํยุตฺตาย จ ปน กถาย กจฺฉมานาย น สุสฺสูสติ, น โสตํ โอทหติ, น อญฺญา จิตฺตํ อุปฏฺฐาเปติ, น จ ตํ ปุริสํ ภชติ, น จ เตน วิตฺติํ อาปชฺชติ.\csromanspacer{} เสยฺยถาปิ สุนกฺขตฺต ปณฺฑุปลาโส พนฺธนา ปวุตฺโต อภพฺโพ หริตตฺตาย.\csromanspacer{} เอวเมว โข สุนกฺขตฺต อาเนญฺชาธิมุตฺตสฺส ปุริสปุคฺคลสฺส เย โลกามิสสํโยชเน, เส ปวุตฺเต.\csromanspacer{} โส เอวมสฺส เวทิตพฺโพ “โลกามิสสํโยชเนน หิ โข วิสํยุตฺโต อาเนญฺชาธิมุตฺโต ปุริสปุคฺคโล”ติ.}

\proseitem{60}{ฐานํ โข ปเนตํ สุนกฺขตฺต วิชฺชติ, ยํ อิเธกจฺโจ ปุริสปุคฺคโล อากิญฺจญฺญายตนาธิมุตฺโต อสฺส, อากิญฺจญฺญายตนาธิมุตฺตสฺส โข สุนกฺขตฺต ปุริสปุคฺคลสฺส ตปฺปติรูปี เจว กถา สณฺฐาติ, ตทนุธมฺมญฺจ อนุวิตกฺเกติ อนุวิจาเรติ, ตญฺจ ปุริสํ ภชติ, เตน จ วิตฺติํ อาปชฺชติ.\csromanspacer{} อาเนญฺชปฏิสํยุตฺตาย จ ปน กถาย กจฺฉมานาย น สุสฺสูสติ, น โสตํ โอทหติ, น อญฺญา จิตฺตํ อุปฏฺฐาเปติ, น จ ตํ ปุริสํ ภชติ, น จ เตน วิตฺติํ อาปชฺชติ.\csromanspacer{} เสยฺยถาปิ สุนกฺขตฺต ปุถุสิลา เทฺวธาภินฺนา อปฺปฏิสนฺธิกา โหติ.\csromanspacer{} เอวเมว โข สุนกฺขตฺต อากิญฺจญฺญายตนาธิมุตฺตสฺส ปุริสปุคฺคลสฺส เย อาเนญฺชสํโยชเน, เส ภินฺเน.\csromanspacer{} โส เอวมสฺส เวทิตพฺโพ “อาเนญฺชสํโยชเนน หิ โข วิสํยุตฺโต อากิญฺจญฺญายตนาธิมุตฺโต ปุริสปุคฺคโล”ติ.}

\proseitem{61}{ฐานํ โข ปเนตํ สุนกฺขตฺต วิชฺชติ, ยํ อิเธกจฺโจ ปุริสปุคฺคโล เนวสญฺญานาสญฺญายตนาธิมุตฺโต อสฺส, เนวสญฺญานาสญฺญายตนาธิมุตฺตสฺส โข สุนกฺขตฺต ปุริสปุคฺคลสฺส ตปฺปติรูปี เจว กถา สณฺฐาติ, ตทนุธมฺมญฺจ อนุวิตกฺเกติ อนุวิจาเรติ, ตญฺจ ปุริสํ ภชติ, เตน จ วิตฺติํ อาปชฺชติ.\csromanspacer{} อากิญฺจญฺญายตนปฏิสํยุตฺตาย จ ปน กถาย กจฺฉมานาย น สุสฺสูสติ, น โสตํ โอทหติ, น อญฺญา จิตฺตํ อุปฏฺฐาเปติ, น จ ตํ ปุริสํ ภชติ, น จ เตน วิตฺติํ อาปชฺชติ.\csromanspacer{} เสยฺยถาปิ สุนกฺขตฺต ปุริโส มนุญฺญโภชนํ ภุตฺตาวี ฉฑฺเฑยฺย\footnote{ฉทฺเทยฺย (?)}.\csromanspacer{} ตํ กิํ มญฺญสิ สุนกฺขตฺต,}

\vspace{\tipitakaparskip}
\csromancenter{สุนกฺขตฺตสุตฺต (105) อปิ นุ ตสฺส ปุริสสฺส ตสฺมิํ ภตฺเต\footnote{วนฺเต (ก-สี), ภุตฺเต (ก-สี, ก)} ปุน โภตฺตุกมฺยตา อสฺสาติ.\csromanspacer{} โน เหตํ ภนฺเต.\csromanspacer{} ตํ กิสฺส เหตุ, อทุํ หิ ภนฺเต ภตฺตํ\footnote{วนฺตํ (สี)} ปฏิกูลสมฺมตนฺติ.\csromanspacer{} เอวเมว โข สุนกฺขตฺต เนวสญฺญานาสญฺญายตนาธิมุตฺตสฺส ปุริสปุคฺคลสฺส เย อากิญฺจญฺญายตนสํโยชเน, เส วนฺเต.\csromanspacer{} โส เอวมสฺส เวทิตพฺโพ “อากิญฺจญฺญายตนสํโยชเนน หิ โข วิสํยุตฺโต เนวสญฺญานาสญฺญายตนาธิมุตฺโต ปุริสปุคฺคโล”ติ.}
\proseitem{62}{\csromanfolio{43}ฐานํ โข ปเนตํ สุนกฺขตฺต วิชฺชติ, ยํ อิเธกจฺโจ ปุริสปุคฺคโล สมฺมา นิพฺพานาธิมุตฺโต อสฺส, สมฺมา นิพฺพานาธิมุตฺตสฺส โข สุนกฺขตฺต ปุริสปุคฺคลสฺส ตปฺปติรูปี เจว กถา สณฺฐาติ, ตทนุธมฺมญฺจ อนุวิตกฺเกติ อนุวิจาเรติ, ตญฺจ ปุริสํ ภชติ, เตน จ วิตฺติํ อาปชฺชติ.\csromanspacer{} เนวสญฺญานาสญฺญายตนปฏิสํยุตฺตาย จ ปน กถาย กจฺฉมานาย น สุสฺสูสติ, น โสตํ โอทหติ, น อญฺญา จิตฺตํ อุปฏฺฐาเปติ, น จ ตํ ปุริสํ ภชติ, น จ เตน วิตฺติํ อาปชฺชติ.\csromanspacer{} เสยฺยถาปิ สุนกฺขตฺต ตาโล มตฺตกจฺฉินฺโน อภพฺโพ ปุน วิรุฬฺหิยา.\csromanspacer{} เอวเมว โข สุนกฺขตฺต สมฺมา นิพฺพานาธิมุตฺตสฺส ปุริสปุคฺคลสฺส เย เนวสญฺญานาสญฺญายตนสํโยชเน, เส อุจฺฉินฺนมูเล ตาลาวตฺถุกเต อนภาวํกเต\footnote{อนภาวกเต (สี, อิ), อนภาวงฺคเต (สฺยา, กํ)} อายติํ อนุปฺปาทธมฺเม.\csromanspacer{} โส เอวมสฺส เวทิตพฺโพ “เนวสญฺญานาสญฺญายตนสํโยชเนน หิ โข วิสํยุตฺโต สมฺมา นิพฺพานาธิมุตฺโต ปุริสปุคฺคโล”ติ.}

\proseitem{63}{ฐานํ โข ปเนตํ สุนกฺขตฺต วิชฺชติ, ยํ อิเธกจฺจสฺส ภิกฺขุโน เอวมสฺส “ตณฺหา โข สลฺลํ สมเณน วุตฺตํ, อวิชฺชาวิสโทโส ฉนฺทราคพฺยาปาเทน รุปฺปติ, ตํ เม ตณฺหาสลฺลํ ปหีนํ, อปนีโต อวิชฺชาวิสโทโส, สมฺมา นิพฺพานาธิมุตฺโตหมสฺมี”ติ, เอวํมานิ\footnote{เอวํมานี (สี, อิ, ก), เอวมาทิ (สฺยา, กํ)} อสฺส อตถํ สมานํ\footnote{อตฺถํ สมานํ (สฺยา, กํ, อิ), อตฺถสมานํ (สี)}.\csromanspacer{} โส ยานิ สมฺมา นิพฺพานาธิมุตฺตสฺส อสปฺปายานิ, ตานิ อนุยุญฺเชยฺย, อสปฺปายํ จกฺขุนา รูปทสฺสนํ อนุยุญฺเชยฺย, อสปฺปายํ โสเตน สทฺทํ อนุยุญฺเชยฺย, อสปฺปายํ ฆาเนน คนฺธํ อนุยุญฺเชยฺย, อสปฺปายํ ชิวฺหาย รสํ อนุยุญฺเชยฺย, อสปฺปายํ กาเยน โผฏฺฐพฺพํ \csromanfolio{44}อนุยุญฺเชยฺย, อสปฺปายํ มนสา ธมฺมํ อนุยุญฺเชยฺย.\csromanspacer{} ตสฺส อสปฺปายํ จกฺขุนา รูปทสฺสนํ อนุยุตฺตสฺส, อสปฺปายํ โสเตน สทฺทํ อนุยุตฺตสฺส, อสปฺปายํ ฆาเนน คนฺธํ อนุยุตฺตสฺส, อสปฺปายํ ชิวฺหาย รสํ อนุยุตฺตสฺส, อสปฺปายํ กาเยน โผฏฺฐพฺพํ อนุยุตฺตสฺส, อสปฺปายํ มนสา ธมฺมํ อนุยุตฺตสฺส ราโค จิตฺตํ อนุทฺธํเสยฺย, โส ราคานุทฺธํสิเตน จิตฺเตน มรณํ วา นิคจฺเฉยฺย มรณมตฺตํ วา ทุกฺขํ.}

\prose{เสยฺยถาปิ สุนกฺขตฺต ปุริโส สลฺเลน วิทฺโธ อสฺส สวิเสน คาฬฺหูปเลปเนน.\csromanspacer{} ตสฺส มิตฺตามจฺจา ญาติสาโลหิตา ภิสกฺกํ สลฺลกตฺตํ อุปฏฺฐาเปยฺยุํ, ตสฺส โส ภิสกฺโก สลฺลกตฺโต สตฺเถน วณมุขํ ปริกนฺเตยฺย, สตฺเถน วณมุขํ ปริกนฺติตฺวา เอสนิยา สลฺลํ เอเสยฺย, เอสนิยา สลฺลํ เอสิตฺวา สลฺลํ อพฺพุเหยฺย, อปเนยฺย วิสโทสํ ส-อุปาทิเสสํ ส-อุปาทิเสโสติ\footnote{อนุปาทิเสโสติ (สพฺพตฺถ) อยํ หิ ตถาคตสฺส วิสโย.} ชานมาโน.\csromanspacer{} โส เอวํ วเทยฺย “อมฺโภ ปุริส อุพฺภตํ โข เต สลฺลํ, อปนีโต วิสโทโส ส-อุปาทิเสโส\footnote{อนุปาทิเสโส (สพฺพตฺถ) อยมฺปิ ตถาคตสฺส วิสโย.}, อนลญฺจ เต อนฺตรายาย, สปฺปายานิ เจว โภชนานิ ภุญฺเชยฺยาสิ, มา เต อสปฺปายานิ โภชนานิ ภุญฺชโต วโณ อสฺสาวี อสฺส.\csromanspacer{} กาเลน กาลํ จ วณํ โธเวยฺยาสิ, กาเลน กาลํ วณมุขํ อาลิมฺเปยฺยาสิ, มา เต น กาเลน กาลํ วณํ โธวโต น กาเลน กาลํ วณมุขํ อาลิมฺปโต ปุพฺพโลหิตํ วณมุขํ ปริโยนนฺธิ.\csromanspacer{} มา จ วาตาตเป จาริตฺตํ อนุยุญฺชิ, มา เต วาตาตเป จาริตฺตํ อนุยุตฺตสฺส รโชสูกํ วณมุขํ อนุทฺธํเสสิ.\csromanspacer{} วณานุรกฺขี จ อมฺโภ ปุริส วิหเรยฺยาสิ วณสาโรปี”ติ\footnote{วณสฺสาโรปีติ (ก) วณ + สํ + โรปี = วณสาโรปี-อิติ ปทวิภาโค.}.\csromanspacer{} ตสฺส เอวมสฺส “อุพฺภตํ โข เม สลฺลํ, อปนีโต วิสโทโส อนุปาทิเสโส, อนลญฺจ เม อนฺตรายายา”ติ.\csromanspacer{} โส อสปฺปายานิ เจว โภชนานิ ภุญฺเชยฺย, ตสฺส อสปฺปายานิ โภชนานิ ภุญฺชโต วโณ อสฺสาวี อสฺส.\csromanspacer{} น จ กาเลน กาลํ วณํ โธเวยฺย, น จ กาเลน กาลํ วณมุขํ อาลิมฺเปยฺย, ตสฺส น กาเลน กาลํ วณํ โธวโต น กาเลน กาลํ วณมุขํ อาลิมฺปโต ปุพฺพโลหิตํ วณมุขํ ปริโยนนฺเธยฺย.\csromanspacer{} วาตาตเป จ จาริตฺตํ อนุยุญฺเชยฺย, ตสฺส วาตาตเป จาริตฺตํ อนุยุตฺตสฺส รโชสูกํ}

\vspace{\tipitakaparskip}
\csromancenter{สุนกฺขตฺตสุตฺต (105) วณมุขํ อนุทฺธํเสยฺย.\csromanspacer{} น จ วณานุรกฺขี วิหเรยฺย น วณสาโรปี.\csromanspacer{} ตสฺส อิมิสฺสา จ อสปฺปายกิริยาย อสุจิ วิสโทโส, อปนีโต ส- อุปาทิเสโส, ตทุภเยน วโณ ปุถุตฺถํ คจฺเฉยฺย, โส ปุถุตฺตํ คเตน วเณน มรณํ วา นิคจฺเฉยฺย มรณมตฺตํ วา ทุกฺขํ.}
\prose{\csromanfolio{45}เอวเมว โข สุนกฺขตฺต ฐานเมตํ วิชฺชติ, ยํ อิเธกจฺจสฺส ภิกฺขุโน เอวมสฺส “ตณฺหา โข สลฺลํ สมเณน วุตฺตํ, อวิชฺชาวิสโทโส ฉนฺทราคพฺยาปาเทน รุปฺปติ, ตํ เม ตณฺหาสลฺลํ ปหีนํ, อปนีโต อวิชฺชาวิสโทโส, สมฺมา นิพฺพานาธิมุตฺโตหมสฺมี”ติ, เอวํมานิ อสฺส อตถํ สมานํ.\csromanspacer{} โส ยานิ สมฺมา นิพฺพานิธิมุตฺตสฺส อสปฺปายานิ, ตานิ อนุยุญฺเชยฺย, อสปฺปายํ จกฺขุนา รูปทสฺสนํ อนุยุญฺเชยฺย, อสปฺปายํ โสเตน สทฺทํ อนุยุญฺเชยฺย, อสปฺปายํ ฆาเนน คนฺธํ อนุยุญฺเชยฺย, อสปฺปายํ ชิวฺหาย รสํ อนุยุญฺเชยฺย, อสปฺปายํ กาเยน โผฏฺฐพฺพํ อนุยุญฺเชยฺย, อสปฺปายํ มนสา ธมฺมํ อนุยุญฺเชยฺย.\csromanspacer{} ตสฺส อสปฺปายํ จกฺขุนา รูปทสฺสนํ อนุยุตฺตสฺส, อสปฺปายํ โสเตน สทฺทํ อนุยุตฺตสฺส, อสปฺปายํ ฆาเนน คนฺธํ อนุยุตฺตสฺส, อสปฺปายํ ชิวฺหาย รสํ อนุยุตฺตสฺส, อสปฺปายํ กาเยน โผฏฺฐพฺพํ อนุยุตฺตสฺส, อสปฺปายํ มนสา ธมฺมํ อนุยุตฺตสฺส ราโค จิตฺตํ อนุทฺธํเสยฺย, โส ราคานุทฺธํสิเตน จิตฺเตน มรณํ วา นิคจฺเฉยฺย มรณมตฺตํ วา ทุกฺขํ.\csromanspacer{} มรณํ เหตํ สุนกฺขตฺต อริยสฺส วินเย, โย สิกฺขํ ปจฺจกฺขาย หีนายาวตฺตติ.\csromanspacer{} มรณมตฺตํ เหตํ สุนกฺขตฺต ทุกฺขํ, ยํ อญฺญตรํ สํกิลิฏฺฐํ อาปตฺติํ อาปชฺชติ.}

\proseitem{64}{ฐานํ โข ปเนตํ สุนกฺขตฺต วิชฺชติ, ยํ อิเธกจฺจสฺส ภิกฺขุโน เอวมสฺส “ตณฺหา โข สลฺลํ สมเณน วุตฺตํ, อวิชฺชาวิสโทโส ฉนฺทราคพฺยาปาเทน รุปฺปติ, ตํ เม ตณฺหาสลฺลํ ปหีนํ, อปนีโต อวิชฺชาวิสโทโส, สมฺมา นิพฺพานาธิมุตฺโตหมสฺมี”ติ สมฺมา นิพฺพานาธิมุตฺตสฺเสว สโต.\csromanspacer{} โส ยานิ สมฺมา นิพฺพานาธิมุตฺตสฺส อสปฺปายานิ ตานิ นานุยุญฺเชยฺย, อสปฺปายํ จกฺขุนา รูปทสฺสนํ นานุยุญฺเชยฺย, อสปฺปายํ โสเตน สทฺทํ นานุยุญฺเชยฺย, อสปฺปายํ ฆาเนน คนฺธํ นานุยุญฺเชยฺย, อสปฺปายํ ชิวฺหาย รสํ นานุยุญฺเชยฺย, อสปฺปายํ กาเยน โผฏฺฐพฺพํ นานุยุญฺเชยฺย, อสปฺปายํ มนสา ธมฺมํ นานุยุญฺเชยฺย.\csromanspacer{} ตสฺส อสปฺปายํ จกฺขุนา รูปทสฺสนํ นานุยุตฺตสฺส, อสปฺปายํ โสเตน สทฺทํ นานุยุตฺตสฺส, อสปฺปายํ ฆาเนน คนฺธํ นานุยุตฺตสฺส, อสปฺปายํ ชิวฺหาย รสํ นานุยุตฺตสฺส, \csromanfolio{46}อสปฺปายํ กาเยน โผฏฺฐพฺพํ นานุยุตฺตสฺส, อสปฺปายํ มนสา ธมฺมํ นานุยุตฺตสฺส ราโค จิตฺตํ นานุทฺธํเสยฺย, โส น ราคานุทฺธํสิเตน จิตฺเตน เนว มรณํ วา นิคจฺเฉยฺย น มรณมตฺตํ วา ทุกฺขํ.}

\prose{เสยฺยถาปิ สุนกฺขตฺต ปุริโส สลฺเลน วิทฺโธ อสฺส สวิเสน คาฬฺหูปเลปเนน.\csromanspacer{} ตสฺส มิตฺตามจฺจา ญาติสาโลหิตา ภิสกฺกํ สลฺลกตฺตํ อุปฏฺฐาเปยฺยุํ, ตสฺส โส ภิสกฺโก สลฺลกตฺโต สตฺเถน วณมุขํ ปริกนฺเตยฺย, สตฺเถน วณมุขํ ปริกนฺติตฺวา เอสนิยา สลฺลํ เอเสยฺย, เอสนิยา สลฺลํ เอสิตฺวา สลฺลํ อพฺพุเหยฺย, อปเนยฺย วิสโทสํ อนุปาทิเสสํ อนุปาทิเสโสติ ชานมาโน.\csromanspacer{} โส เอวํ วเทยฺย “อมฺโภ ปุริส อุพฺภตํ โข เต สลฺลํ, อปนีโต วิสโทโส อนุปาทิเสโส, อนลญฺจ เต อนฺตรายาย, สปฺปายานิ เจว โภชนานิ ภุญฺเชยฺยาสิ, มา เต อสปฺปายานิ โภชนานิ ภุญฺชโต วโณ อสฺสาวี อสฺส.\csromanspacer{} กาเลน กาลํ จ วณํ โธเวยฺยาสิ, กาเลน กาลํ วณมุขํ อาลิมฺเปยฺยาสิ, มา เต น กาเลน กาลํ วณํ โธวโต น กาเลน กาลํ วณมุขํ อาลิมฺปโต ปุพฺพโลหิตํ วณมุขํ ปริโยนนฺธิ.\csromanspacer{} มา จ วาตาตเป จาริตฺตํ อนุยุญฺชิ, มา เต วาตาตเป จาริตฺตํ อนุยุตฺตสฺส รโชสูกํ วณมุขํ อนุทฺธํเสสิ.\csromanspacer{} วณานุรกฺขี จ อมฺโภ ปุริส วิหเรยฺยาสิ วณสาโรปี”ติ.\csromanspacer{} ตสฺส เอวมสฺส “อุพฺภตํ โข เม สลฺลํ, อปนีโต วิสโทโส อนุปาทิเสโส, อนลญฺจ เม อนฺตรายายา”ติ.\csromanspacer{} โส สปฺปายานิ เจว โภชนานิ ภุญฺเชยฺย, ตสฺส สปฺปายานิ โภชนานิ ภุญฺชโต วโณ น อสฺสาวี อสฺส.\csromanspacer{} กาเลน กาลํ จ วณํ โธเวยฺย, กาเลน กาลํ วณํมุขํ อาลิมฺเปยฺย, ตสฺส กาเลน กาลํ วณํ โธวโต กาเลน กาลํ วณมุขํ อาลิมฺปโต น ปุพฺพโลหิตํ วณมุขํ ปริโยนนฺเธยฺย.\csromanspacer{} น จ วาตาตเป จาริตฺตํ อนุยุญฺเชยฺย, ตสฺส วาตาตเป จาริตฺตํ อนนุยุตฺตสฺส รโชสูกํ วณมุขํ นานุทฺธํเสยฺย.\csromanspacer{} วณานุรกฺขี จ วิหเรยฺย วณสาโรปี.\csromanspacer{} ตสฺส อิมิสฺสา จ สปฺปายกิริยาย อสุ จ\footnote{อสุจิ (สพฺพตฺถ) โสจาติ ตพฺพณฺณนา มนสิกาตพฺพา.} วิสโทโส อปนีโต อนุปาทิเสโส.\csromanspacer{} ตทุภเยน วโณ วิรุเหยฺย, โส รุเฬฺหน วเณน สญฺฉวินา เนว มรณํ วา นิคจฺเฉยฺย น มรณมตฺตํ วา ทุกฺขํ.}

\csromanheader{2}{5. สุนกฺขตฺตสุตฺต (105)}
\prose{\csromanfolio{47}เอวเมว โข สุนกฺขตฺต ฐานเมตํ วิชฺชติ, ยํ อิเธกจฺจสฺส ภิกฺขุโน เอวมสฺส “ตณฺหา โข สลฺลํ สมเณน วุตฺตํ, อวิชฺชาวิสโทโส ฉนฺทราคพฺยาปาเทน รุปฺปติ, ตํ เม ตณฺหาสลฺลํ ปหีนํ, อปนีโต อวิชฺชาวิสโทโส, สมฺมา นิพฺพานาธิมุตฺโตหมสฺมี”ติ สมฺมา นิพฺพานาธิมุตฺตสฺเสว สโต.\csromanspacer{} โส ยานิ สมฺมา นิพฺพานาธิมุตฺตสฺส อสปฺปายานิ, ตานิ นานุยุญฺเชยฺย, อสปฺปายํ จกฺขุนา รูปทสฺสนํ นานุยุญฺเชยฺย, อสปฺปายํ โสเตน สทฺทํ นานุยุญฺเชยฺย, อสปฺปายํ ฆาเนน คนฺธํ นานุยุญฺเชยฺย, อสปฺปายํ ชิวฺหาย รสํ นานุยุญฺเชยฺย, อสปฺปายํ กาเยน โผฏฺฐพฺพํ นานุยุญฺเชยฺย, อสปฺปายํ มนสา ธมฺมํ นานุยุญฺเชยฺย, ตสฺส อสปฺปายํ จกฺขุนา รูปทสฺสนํ นานุยุตฺตสฺส, อสปฺปายํ โสเตน สทฺทํ นานุยุตฺตสฺส, อสปฺปายํ ฆาเนน คนฺธํ นานุยุตฺตสฺส, อสปฺปายํ ชิวฺหาย รสํ นานุยุตฺตสฺส, อสปฺปายํ กาเยน โผฏฺฐพฺพํ นานุยุตฺตสฺส, อสปฺปายํ มนสา ธมฺมํ นานุยุตฺตสฺส ราโค จิตฺตํ นานุทฺธํเสยฺย, โส น ราคานุทฺธํสิเตน จิตฺเตน เนว มรณํ วา นิคจฺเฉยฺย น มรณมตฺตํ วา ทุกฺขํ.}

\proseitem{65}{อุปมา โข เม อยํ สุนกฺขตฺต กตา อตฺถสฺส วิญฺญาปนาย, อยํเยเวตฺถ อตฺโถ.\csromanspacer{} วโณติ โข สุนกฺขตฺต ฉนฺเนตํ อชฺฌตฺติกานํ อายตนานํ อธิวจนํ.\csromanspacer{} วิสโทโสติ โข สุนกฺขตฺต อวิชฺชาเยตํ อธิวจนํ.\csromanspacer{} สลฺลนฺติ โข สุนกฺขตฺต ตณฺหาเยตํ อธิวจนํ.\csromanspacer{} เอสนีติ โข สุนกฺขตฺต สติยาเยตํ อธิวจนํ.\csromanspacer{} สตฺถนฺติ โข สุนกฺขตฺต อริยาเยตํ ปญฺญาย อธิวจนํ.\csromanspacer{} ภิสกฺโก สลฺลกตฺโตติ โข สุนกฺขตฺต ตถาคตสฺเสตํ อธิวจนํ อรหโต สมฺมาสมฺพุทฺธสฺส.}

\prose{โส วต สุนกฺขตฺต ภิกฺขุ ฉสุ ผสฺสายตเนสุ สํวุตการี “อุปธิ ทุกฺขสฺส มูลนฺ”ติ อิติ วิทิตฺวา นิรุปธิ อุปธิสงฺขเย วิมุตฺโต อุปธิสฺมิํ วา กายํ อุปสํหริสฺสติ จิตฺตํ วา อุปฺปาเทสฺสตีติ เนตํ ฐานํ วิชฺชติ.\csromanspacer{} เสยฺยถาปิ สุนกฺขตฺต อาปานียกํโส วณฺณสมฺปนฺโน คนฺธสมฺปนฺโน รสสมฺปนฺโน, โส จ โข วิเสน สํสฏฺโฐ.\csromanspacer{} อถ ปุริโส อาคจฺเฉยฺย ชีวิตุกาโม อมริตุกาโม สุขกาโม ทุกฺขปฏิกูโล.\csromanspacer{} ตํ กิํ มญฺญสิ สุนกฺขตฺต, อปิ นุ โส ปุริโส อมุํ อาปานียกํสํ ปิเวยฺย, ยํ ชญฺญา “อิมาหํ ปิวิตฺวา มรณํ วา นิคจฺฉามิ มรณมตฺตํ วา ทุกฺขนฺ”ติ.\csromanspacer{} โน เหตํ ภนฺเต.\csromanspacer{} เอวเมว โข สุนกฺขตฺต โส วต ภิกฺขุ ฉสุ ผสฺสายตเนสุ สํวุตการี “อุปธิ ทุกฺขสฺส มูลนฺ”ติ อิติ วิทิตฺวา นิรุปธิ อุปธิสงฺขเย \csromanfolio{48}วิมุตฺโต อุปธิสฺมิํ วา กายํ อุปสํหริสฺสติ จิตฺตํ วา อุปฺปาเทสฺสตีติ เนตํ ฐานํ วิชฺชติ.\csromanspacer{} เสยฺยถาปิ สุนกฺขตฺต อาสีวิโส\footnote{อาสิวิโส (ก)} โฆรวิโส.\csromanspacer{} อถ ปุริโส อาคจฺเฉยฺย ชีวิตุกาโม อมริตุกาโม สุขกาโม ทุกฺขปฏิกูโล.\csromanspacer{} ตํ กิํ มญฺญสิ สุนกฺขตฺต, อปิ นุ โส ปุริโส อมุสฺส อาสีวิสสฺส โฆรวิสสฺส หตฺถํ วา องฺคุฏฺฐํ วา ทชฺชา\footnote{ยุญฺเชยฺย (ก)} ยํ ชญฺญา “อิมินาหํ ทฏฺโฐ มรณํ วา นิคจฺฉามิ มรณมตฺตํ วา ทุกฺขนฺ”ติ.\csromanspacer{} โน เหตํ ภนฺเต.\csromanspacer{} เอวเมว โข สุนกฺขตฺต โส วต ภิกฺขุ ฉสุ ผสฺสายตเนสุ สํวุตการี “อุปธิ ทุกฺขสฺส มูลนฺ”ติ อิติ วิทิตฺวา นิรุปธิ อุปธิสงฺขเย วิมุตฺโต อุปธิสฺมิํ วา กายํ อุปสํหริสฺสติ จิตฺตํ วา อุปฺปาเทสฺสตีติ เนตํ ฐานํ วิชฺชตีติ.}

\prose{อิทมโวจ ภควา.\csromanspacer{} อตฺตมโน สุนกฺขตฺโต ลิจฺฉวิปุตฺโต ภควโต ภาสิตํ อภินนฺทีติ.}

\nitthitamruled{สุนกฺขตฺตสุตฺตํ นิฏฺฐิตํ ปญฺจมํ.}
\csromanmatikaanchor{mtk.7}
\csromantocmarkref{subsection}{6. อาเนญฺชสปฺปายสุตฺต}{mtk.7}
\bookmark[dest={mtk.7},level=2]{6. อาเนญฺชสปฺปายสุตฺต}
\csromanheader{2}{6. อาเนญฺชสปฺปายสุตฺต}
\proseitem{66}{เอวํ เม สุตํ\csromandash{}เอกํ สมยํ ภควา กุรูสุ วิหรติ กมฺมาสธมฺมํ นาม กุรูนํ นิคโม.\csromanspacer{} ตตฺร โข ภควา ภิกฺขู อามนฺเตสิ “ภิกฺขโว”ติ. “ภทนฺเต”ติ เต ภิกฺขู ภควโต ปจฺจสฺโสสุํ.\csromanspacer{} ภควา เอตทโวจ\csromandash{}อนิจฺจา ภิกฺขเว กามา ตุจฺฉา มุสา โมสธมฺมา.\csromanspacer{} มายากตเม ตํ ภิกฺขเว พาลลาปนํ.\csromanspacer{} เย จ ทิฏฺฐธมฺมิกา กามา, เย จ สมฺปรายิกา กามา, ยา จ ทิฏฺฐธมฺมิกา กามสญฺญา, ยา จ สมฺปรายิกา กามสญฺญา, อุภยเมตํ มารเธยฺยํ, มารสฺเสส\footnote{มารสฺเสว (ก)} วิสโย, มารสฺเสส นิวาโป, มารสฺเสส โคจโร.\csromanspacer{} เอตฺเถเต ปาปกา อกุสลา มานสา อภิชฺฌาปิ พฺยาปาทาปิ สารมฺภาปิ สํวตฺตนฺติ, เตว อริยสาวกสฺส อิธมนุสิกฺขโต อนฺตรายาย สมฺภวนฺติ.\csromanspacer{} ตตฺร ภิกฺขเว อริยสาวโก อิติ ปฏิสญฺจิกฺขติ “เย จ ทิฏฺฐธมฺมิกา กามา, เย จ สมฺปรายิกา กามา, ยา จ ทิฏฺฐธมฺมิกา กามสญฺญา, ยา จ สมฺปรายิกา กามสญฺญา, อุภยเมตํ มารเธยฺยํ, มารสฺเสส วิสโย, มารสฺเสส นิวาโป, มารสฺเสส โคจโร.}

\vspace{\tipitakaparskip}
\csromancenter{อาเนญฺชสปฺปายสุตฺต (106) เอตฺเถเต ปาปกา อกุสลา มานสา อภิชฺฌาปิ พฺยาปาทาปิ สารมฺภาปิ สํวตฺตนฺติ, เตว อริยสาวกสฺส อิธมนุสิกฺขโต อนฺตรายาย สมฺภวนฺติ.\csromanspacer{} ยํนูนาหํ วิปุเลน มหคฺคเตน เจตสา วิหเรยฺยํ อภิภุยฺย โลกํ อธิฏฺฐาย มนสา.\csromanspacer{} วิปุเลน หิ เม มหคฺคเตน เจตสา วิหรโต อภิภุยฺย โลกํ อธิฏฺฐาย มนสา, เย ปาปกา อกุสลา มานสา อภิชฺฌาปิ พฺยาปาทาปิ สารมฺภาปิ, เต น ภวิสฺสนฺติ.\csromanspacer{} เตสํ ปหานา อปริตฺตญฺจ เม จิตฺตํ ภวิสฺสติ อปฺปมาณํ สุภาวิตนฺ”ติ.\csromanspacer{} ตสฺส เอวํปฏิปนฺนสฺส ตพฺพหุลวิหาริโน อายตเน จิตฺตํ ปสีทติ, สมฺปสาเท สติ เอตรหิ วา อาเนญฺชํ สมาปชฺชติ, ปญฺญาย วา อธิมุจฺจติ, กายสฺส เภทา ปรํ มรณา ฐานเมตํ วิชฺชติ, ยํ ตํสํวตฺตนิกํ วิญฺญาณํ อสฺส อาเนญฺชูปคํ.\csromanspacer{} อยํ ภิกฺขเว ปฐมา อาเนญฺชสปฺปายา ปฏิปทา อกฺขายติ.}
\proseitem{67}{\csromanfolio{49}ปุน จปรํ ภิกฺขเว อริยสาวโก อิติ ปฏิสญฺจิกฺขติ “เย จ ทิฏฺฐธมฺมิกา กามา, เย จ สมฺปรายิกา กามา, ยา จ ทิฏฺฐธมฺมิกา กามสญฺญา, ยา จ สมฺปรายิกา กามสญฺญา, ยํ กิญฺจิ รูปํ (สพฺพํ รูปํ) จตฺตาริ จ มหาภูตานิ จตุนฺนญฺจ มหาภูตานํ อุปาทายรูปนฺ”ติ.\csromanspacer{} ตสฺส เอวํปฏินฺนสฺส ตพฺพหุลวิหาริโน อายตเน จิตฺตํ ปสีทติ, สมฺปสาเท สติ เอตรหิ วา อาเนญฺชํ สมาปชฺชติ, ปญฺญาย วา อธิมุจฺจติ, กายสฺส เภทา ปรํ มรณา ฐานเมตํ วิชฺชติ, ยํ ตํสํวตฺตนิกํ วิญฺญาณํ อสฺส อาเนญฺชูปคํ.\csromanspacer{} อยํ ภิกฺขเว ทุติยา อาเนญฺชสปฺปายา ปฏิปทา อกฺขายติ.}

\prose{ปุน จปรํ ภิกฺขเว อริยสาวโก อิติ ปฏิสญฺจิกฺขติ “เย จ ทิฏฺฐธมฺมิกา กามา, เย จ สมฺปรายิกา กามา, ยา จ ทิฏฺฐธมฺมิกา กามสญฺญา, ยา จ สมฺปรายิกา กามสญฺญา, เย จ ทิฏฺฐธมฺมิกา รูปา, เย จ สมฺปรายิกา รูปา, ยา จ ทิฏฺฐธมฺมิกา รูปสญฺญา, ยา จ สมฺปรายิกา รูปสญฺญา, อุภยเมตํ อนิจฺจํ, ยทนิจฺจํ, ตํ นาลํ อภินนฺทิตุํ นาลํ อภิวทิตุํ นาลํ อชฺโฌสิตุนฺ”ติ.\csromanspacer{} ตสฺส เอวํปฏิปนฺนสฺส ตพฺพหุลวิหาริโน อายตเน จิตฺตํ ปสีทติ, สมฺปสาเท สติ เอตรหิ วา อาเนญฺชํ สมาปชฺชติ, ปญฺญาย วา อธิมุจฺจติ.\csromanspacer{} กายสฺส เภทา ปรํ มรณา ฐานเมตํ วิชฺชติ, ยํ ตํสํวตฺตนิกํ วิญฺญาณํ อสฺส อาเนญฺชูปคํ.\csromanspacer{} อยํ ภิกฺขเว ตติยา อาเนญฺชสปฺปายา ปฏิปทา อกฺขายติ.}

\proseitem{68}{\csromanfolio{50}ปุน จปรํ ภิกฺขเว อริยสาวโก อิติ ปฏิสญฺจิกฺขติ “เย จ ทิฏฺฐธมฺมิกา กามา, เย จ สมฺปรายิกา กามา, ยา จ ทิฏฺฐธมฺมิกา กามสญฺญา, ยา จ สมฺปรายิกา กามสญฺญา, เย จ ทิฏฺฐธมฺมิกา รูปา, เย จ สมฺปรายิกา รูปา, ยา จ ทิฏฺฐธมฺมิกา รูปสญฺญา, ยา จ สมฺปรายิกา รูปสญฺญา, ยา จ อาเนญฺชสญฺญา, สพฺพา สญฺญา, ยตฺเถตา อปริเสสา นิรุชฺฌนฺติ.\csromanspacer{} เอตํ สนฺตํ เอตํ ปณีตํ, ยทิทํ อากิญฺจญฺญายตนนฺ”ติ.\csromanspacer{} ตสฺส เอวํปฏิปนฺนสฺส ตพฺพหุลวิหาริโน อายตเน จิตฺตํ ปสีทติ, สมฺปสาเท สติ เอตรหิ วา อากิญฺจญฺญายตนํ สมาปชฺชติ, ปญฺญาย วา อธิมุจฺจติ, กายสฺส เภทา ปรํ มรณา ฐานเมตํ วิชฺชติ, ยํ ตํสํวตฺตนิกํ วิญฺญาณํ อสฺส อากิญฺจญฺญายตนูปคํ.\csromanspacer{} อยํ ภิกฺขเว ปฐมา อากิญฺจญฺญายตนสปฺปายา ปฏิปทา อกฺขายติ.}

\proseitem{69}{ปุน จปรํ ภิกฺขเว อริยสาวโก อรญฺญคโต วา รุกฺขมูลคโต วา สุญฺญาคารคโต วา อิติ ปฏิสญฺจิกฺขติ “สุญฺญมิทํ อตฺเตน วา อตฺตนิเยน วา”ติ.\csromanspacer{} ตสฺส เอวํปฏิปนฺนสฺส ตพฺพหุลวิหาริโน อายตเน จิตฺตํ ปสีทติ, สมฺปสาเท สติ เอตรหิ วา อากิญฺจญฺญายตนํ สมาปชฺชติ.\csromanspacer{} ปญฺญาย วา อธิมุจฺจติ, กายสฺส เภทา ปรํ มรณา ฐานเมตํ วิชฺชติ, ยํ ตํสํวตฺตนิกํ วิญฺญาณํ อสฺส อากิญฺจญฺญายตนูปคํ.\csromanspacer{} อยํ ภิกฺขเว ทุติยา อากิญฺจญฺญายตนสปฺปายา ปฏิปทา อกฺขายติ.}

\proseitem{70}{ปุน จปรํ ภิกฺขเว อริยสาวโก อิติ ปฏิสญฺจิกฺขติ “นาหํ กฺวจนิ\footnote{กฺวจินิ (สฺยา, กํ, สี-ฏฺฐ)} กสฺสจิ กิญฺจนตสฺมิํ\footnote{กิญฺจนตสฺมิ (?)}, น จ มม กฺวจนิ กิสฺมิญฺจิ กิญฺจนํ นตฺถี”ติ.\csromanspacer{} ตสฺส เอวํปฏิปนฺนสฺส ตพฺพหุลวิหาริโน อายตเน จิตฺตํ ปสีทติ, สมฺปสาเท สติ เอตรหิ วา อากิญฺจญฺญายตนํ สมาปชฺชติ, ปญฺญาย วา อธิมุจฺจติ, กายสฺส เภทา ปรํ มรณา ฐานเมตํ วิชฺชติ, ยํ ตํสํวตฺตนิกํ วิญฺญาณํ อสฺส อากิญฺจญฺญายตนูปคํ.\csromanspacer{} อยํ ภิกฺขเว ตติยา อากิญฺจญฺญายตนสปฺปายา ปฏิปทา อกฺขายติ.}

\prose{ปุน จปรํ ภิกฺขเว อริยสาวโก อิติ ปฏิสญฺจิกฺขติ “เย จ ทิฏฺฐธมฺมิกา กามา, เย จ สมฺปรายิกา กามา, ยา จ ทิฏฺฐธมฺมิกา กามสญฺญา, ยา จ สมฺปรายิกา กามสญฺญา, เย จ ทิฏฺฐธมฺมิกา รูปา, เย จ สมฺปรายิกา รูปา, ยา จ ทิฏฺฐธมฺมิกา รูปสญฺญา, ยา จ สมฺปรายิกา}

\vspace{\tipitakaparskip}
\csromancenter{อาเนญฺชสปฺปายสุตฺต (106) รูปสญฺญา, ยา จ อาเนญฺชสญฺญา, ยา จ อากิญฺจญฺญายตนสญฺญา, สพฺพา สญฺญา, ยตฺเถตา อปริเสสา นิรุชฺฌนฺติ.\csromanspacer{} เอตํ สนฺตํ เอตํ ปณีตํ ยทิทํ เนวสญฺญานาสญฺญายตนนฺ”ติ.\csromanspacer{} ตสฺส เอวํปฏิปนฺนสฺส ตพฺพหุลวิหาริโน อายตเน จิตฺตํ ปสีทติ, สมฺปสาเท สติ เอตรหิ วา เนวสญฺญานาสญฺญายตนํ สมาปชฺชติ, ปญฺญาย วา อธิมุจฺจติ, กายสฺส เภทา ปรํ มรณา ฐานเมตํ วิชฺชติ, ยํ ตํสํวตฺตนิกํ วิญฺญาณํ อสฺส เนวสญฺญานาสญฺญายตนูปคํ.\csromanspacer{} อยํ ภิกฺขเว เนวสญฺญานาสญฺญายตนสปฺปายา ปฏิปทา อกฺขายตีติ.}
\proseitem{71}{\csromanfolio{51}เอวํ วุตฺเต อายสฺมา อานนฺโท ภควนฺตํ เอตทโวจ “อิธ ภนฺเต ภิกฺขุ เอวํ ปฏิปนฺโน โหติ ‘โน จสฺส โน จ เม สิยา น ภวิสฺสติ น เม ภวิสฺสติ ยทตฺถิ ยํ ภูตํ, ตํ ปชหามี’ติ เอวํ อุเปกฺขํ ปฏิลภติ.\csromanspacer{} ปรินิพฺพาเยยฺย นุ โข โส ภนฺเต ภิกฺขุ น วา ปรินิพฺพาเยยฺยา”ติ.\csromanspacer{} อเปตฺเถกจฺโจ อานนฺท ภิกฺขุ ปรินิพฺพาเยยฺย, อเปตฺเถกจฺโจ ภิกฺขุ น ปรินิพฺพาเยยฺยาติ.\csromanspacer{} โก นุ โข ภนฺเต เหตุ โก ปจฺจโย, เยนเปตฺเถกจฺโจ ภิกฺขุ ปรินิพฺพาเยยฺย, อเปตฺเถกจฺโจ ภิกฺขุ น ปรินิพฺพาเยยฺยาติ.\csromanspacer{} อิธานนฺท ภิกฺขุ เอวํ ปฏิปนฺโน โหติ “โน จสฺส โน จ เม สิยา น ภวิสฺสติ น เม ภวิสฺสติ ยทตฺถิ ยํ ภูตํ, ตํ ปชหามี”ติ เอวํ อุเปกฺขํ ปฏิลภติ, โส ตํ อุเปกฺขํ อภินนฺทติ อภิวทติ อชฺโฌสาย ติฏฺฐติ, ตสฺส ตํ อุเปกฺขํ อภินนฺทโต อภิวทโต อชฺโฌสาย ติฏฺฐโต ตนฺนิสฺสิตํ โหติ วิญฺญาณํ, ตทุปาทานํ, ส-อุปาทาโน อานนฺท ภิกฺขุ น ปรินิพฺพายตีติ.\csromanspacer{} กหํ ปน โส ภนฺเต ภิกฺขุ อุปาทิยมาโน อุปาทิยตีติ.\csromanspacer{} เนวสญฺญานาสญฺญายตนํ อานนฺทาติ.\csromanspacer{} อุปาทานเสฏฺฐํ กิร โส ภนฺเต ภิกฺขุ อุปาทิยมาโน อุปาทิยตีติ.\csromanspacer{} อุปาทานเสฏฺฐํ หิ โส อานนฺท ภิกฺขุ อุปาทิยมาโน อุปาทิยติ, อุปาทานเสฏฺฐํ เหตํ อานนฺท ยทิทํ เนวสญฺญานาสญฺญายตนํ.}

\proseitem{72}{อิธานนฺท ภิกฺขุ เอวํ ปฏิปนฺโน โหติ “โน จสฺส โน จ เม สิยา น ภวิสฺสติ น เม ภวิสฺสติ ยทตฺถิ ยํ ภูตํ, ตํ ปชหามี”ติ เอวํ อุเปกฺขํ ปฏิลภติ, โส ตํ อุเปกฺขํ นาภินนฺทติ นาภิวทติ น อชฺโฌสาย ติฏฺฐติ, ตสฺส ตํ อุเปกฺขํ อนภินนฺทโต อนภิวทโต \csromanfolio{52}อนชฺโฌสาย ติฏฺฐโต น ตนฺนิสฺสิตํ โหติ วิญฺญาณํ น ตทุปาทานํ, อนุปาทาโน อานนฺท ภิกฺขุ ปรินิพฺพายตีติ.}

\proseitem{73}{อจฺฉริยํ ภนฺเต, อพฺภุตํ ภนฺเต, นิสฺสาย กิร โน ภนฺเต ภควตา โอฆสฺส นิตฺถรณา อกฺขาตา.\csromanspacer{} กตโม ปน ภนฺเต อริโย วิโมกฺโขติ.\csromanspacer{} อิธานนฺท ภิกฺขุ อริยสาวโก อิติ ปฏิสญฺจิกฺขติ “เย จ ทิฏฺฐธมฺมิกา กามา, เย จ สมฺปรายิกา กามา, ยา จ ทิฏฺฐธมฺมิกา กามสญฺญา, ยา จ สมฺปรายิกา กามสญฺญา, เย จ ทิฏฺฐธมฺมิกา รูปา, เย จ สมฺปรายิกา รูปา, ยา จ ทิฏฺฐธมฺมิกา รูปสญฺญา, ยา จ สมฺปรายิกา รูปสญฺญา, ยา จ อาเนญฺชสญฺญา, ยา จ อากิญฺจญฺญายตนสญฺญา, ยา จ เนวสญฺญานาสญฺญายตนสญฺญา, เอส สกฺกาโย, ยาวตา สกฺกาโย.\csromanspacer{} เอตํ อมตํ ยทิทํ อนุปาทา จิตฺตสฺส วิโมกฺโข.\csromanspacer{} อิติ โข อานนฺท เทสิตา มยา อาเนญฺชสปฺปายา ปฏิปทา, เทสิตา อากิญฺจญฺญายตนสปฺปายา ปฏิปทา, เทสิตา เนวสญฺญานาสญฺญายตนสปฺปายา ปฏิปทา, เทสิตา นิสฺสาย นิสฺสาย โอฆสฺส นิตฺถรณา, เทสิโต อริโย วิโมกฺโข.\csromanspacer{} ยํ โข อานนฺท สตฺถารา กรณียํ สาวกานํ หิเตสินา อนุกมฺปเกน อนุกมฺปํ อุปาทาย, กตํ โว ตํ มยา.\csromanspacer{} เอตานิ อานนฺท รุกฺขมูลานิ เอตานิ สุญฺญาคารานิ, ฌายถานนฺท มา ปมาทตฺถ, มา ปจฺฉา วิปฺปฏิสาริโน อหุวตฺถ.\csromanspacer{} อยํ โว อมฺหากํ อนุสาสนีติ.}

\prose{อิทมโวจ ภควา.\csromanspacer{} อตฺตมโน อายสฺมา อานนฺโท ภควโต ภาสิตํ อภินนฺทีติ.}

\nitthitamruled{อาเนญฺชสปฺปายสุตฺตํ นิฏฺฐิตํ ฉฏฺฐํ.}
\csromanmatikaanchor{mtk.8}
\csromantocmarkref{subsection}{7. คณกโมคฺคลฺลานสุตฺต}{mtk.8}
\bookmark[dest={mtk.8},level=2]{7. คณกโมคฺคลฺลานสุตฺต}
\csromanheader{2}{7. คณกโมคฺคลฺลานสุตฺต}
\proseitem{74}{เอวํ เม สุตํ\csromandash{}เอกํ สมยํ ภควา สาวตฺถิยํ วิหรติ ปุพฺพาราเม มิคารมาตุปาสาเท.\csromanspacer{} อถ โข คณกโมคฺคลฺลาโน\footnote{คณกโมคฺคลาโน (ก)} พฺราหฺมโณ เยน ภควา เตนุปงฺกมิ, อุปสงฺกมิตฺวา ภควตา สทฺธิํ สมฺโมทิ, สมฺโมทนียํ}

\vspace{\tipitakaparskip}
\csromancenter{คณกโมคฺคลฺลานสุตฺต (107) กถํ สารณียํ วีติสาเรตฺวา เอกมนฺตํ นิสีทิ, เอกมนฺตํ นิสินฺโน โข คณกโมคฺคลฺลาโน พฺราหฺมโณ ภควนฺตํ เอตทโวจ\csromandash{}}
\prose{\csromanfolio{53}เสยฺยถาปิ โภ โคตม อิมสฺส มิคารมาตุปาสาทสฺส ทิสฺสติ อนุปุพฺพสิกฺขา อนุปุพฺพกิริยา อนุปุพฺพปฏิปทา ยทิทํ ยาว ปจฺฉิมโสปานกเฬวรา.\csromanspacer{} อิเมสํปิ หิ โภ โคตม พฺราหฺมณานํ ทิสฺสติ อนุปุพฺพสิกฺขา อนุปุพฺพกิริยา อนุปุพฺพปฏิปทา ยทิทํ อชฺเฌเน.\csromanspacer{} อิเมสํปิ หิ โภ โคตม อิสฺสาสานํ ทิสฺสติ อนุปุพฺพสิกฺขา อนุปุพฺพกิริยา อนุปุพฺพปฏิปทา ยทิทํ อิสฺสตฺเถ\footnote{อิสฺสตฺเต (ก)}.\csromanspacer{} อมฺหากํปิ หิ โภ โคตม คณกานํ คณนาชีวานํ ทิสฺสติ อนุปุพฺพสิกฺขา อนุปุพฺพกิริยา อนุปุพฺพปฏิปทา ยทิทํ สงฺขาเน.\csromanspacer{} มยํ หิ โภ โคตม อนฺเตวาสิํ ลภิตฺวา ปฐมํ เอวํ คณาเปม “เอกํ เอกกํ, เทฺว ทุกา, ตีณิ ติกา, จตฺตาริ จตุกฺกา, ปญฺจ ปญฺจกา, ฉ ฉกฺกา, สตฺต สตฺตกา, อฏฺฐ อฏฺฐกา, นว นวกา, ทส ทสกา”ติ.\csromanspacer{} สตมฺปิ มยํ โภ โคตม คณาเปม, ภิยฺโยปิ คณาเปม.\csromanspacer{} สกฺกา นุ โข โภ โคตม อิมสฺมิมฺปิ ธมฺมวินเย เอวเมว อนุปุพฺพสิกฺขา อนุปุพฺพกิริยา อนุปุพฺพปฏิปทา ปญฺญเปตุนฺติ.}

\proseitem{75}{สกฺกา พฺราหฺมณ อิมสฺมิมฺปิ ธมฺมวินเย อนุปุพฺพสิกฺขา อนุปุพฺพกิริยา อนุปุพฺพปฏิปทา ปญฺญเปตุํ.\csromanspacer{} เสยฺยถาปิ พฺราหฺมณ ทกฺโข อสฺสทมฺมโก ภทฺทํ อสฺสาชานียํ ลภิตฺวา ปฐเมเนว มุขาธาเน การณํ กาเรติ, อถ อุตฺตริํ การณํ กาเรติ.\csromanspacer{} เอวเมว โข พฺราหฺมณ ตถาคโต ปุริสทมฺมํ ลภิตฺวา ปฐมํ เอวํ วิเนติ “เอหิ ตฺวํ ภิกฺขุ สีลวา โหหิ, ปาติโมกฺขสํวรสํวุโต วิหราหิ อาจารโคจสมฺปนฺโน, อณุมตฺเตสุ วชฺเชสุ ภยทสฺสาวี สมาทาย สิกฺขสฺสุ สิกฺขาปเทสู”ติ.}

\prose{ยโต โข พฺราหฺมณ ภิกฺขุ สีลวา โหติ, ปาติโมกฺขสํวรสํวุโต วิหรติ อาจารโคจรสมฺปนฺโน, อณุมตฺเตสุ วชฺเชสุ ภยทสฺสาวี สมาทาย สิกฺขติ สิกฺขาปเทสุ.\csromanspacer{} ตเมนํ ตถาคโต อุตฺตริํ วิเนติ “เอหิ ตฺวํ ภิกฺขุ อินฺทฺริเยสุ คุตฺตทฺวาโร โหหิ, จกฺขุนา รูปํ ทิสฺวา มา นิมิตฺตคฺคาหี โหหิ มานุพฺยญฺชนคฺคาหี, ยตฺวาธิกรณเมนํ จกฺขุนฺทฺริยํ อสํวุตํ วิหรนฺตํ อภิชฺฌาโทมนสฺสา ปาปกา อกุสลา ธมฺมา อนฺวาสฺสเวยฺยุํ, ตสฺส สํวราย ปฏิปชฺชาหิ, รกฺขาหิ จกฺขุนฺทฺริยํ, จกฺขุนฺทฺริเย \csromanfolio{54}สํวรํ อาปชฺชาหิ.\csromanspacer{} โสเตน สทฺทํ สุตฺวา ฯเปฯ.\csromanspacer{} ฆาเนน คนฺธํ ฆายิตฺวา ฯเปฯ.\csromanspacer{} ชิวฺหาย รสํ สายิตฺวา ฯเปฯ.\csromanspacer{} กาเยน โผฏฺฐพฺพํ ผุสิตฺวา ฯเปฯ.\csromanspacer{} มนสา ธมฺมํ วิญฺญาย มา นิมิตฺตคฺคาหี โหหิ มานุพฺยญฺชนคฺคาหี, ยตฺวาธิกรณเมนํ มนินฺทฺริยํ อสํวุตํ วิหรนฺตํ อภิชฺฌาโทมนสฺสา ปาปกา อกุสลา ธมฺมา อนฺวาสฺสเวยฺยุํ, ตสฺส สํวราย ปฏิปชฺชาหิ, รกฺขาหิ มนินฺทฺริยํ, มนินฺทฺริเย สํวรํ อาปชฺชาหี”ติ.}

\prose{ยโต โข พฺราหฺมณ ภิกฺขุ อินฺทฺริเยสุ คุตฺตทฺวาโร โหติ, ตเมนํ ตถาคโต อุตฺตริํ วิเนติ “เอหิ ตฺวํ ภิกฺขุ โภชเน มตฺตญฺญู โหหิ, ปฏิสงฺขา โยนิโส อาหารํ อาหาเรยฺยาสิ เนว ทวาย น มทาย น มณฺฑนาย น วิภูสนาย, ยาวเทว อิมสฺส กายสฺส ฐิติยา ยาปนาย วิหิํสูปรติยา พฺรหฺมจริยานุคฺคหาย.\csromanspacer{} อิติ ปุราณญฺจ เวทนํ ปฏิหงฺขามิ, นวญฺจ เวทนํ น อุปฺปาเทสฺสามิ, ยาตฺรา จ เม ภวิสฺสติ อนวชฺชตา จ ผาสุวิหาโร จา”ติ.}

\prose{ยโต โข พฺราหฺมณ ภิกฺขุ โภชเน มตฺตญฺญู โหติ, ตเมนํ ตถาคโต อุตฺตริํ วิเนติ “เอหิ ตฺวํ ภิกฺขุ ชาคริยํ อนุยุตฺโต วิหราหิ, ทิวสํ จงฺกเมน นิสชฺชาย อาวรณีเยหิ ธมฺเมหิ จิตฺตํ ปริโสเธหิ, รตฺติยา ปฐมํ ยามํ จงฺกเมน นิสชฺชาย อาวรณีเยหิ ธมฺเมหิ จิตฺตํ ปริโสเธหิ, รตฺติยา มชฺฌิมํ ยามํ ทกฺขิเณน ปสฺเสน สีหเสยฺยํ กปฺเปยฺยาสิ ปาเท ปาทํ อจฺจาธาย สโต สมฺปชาโน อุฏฺฐานสญฺญํ มนสิกริตฺวา, รตฺติยา ปจฺฉิมํ ยามํ ปจฺจุฏฺฐาย จงฺกเมน นิสชฺชาย อาวรณีเยหิ ธมฺเมหิ จิตฺตํ ปริโสเธหี”ติ.}

\prose{ยโต โข พฺราหฺมณ ภิกฺขุ ชาคริยํ อนุยุตฺโต โหติ, ตเมนํ ตถาคโต อุตฺตริํ วิเนติ “เอหิ ตฺวํ ภิกฺขุ สติสมฺปชญฺเญน สมนฺนาคโต โหติ, อภิกฺกนฺเต ปฏิกฺกนฺเต สมฺปชานการี, อาโลกิเต วิโลกิเต สมฺปชานการี, สมิญฺชิเต ปสาริเต สมฺปชานการี, สงฺฆาฏิปตฺตจีวรธารเณ สมฺปชานการี, อสิเต ปีเต ขายิเต สายิเต สมฺปชานการี, อุจฺจารปสฺสาวกมฺเม สมฺปชานการี, คเต ฐิเต นิสินฺเน สุตฺเต ชาคริเต ภาสิเต ตุณฺหีภาเว สมฺปชานการี”ติ.}

\prose{ยโต โข พฺราหฺมณ ภิกฺขุ สติสมฺปชญฺเญน สมนฺนาคโต โหติ, ตเมนํ ตถาคโต อุตฺตริํ วิเนติ “เอหิ ตฺวํ ภิกฺขุ วิวิตฺตํ เสนาสนํ}

\vspace{\tipitakaparskip}
\csromancenter{คณกโมคฺคลฺลานสุตฺต (107) ภชาหิ อรญฺญํ รุกฺขมูลํ ปพฺพตํ กนฺทรํ คิริคุหํ สุสานํ วนปตฺถํ อพฺโภกาสํ ปลาลปุญฺชนฺ”ติ.\csromanspacer{} โส วิวิตฺตํ เสนาสนํ ภชติ อรญฺญํ รุกฺขมูลํ ปพฺพตํ กนฺทรํ คิริคุหํ สุสานํ วนปฺปตฺถํ อพฺโภกาสํ ปลาลปุญฺชํ, โส ปจฺฉาภตฺตํ ปิณฺฑปาตปฏิกฺกนฺโต นิสีทติ ปลฺลงฺกํ อาภุชิตฺวา อุชุํ กายํ ปณิธาย ปริมุขํ สติํ อุปฏฺฐเปตฺวา, โส อภิชฺฌํ โลเก ปหาย วิคตาภิชฺเฌน เจตสา วิหรติ, อภิชฺฌาย จิตฺตํ ปริโสเธติ.\csromanspacer{} พฺยาปาทปโทสํ ปหาย อพฺยาปนฺนจิตฺโต วิหรติ สพฺพปาณภูตหิตานุกมฺปี, พฺยาปาทปโทสา จิตฺตํ ปริโสเธติ.\csromanspacer{} ถินมิทฺธํ\footnote{ถีนมิทฺธํ (สี, สฺยา, กํ, อิ)} ปหาย วิคตถินมิทฺโธ วิหรติ อาโลกสญฺญี สโต สมฺปชาโน, ถินมิทฺธา จิตฺตํ ปริโสเธติ.\csromanspacer{} อุทฺธจฺจกุกฺกุจฺจํ ปหาย อนุทฺธโต วิหรติ อชฺฌตฺตํ วูปสนฺตจิตฺโต, อุทฺธจฺจกุกฺกุจฺจา จิตฺตํ ปริโสเธติ.\csromanspacer{} วิจิกิจฺฉํ ปหาย ติณฺณวิจิกิจฺโฉ วิหรติ อกถํกถี กุสเลสุ ธมฺเมสุ, วิจิกิจฺฉาย จิตฺตํ ปริโสเธติ.}
\proseitem{76}{\csromanfolio{55}โส อิเม ปญฺจ นีวรเณ ปหาย เจตโส อุปกฺกิเลเส ปญฺญาย ทุพฺพลีกรเณ วิวิจฺเจว กาเมหิ วิวิจฺจ อกุสเลหิ ธมฺเมหิ สวิตกฺกํ สวิจารํ วิเวกชํ ปีติสุขํ ปฐมํ ฌานํ อุปสมฺปชฺช วิหรติ.\csromanspacer{} วิตกฺกวิจารานํ วูปสมา อชฺฌตฺตํ สมฺปสาทนํ ฯเปฯ ทุติยํ ฌานํ อุปสมฺปชฺช วิหรติ.\csromanspacer{} ปีติยา จ วิราคา ฯเปฯ ตติยํ ฌานํ อุปสมฺปชฺช วิหรติ.\csromanspacer{} สุขสฺส จ ปหานา ฯเปฯ จตุตฺถํ ฌานํ อุปสมฺปชฺช วิหรติ.}

\prose{เย โข เต พฺราหฺมณ ภิกฺขู เสกฺขา\footnote{เสขา (สพฺพตฺถ)} อปตฺตมานสา อนุตฺตรํ โยคกฺเขมํ ปตฺถยมานา วิหรนฺติ, เตสุ เม อยํ เอวรูปี อนุสาสนี โหติ.\csromanspacer{} เย ปน เต ภิกฺขู อรหนฺโต ขีณาสวา วุสิตวนฺโต กตกรณียา โอหิตภารา อนุปฺปตฺตสทตฺถา ปริกฺขีณภวสํโยชนา สมฺมทญฺญา วิมุตฺตา, เตสํ อิเม ธมฺมา ทิฏฺฐธมฺมสุขวิหาราย เจว สํวตฺตนฺติ สติสมฺปชญฺญาย จาติ.}

\prose{เอวํ วุตฺเต คณกโมคฺคลฺลาโน พฺราหฺมโณ ภควนฺตํ เอตทโวจ “กิํ นุ โข โภโต โคตมสฺส สาวกา โภตา โคตเมน เอวํ โอวทียมานา เอวํ อนุสาสียมานา สพฺเพ อจฺจนฺตํ นิฏฺฐํ นิพฺพานํ อาราเธนฺติ, อุทาหุ เอกจฺเจ นาราเธนฺตี”ติ.\csromanspacer{} อปฺเปกจฺเจ โข พฺราหฺมณ มม สาวกา \csromanfolio{56}มยา เอวํ โอวทีมานา เอวํ อนุสาสียมานา อจฺจนฺตํ นิฏฺฐํ นิพฺพานํ อาราเธนฺติ, เอกจฺเจ นาราเธนฺตีติ.}

\prose{โก นุ โข โภ โคตม เหตุ โก ปจฺจโย, ยํ ติฏฺฐเตว นิพฺพานํ, ติฏฺฐติ นิพฺพานคามี มคฺโค, ติฏฺฐติ ภวํ โคตโม สมาทเปตา.\csromanspacer{} อถ จ ปน โภโต โคตมสฺส สาวกา โภตา โคตเมน เอวํ โอวทียมานา เอวํ อนุสาสียมานา อปฺเปกจฺเจ อจฺจนฺตํ นิฏฺฐํ นิพฺพานํ อาราเธนฺติ, เอกจฺเจ นาราเธนฺตีติ.}

\proseitem{77}{เตน หิ พฺราหฺมณ ตํเยเวตฺถ ปฏิปุจฺฉิสฺสามิ, ยถา เต ขเมยฺย, ตถา นํ พฺยากเรยฺยาสิ.\csromanspacer{} ตํ กิํ มญฺญสิ พฺราหฺมณ, กุสโล ตฺวํ ราชคหคามิสฺส มคฺคสฺสาติ.\csromanspacer{} เอวํ โภ กุสโล อหํ ราชคหคามิสฺส มคฺคสฺสาติ.\csromanspacer{} ตํ กิํ มญฺญสิ พฺราหฺมณ, อิธ ปุริโส อาคจฺเฉยฺย ราชคหํ คนฺตุกาโม, โส ตํ อุปสงฺกมิตฺวา เอวํ วเทยฺย “อิจฺฉามหํ ภนฺเต ราชคหํ คนฺตุํ, ตสฺส เม ราชคหสฺส มคฺคํ อุปทิสา”ติ.\csromanspacer{} ตเมนํ ตฺวํ เอวํ วเทยฺยาสิ “เอหมฺโภ\footnote{เอวํ โภ (สี, อิ)} ปุริส อยํ มคฺโค ราชคหํ คจฺฉติ, เตน มุหุตฺตํ คจฺฉ, เตน มุหุตฺตํ คนฺตฺวา ทกฺขิสฺสสิ อมุกํ นาม คามํ, เตน มุหุตฺตํ คจฺฉ, เตน มุหุตฺตํ คนฺตฺวา ทกฺขิสฺสสิ อมุกํ นาม นิคมํ, เตน มุหุตฺตํ คจฺฉ, เตน มุหุตฺตํ คนฺตฺวา ทกฺขิสฺสสิ ราชคหสฺส อารามรามเณยฺยกํ วนรามเณยฺยกํ ภูมิรามเณยฺยกํ โปกฺขรณีรามเณยฺยกนฺ”ติ.\csromanspacer{} โส ตยา เอวํ โอวทียมาโน เอวํ อนุสาสียมาโน อุมฺมคฺคํ คเหตฺวา ปจฺฉามุโข คจฺเฉยฺย.\csromanspacer{} อถ ทุติโย ปุริโส อาคจฺเฉยฺย ราชคหํ คนฺตุกาโม, โส ตํ อุปสงฺกมิตฺวา เอวํ วเทยฺย “อิจฺฉามหํ ภนฺเต ราชคหํ คนฺตุํ, ตสฺส เม ราชคหสฺส มคฺคํ อุปทิสา”ติ.\csromanspacer{} ตเมนํ ตฺวํ เอวํ วเทยฺยาสิ “เอหมฺโภ ปุริส อยํ มคฺโค ราชคหํ คจฺฉติ, เตน มุหุตฺตํ คจฺฉ, เตน มุหุตฺตํ คนฺตฺวา ทกฺขิสฺสสิ อมุกํ นาม คามํ, เตน มุหุตฺตํ คจฺฉ, เตน มุหุตฺตํ คนฺตฺวา ทกฺขิสฺสสิ อมุกํ นาม นิคมํ, เตน มุหุตฺตํ คจฺฉ, เตน มุหุตฺตํ คนฺตฺวา ทกฺขิสฺสสิ ราชคหสฺส อารามรามเณยฺยกํ วนรามเณยฺยกํ ภูมิรามเณยฺยกํ โปกฺขรณีรามเณยฺยกนฺ”ติ.\csromanspacer{} โส ตยา เอวํ โอวทียมาโน เอวํ อนุสาสียมาโน โสตฺถินา ราชคหํ คจฺเฉยฺย.\csromanspacer{} โก นุ โข พฺราหฺมณ เหตุ โก ปจฺจโย, ยํ ติฏฺฐเตว}

\vspace{\tipitakaparskip}
\csromancenter{คณกโมคฺคลฺลานสุตฺต (107) ราชคหํ, ติฏฺฐติ ราชคหคามี มคฺโค, ติฏฺฐสิ ตฺวํ สมาทเปตา.\csromanspacer{} อถ จ ปน ตยา เอวํ โอวทียมาโน เอวํ อนุสาสียมาโน เอโก ปุริโส อุมฺมคฺคํ คเหตฺวา ปจฺฉามุโข คจฺเฉยฺย, เอโก โสตฺถินา ราชคหํ คจฺเฉยฺยาติ.\csromanspacer{} เอตฺถ กฺยาหํ โภ โคตม กโรมิ มคฺคกฺขายีหํ โภ โคตมาติ.}
\prose{\csromanfolio{57}เอวเมว โข พฺราหฺมณ ติฏฺฐเตว นิพฺพานํ, ติฏฺฐติ นิพฺพานคามี มคฺโค, ติฏฺฐามหํ สมาทเปตา.\csromanspacer{} อถ จ ปน มม สาวกา มยา เอวํ โอวทียมานา เอวํ อนุสาสียมานา อปฺเปกจฺเจ อจฺจนฺตํ นิฏฺฐํ นิพฺพานํ อาราเธนฺติ, เอกจฺเจ นาราเธนฺติ.\csromanspacer{} เอตฺถ กฺยาหํ พฺราหฺมณ กโรมิ มคฺคกฺขายีหํ พฺราหฺมณ ตถาคโตติ.}

\proseitem{78}{เอวํ วุตฺเต คณกโมคฺคลฺลาโน พฺราหฺมโณ ภควนฺตํ เอตทโวจ “เยเม โภ โคตม ปุคฺคลา อสฺสทฺธา ชีวิกตฺถา น สทฺธา อคารสฺมา อนคาริยํ ปพฺพชิตา สฐา มายาวิโน เกตพิโน\footnote{เกฏุภิโน (สี, สฺยา, กํ, อิ)} อุทฺธตา อุนฺนฬา จปลา มุขรา วิกิณฺณวาจา อินฺทฺริเยสุ อคุตฺตทฺวารา โภชเน อมตฺตญฺญุโน ชาคริยํ อนนุยุตฺตา สามญฺเญ อนเปกฺขวนฺโต สิกฺขาย น ติพฺพคารวา พาหุลิกา\footnote{พาหุลฺลิกา (สฺยา, กํ)} สาถลิกา โอกฺกมเน ปุพฺพงฺคมา ปวิเวเก นิกฺขิตฺตธุรา กุสีตา หีนวีริยา มุฏฺฐสฺสติโน อสมฺปชานา อสมาหิตา วิพฺภนฺตจิตฺตา ทุปฺปญฺญา เอฬมูคา, น เตหิ ภวํ โคตโม สทฺธิํ สํวสติ.}

\prose{เย ปน เต กุลปุตฺตา สทฺธา อคารสฺมา อนคาริยํ ปพฺพชิตา อสฐา อมายาวิโน อเกตพิโน อนุทฺธตา อนุนฺนฬา อจปลา อมุขรา อวิกิณฺณวาจา อินฺทฺริเยสุ คุตฺตทฺวารา โภชเน มตฺตญฺญุโน ชาคริยํ อนุยุตฺตา สามญฺเญ อเปกฺขวนฺโต สิกฺขาย ติพฺพคารวา นพาหุลิกา นสาถลิกา โอกฺกมเน นิกฺขิตฺตธุรา ปวิเวเก ปุพฺพงฺคมา อารทฺธวีริยา ปหิตตฺตา อุปฏฺฐิตสฺสติโน สมฺปชานา สมาหิตา เอกคฺคจิตฺตา ปญฺญวนฺโต อเนฬมูคา, เตหิ ภวํ โคตโม สทฺธิํ สํวสติ.}

\prose{เสยฺยถาปิ โภ โคตม เย เกจิ มูลคนฺธา, กาลานุสาริ เตสํ อคฺคมกฺขายติ.\csromanspacer{} เย เกจิ สารคนฺธา, โลหิตจนฺทนํ เตสํ อคฺคมกฺขายติ.\csromanspacer{} เย เกจิ ปุปฺผคนฺธา, วสฺสิกํ เตสํ อคฺคมกฺขายติ.\csromanspacer{} เอวเมว โภโต โคตมสฺส โอวาโท ปรมชฺชธมฺเมสุ.}

\prose{\csromanfolio{58}อภิกฺกนฺตํ โภ โคตม, อภิกฺกนฺตํ โภ โคตม, เสยฺยถาปิ โภ โคตม นิกฺกุชฺชิตํ วา อุกฺกุชฺเชยฺย, ปฏิจฺฉนฺนํ วา วิวเรยฺย, มูฬฺหสฺส วา มคฺคํ อาจิกฺเขยฺย, อนฺธกาเร วา เตลปชฺโชตํ ธาเรยฺย “จกฺขุมนฺโต รูปานิ ทกฺขนฺตี”ติ, เอวเมวํ โภตา โคตเมน อเนกปริยาเยน ธมฺโม ปกาสิโต, เอสาหํ ภวนฺตํ โคตมํ สรณํ คจฺฉามิ ธมฺมญฺจ ภิกฺขุสํฆญฺจ, อุปาสกํ มํ ภวํ โคตโม ธาเรตุ อชฺชตคฺเค ปาณุเปตํ สรณํ คตนฺติ.}

\nitthitamruled{คณกโมคฺคลฺลานสุตฺตํ นิฏฺฐิตํ สตฺตมํ.}
\csromanmatikaanchor{mtk.9}
\csromantocmarkref{subsection}{8. โคปกโมคฺคลฺลานสุตฺต}{mtk.9}
\bookmark[dest={mtk.9},level=2]{8. โคปกโมคฺคลฺลานสุตฺต}
\csromanheader{2}{8. โคปกโมคฺคลฺลานสุตฺต}
\proseitem{79}{เอวํ เม สุตํ\csromandash{}เอกํ สมยํ อายสฺมา อานนฺโท ราชคเห วิหรติ เวฬุวเน กลนฺทกนิวาเป อจิรปนิพฺพุเต ภควติ.\csromanspacer{} เตน โข ปน สมเยน ราชา มาคโธ อชาตสตฺตุ เวเทหิปุตฺโต ราชคหํ ปฏิสงฺขาราเปติ รญฺโญ ปชฺโชตสฺส อาสงฺกมาโน.\csromanspacer{} อถ โข อายสฺมา อานนฺโท ปุพฺพณฺหสมยํ นิวาเสตฺวา ปตฺตจีวรมาทาย ราชคหํ ปิณฺฑาย ปาวิสิ.\csromanspacer{} อถ โข อายสฺมโต อานนฺทสฺส เอตทโหสิ “อติปฺปโค โข ตาว ราชคเห ปิณฺฑาย จริตุํ, ยํนูนาหํ เยน โคปกโมคฺคลฺลานสฺส พฺราหฺมณสฺส กมฺมนฺโต, เยน โคปกโมคฺคลฺลาโน พฺราหฺมโณ เตนุปสงฺกเมยฺยนฺ”ติ.}

\prose{อถ โข อายสฺมา อานนฺโท เยน โคปกโมคฺคลฺลานสฺส พฺราหฺมณสฺส กมฺมนฺโต, เยน โคปกโมคฺคลฺลาโน พฺราหฺมโณ เตนุปสงฺกมิ.\csromanspacer{} อทฺทสา โข โคปกโมคฺคลฺลาโน พฺราหฺมโณ อายสฺมนฺตํ อานนฺทํ ทูรโตว อาคจฺฉนฺตํ, ทิสฺวาน อายสฺมนฺตํ อานนฺทํ เอตทโวจ “เอตุ โข ภวํ อานนฺโท, สฺวาคตํ โภโต อานนฺทสฺส, จิรสฺสํ โข ภวํ อานนฺโท อิมํ ปริยายมกาสิ ยทิทํ อิธาคมนาย, นิสีทตุ ภวํ อานนฺโท, อิทมาสนํ ปญฺญตฺตนฺ”ติ.\csromanspacer{} นิสีทิ โข อายสฺมา อานนฺโท ปญฺญตฺเต อาสเน.\csromanspacer{} โคปกโมคฺคลฺลาโนปิ โข พฺราหฺมโณ อญฺญตรํ นีจํ อาสนํ คเหตฺวา เอกมนฺตํ นิสีทิ, เอกมนฺตํ นิสินฺโน โข โคปกโมคฺคลฺลาโน พฺราหฺมโณ อายสฺมนฺตํ อานนฺทํ เอตทโวจ “อตฺถิ นุ โข โภ อานนฺท เอกภิกฺขุปิ}

\vspace{\tipitakaparskip}
\csromancenter{โคปกโมคฺคลฺลานสุตฺต (108) เตหิ ธมฺเมหิ สพฺเพนสพฺพํ สพฺพถาสพฺพํ สมนฺนาคโต, เยหิ ธมฺเมหิ สมนฺนาคโต โส ภวํ โคตโม อโหสิ อรหํ สมฺมาสมฺพุทฺโธ”ติ.\csromanspacer{} นตฺถิ โข พฺราหฺมณ เอกภิกฺขุปิ เตหิ ธมฺเมหิ สพฺเพนสพฺพํ สพฺพถาสพฺพํ สมนฺนาคโต, เยหิ ธมฺเมหิ สมนฺนาคโต โส ภควา อโหสิ อรหํ สมฺมาสมฺพุทฺโธ.\csromanspacer{} โส หิ พฺราหฺมณ ภควา อนุปฺปนฺนสฺส มคฺคสฺส อุปฺปาเทตา, อสญฺชาตสฺส มคฺคสฺส สญฺชเนตา, อนกฺขาตสฺส มคฺคสฺส อกฺขาตา, มคฺคญฺญู มคฺควิทู มคฺคโกวิโท.\csromanspacer{} มคฺคานุคา จ ปน เอตรหิ สาวกา วิหรนฺติ ปจฺฉา สมนฺนาคตาติ.\csromanspacer{} อยญฺจ หิทํ อายสฺมโต อานนฺทสฺส โคปกโมคฺคลฺลาเนน พฺราหฺมเณน สทฺธิํ อนฺตรากถา วิปฺปกตา อโหสิ.}
\prose{\csromanfolio{59}อถ โข วสฺสกาโร พฺราหฺมโณ มคธมหามตฺโต ราชคเห กมฺมนฺเต อนุสญฺญายมาโน เยน โคปกโมคฺคลฺลานสฺส พฺราหฺมณสฺส กมฺมนฺโต, เยนายสฺมา อานนฺโท เตนุปสงฺกมิ, อุปสงฺกมิตฺวา อายสฺมตา อานนฺเทน สทฺธิํ สมฺโมทิ, สมฺโมทนียํ กถํ สารณียํ วีติสาเรตฺวา เอกมนฺตํ นิสีทิ, เอกมนฺตํ นิสินฺโน โข วสฺสกาโร พฺราหฺมโณ มคธมหามตฺโต อายสฺมนฺตํ อานนฺทํ เอตทโวจ “กาย นุตฺถ โภ อานนฺท เอตรหิ กถาย สนฺนิสินฺนา, กา จ ปน โว อนฺตรากถา วิปฺปกตา”ติ.\csromanspacer{} อิธ มํ พฺราหฺมณ โคปกโมคฺคลฺลาโน พฺราหฺมโณ เอวมาห “อตฺถิ นุ โข โภ อานนฺท เอกภิกฺขุปิ เตหิ ธมฺเมหิ สพฺเพนสพฺพํ สพฺพถาสพฺพํ สมนฺนาคโต, เยหิ ธมฺเมหิ สมนฺนาคโต โส ภวํ โคตโม อโหสิ อรหํ สมฺมาสมฺพุทฺโธ”ติ.\csromanspacer{} เอวํ วุตฺเต อหํ พฺราหฺมณ โคปกโมคฺคลฺลานํ พฺราหฺมณํ เอตทโวจํ “นตฺถิ โข พฺราหฺมณ เอกภิกฺขุปิ เตหิ ธมฺเมหิ สพฺเพนสพฺพํ สพฺพถาสพฺพํ สมนฺนาคโต, เยหิ ธมฺเมหิ สมนฺนาคโต โส ภควา อโหสิ อรหํ สมฺมาสมฺพุทฺโธ.\csromanspacer{} โส หิ พฺราหฺมณ ภควา อนุปฺปนฺนสฺส มคฺคสฺส อุปฺปาเทตา, อสญฺชาตสฺส มคฺคสฺส สญฺชเนตา, อนกฺขาตสฺส มคฺคสฺส อกฺขาตา, มคฺคญฺญู มคฺควิทู มคฺคโกวิโท.\csromanspacer{} มคฺคานุคา จ ปน เอตรหิ สาวกา วิหรนฺติ ปจฺฉา สมนฺนาคตา”ติ.\csromanspacer{} อยํ โข โน พฺราหฺมณ โคปกโมคฺคลฺลาเนน พฺราหฺมเณน สทฺธิํ อนฺตรากถา วิปฺปกตา, อถ ตฺวํ อนุปฺปตฺโตติ.}

\proseitem{80}{อตฺถิ นุ โข โภ อานนฺท เอกภิกฺขุปิ เตน โภตา โคตเมน ฐปิโต “อยํ โว มมจฺจเยน ปฏิสรณํ ภวิสฺสตี”ติ, ยํ ตุเมฺห \csromanfolio{60}เอตรหิ ปฏิปาเทยฺยาถาติ\footnote{ปฏิธาเวยฺยาถาติ (สี, สฺยา, กํ, อิ)}.\csromanspacer{} นตฺถิ โข พฺราหฺมณ เอกภิกฺขุปิ เตน ภควตา ชานตา ปสฺสตา อรหตา สมฺมาสมฺพุทฺเธน ฐปิโต “อยํ โว มมจฺจเยน ปฏิสรณํ ภวิสฺสตี”ติ, ยํ มยํ เอตรหิ ปฏิปาเทยฺยามาติ.\csromanspacer{} อตฺถิ ปน โภ อานนฺท เอกภิกฺขุปิ สํเฆน สมฺมโต สมฺพหุเลหิ เถเรหิ ภิกฺขูหิ ฐปิโต “อยํ โน ภควโต อจฺจเยน ปฏิสรณํ ภวิสฺสตี”ติ, ยํ ตุเมฺห เอตรหิ ปฏิปาเทยฺยาถาติ.\csromanspacer{} นตฺถิ โข พฺราหฺมณ เอกภิกฺขุปิ สํเฆน สมฺมโต สมฺพหุเลหิ เถเรหิ ภิกฺขูหิ ฐปิโต “อยํ โน ภควโต อจฺจเยน ปฏิสรณํ ภวิสฺสตี”ติ, ยํ มยํ เอตรหิ ปฏิปาเทยฺยามาติ.\csromanspacer{} เอวํ อปฺปฏิสรเณ จ ปน โภ อานนฺท โก เหตุ สามคฺคิยาติ.\csromanspacer{} น โข มยํ พฺราหฺมณ อปฺปฏิสรณา, สปฺปฏิสรณา มยํ พฺราหฺมณ ธมฺมปฺปฏิสรณาติ.}

\prose{“อตฺถิ นุ โข โภ อานนฺท เอกภิกฺขุปิ เตน โภตา โคตเมน ฐปิโต ‘อยํ โว มมจฺจเยน ปฏิสรณํ ภวิสฺสตี’ติ, ยํ ตุเมฺห เอตรหิ ปฏิปาเทยฺยาถา”ติ อิติ ปุฏฺโฐ สมาโน “นตฺถิ โข พฺราหฺมณ เอกภิกฺขุปิ เตน ภควตา ชานตา ปสฺสตา อรหตา สมฺมาสมฺพุทฺเธน ฐปิโต ‘อยํ โว มมจฺจเยน ปฏิสรณํ ภวิสฺสตี’ติ, ยํ มยํ เอตรหิ ปฏิปาเทยฺยามา”ติ วเทสิ. “อตฺถิ ปน โภ อานนฺท เอกภิกฺขุปิ สํเฆน สมฺมโต สมฺพหุเลหิ เถเรหิ ภิกฺขูหิ ฐปิโต ‘อยํ โน ภควโต อจฺจเยน ปฏิสรณํ ภวิสฺสตี’ติ, ยํ ตุเมฺห เอตรหิ ปฏิปาเทยฺยาถา”ติ อิติ ปุฏฺโฐ สมาโน “นตฺถิ โข พฺราหฺมณ เอกภิกฺขุปิ สํเฆน สมฺมโต สมฺพหุเลหิ เถเรหิ ภิกฺขูหิ ฐปิโต ‘อยํ โน ภควโต อจฺจเยน ปฏิสรณํ ภวิสฺสตี’ติ, ยํ มยํ เอตรหิ ปฏิปาเทยฺยามา”ติ วเทสิ.\csromanspacer{} เอวํ อปฺปฏิสรเณ จ ปน โภ อานนฺท “โก เหตุ สามคฺคิยา”ติ อิติ ปุฏฺโฐ สมาโน “น โข มยํ พฺราหฺมณ อปฺปฏิสรณา, สปฺปฏิสรณา มยํ พฺราหฺมณ ธมฺมปฺปฏิสรณา”ติ วเทสิ.\csromanspacer{} อิมสฺส ปน โภ อานนฺท ภาสิตสฺส กถํ อตฺโถ ทฏฺฐพฺโพติ.}

\proseitem{81}{อตฺถิ โข พฺราหฺมณ เตน ภควตา ชานตา ปสฺสตา อรหตา สมฺมาสมฺพุทฺเธน ภิกฺขูนํ สิกฺขาปทํ ปญฺญตฺตํ, ปาติโมกฺขํ อุทฺทิฏฺฐํ, เต มยํ ตทหุโปสเถ ยาวติกา เอกํ คามเขตฺตํ อุปนิสฺสาย วิหราม, เต}

\vspace{\tipitakaparskip}
\csromancenter{โคปกโมคฺคลฺลานสุตฺต (108) สพฺเพ เอกชฺฌํ สนฺนิปตาม, สนฺนิปติตฺวา ยสฺส ตํ ปวตฺตติ, ตํ อชฺเฌสาม, ตสฺมิํ เจ ภญฺญมาเน โหติ ภิกฺขุสฺส อาปตฺติ, โหติ วีติกฺกโม, ตํ มยํ ยถาธมฺมํ ยถานุสิฏฺฐํ กาเรมาติ.}
\prose{\csromanfolio{61}น กิร โน ภวนฺโต กาเรนฺติ, ธมฺโม โน กาเรติ.\csromanspacer{} อตฺถิ นุ โข โภ อานนฺท เอกภิกฺขุปิ, ยํ ตุเมฺห เอตรหิ สกฺกโรถ ครุํ กโรถ\footnote{ครุกโรถ (สี, สฺยา,กํ, อิ)} มาเนถ ปูเชถ, สกฺกตฺวา ครุํ กตฺวา\footnote{ครุกตฺวา (สี, สฺยา, กํ, อิ)} อุปนิสฺสาย วิหรถาติ.\csromanspacer{} นตฺถิ โข พฺราหฺมณ เอกภิกฺขุปิ, ยํ มยํ เอตรหิ สกฺกโรม ครุํ กโรม มาเนม ปูเชม, สกฺกตฺวา ครุํ กตฺวา อุปนิสฺสาย วิหรามาติ.}

\prose{“อตฺถิ นุ โข โภ อานนฺท เอกภิกฺขุปิ เตน โภตา โคตเมน ฐปิโต ‘อยํ โว มมจฺจเยน ปฏิสรณํ ภวิสฺสตี’ติ, ยํ ตุเมฺห เอตรหิ ปฏิปาเทยฺยาถา”ติ อิติ ปุฏฺโฐ สมาโน “นตฺถิ โข พฺราหฺมณ เอกภิกฺขุปิ เตน ภควตา ชานตา ปสฺสตา อรหตา สมฺมาสมฺพุทฺเธน ฐปิโต ‘อยํ โว มมจฺจเยน ปฏิสรณํ ภวิสฺสตี’ติ, ยํ มยํ เอตรหิ ปฏิปาเทยฺยามา”ติ วเทสิ. “อตฺถิ ปน โภ อานนฺท เอกภิกฺขุปิ สํเฆน สมฺมโต สมฺพหุเลหิ เถเรหิ ภิกฺขูหิ ฐปิโต ‘อยํ โน ภควโต อจฺจเยน ปฏิสรณํ ภวิสฺสตี’ติ, ยํ ตุเมฺห เอตรหิ ปฏิปาเทยฺยาถา”ติ อิติ ปุฏฺโฐ สมาโน “นตฺถิ โข พฺราหฺมณ เอกภิกฺขุปิ สํเฆน สมฺมโต สมฺพหุเลหิ เถเรหิ ภิกฺขูหิ ฐปิโต ‘อยํ โน ภควโต อจฺจเยน ปฏิสรณํ ภวิสฺสตี’ติ, ยํ มยํ เอตรหิ ปฏิปาเทยฺยามา”ติ วเทสิ. “อตฺถิ นุ โข โภ อานนฺท เอกภิกฺขุปิ ‘ยํ ตุเมฺห เอตรหิ สกฺกโรถ ครุํ กโรถ มาเนถ ปูเชถ, สกฺกตฺวา ครุํ กตฺวา อุปนิสฺสาย วิหรถา’ติ อิติ ปุฏฺโฐ สมาโน “นตฺถิ โข พฺราหฺมณ เอกภิกฺขุปิ ยํ มยํ เอตรหิ สกฺกโรม ครุํ กโรม มาเนม ปูเชม, สกฺกตฺวา ครุํ กตฺวา อุปนิสฺสาย วิหรามา”ติ วเทสิ.\csromanspacer{} อิมสฺส ปน โภ อานนฺท ภาสิตสฺส กถํ อตฺโถ ทฏฺฐพฺโพติ.}

\proseitem{82}{อตฺถิ โข พฺราหฺมณ เตน ภควตา ชานตา ปสฺสตา อรหตา สมฺมาสมฺพุทฺเธน ทส ปสาทนียา ธมฺมา อกฺขาตา, ยสฺมิํ โน อิเม ธมฺมา สํวิชฺชนฺติ, ตํ มยํ เอตรหิ สกฺกโรม ครุํ กโรม มาเนม ปูเชม, สกฺกตฺวา ครุํ กตฺวา อุปนิสฺสาย วิหราม.\csromanspacer{} กตเม ทส.}

\prose{\csromanfolio{62}อิธ พฺราหฺมณ ภิกฺขุ สีลวา โหติ, ปาติโมกฺขสํวรสํวุโต วิหรติ อาจารโคจรสมฺปนฺโน อณุมตฺเตสุ วชฺเชสุ ภยทสฺสาวี, สมาทาย สิกฺขติ สิกฺขาปเทสุ. (1)}

\prose{พหุสฺสุโต โหติ สุตธโร สุตสนฺนิจโย, เย เต ธมฺมา อาทิกลฺยาณา มชฺเฌกลฺยาณา ปริโยสานกลฺยาณา สาตฺถํ สพฺยญฺชนํ\footnote{สาตฺถา สพฺยญฺชนา (สี, สฺยา, กํ)} เกวลปริปุณฺณํ ปริสุทฺธํ พฺรหฺมจริยํ อภิวทนฺติ, ตถารูปาสฺส ธมฺมา พหุสฺสุตา โหนฺติ ธาตา\footnote{ธตา (สี, สฺยา, กํ, อิ)} วจสา ปริจิตา มนสานุเปกฺขิตา ทิฏฺฐิยา สุปฺปฏิวิทฺธา. (2)}

\prose{สนฺตุฏฺโฐ โหติ ( )\footnote{(อิตรีตเรหิ) ที 3. 223 ปิฏฺเฐ ทิสฺสติ.} จีวรปิณฺฑปาตเสนาสนคิลานปฺปจฺจยเภสชฺชปริกฺขาเรหิ. (3)}

\prose{จตุนฺนํ ฌานานํ อาภิเจตสิกานํ ทิฏฺฐธมฺมสุขวิหารานํ นิกามลาภี โหติ อกิจฺฉลาภี อกสิรลาภี. (4)}

\prose{อเนกวิหิตํ อิทฺธิวิธํ ปจฺจนุโภติ, เอโกปิ หุตฺวา พหุธา โหติ, พหุธาปิ หุตฺวา เอโก โหติ, อาวิภาวํ ติโรภาวํ ติโรกุฏฺฏํ\footnote{ติโรกุฑฺฑํ (สี, สฺยา, กํ, อิ)} ติโรปาการํ ติโรปพฺพตํ อสชฺชมาโน คจฺฉติ เสยฺยถาปิ อากาเส, ปถวิยาปิ อุมฺมุชฺชนิมุชฺชํ กโรติ เสยฺยถาปิ อุทเก, อุทเกปิ อภิชฺชมาเน คจฺฉติ เสยฺยถาปิ ปถวิยํ, อากาเสปิ ปลฺลงฺเกน กมติ เสยฺยถาปิ ปกฺขี สกุโณ, อิเมปิ จนฺทิมสูริเย เอวํมหิทฺธิเก เอวํมหานุภาเว ปาณินา ปริมสติ\footnote{ปรามสติ (ก)} ปริมชฺชติ, ยาว พฺรหฺมโลกาปิ กาเยน วสํ วตฺเตติ. (5)}

\prose{ทิพฺพาย โสตธาตุยา วิสุทฺธาย อติกฺกนฺตมานุสิกาย อุโภ สทฺเท สุณาติ ทิพฺเพ จ มานุเส จ เย ทูเร สนฺติเก จ. (6)}

\prose{ปรสตฺตานํ ปรปุคฺคลานํ เจตสา เจโต ปริจฺจ ปชานาติ, สราคํ วา จิตฺตํ สราคํ จิตฺตนฺติ ปชานาติ, วีตราคํ วา จิตฺตํ วีตราคํ จิตฺตนฺติ ปชานาติ, สโทสํ วา จิตฺตํ สโทสํ จิตฺตนฺติ ปชานาติ, วีตโทสํ วา จิตฺตํ วีตโทสํ จิตฺตนฺติ ปชานาติ, สโมหํ วา จิตฺตํ สโมหํ จิตฺตนฺติ ปชานาติ, วีตโมหํ วา จิตฺตํ วีตโมหํ จิตฺตนฺติ ปชานาติ, สํขิตฺตํ วา จิตฺตํ สํขิตฺตํ จิตฺตนฺติ ปชานาติ, วิกฺขิตฺตํ วา จิตฺตํ วิกฺขิตฺตํ จิตฺตนฺติ ปชานาติ, มหคฺคตํ วา จิตฺตํ}

\vspace{\tipitakaparskip}
\csromancenter{โคปกโมคฺคลฺลานสุตฺต (108) มหคฺคตํ จิตฺตนฺติ ปชานาติ, อมหคฺคตํ วา จิตฺตํ อมหคฺคตํ จิตฺตนฺติ ปชานาติ, ส-อุตฺตรํ วา จิตฺตํ ส-อุตฺตรํ จิตฺตนฺติ ปชานาติ, อนุตฺตรํ วา จิตฺตํ อนุตฺตรํ จิตฺตนฺติ ปชานาติ, สมาหิตํ วา จิตฺตํ สมาหิตํ จิตฺตนฺติ ปชานาติ, อสมาหิตํ วา จิตฺตํ อสมาหิตํ จิตฺตนฺติ ปชานาติ, วิมุตฺตํ วา จิตฺตํ วิมุตฺตํ จิตฺตนฺติ ปชานาติ, อวิมุตฺตํ วา จิตฺตํ อวิมุตฺตํ จิตฺตนฺติ ปชานาติ. (7)}
\prose{\csromanfolio{63}อเนกวิหิตํ ปุพฺเพนิวาสํ อนุสฺสรติ.\csromanspacer{} เสยฺยถิทํ, เอกมฺปิ ชาติํ เทฺวปิ ชาติโย ติสฺโสปิ ชาติโย จตสฺโสปิ ชาติโย ปญฺจปิ ชาติโย ทสปิ ชาติโย วีสมฺปิ ชาติโย ติํสมฺปิ ชาติโย จตฺตารีสมฺปิ ชาติโย ปญฺญาสมฺปิ ชาติโย ชาติสตมฺปิ ชาติสหสฺสมฺปิ ชาติสตสหสฺสมฺปิ อเนเกปิ สํวฏฺฏกปฺเป อเนเกปิ วิวฏฺฏกปฺเป อเนเกปิ สํวฏฺฏวิวฏฺฏกปฺเป, “อมุตฺราสิํ เอวํนาโม เอวํโคตฺโต เอวํวณฺโณ เอวมาหาโร เอวํสุขทุกฺขปฺปฏิสํเวที เอวมายุปริยนฺโต, โส ตโต จุโต อมุตฺร อุทปาทิํ, ตตฺราปาสิํ เอวํนาโม เอวํโคตฺโต เอวํวณฺโณ เอวมาหาโร เอวํสุขทุกฺขปฺปฏิสํเวที เอวมายุปริยนฺโต, โส ตโต จุโต อิธูปปนฺโน”ติ, อิติ สาการํ ส-อุทฺเทสํ อเนกวิหิตํ ปุพฺเพนิวาสํ อนุสฺสรติ. (8)}

\prose{ทิพฺเพน จกฺขุนา วิสุทฺเธน อติกฺกนฺตมานุสเกน สตฺเต ปสฺสติ จวมาเน อุปปชฺชมาเน หีเน ปณีเต สุวณฺเณ ทุพฺพณฺเณ สุคเต ทุคฺคเต, ยถากมฺมูปเค สตฺเต ปชานาติ. (9)}

\prose{อาสวานํ ขยา อนาสวํ เจโตวิมุตฺติํ ปญฺญาวิมุตฺติํ ทิฏฺเฐว ธมฺเม สยํ อภิญฺญา สจฺฉิกตฺวา อุปสมฺปชฺช วิหรติ. (10)}

\prose{อิเม โข พฺราหฺมณ เตน ภควตา ชานตา ปสฺสตา อรหตา สมฺมาสมฺพุทฺเธน ทส ปสาทนียา ธมฺมา อกฺขาตา.\csromanspacer{} ยสฺมิํ โน อิเม ธมฺมา สํวิชฺชนฺติ, ตํ มยํ เอตรหิ สกฺกโรม ครุํ กโรม มาเนม ปูเชม, สกฺกตฺวา ครุํ กตฺวา อุปนิสฺสาย วิหรามาติ.}

\proseitem{83}{เอวํ วุตฺเต วสฺสกาโร พฺราหฺมโณ มคธมหามตฺโต อุปนนฺทํ เสนาปติํ อามนฺเตสิ “ตํ กิํ มญฺญติ ภวํ เสนาปติ\footnote{มญฺญสิ เอวํ เสนาปติ (สฺยา, กํ, อิ), มญฺญสิ เสนาปติ (สี), มญฺญสิ ภวํ เสนาปติ (ก)}, ยทิเม โภนฺโต สกฺกาตพฺพํ สกฺกโรนฺติ, ครุํ กาตพฺพํ ครุํ กโรนฺติ, มาเนตพฺพํ \csromanfolio{64}มาเนนฺติ, ปูเชตพฺพํ ปูเชนฺติ.\csromanspacer{} ตคฺฆิเม\footnote{ตคฺฆ เม (ก)} โภนฺโต สกฺกาตพฺพํ สกฺกโรนฺติ, ครุํ กาตพฺพํ ครุํ กโรนฺติ, มาเนตพฺพํ มาเนนฺติ, ปูเชตพฺพํ ปูเชนฺติ.\csromanspacer{} อิมญฺจ หิ เต โภนฺโต น สกฺกเรยฺยุํ น ครุํ กเรยฺยุํ น มาเนยฺยุํ น ปูเชยฺยุํ, อถ กิญฺจรหิ เต โภนฺโต สกฺกเรยฺยุํ ครุํ กเรยฺยุํ มาเนยฺยุํ ปูเชยฺยุํ, สกฺกตฺวา ครุํ กตฺวา มาเนตฺวา ปูเชตฺวา อุปนิสฺสาย วิหเรยฺยุนฺติ.\csromanspacer{} อถ โข วสฺสกาโร พฺราหฺมโณ มคธมหามตฺโต อายสฺมนฺตํ อานนฺทํ เอตทโวจ “กหํ ปน ภวํ อานนฺโท เอตรหิ วิหรตี”ติ.\csromanspacer{} เวฬุวเน โขหํ พฺราหฺมณ เอตรหิ วิหรามีติ.\csromanspacer{} กจฺจิ ปน โภ อานนฺท เวฬุวนํ รมณียญฺเจว อปฺปสทฺทญฺจ อปฺปนิคฺโฆสญฺจ วิชนวาตํ มนุสฺสราหสฺเสยฺยกํ\footnote{มนุสฺสราหเสยฺยกํ (สี, สฺยา, กํ, อิ)} ปฏิสลฺลานสารุปฺปนฺติ.\csromanspacer{} ตคฺฆ พฺราหฺมณ เวฬุวนํ รมณียญฺเจว อปฺปสทฺทญฺจ อปฺปนิคฺโฆสญฺจ วิชนวาตํ มนุสฺสราหสฺเสยฺยกํ ปฏิสลฺลานสารุปฺปํ.\csromanspacer{} ยถา ตํ ตุมฺหาทิเสหิ รกฺขเกหิ โคปเกหีติ.\csromanspacer{} ตคฺฆ โภ อานนฺท เวฬุวนํ รมณียญฺเจว อปฺปสทฺทญฺจ อปฺปนิคฺโฆสญฺจ วิชนวาตํ มนุสฺสราหสฺเสยฺยกํ ปฏิสลฺลานสารุปฺปํ, ยถา ตํ ภวนฺเตหิ ฌายีหิ ฌานสีลีหิ, ฌายิโน เจว ภวนฺโต ฌานสีลิโน จ.}

\prose{เอกมิทาหํ โภ อานนฺท สมยํ โส ภวํ โคตโม เวสาลิยํ วิหรติ มหาวเน กูฏาคารสาลายํ.\csromanspacer{} อถ ขฺวาหํ โภ อานนฺท เยน มหาวนํ กูฏาคารสาลา, เยน โส ภวํ โคตโม เตนุปสงฺกมิํ.\csromanspacer{} ตตฺร จ ปน โส\footnote{ตตฺร จ โส (สี, อิ)} ภวํ โคตโม อเนกปริยาเยน ฌานกถํ กเถสิ, ฌายี เจว โส ภวํ โคตโม อโหสิ ฌานสีลี จ, สพฺพญฺจ ปน โส ภวํ โคตโม ฌานํ วณฺเณสีติ.}

\proseitem{84}{น จ โข พฺราหฺมณ โส ภควา สพฺพํ ฌานํ วณฺเณสิ, นปิ โส ภควา สพฺพํ ฌานํ น วณฺเณสีติ.\csromanspacer{} กถํ รูปญฺจ พฺราหฺมณ โส ภควา ฌานํ น วณฺเณสิ.\csromanspacer{} อิธ พฺราหฺมณ เอกจฺโจ กามราคปริยุฏฺฐิเตน เจตสา วิหรติ กามราคปเรเตน, อุปฺปนฺนสฺส จ กามราคสฺส นิสฺสรณํ ยถาภูตํ นปฺปชานาติ, โส กามราคํเยว อนฺตรํ กริตฺวา ฌายติ ปชฺฌายติ นิชฺฌายติ อปชฺฌายติ.\csromanspacer{} พฺยาปาทปริยุฏฺฐิเตน เจตสา วิหรติ พฺยาปาทปเรเตน, อุปฺปนฺนสฺส จ พฺยาปาทสฺส นิสฺสรณํ ยถาภูตํ นปฺปชานาติ, โส พฺยาปาทํเยว อนฺตรํ กริตฺวา ฌายติ ปชฺฌายติ นิชฺฌายติ อปชฺฌายติ.}

\vspace{\tipitakaparskip}
\csromancenter{โคปกโมคฺคลฺลานสุตฺต (108) ถินมิทฺธปริยุฏฺฐิเตน เจตสา วิหรติ ถินมิทฺธปเรเตน, อุปฺปนฺนสฺส จ ถินมิทฺธสฺส นิสฺสรณํ ยถาภูตํ นปฺปชานาติ, โส ถินมิทฺธํเยว อนฺตรํ กริตฺวา ฌายติ ปชฺฌายติ นิชฺฌายติ อปชฺฌายติ.\csromanspacer{} อุทฺธจฺจกุกฺกุจฺจปริยุฏฺฐิเตน เจตสา วิหรติ อุทฺธจฺจกุกฺกุจฺจปเรเตน, อุปฺปนฺนสฺส จ อุทฺธจฺจกุกฺกุจฺจสฺส นิสฺสรณํ ยถาภูตํ นปฺปชานาติ, โส อุทฺธจฺจกุกฺกุจฺจํเยว อนฺตรํ กริตฺวา ฌายติ ปชฺฌายติ นิชฺฌายติ อปชฺฌายติ.\csromanspacer{} วิจิกิจฺฉาปริยุฏฺฐิเตน เจตสา วิหรติ วิจิกิจฺฉาปเรเตน, อุปฺปนฺนาย จ วิจิกิจฺฉาย นิสฺสรณํ ยถาภูตํ นปฺปชานาติ, โส วิจิกิจฺฉํเยว อนฺตรํ กริตฺวา ฌายติ ปชฺฌายติ นิชฺฌายติ อปชฺฌายติ.\csromanspacer{} เอวรูปํ โข พฺราหฺมณ โส ภควา ฌานํ น วณฺเณสิ.}
\prose{\csromanfolio{65}กถํ รูปญฺจ พฺราหฺมณ โส ภควา ฌานํ วณฺเณสิ.\csromanspacer{} อิธ พฺราหฺมณ ภิกฺขุ วิวิจฺเจว กาเมหิ วิวิจฺจ อกุสเลหิ ธมฺเมหิ สวิตกฺกํ สวิจารํ วิเวกชํ ปีติสุขํ ปฐมํ ฌานํ อุปสมฺปชฺช วิหรติ.\csromanspacer{} วิตกฺกวิจารานํ วูปสมา อชฺฌตฺตํ สมฺปสาทนํ เจตโส เอโกทิภาวํ อวิตกฺกํ อวิจารํ สมาธิชํ ปีติสุขํ ทุติยํ ฌานํ.\csromanspacer{} ตติยํ ฌานํ.\csromanspacer{} จตุตฺถํ ฌานํ อุปสมฺปชฺช วิหรติ.\csromanspacer{} เอวรูปํ โข พฺราหฺมณ โส ภควา ฌานํ วณฺเณสีติ.}

\prose{คารยฺหํ กิร โภ อานนฺท โส ภวํ โคตโม ฌานํ ครหิ, ปาสํสํ ปสํสิ.\csromanspacer{} หนฺท จ ทานิ มยํ โภ อานนฺท คจฺฉาม, พหุกิจฺจา มยํ พหุกรณียาติ.\csromanspacer{} ยสฺสทานิ ตฺวํ พฺราหฺมณ กาลํ มญฺญสีติ.\csromanspacer{} อถ โข วสฺสกาโร พฺราหฺมโณ มคธมหามตฺโต อายสฺมโต อานนฺทสฺส ภาสิตํ อภินนฺทิตฺวา อนุโมทิตฺวา อุฏฺฐายาสนา ปกฺกามิ.}

\prose{อถ โข โคปกโมคฺคลฺลาโน พฺราหฺมโณ อจิรปกฺกนฺเต วสฺสกาเร พฺราหฺมเณ มคธมหามตฺเต อายสฺมนฺตํ อานนฺทํ เอตทโวจ “ยํ โน มยํ ภวนฺตํ อานนฺทํ อปุจฺฉิมฺหา, ตํ โน ภวํ อานนฺโท น พฺยากาสี”ติ.\csromanspacer{} นนุ เต พฺราหฺมณ อโวจุมฺหา “นตฺถิ โข พฺราหฺมณ เอกภิกฺขุปิ เตหิ ธมฺเมหิ สพฺเพนสพฺพํ สพฺพถาสพฺพํ สมนฺนาคโต, เยหิ ธมฺเมหิ สมนฺนาคโต โส ภควา อโหสิ อรหํ สมฺมาสมฺพุทฺโธ.\csromanspacer{} โส หิ พฺราหฺมณ ภควา อนุปฺปนฺนสฺส มคฺคสฺส อุปฺปาเทตา, อสญฺชาตสฺส มคฺคสฺส สญฺชเนตา, อนกฺขาตสฺส มคฺคสฺส อกฺขาตา, มคฺคญฺญู มคฺควิทู \csromanfolio{66}มคฺคโกวิโท.\csromanspacer{} มคฺคานุคา จ ปน เอตรหิ สาวกา วิหรนฺติ ปจฺฉา สมนฺนาคตา”ติ.}

\nitthitamruled{โคปกโมคฺคลฺลานสุตฺตํ นิฏฺฐิตํ อฏฺฐมํ.}
\csromanmatikaanchor{mtk.10}
\csromantocmarkref{subsection}{9. มหาปุณฺณมสุตฺต}{mtk.10}
\bookmark[dest={mtk.10},level=2]{9. มหาปุณฺณมสุตฺต}
\csromanheader{2}{9. มหาปุณฺณมสุตฺต}
\proseitem{85}{เอวํ เม สุตํ\csromandash{}เอกํ สมยํ ภควา สาวตฺถิยํ วิหรติ ปุพฺพาราเม มิคารมาตุปาสาเท.\csromanspacer{} เตน โข ปน สมเยน ภควา ตทหุโปสเถ ปนฺนรเส ปุณฺณาย ปุณฺณมาย รตฺติยา ภิกฺขุสํฆปริวุโต อพฺโภกาเส นิสินฺโน โหติ.\csromanspacer{} อถ โข อญฺญตโร ภิกฺขุ อุฏฺฐายาสนา เอกํสํ จีวรํ กตฺวา เยน ภควา เตนญฺชลิํ ปณาเมตฺวา ภควนฺตํ เอตทโวจ\csromandash{}}

\prose{ปุจฺเฉยฺยาหํ ภนฺเต ภควนฺตํ กิญฺจิเทว เทสํ, สเจ เม ภควา โอกาสํ กโรติ ปญฺหสฺส เวยฺยากรณายาติ.\csromanspacer{} เตน หิ ตฺวํ ภิกฺขุ สเก อาสเน นิสีทิตฺวา ปุจฺฉ, ยทากงฺขสีติ.}

\proseitem{86}{อถ โข โส ภิกฺขุ สเก อาสเน นิสีทิตฺวา ภควนฺตํ เอตทโวจ “อิเม นุ โข ภนฺเต ปญฺจุปาทานกฺขนฺธา.\csromanspacer{} เสยฺยถิทํ, รูปุปาทานกฺขนฺโธ เวทนุปาทานกฺขนฺโธ สญฺญุปาทานกฺขนฺโธ สงฺขารุปาทานกฺขนฺโธ วิญฺญาณุปาทานกฺขนฺโธติ.\csromanspacer{} อิเม โข ภิกฺขุ ปญฺจุปาทานกฺขนฺธา.\csromanspacer{} เสยฺยถิทํ, รูปุปาทานกฺขนฺโธ เวทนุปาทานกฺขนฺโธ สญฺญุปาทานกฺขนฺโธ สงฺขารุปาทานกฺขนฺโธ วิญฺญาณุปาทานกฺขนฺโธติ.}

\prose{“สาธุ ภนฺเต”ติ โข โส ภิกฺขุ ภควโต ภาสิตํ อภินนฺทิตฺวา อนุโมทิตฺวา ภควนฺตํ อุตฺตริํ ปญฺหํ ปุจฺฉิ “อิเม ปน ภนฺเต ปญฺจุปาทานกฺขนฺธา กิํมูลกา”ติ.\csromanspacer{} อิเม โข ภิกฺขุ ปญฺจุปาทานกฺขนฺธา ฉนฺทมูลกาติ.\csromanspacer{} ตํเยว นุ โข ภนฺเต อุปาทานํ เต ปญฺจุปาทานกฺขนฺธา, อุทาหุ อญฺญตฺร ปญฺจหุปาทานกฺขนฺเธหิ อุปาทานนฺติ.\csromanspacer{} น โข ภิกฺขุ ตํเยว อุปาทานํ เต ปญฺจุปาทานกฺขนฺเธหิ, นาปิ อญฺญตฺร ปญฺจหุปาทานกฺขนฺเธหิ อุปาทานํ.\csromanspacer{} โย โข ภิกฺขุ ปญฺจสุ อุปาทานกฺขนฺเธสุ ฉนฺทราโค, ตํ ตตฺถ อุปาทานนฺติ.}

\prose{สิยา ปน ภนฺเต ปญฺจสุ อุปาทานกฺขนฺเธสุ ฉนฺทราคเวมตฺตตาติ. “สิยา ภิกฺขู”ติ ภควา อโวจ.\csromanspacer{} อิธ ภิกฺขุ เอกจฺจสฺส เอวํ โหติ}

\vspace{\tipitakaparskip}
\csromancenter{มหาปุณฺณมสุตฺต (109) “เอวํรูโป สิยํ อนาคตมทฺธานํ, เอวํเวทโน สิยํ อนาคตมทฺธานํ, เอวํสญฺโญ สิยํ อนาคตมทฺธานํ, เอวํสงฺขาโร สิยํ อนาคตมทฺธานํ, เอวํวิญฺญาโณ สิยํ อนาคตมทฺธานนฺ”ติ.\csromanspacer{} เอวํ โข ภิกฺขุ สิยา ปญฺจสุ อุปาทานกฺขนฺเธสุ ฉนฺทราคเวมตฺตตาติ.}
\prose{\csromanfolio{67}กิตฺตาวตา ปน ภนฺเต ขนฺธานํ ขนฺธาธิวจนํ โหตีติ.\csromanspacer{} ยํ กิญฺจิ ภิกฺขุ รูปํ อตีตานาคตปจฺจุปฺปนฺนํ อชฺฌตฺตํ วา พหิทฺธา วา โอฬาริกํ วา สุขุมํ วา หีนํ วา ปณีตํ วา ยํ ทูเร สนฺติเก วา, อยํ รูปกฺขนฺโน.\csromanspacer{} ยา กาจิ เวทนา อตีตานาคตปจฺจุปฺปนฺนา อชฺฌตฺตํ วา พหิทฺธา วา โอฬาริกา วา สุขุมา วา หีนา วา ปณีตา วา ยา ทูเร สนฺติเก วา, อยํ เวทนากฺขนฺโธ.\csromanspacer{} ยา กาจิ สญฺญา อตีตานาคตปจฺจุปฺปนฺนา ฯเปฯ.\csromanspacer{} ยา ทูเร สนฺติเก วา, อยํ สญฺญากฺขนฺโธ.\csromanspacer{} เย เกจิ สงฺขารา อตีตานาคตปจฺจุปฺปนฺนา อชฺฌตฺตํ วา พหิทฺธา วา โอฬาริกา วา สุขุมา วา หีนา วา ปณีตา วา เย ทูเร สนฺติเก วา, อยํ สงฺขารากฺขนฺโธ.\csromanspacer{} ยํ กิญฺจิ วิญฺญาณํ อตีตานาคตปจฺจุปฺปนฺนํ อชฺฌตฺตํ วา พหิทฺธา วา โอฬาริกํ วา สุขุมํ วา หีนํ วา ปณีตํ วา ยํ ทูเร สนฺติเก วา, อยํ วิญฺญาณกฺขนฺโธ.\csromanspacer{} เอตฺตาวตา โข ภิกฺขุ ขนฺธานํ ขนฺธาธิวจนํ โหตีติ.}

\prose{โก นุ โข ภนฺเต เหตุ โก ปจฺจโย รูปกฺขนฺธสฺส ปญฺญาปนาย, โก เหตุ โก ปจฺจโย เวทนากฺขนฺธสฺส ปญฺญาปนาย, โก เหตุ โก ปจฺจโย สญฺญากฺขนฺธสฺส ปญฺญาปนาย, โก เหตุ โก ปจฺจโย สงฺขารกฺขนฺธสฺส ปญฺญาปนาย, โก เหตุ โก ปจฺจโย วิญฺญาณกฺขนฺธสฺส ปญฺญาปนายาติ.}

\prose{จตฺตาโร โข ภิกฺขุ มหาภูตา เหตุ จตฺตาโร มหาภูตา ปจฺจโย รูปกฺขนฺธสฺส ปญฺญาปนาย, ผสฺโส เหตุ ผสฺโส ปจฺจโย เวทนากฺขนฺธสฺส ปญฺญาปนาย, ผสฺโส เหตุ ผสฺโส ปจฺจโย สญฺญากฺขนฺธสฺส ปญฺญาปนาย, ผสฺโส เหตุ ผสฺโส ปจฺจโย สงฺขารกฺขนฺธสฺส ปญฺญาปนาย, นามรูปํ โข ภิกฺขุ เหตุ นามรูปํ ปจฺจโย วิญฺญาณกฺขนฺธสฺส ปญฺญาปนายาติ.}

\proseitem{87}{กถํ ปน ภนฺเต สกฺกายทิฏฺฐิ โหตีติ.\csromanspacer{} อิธ ภิกฺขุ อสฺสุตวา ปุถุชฺชโน อริยานํ อทสฺสาวี อริยธมฺมสฺส อโกวิโท อริยธมฺเม \csromanfolio{68}อวินีโต สปฺปุริสานํ อทสฺสาวี สปฺปุริสธมฺมสฺส อโกวิโท สปฺปุริสธมฺเม อวินีโต รูปํ อตฺตโต สมนุปสฺสติ, รูปวนฺตํ วา อตฺตานํ, อตฺตนิ วา รูปํ, รูปสฺมิํ วา อตฺตานํ.\csromanspacer{} เวทนํ อตฺตโต สมนุปสฺสติ, เวทนาวนฺตํ วา อตฺตานํ, อตฺตนิ วา เวทนํ, เวทนาย วา อตฺตานํ.\csromanspacer{} สญฺญํ อตฺตโต สมนุปสฺสติ, สญฺญาวนฺตํ วา อตฺตานํ, อตฺตนิ วา สญฺญํ, สญฺญาย วา อตฺตานํ.\csromanspacer{} สงฺขาเร อตฺตโต สมนุปสฺสติ, สงฺขารวนฺตํ วา อตฺตานํ, อตฺตนิ วา สงฺขาเร, สงฺขาเรสุ วา อตฺตานํ.\csromanspacer{} วิญฺญาณํ อตฺตโต สมนุปสฺสติ, วิญฺญาณวนฺตํ วา อตฺตานํ, อตฺตนิ วา วิญฺญาณํ, วิญฺญาณสฺมิํ วา อตฺตานํ.\csromanspacer{} เอวํ โข ภิกฺขุ สกฺกายทิฏฺฐิ โหตีติ.}

\prose{กถํ ปน ภนฺเต สกฺกายทิฏฺฐิ น โหตีติ.\csromanspacer{} อิธ ภิกฺขุ สุตวา อริยสาวโก อริยานํ ทสฺสาวี อริยธมฺมสฺส โกวิโท อริยธมฺเม สุวินีโต สปฺปุริสานํ ทสฺสาวี สปฺปุริสธมฺมสฺส โกวิโท สปฺปุริสธมฺเม สุวินีโต น รูปํ อตฺตโต สมนุปสฺสติ, น รูปวนฺตํ วา อตฺตานํ, น อตฺตนิ วา รูปํ, น รูปสฺมิํ วา อตฺตานํ.\csromanspacer{} น เวทนํ อตฺตโต สมนุปสฺสติ, น เวทนาวนฺตํ วา อตฺตานํ, น อตฺตนิ วา เวทนํ, น เวทนาย วา อตฺตานํ.\csromanspacer{} น สญฺญํ อตฺตโต สมนุปสฺสติ, น สญฺญาวนฺตํ วา อตฺตานํ, น อตฺตนิ วา สญฺญํ, น สญฺญาย วา อตฺตานํ.\csromanspacer{} น สงฺขาเร อตฺตโต สมนุปสฺสติ, น สงฺขารวนฺตํ วา อตฺตานํ, น อตฺตนิ วา สงฺขาเร, น สงฺขาเรสุ วา อตฺตานํ.\csromanspacer{} น วิญฺญาณํ อตฺตโต สมนุปสฺสติ, น วิญฺญาณวนฺตํ วา อตฺตานํ, น อตฺตนิ วา วิญฺญาณํ, น วิญฺญาณสฺมิํ วา อตฺตานํ.\csromanspacer{} เอวํ โข ภิกฺขุ สกฺกายทิฏฺฐิ น โหตีติ.}

\proseitem{88}{โก นุ โข ภนฺเต รูเป อสฺสาโท, โก อาทีนโว, กิํ นิสฺสรณํ.\csromanspacer{} โก เวทนาย อสฺสาโท, โก อาทีนโว, กิํ นิสฺสรณํ.\csromanspacer{} โก สญฺญาย อสฺสาโท, โก อาทีนโว, กิํ นิสฺสรณํ.\csromanspacer{} โก สงฺขาเรสุ อสฺสาโท, โก อาทีนโว, กิํ นิสฺสรณํ.\csromanspacer{} โก วิญฺญาเณ อสฺสาโท, โก อาทีนโว, กิํ นิสฺสรณนฺติ.\csromanspacer{} ยํ โข ภิกฺขุ รูปํ ปฏิจฺจ อุปฺปชฺชติ สุขํ โสมนสฺสํ, อยํ รูเป อสฺสาโท.\csromanspacer{} ยํ รูปํ อนิจฺจํ ทุกฺขํ วิปริณามธมฺมํ, อยํ รูเป อาทีนโว.\csromanspacer{} โย รูเป ฉนฺทราควินโย ฉนฺทราคปฺปหานํ, อิทํ รูเป นิสฺสรณํ.\csromanspacer{} ยํ โข\footnote{ยญฺจ (สฺยา, กํ)} ภิกฺขุ เวทนํ ปฏิจฺจ.}

\vspace{\tipitakaparskip}
\csromancenter{มหาปุณฺณมสุตฺต (109) สญฺญํ ปฏิจฺจ.\csromanspacer{} สงฺขาเร ปฏิจฺจ.\csromanspacer{} วิญฺญาณํ ปฏิจฺจ อุปฺปชฺชติ สุขํ โสมนสฺสํ, อยํ วิญฺญาเณ อสฺสาโท.\csromanspacer{} ยํ วิญฺญาณํ อนิจฺจํ ทุกฺขํ วิปริณามธมฺมํ, อยํ วิญฺญาเณ อาทีนโว.\csromanspacer{} โย วิญฺญาเณ ฉนฺทราควินโย ฉนฺทราคปฺปหานํ, อิทํ วิญฺญาเณ นิสฺสรณนฺติ.}
\proseitem{89}{\csromanfolio{69}กถํ ปน ภนฺเต ชานโต กถํ ปสฺสโต อิมสฺมิํ จ สวิญฺญาณเก กาเย พหิทฺธา จ สพฺพนิมิตฺเตสุ อหํการมมํการมานานุสยา น โหนฺตีติ.\csromanspacer{} ยํ กิญฺจิ ภิกฺขุ รูปํ อตีตานาคตปจฺจุปฺปนฺนํ อชฺฌตฺตํ วา พหิทฺธา วา โอฬาริกํ วา สุขุมํ วา หีนํ วา ปณีตํ วา ยํ ทูเร สนฺติเก วา, สพฺพํ รูปํ “เนตํ มม, เนโสหมสฺมิ, น เมโส อตฺตา”ติ เอวเมตํ ยถาภูตํ สมฺมปฺปญฺญาย ปสฺสติ.\csromanspacer{} ยา กาจิ เวทนา.\csromanspacer{} ยา กาจิ สญฺญา.\csromanspacer{} เย เกจิ สงฺขารา.\csromanspacer{} ยํ กิญฺจิ วิญฺญาณํ อตีตานาคตปจฺจุปฺปนฺนํ อชฺฌตฺตํ วา พหิทฺธา วา โอฬาริกํ วา สุขุมํ วา หีนํ วา ปณีตํ วา ยํ ทูเร สนฺติเก วา, สพฺพํ วิญฺญาณํ “เนตํ มม, เนโสหมสฺมิ, น เมโส อตฺตา”ติ เอวเมตํ ยถาภูตํ สมฺมปฺปญฺญาย ปสฺสติ.\csromanspacer{} เอวํ โข ภิกฺขุ ชานโต เอวํ ปสฺสโต อิมสฺมิํ จ สวิญฺญาณเก กาเย พหิทฺธา จ สพฺพนิมิตฺเตสุ อหํการมมํการมานุสยา น โหนฺตีติ.}

\proseitem{90}{อถ โข อญฺญตรสฺส ภิกฺขุโน เอวํ เจตโส ปริวิตกฺโก อุทปาทิ “อิติ กิร โภ รูปํ อนตฺตา, เวทนา อนตฺตา, สญฺญา อนตฺตา, สงฺขารา อนตฺตา, วิญฺญาณํ อนตฺตา, อนตฺตกตานิ กมฺมานิ กมตฺตานํ\footnote{กถมตฺตานํ (สํ 2. 85 ปิฏฺเฐ)} ผุสิสฺสนฺตี”ติ.\csromanspacer{} อถ โข ภควา ตสฺส ภิกฺขุโน เจตสา เจโตปริวิตกฺกมญฺญาย ภิกฺขู อามนฺเตสิ “ถานํ โข ปเนตํ ภิกฺขเว วิชฺชติ, ยํ อิเธกจฺโจ โมฆปุริโส อวิทฺวา อวิชฺชาคโต ตณฺหาธิปเตยฺเยน เจตสา สตฺถุ สาสนํ อติธาวิตพฺพํ มญฺเญยฺย ‘อิติ กิร โภ รูปํ อนตฺตา, เวทนา อนตฺตา, สญฺญา อนตฺตา, สงฺขารา อนตฺตา, วิญฺญาณํ อนตฺตา, อนตฺตกตานิ กมฺมานิ กมตฺตานํ ผุสิสฺสนฺตี’ติ”.\csromanspacer{} ปฏิวินีตา\footnote{ปฏิจฺจ วินีตา (สี, อิ), ปฏิปุจฺฉามิ วินีตา (สฺยา, กํ)} โข เม ตุเมฺห ภิกฺขเว ตตฺร ตตฺร ธมฺเมสุ.}

\prose{ตํ กิํ มญฺญถ ภิกฺขเว, รูปํ นิจฺจํ วา อนิจฺจํ วาติ.\csromanspacer{} อนิจฺจํ ภนฺเต.\csromanspacer{} ยํ ปนานิจฺจํ, ทุกฺขํ วา ตํ สุขํ วาติ.\csromanspacer{} ทุกฺขํ ภนฺเต.\csromanspacer{} ยํ ปนานิจฺจํ ทุกฺขํ วิปริณามธมฺมํ, กลฺลํ นุ ตํ สมนุปสฺสิตุํ “เอตํ มม, เอโสหมสฺมิ, เอโส เม อตฺตา”ติ.\csromanspacer{} โน \csromanfolio{70}เหตํ ภนฺเต.\csromanspacer{} ตํ กิํ มญฺญถ ภิกฺขเว, เวทนา.\csromanspacer{} สญฺญา.\csromanspacer{} สงฺขารา.\csromanspacer{} วิญฺญาณํ นิจฺจํ วา อนิจฺจํ วาติ.\csromanspacer{} อนิจฺจํ ภนฺเต.\csromanspacer{} ยํ ปนานิจฺจํ, ทุกฺขํ วา ตํ สุขํ วาติ.\csromanspacer{} ทุกฺขํ ภนฺเต.\csromanspacer{} ยํ ปนานิจฺจํ ทุกฺขํ วิปริณามธมฺมํ, กลฺลํ นุ ตํ สมนุปสฺสิตุํ “เอตํ มม, เอโสหมสฺมิ, เอโส เม อตฺตา”ติ.\csromanspacer{} โน เหตํ ภนฺเต.\csromanspacer{} ตสฺมาติห ภิกฺขเว ยํ กิญฺจิ รูปํ อตีตานาคตปจฺจุปฺปนฺนํ อชฺฌตฺตํ วา พหิทฺธา วา โอฬาริกํ วา สุขุมํ วา หีนํ วา ปณีตํ วา ยํ ทูเร สนฺติเก วา, สพฺพํ รูปํ “เนตํ มม, เนโสหมสฺมิ, น เมโส อตฺตา”ติ เอวเมตํ ยถาภูตํ สมฺมปฺปญฺญาย ทฏฺฐพฺพํ.\csromanspacer{} ยา กาจิ เวทนา.\csromanspacer{} ยา กาจิ สญฺญา.\csromanspacer{} เย เกจิ สงฺขารา.\csromanspacer{} ยํ กิญฺจิ วิญฺญาณํ อตีตานาคตปจฺจุปฺปนฺนํ อชฺฌตฺตํ วา พหิทฺธา วา โอฬาริกํ วา สุขุมํ วา หีนํ วา ปณีตํ วา ยํ ทูเร สนฺติเก วา, สพฺพํ วิญฺญาณํ “เนตํ มม, เนโสหมสฺมิ, น เมโส อตฺตา”ติ เอวเมตํ ยถาภูตํ สมฺมปฺปญฺญาย ทฏฺฐพฺพํ.\csromanspacer{} เอวํ ปสฺสํ ภิกฺขเว สุตวา อริยสาวโก รูปสฺมิํปิ นิพฺพินฺทติ, เวทนายปิ นิพฺพินฺทติ, สญฺญายปิ นิพฺพินฺทติ, สงฺขาเรสุปิ นิพฺพินฺทติ, วิญฺญาณสฺมิํปิ นิพฺพินฺทติ.\csromanspacer{} นิพฺพินฺทํ วิรชฺชติ, วิราคา วิมุจฺจติ, วิมุตฺตสฺมิํ “วิมุตฺตมฺ”อิติ ญาณํ โหติ, “ขีณา ชาติ, วุสิตํ พฺรหฺมจริยํ, กตํ กรณียํ, นาปรํ อิตฺถตฺตายา”ติ ปชานาตีติ.}

\prose{อิทมโวจ ภควา.\csromanspacer{} อตฺตมนา เต ภิกฺขู ภควโต ภาสิตํ อภินนฺทุนฺติ.\csromanspacer{} อิมสฺมิํ จ ปน เวยฺยากรณสฺมิํ ภญฺญมาเน สฏฺฐิมตฺตานํ ภิกฺขูนํ อนุปาทาย อาสเวหิ จิตฺตานิ วิมุจฺจิํสูติ.}

\nitthitamruled{มหาปุณฺณมสุตฺตํ นิฏฺฐิตํ นวมํ.}
\csromanmatikaanchor{mtk.11}
\csromantocmarkref{subsection}{10. จูฬปุณฺณมสุตฺต}{mtk.11}
\bookmark[dest={mtk.11},level=2]{10. จูฬปุณฺณมสุตฺต}
\csromanheader{2}{10. จูฬปุณฺณมสุตฺต}
\proseitem{91}{เอวํ เม สุตํ\csromandash{}เอกํ สมยํ ภควา สาวตฺถิยํ วิหรติ ปุพฺพาราเม มิคารมาตุปาสาเท.\csromanspacer{} เตน โข ปน สมเยน ภควา ตทหุโปสเถ ปนฺนรเส ปุณฺณาย ปุณฺณมาย รตฺติยา ภิกฺขุสํฆปริวุโต อพฺโภกาเส นิสินฺโน โหติ.\csromanspacer{} อถ โข ภควา ตุณฺหีภูตํ ตุณฺหีภูตํ ภิกฺขุสํฆํ อนุวิโลเกตฺวา ภิกฺขู อามนฺเตสิ “ชาเนยฺย นุ โข ภิกฺขเว อสปฺปุริโส อสปฺปุริสํ ‘อสปฺปุริโส อยํ ภวนฺ’ติ”.\csromanspacer{} โน เหตํ ภนฺเต.\csromanspacer{} สาธุ ภิกฺขเว, อฏฺฐานเมตํ ภิกฺขเว}

\vspace{\tipitakaparskip}
\csromancenter{จูฬปุณฺณมสุตฺต (110) อนวกาโส, ยํ อสปฺปุริโส อสปฺปุริสํ ชาเนยฺย “อสปฺปุริโส อยํ ภวนฺ”ติ.\csromanspacer{} ชาเนยฺย ปน ภิกฺขเว อสปฺปุริโส สปฺปุริสํ “สปฺปุริโส อยํ ภวนฺ”ติ.\csromanspacer{} โน เหตํ ภนฺเต.\csromanspacer{} สาธุ ภิกฺขเว, เอตมฺปิ โข ภิกฺขเว อฏฺฐานํ อนวกาโส, ยํ อสปฺปุริโส สปฺปุริสํ ชาเนยฺย “สปฺปุริโส อยํ ภวนฺ”ติ.\csromanspacer{} อสปฺปุริโส ภิกฺขเว อสฺสทฺธมฺมสมนฺนาคโต โหติ, อสปฺปุริสภตฺติ\footnote{อสปฺปุริสภตฺตี (สพฺพตฺถ)} โหติ, อสปฺปุริสจินฺตี โหติ, อสปฺปุริสมนฺตี โหติ, อสปฺปุริสวาโจ โหติ, อสปฺปุริสกมฺมนฺโต โหติ, อสปฺปุริสทิฏฺฐิ\footnote{อสปฺปุริสทิฏฺฐี (สพฺพตฺถ)} โหติ, อสปฺปุริสทานํ เทติ.}
\prose{\csromanfolio{71}กถญฺจ ภิกฺขเว อสปฺปุริโส อสฺสทฺธมฺมสมนฺนาคโต โหติ, อิธ ภิกฺขเว อสปฺปุริโส อสฺสทฺโธ โหติ, อหิริโต โหติ, อโนตฺตปฺปี โหติ, อปฺปสฺสุโต โหติ, กุสีโต โหติ, มุฏฺฐสฺสติ โหติ, ทุปฺปญฺโญ โหติ.\csromanspacer{} เอวํ โข ภิกฺขเว อสปฺปุริโส อสฺสทฺธมฺมสมนฺนาคโต โหติ. (1)}

\prose{กถญฺจ ภิกฺขเว อสปฺปุริโส อสปฺปุริสภตฺติ โหติ, อิธ ภิกฺขเว อสปฺปุริสสฺส เย เต สมณพฺราหฺมณา อสฺสทฺธา อหิริกา อโนตฺตปฺปิโน อปฺปสฺสุตา กุสีตา มุฏฺฐสฺสติโน ทุปฺปญฺญา, ตฺยาสฺส มิตฺตา โหนฺติ เต สหายา.\csromanspacer{} เอวํ โข ภิกฺขเว อสปฺปุริโส อสปฺปุริสภตฺติ โหติ. (2)}

\prose{กถญฺจ ภิกฺขเว อสปฺปุริโส อสปฺปุริสจินฺตี โหติ, อิธ ภิกฺขเว อสปฺปุริโส อตฺตพฺยาพาธายปิ เจเตติ, ปรพฺยาพาธายปิ เจเตติ, อุภยพฺยาพาธายปิ เจเตติ.\csromanspacer{} เอวํ โข ภิกฺขเว อสปฺปุริโส อสปฺปุริสจินฺตี โหติ. (3)}

\prose{กถญฺจ ภิกฺขเว อสปฺปุริโส อสปฺปุริสมนฺตี โหติ, อิธ ภิกฺขเว อสปฺปุริโส อตฺตพฺยาพาธายปิ มนฺเตติ, ปรพฺยาพาธายปิ มนฺเตติ, อุภยพฺยาธายปิ มนฺเตติ.\csromanspacer{} เอวํ โข ภิกฺขเว อสปฺปุริโส อสปฺปุริสมนฺตี โหติ. (4)}

\prose{กถญฺจ ภิกฺขเว อสปฺปุริโส อสปฺปุริสวาโจ โหติ, อิธ ภิกฺขเว อสปฺปุริโส มุสาวาที โหติ, ปิสุณวาโจ โหติ, ผรุสวาโจ \csromanfolio{72}โหติ, สมฺผปฺปลาปี โหติ.\csromanspacer{} เอวํ โข ภิกฺขเว อสปฺปุริโส อสปฺปุริสวาโจ โหติ. (5)}

\prose{กถญฺจ ภิกฺขเว อสปฺปุริโส อสปฺปุริสกมฺมนฺโต โหติ, อิธ ภิกฺขเว อสปฺปุริโส ปาณาติปาตี โหติ, อทินฺนาทายี โหติ, กาเมสุมิจฺฉาจารี โหติ.\csromanspacer{} เอวํ โข ภิกฺขเว อสปฺปุริโส อสปฺปุริสกมฺมนฺโต โหติ. (6)}

\prose{กถญฺจ ภิกฺขเว อสปฺปุริโส อสปฺปุริสทิฏฺฐิ โหติ, อิธ ภิกฺขเว อสปฺปุริโส เอวํทิฏฺฐิ\footnote{เอวํทิฏฺฐี (สี, อิ), เอวํทิฏฺฐิโก (สฺยา, กํ)} โหติ “นตฺถิ ทินฺนํ, นตฺถิ ยิฏฺฐํ, นตฺถิ หุตํ, นตฺถิ สุกตทุกฺกฏานํ\footnote{สุกฺกฏทุกฺกฏานํ (สี, อิ)} กมฺมานํ ผลํ วิปาโก, นตฺถิ อยํ โลโก, นตฺถิ ปโร โลโก, นตฺถิ มาตา, นตฺถิ ปิตา, นตฺถิ สตฺตา โอปปาติกา, นตฺถิ โลเก สมณพฺราหฺมณา สมฺมคฺคตา\footnote{สมคฺคตา (ก)} สมฺมาปฏิปนฺนา, เย อิมญฺจ โลกํ ปรญฺจ โลกํ สยํ อภิญฺญา สจฺฉิกตฺวา ปเวเทนฺตี”ติ.\csromanspacer{} เอวํ โข ภิกฺขเว อสปฺปุริโส อสปฺปุริสทิฏฺฐิ โหติ. (7)}

\prose{กถญฺจ ภิกฺขเว อสปฺปุริโส อสปฺปุริสทานํ เทติ, อิธ ภิกฺขเว อสปฺปุริโส อสกฺกจฺจํ ทานํ เทติ, อสหตฺถา ทานํ เทติ, อจิตฺตีกตฺวา ทานํ เทติ, อปวิฏฺฐํ ทานํ เทติ, อนาคมนทิฏฺฐิโก ทานํ เทติ.\csromanspacer{} เอวํ โข ภิกฺขเว อสปฺปุริโส อสปฺปุริสทานํ เทติ. (8)}

\prose{โส ภิกฺขเว อสปฺปุริโส เอวํ อสฺสทฺธมฺมสมนฺนาคโต เอวํ อสปฺปุริสภตฺติ เอวํ อสปฺปุริสจินฺตี เอวํ อสปฺปุริสมนฺตี เอวํ อสปฺปุริสวาโจ เอวํ อสปฺปุริสกมฺมนฺโต เอวํ อสปฺปุริสทิฏฺฐิ เอวํ อสปฺปุริสทานํ ทตฺวา กายสฺส เภทา ปรํ มรณา ยา อสปฺปุริสานํ คติ, ตตฺถ อุปปชฺชติ.\csromanspacer{} กา จ ภิกฺขเว อสปฺปุริสานํ คติ, นิรโย วา ติรจฺฉานโยนิ วา.}

\proseitem{92}{ชาเนยฺย นุ โข ภิกฺขเว สปฺปุริโส สปฺปุริสํ “สปฺปุริโส อยํ ภวนฺ”ติ.\csromanspacer{} เอวํ ภนฺเต.\csromanspacer{} สาธุ ภิกฺขเว, ฐานเมตํ ภิกฺขเว วิชฺชติ, ยํ สปฺปุริโส สปฺปุริสํ ชาเนยฺย “สปฺปุริโส อยํ ภวนฺ”ติ.\csromanspacer{} ชาเนยฺย ปน ภิกฺขเว สปฺปุริโส อสปฺปุริสํ “อสปฺปุริโส อยํ ภวนฺ”ติ.\csromanspacer{} เอวํ ภนฺเต.\csromanspacer{} สาธุ ภิกฺขเว, เอตมฺปิ โข ภิกฺขเว ฐานํ วิชฺชติ, ยํ สปฺปุริโส อสปฺปุริสํ ชาเนยฺย “อสปฺปุริโส อยํ ภวนฺ”ติ.\csromanspacer{} สปฺปุริโส ภิกฺขเว สทฺธมฺมสมนฺนาคโต โหติ, สปฺปุริสภตฺติ โหติ,}

\vspace{\tipitakaparskip}
\csromancenter{จูฬปุณฺณมสุตฺต (110) สปฺปุริสจินฺตี โหติ, สปฺปุริสมนฺตี โหติ, สปฺปุริสวาโจ โหติ, สปฺปุริสกมฺมนฺโต โหติ, สปฺปุริสทิฏฺฐิ โหติ, สปฺปุริสทานํ เทติ.}
\prose{\csromanfolio{73}กถญฺจ ภิกฺขเว สปฺปุริโส สทฺธมฺมสมนฺนาคโต โหติ, อิธ ภิกฺขเว สปฺปุริโส สทฺโธ โหติ, หิริมา โหติ, โอตฺตปฺปี โหติ, พหุสฺสุโต โหติ, อารทฺธวีริโย โหติ, อุปฏฺฐิตสฺสติ โหติ, ปญฺญวา โหติ.\csromanspacer{} เอวํ โข ภิกฺขเว สปฺปุริโส สทฺธมฺมสมนฺนาคโต โหติ. (1)}

\prose{กถญฺจ ภิกฺขเว สปฺปุริโส สปฺปุริสภตฺติ โหติ, อิธ ภิกฺขเว สปฺปุริสสฺส เย เต สมณพฺราหฺมณา สทฺธา หิริมนฺโต โอตฺตปฺปิโน พหุสฺสุตา อารทฺธวีริยา อุปฏฺฐิตสฺสติโน ปญฺญวนฺโต, ตฺยาสฺส มิตฺตา โหนฺติ เต สหายา.\csromanspacer{} เอวํ โข ภิกฺขเว สปฺปุริโส สปฺปุริสภตฺติ โหติ. (2)}

\prose{กถญฺจ ภิกฺขเว สปฺปุริโส สปฺปุริสจินฺตี โหติ, อิธ ภิกฺขเว สปฺปุริโส เนวตฺตพฺยาพาธาย เจเตติ, น ปรพฺยาพาธาย เจเตติ, น อุภยพฺยาพาธาย เจเตติ.\csromanspacer{} เอวํ โข ภิกฺขเว สปฺปุริโส สปฺปุริสจินฺตี โหติ. (3)}

\prose{กถญฺจ ภิกฺขเว สปฺปุริโส สปฺปุริสมนฺตี โหติ, อิธ ภิกฺขเว สปฺปุริโส เนวตฺตพฺยาธาย มนฺเตติ, น ปรพฺยาพาธาย มนฺเตติ, น อุภยพฺยาพาธาย มนฺเตติ.\csromanspacer{} เอวํ โข ภิกฺขเว สปฺปุริโส สปฺปุริสมนฺตี โหติ. (4)}

\prose{กถญฺจ ภิกฺขเว สปฺปุริโส สปฺปุริสวาโจ โหติ, อิธ ภิกฺขเว สปฺปุริโส มุสาวาทา ปฏิวิรโต โหติ, ปิสุณาย วาจาย ปฏิวิรโต โหติ, ผรุสาย วาจาย ปฏิวิรโต โหติ, สมฺผปฺปลาปา ปฏิวิรโต โหติ.\csromanspacer{} เอวํ โข ภิกฺขเว สปฺปุริโส สปฺปุริสวาโจ โหติ. (5)}

\prose{กถญฺจ ภิกฺขเว สปฺปุริโส สปฺปุริสกมฺมนฺโต โหติ, อิธ ภิกฺขเว สปฺปุริโส ปาณาติปาตา ปฏิวิรโต โหติ, อทินฺนาทานา ปฏิวิรโต โหติ, กาเมสุมิจฺฉาจารา ปฏิวิรโต โหติ.\csromanspacer{} เอวํ โข ภิกฺขเว สปฺปุริโส สปฺปุริสกมฺมนฺโต โหติ. (6)}

\csromanmatikaanchor{mtk.12}
\csromantocmarkref{subsection}{อุทฺทานคาถา}{mtk.12}
\bookmark[dest={mtk.12},level=2]{อุทฺทานคาถา}
\prose{\csromanfolio{74}กถญฺจ ภิกฺขเว สปฺปุริโส สปฺปุริสทิฏฺฐิ โหติ, อิธ ภิกฺขเว สปฺปุริโส เอวํทิฏฺฐิ โหติ “อตฺถิ ทินฺนํ, อตฺถิ ยิฏฺฐํ, อตฺถิ หุตํ, อตฺถิ สุกตทุกฺกฏานํ กมฺมานํ ผลํ วิปาโก, อตฺถิ อยํ โลโก, อตฺถิ ปโร โลโก, อตฺถิ มาตา, อตฺถิ ปิตา, อตฺถิ สตฺตา โอปปาติกา, อตฺถิ โลเก สมณพฺราหฺมณา สมฺมคฺคตา สมฺมาปฏิปนฺนา, เย อิมญฺจ โลกํ ปรญฺจ โลกํ สยํ อภิญฺญา สจฺฉิกตฺวา ปเวเทนฺตี”ติ.\csromanspacer{} เอวํ โข ภิกฺขเว สปฺปุริโส สปฺปุริสทิฏฺฐิ โหติ. (7)}

\prose{กถญฺจ ภิกฺขเว สปฺปุริโส สปฺปุริสทานํ เทติ, อิธ ภิกฺขเว สปฺปุริโส สกฺกจฺจํ ทานํ เทติ, สหตฺถา ทานํ เทติ, จิตฺตีกตฺวา ทานํ เทติ, อนปวิฏฺฐํ ทานํ เทติ, อาคมนทิฏฺฐิโก ทานํ เทติ.\csromanspacer{} เอวํ โข ภิกฺขเว สปฺปุริโส สปฺปุริสทานํ เทติ. (8)}

\prose{โส ภิกฺขเว สปฺปุริโส เอวํ สทฺธมฺมสมนฺนาคโต เอวํ สปฺปุริสภตฺติ เอวํ สปฺปุริสจินฺตี เอวํ สปฺปุริสมนฺตี เอวํ สปฺปุริสวาโจ เอวํ สปฺปุริสกมฺมนฺโต เอวํ สปฺปุริสทิฏฺฐิ เอวํ สปฺปุริสทานํ ทตฺวา กายสฺส เภทา ปรํ มรณา ยา สปฺปุริสานํ คติ, ตตฺถ อุปปชฺชติ.\csromanspacer{} กา จ ภิกฺขเว สปฺปุริสานํ คติ, เทวมหตฺตตา วา มนุสฺสมหตฺตตา วาติ.}

\prose{อิทมโวจ ภควา.\csromanspacer{} อตฺตมนา เต ภิกฺขู ภควโต ภาสิตํ อภินนฺทุนฺติ.}

\nitthitam{จูฬปุณฺณมสุตฺตํ นิฏฺฐิตํ ทสมํ.}
\nitthitam{\textbf{เทวทหวคฺโค นิฏฺฐิโต ปฐโม.}}
\titlehead{\textbf{ตสฺสุทฺทานํ}}
\csromangathameasure{เทวทหํ ปญฺจตฺตยํ, กินฺติ สาม สุนกฺขตฺตํ. \\ สปฺปายคณโคปก, มหาปุณฺณจูฬปุณฺณญฺจาติ.}
\csromangathagroup{\csromangathastanza{เทวทหํ ปญฺจตฺตยํ, กินฺติ สาม สุนกฺขตฺตํ. \\ สปฺปายคณโคปก, มหาปุณฺณจูฬปุณฺณญฺจาติ.}}

\csromanmatikaanchor{mtk.13}
\csromantocmarknopage{chapter}{2. อนุปทวคฺค}{mtk.13}
\bookmark[dest={mtk.13},level=0]{2. อนุปทวคฺค}
\chapterheadpageread{2. อนุปทวคฺค}
\csromanmatikaanchor{mtk.14}
\csromantocmarkref{subsection}{1. อนุปทสุตฺต}{mtk.14}
\bookmark[dest={mtk.14},level=2]{1. อนุปทสุตฺต}
\csromanheader{2}{1. อนุปทสุตฺต}
\proseitem{93}{\csromanfolio{75}เอวํ เม สุตํ\csromandash{}เอกํ สมยํ ภควา สาวตฺถิยํ วิหรติ เชตวเน อนาถปิณฺฑิกสฺส อาราเม.\csromanspacer{} ตตฺร โข ภควา ภิกฺขู อามนฺเตสิ “ภิกฺขโว”ติ. “ภทนฺเต”ติ เต ภิกฺขู ภควโต ปจฺจสฺโสสุํ.\csromanspacer{} ภควา เอตทโวจ\csromandash{}}

\prose{ปณฺฑิโต ภิกฺขเว สาริปุตฺโต, มหาปญฺโญ ภิกฺขเว สาริปุตฺโต, ปุถุปญฺโญ ภิกฺขเว สาริปุตฺโต, หาสปญฺโญ\footnote{หาสุปญฺโญ (สี, อิ)} ภิกฺขเว สาริปุตฺโต, ชวนปญฺโญ ภิกฺขเว สาริปุตฺโต, ติกฺขปญฺโญ ภิกฺขเว สาริปุตฺโต, นิพฺเพธิกปญฺโญ ภิกฺขเว สาริปุตฺโต.\csromanspacer{} สาริปุตฺโต ภิกฺขเว อฑฺฒมาสํ อนุปทธมฺมวิปสฺสนํ วิปสฺสติ.\csromanspacer{} ตตฺริทํ ภิกฺขเว สาริปุตฺตสฺส อนุปทธมฺมวิปสฺสนาย โหติ.}

\proseitem{94}{อิธ ภิกฺขเว สาริปุตฺโต วิวิจฺเจว กาเมหิ วิวิจฺจ อกุสเลหิ ธมฺเมหิ สวิตกฺกํ สวิจารํ วิเวกชํ ปีติสุขํ ปฐมํ ฌานํ อุปสมฺปชฺช วิหรติ.\csromanspacer{} เย จ ปฐเม ฌาเน\footnote{ปฐมชฺฌาเน (ก-สี, อิ, ก)} ธมฺมา วิตกฺโก วิจาโร จ ปีติ จ สุขญฺจ จิตฺเตกคฺคตา จ ผสฺโส เวทนา สญฺญา เจตนา จิตฺตํ ฉนฺโท อธิโมกฺโข วีริยํ สติ อุเปกฺขา มนสิกาโร.\csromanspacer{} ตฺยาสฺส ธมฺมา อนุปทววตฺถิตา โหนฺติ, ตฺยาสฺส ธมฺมา วิทิตา อุปฺปชฺชนฺติ, วิทิตา อุปฏฺฐหนฺติ, วิทิตา อพฺภตฺถํ คจฺฉนฺติ.\csromanspacer{} โส เอวํ ปชานาติ “เอวํ กิรเม ธมฺมา อหุตฺวา สมฺโภนฺติ, หุตฺวา ปฏิเวนฺตี”ติ.\csromanspacer{} โส เตสุ ธมฺเมสุ อนุปาโย อนปาโย อนิสฺสิโต อปฺปฏิพทฺโธ\footnote{อปฺปฏิพนฺโธ (ก)} วิปฺปมุตฺโต วิสํยุตฺโต วิมริยาทีกเตน เจตสา วิหรติ, โส “อตฺถิ อุตฺตริ นิสฺสรณนฺ”ติ ปชานาติ.\csromanspacer{} ตพฺพหุลีการา อตฺถิเตฺววสฺส\footnote{อตฺถิเตวสฺส (สี, อิ)} โหติ. (1)}

\prose{ปุน จปรํ ภิกฺขเว สาริปุตฺโต วิตกฺกวิจารานํ วูปสมา อชฺฌตฺตํ สมฺปสาทนํ เจตโส เอโกทิภาวํ อวิตกฺกํ อวิจารํ สมาธิชํ ปีติสุขํ ทุติยํ ฌานํ อุปสมฺปชฺช วิหรติ.\csromanspacer{} เย จ ทุติเย ฌาเน ธมฺมา อชฺฌตฺตํ สมฺปสาโท จ ปีติ จ สุขญฺจ จิตฺเตกคฺคตา จ ผสฺโส เวทนา สญฺญา \csromanfolio{76}เจตนา จิตฺตํ ฉนฺโท อธิโมกฺโข วีริยํ สติ อุเปกฺขา มนสิกาโร.\csromanspacer{} ตฺยาสฺส ธมฺมา อนุปทววตฺถิตา โหนฺติ, ตฺยาสฺส ธมฺมา วิทิตา อุปฺปชฺชนฺติ, วิทิตา อุปฏฺฐหนฺติ, วิทิตา อพฺภตฺถํ คจฺฉนฺติ.\csromanspacer{} โส เอวํ ปชานาติ “เอวํ กิร เม ธมฺมา อหุตฺวา สมฺโภนฺติ, หุตฺวา ปฏิเวนฺตี”ติ.\csromanspacer{} โส เตสุ ธมฺเมสุ อนุปาโย อนปาโย อนิสฺสิโต อปฺปฏิพทฺโธ วิปฺปมุตฺโต วิสํยุตฺโต วิมริยาทีกเตน เจตสา วิหรติ, โส “อตฺถิ อุตฺตริ นิสฺสรณนฺ”ติ ปชานาติ.\csromanspacer{} ตพฺพหุลีการา อตฺถิเตฺววสฺส โหติ. (2)}

\prose{ปุน จปรํ ภิกฺขเว สาริปุตฺโต ปีติยา จ วิราคา อุเปกฺขโก จ วิหรติ, สโต จ สมฺปชาโน สุขญฺจ กาเยน ปฏิสํเวเทติ.\csromanspacer{} ยํ ตํ อริยา อาจิกฺขนฺติ “อุเปกฺขโก สติมา สุขวิหารี”ติ, ตติยํ ฌานํ อุปสมฺปชฺช วิหรติ.\csromanspacer{} เย จ ตติเย ฌาเน ธมฺมา สุขญฺจ สติ จ สมฺปชญฺญญฺจ จิตฺเตกคฺคตา จ ผสฺโส เวทนา สญฺญา เจตนา จิตฺตํ ฉนฺโท อธิโมกฺโข วีริยํ สติ อุเปกฺขา มนสิกาโร.\csromanspacer{} ตฺยาสฺส ธมฺมา อนุปทววตฺถิตา โหนฺติ, ตฺยาสฺส ธมฺมา วิทิตา อุปฺปชฺชนฺติ, วิทิตา อุปฏฺฐหนฺติ, วิทิตา อพฺภตฺถํ คจฺฉนฺติ.\csromanspacer{} โส เอวํ ปชานาติ “เอวํ กิรเม ธมฺมา อหุตฺวา สมฺโภนฺติ, หุตฺวา ปฏิเวนฺตี”ติ.\csromanspacer{} โส เตสุ ธมฺเมสุ อนุปาโย อนปาโย อนิสฺสิโต อปฺปฏิพทฺโธ วิปฺปมุตฺโต วิสํยุตฺโต วิมริยาทีกเตน เจตสา วิหรติ, โส “อตฺถิ อุตฺตริ นิสฺสรณนฺ”ติ ปชานาติ.\csromanspacer{} ตพฺพหุลีการา อตฺถิเตฺววสฺส โหติ. (3)}

\prose{ปุน จปรํ ภิกฺขเว สาริปุตฺโต สุขสฺส จ ปหานา ทุกฺขสฺส จ ปหานา ปุพฺเพว โสมนสฺสโทมนสฺสานํ อตฺถงฺคมา อทุกฺขมสุขํ อุเปกฺขา สติปาริสุทฺธิํ จตุตฺถํ ฌานํ อุปสมฺปชฺช วิหรติ.\csromanspacer{} เย จ จตุตฺเถ ฌาเน ธมฺมา อุเปกฺขา อทุกฺขมสุขา เวทนา ปสฺสทฺธตฺตา เจตโส อนาโภโค สติปาริสุทฺธิ จิตฺเตกคฺคตา จ ผสฺโส เวทนา สญฺญา เจตนา จิตฺตํ ฉนฺโท อธิโมกฺโข วีริยํ สติ อุเปกฺขา มนสิกาโร.\csromanspacer{} ตฺยาสฺส ธมฺมา อนุปทววตฺถิตา โหนฺติ, ตฺยาสฺส ธมฺมา วิทิตา อุปฺปชฺชนฺติ, วิทิตา อุปฏฺฐหนฺติ, วิทิตา อพฺภตฺถํ คจฺฉนฺติ.\csromanspacer{} โส เอวํ ปชานาติ “เอวํ กิรเม ธมฺมา อหุตฺวา สมฺโภนฺติ, หุตฺวา ปฏิเวนฺตี”ติ.\csromanspacer{} โส เตสุ ธมฺเมสุ อนุปาโย อนปาโย อนิสฺสิโต อปฺปฏิพทฺโธ วิปฺปมุตฺโต วิสํยุตฺโต วิมริยาทิกเตน เจตสา วิหรติ, โส “อตฺถิ อุตฺตริ นิสฺสรณนฺ”ติ ปชานาติ.\csromanspacer{} ตพฺพหุลีการา อตฺถิเตฺววสฺส โหติ. (4)}

\csromanheader{2}{1. อนุปทสุตฺต (111)}
\prose{\csromanfolio{77}ปุน จปรํ ภิกฺขเว สาริปุตฺโต สพฺพโส รูปสญฺญานํ สมติกฺกมา ปฏิฆสญฺญานํ อตฺถงฺคมา นานตฺตสญฺญานํ อมนสิการา “อนนฺโต อากาโส”ติ อากาสานญฺจายตนํ อุปสมฺปชฺช วิหรติ.\csromanspacer{} เย จ อากาสานญฺจายตเน ธมฺมา อากาสานญฺจายตสญฺญา จ จิตฺเตกคฺคตา จ ผสฺโส เวทนา สญฺญา เจตนา จิตฺตํ ฉนฺโท อธิโมกฺโข วีริยํ สติ อุเปกฺขา มนสิกาโร.\csromanspacer{} ตฺยาสฺส ธมฺมา อนุปทววตฺถิตา โหนฺติ, ตฺยาสฺส ธมฺมา วิทิตา อุปฺปชฺชนฺติ, วิทิตา อุปฏฺฐหนฺติ, วิทิตา อพฺภตฺถํ คจฺฉนฺติ.\csromanspacer{} โส เอวํ ปชานาติ “เอวํ กิรเม ธมฺมา อหุตฺวา สมฺโภนฺติ, หุตฺวา ปฏิเวนฺตี”ติ.\csromanspacer{} โส เตสุ ธมฺเมสุ อนุปาโย อนปาโย อนิสฺสิโต อปฺปฏิพทฺโธ วิปฺปมุตฺโต วิสํยุตฺโต วิมริยาทีกเตน เจตสา วิหรติ, โส “อตฺถิ อุตฺตริ นิสฺสรณนฺ”ติ ปชานาติ.\csromanspacer{} ตพฺพหุลีการา อตฺถิเตฺววสฺส โหติ. (5)}

\prose{ปุน จปรํ ภิกฺขเว สาริปุตฺโต สพฺพโส อากาสานญฺจายตนํ สมติกฺกมฺม “อนนฺตํ วิญฺญาณนฺ”ติ วิญฺญาณญฺจายตนํ อุปสมฺปชฺช วิหรติ.\csromanspacer{} เย จ วิญฺญาณญฺจายตเน ธมฺมา วิญฺญาณญฺจายตนสญฺญา จ จิตฺเตกคฺคตา จ ผสฺโส เวทนา สญฺญา เจตนา จิตฺตํ ฉนฺโท อธิโมกฺโข วีริยํ สติ อุเปกฺขา มนสิกาโร.\csromanspacer{} ตฺยาสฺส ธมฺมา อนุปทววตฺถิตา โหนฺติ, ตฺยาสฺส ธมฺมา วิทิตา อุปฺปชฺชนฺติ, วิทิตา อุปฏฺฐหนฺติ, วิทิตา อพฺภตฺถํ คจฺฉนฺติ.\csromanspacer{} โส เอวํ ปชานาติ “เอวํ กิรเม ธมฺมา อหุตฺวา สมฺโภนฺติ, หุตฺวา ปฏิเวนฺตี”ติ.\csromanspacer{} โส เตสุ ธมฺเมสุ อนุปาโย อนปาโย อนิสฺสิโต อปฺปฏิพทฺโธ วิปฺปมุตฺโต วิสํยุตฺโต วิมริยาทีกเตน เจตสา วิหรติ, โส “อตฺถิ อุตฺตริ นิสฺสรณนฺ”ติ ปชานาติ.\csromanspacer{} ตพฺพหุลีการา อตฺถิเตฺววสฺส โหติ.(6)}

\prose{ปุน จปรํ ภิกฺขเว สาริปุตฺโต สพฺพโส วิญฺญาณญฺจายตนํ สมติกฺกมฺม “นตฺถิ กิญฺจี”ติ อากิญฺจญฺญายตนํ อุปสมฺปชฺช วิหรติ.\csromanspacer{} เย จ อากิญฺจญฺญายตเน ธมฺมา อากิญฺจญฺญายตนสญฺญา จ จิตฺเตกคฺคตา จ ผสฺโส เวทนา สญฺญา เจตนา จิตฺตํ ฉนฺโท อธิโมกฺโข วีริยํ สติ อุเปกฺขา มนสิกาโร.\csromanspacer{} ตฺยาสฺส ธมฺมา อนุปทววตฺถิตา โหนฺติ, ตฺยาสฺส ธมฺมา วิทิตา อุปฺปชฺชนฺติ, วิทิตา อุปฏฺฐหนฺติ, วิทิตา อพฺภตฺถํ คจฺฉนฺติ.\csromanspacer{} โส เอวํ ปชานาติ “เอวํ กิรเม ธมฺมา อหุตฺวา สมฺโภนฺติ, หุตฺวา ปฏิเวนฺตี”ติ.\csromanspacer{} โส เตสุ ธมฺเมสุ อนุปาโย อนปาโย อนิสฺสิโต อปฺปฏิพทฺโธ}

\prose{\csromanfolio{78}วิปฺปมุตฺโต วิสํยุตฺโต วิมริยาทีกเตน เจตสา วิหรติ, โส “อตฺถิ อุตฺตริ นิสฺสรณนฺ”ติ ปชานาติ.\csromanspacer{} ตพฺพหุลีการา อตฺถิเตฺววสฺส โหติ. (7)}

\proseitem{95}{ปุน จปรํ ภิกฺขเว สาริปุตฺโต สพฺพโส อากิญฺจญฺญายตนํ สมติกฺกมฺม เนวสญฺญานาสญฺญายตนํ อุปสมฺปชฺช วิหรติ.\csromanspacer{} โส ตาย สมาปตฺติยา สโต วุฏฺฐหติ, โส ตาย สมาปตฺติยา สโต วุฏฺฐหิตฺวา เย ธมฺมา\footnote{เย เต ธมฺมา (สี)} อตีตา นิรุทฺธา วิปริณตา, เต ธมฺเม สมนุปสฺสติ “เอวํ กิรเม ธมฺมา อหุตฺวา สมฺโภนฺติ, หุตฺวา ปฏิเวนฺตี”ติ.\csromanspacer{} โส เตสุ ธมฺเมสุ อนุปาโย อนปาโย อนิสฺสิโต อปฺปฏิพทฺโธ วิปฺปมุตฺโต วิสํยุตฺโต วิมริยาทีกเตน เจตสา วิหรติ, โส “อตฺถิ อุตฺตริ นิสฺสรณนฺ”ติ ปชานาติ.\csromanspacer{} ตพฺพหุลีการา อตฺถิเตฺววสฺส โหติ. (8)}

\proseitem{96}{ปุน จปรํ ภิกฺขเว สาริปุตฺโต สพฺพโส เนวสญฺญานาสญฺญายตนํ สมติกฺกมฺม สญฺญาเวทยิตนิโรธํ อุปสมฺปชฺช วิหรติ, ปญฺญาย จสฺส ทิสฺวา อาสวา ปริกฺขีณา โหนฺติ.\csromanspacer{} โส ตาย สมาปตฺติยา สโต วุฏฺฐหติ, โส ตาย สมาปตฺติยา สโต วุฏฺฐหิตฺวา เย ธมฺมา1 อตีตา นิรุทฺธา วิปริณตา, เต ธมฺเม สมนุปสฺสติ “เอวํ กิรเม ธมฺมา อหุตฺวา สมฺโภนฺติ, หุตฺวา ปฏิเวนฺตี”ติ.\csromanspacer{} โส เตสุ ธมฺเมสุ อนุปาโย อนปาโย อนิสฺสิโต อปฺปฏิพทฺโธ วิปฺปมุตฺโต วิสํยุตฺโต วิมริยาทีกเตน เจตสา วิหรติ, โส “นตฺถิ อุตฺตริ นิสฺสรณนฺ”ติ ปชานาติ.\csromanspacer{} ตพฺพหุลีการา นตฺถิเตฺววสฺส โหติ. (9)}

\proseitem{97}{ยํ โข ตํ ภิกฺขเว สมฺมา วทมาโน วเทยฺย “วสิปฺปตฺโต ปารมิปฺปตฺโต อริยสฺมิํ สีลสฺมิํ, วสิปฺปตฺโต ปารมิปฺปตฺโต อริยสฺมิํ สมาธิสฺมิํ, วสิปฺปตฺโต ปารมิปฺปตฺโต อริยาย ปญฺญาย, วสิปฺปตฺโต ปารมิปฺปตฺโต อริยาย วิมุตฺติยา”ติ.\csromanspacer{} สาริปุตฺตเมว ตํ สมฺมา วทมาโน วเทยฺย “วสิปฺปตฺโต ปารมิปฺปตฺโต อริยสฺมิํ สีลสฺมิํ, วสิปฺปตฺโต ปารมิปฺปตฺโต อริยสฺมิํ สมาธิสฺมิํ, วสิปฺปตฺโต ปารมิปฺปตฺโต อริยาย ปญฺญาย, วสิปฺปตฺโต ปารมิปฺปตฺโต อริยาย วิมุตฺติยา”ติ.\csromanspacer{} ยํ โข ตํ}

\vspace{\tipitakaparskip}
\csromancenter{ฉพฺพิโสธนสุตฺต (112) ภิกฺขเว สมฺมา วทมาโน วเทยฺย “ภควโต ปุตฺโต โอรโส มุขโต ชาโต ธมฺมโช ธมฺมนิมฺมิโต ธมฺมทายาโท, โน อามิสทายาโท”ติ.\csromanspacer{} สาริปุตฺตเมว ตํ สมฺมา วทมาโน วเทยฺย “ภควโต ปุตฺโต โอรโส มุขโต ชาโต ธมฺมโช ธมฺมนิมฺมิโต ธมฺมทายาโท, โน อามิสทายาโท”ติ.\csromanspacer{} สาริปุตฺโต ภิกฺขเว ตถาคเตน อนุตฺตรํ ธมฺมจกฺกํ ปวตฺติตํ สมฺมเทว อนุปฺปวตฺเตตีติ.}
\csromancenter{อิทมโวจ ภควา.\csromanspacer{} อตฺตมนา เต ภิกฺขู ภควโต ภาสิตํ อภินนฺทุนฺติ.}
\nitthitamruled{อนุปทสุตฺตํ นิฏฺฐิตํ ปฐมํ.}
\csromanmatikaanchor{mtk.15}
\csromantocmarkref{subsection}{2. ฉพฺพิโสธนสุตฺต}{mtk.15}
\bookmark[dest={mtk.15},level=2]{2. ฉพฺพิโสธนสุตฺต}
\csromanheader{2}{2. ฉพฺพิโสธนสุตฺต}
\proseitem{98}{\csromanfolio{79}เอวํ เม สุตํ\csromandash{}เอกํ สมยํ ภควา สาวตฺถิยํ วิหรติ เชตวเน อนาถปิณฺฑิกสฺส อาราเม.\csromanspacer{} ตตฺร โข ภควา ภิกฺขู อามนฺเตสิ “ภิกฺขโว”ติ. “ภทนฺเต”ติ เต ภิกฺขู ภควโต ปจฺจสฺโสสุํ.\csromanspacer{} ภควา เอตทโวจ\csromandash{}}

\prose{อิธ ภิกฺขเว ภิกฺขุ อญฺญํ พฺยากโรติ “ขีณา ชาติ, วุสิตํ พฺรหฺมจริยํ, กตํ กรณียํ, ‘นาปรํ อิตฺถตฺตายา’ติ ปชานามี”ติ.\csromanspacer{} ตสฺส ภิกฺขเว ภิกฺขุโน ภาสิตํ เนว อภินนฺทิตพฺพํ นปฺปฏิกฺโกสิตพฺพํ.\csromanspacer{} อนภินนฺทิตฺวา อปฺปฏิกฺโกสิตฺวา ปโญฺห ปุจฺฉิตพฺโพ.\csromanspacer{} จตฺตาโรเม อาวุโส โวหารา เตน ภควตา ชานตา ปสฺสตา อรหตา สมฺมาสมฺพุทฺเธน สมฺมทกฺขาตา.\csromanspacer{} กตเม จตฺตาโร, ทิฏฺเฐ ทิฏฺฐวาทิตา สุเต สุตวาทิตา มุเต มุตวาทิตา วิญฺญาเต วิญฺญาตวาทิตา.\csromanspacer{} อิเม โข อาวุโส จตฺตาโร โวหารา เตน ภควตา ชานตา ปสฺสตา อรหตา สมฺมาสมฺพุทฺเธน สมฺมทกฺขาตา.\csromanspacer{} กถํ ชานโต ปนายสฺมโต กถํ ปสฺสโต อิเมสุ จตูสุ โวหาเรสุ อนุปาทาย อาสเวหิ จิตฺตํ วิมุตฺตนฺติ.\csromanspacer{} ขีณาสวสฺส ภิกฺขเว ภิกฺขุโน วุสิตวโต กตกรณียสฺส โอหิตภารสฺส อนุปฺปตฺตสทตฺถสฺส ปริกฺขีณภวสํโยชนสฺส สมฺมทญฺญาวิมุตฺตสฺส อยมนุธมฺโม โหติ เวยฺยากรณาย.\csromanspacer{} ทิฏฺเฐ โข อหํ อาวุโส อนุปาโย อนปาโย \csromanfolio{80}อนิสฺสิโต อปฺปฏิพทฺโธ วิปฺปมุตฺโต วิสํยุตฺโต วิมริยาทีกเตน เจตสา วิหรามิ.\csromanspacer{} สุเต โข อหํ อาวุโส ฯเปฯ.\csromanspacer{} มุเต โข อหํ อาวุโส ฯเปฯ.\csromanspacer{} วิญฺญาเต โข อหํ อาวุโส อนุปาโย อนปาโย อนิสฺสิโต อปฺปฏิพทฺโธ วิปฺปมุตฺโต วิสํยุตฺโต วิมริยาทีกเตน เจตสา วิหรามิ.\csromanspacer{} เอวํ โข เม อาวุโส ชานโต เอวํ ปสฺสโต อิเมสุ จตูสุ โวหาเรสุ อนุปาทาย อาสเวหิ จิตฺตํ วิมุตฺตนฺติ.\csromanspacer{} ตสฺส ภิกฺขเว ภิกฺขุโน “สาธู”ติ ภาสิตํ อภินนฺทิตพฺพํ อนุโมทิตพฺพํ, “สาธู”ติ ภาสิตํ อภินนฺทิตฺวา อนุโมทิตฺวา อุตฺตริํ ปโญฺห ปุจฺฉิตพฺโพ.}

\proseitem{99}{ปญฺจิเม อาวุโส อุปาทานกฺขนฺธา เตน ภควตา ชานตา ปสฺสตา อรหตา สมฺมาสมฺพุทฺเธน สมฺมทกฺขาตา.\csromanspacer{} กตเม ปญฺจ, เสยฺยถิทํ, รูปุปาทานกฺขนฺโธ เวทนุปาทานกฺขนฺโธ สญฺญุปาทานกฺขนฺโธ สงฺขารุปาทานกฺขนฺโธ วิญฺญาณุปาทานกฺขนฺโธ.\csromanspacer{} อิเม โข อาวุโส ปญฺจุปาทานกฺขนฺธา เตน ภควตา ชานตา ปสฺสตา อรหตา สมฺมาสมฺพุทฺเธน สมฺมทกฺขาตา.\csromanspacer{} กถํ ชานโต ปนายสฺมโต กถํ ปสฺสโต อิเมสุ ปญฺจสุ อุปาทานกฺขนฺเธสุ อนุปาทาย อาสเวหิ จิตฺตํ วิมุตฺตนฺติ.\csromanspacer{} ขีณาสวสฺส ภิกฺขเว ภิกฺขุโน วุสิตวโต กตกรณียสฺส โอหิตภารสฺส อนุปฺปตฺตสทตฺถสฺส ปริกฺขีณภวสํโยชนสฺส สมฺมทญฺญาวิมุตฺตสฺส อยมนุธมฺโม โหติ เวยฺยากรณาย.\csromanspacer{} รูปํ โข อหํ อาวุโส อพลํ วิราคุนํ\footnote{วิราคํ (สี, อิ), วิราคุตํ (ฏีกา)} อนสฺสาสิกนฺติ วิทิตฺวา เย รูเป อุปายูปาทานา\footnote{อุปยูปาทานา (ก)} เจตโส อธิฏฺฐานาภินิเวสานุสยา, เตสํ ขยา วิราคา นิโรธา จาคา ปฏินิสฺสคฺคา “วิมุตฺตํ เม จิตฺตนฺ”ติ ปชานามิ.\csromanspacer{} เวทนํ โข อหํ อาวุโส ฯเปฯ.\csromanspacer{} สญฺญํ โข อหํ อาวุโส.\csromanspacer{} สงฺขาเร โข อหํ อาวุโส.\csromanspacer{} วิญฺญาณํ โข อหํ อาวุโส อพลํ วิราคุนํ อนสฺสาสิกนฺติ วิทิตฺวา เย วิญฺญาเณ อุปายูปาทานา เจตโส อธิฏฺฐานาภินิเวสานุสยา, เตสํ ขยา วิราคา นิโรธา จาคา ปฏิสฺสคฺคา “วิมุตฺตํ เม จิตฺตนฺ”ติ ปชานามิ.\csromanspacer{} เอวํ โข เม อาวุโส ชานโต เอวํ ปสฺสโต อิเมสุ ปญฺจสุ อุปาทานกฺขนฺเธสุ อนุปาทาย อาสเวหิ จิตฺตํ วิมุตฺตนฺติ.\csromanspacer{} ตสฺส ภิกฺขเว ภิกฺขุโน “สาธู”ติ ภาสิตํ อภินนฺทิตพฺพํ อนุโมทิตพฺพํ, “สาธู”ติ ภาสิตํ อภินนฺทิตฺวา อนุโมทิตฺวา อุตฺตริํ ปโญฺห ปุจฺฉิตพฺโพ.}

\csromanheader{2}{2. ฉพฺพิโสธนสุตฺต (112)}
\proseitem{100}{\csromanfolio{81}ฉยิมา อาวุโส ธาตุโย เตน ภควตา ชานตา ปสฺสตา อรหตา สมฺมาสมฺพุทฺเธน สมฺมทกฺขาตา.\csromanspacer{} กตมา ฉ, ปถวีธาตุ อาโปธาตุ เตโชธาตุ วาโยธาตุ อากาสธาตุ วิญฺญาณธาตุ.\csromanspacer{} อิมา โข อาวุโส ฉ ธาตุโย เตน ภควตา ชานตา ปสฺสตา อรหตา สมฺมาสมฺพุทฺเธน สมฺมทกฺขาตา.\csromanspacer{} กถํ ชานโต ปนายสฺมโต กถํ ปสฺสโต อิมาสุ ฉสุ ธาตูสุ อนุปาทาย อาสเวหิ จิตฺตํ วิมุตฺตนฺติ.\csromanspacer{} ขีณาสวสฺส ภิกฺขเว ภิกฺขุโน วุสิตวโต กตกรณียสฺส โอหิตภารสฺส อนุปฺปตฺตสทตฺถสฺส ปริกฺขีณภวสํโยชนสฺส สมฺปทญฺญาวิมุตฺตสฺส อยมนุธมฺโม โหติ เวยฺยากรณาย.\csromanspacer{} ปถวีธาตุํ โข อหํ อาวุโส น อตฺตโต อุปคจฺฉิํ, น จ ปถวีธาตุนิสฺสิตํ อตฺตานํ.\csromanspacer{} เย จ ปถวีธาตุนิสฺสิตา อุปายูปาทานา เจตโส อธิฏฺฐานาภินิเวสานุสยา, เตสํ ขยา วิราคา นิโรธา จาคา ปฏินิสฺสคฺคา “วิมุตฺตํ เม จิตฺตนฺ”ติ ปชานามิ.\csromanspacer{} อาโปธาตุํ โข อหํ อาวุโส ฯเปฯ.\csromanspacer{} เตโชธาตุํ โข อหํ อาวุโส.\csromanspacer{} วาโยธาตุํ โข อหํ อาวุโส.\csromanspacer{} อากาสธาตุํ โข อหํ อาวุโส.\csromanspacer{} วิญฺญาณธาตุํ โข อหํ อาวุโส น อตฺตโต อุปคจฺฉิํ, น จ วิญฺญาณธาตุนิสฺสิตํ อตฺตานํ.\csromanspacer{} เย จ วิญฺญาณธาตุนิสฺสิตา อุปายูปาทานา เจตโส อธิฏฺฐานาภินิเวสานุสยา, เตสํ ขยา วิราคา นิโรธา จาคา ปฏินิสฺสคฺคา “วิมุตฺตํ เม จิตฺตนฺ”ติ ปชานามิ.\csromanspacer{} เอวํ โข เม อาวุโส ชานโต เอวํ ปสฺสโต อิมาสุ ฉสุ ธาตูสุ อนุปาทาย อาสเวหิ จิตฺตํ วิมุตฺตนฺติ.\csromanspacer{} ตสฺส ภิกฺขเว ภิกฺขุโน “สาธู”ติ ภาสิตํ อภินนฺทิตพฺพํ อนุโมทิตพฺพํ, “สาธู”ติ ภาสิตํ อภินนฺทิตฺวา อนุโมทิตฺวา อุตฺตริํ ปโญฺห ปุจฺฉิตพฺโพ.}

\proseitem{101}{ฉ โข ปนิมานิ อาวุโส อชฺฌตฺติกพาหิรานิ\footnote{อชฺฌตฺติกานิ พาหิรานิ (สฺยา, กํ, อิ)} อายตนานิ เตน ภควตา ชานตา ปสฺสตา อรหตา สมฺมาสมฺพุทฺเธน สมฺมทกฺขาตานิ.\csromanspacer{} กตมานิ ฉ, จกฺขุ เจว รูปา จ โสตญฺจ สทฺทา จ ฆานญฺจ คนฺธา จ ชิวฺหา จ รสา จ กาโย จ โผฏฺฐพฺพา จ มโน จ ธมฺมา จ.\csromanspacer{} อิมานิ โข อาวุโส ฉ อชฺฌตฺติกพาหิรานิ อายตนานิ เตน ภควตา ชานตา ปสฺสตา อรหตา สมฺมาสมฺพุทฺเธน สมฺมทกฺขาตานิ.\csromanspacer{} กถํ ชานโต ปนายสฺมโต กถํ ปสฺสโต อิเมสุ ฉสุ อชฺฌตฺติกพาหิเรสุ \csromanfolio{82}อายตเนสุ อนุปาทาย อาสเวหิ จิตฺตํ วิมุตฺตนฺติ.\csromanspacer{} ขีณาสวสฺส ภิกฺขเว ภิกฺขุโน วุสิตวโต กตกรณียสฺส โอหิตภารสฺส อนุปฺปตฺตสทตฺถสฺส ปริกฺขีณภวสํโยชนสฺส สมฺมทญฺญาวิมุตฺตสฺส อยมนุธมฺโม โหติ เวยฺยากรณาย.\csromanspacer{} จกฺขุสฺมิํ อาวุโส รูเป จกฺขุวิญฺญาเณ จกฺขุวิญฺญาณวิญฺญาตพฺเพสุ ธมฺเมสุ โย ฉนฺโท โย ราโค ยา นนฺที\footnote{นนฺทิ (สี, สฺยา, กํ, อิ)} ยา ตณฺหา เย จ อุปายูปาทานา เจตโส อธิฏฺฐานาภินิเวสานุสยา, เตสํ ขยา วิราคา นิโรธา จาคา ปฏินิสฺสคฺคา “วิมุตฺตํ เม จิตฺตนฺ”ติ ปชานามิ.\csromanspacer{} โสตสฺมิํ อาวุโส สทฺเท โสตวิญฺญาเณ ฯเปฯ.\csromanspacer{} ฆานสฺมิํ อาวุโส คนฺเธ ฆานวิญฺญาเณ.\csromanspacer{} ชิวฺหาย อาวุโส รเส ชิวฺหาวิญฺญาเณ.\csromanspacer{} กายสฺมิํ อาวุโส โผฏฺฐพฺเพ กายวิญฺญาเณ.\csromanspacer{} มนสฺมิํ อาวุโส ธมฺเม มโนวิญฺญาเณ มโนวิญฺญาณวิญฺญาตพฺเพสุ ธมฺเมสุ โย ฉนฺโท โย ราโค ยา นนฺที ยา ตณฺหา เย จ อุปายูปาทานา เจตโส อธิฏฺฐานาภินิเวสานุสยา, เตสํ ขยา วิราคา นิโรธา จาคา ปฏินิสฺสคฺคา “วิมุตฺตํ เม จิตฺตนฺ”ติ ปชานามิ.\csromanspacer{} เอวํ โข เม อาวุโส ชานโต เอวํ ปสฺสโต อิเมสุ ฉสุ อชฺฌตฺติกพาหิเรสุ อายตเนสุ อนุปาทาย อาสเวหิ จิตฺตํ วิมุตฺตนฺติ.\csromanspacer{} ตสฺส ภิกฺขเว ภิกฺขุโน “สาธู”ติ ภาสิตํ อภินนฺทิตพฺพํ อนุโมทิตพฺพํ, “สาธู”ติ ภาสิตํ อภินนฺทิตฺวา อนุโมทิตฺวา อุตฺตริํ ปโญฺห ปุจฺฉิตพฺโพ.}

\proseitem{102}{กถํ ชานโต ปนายสฺมโต กถํ ปสฺสโต อิมสฺมิํ จ สวิญฺญาณเก กาเย พหิทฺธา จ สพฺพนิมิตฺเตสุ อหงฺการมมงฺการมานานุสยา สมูหตาติ\footnote{สุสมูหตาติ (สี, สฺยา, กํ, อิ)}.\csromanspacer{} ขีณาสวสฺส ภิกฺขเว ภิกฺขุโน วุสิตวโต กตกรณียสฺส โอหิตภารสฺส อนุปฺปตฺตสทตฺถสฺส ปริกฺขีณภวสํโยชนสฺส สมฺมทญฺญาวิมุตฺตสฺส อยมนุธมฺโม โหติ เวยฺยากรณาย.\csromanspacer{} ปุพฺเพ โข อหํ อาวุโส อคาริยภูโต สมาโน อวิทฺทสุ อโหสิํ, ตสฺส เม ตถาคโต วา ตถาคตสาวโก วา ธมฺมํ เทเสสิ, ตาหํ ธมฺมํ สุตฺวา ตถาคเต สทฺธํ ปฏิลภิํ, โส เตน สทฺธาปฏิลาเตน สมนฺนาคโต อิติ ปฏิสญฺจิกฺขิํ “สมฺพาโธ ฆราวาโส รชาปโถ, อพฺโภกาโส ปพฺพชฺชา, น ยิทํ สุกรํ อคารํ อชฺฌาวสตา เอกนฺตปริปุณฺณํ เอกนฺตปริสุทฺธํ สงฺขลิขิตํ พฺรหฺมจริยํ จริตุํ, ยํนูนาหํ เกสมสฺสุํ โอหาเรตฺวา กาสายานิ วตฺถานิ อจฺฉาเทตฺวา}

\vspace{\tipitakaparskip}
\csromancenter{ฉพฺพิโสธนสุตฺต (112) อคารสฺมา อนคาริยํ ปพฺพเชยฺยนฺ”ติ.\csromanspacer{} โส โข อหํ อาวุโส อปเรน สมเยน อปฺปํ วา โภคกฺขนฺธํ ปหาย มหนฺตํ วา โภคกฺขนฺธํ ปหาย อปฺปํ วา ญาติปริวฏฺฏํ ปหาย มหนฺตํ วา ญาติปริวฏฺฏํ ปหาย เกสมสฺสุํ โอหาเรตฺวา กาสายานิ วตฺถานิ อจฺฉาเทตฺวา อคารสฺมา อนคาริยํ ปพฺพชิํ, โส เอวํ ปพฺพชิโต สมาโน ภิกฺขูนํ สิกฺขาสาชีวสมาปนฺโน ปาณาติปาตํ ปหาย ปาณาติปาตา ปฏิวิรโต อโหสิํ นิหิตทณฺโฑ นิหิตสตฺโถ ลชฺชี ทยาปนฺโน สพฺพปาณภูตหิตานุกมฺปี วิหาสิํ.\csromanspacer{} อทินฺนาทานํ ปหาย อทินฺนาทานา ปฏิวิรโต อโหสิํ ทินฺนาทายี ทินฺนปาฏิกงฺขี, อเถเนน สุจิภูเตน อตฺตนา วิหาสิํ.\csromanspacer{} อพฺรหฺมจริยํ ปหาย พฺรหฺมจารี อโหสิํ อาราจารี วิรโต เมถุนา คามธมฺมา.\csromanspacer{} มุสาวาทํ ปหาย มุสาวาทา ปฏิวิรโต อโหสิํ สจฺจวาที สจฺจสนฺโธ เถโต ปจฺจยิโก อวิสํวาทโก โลกสฺส.\csromanspacer{} ปิสุณํวาจํ ปหาย ปิสุณาย วาจาย ปฏิวิรโต อโหสิํ, อิโต สุตฺวา น อมุตฺร อกฺขาตา อิเมสํ เภทาย, อมุตฺร วา สุตฺวา น อิเมสํ อกฺขาตา อมูสํ เภทาย, อิติ ภินฺนานํ วา สนฺธาตา, สหิตานํ วา อนุปฺปทาตา, สมคฺคาราโม สมคฺครโต สมคฺคนนฺที สมคฺคกรณิํ วาจํ ภาสิตา อโหสิํ.\csromanspacer{} ผรุสํ วาจํ ปหาย ผรุสาย วาจาย ปฏิวิรโต อโหสิํ, ยา สา วาจา เนลา กณฺณสุขา เปมนียา หทยงฺคมา โปรี พหุชนกนฺตา พหุชนมนาปา, ตถารูปิํ วาจํ ภาสิตา อโหสิํ.\csromanspacer{} สมฺผปฺปลาปํ ปหาย สมฺปปฺปลาปา ปฏิวิรโต อโหสิํ กาลวาที ภูตวาที อตฺถวาที ธมฺมวาที วินยวาที, นิธานวติํ วาจํ ภาสิตา อโหสิํ กาเลน สาปเทสํ ปริยนฺตวติํ อตฺถสํหิตํ.}
\prose{\csromanfolio{83}โส พีชคามภูตคามสมารมฺภา ปฏิวิรโต อโหสิํ.\csromanspacer{} เอกภตฺติโก อโหสิํ รตฺตูปรโต วิรโต วิกาลโภชนา.\csromanspacer{} นจฺจคีตวาทิตวิสูกทสฺสนา ปฏิวิรโต อโหสิํ.\csromanspacer{} มาลาคนฺธวิเลปนธารณมณฺฑนวิภูสนฏฺฐานา ปฏิวิรโต อโหสิํ.\csromanspacer{} อุจฺจาสยนมหาสยนา ปฏิวิรโต อโหสิํ.\csromanspacer{} ชาตรูปรชตปฏิคฺคหณา ปฏิวิรโต อโหสิํ.\csromanspacer{} อามกธญฺญปฏิคฺคหณา ปฏิวิรโต อโหสิํ.\csromanspacer{} อามกมํสปฏิคฺคหณา ปฏิวิรโต อโหสิํ.\csromanspacer{} อิตฺถิกุมาริกปฏิคฺคหณา ปฏิวิรโต อโหสิํ.\csromanspacer{} ทาสิทาสปฏิคฺคหณา ปฏิวิรโต อโหสิํ.\csromanspacer{} อเชฬกปฏิคฺคหณา ปฏิวิรโต อโหสิํ.\csromanspacer{} กุกฺกุฏสูกรปฏิคฺคหณา ปฏิวิรโต \csromanfolio{84}อโหสิํ.\csromanspacer{} หตฺถิควสฺสวฬวปฏิคฺคหณา ปฏิวิรโต อโหสิํ.\csromanspacer{} เขตฺตวตฺถุปฏิคฺคหณา ปฏิวิรโต อโหสิํ.\csromanspacer{} ทูเตยฺยปหิณคมนานุโยคา ปฏิวิรโต อโหสิํ.\csromanspacer{} กยวิกฺกยา ปฏิวิรโต อโหสิํ.\csromanspacer{} ตุลากูฏกํสกูฏมานกูฏา ปฏิวิรโต อโหสิํ.\csromanspacer{} อุกฺโกฏนวญฺจนนิกติสาจิโยคา ปฏิวิรโต อโหสิํ, เฉทนวธพนฺธนวิปราโมส-อาโลปสหสาการา ปฏิวิรโต อโหสิํ.}

\prose{โส สนฺตุฏฺโฐ อโหสิํ กายปริหาริเกน จีวเรน กุจฺฉิปริหาริเกน ปิณฺฑปาเตน, โส เยน เยเนว\footnote{เยน เยน จ (ก)} ปกฺกมิํ, สมาทาเยว ปกฺกมิํ.\csromanspacer{} เสยฺยถาปิ นาม ปกฺขี สกุโณ เยน เยเนว เฑติ, สปตฺตภาโรว เฑติ.\csromanspacer{} เอวเมว โข อหํ อาวุโส สนฺตุฏฺโฐ อโหสิํ กายปริหาริเกน จีวเรน กุจฺฉิปริหาริเกน ปิณฺฑปาเตน, โส เยน เยเนว ปกฺกมิํ, สมาทาเยว ปกฺกมิํ.\csromanspacer{} โส อิมินา อริเยน สีลกฺขนฺเธน สมนฺนาคโต อชฺฌตฺตํ อนวชฺชสุขํ ปฏิสํเวเทสิํ.}

\proseitem{103}{โส จกฺขุนา รูปํ ทิสฺวา น นิมิตฺตคฺคาหี อโหสิํ นานุพฺยญฺชนคฺคาหี, ยตฺวาธิกรณเมนํ จกฺขุนฺทฺริยํ อสํวุตํ วิหรนฺตํ อภิชฺฌาโทมนสฺสา ปาปกา อกุสลา ธมฺมา อนฺวาสฺสเวยฺยุํ, ตสฺส สํวราย ปฏิปชฺชิํ, รกฺขิํ จกฺขุนฺทฺริยํ, จกฺขุนฺทฺริเย สํวรํ อาปชฺชิํ.\csromanspacer{} โสเตน สทฺทํ สุตฺวา ฯเปฯ.\csromanspacer{} ฆาเนน คนฺธํ ฆายิตฺวา ฯเปฯ.\csromanspacer{} ชิวฺหาย รสํ สายิตฺวา ฯเปฯ.\csromanspacer{} กาเยน โผฏฺฐพฺพํ ผุสิตฺวา ฯเปฯ.\csromanspacer{} มนสา ธมฺมํ วิญฺญาย น นิมิตฺตคฺคาหี อโหสิํ นานุพฺยญฺชนคฺคาหี, ยตฺวาธิกรณเมนํ มนินฺทฺริยํ อสํวุตํ วิหรนฺตํ อภิชฺฌาโทมนสฺสา ปาปกา อกุสลา ธมฺมา อนฺวาสฺสเวยฺยุํ, ตสฺส สํวราย ปฏิปชฺชิํ, รกฺขิํ มนินฺทฺริยํ, มนินฺทฺริเย สํวรํ อาปชฺชิํ.\csromanspacer{} โส อิมินา อริเยน อินฺทฺริยสํวเรน สมนฺนาคโต อชฺฌตฺตํ อพฺยาเสกสุขํ ปฏิสํเวเทสิํ.}

\prose{โส อภิกฺกนฺเต ปฏิกฺกนฺเต สมฺปชานการี อโหสิํ, อาโลกิเต วิโลกิเต สมฺปชานการี อโหสิํ, สมิญฺชิเต ปสาริเต สมฺปชานการี อโหสิํ, สํฆาฏิปตฺตจีวรธารเณ สมฺปชานการี อโหสิํ, อสิเต ปีเต ขายิเต สายิเต สมฺปชานการี อโหสิํ, อุจฺจารปสฺสาวกมฺเม สมฺปชานการี อโหสิํ, คเต ฐิเต นิสินฺเน สุตฺเต ชาคริเต ภาสิเต ตุณฺหีภาเว สมฺปชานการี อโหสิํ.}

\csromanheader{2}{2. ฉพฺพิโสธนสุตฺต (112)}
\prose{\csromanfolio{85}โส อิมินา จ อริเยน สีลกฺขนฺเธน สมนฺนาคโต (อิมาย จ อริยาย สนฺตุฏฺฐิยา สมนฺนาคโต)\footnote{ปสฺส ม 1 จูฬหตฺถิปโทปเม (239) ปิฏฺเฐ.} อิมินา จ อริเยน อินฺทฺริยสํวเรน สมนฺนาคโต อิมินา จ อริเยน สติสมฺปชญฺเญน สมนฺนาคโต วิวิตฺตํ เสนาสนํ ภชิํ อรญฺญํ รุกฺขมูลํ ปพฺพตํ กนฺทรํ คิริคุหํ สุสานํ วนปตฺถํ อพฺโภกาสํ ปลาลปุญฺชํ.\csromanspacer{} โส ปจฺฉาภตฺตํ ปิณฺฑปาตปฏิกฺกนฺโต นิสีทิํ ปลฺลงฺกํ อาภุชิตฺวา อุชุํ กายํ ปณิธาย ปริมุขํ สติํ อุปฏฺฐเปตฺวา.}

\prose{โส อภิชฺฌํ โลเก ปหาย วิคตาภิชฺเฌน เจตสา วิหาสิํ, อภิชฺฌาย จิตฺตํ ปริโสเธสิํ.\csromanspacer{} พฺยาปาทปโทสํ ปหาย อพฺยาปนฺนจิตฺโต วิหาสิํ สพฺพปาณภูตหิตานุกมฺปี, พฺยาปาทปโทสา จิตฺตํ ปริโสเธสิํ.\csromanspacer{} ถินมิทฺธํ ปหาย วิคตถินมิทฺโธ วิหาสิํ อาโลกสญฺญี สโต สมฺปชาโน, ถินมิทฺธา จิตฺตํ ปริโสเธสิํ.\csromanspacer{} อุทฺธจฺจกุกฺกุจฺจํ ปหาย อนุทฺธโต วิหาสิํ อชฺฌตฺตํ วูปสนฺตจิตฺโต, อุทฺธจฺจกุกฺกุจฺจา จิตฺตํ ปริโสเธสิํ.\csromanspacer{} วิจิกิจฺฉํ ปหาย ติณฺณวิจิกิจฺโฉ วิหาสิํ อกถํกถี กุสเลสุ ธมฺเมสุ, วิจิกิจฺฉาย จิตฺตํ ปริโสเธสิํ.}

\proseitem{104}{โส อิเม ปญฺจ นีวรเณ ปหาย เจตโส อุปกฺกิเลเส ปญฺญาย ทุพฺพลีกรเณ วิวิจฺเจว กาเมหิ วิวิจฺจ อกุสเลหิ ธมฺเมหิ สวิตกฺกํ สวิจารํ วิเวกชํ ปีติสุขํ ปฐมํ ฌานํ อุปสมฺปชฺช วิหาสิํ.\csromanspacer{} วิตกฺกวิจารานํ วูปสมา อชฺฌตฺตํ สมฺปสาทนํ เจตโส เอโกทิภาวํ อวิตกฺกํ อวิจารํ สมาธิชํ ปีติสุขํ ทุติยํ ฌานํ ฯเปฯ ตติยํ ฌานํ ฯเปฯ จตุตฺถํ ฌานํ อุปสมฺปชฺช วิหาสิํ.}

\prose{โส เอวํ สมาหิเต จิตฺเต ปริสุทฺเธ ปริโยทาเต อนงฺคเณ วิคตูปกฺกิเลเส มุทุภูเต กมฺมนิเย ฐิเต อาเนญฺชปฺปตฺเต อาสวานํ ขยญาณาย จิตฺตํ อภินินฺนาเมสิํ, โส “อิทํ ทุกฺขนฺ”ติ ยถาภูตํ อพฺภญฺญาสิํ, “อยํ ทุกฺขสมุทโย”ติ ยถาภูตํ อพฺภญฺญาสิํ, “อยํ ทุกฺขนิโรธา”ติ ยถาภูตํ อพฺภญฺญาสิํ, “อยํ ทุกฺขนิโรธคามินี ปฏิปทา”ติ ยถาภูตํ อพฺภญฺญาสิํ. “อิเม อาสวา”ติ ยถาภูตํ อพฺภญฺญาสิํ, “อยํ อาสวสมุทโย”ติ ยถาภูตํ อพฺภญฺญาสิํ, “อยํ อาสวนิโรโธ”ติ ยถาภูตํ อพฺภญฺญาสิํ, “อยํ อาสวนิโรธคามินี ปฏิปทา”ติ ยถาภูตํ อพฺภญฺญาสิํ, ตสฺส เม เอวํ ชานโต \csromanfolio{86}เอวํ ปสฺสโต กามาสวาปิ จิตฺตํ วิมุจฺจิตฺถ, ภวาสวาปิ จิตฺตํ วิมุจฺจิตฺถ, อวิชฺชาสวาปิ จิตฺตํ วิมุจฺจิตฺถ, วิมุตฺตสฺมิํ “วิมุตฺตมฺ”อิติ ญาณํ อโหสิ, “ขีณา ชาติ, วุสิตํ พฺรหฺมจริยํ, กตํ กรณียํ, นาปรํ อิตฺถตฺตายา”ติ อพฺภญฺญาสิํ.\csromanspacer{} เอวํ โข เม อาวุโส ชานโต เอวํ ปสฺสโต อิมสฺมิํ จ สวิญฺญาณเก กาเย พหิทฺธา จ สพฺพนิมิตฺเตสุ อหงฺการมมงฺการมานานุสยา สมูหตาติ.\csromanspacer{} ตสฺส ภิกฺขเว ภิกฺขุโน “สาธู”ติ ภาสิตํ อภินนฺทิตพฺพํ อนุโมทิตพฺพํ. “สาธู”ติ ภาสิตํ อภินนฺทิตฺวา อนุโมทิตฺวา เอวมสฺส วจนีโย “ลาภา โน อาวุโส, สุลทฺธํ โน อาวุโส, เย มยํ อายสฺมนฺตํ ตาทิสํ สพฺรหฺมจาริํ สมนุปสฺสามา”ติ\footnote{ปสฺสามาติ (สี)}.}

\prose{อิทมโวจ ภควา.\csromanspacer{} อตฺตมนา เต ภิกฺขู ภควโต ภาสิตํ อภินนฺทุนฺติ.}

\nitthitamruled{ฉพฺพิโสธนสุตฺตํ นิฏฺฐิตํ ทุติยํ.}
\csromanmatikaanchor{mtk.16}
\csromantocmarkref{subsection}{3. สปฺปุริสสุตฺต}{mtk.16}
\bookmark[dest={mtk.16},level=2]{3. สปฺปุริสสุตฺต}
\csromanheader{2}{3. สปฺปุริสสุตฺต}
\proseitem{105}{เอวํ เม สุตํ\csromandash{}เอกํ สมยํ ภควา สาวตฺถิยํ วิหรติ เชตวเน อนาถปิณฺฑิกสฺส อาราเม.\csromanspacer{} ตตฺร โข ภควา ภิกฺขู อามนฺเตสิ “ภิกฺขโว”ติ. “ภทนฺเต”ติ เต ภิกฺขู ภควโต ปจฺจสฺโสสุํ.\csromanspacer{} ภควา เอตทโวจ “สปฺปุริสธมฺมญฺจ โว ภิกฺขเว เทเสสฺสามิ อสปฺปุริสธมฺมญฺจ, ตํ สุณาถ สาธุกํ มนสิ กโรถ ภาสิสฺสามี”ติ. “เอวํ ภนฺเต”ติ โข เต ภิกฺขู ภควโต ปจฺจสฺโสสุํ, ภควา เอตทโวจ\csromandash{}}

\prose{กตโม จ ภิกฺขเว อสปฺปุริสธมฺโม, อิธ ภิกฺขเว อสปฺปุริโส อุจฺจากุลา ปพฺพชิโต โหติ, โส อิติ ปฏิสญฺจิกฺขติ “อหํ โขมฺหิ อุจฺจากุลา ปพฺพชิโต, อิเม ปนญฺเญ ภิกฺขู น อุจฺจากุลา ปพฺพชิตา”ติ.\csromanspacer{} โส ตาย อุจฺจากุลีนตาย อตฺตานุกฺกํเสติ ปรํ วมฺเภติ.\csromanspacer{} อยํ\footnote{อยมฺปิ (สี, อิ)} ภิกฺขเว อสปฺปุริสธมฺโม.\csromanspacer{} สปฺปุริโส จ โข ภิกฺขเว อิติ ปฏิสญฺจิกฺขติ “น โข อุจฺจากุลีนตาย โลภธมฺมา วา ปริกฺขยํ คจฺฉนฺติ, โทสธมฺมา วา ปริกฺขยํ}

\vspace{\tipitakaparskip}
\csromancenter{สปฺปุริสสุตฺต (113) คจฺฉนฺติ, มหธมฺมา วา ปริกฺขยํ คจฺฉนฺติ.\csromanspacer{} โน เจปิ อุจฺจากุลา ปพฺพชิโต โหติ, โส จ โหติ ธมฺมานุธมฺมปฺปฏิปนฺโน สามีจิปฺปฏิปนฺโน อนุธมฺมจารี, โส ตตฺถ ปุชฺโช, โส ตตฺถ ปาสํโส”ติ.\csromanspacer{} โส ปฏิปทํเยว อนฺตรํ กริตฺวา ตาย อุจฺจากุลีนตาย เนวตฺตานุกฺกํเสติ น ปรํ วมฺเภติ.\csromanspacer{} อยํ ภิกฺขเว สปฺปุริสธมฺโม. (1)}
\prose{\csromanfolio{87}ปุน จปรํ ภิกฺขเว อสปฺปุริโส มหากุลา ปพฺพชิโต โหติ ฯเปฯ\footnote{ยถา อุจฺจากุลวาเร, ตถา วิตฺถาเรตพฺพํ.} มหาโภคกุลา ปพฺพชิโต โหติ ฯเปฯ อุฬารโภคกุลา ปพฺพชิโต โหติ, โส อิติ ปฏิสญฺจิกฺขติ “อหํ โขมฺหิ อุฬารโภคกุลา ปพฺพชิโต, อิเม ปนญฺเญ ภิกฺขู น อุฬารโภคกุลา ปพฺพชิตา”ติ.\csromanspacer{} โส ตาย อุฬารโภคตาย อตฺตานุกฺกํเสติ ปรํ วมฺเภติ.\csromanspacer{} อยมฺปิ ภิกฺขเว อสปฺปุริสธมฺโม.\csromanspacer{} สปฺปุริโส จ โข ภิกฺขเว อิติ ปฏิสญฺจิกฺขติ “น โข อุฬารโภคตาย โลภธมฺมา วา ปริกฺขยํ คจฺฉนฺติ, โทสธมฺมา วา ปริกฺขยํ คจฺฉนฺติ, โมหธมฺมา วา ปริกฺขยํ คจฺฉนฺติ.\csromanspacer{} โน เจปิ อุฬารโภคกุลา ปพฺพชิโต โหติ, โส จ โหติ ธมฺมานุธมฺมปฺปฏิปนฺโน สามีจิปฺปฏิปนฺโน อนุธมฺมจารี, โส ตตฺถ ปุชฺโช, โส ตตฺถ ปาสํโส”ติ.\csromanspacer{} โส ปฏิปทํเยว อนฺตรํ กริตฺวา ตาย อุฬารโภคตาย เนวตฺตานุกฺกํเสติ น ปรํ วมฺเภติ.\csromanspacer{} อยมฺปิ ภิกฺขเว สปฺปุริสธมฺโม. (\footnote{ญาเตน (สี, ก), ญาตตฺเตน (สฺยา, กํ, อิ)}-4)}

\proseitem{106}{ปุน จปรํ ภิกฺขเว อสปฺปุริโส ญาโต โหติ ยสสฺสี, โส อิติ ปฏิสญฺจิกฺขติ “อหํ โขมฺหิ ญาโต ยสสฺสี, อิเม ปนญฺเญ ภิกฺขู อปฺปญฺญาตา อปฺเปสกฺขา”ติ.\csromanspacer{} โส เตน ญตฺเตน2 อตฺตานุกฺกํเสติ ปรํ วมฺเภติ.\csromanspacer{} อยมฺปิ ภิกฺขเว อสปฺปุริสธมฺโม.\csromanspacer{} สปฺปุริโส จ โข ภิกฺขเว อิติ ปฏิสญฺจิกฺขติ “น โข ญตฺเตน โลภธมฺมา วา ปริกฺขยํ คจฺฉนฺติ, โทสธมฺมา วา ปริกฺขยํ คจฺฉนฺติ, โมหธมฺมา วา ปริกฺขยํ คจฺฉนฺติ.\csromanspacer{} โน เจปิ ญาโต โหติ ยสสฺสี, โส จ โหติ ธมฺมานุธมฺมปฺปฏิปนฺโน สามีจิปฺปฏิปนฺโน อนุธมฺมจารี, โส ตตฺถ ปุชฺโช, โส ตตฺถ ปาสํโส”ติ.\csromanspacer{} โส ปฏิปทํเยว อนฺตรํ กริตฺวา เตน ญาตฺเตน เนวตฺตานุกฺกํเสติ น ปรํ วมฺเภติ.\csromanspacer{} อยมฺปิ ภิกฺขเว สปฺปุริสธมฺโม. (5)}

\prose{ปุน จปรํ ภิกฺขเว อสปฺปุริโส ลาภี โหติ จีวรปิณฺฑปาตเสนาสนคิลานปฺปจฺจยเภสชฺชปริกฺขารานํ, โส อิติ ปฏิสญฺจิกฺขติ “อหํ \csromanfolio{88}โขมฺหิ ลาภี อีวรปิณฺฑปาตเสนาสนคิลานปฺปจฺจยเภสชฺชปริกฺขารานํ, อิเม ปนญฺเญ ภิกฺขู น ลาภิโน จีวรปิณฺฑปาตเสนาสนคิลานปฺปจฺจยเภสชฺชปริกฺขารานนฺ”ติ.\csromanspacer{} โส เตน ลาเภน อตฺตานุกฺกํเสติ ปรํ วมฺเภติ.\csromanspacer{} อยมฺปิ ภิกฺขเว อสปฺปุริสธมฺโม.\csromanspacer{} สปฺปุริโส จ โข ภิกฺขเว อิติ ปฏิสญฺจิกฺขติ “น โข ลาเภน โลภธมฺมา วา ปริกฺขยํ คจฺฉนฺติ, โทสธมฺมา วา ปริกฺขยํ คจฺฉนฺติ, โมหธมฺมา วา ปริกฺขยํ คจฺฉนฺติ.\csromanspacer{} โน เจปิ ลาภี โหติ จีวรปิณฺฑปาตเสนาสนคิลานปฺปจฺจยเภสชฺชปริกฺขารานํ, โส จ โหติ ธมฺมานุธมฺมปฺปฏิปนฺโน สามีจิปฺปฏิปนฺโน อนุธมฺมจารี, โส ตตฺถ ปุชฺโช, โส ตตฺถ ปาสํโส”ติ.\csromanspacer{} โส ปฏิปทํเยว อนฺตรํ กริตฺวา เตน ลาเภน เนวตฺตานุกฺกํเสติ น ปรํ วมฺเภติ.\csromanspacer{} อยมฺปิ ภิกฺขเว สปฺปุริสธมฺโม. (6)}

\prose{ปุน จปรํ ภิกฺขเว อสปฺปุริโส พหุสฺสุโต โหติ, โส อิติ ปฏิสญฺจิกฺขติ “อหํ โขมฺหิ พหุสฺสุโต, อิเม ปนญฺเญ ภิกฺขู น พหุสฺสุตา”ติ.\csromanspacer{} โส เตน พาหุสจฺเจน อตฺตานุกฺกํเสติ ปรํ วมฺเภติ.\csromanspacer{} อยมฺปิ ภิกฺขเว อสปฺปุริสธมฺโม.\csromanspacer{} สปฺปุริโส จ โข ภิกฺขเว อิติ ปฏิสญฺจิกฺขติ “น โข พาหุสจฺเจน โลภธมฺมา วา ปริกฺขยํ คจฺฉนฺติ, โทสธมฺมา วา ปริกฺขยํ คจฺฉนฺติ, โมหธมฺมา วา ปริกฺขยํ คจฺฉนฺติ.\csromanspacer{} โน เจปิ พหุสฺสุโต โหติ, โส จ โหติ ธมฺมานุธมฺมปฺปฏิปนฺโน สามีจิปฺปฏิปนฺโน อนุธมฺมจารี, โส ตตฺถ ปุชฺโช, โส ตตฺถ ปาสํโส”ติ.\csromanspacer{} โส ปฏิปทํเยว อนฺตรํ กริตฺวา เตน พาหุสจฺเจน เนวตฺตานุกฺกํเสติ น ปรํ วมฺเภติ.\csromanspacer{} อยมฺปิ ภิกฺขเว สปฺปุริสธมฺโม. (7)}

\prose{ปุน จปรํ ภิกฺขเว อสปฺปุริโส วินยธโร โหติ, โส อิติ ปฏิสญฺจิกฺขติ “อหํ โขมฺหิ วินยธโร, อิเม ปนญฺเญ ภิกฺขู น วินยธรา”ติ.\csromanspacer{} โส เตน วินยธรตฺเตน อตฺตานุกฺกํเสติ ปรํ วมฺเภติ.\csromanspacer{} อยมฺปิ ภิกฺขเว อสปฺปุริสธมฺโม.\csromanspacer{} สปฺปุริโส จ โข ภิกฺขเว อิติ ปฏิสญฺจิกฺขติ “น โข วินยธรตฺเตน โลภธมฺมา วา ปริกฺขยํ คจฺฉนฺติ, โทสธมฺมา วา ปริกฺขยํ คจฺฉนฺติ, โมหธมฺมา วา ปริกฺขยํ คจฺฉนฺติ.\csromanspacer{} โน เจปิ วินยธโร โหติ, โส จ โหติ ธมฺมานุธมฺมปฺปฏิปนฺโน สามีจิปฺปฏิปนฺโน อนุธมฺมจารี, โส ตตฺถ ปุชฺโช, โส ตตฺถ ปาสํโส”ติ.\csromanspacer{} โส ปฏิปทํเยว อนฺตรํ กริตฺวา เตน วินยธรตฺเตน เนวตฺตานุกฺกํเสติ น ปรํ วมฺเภติ.\csromanspacer{} อยมฺปิ ภิกฺขเว สปฺปุริสธมฺโม. (8)}

\csromanheader{2}{3. สปฺปุริสสุตฺต (113)}
\prose{\csromanfolio{89}ปุน จปรํ ภิกฺขเว อสปฺปุริโส ธมฺมกถิโก โหติ, โส อิติ ปฏิสญฺจิกฺขติ “อหํ โขมฺหิ ธมฺมกถิโก, อิเม ปนญฺเญ ภิกฺขู น ธมฺมกถิกา”ติ.\csromanspacer{} โส เตน ธมฺมกถิกตฺเตน อตฺตานุกฺกํเสติ ปรํ วมฺเภติ.\csromanspacer{} อยมฺปิ ภิกฺขเว อสปฺปุริสธมฺโม.\csromanspacer{} สปฺปุริโส จ โข ภิกฺขเว อิติ ปฏิสญฺจิกฺขติ “น โข ธมฺมกถิกตฺเตน โลภธมฺมา วา ปริกฺขยํ คจฺฉนฺติ, โทสธมฺมา วา ปริกฺขยํ คจฺฉนฺติ, โมหธมฺมา วา ปริกฺขยํ คจฺฉนฺติ.\csromanspacer{} โน เจปิ ธมฺมกถิโก โหติ, โส จ โหติ ธมฺมานุธมฺมปฺปฏิปนฺโน สามีจิปฺปฏิปนฺโน อนุธมฺมจารี, โส ตตฺถ ปุชฺโช, โส ตตฺถ ปาสํโส”ติ.\csromanspacer{} โส ปฏิปทํเยว อนฺตรํ กริตฺวา เตน ธมฺมกถิกตฺเตน เนวตฺตานุกฺกํเสติ น ปรํ วมฺเภติ.\csromanspacer{} อยมฺปิ ภิกฺขเว สปฺปุริสธมฺโม. (9)}

\proseitem{107}{ปุน จปรํ ภิกฺขเว อสปฺปุริโส อารญฺญิโก โหติ, โส อิติ ปฏิสญฺจิกฺขติ “อหํ โขมฺหิ อารญฺญิโก, อิเม ปนญฺเญ ภิกฺขู น อารญฺญิกา”ติ.\csromanspacer{} โส เตน อารญฺญิกตฺเตน อตฺตานุกฺกํเสติ ปรํ วมฺเภติ.\csromanspacer{} อยมฺปิ ภิกฺขเว อสปฺปุริสธมฺโม.\csromanspacer{} สปฺปุริโส จ โข ภิกฺขเว อิติ ปฏิสญฺจิกฺขติ “น โข อารญฺญิกตฺเตน โลภธมฺมา วา ปริกฺขยํ คจฺฉนฺติ, โทสธมฺมา วา ปริกฺขยํ คจฺฉนฺติ, โมหธมฺมา วา ปริกฺขยํ คจฺฉนฺติ.\csromanspacer{} โน เจปิ อารญฺญิโก โหติ, โส จ โหติ ธมฺมานุธมฺมปฺปฏิปนฺโน สามีจิปฺปฏิปนฺโน อนุธมฺมจารี, โส ตตฺถ ปุชฺโช, โส ตตฺถ ปาสํโส”ติ.\csromanspacer{} โส ปฏิปทํเยว อนฺตรํ กริตฺวา เตน อารญฺญิกตฺเตน เนวตฺตานุกฺกํเสติ น ปรํ วมฺเภติ.\csromanspacer{} อยมฺปิ ภิกฺขเว สปฺปุริสธมฺโม. (10)}

\prose{ปุน จปรํ ภิกฺขเว อสปฺปุริโส ปํสุกูลิโก โหติ, โส อิติ ปฏิสญฺจิกฺขติ “อหํ โขมฺหิ ปํสุกูลิโก, อิเม ปนญฺเญ ภิกฺขู น ปํสุกูลิกา”ติ.\csromanspacer{} โส เตน ปํสุกูลิกตฺเตน อตฺตานุกฺกํเสติ ปรํ วมฺเภติ.\csromanspacer{} อยมฺปิ ภิกฺขเว อสปฺปุริสธมฺโม.\csromanspacer{} สปฺปุริโส จ โข ภิกฺขเว อิติ ปฏิสญฺจิกฺขติ “น โข ปํสุกูลิกตฺเตน โลภธมฺมา วา ปริกฺขยํ คจฺฉนฺติ, โทสธมฺมา วา ปริกฺขยํ คจฺฉนฺติ, โมหธมฺมา วา ปริกฺขยํ คจฺฉนฺติ.\csromanspacer{} โน เจปิ ปํสุกูลิโก โหติ.\csromanspacer{} โส จ โหติ ธมฺมานุธมฺมปฺปฏิปนฺโน สามีจิปฺปฏิปนฺโน อนุธมฺมจารี, โส ตตฺถ ปุชฺโช, โส ตตฺถ ปาสํโส”ติ.\csromanspacer{} โส ปฏิปทํเยว อนฺตรํ กริตฺวา เตน ปํสุกูลิกตฺเตน เนวตฺตานุกฺกํเสติ น ปรํ วมฺเภติ.\csromanspacer{} อยมฺปิ ภิกฺขเว สปฺปุริสธมฺโม. (11)}

\prose{\csromanfolio{90}ปุน จปรํ ภิกฺขเว อสปฺปุริโส ปิณฺฑปาติโก โหติ, โส อิติ ปฏิสญฺจิกฺขติ “อหํ โขมฺหิ ปิณฺฑปาติโก, อิเม ปนญฺเญ ภิกฺขู น ปิณฺฑปาติกา”ติ.\csromanspacer{} โส เตน ปิณฺฑปาติกตฺเตน อตฺตานุกฺกํเสติ ปรํ วมฺเภติ.\csromanspacer{} อยมฺปิ ภิกฺขเว อสปฺปุริสธมฺโม.\csromanspacer{} สปฺปุริโส จ โข ภิกฺขเว อิติ ปฏิสญฺจิกฺขติ “น โข ปิณฺฑปาติกตฺเตน โลภธมฺมา วา ปริกฺขยํ คจฺฉนฺติ, โทสธมฺมา วา ปริกฺขยํ คจฺฉนฺติ, โมหธมฺมา วา ปริกฺขยํ คจฺฉนฺติ.\csromanspacer{} โน เจปิ ปิณฺฑปาติโก โหติ.\csromanspacer{} โส จ โหติ ธมฺมานุธมฺมปฺปฏิปนฺโน สามีจิปฺปฏิปนฺโน อนุธมฺมจารี, โส ตตฺถ ปุชฺโช, โส ตตฺถ ปาสํโส”ติ.\csromanspacer{} โส ปฏิปทํเยว อนฺตรํ กริตฺวา เตน ปิณฺฑปาติกตฺเตน เนวตฺตานุกฺกํเสติ น ปรํ วมฺเภติ.\csromanspacer{} อยมฺปิ ภิกฺขเว สปฺปุริสธมฺโม. (12)}

\prose{ปุน จปรํ ภิกฺขเว อสปฺปุริโส รุกฺขมูลิโก โหติ, โส อิติ ปฏิสญฺจิกฺขติ “อหํ โขมฺหิ รุกฺขมูลิโก, อิเม ปนญฺเญ ภิกฺขู น รุกฺขมูลิกา”ติ.\csromanspacer{} โส เตน รุกฺขมูลิกตฺเตน อตฺตานุกฺกํเสติ ปรํ วมฺเภติ.\csromanspacer{} อยมฺปิ ภิกฺขเว อสปฺปุริสธมฺโม.\csromanspacer{} สปฺปุริโส จ โข ภิกฺขเว อิติ ปฏิสญฺจิกฺขติ “น โข รุกฺขมูลิกตฺเตน โลภธมฺมา วา ปริกฺขยํ คจฺฉนฺติ, โทสธมฺมา วา ปริกฺขยํ คจฺฉนฺติ, โมหธมฺมา วา ปริกฺขยํ คจฺฉนฺติ.\csromanspacer{} โน เจปิ รุกฺขมูลิโก โหติ, โส จ โหติ ธมฺมานุธมฺมปฺปฏิปนฺโน สามีจิปฺปฏิปนฺโน อนุธมฺมจารี, โส ตตฺถ ปุชฺโช, โส ตตฺถ ปาสํโส”ติ.\csromanspacer{} โส ปฏิปทํเยว อนฺตรํ กริตฺวา เตน รุกฺขมูลิกตฺเตน เน วตฺตานุกฺกํเสติ น ปรํ วมฺเภติ.\csromanspacer{} อยมฺปิ ภิกฺขเว สปฺปุริสธมฺโม. (13)}

\prose{ปุน จปรํ ภิกฺขเว อสปฺปุริโส โสสานิโก โหติ ฯเปฯ อพฺโภกาสิโก โหติ ฯเปฯ เนสชฺชิโก โหติ ฯเปฯ ยถาสนฺถติโก โหติ ฯเปฯ เอกาสนิโก โหติ, โส อิติ ปฏิสญฺจิกฺขติ “อหํ โขมฺหิ เอกาสนิโก, อิเม ปนญฺเญ ภิกฺขู น เอกาสนิกา”ติ.\csromanspacer{} โส เตน เอกาสนิกตฺเตน อตฺตานุกฺกํเสติ ปรํ วมฺเภติ.\csromanspacer{} อยมฺปิ ภิกฺขเว อสปฺปุริสธมฺโม.\csromanspacer{} สปฺปุริโส จ โข ภิกฺขเว อิติ ปฏิสญฺจิกฺขติ “น โข เอกาสนิกตฺเตน โลภธมฺมา วา ปริกฺขยํ คจฺฉนฺติ, โทสธมฺมา วา ปริกฺขยํ คจฺฉนฺติ, โมหธมฺมา วา ปริกฺขยํ คจฺฉนฺติ.\csromanspacer{} โน เจปิ เอกาสนิโก โหติ, โส จ โหติ ธมฺมานุธมฺมปฺปฏิปนฺโน สามีจิปฺปฏิปนฺโน อนุธมฺมจารี, โส ตตฺถ ปุชฺโช, โส ตตฺถ ปาสํโส”ติ.\csromanspacer{} โส ปฏิปทํเยว อนฺตรํ}

\vspace{\tipitakaparskip}
\csromancenter{สปฺปุริสสุตฺต (113) กริตฺวา เตน เอกาสนิกตฺเตน เนวตฺตานุกฺกํเสติ น ปรํ วมฺเภติ.\csromanspacer{} อยมฺปิ ภิกฺขเว สปฺปุริสธมฺโม. (14-18)}
\proseitem{108}{\csromanfolio{91}ปุน จปรํ ภิกฺขเว อสปฺปุริโส วิวิจฺเจว กาเมหิ วิวิจฺจ อกุสเลหิ ธมฺเมหิ สวิตกฺกํ สวิจารํ วิเวกชํ ปีติสุขํ ปฐมํ ฌานํ อุปสมฺปชฺช วิหรติ, โส อิติ ปฏิสญฺจิกฺขติ “อหํ โขมฺหิ ปฐมชฺฌานสมาปตฺติยา ลาภี, อิเม ปนญฺเญ ภิกฺขู ปฐมชฺฌานสมาปตฺติยา น ลาภิโน”ติ.\csromanspacer{} โส ตาย ปฐมชฺฌานสมาปตฺติยา อตฺตานุกฺกํเสติ ปรํ วมฺเภติ.\csromanspacer{} อยมฺปิ ภิกฺขเว อสปฺปุริสธมฺโม.\csromanspacer{} สปฺปุริโส จ โข ภิกฺขเว อิติ ปฏิสญฺจิกฺขติ “ปฐมชฺฌานสมาปตฺติยาปิ โข อตมฺมยตา วุตฺตา ภควตา, เยน เยน หิ มญฺญนฺติ, ตโต ตํ โหติ อญฺญถา”ติ.\csromanspacer{} โส อตมฺมยตญฺเญว อนฺตรํ กริตฺวา ตาย ปฐมชฺฌานสมาปตฺติยา เนวตฺตานุกฺกํเสติ น ปรํ วมฺเภติ.\csromanspacer{} อยมฺปิ ภิกฺขเว สปฺปุริสธมฺโม. (19)}

\prose{ปุน จปรํ ภิกฺขเว อสปฺปุริโส วิตกฺกวิจารานํ วูปสมา อชฺฌตฺตํ สมฺปสาทนํ เจตโส เอโกทิภาวํ อวิตกฺกํ อวิจารํ สมาธิชํ ปีติสุขํ ทุติยํ ฌานํ ฯเปฯ.\csromanspacer{} ตติยํ ฌานํ.\csromanspacer{} จตุตฺถํ ฌานํ อุปสมฺปชฺช วิหรติ, โส อิติ ปฏิสญฺจิกฺขติ “อหํ โขมฺหิ จตุตฺถชฺฌานสมาปตฺติยา ลาภี, อิเม ปนญฺเญ ภิกฺขู จตุตฺถชฺฌานสมาปตฺติยา น ลาภิโน”ติ.\csromanspacer{} โส ตาย จตุตฺถชฺฌานสมาปตฺติยา อตฺตานุกฺกํเสติ ปรํ วมฺเภติ.\csromanspacer{} อยมฺปิ ภิกฺขเว อสปฺปุริสธมฺโม.\csromanspacer{} สปฺปุริโส จ โข ภิกฺขเว อิติ ปฏิสญฺจิกฺขติ “จตุตฺถชฺฌานสมาปตฺติยาปิ โข อตมฺมยตา วุตฺตา ภควตา, เยน เยน หิ มญฺญนฺติ, ตโต ตํ โหติ อญฺญถา”ติ.\csromanspacer{} โส อตมฺมยตญฺเญว อนฺตรํ กริตฺวา ตาย จตุตฺถชฺฌานสมาปตฺติยา เนวตฺตานุกฺกํเสติ น ปรํ วมฺเภติ.\csromanspacer{} อยมฺปิ ภิกฺขเว สปฺปุริสธมฺโม. (20-22)}

\prose{ปุน จปรํ ภิกฺขเว อสปฺปุริโส สพฺพโส รูปสญฺญานํ สมติกฺกมา ปฏิฆสญฺญานํ อตฺถงฺคมา นานตฺตสญฺญานํ อมนสิการา “อนนฺโต อากาโส”ติ อากาสานญฺจายตนํ อุปสมฺปชฺช วิหรติ, โส อิติ ปฏิสญฺจิกฺขติ “อหํ โขมฺหิ อากาสานญฺจายตนสมาปตฺติยา ลาภี, อิเม ปนญฺเญ ภิกฺขู อากาสานญฺจายตนสมาปตฺติยา น ลาภิโน”ติ.\csromanspacer{} โส ตาย อากาสานญฺจายตนสมาปตฺติยา อตฺตานุกฺกํเสติ ปรํ วมฺเภติ.\csromanspacer{} อยมฺปิ ภิกฺขเว อสปฺปุริสธมฺโม.\csromanspacer{} สปฺปุริโส จ โข ภิกฺขเว อิติ \csromanfolio{92}ปฏิสญฺจิกฺขติ “อากาสานญฺจายตนสมาปตฺติยาปิ โข อตมฺมยตา วุตฺตา ภควตา, เยน เยน หิ มญฺญนฺติ, ตโต ตํ โหติ อญฺญถา”ติ.\csromanspacer{} โส อตมฺมยตญฺเญว อนฺตรํ กริตฺวา ตาย อากาสานญฺจายตนสมาปตฺติยา เนวตฺตานุกฺกํเสติ น ปรํ วมฺเภติ.\csromanspacer{} อยมฺปิ ภิกฺขเว สปฺปุริสธมฺโม. (23)}

\prose{ปุน จปรํ ภิกฺขเว อสปฺปุริโส สพฺพโส อากาสานญฺจายตนํ สมติกฺกมฺม “อนนฺตํ วิญฺญาณนฺ”ติ วิญฺญาณญฺจายตนํ อุปสมฺปชฺช วิหรติ, โส อิติ ปฏิสญฺจิกฺขติ “อหํ โขมฺหิ วิญฺญาณญฺจายตนสมาปตฺติยา ลาภี, อิเม ปนญฺเญ ภิกฺขู วิญฺญาณญฺจายตนสมาปตฺติยา น ลภิโน”ติ.\csromanspacer{} โส ตาย วิญฺญาณญฺจายตนสมาปตฺติยา อตฺตานุกฺกํเสติ ปรํ วมฺเภติ.\csromanspacer{} อยมฺปิ ภิกฺขเว อสปฺปุริสธมฺโม.\csromanspacer{} สปฺปุริโส จ โข ภิกฺขเว อิติ ปฏิสญฺจิกฺขติ “วิญฺญาณญฺจายตนสมาปตฺติยาปิ โข อตมฺมยตา วุตฺตา ภควตา, เยน เยน หิ มญฺญนฺติ, ตโต ตํ โหติ อญฺญถา”ติ.\csromanspacer{} โส อตมฺมยตญฺเญว อนฺตรํ กริตฺวา ตาย วิญฺญาณญฺจายตนสมาปตฺติยา เนวตฺตานุกฺกํเสติ น ปรํ วมฺเภติ.\csromanspacer{} อยมฺปิ ภิกฺขเว สปฺปุริสธมฺโม. (24)}

\prose{ปุน จปรํ ภิกฺขเว อสปฺปุริโส สพฺพโส วิญฺญาณญฺจายตนํ สมติกฺกมฺม “นตฺถิ กิญฺจี”ติ อากิญฺจญฺญายตนํ อุปสมฺปชฺช วิหรติ, โส อิติ ปฏิสญฺจิกฺขติ “อหํ โขมฺหิ อากิญฺจญฺญายตนสมาปตฺติยา ลาภี, อิเม ปนญฺเญ ภิกฺขู อากิญฺจญฺญายตนสมาปตฺติยา น ลาภิโน”ติ.\csromanspacer{} โส ตาย อากิญฺจญฺญายตนสมาปตฺติยา อตฺตานุกฺกํเสติ ปรํ วมฺเภติ.\csromanspacer{} อยมฺปิ ภิกฺขเว อสปฺปุริสธมฺโม.\csromanspacer{} สปฺปุริโส จ โข ภิกฺขเว อิติ ปฏิสญฺจิกฺขติ “อากิญฺจญฺญายตนสมาปตฺติยาปิ โข อตมฺมยตา วุตฺตา ภควตา, เยน เยน หิ มญฺญนฺติ, ตโต ตํ โหติ อญฺญาถา”ติ.\csromanspacer{} โส อตมฺมยตญฺเญว อนฺตรํ กริตฺวา ตาย อากิญฺจญฺญายตนสมาปตฺติยา เนวตฺตานุกฺกํเสติ น ปรํ วมฺเภติ.\csromanspacer{} อยมฺปิ ภิกฺขเว สปฺปุริสธมฺโม. (25)}

\prose{ปุน จปรํ ภิกฺขเว อสปฺปุริโส สพฺพโส อากิญฺจญฺญายตนํ สมติกฺกมฺม เนวสญฺญานาสญฺญายตนํ อุปสมฺปชฺช วิหรติ, โส อิติ ปฏิสญฺจิกฺขติ “อหํ โขมฺหิ เนวสญฺญานาสญฺญายตนสมาปตฺติยา ลาภี, อิเม ปนญฺเญ ภิกฺขู เนวสญฺญานาสญฺญายตนสมาปตฺติยา น ลาภิโน”ติ.\csromanspacer{} โส ตาย เนวสญฺญานาสญฺญายตนสมาปตฺติยา}

\vspace{\tipitakaparskip}
\csromancenter{เสวิตพฺพาเสวิตพฺพสุตฺต (114) อตฺตานุกฺกํเสติ ปรํ วมฺเภติ.\csromanspacer{} อยมฺปิ ภิกฺขเว อสปฺปุริสธมฺโม.\csromanspacer{} สปฺปุริโส จ โข ภิกฺขเว อิติ ปฏิสญฺจิกฺขติ “เนวสญฺญานาสญฺญายตน- สมาปตฺติยาปิ โข อตมฺมยตา วุตฺตา ภควตา, เยน เยน หิ มญฺญนฺติ, ตโต ตํ โหติ อญฺญถา”ติ.\csromanspacer{} โส อตมฺมยตญฺเญว อนฺตรํ กริตฺวา ตาย เนวสญฺญานาสญฺญายตนสมาปตฺติยา เนวตฺตานุกฺกํเสติ น ปรํ วมฺเภติ.\csromanspacer{} อยมฺปิ ภิกฺขเว สปฺปุริสธมฺโม. (26)}
\csromanmatikaanchor{mtk.17}
\csromantocmarkref{subsection}{4. เสวิตพฺพาเสวิตพฺพสุตฺต}{mtk.17}
\bookmark[dest={mtk.17},level=2]{4. เสวิตพฺพาเสวิตพฺพสุตฺต}
\prose{\csromanfolio{93}ปุน จปรํ ภิกฺขเว สปฺปุริโส สพฺพโส เนวสญฺญานาสญฺญายตนํ สมติกฺกมฺม สญฺญาเวทยิตนิโรธํ อุปสมฺปชฺช วิหรติ, ปญฺญาย จสฺส ทิสฺวา อาสวา\footnote{เอกจฺเจ อาสวา (ก)} ปริกฺขีณา โหนฺติ.\csromanspacer{} อยํ\footnote{อยํ โข (สฺยา, กํ)} ภิกฺขเว ภิกฺขุ น กิญฺจิ มญฺญติ น กุหิญฺจิ มญฺญติ น เกนจิ มญฺญตีติ. (27)}

\prose{อิทมโวจ ภควา.\csromanspacer{} อตฺตมนา เต ภิกฺขู ภควโต ภาสิตํ อภินนฺทุนฺติ.\csromanspacer{} สปฺปุริสสุตฺตํ นิฏฺฐิตํ ตติยํ.\textbf{4.}\csromanspacer{}\textbf{ เสวิตพฺพาเสวิตพฺพสุตฺต}}
\csromansectionrule

\proseitem{109}{เอวํ เม สุตํ\csromandash{}เอกํ สมยํ ภควา สาวตฺถิยํ วิหรติ เชตวเน อนาถปิณฺฑิกสฺส อาราเม.\csromanspacer{} ตตฺร โข ภควา ภิกฺขู อามนฺเตสิ “ภิกฺขโว”ติ. “ภทนฺเต”ติ เต ภิกฺขู ภควโต ปจฺจสฺโสสุํ.\csromanspacer{} ภควา เอตทโปจ “เสวิตพฺพาเสวิตพฺพํ โว ภิกฺขเว ธมฺมปริยายํ เทเสสฺสามิ, ตํ สุณาถ สาธุกํ มนสิ กโรถ ภาสิสฺสามี”ติ. “เอวํ ภนฺเต” โข เต ภิกฺขู ภควโต ปจฺจสฺโสสุํ.\csromanspacer{} ภควา เอตทโวจ\csromandash{}}

\prose{กายสมาจารํปาหํ\footnote{ปหํ (สพฺพตฺถ)} ภิกฺขเว ทุวิเธน วทามิ เสวิตพฺพมฺปิ อเสวิตพฺพมฺปิ, ตญฺจ อญฺญมญฺญํ กายสมาจารํ.\csromanspacer{} วจีสมาจารํปาหํ ภิกฺขเว ทุวิเธน วทามิ เสวิตพฺพมฺปิ อเสวิตพฺพมฺปิ, ตญฺจ อญฺญมญฺญํ วจีสมาจารํ.\csromanspacer{} มโนสมาจารํปาหํ ภิกฺขเว ทุวิเธน วทามิ เสวิตพฺพมฺปิ อเสวิตพฺพมฺปิ, ตญฺจ อญฺญมญฺญํ มโนสมาจารํ.\csromanspacer{} จิตฺตุปฺปาทํปาหํ ภิกฺขเว ทุวิเธน วทามิ เสวิตพฺพมฺปิ อเสวิตพฺพมฺปิ, ตญฺจ อญฺญมญฺญํ จิตฺตุปฺปาทํ.\csromanspacer{} สญฺญาปฏิลาภํปาหํ ภิกฺขเว ทุวิเธน วทามิ \csromanfolio{94}เสวิตพฺพมฺปิ อเสวิตพฺพมฺปิ, ตญฺจ อญฺญมญฺญํ สญฺญาปฏิลาภํ.\csromanspacer{} ทิฏฺฐิปฏิลาภํปาหํ ภิกฺขเว ทุวิเธน วทามิ เสวิตพฺพมฺปิ อเสวิตพฺพมฺปิ, ตญฺจ อญฺญมญฺญํ ทิฏฺฐิปฏิลาภํ.\csromanspacer{} อตฺตภาวปฏิลาภํปาหํ ภิกฺขเว ทุวิเธน วทามิ เสวิตพฺพมฺปิ อเสวิตพฺพมฺปิ, ตญฺจ อญฺญมญฺญํ อตฺตภาวปฏิลาภนฺติ.}

\prose{เอวํ วุตฺเต อายสฺมา สาริปุตฺโต ภควนฺตํ เอตทโวจ\csromandash{}อิมสฺส โข อหํ ภนฺเต ภควตา สํขิตฺเตน ภาสิตสฺส วิตฺถาเรน อตฺถํ อวิภตฺตสฺส เอวํ วิตฺถาเรน อตฺถํ อาชานามิ\csromandash{}}

\proseitem{110}{“กายสมาจารํปาหํ ภิกฺขเว ทุวิเธน วทามิ เสวิตพฺพมฺปิ อเสวิตพฺพมฺปิ, ตญฺจ อญฺญมญฺญํ กายสมาจารนฺ”ติ อิติ โข ปเนตํ วุตฺตํ ภควตา, กิญฺเจตํ ปฏิจฺจ วุตฺตํ.\csromanspacer{} ยถารูปํ ภนฺเต กายสมาจารํ เสวโต อกุสลา ธมฺมา อภิวฑฺฒนฺติ, กุสลา ธมฺมา ปริหายนฺติ, เอวรูโป กายสมาจาโร น เสวิตพฺโพ.\csromanspacer{} ยถารูปญฺจ โข ภนฺเต กายสมาจารํ เสวโต อกุสลา ธมฺมา ปริหายนฺติ, กุสลา ธมฺมา อภิวฑฺฒนฺติ, เอวรูโป กายสมาจาโร เสวิตพฺโพ.}

\proseitem{111}{กถํรูปํ ภนฺเต กายสมาจารํ เสวโต อกุสลา ธมฺมา อภิวฑฺฒนฺติ, กุสลา ธมฺมา ปริหายนฺติ.\csromanspacer{} อิธ ภนฺเต เอกจฺโจ ปาณาติปาตี โหติ ลุทฺโท โลหิตปาณิ หตปฺปหเต นิวิฏฺโฐ อทยาปนฺโน ปาณภูเตสุ.\csromanspacer{} อทินฺนาทายี โข ปน โหติ, ยํ ตํ ปรสฺส ปรวิตฺตูปกรณํ คามคตํ วา อรญฺญคตํ วา, ตํ อทินฺนํ เถยฺยสงฺขาตํ อาทาตา โหติ.\csromanspacer{} กาเมสุมิจฺฉาจารี โข ปน โหติ, ยา ตา มาตุรกฺขิตา ปิตุรกฺขิตา มาตาปิตุรกฺขิตา ภาตุรกฺขิตา ภคินิรกฺขิตา ญาติรกฺขิตา โคตฺตรกฺขิตา ธมฺมรกฺขิตา สสฺสามิกา สปริทณฺฑา อนฺตมโส มาลาคุฬปริกฺขิตฺตาปิ, ตถารูปาสุ จาริตฺตํ อาปชฺชิตา โหติ.\csromanspacer{} เอวรูปํ ภนฺเต กายสมาจารํ เสวโต อกุสลา ธมฺมา อภิวฑฺฒนฺติ, กุสลา ธมฺมา ปริหายนฺติ.}

\prose{กถํรูปํ ภนฺเต กายสมาจารํ เสวโต อกุสลา ธมฺมา ปริหายนฺติ, กุสลา ธมฺมา อภิวฑฺฒนฺติ.\csromanspacer{} อิธ ภนฺเต เอกจฺโจ ปาณาติปาตํ ปหาย ปาณาติปาตา ปฏิวิรโต โหติ, นิหิตทณฺโฑ นิหิตสตฺโถ ลชฺชี ทยาปนฺโน สพฺพปาณภูตหิตานุกมฺปี วิหรติ.\csromanspacer{} อทินฺนาทานํ ปหาย อทินฺนาทานา ปฏิวิรโต โหติ, ยํ ตํ ปรสฺส}

\vspace{\tipitakaparskip}
\csromancenter{เสวิตพฺพาเสวิตพฺพสุตฺต (114) ปรวิตฺตูปกรณํ คามคตํ วา อรญฺญคตํ วา, ตํ นาทินฺนํ เถยฺยสงฺขาตํ อาทาตา โหติ.\csromanspacer{} กาเมสุมิจฺฉาจารํ ปหาย กาเมสุมิจฺฉาจารา ปฏิวิรโต โหติ, ยา ตา มาตุรกฺขิตา ปิตุรกฺขิตา มาตาปิตุรกฺขิตา ภาตุรกฺขิตา ภคินิรกฺขิตา ญาติรกฺขิตา โคตฺตรกฺขิตา ธมฺมรกฺขิตา สสฺสามิกา สปริทณฺฑา อนฺตมโส มาลาคุฬปริกฺขิตฺตาปิ, ตถารูปาสุ น จาริตฺตํ อาปชฺชิตา โหติ.\csromanspacer{} เอวรูปํ ภนฺเต กายสมาจารํ เสวโต อกุสลา ธมฺมา ปริหายนฺติ, กุสลา ธมฺมา อภิวฑฺฒนฺติ. “กายสมาจารํปาหํ ภิกฺขเว ทุวิเธน วทามิ เสวิตพฺพมฺปิ อเสวิตพฺพมฺปิ, ตญฺจ อญฺญมญฺญํ กายสมาจารนฺ”ติ อิติ ยํ ตํ วุตฺตํ ภควตา, อิทเมตํ ปฏิจฺจ วุตฺตํ.}
\prose{\csromanfolio{95}“วจีสมาจารํปาหํ ภิกฺขเว ทุวิเธน วทามิ เสวิตพฺพมฺปิ อเสวิตพฺพมฺปิ, ตญฺจ อญฺญมญฺญํ วจีสมาจารนฺ”ติ อิติ โข ปเนตํ วุตฺตํ ภควตา, กิญฺเจตํ ปฏิจฺจ วุตฺตํ.\csromanspacer{} ยถารูปํ ภนฺเต วจีสมาจารํ เสวโต อกุสลา ธมฺมา อภิวฑฺฒนฺติ, กุสลา ธมฺมา ปริหายนฺติ, เอวรูโป วจีสมาจาโร น เสวิตพฺโพ.\csromanspacer{} ยถารูปญฺจ โข ภนฺเต วจีสมาจารํ เสวโต อกุสลา ธมฺมา ปริหายนฺติ, กุสลา ธมฺมา อภิวฑฺฒนฺติ, เอวรูโป วจีสมาจาโร เสวิตพฺโพ.}

\proseitem{112}{กถํรูปํ ภนฺเต วจีสมาจารํ เสวโต อกุสลา ธมฺมา อภิวฑฺฒนฺติ, กุสลา ธมฺมา ปริหายนฺติ.\csromanspacer{} อิธ ภนฺเต เอกจฺโจ มุสาวาที โหติ สภาคโต\footnote{สภคฺคโต (พหูสุ)} วา ปริสาคโต\footnote{ปริสคฺคโต (พหูสุ)} วา ญาติมชฺฌคโต วา ปูคมชฺฌคโต วา ราชกุลมชฺฌคโต วา อภินีโต สกฺขิปุฏฺโฐ “เอหมฺโภ ปุริส ยํ ชานาสิ, ตํ วเทหี”ติ โส อชานํ วา อาห “ชานามี”ติ, ชานํ วา อาห “น ชานามี”ติ, อปสฺสํ วา อาห “ปสฺสามี”ติ, ปสฺสํ วา อาห “น ปสฺสามี”ติ.\csromanspacer{} อิติ\footnote{ปสฺส ม 1 สาเลยฺยกสุตฺเต (356) ปิฏฺเฐ.} อตฺตเหตุ วา ปรเหตุ วา อามิสกิญฺจิกฺขเหตุ\footnote{กิญฺจกฺขเหตุ (สี)} วา สมฺปชานมุสา ภาสิตา โหติ.\csromanspacer{} ปิสุณวาโจ โข ปน โหติ, อิโต สุตฺวา อมุตฺร อกฺขาตา อิเมสํ เภทาย, อมุตฺร วา สุตฺวา อิเมสํ อกฺขาตา อมูสํ เภทาย, อิติ สมคฺคานํ วา เภตฺตา ภินฺนานํ วา อนุปฺปทาตา วคฺคาราโม วคฺครโต วคฺคนนฺที วคฺคกรณิํ วาจํ \csromanfolio{96}ภาสิตา โหติ.\csromanspacer{} ผรุสวาโจ โข ปน โหติ, ยา สา วาจา กณฺฑกา กกฺกสา ผรุสา ปรกฏุกา ปราภิสชฺชนี โกธสามนฺตา อสมาธิสํวตฺตนิกา, ตถารูปิํ วาจํ ภาสิตา โหติ.\csromanspacer{} สมฺผปฺปลาปี โข ปน โหติ, อกาลวาที อภูตวาที อนตฺถวาที อธมฺมวาที อวินยวาที, อนิธานวติํ วาจํ ภาสิตา โหติ อกาเลน อนปเทสํ อปริยนฺตวติํ อนตฺถสํหิตํ.\csromanspacer{} เอวรูปํ ภนฺเต วจีสมาจารํ เสวโต อกุสลา ธมฺมา อภิวฑฺฒนฺติ, กุสลา ธมฺมา ปริหายนฺติ.}

\prose{กถํรูปํ ภนฺเต วจีสมาจารํ เสวโต อกุสลา ธมฺมา ปริหายนฺติ, กุสลา ธมฺมา อภิวฑฺฒนฺติ.\csromanspacer{} อิธ ภนฺเต เอกจฺโจ มุสาวาทํ ปหาย มุสาวาทา ปฏิวิรโต โหติ สภาคโต วา ปริสาคโต วา ญาติมชฺฌคโต วา ปูคมชฺฌคโต วา ราชกุลมชฺฌคโต วา อภินีโต สกฺขิปุฏฺโฐ “เอหมฺโภ ปุริส ยํ ชานาสิ, ตํ วเทหี”ติ โส อชานํ วา อาห “น ชานามี”ติ, ชานํ วา อาห “ชานามี”ติ, อปสฺสํ วา อาห “น ปสฺสามี”ติ, ปสฺสํ วา อาห ‘ปสฺสามี’ติ.\csromanspacer{} อิติ อตฺตเหตุ วา ปรเหตุ วา อามิสกิญฺจิกฺขเหตุ วา น สมฺปชานมุสา ภาสิตา โหติ.\csromanspacer{} ปิสุณํ วาจํ ปหาย ปิสุณาย วาจาย ปฏิวิรโต โหติ, อิโต สุตฺวา น อมุตฺร อกฺขาตา อิเมสํ เภทาย, อมุตฺร วา สุตฺวา น อิเมสํ อกฺขาตา อมูสํ เภทาย, อิติ ภินฺนานํ วา สนฺธาตา สหิตานํ วา อนุปฺปทาตา สมคฺคาราโม สมคฺครโต สมคฺคนนฺที สมคฺคกรณิํ วาจํ ภาสิตา โหติ.\csromanspacer{} ผรุสํ วาจํ ปหาย ผรุสาย วาจาย ปฏิวิรโต โหติ, ยา สา วาจา เนลา กณฺณสุขา เปมนียา หทยงฺคมา โปรี พหุชนกนฺตา พหุชนมนาปา, ตถารูปิํ วาจํ ภาสิตา โหติ.\csromanspacer{} สมฺผปฺปลาปํ ปหาย สมฺผปฺปลาปา ปฏิวิรโต โหติ, กาลวาที ภูตวาที อตฺถวาที ธมฺมวาที วินยวาที, นิธานวติํ วาจํ ภาสิตา โหติ กาเลน สาปเทสํ ปริยนฺตวติํ อตฺถสํหิตํ.\csromanspacer{} เอวรูปํ ภนฺเต วจีสมาจารํ เสวโต อกุสลา ธมฺมา ปริหายนฺติ, กุสลา ธมฺมา อภิวฑฺฒนฺติ. “วจีสมาจารํปาหํ ภิกฺขเว ทุวิเธน วทามิ เสวิตพฺพมฺปิ อเสวิตพฺพมฺปิ, ตญฺจ อญฺญมญฺญํ วจีสมาจารนฺ”ติ อิติ ยํ ตํ วุตฺตํ ภควตา, อิทเมตํ ปฏิจฺจ วุตฺตํ.}

\prose{“มโนสมาจารํปาหํ ภิกฺขเว ทุวิเธน วทามิ เสวิตพฺพมฺปิ อเสวิตพฺพมฺปิ, ตญฺจ อญฺญมญฺญํ มโนสมาจารนฺ”ติ อิติ โข ปเนตํ วุตฺตํ ภควตา, กิญฺเจตํ ปฏิจฺจ วุตฺตํ.\csromanspacer{} ยถารูปํ ภนฺเต มโนสมาจารํ เสวโต}

\vspace{\tipitakaparskip}
\csromancenter{เสวิตพฺพาเสวิตพฺพสุตฺต (114) อกุสลา ธมฺมา อติวฑฺฒนฺติ, กุสลา ธมฺมา ปริหายนฺติ, เอวรูโป มโนสมาจาโร น เสวิตพฺโพ.\csromanspacer{} ยถารูปญฺจ โข ภนฺเต มโนสมาจารํ เสวโต อกุสลา ธมฺมา ปริหายนฺติ, กุสลา ธมฺมา อภิวฑฺฒนฺติ, เอวรูโป มโนสมาจาโร เสวิตพฺโพ.}
\proseitem{113}{\csromanfolio{97}กถํรูปํ ภนฺเต มโนสมาจารํ เสวโต อกุสลา ธมฺมา อภิวฑฺฒนฺติ, กุสลา ธมฺมา ปริหายนฺติ.\csromanspacer{} อิธ ภนฺเต เอกจฺโจ อภิชฺฌาลุ โหติ.\csromanspacer{} ยํ ตํ ปรสฺส ปรวิตฺตูปกรณํ, ตํ อภิชฺฌาตา โหติ “อโห วต ยํ ปรสฺส, ตํ มมสฺสา”ติ.\csromanspacer{} พฺยาปนฺนจิตฺโต โข ปน โหติ ปทุฏฺฐมนสงฺกปฺโป “อิเม สตฺตา หญฺญนฺตุ วา วชฺฌนฺตุ วา อุจฺฉิชฺชนฺตุ วา วินสฺสนฺตุ วา มา วา อเหสุนฺ”ติ.\csromanspacer{} เอวรูปํ ภนฺเต มโนสมาจารํ เสวโต อกุสลา ธมฺมา อภิวฑฺฒนฺติ, กุสลา ธมฺมา ปริหายนฺติ.}

\prose{กถํรูปํ ภนฺเต มโนสมาจารํ เสวโต อกุสลา ธมฺมา ปริหายนฺติ, กุสลา ธมฺมา อภิวฑฺฒนฺติ.\csromanspacer{} อิธ ภนฺเต เอกจฺโจ อนภิชฺฌาลุ โหติ, ยํ ตํ ปรสฺส ปรวิตฺตูปกรณํ, ตํ นาภิชฺฌาตา โหติ “อโหวต ยํ ปรสฺส, ตํ มมสฺสา”ติ.\csromanspacer{} อพฺยาปนฺนจิตฺโต โข ปน โหติ อปฺปทุฏฺฐมนสงฺกปฺโป “อิเม สตฺตา อเวรา อพฺยาพชฺฌา\footnote{อพฺยาปชฺฌา (สี, สฺยา, กํ, อิ, ก)} อนีฆา สุขี อตฺตานํ ปริหรนฺตู”ติ.\csromanspacer{} เอวรูปํ ภนฺเต มโนสมาจารํ เสวโต อกุสลา ธมฺมา ปริหายนฺติ, กุสลา ธมฺมา อภิวฑฺฒนฺติ. “มโนสมาจารํปาหํ ภิกฺขเว ทุวิเธน วทามิ เสวิตพฺพมฺปิ อเสวิตพฺพมฺปิ, ตญฺจ อญฺญมญฺญํ มโนสมาจารนฺ”ติ อิติ ยํ ตํ วุตฺตํ ภควตา, อิทเมตํ ปฏิจฺจ วุตฺตํ.}

\proseitem{114}{“จิตฺตุปฺปาทํปาหํ ภิกฺขเว ทุวิเธน วทามิ เสวิตพฺพมฺปิ อเสวิตพฺพมฺปิ ตญฺจ อญฺญมญฺญํ จิตฺตุปฺปาทนฺ”ติ อิติ โข ปเนตํ วุตฺตํ ภควตา, กิญฺเจตํ ปฏิจฺจ วุตฺตํ.\csromanspacer{} ยถารูปํ ภนฺเต จิตฺตุปฺปาทํ เสวโต อกุสลา ธมฺมา อภิวฑฺฒนฺติ, กุสลา ธมฺมา ปริหายนฺติ, เอวรูโป จิตฺตุปฺปาโท น เสวิตพฺโพ.\csromanspacer{} ยถารูปญฺจ โข ภนฺเต จิตฺตุปฺปาทํ เสวโต อกุสลา ธมฺมา ปริหายนฺติ, กุสลา ธมฺมา อภิวฑฺฒนฺติ, เอวรูโป จิตฺตุปฺปาโท เสวิตพฺโพ.}

\prose{กถํรูปํ ภนฺเต จิตฺตุปฺปาทํ เสวโต อกุสลา ธมฺมา อภิวฑฺฒนฺติ, กุสลา ธมฺมา ปริหายนฺติ.\csromanspacer{} อิธ ภนฺเต เอกจฺโจ อภิชฺฌาลุ โหติ, \csromanfolio{98}อภิชฺฌาสหคเตน เจตสา วิหรติ, พฺยาปาทวา โหติ, พฺยาปาทสหคเตน เจตสา วิหรติ, วิเหสวา โหติ, วิเหสาสหคเตน เจตสา วิหรติ.\csromanspacer{} เอวรูปํ ภนฺเต จิตฺตุปฺปาทํ เสวโต อกุสลา ธมฺมา อภิวฑฺฒนฺติ, กุสลา ธมฺมา ปริหายนฺติ.}

\prose{กถํรูปํ ภนฺเต จิตฺตุปฺปาทํ เสวโต อกุสลา ธมฺมา ปริหายนฺติ, กุสลา ธมฺมา อภิวฑฺฒนฺติ.\csromanspacer{} อิธ ภนฺเต เอกจฺโจ อนภิชฺฌาลุ โหติ, อนภิชฺฌาสหคเตน เจตสา วิหรติ, อพฺยาปาทวา โหติ, อพฺยาปาทสหคเตน เจตสา วิหรติ, อวิเหสวา โหติ, อวิเหสาสหคเตน เจตสา วิหรติ.\csromanspacer{} เอวรูปํ ภนฺเต จิตฺตุปฺปาทํ เสวโต อกุสลา ธมฺมา ปริหายนฺติ, กุสลา ธมฺมา อภิวฑฺฒนฺติ. “จิตฺตุปฺปาทํปาหํ ภิกฺขเว ทุวิเธน วทามิ เสวิตพฺพมฺปิ อเสวิตพฺพมฺปิ, ตญฺจ อญฺญมญฺญํ จิตฺตุปฺปาทนฺ”ติ อิติ ยํ ตํ วุตฺตํ ภควตา, อิทเมตํ ปฏิจฺจ วุตฺตํ.}

\proseitem{115}{“สญฺญาปฏิลาภํปาหํ ภิกฺขเว ทุวิเธน วทามิ เสวิตพฺพมฺปิ อเสวิตพฺพมฺปิ, ตญฺจ อญฺญมญฺญํ สญฺญาปฏิลาภนฺ”ติ อิติ โข ปเนตํ วุตฺตํ ภควตา, กิญฺเจตํ ปฏิจฺจ วุตฺตํ.\csromanspacer{} ยถารูปํ ภนฺเต สญฺญาปฏิลาภํ เสวโต อกุสลา ธมฺมา อภิวฑฺฒนฺติ, กุสลา ธมฺมา ปริหายนฺติ, เอวรูโป สญฺญาปฏิลาโภ น เสวิตพฺโพ.\csromanspacer{} ยถารูปญฺจ โข ภนฺเต สญฺญาปฏิลาภํ เสวโต อกุสลา ธมฺมา ปริหายนฺติ, กุสลา ธมฺมา อภิวฑฺฒนฺติ, เอวรูโป สญฺญาปฏิลาโภ เสวิตพฺโพ.}

\prose{กถํรูปํ ภนฺเต สญฺญาปฏิลาภํ เสวโต อกุสลา ธมฺมา อภิวฑฺฒนฺติ, กุสลา ธมฺมา ปริหายนฺติ.\csromanspacer{} อิธ ภนฺเต เอกจฺโจ อภิชฺฌาลุ โหติ, อภิชฺฌาสหคตาย สญฺญาย วิหรติ, พฺยาปาทวา โหติ, พฺยาปาทสหคตาย สญฺญาย วิหรติ, วิเหสวา โหติ, วิเหสาสหคตาย สญฺญาย วิหรติ.\csromanspacer{} เอวรูปํ ภนฺเต สญฺญาปฏิลาภํ เสวโต อกุสลา ธมฺมา อภิวฑฺฒนฺติ, กุสลา ธมฺมา ปริหายนฺติ.}

\prose{กถํรูปํ ภนฺเต สญฺญาปฏิลาภํ เสวโต อกุสลา ธมฺมา ปริหายนฺติ, กุสลา ธมฺมา อภิวฑฺฒนฺติ.\csromanspacer{} อิธ ภนฺเต เอกจฺโจ อนภิชฺฌาลุ โหติ, อนภิชฺฌาสหคตาย สญฺญาย วิหรติ, อพฺยาปาทวา โหติ, อพฺยาปาทสหคตาย สญฺญาย วิหรติ, อวิเหสวา โหติ, อวิเหสาสหคตาย สญฺญาย วิหรติ.\csromanspacer{} เอวรูปํ ภนฺเต สญฺญาปฏิลาภํ}

\vspace{\tipitakaparskip}
\csromancenter{เสวิตพฺพาเสวิตพฺพสุตฺต (114) เสวโต อกุสลา ธมฺมา ปริหายนฺติ, กุสลา ธมฺมา อภิวฑฺฒนฺติ. “สญฺญาปฏิลาภํปาหํ ภิกฺขเว ทุวิเธน วทามิ เสวิตพฺพมฺปิ อเสวิตพฺพมฺปิ, ตญฺจ อญฺญมญฺญํ สญฺญาปฏิลาภนฺ”ติ อิติ ยํ ตํ วุตฺตํ ภควตา, อิทเมตํ ปฏิจฺจ วุตฺตํ.}
\proseitem{116}{\csromanfolio{99}“ทิฏฺฐิปฏิลาภํปาหํ ภิกฺขเว ทุวิเธน วทามิ เสวิตพฺพมฺปิ อเสวิตพฺพมฺปิ, ตญฺจ อญฺญมญฺญํ ทิฏฺฐิปฏิลาภนฺ”ติ อิติ โข ปเนตํ วุตฺตํ ภควตา, กิญฺเจตํ ปฏิจฺจ วุตฺตํ.\csromanspacer{} ยถารูปํ ภนฺเต ทิฏฺฐิปฏิลาภํ เสวโต อกุสลา ธมฺมา อภิวฑฺฒนฺติ, กุสลา ธมฺมา ปริหายนฺติ, เอวรูโป ทิฏฺฐิปฏิลาโภ น เสวิตพฺโพ.\csromanspacer{} ยถารูปญฺจ โข ภนฺเต ทิฏฺฐิปฏิลาภํ เสวโต อกุสลา ธมฺมา ปริหายนฺติ, กุสลา ธมฺมา อภิวฑฺฒนฺติ, เอวรูโป ทิฏฺฐิปฏิลาโภ เสวิตพฺโพ.}

\prose{กถํรูปํ ภนฺเต ทิฏฺฐิปฏิลาภํ เสวโต อกุสลา ธมฺมา อภิวฑฺฒนฺติ, กุสลา ธมฺมา ปริหายนฺติ.\csromanspacer{} อิธ ภนฺเต เอกจฺโจ เอวํทิฏฺฐิโก โหติ “นตฺถิ ทินฺนํ, นตฺถิ ยิฏฺฐํ, นตฺถิ หุตํ, นตฺถิ สุกตทุกฺกฏานํ กมฺมานํ ผลํ วิปาโก, นตฺถิ อยํ โลโก, นตฺถิ ปโร โลโก, นตฺถิ มาตา, นตฺถิ ปิตา, นตฺถิ สตฺตา โอปปาติกา, นตฺถิ โลเก สมณพฺราหฺมณา สมฺมคฺคตา สมฺมาปฏิปนฺนา, เย อิมญฺจ โลกํ ปรญฺจ โลกํ สยํ อภิญฺญา สจฺฉิกตฺวา ปเวเทนฺตี”ติ.\csromanspacer{} เอวรูปํ ภนฺเต ทิฏฺฐิปฏิลาภํ เสวโต อกุสลา ธมฺมา อภิวฑฺฒนฺติ, กุสลา ธมฺมา ปริหายนฺติ.}

\prose{กถํรูปํ ภนฺเต ทิฏฺฐิปฏิลาภํ เสวโต อกุสลา ธมฺมา ปริหายนฺติ, กุสลา ธมฺมา อภิวฑฺฒนฺติ.\csromanspacer{} อิธ ภนฺเต เอกจฺโจ เอวํทิฏฺฐิโก โหติ “อตฺถิ ทินฺนํ, อตฺถิ ยิฏฺฐํ, อตฺถิ หุตํ, อตฺถิ สุกตทุกฺกฏานํ กมฺมานํ ผลํ วิปาโก, อตฺถิ อยํ โลโก, อตฺถิ ปโร โลโก, อตฺถิ มาตา, อตฺถิ ปิตา, อตฺถิ สตฺตา โอปปาติกา, อตฺถิ โลเก สมณพฺราหฺมณา สมฺมคฺคตา สมฺมาปฏิปนฺนา, เย อิมญฺจ โลกํ ปรญฺจ โลกํ สยํ อภิญฺญา สจฺฉิกตฺวา ปเวเทนฺตี”ติ.\csromanspacer{} เอวรูปํ ภนฺเต ทิฏฺฐิปฏิลาภํ เสวโต อกุสลา ธมฺมา ปริหายนฺติ, กุสลา ธมฺมา อภิวฑฺฒนฺติ. “ทิฏฺฐิปฏิลาภํปาหํ ภิกฺขเว ทุวิเธน วทามิ เสวิตพฺพมฺปิ อเสวิตพฺพมฺปิ, ตญฺจ อญฺญมญฺญํ ทิฏฺฐิปฏิลาภนฺ”ติ อิติ ยํ ตํ วุตฺตํ ภควตา, อิทเมตํ ปฏิจฺจ วุตฺตํ.}

\proseitem{117}{\csromanfolio{100}“อตฺตภาวปฏิลาภํปาหํ ภิกฺขเว ทุวิเธน วทามิ เสวิตพฺพมฺปิ อเสวิตพฺพมฺปิ, ตญฺจ อญฺญมญฺญํ อตฺตภาวปฏิลาภนฺ”ติ อิติ โข ปเนตํ วุตฺตํ ภควตา, กิญฺเจตํ ปฏิจฺจ วุตฺตํ.\csromanspacer{} ยถารูปํ ภนฺเต อตฺตภาวปฏิลาภํ เสวโต อกุสลา ธมฺมา อภิวฑฺฒนฺติ, กุสลา ธมฺมา ปริหายนฺติ, เอวรูโป อตฺตภาวปฏิลาโภ น เสวิตพฺโพ.\csromanspacer{} ยถารูปญฺจ โข ภนฺเต อตฺตภาวปฏิลาภํ เสวโต อกุสลา ธมฺมา ปริหายนฺติ, กุสลา ธมฺมา อภิวฑฺฒนฺติ, เอวรูโป อตฺตภาวปฏิลาโภ เสวิตพฺโพ.}

\prose{กถํรูปํ ภนฺเต อตฺตภาวปฏิลาภํ เสวโต อกุสลา ธมฺมา อภิวฑฺฒนฺติ, กุสลา ธมฺมา ปริหายนฺติ.\csromanspacer{} สพฺยาพชฺฌํ\footnote{สพฺยาปชฺฌํ (สี, สฺยา, กํ, อิ, ก)} ภนฺเต อตฺตภาวปฏิลาภํ อภินิพฺพตฺตยโต อปรินิฏฺฐิตภาวาย อกุสลา ธมฺมา อภิวฑฺฒนฺติ, กุสลา ธมฺมา ปริหายนฺติ.\csromanspacer{} อพฺยาพชฺฌํ ภนฺเต อตฺตภาวปฏิลาภํ อภินิพฺพตฺตยโต ปรินิฏฺฐิตภาวาย อกุสลา ธมฺมา ปริหายนฺติ, กุสลา ธมฺมา อภิวฑฺฒนฺติ. “อตฺตภาวปฏิลาภํปาหํ ภิกฺขเว ทุวิเธน วทามิ เสวิตพฺพมฺปิ อเสวิตพฺพมฺปิ, ตญฺจ อญฺญมญฺญํ อตฺตภาวปฏิลาภนฺ”ติ อิติ ยํ ตํ วุตฺตํ ภควตา, อิทเมตํ ปฏิจฺจ วุตฺตํ.}

\prose{อิมสฺส โข อหํ ภนฺเต ภควตา สํขิตฺเตน ภาสิตสฺส วิตฺถาเรน อตฺถํ อวิภตฺตสฺส เอวํ วิตฺถาเรน อตฺถํ อาชานามีติ.}

\proseitem{118}{สาธุ สาธุ สาริปุตฺต, สาธุ โข ตฺวํ สาริปุตฺต อิมสฺส มยา สํขิตฺเตน ภาสิตสฺส วิตฺถาเรน อตฺถํ อวิภตฺตสฺส เอวํ วิตฺถาเรน อตฺถํ อาชานาสิ\csromandash{}}

\prose{“กายสมาจารํปาหํ ภิกฺขเว ทุวิเธน วทามิ เสวิตพฺพมฺปิ อเสวิตพฺพมฺปิ, ตญฺจ อญฺญมญฺญํ กายสมาจารนฺ”ติ อิติ โข ปเนตํ วุตฺตํ มยา, กิญฺเจตํ ปฏิจฺจ วุตฺตํ.\csromanspacer{} ยถารูปํ สาริปุตฺต กายสมาจารํ เสวโต อกุสลา ธมฺมา อภิวฑฺฒนฺติ, กุสลา ธมฺมา ปริหายนฺติ, เอวรูโป กายสมาจาโร น เสวิตพฺโพ.\csromanspacer{} ยถารูปญฺจ โข สาริปุตฺต กายสมาจารํ เสวโต อกุสลา ธมฺมา ปริหายนฺติ, กุสลา ธมฺมา อภิวฑฺฒนฺติ, เอวรูโป กายสมาจาโร เสวิตพฺโพ.}

\prose{กถํรูปํ สาริปุตฺต กายสมาจารํ เสวโต อกุสลา ธมฺมา อภิวฑฺฒนฺติ, กุสลา ธมฺมา ปริหายนฺติ.\csromanspacer{} อิธ สาริปุตฺต เอกจฺโจ ปาณาติปาตี}

\vspace{\tipitakaparskip}
\csromancenter{เสวิตพฺพาเสวิตพฺพสุตฺต (114) โหติ, ลุทฺโท โลหิตปาณิ หตปฺปหเต นิวิฏฺโฐ อทยาปนฺโน ปาณภูเตสุ.\csromanspacer{} อทินฺนาทายี โข ปน โหติ, ยํ ตํ ปรสฺส ปรวิตฺตูปกรณํ คามคตํ วา อรญฺญคตํ วา, ตํ อทินฺนํ เถยฺยสงฺขาตํ อาทาตา โหติ.\csromanspacer{} กาเมสุมิจฺฉาจารี โข ปน โหติ, ยา ตา มาตุรกฺขิตา ปิตุรกฺขิตา มาตาปิตุรกฺขิตา ภาตุรกฺขิตา ภคินิรกฺขิตา ญาติรกฺขิตา โคตฺตรกฺขิตา ธมฺมรกฺขิตา สสฺสามิกา สปริทณฺฑา อนฺตมโส มาลาคุฬปริกฺขิตฺตาปิ, ตถารูปาสุ จาริตฺตํ อาปชฺชิตา โหติ.\csromanspacer{} เอวรูปํ สาริปุตฺต กายสมาจารํ เสวโต อกุสลา ธมฺมา อภิวฑฺฒนฺติ, กุสลา ธมฺมา ปริหายนฺติ.}
\prose{\csromanfolio{101}กถํรูปํ สาริปุตฺต กายสมาจารํ เสวโต อกุสลา ธมฺมา ปริหายนฺติ, กุสลา ธมฺมา อภิวฑฺฒนฺติ.\csromanspacer{} อิธ สาริปุตฺต เอกจฺโจ ปาณาติปาตํ ปหาย ปาณาติปาตา ปฏิวิรโต โหติ นิหิตทณฺโฑ นิหิตสตฺโถ ลชฺชี ทยาปนฺโน สพฺพปาณภูตหิตานุกมฺปี วิหรติ.\csromanspacer{} อทินฺนาทานํ ปหาย อทินฺนาทานา ปฏิวิรโต โหติ, ยํ ตํ ปรสฺส ปรวิตฺตูปกรณํ คามคตํ วา อรญฺญคตํ วา, ตํ นาทินฺนํ เถยฺยสงฺขาตํ อาทาตา โหติ.\csromanspacer{} กาเมสุมิจฺฉาจารํ ปหาย กาเมสุมิจฺฉาจารา ปฏิวิรโต โหติ, ยา ตา มาตุรกฺขิตา ปิตุรกฺขิตา มาตาปิตุรกฺขิตา ภาตุรกฺขิตา ภคินิรกฺขิตา ญาติรกฺขิตา โคตฺตรกฺขิตา ธมฺมรกฺขิตา สสฺสามิกา สปริทณฺฑา อนฺตมโส มาลาคุฬปริกฺขิตฺตาปิ, ตถารูปาสุ น จาริตฺตํ อาปชฺชิตา โหติ.\csromanspacer{} เอวรูปํ สาริปุตฺต กายสมาจารํ เสวโต อกุสลา ธมฺมา ปริหายนฺติ, กุสลา ธมฺมา อภิวฑฺฒนฺติ. “กายสมาจารํปาหํ ภิกฺขเว ทุวิเธน วทามิ เสวิตพฺพมฺปิ อเสวิตพฺพมฺปิ, ตญฺจ อญฺญมญฺญํ กายสมาจารนฺ”ติ อิติ ยํ ตํ วุตฺตํ มยา, อิทเมตํ ปฏิจฺจ วุตฺตํ.}

\prose{วจีสมาจารํปาหํ ภิกฺขเว ทุวิเธน วทามิ ฯเปฯ.\csromanspacer{} มโนสมาจารํปาหํ ภิกฺขเว ทุวิเธน วทามิ ฯเปฯ.\csromanspacer{} จิตฺตุปฺปาทํปาหํ ภิกฺขเว ทุวิเธน วทามิ ฯเปฯ.\csromanspacer{} สญฺญาปฏิลาภํปาหํ ภิกฺขเว ทุวิเธน วทามิ ฯเปฯ.\csromanspacer{} ทิฏฺฐิปฏิลาภํปาหํ ภิกฺขเว ทุวิเธน วทามิ ฯเปฯ.}

\prose{“อตฺตภาวปฏิลาภํปาหํ ภิกฺขเว ทุวิเธน วทามิ เสวิตพฺพมฺปิ อเสวิตพฺพมฺปิ, ตญฺจ อญฺญมญฺญํ อตฺตภาวปฏิลาภนฺ”ติ อิติ โข ปเนตํ วุตฺตํ มยา, \csromanfolio{102}กิญฺเจตํ ปฏิจฺจ วุตฺตํ.\csromanspacer{} ยถารูปํ สาริปุตฺต อตฺตภาวปฏิลาภํ เสวโต อกุสลา ธมฺมา อภิวฑฺฒนฺติ, กุสลา ธมฺมา ปริหายนฺติ, เอวรูโป อตฺตภาวปฏิลาโภ น เสวิตพฺโพ.\csromanspacer{} ยถารูปญฺจ โข สาริปุตฺต อตฺตภาวปฏิลาภํ เสวโต อกุสลา ธมฺมา ปริหายนฺติ, กุสลา ธมฺมา อภิวฑฺฒนฺติ, เอวรูโป อตฺตภาวปฏิลาโภ เสวิตพฺโพ.}

\prose{กถํรูปํ สาริปุตฺต อตฺตภาวปฏิลาภํ เสวโต อกุสลา ธมฺมา อภิวฑฺฒนฺติ, กุสลา ธมฺมา ปริหายนฺติ.\csromanspacer{} สพฺยาพชฺฌํ สาริปุตฺต อตฺตภาวปฏิลาภํ อภินิพฺพตฺตยโต อปรินิฏฺฐิตภาวาย อกุสลา ธมฺมา อภิวฑฺฒนฺติ, กุสลา ธมฺมา ปริหายนฺติ.\csromanspacer{} อพฺยาพชฺฌํ สาริปุตฺต อตฺตภาวปฏิลาภํ อภินิพฺพตฺตยโต ปรินิฏฺฐิตภาวาย อกุสลา ธมฺมา ปริหายนฺติ, กุสลา ธมฺมา อภิวฑฺฒนฺติ. “อตฺตภาวปฏิลาภํปาหํ ภิกฺขเว ทุวิเธน วทามิ เสวิตพฺพมฺปิ อเสวิตพฺพมฺปิ, ตญฺจ อญฺญมญฺญํ อตฺตภาวปฏิลาภนฺ”ติ อิติ ยํ ตํ วุตฺตํ มยา, อิทเมตํ ปฏิจฺจ วุตฺตํ.\csromanspacer{} อิมสฺส โข สาริปุตฺต มยา สํขิตฺเตน ภาสิตสฺส เอวํ วิตฺถาเรน อตฺโถ ทฏฺฐพฺโพ.}

\proseitem{119}{จกฺขุวิญฺเญยฺยํ รูปํปาหํ สาริปุตฺต ทุวิเธน วทามิ เสวิตพฺพมฺปิ อเสวิตพฺพมฺปิ.\csromanspacer{} โสตวิญฺเญยฺยํ สทฺทํปาหํ สาริปุตฺต ทุวิเธน วทามิ เสวิตพฺพมฺปิ อเสวิตพฺพมฺปิ.\csromanspacer{} ฆานวิญฺเญยฺยํ คนฺธํปาหํ สาริปุตฺต ทุวิเธน วทามิ เสวิตพฺพมฺปิ อเสวิตพฺพมฺปิ.\csromanspacer{} ชิวฺหาวิญฺเญยฺยํ รสํปาหํ สาริปุตฺต ทุวิเธน วทามิ เสวิตพฺพมฺปิ อเสวิตพฺพมฺปิ.\csromanspacer{} กายวิญฺเญยฺยํ โผฏฺฐพฺพํปาหํ สาริปุตฺต ทุวิเธน วทามิ เสวิตพฺพมฺปิ อเสวิตพฺพมฺปิ.\csromanspacer{} มโนวิญฺเญยฺยํ ธมฺมํปาหํ สาริปุตฺต ทุวิเธน วทามิ เสวิตพฺพมฺปิ อเสวิตพฺพมฺปีติ.}

\prose{เอวํ วุตฺเต อายสฺมา สาริปุตฺโต ภควนฺตํ เอตทโวจ\csromandash{}อิมสฺส โข อหํ ภนฺเต ภควตา สํขิตฺเตน ภาสิตสฺส วิตฺถาเรน อตฺถํ อวิภตฺตสฺส เอวํ วิตฺถาเรน อตฺถํ อาชานามิ\csromandash{}“จกฺขุวิญฺเญยฺยํ รูปํปาหํ สาริปุตฺต ทุวิเธน วทามิ เสวิตพฺพมฺปิ อเสวิตพฺพมฺปี”ติ อิติ โข ปเนตํ วุตฺตํ ภควตา, กิญฺเจตํ ปฏิจฺจ วุตฺตํ.\csromanspacer{} ยถารูปํ ภนฺเต จกฺขุวิญฺเญยฺยํ รูปํ เสวโต อกุสลา ธมฺมา อภิวฑฺฒนฺติ, กุสลา ธมฺมา ปริหายนฺติ, เอวรูปํ จกฺขุวิญฺเญยฺยํ รูปํ น เสวิตพฺพํ.\csromanspacer{} ยถารูปญฺจ โข ภนฺเต จกฺขุวิญฺเญยฺยํ รูปํ เสวโต อกุสลา ธมฺมา ปริหายนฺติ, กุสลา ธมฺมา อภิวฑฺฒนฺติ, เอวรูปํ จกฺขุวิญฺเญยฺยํ รูปํ เสวิตพฺพํ.}

\vspace{\tipitakaparskip}
\csromancenter{เสวิตพฺพาเสวิตพฺพสุตฺต (114) “จกฺขุวิญฺเญยฺยํ รูปํปาหํ สาริปุตฺต ทุวิเธน วทามิ เสวิตพฺพมฺปิ อเสวิตพฺพมฺปี”ติ อิติ ยํ ตํ วุตฺตํ ภควตา, อิทเมตํ ปฏิจฺจ วุตฺตํ.}
\prose{\csromanfolio{103}โสตวิญฺเญยฺยํ สทฺทํปาหํ สาริปุตฺต ฯเปฯ เอวรูโป โสตวิญฺเญยฺโย สทฺโท น เสวิตพฺโพ.\csromanspacer{} เอวรูโป โสตวิญฺเญยฺโย สทฺโท เสวิตพฺโพ.\csromanspacer{} เอวรูโป ฆานวิญฺเญยฺโย คนฺโธ น เสวิตพฺโพ.\csromanspacer{} เอวรูโป ฆานวิญฺเญยฺโย คนฺโธ เสวิตพฺโพ.\csromanspacer{} เอวรูโป ชิวฺหาวิญฺเญยฺโย รโส น เสวิตพฺโพ.\csromanspacer{} เอวรูโป ชิวฺหาวิญฺเญยฺโย รโส เสวิตพฺโพ.\csromanspacer{} กายวิญฺเญยฺยํ โผฏฺฐพฺพํปาหํ สาริปุตฺต ฯเปฯ เอวรูโป กายวิญฺเญยฺโย โผฏฺฐพฺโพ น เสวิตพฺโพ.\csromanspacer{} เอวรูโป กายวิญฺเญยฺโย โผฏฺฐพฺโพ เสวิตพฺโพ.}

\prose{“มโนวิญฺเญยฺยํ ธมฺมํปาหํ สาริปุตฺต ทุวิเธน วทามิ เสวิตพฺพมฺปิ อเสวิตพฺพมฺปี”ติ อิติ โข ปเนตํ วุตฺตํ ภควตา, กิญฺเจตํ ปฏิจฺจ วุตฺตํ.\csromanspacer{} ยถารูปํ ภนฺเต มโนวิญฺเญยฺยํ ธมฺมํ เสวโต อกุสลา ธมฺมา อภิวฑฺฒนฺติ, กุสลา ธมฺมา ปริหายนฺติ, เอวรูโป มโนวิญฺเญยฺโย ธมฺโม น เสวิตพฺโพ.\csromanspacer{} ยถารูปญฺจ โข ภนฺเต มโนวิญฺเญยฺยํ ธมฺมํ เสวโต อกุสลา ธมฺมา ปริหายนฺติ, กุสลา ธมฺมา อภิวฑฺฒนฺติ, เอวรูโป มโนวิญฺเญยฺโย ธมฺโม เสวิตพฺโพ. “มโนวิญฺเญยฺยํ ธมฺมํปาหํ สาริปุตฺต ทุวิเธน วทามิ เสวิตพฺพมฺปิ อเสวิตพฺพมฺปี”ติ อิติ ยํ ตํ วุตฺตํ ภควตา, อิทเมตํ ปฏิจฺจ วุตฺตํ.\csromanspacer{} อิมสฺส โข อหํ ภนฺเต ภควตา สํขิตฺเตน ภาสิตสฺส วิตฺถาเรน อตฺถํ อวิภตฺตสฺส เอวํ วิตฺถาเรน อตฺถํ อาชานามีติ.}

\proseitem{120}{สาธุ สาธุ สาริปุตฺต, สาธุ โข ตฺวํ สาริปุตฺต อิมสฺส มยา สํขิตฺเตน ภาสิตสฺส วิตฺถาเรน อตฺถํ อวิภตฺตสฺส เอวํ วิตฺถาเรน อตฺถํ อาชานาสิ\csromandash{}“จกฺขุวิญฺเญยฺยํ รูปํปาหํ สาริปุตฺต ทุวิเธน วทามิ เสวิตพฺพมฺปิ อเสวิตพฺพมฺปี”ติ อิติ โข ปเนตํ วุตฺตํ มยา, กิญฺเจตํ ปฏิจฺจ วุตฺตํ.\csromanspacer{} ยถารูปํ สาริปุตฺต จกฺขุวิญฺเญยฺยํ รูปํ เสวโต อกุสลา ธมฺมา อภิวฑฺฒนฺติ, กุสลา ธมฺมา ปริหายนฺติ, เอวรูปํ จกฺขุวิญฺเญยฺยํ รูปํ น เสวิตพฺพํ.\csromanspacer{} ยถารูปญฺจ โข สาริปุตฺต จกฺขุวิญฺเญยฺยํ รูปํ เสวโต อกุสลา ธมฺมา ปริหายนฺติ, กุสลา ธมฺมา อภิวฑฺฒนฺติ, เอวรูปํ จกฺขุวิญฺเญยฺยํ รูปํ เสวิตพฺพํ. “จกฺขุวิญฺเญยฺยํ รูปํปาหํ สาริปุตฺต ทุวิเธน วทามิ เสวิตพฺพมฺปิ อเสวิตพฺพมฺปี”ติ อิติ ยํ ตํ วุตฺตํ มยา, อิทเมตํ ปฏิจฺจ วุตฺตํ.}

\prose{\csromanfolio{104}โสตวิญฺเญยฺยํ สทฺทํปาหํ สาริปุตฺต ฯเปฯ เอวรูโป โสตวิญฺเญยฺโย สทฺโท น เสวิตพฺโพ.\csromanspacer{} เอวรูโป โสตวิญฺเญยฺโย สทฺโท เสวิตพฺโพ.\csromanspacer{} เอวรูโป ฆานวิญฺเญยฺโย คนฺโธ น เสวิตพฺโพ.\csromanspacer{} เอวรูโป ฆานวิญฺเญยฺโย คนฺโธ เสวิตพฺโพ.\csromanspacer{} เอวรูโป ชิวฺหาวิญฺเญยฺโย รโส น เสวิตพฺโพ.\csromanspacer{} เอวรูโป ชิวฺหาวิญฺเญยฺโย รโส เสวิตพฺโพ.\csromanspacer{} เอวรูโป กายวิญฺเญยฺโย โผฏฺฐพฺโพ น เสวิตพฺโพ.\csromanspacer{} เอวรูโป กายวิญฺเญยฺโย โผฏฺฐพฺโพ เสวิตพฺโพ.}

\prose{มโนวิญฺเญยฺยํ ธมฺมํปาหํ สาริปุตฺต ฯเปฯ เอวรูโป มโนวิญฺเญยฺโย ธมฺโม น เสวิตพฺโพ.\csromanspacer{} เอวรูโป มโนวิญฺเญยฺโย ธมฺโม เสวิตพฺโพ. “มโนวิญฺเญยฺยํ ธมฺมํปาหํ สาริปุตฺต ทุวิเธน วทามิ เสวิตพฺพมฺปิ อเสวิตพฺพมฺปี”ติ อิติ ยํ ตํ วุตฺตํ มยา, อิทเมตํ ปฏิจฺจ วุตฺตํ.\csromanspacer{} อิมสฺส โข สาริปุตฺต มยา สํขิตฺเตน ภาสิตสฺส เอวํ วิตฺถาเรน อตฺโถ ทฏฺฐพฺโพ.}

\proseitem{121}{จีวรํปาหํ สาริปุตฺต ทุวิเธน วทามิ เสวิตพฺพมฺปิ อเสวิตพฺพมฺปิ ฯเปฯ.\csromanspacer{} ปิณฺฑปาตํปาหํ สาริปุตฺต.\csromanspacer{} เสนาสนํปาหํ สาริปุตฺต.\csromanspacer{} คามํปาหํ สาริปุตฺต.\csromanspacer{} นิคมํปาหํ สาริปุตฺต.\csromanspacer{} นครํปาหํ สาริปุตฺต.\csromanspacer{} ชนปทํปาหํ สาริปุตฺต.\csromanspacer{} ปุคฺคลํปาหํ สาริปุตฺต ทุวิเธน วทามิ เสวิตพฺพมฺปิ อเสวิตพฺพมฺปีติ.}

\prose{เอวํ วุตฺเต อายสฺมา สาริปุตฺโต ภควนฺตํ เอตทโวจ\csromandash{}อิมสฺส โข อหํ ภนฺเต ภควตา สํขิตฺเตน ภาสิตสฺส วิตฺถาเรน อตฺถํ อวิภตฺตสฺส เอวํ วิตฺถาเรน อตฺถํ อาชานามิ\csromandash{}“จีวรํปาหํ สาริปุตฺต ทุวิเธน วทามิ เสวิตพฺพมฺปิ อเสวิตพฺพมฺปี”ติ อิติ โข ปเนตํ วุตฺตํ ภควตา, กิญฺเจตํ ปฏิจฺจ วุตฺตํ.\csromanspacer{} ยถารูปํ ภนฺเต จีวรํ เสวโต อกุสลา ธมฺมา อภิวฑฺฒนฺติ, กุสลา ธมฺมา ปริหายนฺติ, เอวรูปํ จีวรํ น เสวิตพฺพํ.\csromanspacer{} ยถารูปญฺจ โข ภนฺเต จีวรํ เสวโต อกุสลา ธมฺมา ปริหายนฺติ, กุสลา ธมฺมา อภิวฑฺฒนฺติ, เอวรูปํ จีวรํ เสวิตพฺพํ. “จีวรํปาหํ สาริปุตฺต ทุวิเธน วทามิ เสวิตพฺพมฺปิ อเสวิตพฺพมฺปี”ติ อิติ ยํ ตํ วุตฺตํ ภควตา, อิทเมตํ ปฏิจฺจ วุตฺตํ.}

\prose{ปิณฺฑปาตํปาหํ สาริปุตฺต ฯเปฯ เอวรูโป ปิณฺฑปาโต น เสวิตพฺโพ.\csromanspacer{} เอวรูโป ปิณฺฑปาโต เสวิตพฺโพ.\csromanspacer{} เสนาสนํปาหํ สาริปุตฺต ฯเปฯ เอวรูปํ เสนาสนํ น เสวิตพฺพํ.\csromanspacer{} เอวรูปํ เสนาสนํ เสวิตพฺพํ.\csromanspacer{} คามํปาหํ}

\vspace{\tipitakaparskip}
\csromancenter{เสวิตพฺพาเสวิตพฺพสุตฺต (114) สาริปุตฺต ฯเปฯ เอวรูโป คาโม น เสวิตพฺโพ.\csromanspacer{} เอวรูโป คาโม เสวิตพฺโพ.\csromanspacer{} เอวรูโป นิคโม น เสวิตพฺโพ.\csromanspacer{} เอวรูโป นิคโม เสวิตพฺโพ.\csromanspacer{} เอวรูปํ นครํ น เสวิตพฺพํ.\csromanspacer{} เอวรูปํ นครํ เสวิตพฺพํ.\csromanspacer{} เอวรูโป ชนปโท น เสวิตพฺโพ.\csromanspacer{} เอวรูโป ชนปโท เสวิตพฺโพ.}
\prose{\csromanfolio{105}“ปุคฺคลํปาหํ สาริปุตฺต ทุวิเธน วทามิ เสวิตพฺพมฺปิ อเสวิตพฺพมฺปี”ติ อิติ โข ปเนตํ วุตฺตํ ภควตา, กิญฺเจตํ ปฏิจฺจ วุตฺตํ.\csromanspacer{} ยถารูปํ ภนฺเต ปุคฺคลํ เสวโต อกุสลา ธมฺมา อภิวฑฺฒนฺติ, กุสลา ธมฺมา ปริหายนฺติ, เอวรูโป ปุคฺคโล น เสวิตพฺโพ.\csromanspacer{} ยถารูปญฺจ โข ภนฺเต ปุคฺคลํ เสวโต อกุสลา ธมฺมา ปริหายนฺติ, กุสลา ธมฺมา อภิวฑฺฒนฺติ, เอวรูโป ปุคฺคโล เสวิตพฺโพ. “ปุคฺคลํปาหํ สาริปุตฺต ทุวิเธน วทามิ เสวิตพฺพมฺปิ อเสวิตพฺพมฺปี”ติ อิติ ยํ ตํ วุตฺตํ ภควตา, อิทเมตํ ปฏิจฺจ วุตฺตนฺติ.\csromanspacer{} อิมสฺส โข อหํ ภนฺเต ภควตา สํขิตฺเตน ภาสิตสฺส วิตฺถาเรน อตฺถํ อวิภตฺตสฺส เอวํ วิตฺถาเรน อตฺถํ อาชานามีติ.}

\proseitem{122}{สาธุ สาธุ สาริปุตฺต, สาธุ โข ตฺวํ สาริปุตฺต อิมสฺส มยา สํขิตฺเตน ภาสิตสฺส วิตฺถาเรน อตฺถํ อวิภตฺตสฺส เอวํ วิตฺถาเรน อตฺถํ อาชานาสิ\csromandash{}“จีวรํปาหํ สาริปุตฺต ทุวิเธน วทามิ เสวิตพฺพมฺปิ อเสวิตพฺพมฺปี”ติ อิติ โข ปเนตํ วุตฺตํ มยา, กิญฺเจตํ ปฏิจฺจ วุตฺตํ.\csromanspacer{} ยถารูปํ สาริปุตฺต จีวรํ เสวโต อกุสลา ธมฺมา อภิวฑฺฒนฺติ, กุสลา ธมฺมา ปริหายนฺติ.\csromanspacer{} เอวรูปํ จีวรํ น เสวิตพฺพํ.\csromanspacer{} ยถารูปญฺจ โข สาริปุตฺต จีวรํ เสวโต อกุสลา ธมฺมา ปริหายนฺติ, กุสลา ธมฺมา อภิวฑฺฒนฺติ, เอวรูปํ จีวรํ เสวิตพฺพํ. “จีวรํปาหํ สาริปุตฺต ทุวิเธน วทามิ เสวิตพฺพมฺปิ อเสวิตพฺพมฺปี”ติ อิติ ยํ ตํ วุตฺตํ มยา, อิทเมตํ ปฏิจฺจ วุตฺตํ. (ยถา ปฐมํ, ตถา วิตฺถาเรตพฺพํ.) เอวรูโป ปิณฺฑปาโต.\csromanspacer{} เอวรูปํ เสนาสนํ.\csromanspacer{} เอวรูโป คาโม.\csromanspacer{} เอวรูโป นิคโม.\csromanspacer{} เอวรูปํ นครํ.\csromanspacer{} เอวรูโป ชนปโท.}

\prose{“ปุคฺคลํปาหํ สาริปุตฺต ทุวิเธน วทามิ เสวิตพฺพมฺปิ อเสวิตพฺพมฺปี”ติ อิติ โข ปเนตํ วุตฺตํ มยา, กิญฺเจตํ ปฏิจฺจ วุตฺตํ.\csromanspacer{} ยถารูปํ สาริปุตฺต ปุคฺคลํ เสวโต อกุสลา ธมฺมา อภิวฑฺฒนฺติ, กุสลา ธมฺมา ปริหายนฺติ, เอวรูโป ปุคฺคโล น เสวิตพฺโพ.\csromanspacer{} ยถารูปญฺจ โข สาริปุตฺต ปุคฺคลํ \csromanfolio{106}เสวโต อกุสลา ธมฺมา ปริหายนฺติ, กุสลา ธมฺมา อภิวฑฺฒนฺติ, เอวรูโป ปุคฺคโล เสวิตพฺโพ. “ปุคฺคลํปาหํ สาริปุตฺต ทุวิเธน วทามิ เสวิตพฺพมฺปิ อเสวิตพฺพมฺปี”ติ อิติ ยํ ตํ วุตฺตํ มยา, อิทเมตํ ปฏิจฺจ วุตฺตํ.\csromanspacer{} อิมสฺส โข สาริปุตฺต มยา สํขิตฺเตน ภาสิตสฺส เอวํ วิตฺถาเรน อตฺโถ ทฏฺฐพฺโพ.}

\proseitem{123}{สพฺเพปิ เจ สาริปุตฺต ขตฺติยา อิมสฺส มยา สํขิตฺเตน ภาสิตสฺส เอวํ วิตฺถาเรน อตฺถํ อาชาเนยฺยุํ สพฺเพสานมฺปิสฺส ขตฺติยานํ ทีฆรตฺตํ หิตาย สุขาย.\csromanspacer{} สพฺเพปิ เจ สาริปุตฺต พฺราหฺมณา ฯเปฯ.\csromanspacer{} สพฺเพปิ เจ สาริปุตฺต เวสฺสา.\csromanspacer{} สพฺเพปิ เจ สาริปุตฺต สุทฺทา อิมสฺส มยา สํขิตฺเตน ภาสิตสฺส เอวํ วิตฺถาเรน อตฺถํ อาชาเนยฺยุํ สพฺเพสานมฺปิสฺส สุทฺทานํ ทีฆรตฺตํ หิตาย สุขาย.\csromanspacer{} สเทวโกปิ เจ สาริปุตฺต โลโก สมารโก สพฺรหฺมโก สสฺสมณพฺราหฺมณี ปชา สเทวมนุสฺสา อิมสฺส มยา สํขิตฺเตน ภาสิตสฺส เอวํ วิตฺถาเรน อตฺถํ อาชาเนยฺย สเทวกสฺสปิสฺส โลกสฺส สมารกสฺส สพฺรหฺมกสฺส สสฺสมณพฺราหฺมณิยา ปชาย สเทวมนุสฺสาย ทีฆรตฺตํ หิตาย สุขายาติ.}

\prose{อิทมโวจ ภควา.\csromanspacer{} อตฺตมโน อายสฺมา สาริปุตฺโต ภควโต ภาสิตํ อภินนฺทีติ.}

\nitthitamruled{เสวิตพฺพาเสวิตพฺพสุตฺตํ นิฏฺฐิตํ จตุตฺถํ.}
\csromanmatikaanchor{mtk.18}
\csromantocmarkref{subsection}{5. พหุธาตุกสุตฺต}{mtk.18}
\bookmark[dest={mtk.18},level=2]{5. พหุธาตุกสุตฺต}
\csromanheader{2}{5. พหุธาตุกสุตฺต}
\proseitem{124}{เอวํ เม สุตํ\csromandash{}เอกํ สมยํ ภควา สาวตฺถิยํ วิหรติ เชตวเน อนาถปิณฺฑิกสฺส อาราเม.\csromanspacer{} ตตฺร โข ภควา ภิกฺขู อามนฺเตสิ “ภิกฺขโว”ติ. “ภทนฺเต”ติ เต ภิกฺขู ภควโต ปจฺจสฺโสสุํ.\csromanspacer{} ภควา เอตทโวจ\csromandash{}}

\prose{ยานิ กานิจิ ภิกฺขเว ภยานิ อุปฺปชฺชนฺติ, สพฺพานิ ตานิ พาลโต อุปฺปชฺชนฺติ, โน ปณฺฑิตโต.\csromanspacer{} เย เกจิ อุปทฺทวา อุปฺปชฺชนฺติ, สพฺเพ เต พาลโต อุปฺปชฺชนฺติ, โน ปณฺฑิตโต.\csromanspacer{} เย เกจิ อุปสคฺคา อุปฺปชฺชนฺติ, สพฺเพ เต พาลโต อุปฺปชฺชนฺติ, โน ปณฺฑิตโต.\csromanspacer{} เสยฺยถาปิ}

\vspace{\tipitakaparskip}
\csromancenter{พหุธาตุกสุตฺต (115) ภิกฺขเว นฬาคารา วา ติณาคารา วา อคฺคิ มุตฺโต\footnote{อคฺคิมุกฺโก (สี, อิ)} กูฏาคารานิปิ ทหติ อุลฺลิตฺตาวลิตฺตานิ นิวาตานิ ผุสิตคฺคฬานิ ปิหิตวาตปานานิ.\csromanspacer{} เอวเมว โข ภิกฺขเว ยานิ กานิจิ ภยานิ อุปฺปชฺชนฺติ, สพฺพานิ ตานิ พาลโต อุปฺปชฺชนฺติ, โน ปณฺฑิตโต.\csromanspacer{} เย เกจิ อุปทฺทวา อุปฺปชฺชนฺติ, สพฺเพ เต พาลโต อุปฺปชฺชนฺติ, โน ปณฺฑิตโต.\csromanspacer{} เย เกจิ อุปสคฺคา อุปฺปชฺชนฺติ, สพฺเพ เต พาลโต อุปฺปชฺชนฺติ, โน ปณฺฑิตโต.\csromanspacer{} อิติ โข ภิกฺขเว สปฺปฏิภโย พาโล, อปฺปฏิภโย ปณฺฑิโต.\csromanspacer{} ส- อุปทฺทโว พาโล, อนุปทฺทโว ปณฺฑิโต.\csromanspacer{} ส-อุปสคฺโค พาโล, อนุปสคฺโค ปณฺฑิโต.\csromanspacer{} นตฺถิ ภิกฺขเว ปณฺฑิตโต ภยํ, นตฺถิ ปณฺฑิตโต อุปทฺทโว, นตฺถิ ปณฺฑิตโต อุปสคฺโค.\csromanspacer{} ตสฺมาติห ภิกฺขเว ปณฺฑิตา ภวิสฺสาม วีมํสกาติ, เอวํ หิ โว ภิกฺขเว สิกฺขิตพฺพนฺติ.}
\prose{\csromanfolio{107}เอวํ วุตฺเต อายสฺมา อานนฺโท ภควนฺตํ เอตทโวจ “กิตฺตาวตา นุ โข ภนฺเต ‘ปณฺฑิโต ภิกฺขุ วีมํสโก’ติ อลํ วจนายา”ติ.\csromanspacer{} ยโต โข อานนฺท ภิกฺขุ ธาตุกุสโล จ โหติ, อายตนกุสโล จ โหติ, ปฏิจฺจสมุปฺปาทกุสโล จ โหติ, ฐานาฐานกุสโล จ โหติ, เอตฺตาวตา โข อานนฺท “ปณฺฑิโต ภิกฺขุ วีมํสโก”ติ อลํ วจนายาติ.}

\proseitem{125}{กิตฺตาวตา ปน ภนฺเต “ธาตุกุสโล ภิกฺขู”ติ อลํ วจนายาติ.\csromanspacer{} อฏฺฐารส โข อิมา อานนฺท ธาตุโย, จกฺขุธาตุ รูปธาตุ จกฺขุวิญฺญาณธาตุ, โสตธาตุ สทฺทธาตุ โสตวิญฺญาณธาตุ, ฆานธาตุ คนฺธธาตุ ฆานวิญฺญาณธาตุ, ชิวฺหาธาตุ รสธาตุ ชิวฺหาวิญฺญาณธาตุ, กายธาตุ โผฏฺฐพฺพธาตุ กายวิญฺญาณธาตุ, มโนธาตุ ธมฺมธาตุ มโนวิญฺญาณธาตุ.\csromanspacer{} อิมา โข อานนฺท อฏฺฐารส ธาตุโย ยโต ชานาติ ปสฺสติ, เอตฺตาวตาปิ โข อานนฺท “ธาตุกุสโล ภิกฺขู”ติ อลํ วจนายาติ. (1)}

\prose{สิยา ปน ภนฺเต อญฺโญปิ ปริยาโย ยถา “ธาตุกุสโล ภิกฺขู”ติ อลํ วจนายาติ.\csromanspacer{} สิยา อานนฺท.\csromanspacer{} ฉยิมา อานนฺท ธาตุโย, ปถวีธาตุ อาโปธาตุ เตโชธาตุ วาโยธาตุ อากาสธาตุ วิญฺญาณธาตุ.\csromanspacer{} อิมา โข อานนฺท ฉ ธาตุโย ยโต \csromanfolio{108}ชานาติ ปสฺสติ, เอตฺตาวตาปิ โข อานนฺท “ธาตุกุสโล ภิกฺขู”ติ อลํ วจนายาติ. (2)}

\prose{สิยา ปน ภนฺเต อญฺโญปิ ปริยาโย ยถา “ธาตุกุสโล ภิกฺขู”ติ อลํ วจนายาติ.\csromanspacer{} สิยา อานนฺท.\csromanspacer{} ฉยิมา อานนฺท ธาตุโย, สุขธาตุ ทุกฺขธาตุ โสมนสฺสธาตุ โทมนสฺสธาตุ อุเปกฺขาธาตุ อวิชฺชาธาตุ.\csromanspacer{} อิมา โข อานนฺท ฉ ธาตุโย ยโต ชานาติ ปสฺสติ, เอตฺตาวตาปิ โข อานนฺท “ธาตุกุสโล ภิกฺขู”ติ อลํ วจนายาติ. (3)}

\prose{สิยา ปน ภนฺเต อญฺโญปิ ปริยาโย ยถา “ธาตุกุสโล ภิกฺขู”ติ อลํ วจนายาติ.\csromanspacer{} สิยา อานนฺท.\csromanspacer{} ฉ ยิมา อานนฺท ธาตุโย, กามธาตุ เนกฺขมฺมธาตุ พฺยาปาทธาตุ อพฺยาปาทธาตุ วิหิํสาธาตุ อวิหิํสาธาตุ.\csromanspacer{} อิมา โข อานนฺท ฉ ธาตุโย ยโต ชานาติ ปสฺสติ, เอตฺตาวตาปิ โข อานนฺท “ธาตุกุสโล ภิกฺขู”ติ อลํ วจนายาติ. (4)}

\prose{สิยา ปน ภนฺเต อญฺโญปิ ปริยาโย ยถา “ธาตุกุสโล ภิกฺขู”ติ อลํ วจนายาติ.\csromanspacer{} สิยา อานนฺท.\csromanspacer{} ติสฺโส อิมา อานนฺท ธาตุโย, กามธาตุ รูปธาตุ อรูปธาตุ.\csromanspacer{} อิมา โข อานนฺท ติสฺโส ธาตุโย ยโต ชานาติ ปสฺสติ, เอตฺตาวตาปิ โข อานนฺท “ธาตุกุสโล ภิกฺขู”ติ อลํ วจนายาติ. (5)}

\prose{สิยา ปน ภนฺเต อญฺโญปิ ปริยาโย ยถา “ธาตุกุสโล ภิกฺขู”ติ อลํ วจนายาติ.\csromanspacer{} สิยา อานนฺท.\csromanspacer{} เทฺว อิมา อานนฺท ธาตุโย, สงฺขตาธาตุ อสงฺขตาธาตุ.\csromanspacer{} อิมา โข อานนฺท เทฺว ธาตุโย ยโต ชานาติ ปสฺสติ, เอตฺตาวตาปิ โข อานนฺท “ธาตุกุสโล ภิกฺขู”ติ อลํ วจนายาติ. (6)}

\proseitem{126}{กิตฺตาวตา ปน ภนฺเต “อายตนกุสโล ภิกฺขู”ติ อลํ วจนายาติ.\csromanspacer{} ฉ โข ปนิมานิ อานนฺท อชฺฌตฺติกพาหิรานิ อายตนานิ.\csromanspacer{} จกฺขุเจว รูปา จ โสตญฺจ สทฺทา จ ฆานญฺจ คนฺธา จ ชิวฺหา จ รสา จ กาโย จ โผฏฺฐพฺพา จ มโน จ ธมฺมา จ.\csromanspacer{} อิมานิ โข อานนฺท ฉ อชฺฌตฺติกพาหิรานิ อายตนานิ ยโต ชานาติ ปสฺสติ, เอตฺตาวตา โข อานนฺท “อายตนกุสโล ภิกฺขู”ติ อลํ วจนายาติ.}

\csromanheader{2}{5. พหุธาตุกสุตฺต (115)}
\prose{\csromanfolio{109}กิตฺตาวตา ปน ภนฺเต “ปฏิจฺจสมุปฺปาทกุสโล ภิกฺขู”ติ อลํ วจนายาติ.\csromanspacer{} อิธานนฺท ภิกฺขุ เอวํ ปชานาติ\csromandash{}อิมสฺมิํ สติ อิทํ โหติ, อิมสฺสุปฺปาทา อิทํ อุปฺปชฺชติ, อิมสฺมิํ อสติ อิทํ น โหติ, อิมสฺส นิโรธา อิทํ นิรุชฺฌติ, ยทิทํ อวิชฺชาปจฺจยา สงฺขารา, สงฺขารปจฺจยา วิญฺญาณํ, วิญฺญาณปจฺจยา นามรูปํ, นามรูปปจฺจยา สฬายตนํ, สฬายตนปจฺจยา ผสฺโส, ผสฺสปจฺจยา เวทนา, เวทนาปจฺจยา ตณฺหา, ตณฺหาปจฺจยา อุปาทานํ, อุปาทานปจฺจยา ภโว, ภวปจฺจยา ชาติ, ชาติปจฺจยา ชรามรณํ โสกปริเทวทุกฺขโทมนสฺสูปายาสา สมฺภวนฺติ, เอวเมตสฺส เกวลสฺส ทุกฺขกฺขนฺธสฺส สมุทโย โหติ.\csromanspacer{} อวิชฺชายเตฺวว อเสสวิราคนิโรธา สงฺขารนิโรโธ, สงฺขารนิโรธา วิญฺญาณนิโรโธ, วิญฺญาณนิโรธา นามรูปนิโรโธ, นามรูปนิโรธา สฬายตนนิโรโธ, สฬายตนนิโรธา ผสฺสนิโรโธ, ผสฺสนิโรธา เวทนานิโรโธ, เวทนานิโรธา ตณฺหานิโรโธ, ตณฺหานิโรธา อุปาทานนิโรโธ, อุปาทานนิโรธา ภวนิโรโธ, ภวนิโรธา ชาตินิโรโธ, ชาตินิโรธา ชรามรณํ โสกปริเทวทุกฺขโทมนสฺสูปายาสา นิรุชฺฌนฺติ, เอวเมตสฺส เกวลสฺส ทุกฺขกฺขนฺธสฺส นิโรโธ โหติ.\csromanspacer{} เอตฺตาวตา โข อานนฺท “ปฏิจฺจสมุปฺปาทกุสโล ภิกฺขู”ติ อลํ วจนายาติ.}

\proseitem{127}{กิตฺตาวตา ปน ภนฺเต “ฐานาฐานกุสโล ภิกฺขู”ติ อลํ วจนายาติ.\csromanspacer{} อิธานนฺท ภิกฺขุ อฏฺฐานเมตํ อนวกาโส, ยํ ทิฏฺฐิสมฺปนฺโน ปุคฺคโล กิญฺจิ\footnote{กิญฺจิ (สฺยา, กํ, ก)} สงฺขารํ นิจฺจโต อุปคจฺเฉยฺย, เนตํ ฐานํ วิชฺชตีติ ปชานาติ.\csromanspacer{} ฐานญฺจ โข เอตํ วิชฺชติ, ยํ ปุถุชฺชโน กญฺจิ สงฺขารํ นิจฺจโต อุปคจฺเฉยฺย, ฐานเมตํ วิชฺชตีติ ปชานาติ.\csromanspacer{} อฏฺฐานเมตํ อนวกาโส, ยํ ทิฏฺฐิสมนฺโน ปุคฺคโล กญฺจิ สงฺขารํ สุขโต อุปคจฺเฉยฺย, เนตํ ฐานํ วิชฺชตีติ ปชานาติ.\csromanspacer{} ฐานญฺจ โข เอตํ วิชฺชติ, ยํ ปุถุชฺชโน กญฺจิ สงฺขารํ สุขโต อุปคจฺเฉยฺย, ฐานเมตํ วิชฺชตีติ ปชานาติ.\csromanspacer{} อฏฺฐานเมตํ อนวกาโส, ยํ ทิฏฺฐิสมฺปนฺโน ปุคฺคโล กญฺจิ ธมฺมํ อตฺตโต อุปคจฺเฉยฺย, เนตํ ฐานํ วิชฺชตีติ ปชานาติ.\csromanspacer{} ฐานญฺจ โข เอตํ วิชฺชติ, ยํ ปุถุชฺชโน กญฺจิ ธมฺมํ อตฺตโต อุปคจฺเฉยฺย, ฐานเมตํ วิชฺชตีติ ปชานาติ.}

\proseitem{128}{\csromanfolio{110}อฏฺฐานเมตํ อนวกาโส, ยํ ทิฏฺฐิสมฺปนฺโน ปุคฺคโล มาตรํ ชีวิตา โวโรเปยฺย, เนตํ ฐานํ วิชฺชตีติ ปชานาติ.\csromanspacer{} ฐานญฺจ โข เอตํ วิชฺชติ, ยํ ปุถุชฺชโน มาตรํ ชีวิตา โวโรเปยฺย, ฐานเมตํ วิชฺชตีติ ปชานาติ.\csromanspacer{} อฏฺฐานเมตํ อนวกาโส, ยํ ทิฏฺฐิสมฺปนฺโน ปุคฺคโล ปิตรํ ชีวิตา โวโรเปยฺย ฯเปฯ อรหนฺตํ ชีวิตา โวโรเปยฺย, ฐานเมตํ วิชฺชตีติ ปชานาติ.\csromanspacer{} อฏฺฐานเมตํ อนวกาโส, ยํ ทิฏฺฐิสมฺปนฺโน ปุคฺคโล ทุฏฺฐจิตฺโต ตถาคตสฺส โลหิตํ อุปฺปาเทยฺย, เนตํ ฐานํ วิชฺชตีติ ปชานาติ.\csromanspacer{} ฐานญฺจ โข เอตํ วิชฺชติ, ยํ ปุถุชฺชโน ทุฏฺฐจิตฺโต ตถาคตสฺส โลหิตํ อุปฺปาเทยฺย, ฐานเมตํ วิชฺชตีติ ปชานาติ.\csromanspacer{} อฏฺฐานเมตํ อนวกาโส, ยํ ทิฏฺฐิสมฺปนฺโน ปุคฺคโล สํฆํ ภินฺเทยฺย, เนตํ ฐานํ วิชฺชตีติ ปชานาติ.\csromanspacer{} ฐานญฺจ โข เอตํ วิชฺชติ, ยํ ปุถุชฺชโน สํฆํ ภินฺเทยฺย, ฐานเมตํ วิชฺชตีติ ปชานาติ.\csromanspacer{} อฏฺฐานเมตํ อนวกาโส, ยํ ทิฏฺฐิสมฺปนฺโน ปุคฺคโล อญฺญํ สตฺถารํ อุทฺทิเสยฺย, เนตํ ฐานํ วิชฺชตีติ ปชานาติ.\csromanspacer{} ฐานญฺจ โข เอตํ วิชฺชติ, ยํ ปุถุชฺชโน อญฺญํ สตฺถารํ อุทฺทิเสยฺย, ฐานเมตํ วิชฺชตีติ ปชานาติ.}

\proseitem{129}{อฏฺฐานเมตํ อนวกาโส, ยํ เอกิสฺสา โลกธาตุยา เทฺว อรหนฺโต สมฺมาสมฺพุทฺธา อปุพฺพํ อจริมํ อุปฺปชฺเชยฺยุํ, เนตํ ฐานํ วิชฺชตีติ ปชานาติ.\csromanspacer{} ฐานญฺจ โข เอตํ วิชฺชติ, ยํ เอกิสฺสา โลกธาตุยา เอโก อรหํ สมฺมาสมฺพุทฺโธ อุปฺปชฺเชยฺย, ฐานเมตํ วิชฺชตีติ ปชานาติ.\csromanspacer{} อฏฺฐานเมตํ อนวกาโส, ยํ เอกิสฺสา โลกธาตุยา เทฺว ราชาโน จกฺกวตฺติโน อปุพฺพํ อจริมํ อุปฺปชฺเชยฺยุํ, เนตํ ฐานํ วิชฺชตีติ ปชานาติ.\csromanspacer{} ฐานญฺจ โข เอตํ วิชฺชติ, ยํ เอกิสฺสา โลกธาตุยา เอโก ราชา จกฺกวตฺตี อุปฺปชฺเชยฺย, ฐานเมตํ วิชฺชตีติ ปชานาติ.}

\proseitem{130}{อฏฺฐานเมตํ อนวกาโส, ยํ อิตฺถี อรหํ อสฺส สมฺมาสมฺพุทฺโธ, เนตํ ฐานํ วิชฺชตีติ ปชานาติ.\csromanspacer{} ฐานญฺจ โข เอตํ วิชฺชติ, ยํ ปุริโส อรหํ อสฺส สมฺมาสมฺพุทฺโธ, ฐานเมตํ วิชฺชตีติ ปชานาติ.\csromanspacer{} อฏฺฐานเมตํ อนวกาโส, ยํ อิตฺถี ราชา อสฺส จกฺกวตฺตี, เนตํ ฐานํ วิชฺชตีติ ปชานาติ.\csromanspacer{} ฐานญฺจ โข เอตํ วิชฺชติ, ยํ ปุริโส ราชา อสฺส จกฺกวตฺตี, ฐานเมตํ วิชฺชตีติ ปชานาติ.\csromanspacer{} อฏฺฐานเมตํ อนวกาโส, ยํ อิตฺถี สกฺกตฺตํ กเรยฺย.\csromanspacer{} มารตฺตํ กเรยฺย.\csromanspacer{} พฺรหฺมตฺตํ กเรยฺย, เนตํ ฐานํ วิชฺชตีติ ปชานาติ.\csromanspacer{} ฐานญฺจ โข เอตํ วิชฺชติ, ยํ ปุริโส}

\vspace{\tipitakaparskip}
\csromancenter{พหุธาตุกสุตฺต (115) สกฺกตฺตํ กเรยฺย.\csromanspacer{} มารตฺตํ กเรยฺย.\csromanspacer{} พฺรหฺมตฺตํ กเรยฺย, ฐานเมตํ วิชฺชตีติ ปชานาติ.}
\proseitem{131}{\csromanfolio{111}อฏฺฐานเมตํ อนวกาโส, ยํ กายทุจฺจริตสฺส อิฏฺโฐ กนฺโต มนาโป วิปาโก นิพฺพตฺเตยฺย, เนตํ ฐาตํ วิชฺชตีติ ปชานาติ.\csromanspacer{} ฐานญฺจ โข เอตํ วิชฺชติ, ยํ กายทุจฺจริตสฺส อนิฏฺโฐ อกนฺโต อมนาโป วิปาโก นิพฺพตฺเตยฺย, ฐานเมตํ วิชฺชตีติ ปชานาติ.\csromanspacer{} อฏฺฐานเมตํ อนวกาโส, ยํ วจีทุจฺจริตสฺส ฯเปฯ ยํ มโนทุจฺจริตสฺส อิฏฺโฐ กนฺโต มนาโป วิปาโก นิพฺพตฺเตยฺย, เนตํ ฐานํ วิชฺชตีติ ปชานาติ.\csromanspacer{} ฐานญฺจ โข เอตํ วิชฺชติ, ยํ วจีทุจฺจริตสฺส ฯเปฯ ยํ มโนทุจฺจริตสฺส อนิฏฺโฐ อกนฺโต อมนาโป วิปาโก นิพฺพตฺเตยฺย, ฐานเมตํ วิชฺชตีติ ปชานาติ.\csromanspacer{} อฏฺฐานเมตํ อนวกาโส, ยํ กายสุจริตสฺส อนิฏฺโฐ อกนฺโต อมนาโป วิปาโก นิพฺพตฺเตยฺย, เนตํ ฐานํ วิชฺชตีติ ปชานาติ.\csromanspacer{} ฐานญฺจ โข เอตํ วิชฺชติ, ยํ กายสุจริตสฺส อิฏฺโฐ กนฺโต มนาโป วิปาโก นิพฺพตฺเตยฺย, ฐานเมตํ วิชฺชตีติ ปชานาติ.\csromanspacer{} อฏฺฐานเมตํ อนวกาโส, ยํ วจีสุจริตสฺส ฯเปฯ ยํ มโนสุจริตสฺส อนิฏฺโฐ อกนฺโต อมนาโป วิปาโก นิพฺพตฺเตยฺย, เนตํ ฐานํ วิชฺชตีติ ปชานาติ.\csromanspacer{} ฐานญฺจ โข เอตํ วิชฺชติ, ยํ วจีสุจริตสฺส ฯเปฯ ยํ มโนสุจริตสฺส อิฏฺโฐ กนฺโต มนาโป วิปาโก นิพฺพตฺเตยฺย, ฐานเมตํ วิชฺชตีติ ปชานาติ.}

\prose{อฏฺฐานเมตํ อนวกาโส, ยํ กายทุจฺจริตสมงฺคี ตํนิทานา ตปฺปจฺจยา กายสฺส เภทา ปรํ มรณา สุคติํ สคฺคํ โลกํ อุปปชฺเชยฺย, เนตํ ฐานํ วิชฺชตีติ ปชานาติ.\csromanspacer{} ฐานญฺจ โข เอตํ วิชฺชติ, ยํ กายทุจฺจริตสมงฺคี ตํนิทานา ตปฺปจฺจยา กายสฺส เภทา ปรํ มรณา อปายํ ทุคฺคติํ วินิปาตํ นิรยํ อุปปชฺเชยฺย, ฐานเมตํ วิชฺชตีติ ปชานาติ.\csromanspacer{} อฏฺฐานเมตํ อนวกาโส, ยํ วจีทุจฺจริตสมงฺคี ฯเปฯ ยํ มโนทุจฺจริตสมงฺคี ตํนิทานา ตปฺปจฺจยา กายสฺส เภทา ปรํ มรณา สุคติํ สคฺคํ โลกํ อุปปชฺเชยฺย, เนตํ ฐานํ วิชฺชตีติ ปชานาติ.\csromanspacer{} ฐานญฺจ โข เอตํ วิชฺชติ, ยํ วจีทุจฺจริตสมงฺคี ฯเปฯ ยํ มโนทุจฺจริตสมงฺคี ตํนิทานา ตปฺปจฺจยา กายสฺส เภทา ปรํ มรณา อปายํ ทุคฺคติํ วินิปาตํ นิรยํ อุปปชฺเชยฺย, ฐานเมตํ วิชฺชตีติ ปชานาติ.\csromanspacer{} อฏฺฐานเมตํ อนวกาโส, ยํ กายสุจริตสมงฺคี ตํนิทานา ตปฺปจฺจยา กายสฺส เภทา ปรํ มรณา อปายํ \csromanfolio{112}ทุคฺคติํ วินิปาตํ นิรยํ อุปปชฺเชยฺย, เนตํ ฐานํ วิชฺชตีติ ปชานาติ.\csromanspacer{} ฐานญฺจ โข เอตํ วิชฺชติ, ยํ กายสุจริตสมงฺคี ตํนิทานา ตปฺปจฺจยา กายสฺส เภทา ปรํ มรณา สุคติํ สคฺคํ โลกํ อุปปชฺเชยฺย, ฐานเมตํ วิชฺชตีติ ปชานาติ.\csromanspacer{} อฏฺฐานเมตํ อนวกาโส, ยํ วจีสุจริตสมงฺคี ฯเปฯ ยํ มโนสุจริตสมงฺคี ตํนิทานา ตปฺปจฺจยา กายสฺส เภทา ปรํ มรณา อปายํ ทุคฺคติํ วินิปาตํ นิรยํ อุปปชฺเชยฺย, เนตํ ฐานํ วิชฺชตีติ ปชานาติ.\csromanspacer{} ฐานญฺจ โข เอตํ วิชฺชติ, ยํ วจีสุจริตสมงฺคี ฯเปฯ ยํ มโนสุจริตสมงฺคี ตํนิทานา ตปฺปจฺจยา กายสฺส เภทา ปรํ มรณา สุคติํ สคฺคํ โลกํ อุปปชฺเชยฺย, ฐานเมตํ วิชฺชตีติ ปชานาติ.\csromanspacer{} เอตฺตาวตา โข อานนฺท “ฐานาฐานกุสโล ภิกฺขู”ติ อลํ วจนายาติ.}

\proseitem{132}{เอวํ วุตฺเต อายสฺมา อานนฺโท ภควนฺตํ เอตทโวจ “อจฺฉริยํ ภนฺเต, อพฺภุตํ ภนฺเต, โกนาโม อยํ ภนฺเต ธมฺมปริยาโย”ติ.\csromanspacer{} ตสฺมาติห ตฺวํ อานนฺท อิมํ ธมฺมปริยายํ พหุธาตุโกติปิ นํ ธาเรหิ, จตุปริวฏฺโฏติปิ นํ ธาเรหิ, ธมฺมาทาโสติปิ นํ ธาเรหิ, อมตทุนฺทุภีติปิ\footnote{ทุทฺรภีติปิ (ก)} นํ ธาเรหิ, อนุตฺตโร สงฺคามวิชโยติปิ นํ ธาเรหีติ.}

\prose{อิทมโวจ ภควา.\csromanspacer{} อตฺตมโน อายสฺมา อานนฺโท ภควโต ภาสิตํ อภินนฺทีติ.}

\nitthitamruled{พหุธาตุกสุตฺตํ นิฏฺฐิตํ ปญฺจมํ.}
\csromanmatikaanchor{mtk.19}
\csromantocmarkref{subsection}{6. อิสิคิลิสุตฺต}{mtk.19}
\bookmark[dest={mtk.19},level=2]{6. อิสิคิลิสุตฺต}
\csromanheader{2}{6. อิสิคิลิสุตฺต}
\proseitem{133}{เอวํ เม สุตํ\csromandash{}เอกํ สมยํ ภควา ราชคเห วิหรติ อิสิคิลิสฺมิํ ปพฺพเต.\csromanspacer{} ตตฺร โข ภควา ภิกฺขู อามนฺเตสิ “ภิกฺขโว”ติ. “ภทนฺเต”ติ เต ภิกฺขู ภควโต ปจฺจสฺโสสุํ.\csromanspacer{} ภควา เอตทโวจ\csromandash{}}

\prose{ปสฺสถ โน ตุเมฺห ภิกฺขเว เอตํ เวภารํ ปพฺพตนฺติ.\csromanspacer{} เอวํ ภนฺเต.\csromanspacer{} เอตสฺสปิ โข ภิกฺขเว เวภารสฺส ปพฺพตสฺส อญฺญาว สมญฺญา อโหสิ อญฺญา ปญฺญตฺติ.}

\csromanheader{2}{6. อิสิคิลิสุตฺต (116)}
\prose{\csromanfolio{113}ปสฺสถ โน ตุเมฺห ภิกฺขเว เอตํ ปณฺฑวํ ปพฺพตนฺติ.\csromanspacer{} เอวํ ภนฺเต.\csromanspacer{} เอตสฺสปิ โข ภิกฺขเว ปณฺฑวสฺส ปพฺพตสฺส อญฺญาว สมญฺญา อโหสิ อญฺญา ปญฺญตฺติ.}

\prose{ปสฺสถ โน ตุเมฺห ภิกฺขเว เอตํ เวปุลฺลํ ปพฺพตนฺติ.\csromanspacer{} เอวํ ภนฺเต.\csromanspacer{} เอตสฺสปิ โข ภิกฺขเว เวปุลฺลสฺส ปพฺพตสฺส อญฺญาว สมญฺญา อโหสิ อญฺญา ปญฺญตฺติ.}

\prose{ปสฺสถ โน ตุเมฺห ภิกฺขเว เอตํ คิชฺฌกูฏํ ปพฺพตนฺติ.\csromanspacer{} เอวํ ภนฺเต.\csromanspacer{} เอตสฺสปิ โข ภิกฺขเว คิชฺฌกูฏสฺส ปพฺพตสฺส อญฺญาว สมญฺญา อโหสิ อญฺญา ปญฺญตฺติ.}

\prose{ปสฺสถ โน ตุเมฺห ภิกฺขเว อิมํ อิสิคิลิํ ปพฺพตนฺติ.\csromanspacer{} เอวํ ภนฺเต.\csromanspacer{} อิมสฺส โข ปน ภิกฺขเว อิสิคิลิสฺส ปพฺพตสฺส เอสาว สมญฺญา อโหสิ เอสา ปญฺญตฺติ.}

\prose{ภูตปุพฺพํ ภิกฺขเว ปญฺจ ปจฺเจกพุทฺธสตานิ อิมสฺมิํ อิสิคิลิสฺมิํ ปพฺพเต จิรนิวาสิโน อเหสุํ.\csromanspacer{} เต อิมํ ปพฺพตํ ปวิสนฺตา ทิสฺสนฺติ, ปวิฏฺฐา น ทิสฺสนฺติ.\csromanspacer{} ตเมนํ มนุสฺสา ทิสฺวา เอวมาหํสุ “อยํ ปพฺพโต อิเม อิสี\footnote{อิสโย (ก)} คิลตี”ติ อิสิคิลิ อิสิคิลิเตฺวว สมญฺญา อุทปาทิ.\csromanspacer{} อาจิกฺขิสฺสามิ\footnote{อาจิกฺขิสฺสามิ โว (ก)} ภิกฺขเว ปจฺเจกพุทฺธานํ นามานิ, กิตฺตยิสฺสามิ ภิกฺขเว ปจฺเจกพุทฺธานํ นามานิ, เทเสสฺสามิ ภิกฺขเว ปจฺเจกพุทฺธานํ นามานิ, ตํ สุณาถ สาธุกํ มนสิ กโรถ, ภาสิสฺสามีติ. “เอวํ ภนฺเต”ติ โข เต ภิกฺขู ภควโต ปจฺจสฺโสสุํ.\csromanspacer{} ภควา เอตทโวจ\csromandash{}}

\proseitem{134}{อริฏฺโฐ นาม ภิกฺขเว ปจฺเจกสมฺพุทฺโธ\footnote{ปจฺเจกพุทฺโธ (ก-สี, อิ)} อิมสฺมิํ อิสิคิลิสฺมิํ ปพฺพเต จิรนิวาสี อโหสิ, อุปริฏฺโฐ นาม ภิกฺขเว ปจฺเจกสมฺพุทฺโธ อิมสฺมิํ อิสิคิลิสฺมิํ ปพฺพเต จิรนิวาสี อโหสิ, ตครสิขี\footnote{ตคฺครสิขี (ก)} นาม ภิกฺขเว ปจฺเจกสมฺพุทฺโธ อิมสฺมิํ อิสิคิลิสฺมิํ ปพฺพเต จิรนิวาสี อโหสิ, ยสสฺสี นาม ภิกฺขเว ปจฺเจกสมฺพุทฺโธ อิมสฺมิํ อิสิคิลิสฺมิํ ปพฺพเต จิรนิวาสี อโหสิ, สุทสฺสโน นาม ภิกฺขเว ปจฺเจกสมฺพุทฺโธ อิมสฺมิํ อิสิคิลิสฺมิํ ปพฺพเต จิรนิวาสี อโหสิ, ปิยทสฺสี นาม ภิกฺขเว \csromanfolio{114}ปจฺเจกสมฺพุทฺโธ อิมสฺมิํ อิสิคิลิสฺมิํ ปพฺพเต จิรนิวาสี อโหสิ, คนฺธาโร นาม ภิกฺขเว ปจฺเจกสมฺพุทฺโธ อิมสฺมิํ อิสิคิลิสฺมิํ ปพฺพเต จิรนิวาสี อโหสิ, ปิณฺโฑโล นาม ภิกฺขเว ปจฺเจกสมฺพุทฺโธ อิมสฺมิํ อิสิคิลิสฺมิํ ปพฺพเต จิรนิวาสี อโหสิ, อุปาสโภ นาม ภิกฺขเว ปจฺเจกสมฺพุทฺโธ อิมสฺมิํ อิสิคิลิสฺมิํ ปพฺพเต จิรนิวาสี อโหสิ, นีโต นาม ภิกฺขเว ปจฺเจกสมฺพุทฺโธ อิมสฺมิํ อิสิคิลิสฺมิํ ปพฺพเต จิรนิวาสี อโหสิ, ตโถ นาม ภิกฺขเว ปจฺเจกสมฺพุทฺโธ อิมสฺมิํ อิสิคิลิสฺมิํ ปพฺพเต จิรนิวาสี อโหสิ, สุตวา นาม ภิกฺขเว ปจฺเจกสมฺพุทฺโธ อิมสฺมิํ อิสิคิลิสฺมิํ ปพฺพเต จิรนิวาสี อโหสิ, ภาวิตตฺโต นาม ภิกฺขเว ปจฺเจกสมฺพุทฺโธ อิมสฺมิํ อิสิคิลิสฺมิํ ปพฺพเต จิรนิวาสี อโหสิ.}

\csromanhangingbegin
\hangingheaditem{135}{เย สตฺตสารา อนีฆา นิราสา,}
\hangingbody{ปจฺเจกเมวชฺฌคมํสุ โพธิํ1.}
\hangingbody{เตสํ วิสลฺลาน นรุตฺตมานํ,}
\hangingbody{นามานิ เม กิตฺตยโต สุณาถ.}
\hangingbody{อริฏฺโฐ อุปริฏฺโฐ ตครสิขี ยสสฺสี,}
\hangingbody{สุทสฺสโน ปิยทสฺสี จ สุสมฺพุทฺโธ2.}
\hangingbody{คนฺธาโร ปิณฺโฑโล อุปาสโภ จ,}
\hangingbody{นีโต ตโถ สุตวา ภาวิตตฺโต.}
\hangingbody{สุมฺโภ สุโภ มตุโล3 อฏฺฐโม จ,}
\hangingbody{อถสฺสุเมโฆ4 อนีโฆ สุทาโฐ.}
\hangingbody{ปจฺเจกพุทฺธา ภวเนตฺติขีณา,}
\hangingbody{หิงฺคู จ หิงฺโค จ มหานุภาวา.}
\hangingbody{เทฺว ชาลิโน มุนิโน อฏฺฐโก จ,}
\hangingbody{อถ โกสลฺโล พุทฺโธ อโถ สุพาหุ.}
\hangingbody{อุปเนมิโส เนมิโส สนฺตจิตฺโต,}
\hangingbody{สจฺโจ ตโถ วิรโช ปณฺฑิโต จ.}
\csromanhangingend

\csromanheader{2}{6. อิสิคิลิสุตฺต (116)}
\csromangathameasure{กาฬูปกาฬา วิชิโต ชิโต จ, \\ องฺโค จ ปงฺโค จ คุตฺติชิโต จ. \\ ปสฺสิ ชหิ อุปธิทุกฺขมูลํ, \\ อปราชิโต มารพลํ อเชสิ. \\ สตฺถา ปวตฺตา สรภงฺโค โลมหํโส, \\ อุจฺจงฺคมาโย อสิโต อนาสโว. \\ มโนมโย มานจฺฉิโท จ พนฺธุมา, \\ ตทาธิมุตฺโต วิมโล จ เกตุมา. \\ เกตุมฺภราโค จ มาตงฺโค อริโย, \\ อถจฺจุโต อจฺจุตคามพฺยามโก. \\ สุมงฺคโล ทพฺพิโล สุปติฏฺฐิโต, \\ อสโยฺห เขมาภิรโต จ โสรโต. \\ ทุรนฺนโย สํโฆ อโถปิ อุชฺชโย, \\ อปโร มุนิ สโยฺห อโนมนิกฺกโม. \\ อานนฺโท นนฺโท อุปนนฺโท ทฺวาทส, \\ ภารทฺวาโช อนฺติมเทหธารี. \\ โพธิ มหานาโม อโถปิ อุตฺตโร, \\ เกสี สิขี สุนฺทโร ทฺวารภาโช. \\ ติสฺสูปติสฺสา ภวพนฺธนจฺฉิทา, \\ อุปสิขิ ตณฺหจฺฉิโท จ สิขริ. \\ พุทฺโธ อหุ มงฺคโล วีตราโค, \\ อุสภจฺฉิทา ชาลินิํ ทุกฺขมูลํ. \\ สนฺตํ ปทํ อชฺฌคโมปนีโต, \\ อุโปสโถ สุนฺทโร สจฺจนาโม. \\ เชโต ชยนฺโต ปทุโม อุปฺปโล จ, \\ ปทุมุตฺตโร รกฺขิโต ปพฺพโต จ. \\ มานตฺถทฺโธ โสภิโต วีตราโค, \\ กโณฺห จ พุทฺโธ สุวิมุตฺตจิตฺโต. \\ เอเต จ อญฺเญ จ มหานุภาวา, \\ ปจฺเจกพุทฺธา ภวเนตฺติขีณา. \\ เต สพฺพสงฺคาติคเต มเหสี, \\ ปรินิพฺพุเต วนฺทถ อปฺปเมยฺเยติ.}
\csromangathagroup{\csromangathastanza{\csromanfolio{115}กาฬูปกาฬา วิชิโต ชิโต จ, \\ องฺโค จ ปงฺโค จ คุตฺติชิโต จ. \\ ปสฺสิ ชหิ อุปธิทุกฺขมูลํ\footnote{ปสฺสี ชหี อุปธิํ ทุกฺขมูลํ (สี, สฺยา, กํ, อิ)}, \\ อปราชิโต มารพลํ อเชสิ.}\csromangathastanza{สตฺถา ปวตฺตา สรภงฺโค โลมหํโส, \\ อุจฺจงฺคมาโย อสิโต อนาสโว. \\ มโนมโย มานจฺฉิโท จ พนฺธุมา, \\ ตทาธิมุตฺโต วิมโล จ เกตุมา.}\csromangathastanza{เกตุมฺภราโค จ มาตงฺโค อริโย, \\ อถจฺจุโต อจฺจุตคามพฺยามโก. \\ สุมงฺคโล ทพฺพิโล สุปติฏฺฐิโต, \\ อสโยฺห เขมาภิรโต จ โสรโต.}\csromangathastanza{ทุรนฺนโย สํโฆ อโถปิ อุชฺชโย, \\ อปโร มุนิ สโยฺห อโนมนิกฺกโม. \\ อานนฺโท นนฺโท อุปนนฺโท ทฺวาทส, \\ ภารทฺวาโช อนฺติมเทหธารี\footnote{อนฺติมเทหธาริ (สี)}.}\csromangathastanza{โพธิ มหานาโม อโถปิ อุตฺตโร, \\ เกสี สิขี สุนฺทโร ทฺวารภาโช. \\ ติสฺสูปติสฺสา ภวพนฺธนจฺฉิทา, \\ อุปสิขิ ตณฺหจฺฉิโท จ สิขริ\footnote{อุปสีทรี ตณฺหจฺฉิโท จ สีทรี (สี, สฺยา, กํ, อิ)}.}\csromangathastanza{พุทฺโธ อหุ มงฺคโล วีตราโค, \\ อุสภจฺฉิทา ชาลินิํ ทุกฺขมูลํ. \\ สนฺตํ ปทํ อชฺฌคโมปนีโต, \\ อุโปสโถ สุนฺทโร สจฺจนาโม.}\csromangathastanza{เชโต ชยนฺโต ปทุโม อุปฺปโล จ, \\ ปทุมุตฺตโร รกฺขิโต ปพฺพโต จ. \\ มานตฺถทฺโธ โสภิโต วีตราโค, \\ กโณฺห จ พุทฺโธ สุวิมุตฺตจิตฺโต.}\csromangathastanza{\csromanfolio{116}เอเต จ อญฺเญ จ มหานุภาวา, \\ ปจฺเจกพุทฺธา ภวเนตฺติขีณา. \\ เต สพฺพสงฺคาติคเต มเหสี, \\ ปรินิพฺพุเต วนฺทถ อปฺปเมยฺเยติ.}}

\nitthitamruled{อิสิคิลิสุตฺตํ นิฏฺฐิตํ ฉฏฺฐํ.}
\csromanmatikaanchor{mtk.20}
\csromantocmarkref{subsection}{7. มหาจตฺตารีสกสุตฺต}{mtk.20}
\bookmark[dest={mtk.20},level=2]{7. มหาจตฺตารีสกสุตฺต}
\csromanheader{2}{7. มหาจตฺตารีสกสุตฺต}
\proseitem{136}{เอวํ เม สุตํ\csromandash{}เอกํ สมยํ ภควา สาวตฺถิยํ วิหรติ เชตวเน อนาถปิณฺฑิกสฺส อาราเม.\csromanspacer{} ตตฺร โข ภควา ภิกฺขู อามนฺเตสิ “ภิกฺขโว”ติ. “ภทนฺเต”ติ เต ภิกฺขู ภควโต ปจฺจสฺโสสุํ.\csromanspacer{} ภควา เอตทโวจ “อริยํ โว ภิกฺขเว สมฺมาสมาธิํ เทเสสฺสามิ ส-อุปนิสํ สปริกฺขารํ, ตํ สุณาถ สาธุกํ มนสิ กโรถ, ภาสิสฺสามี”ติ. “เอวํ ภนฺเต”ติ โข เต ภิกฺขู ภควโต ปจฺจสฺโสสุํ.\csromanspacer{} ภควา เอตทโวจ\csromandash{}}

\prose{กตโม จ ภิกฺขเว อริโย สมฺมาสมาธิ ส-อุปนิโส สปริกฺขาโร.\csromanspacer{} เสยฺยถิทํ, สมฺมาทิฏฺฐิ สมฺมาสงฺกปฺโป สมฺมาวาจา สมฺมากมฺมนฺโต สมฺมา-อาชีโว สมฺมาวายาโม สมฺมาสติ.\csromanspacer{} ยา โข ภิกฺขเว อิเมหิ สตฺตหงฺเคหิ จิตฺตสฺส เอกคฺคตา ปริกฺขตา, อยํ วุจฺจติ ภิกฺขเว อริโย สมฺมาสมาธิ ส-อุปนิโส อิติปิ สปริกฺขาโร อิติปิ.\csromanspacer{} ตตฺร ภิกฺขเว สมฺมาทิฏฺฐิ ปุพฺพงฺคมา โหติ.\csromanspacer{} กถญฺจ ภิกฺขเว สมฺมาทิฏฺฐิ ปุพฺพงฺคมา โหติ, มิจฺฉาทิฏฺฐิํ “มิจฺฉาทิฏฺฐี”ติ ปชานาติ.\csromanspacer{} สมฺมาทิฏฺฐิํ “สมฺมาทิฏฺฐี”ติ ปชานาติ.\csromanspacer{} สาสฺส โหติ สมฺมาทิฏฺฐิ.}

\prose{กตมา จ ภิกฺขเว มิจฺฉาทิฏฺฐิ.\csromanspacer{} นตฺถิ ทินฺนํ, นตฺถิ ยิฏฺฐํ, นตฺถิ หุตํ, นตฺถิ สุกตทุกฺกฏานํ กมฺมานํ ผลํ วิปาโก, นตฺถิ อยํ โลโก, นตฺถิ ปโร โลโก, นตฺถิ มาตา, นตฺถิ ปิตา, นตฺถิ สตฺตา โอปปาติกา, นตฺถิ โลเก สมณพฺราหฺมณา สมฺมคฺคตา สมฺมาปฏิปนฺนา, เย อิมญฺจ โลกํ ปรญฺจ โลกํ สยํ อภิญฺญา สจฺฉิกตฺวา ปเวเทนฺตีติ.\csromanspacer{} อยํ ภิกฺขเว มิจฺฉาทิฏฺฐิ.}

\csromanheader{2}{7. มหาจตฺตารีสกสุตฺต (117)}
\prose{\csromanfolio{117}กตมา จ ภิกฺขเว สมฺมาทิฏฺฐิ.\csromanspacer{} สมฺมาทิฏฺฐิํปหํ\footnote{สมฺมาทิฏฺฐิมหํ (ก) เอวํ สมฺมาสงฺกปฺปํปหํตฺยาทีสุปิ.} ภิกฺขเว ทฺวายํ\footnote{ทฺวยํ (สี, สฺยา, กํ, อิ) ฏีกา โอโลเกตพฺพา.} วทามิ\csromandash{}อตฺถิ ภิกฺขเว สมฺมาทิฏฺฐิ สาสวา ปุญฺญภาคิยา อุปธิเวปกฺกา, อตฺถิ ภิกฺขเว สมฺมาทิฏฺฐิ อริยา อนาสวา โลกุตฺตรา มคฺคงฺคา.\csromanspacer{} กตมา จ ภิกฺขเว สมฺมาทิฏฺฐิ สาสวา ปุญฺญาภาคิยา อุปธิเวปกฺกา.\csromanspacer{} อตฺถิ ทินฺนํ, อตฺถิ ยิฏฺฐํ, อตฺถิ หุตํ, อตฺถิ สุกตทุกฺกฏานํ กมฺมานํ ผลํ วิปาโก, อตฺถิ อยํ โลโก, อตฺถิ ปโร โลโก, อตฺถิ มาตา, อตฺถิ ปิตา, อตฺถิ สตฺตา โอปปาติกา, อตฺถิ โลเก สมณพฺราหฺมณา สมฺมคฺคตา สมฺมาปฏิปนฺนา, เย อิมญฺจ โลกํ ปรญฺจ โลกํ สยํ อภิญฺญา สจฺฉิกตฺวา ปเวเทนฺตีติ.\csromanspacer{} อยํ ภิกฺขเว สมฺมาทิฏฺฐิ สาสวา ปุญฺญภาคิยา อุปธิเวปกฺกา.}

\prose{กตมา จ ภิกฺขเว สมฺมาทิฏฺฐิ อริยา อนาสวา โลกุตฺตรา มคฺคงฺคา.\csromanspacer{} ยา โข ภิกฺขเว อริยจิตฺตสฺส อนาสวจิตฺตสฺส อริยมคฺคสมงฺคิโน อริยมคฺคํ ภาวยโต ปญฺญา ปญฺญินฺทฺริยํ ปญฺญาพลํ ธมฺมวิจยสมฺโพชฺฌงฺโค สมฺมาทิฏฺฐิ มคฺคงฺคํ\footnote{มคฺคงฺคา (สี, อิ)}.\csromanspacer{} อยํ วุจฺจติ ภิกฺขเว สมฺมาทิฏฺฐิ อริยา อนาสวา โลกุตฺตรา มคฺคงฺคา.\csromanspacer{} โส มิจฺฉาทิฏฺฐิยา ปหานาย วายมติ สมฺมาทิฏฺฐิยา อุปสมฺปทาย, สฺวาสฺส\footnote{สฺวายํ (ก)} โหติ สมฺมาวายาโม.\csromanspacer{} โส สโต มิจฺฉาทิฏฺฐิํ ปชหติ, สโต สมฺมาทิฏฺฐิํ อุปสมฺปชฺช วิหรติ, สาสฺส\footnote{สายํ (ก)} โหติ สมฺมาสติ.\csromanspacer{} อิติยิเม\footnote{อิติเม (สี), อิติสฺสิเม (สฺยา, กํ, อิ)} ตโย ธมฺมา สมฺมาทิฏฺฐิํ อนุปริธาวนฺติ อนุปริวตฺตนฺติ.\csromanspacer{} เสยฺยถิทํ, สมฺมาทิฏฺฐิ สมฺมาวายาโม สมฺมาสติ.}

\proseitem{137}{ตตฺร ภิกฺขเว สมฺมาทิฏฺฐิ ปุพฺพงฺคมา โหติ.\csromanspacer{} กถญฺจ ภิกฺขเว สมฺมาทิฏฺฐิ ปุพฺพงฺคมา โหติ มิจฺฉาสงฺกปฺปํ “มิจฺฉาสงฺกปฺโป”ติ ปชานาติ, สมฺมาสงฺกปฺปํ “สมฺมาสงฺกปฺโป”ติ ปชานาติ, สาสฺส โหติ สมฺมาทิฏฺฐิ.\csromanspacer{} กตโม จ ภิกฺขเว มิจฺฉาสงฺกปฺโป.\csromanspacer{} กามสงฺกปฺโป พฺยาปาทสงฺกปฺโป วิหิํสาสงฺกปฺโป.\csromanspacer{} อยํ ภิกฺขเว มิจฺฉาสงฺกปฺโป.}

\prose{กตโม จ ภิกฺขเว สมฺมาสงฺกปฺโป.\csromanspacer{} สมฺมาสงฺกปฺปํปหํ ภิกฺขเว ทฺวายํ วทามิ\csromandash{}อตฺถิ ภิกฺขเว สมฺมาสงฺกปฺโป สาสโว ปุญฺญภาคิโย อุปธิเวปกฺโก, อตฺถิ ภิกฺขเว สมฺมาสงฺกปฺโป อริโย อนาสโว โลกุตฺตโร มคฺคงฺโค.\csromanspacer{} กตโม จ ภิกฺขเว สมฺมาสงฺกปฺโป สาสโว \csromanfolio{118}ปุญฺญภาคิโย อุปธิเวปกฺโก.\csromanspacer{} เนกฺขมฺมสงฺกปฺโป อพฺยาปาทสงฺกปฺโป อวิหิํสาสงฺกปฺโป.\csromanspacer{} อยํ ภิกฺขเว สมฺมาสงฺกปฺโป สาสโว ปุญฺญภาคิโย อุปธิเวปกฺโก.}

\prose{กตโม จ ภิกฺขเว สมฺมาสงฺกปฺโป อริโย อนาสโว โลกุตฺตโร มคฺคงฺโค.\csromanspacer{} โย โข ภิกฺขเว อริยจิตฺตสฺส อนาสวจิตฺตสฺส อริยมคฺคสมงฺคิโน อริยมคฺคํ ภาวยโต ตกฺโก วิตกฺโก สงฺกปฺโป อปฺปนา พฺยปฺปนา เจตโส อภินิโรปนา วจีสงฺขาโร.\csromanspacer{} อยํ ภิกฺขเว สมฺมาสงฺกปฺโป อริโย อนาสโว โลกุตฺตโร มคฺคงฺโค.\csromanspacer{} โส มิจฺฉาสงฺกปฺปสฺส ปหานาย วายมติ สมฺมาสงฺกปฺปสฺส อุปสมฺปทาย, สฺวาสฺส โหติ สมฺมาวายาโม.\csromanspacer{} โส สโต มิจฺฉาสงฺกปฺปํ ปชหติ, สโต สมฺมาสงฺกปฺปํ อุปสมฺปชฺช วิหรติ, สาสฺส โหติ สมฺมาสติ.\csromanspacer{} อิติยิเม ตโย ธมฺมา สมฺมาสงฺกปฺปํ อนุปริธาวนฺติ อนุปริวตฺตนฺติ.\csromanspacer{} เสยฺยถิทํ, สมฺมาทิฏฺฐิ สมฺมาวายาโม สมฺมาสติ.}

\proseitem{138}{ตตฺร ภิกฺขเว สมฺมาทิฏฺฐิ ปุพฺพงฺคมา โหติ.\csromanspacer{} กถญฺจ ภิกฺขเว สมฺมาทิฏฺฐิ ปุพฺพงฺคมา โหติ, มิจฺฉาวาจํ “มิจฺฉาวาจา”ติ ปชานาติ, สมฺมาวาจํ “สมฺมาวาจา”ติ ปชานาติ, สาสฺส โหติ สมฺมาทิฏฺฐิ.\csromanspacer{} กตมา จ ภิกฺขเว มิจฺฉาวาจา.\csromanspacer{} มุสาวาโท ปิสุณา วาจา ผรุสา วาจา สมฺผปฺปลาโป.\csromanspacer{} อยํ ภิกฺขเว มิจฺฉาวาจา.\csromanspacer{} กตมา จ ภิกฺขเว สมฺมาวาจา.\csromanspacer{} สมฺมาวาจํปหํ ภิกฺขเว ทฺวายํ วทามิ\csromandash{}อตฺถิ ภิกฺขเว สมฺมาวาจา สาสวา ปุญฺญภาคิยา อุปธิเวปกฺกา, อตฺถิ ภิกฺขเว สมฺมาวาจา อริยา อนาสวา โลกุตฺตรา มคฺคงฺคา.\csromanspacer{} กตมา จ ภิกฺขเว สมฺมาวาจา สาสวา ปุญฺญภาคิยา อุปธิเวปกฺกา.\csromanspacer{} มุสาวาทา เวรมณี ปิสุณาย วาจาย เวรมณี ผรุสาย วาจาย เวรมณี สมฺผปฺปลาปา เวรมณี.\csromanspacer{} อยํ ภิกฺขเว สมฺมาวาจา สาสวา ปุญฺญภาคิยา อุปธิเวปกฺกา.\csromanspacer{} กตมา จ ภิกฺขเว สมฺมาวาจา อริยา อนาสวา โลกุตฺตรา มคฺคงฺคา.\csromanspacer{} ยา โข ภิกฺขเว อริยจิตฺตสฺส อนาสวจิตฺตสฺส อริยมคฺคสมงฺคิโน อริยมคฺคํ ภาวยโต จตูหิ วจีทุจฺจริเตหิ อารติ วิรติ ปฏิวิรติ เวรมณี.\csromanspacer{} อยํ ภิกฺขเว สมฺมาวาจา อริยา อนาสวา โลกุตฺตรา มคฺคงฺคา.\csromanspacer{} โส มิจฺฉาวาจาย ปหานาย วายมติ สมฺมาวาจาย อุปสมฺปทาย, สฺวาสฺส โหติ สมฺมาวายาโม.\csromanspacer{} โส สโต มิจฺฉาวาจํ ปชหติ, สโต สมฺมาวาจํ อุปสมฺปชฺช วิหรติ, สาสฺส}

\vspace{\tipitakaparskip}
\csromancenter{มหาจตฺตารีสกสุตฺต (117) โหติ สมฺมาสติ.\csromanspacer{} อิติยิเม ตโย ธมฺมา สมฺมาวาจํ อนุปริธาวนฺติ อนุปริวตฺตนฺติ.\csromanspacer{} เสยฺยถิทํ, สมฺมาทิฏฺฐิ สมฺมาวายาโม สมฺมาสติ.}
\proseitem{139}{\csromanfolio{119}ตตฺร ภิกฺขเว สมฺมาทิฏฺฐิ ปุพฺพงฺคมา โหติ.\csromanspacer{} กถญฺจ ภิกฺขเว สมฺมาทิฏฺฐิ ปุพฺพงฺคมา โหติ มิจฺฉากมฺมนฺตํ “มิจฺฉากมฺมนฺโต”ติ ปชานาติ, สมฺมากมฺมนฺตํ “สมฺมากมฺมนฺโต”ติ ปชานาติ, สาสฺส โหติ สมฺมาทิฏฺฐิ.\csromanspacer{} กตโม จ ภิกฺขเว มิจฺฉากมฺมนฺโต.\csromanspacer{} ปาณาติปาโต อทินฺนาทานํ กาเมสุมิจฺฉาจาโร.\csromanspacer{} อยํ ภิกฺขเว มิจฺฉากมฺมนฺโต.\csromanspacer{} กตโม จ ภิกฺขเว สมฺมากมฺมนฺโต.\csromanspacer{} สมฺมากมฺมํนฺตํปหํ ภิกฺขเว ทฺวายํ วทามิ\csromandash{}อตฺถิ ภิกฺขเว สมฺมากมฺมนฺโต สาสโว ปุญฺญภาคิโย อุปธิเวปกฺโก, อตฺถิ ภิกฺขเว สมฺมากมฺมนฺโต อริโย อนาสโว โลกุตฺตโร มคฺคงฺโค.\csromanspacer{} กตโม จ ภิกฺขเว สมฺมากมฺมนฺโต สาสโว ปุญฺญภาคิโย อุปธิเวปกฺโก.\csromanspacer{} ปาณาติปาตา เวรมณี อทินฺนาทานา เวรมณี กาเมสุมิจฺฉาจารา เวรมณี.\csromanspacer{} อยํ ภิกฺขเว สมฺมากมฺมนฺโต สาสโว ปุญฺญภาคิโย อุปธิเวปกฺโก.\csromanspacer{} กตโม จ ภิกฺขเว สมฺมากมฺมนฺโต อริโย อนาสโว โลกุตฺตโร มคฺคงฺโค.\csromanspacer{} ยา โข ภิกฺขเว อริยจิตฺตสฺส อนาสวจิตฺตสฺส อริยมคฺคสมงฺคิโน อริยมคฺคํ ภาวยโต ตีหิ กายทุจฺจริเตหิ อารติ วิรติ ปฏิวิรติ เวรมณี.\csromanspacer{} อยํ ภิกฺขเว สมฺมากมฺมนฺโต อริโย อนาสโว โลกุตฺตโร มคฺคงฺโค.\csromanspacer{} โส มิจฺฉากมฺมนฺตสฺส ปหานาย วายมติ สมฺมากมฺมนฺตสฺส อุปสมฺปทาย, สฺวาสฺส โหติ สมฺมาวายาโม.\csromanspacer{} โส สโต มิจฺฉากมฺมนฺตํ ปชหติ, สโต สมฺมากมฺมนฺตํ อุปสมฺปชฺช วิหรติ, สาสฺส โหติ สมฺมาสติ.\csromanspacer{} อิติยิเม ตโย ธมฺมา สมฺมากมฺมนฺตํ อนุปริธาวนฺติ อนุปริวตฺตนฺติ.\csromanspacer{} เสยฺยถิทํ, สมฺมาทิฏฺฐิ สมฺมาวายาโม สมฺมาสติ.}

\proseitem{140}{ตตฺร ภิกฺขเว สมฺมาทิฏฺฐิ ปุพฺพงฺคมา โหติ.\csromanspacer{} กถญฺจ ภิกฺขเว สมฺมาทิฏฺฐิ ปุพฺพงฺคมา โหติ มิจฺฉา-อาชีวํ “มิจฺฉา-อาชีโว”ติ ปชานาติ, สมฺมา-อาชีวํ “สมฺมา-อาชีโว”ติ ปชานาติ, สาสฺส โหติ สมฺมาทิฏฺฐิ.\csromanspacer{} กตโม จ ภิกฺขเว มิจฺฉา-อาชีโว.\csromanspacer{} กุหนา ลปนา เนมิตฺติกตา นิปฺเปสิกตา ลาเภน ลาภํ นิชิคีสนตา\footnote{นิชิคิํสนตา (สี, สฺยา, กํ, อิ)}.\csromanspacer{} อยํ ภิกฺขเว มิจฺฉา-อาชีโว.\csromanspacer{} กตโม จ ภิกฺขเว สมฺมา-อาชีโว.\csromanspacer{} สมฺมา-อาชีวํปหํ \csromanfolio{120}ภิกฺขเว ทฺวายํ วทามิ\csromandash{}อตฺถิ ภิกฺขเว สมฺมา-อาชีโว สาสโว ปุญฺญภาคิโย อุปธิเวปกฺโก, อตฺถิ ภิกฺขเว สมฺมา-อาชีโว อริโย อนาสโว โลกุตฺตโร มคฺคงฺโค.\csromanspacer{} กตโม จ ภิกฺขเว สมฺมา-อาชีโว สาสโว ปุญฺญภาคิโย อุปธิเวปกฺโก.\csromanspacer{} อิธ ภิกฺขเว อริยสาวโก มิจฺฉา-อาชีวํ ปหาย สมฺมา-อาชีเวน ชีวิกํ กปฺเปติ.\csromanspacer{} อยํ ภิกฺขเว สมฺมา-อาชีโว สาสโว ปุญฺญภาคิโย อุปธิเวปกฺโก.\csromanspacer{} กตโม จ ภิกฺขเว สมฺมา-อาชีโว อริโย อนาสโว โลกุตฺตโร มคฺคงฺโค.\csromanspacer{} ยา โข ภิกฺขเว อริยจิตฺตสฺส อนาสวจิตฺตสฺส อริยมคฺคสมงฺคิโน อริยมคฺคํ ภาวยโต มิจฺฉา-อาชีวา อารติ วิรติ ปฏิวิรติ เวรมณี.\csromanspacer{} อยํ ภิกฺขเว สมฺมา-อาชีโว อริโย อนาสโว โลกุตฺตโร มคฺคงฺโค.\csromanspacer{} โส มิจฺฉา-อาชีวสฺส ปหานาย วายมติ สมฺมา-อาชีวสฺส อุปสมฺปทาย, สฺวาสฺส โหติ สมฺมาวายาโม.\csromanspacer{} โส สโต มิจฺฉา-อาชีวํ ปชหติ, สโต สมฺมา-อาชีวํ อุปสมฺปชฺช วิหรติ, สาสฺส โหติ สมฺมาสติ.\csromanspacer{} อิติยิเม ตโย ธมฺมา สมฺมา-อาชีวํ อนุปริธาวนฺติ อนุปริวตฺตนฺติ.\csromanspacer{} เสยฺยถิทํ, สมฺมาทิฏฺฐิ สมฺมาวายาโม สมฺมาสติ.}

\proseitem{141}{ตตฺร ภิกฺขเว สมฺมาทิฏฺฐิ ปุพฺพงฺคมา โหติ.\csromanspacer{} กถญฺจ ภิกฺขเว สมฺมาทิฏฺฐิ ปุพฺพงฺคมา โหติ.\csromanspacer{} สมฺมาทิฏฺฐิสฺส ภิกฺขเว สมฺมาสงฺกปฺโป ปโหติ, สมฺมาสงฺกปฺปสฺส สมฺมาวาจา ปโหติ, สมฺมาวาจสฺส สมฺมากมฺมนฺโต ปโหติ, สมฺมากมฺมนฺตสฺส สมฺมาชีโว ปโหติ, สมฺมา-อาชีวสฺส สมฺมาวายาโม ปโหติ, สมฺมาวายามสฺส สมฺมาสติ ปโหติ, สมฺมาสติสฺส สมฺมาสมาธิ ปโหติ, สมฺมาสมาธิสฺส สมฺมาญาณํ ปโหติ, สมฺมาญาณสฺส สมฺมาวิมุตฺติ ปโหติ.\csromanspacer{} อิติ โข ภิกฺขเว อฏฺฐงฺคสมนฺนาคโต เสกฺโข\footnote{อฏฺฐงฺคสมนฺนาคตา เสขา ปฏิปทา (สี), อฏฺฐงฺคสมนฺนาคโต เสโข ปาฏิปโท (อิ, ก) ( ) นตฺถิ สี-สฺยา-กํ-อิ-โปตฺถเกสุ.} ทสงฺคมนฺนาคโต อรหา โหติ. (ตตฺรปิ สมฺมาญาเณน อเนเก ปาปกา อกุสลา ธมฺมา วิคตา ภาวนาปาริปูริํ คจฺฉนฺติ.)}

\proseitem{142}{ตตฺร ภิกฺขเว สมฺมาทิฏฺฐิ ปุพฺพงฺคมา โหติ.\csromanspacer{} กถญฺจ ภิกฺขเว สมฺมาทิฏฺฐิ ปุพฺพงฺคมา โหติ.\csromanspacer{} สมฺมาทิฏฺฐิสฺส ภิกฺขเว มิจฺฉาทิฏฺฐิ นิชฺชิณฺณา โหติ, เย จ มิจฺฉาทิฏฺฐิปจฺจยา อเนเก ปาปกา อกุสลา ธมฺมา สมฺภวนฺติ,}

\vspace{\tipitakaparskip}
\csromancenter{มหาจตฺตารีสกสุตฺต (117) เต จสฺส นิชฺชิณฺณา โหนฺติ.\csromanspacer{} สมฺมาทิฏฺฐิปจฺจยา อเนเก กุสลา ธมฺมา ภาวนาปาริปูริํ คจฺฉนฺติ, สมฺมาสงฺกปฺปสฺส ภิกฺขเว มิจฺฉาสงฺกปฺโป นิชฺชิณฺโณ โหติ ฯเปฯ.\csromanspacer{} สมฺมาวาจสฺส ภิกฺขเว มิจฺฉาวาจา นิชิณฺณา โหติ.\csromanspacer{} สมฺมากมฺมนฺตสฺส ภิกฺขเว มิจฺฉากมฺมนฺโต นิชฺชิณฺโณ โหติ.\csromanspacer{} สมฺมา-อาชีวสฺส ภิกฺขเว มิจฺฉา-อาชีโว นิชฺชิณฺโณ โหติ.\csromanspacer{} สมฺมาวายามสฺส ภิกฺขเว มิจฺฉาวายาโม นิชฺชิณฺโณ โหติ.\csromanspacer{} สมฺมาสติสฺส ภิกฺขเว มิจฺฉาสติ นิชฺชิณฺณา โหติ.\csromanspacer{} สมฺมาสมาธิสฺส ภิกฺขเว มิจฺฉาสมาธิ นิชฺชิณฺโณ โหติ.\csromanspacer{} สมฺมาญาณสฺส ภิกฺขเว มิจฺฉาญาณํ นิชฺชิณฺณํ โหติ.\csromanspacer{} สมฺมาวิมุตฺตสฺส ภิกฺขเว มิจฺฉาวิมุตฺติ นิชฺชิณฺณา โหติ.\csromanspacer{} เย จ มิจฺฉาวิมุตฺติปจฺจยา อเนเก ปาปกา อกุสลา ธมฺมา สมฺภวนฺติ, เต จสฺส นิชฺชิณฺณา โหนฺติ.\csromanspacer{} สมฺมาวิมุตฺติปจฺจยา จ อเนเก กุสลา ธมฺมา ภาวนาปาริปูริํ คจฺฉนฺติ.}
\prose{\csromanfolio{121}อิติ โข ภิกฺขเว วีสติ กุสลปกฺขา วีสติ อกุสลปกฺขา มหาจตฺตารีสโก ธมฺมปริยาโย ปวตฺติโต อปฺปฏิวตฺติโย สมเณน วา พฺราหฺมเณน วา เทเวน วา มาเรน วา พฺรหฺมุนา วา เกนจิ วา โลกสฺมิํ.}

\proseitem{143}{โย หิ โกจิ ภิกฺขเว สมโณ วา พฺราหฺมโณ วา อิมํ มหาจตฺตารีสกํ ธมฺมปริยายํ ครหิตพฺพํ ปฏิกฺโกสิตพฺพํ มญฺเญยฺย, ตสฺส ทิฏฺเฐว ธมฺเม ทสสหธมฺมิกา วาทานุวาทา คารยฺหํ ฐานํ อาคจฺฉนฺติ.\csromanspacer{} สมฺมาทิฏฺฐิํ เจ ภวํ ครหติ, เย จ มิจฺฉาทิฏฺฐี สมณพฺราหฺมณา, เต โภโต ปุชฺชา, เต โภโต ปาสํสา.\csromanspacer{} สมฺมาสงฺกปฺปํ เจ ภวํ ครหติ, เย จ มิจฺฉาสงฺกปฺปา สมณพฺราหฺมณา, เต โภโต ปุชฺชา, เต โภโต ปาสํสา.\csromanspacer{} สมฺมาวาจํ เจ ภวํ ครหติ ฯเปฯ.\csromanspacer{} สมฺมากมฺมนฺตํ เจ ภวํ ครหติ.\csromanspacer{} สมฺมา- อาชีวํ เจ ภวํ ครหติ.\csromanspacer{} สมฺมาวายามํ เจ ภวํ ครหติ.\csromanspacer{} สมฺมาสติํ เจ ภวํ ครหติ.\csromanspacer{} สมฺมาสมาธิํ เจ ภวํ ครหติ.\csromanspacer{} สมฺมาญาณํ เจ ภวํ ครหติ.\csromanspacer{} สมฺมาวิมุตฺติํ เจ ภวํ ครหติ, เย จ มิจฺฉาวิมุตฺตี สมณพฺราหฺมณา, เต โภโต ปุชฺชา, เต โภโต ปาสํสา.\csromanspacer{} โย โกจิ ภิกฺขเว สมโณ วา พฺราหฺมโณ วา อิมํ มหาจตฺตารีสกํ ธมฺมปริยายํ ครหิตพฺพํ ปฏิกฺโกสิตพฺพํ มญฺเญยฺย, ตสฺส ทิฏฺเฐว ธมฺเม อิเม ทสสหธมฺมิกา วาทานุวาทา คารยฺหํ ฐานํ อาคจฺฉนฺติ.\csromanspacer{} เยปิ เต ภิกฺขเว อเหสุํ โอกฺกลา วสฺสภญฺญา\footnote{วยภิญฺญา (ก) สํ 2. 60; อํ 1. 340 ปิฏฺเฐสุ ปสฺสิตพฺพํ.} อเหตุวาทา อกิริยวาทา นตฺถิกวาทา เตปิ มหาจตฺตารีสกํ ธมฺมปริยายํ น \csromanfolio{122}ครหิตพฺพํ นปฏิกฺโกสิตพฺพํ อมญฺญิํสุ\footnote{มญฺเญยฺยุํ (ก)}.\csromanspacer{} ตํ กิสฺส เหตุ, นินฺทาพฺยาโรส-อุปารมฺภภยาติ.}

\prose{อิทมโวจ ภควา.\csromanspacer{} อตฺตมนา เต ภิกฺขู ภควโต ภาสิตํ อภินนฺทุนฺติ.}

\nitthitamruled{มหาจตฺตารีสกสุตฺตํ นิฏฺฐิตํ สตฺตมํ.}
\csromanmatikaanchor{mtk.21}
\csromantocmarkref{subsection}{8. อานาปานสฺสติสุตฺต}{mtk.21}
\bookmark[dest={mtk.21},level=2]{8. อานาปานสฺสติสุตฺต}
\csromanheader{2}{8. อานาปานสฺสติสุตฺต}
\proseitem{144}{เอวํ เม สุตํ\csromandash{}เอกํ สมยํ ภควา สาวตฺถิยํ วิหรติ ปุพฺพาราเม มิคารมาตุปาสาเท สมฺพหุเลหิ อภิญฺญาเตหิ อภิญฺญาเตหิ เถเรหิ สาวเกหิ สทฺธิํ, อายสฺมตา จ สาริปุตฺเตน อายสฺมตา จ มหาโมคฺคลฺลาเนน\footnote{มหาโมคฺคลาเนน (ก)} อายสฺมตา จ มหากสฺสเปน อายสฺมตา จ มหากจฺจายเนน อายสฺมตา จ มหาโกฏฺฐิเกน อายสฺมตา จ มหากปฺปิเนน อายสฺมตา จ มหาจุนฺเทน อายสฺมตา จ อนุรุทฺเธน อายสฺมตา จ เรวเตน อายสฺมตา จ อานนฺเทน อญฺเญหิ จ อภิญฺญาเตหิ อภิญฺญาเตหิ เถเรหิ สาวเกหิ สทฺธิํ.}

\prose{เตน โข ปน สมเยน เถรา ภิกฺขู นเว ภิกฺขู โอวทนฺติ อนุสาสนฺติ, อปฺเปกจฺเจ เถรา ภิกฺขู ทสปิ ภิกฺขู โอวทนฺติ อนุสาสนฺติ, อปฺเปกจฺเจ เถรา ภิกฺขู วีสมฺปิ ภิกฺขู โอวทนฺติ อนุสาสนฺติ, อปฺเปกจฺเจ เถรา ภิกฺขู ติํสมฺปิ ภิกฺขู โอวทนฺติ อนุสาสนฺติ, อปฺเปกจฺเจ เถรา ภิกฺขู จตฺตารีสมฺปิ ภิกฺขู โอวทนฺติ อนุสาสนฺติ.\csromanspacer{} เต จ นวา ภิกฺขู เถเรหิ ภิกฺขูหิ โอวทิยมานา อนุสาสิยมานา อุฬารํ ปุพฺเพนาปรํ วิเสสํ ชานนฺติ\footnote{ปชานนฺติ (สฺยา, กํ), สญฺชานนฺติ (ก)}.}

\proseitem{145}{เตน โข ปน สมเยน ภควา ตทหุโปสเถ ปนฺนรเส ปวารณาย ปุณฺณาย ปุณฺณมาย รตฺติยา ภิกฺขุสํฆปริวุโต อพฺโภกาเส นิสินฺโน โหติ.\csromanspacer{} อถ โข ภควา ตุณฺหีภูตํ ตุณฺหีภูตํ ภิกฺขุสํฆํ อนุวิโลเกตฺวา ภิกฺขู อามนฺเตสิ “อารทฺโธสฺมิ ภิกฺขเว อิมาย ปฏิปทาย, อารทฺธจิตฺโตสฺมิ ภิกฺขเว อิมาย ปฏิปทาย.\csromanspacer{} ตสฺมาติห ภิกฺขเว ภิยฺโยโส มตฺตาย วีริยํ อารภถ อปฺปตฺตสฺส ปตฺติยา อนธิคตสฺส}

\vspace{\tipitakaparskip}
\csromancenter{อานาปานสฺสติสุตฺต (118) อธิคมาย อสจฺฉิกตสฺส สจฺฉิกิริยาย.\csromanspacer{} อิเธวาหํ สาวตฺถิยํ โกมุทิํ จาตุมาสินิํ อาคเมสฺสามี”ติ.\csromanspacer{} อสฺโสสุํ โข ชานปทา ภิกฺขู “ภควา กิร ตตฺเถว สาวตฺถิยํ โกมุทิํ จาตุมาสินิํ อาคเมสฺสตี”ติ.\csromanspacer{} เต ชานปทา ภิกฺขู สาวตฺถิํ\footnote{สาวตฺถิยํ (สฺยา, กํ, อิ, ก)} โอสรนฺติ ภควนฺตํ ทสฺสนาย.\csromanspacer{} เต จ โข เถรา ภิกฺขู ภิยฺโยโส มตฺตาย นเว ภิกฺขู โอวทนฺติ อนุสาสนฺติ, อปฺเปกจฺเจ เถรา ภิกฺขู ทสปิ ภิกฺขู โอวทนฺติ อนุสาสนฺติ, อปฺเปกจฺเจ เถรา ภิกฺขู วีสมฺปิ ภิกฺขู โอวทนฺติ อนุสาสนฺติ, อปฺเปกจฺเจ เถรา ภิกฺขู ติํสมฺปิ ภิกฺขู โอวทนฺติ อนุสาสนฺติ, อปฺเปกจฺเจ เถรา ภิกฺขู จตฺตารีสมฺปิ ภิกฺขู โอวทนฺติ อนุสาสนฺติ.\csromanspacer{} เต จ นวา ภิกฺขู เถเรหิ ภิกฺขูหิ โอวทิยมานา อนุสาสิยมานา อุฬารํ ปุพฺเพนาปรํ วิเสสํ ชานนฺติ.}
\proseitem{146}{\csromanfolio{123}เตน โข ปน สมเยน ภควา ตทหุโปสเถ ปนฺนรเส โกมุทิยา จาตุมาสินิยา ปุณฺณาย ปุณฺณมาย รตฺติยา ภิกฺขุสํฆปริวุโต อพฺโภกาเส นิสินฺโน โหติ.\csromanspacer{} อถ โข ภควา ตุณฺหีภูตํ ตุณฺหีภูตํ ภิกฺขุสํฆํ อนุวิโลเกตฺวา ภิกฺขู อามนฺเตสิ “อปลาปายํ ภิกฺขเว ปริสา, นิปฺปลาปายํ ภิกฺขเว ปริสา สุทฺธา สาเร\footnote{สุทฺธสาเร ปติฏฺฐิตา (สฺยา, กํ, อิ)} ปติฏฺฐิตา.\csromanspacer{} ตถารูโป อยํ ภิกฺขเว ภิกฺขุสํโฆ.\csromanspacer{} ตถรูปา อยํ ภิกฺขเว ปริสา, ยถารูปา ปริสา อาหุเนยฺยา ปาหุเนยฺยา ทิกฺขิเณยฺยา อญฺชลิกรณียา อนุตฺตรํ ปุญฺญกฺเขตฺตํ โลกสฺส.\csromanspacer{} ตถารูโป อยํ ภิกฺขเว ภิกฺขุสํโฆ, ตถารูปา อยํ ภิกฺขเว ปริสา, ยถารูปาย ปริสาย อปฺปํ ทินฺนํ พหุ โหติ, พหุ ทินฺนํ พหุตรํ.\csromanspacer{} ตถารูโป อยํ ภิกฺขเว ภิกฺขุสํโฆ, ตถารูปา อยํ ภิกฺขเว ปริสา, ยถารูปา ปริสา ทุลฺลภา ทสฺสนาย โลกสฺส.\csromanspacer{} ตถารูโป อยํ ภิกฺขเว ภิกฺขุสํโฆ, ตถารูปา อยํ ภิกฺขเว ปริสา, ยถารูปํ ปริสํ อลํ โยชนคณนานิ ทสฺสนาย คนฺตุํ ปุโฏเสนาปิ\footnote{ปุโฏเสนาปิ, ตถารูโป อยํ ภิกฺขเว ภิกฺขุสํโฆ, ตถารูปา อยํ ปริสา (สี, อิ, ก)}.}

\proseitem{147}{สนฺติ ภิกฺขเว ภิกฺขู อิมสฺมิํ ภิกฺขุสํเฆ อรหนฺโต ขีณาสวา วุสิตวนฺโต กตกรณียา โอหิตภารา อนุปฺปตฺตสทตฺถา ปริกฺขีณภวสํโยชนา สมฺมทญฺญาวิมุตฺตา, เอวรูปาปิ ภิกฺขเว สนฺติ ภิกฺขู อิมสฺมิํ \csromanfolio{124}ภิกฺขุสํเฆ.\csromanspacer{} สนฺติ ภิกฺขเว ภิกฺขู อิมสฺมิํ ภิกฺขุสํเฆ ปญฺจนฺนํ โอรมฺภาคิยานํ สํโยชนานํ ปริกฺขยา โอปปาติกา ตตฺถ ปรินิพฺพายิโน อนาวตฺติธมฺมา ตสฺมา โลกา, เอวรูปาปิ ภิกฺขเว สนฺติ ภิกฺขู อิมสฺมิํ ภิกฺขุสํเฆ.\csromanspacer{} สนฺติ ภิกฺขเว ภิกฺขู อิมสฺมิํ ภิกฺขุสํเฆ ติณฺณํ สํโยชนานํ ปริกฺขยา ราคโทสโมหานํ ตนุตฺตา สกทาคามิโน สกิเทว\footnote{สกิํเทว (ก)} อิมํ โลกํ อาคนฺตฺวา ทุกฺขสฺสนฺตํ กริสฺสนฺติ, เอวรูปาปิ ภิกฺขเว สนฺติ ภิกฺขู อิมสฺมิํ ภิกฺขุสํเฆ.\csromanspacer{} สนฺติ ภิกฺขเว ภิกฺขู อิมสฺมิํ ภิกฺขูสํเฆ ติณฺณํ สํโยชนานํ ปริกฺขยา โสตาปนฺนา อวินิปาตธมฺมา นิยตา สมฺโพธิปรายนา, เอวรูปาปิ ภิกฺขเว สนฺติ ภิกฺขู อิมสฺมิํ ภิกฺขุสํเฆ. สนฺติ ภิกฺขเว ภิกฺขู อิมสฺมิํ ภิกฺขุสํเฆ จตุนฺนํ สติปฏฺฐานานํ ภาวนานุโยคมนุยุตฺตา วิหรนฺติ, เอวรูปาปิ ภิกฺขเว สนฺติ ภิกฺขู อิมสฺมิํ ภิกฺขุสํเฆ.\csromanspacer{} สนฺติ ภิกฺขเว ภิกฺขู อิมสฺมิํ ภิกฺขุสํเฆ จตุนฺนํ สมฺมปฺปธานานํ ภาวนานุโยคมนุยุตฺตา วิหรนฺติ ฯเปฯ จตุนฺนํ อิทฺธิปาทานํ.\csromanspacer{} ปญฺจนฺนํ อินฺทฺริยานํ.\csromanspacer{} ปญฺจนฺนํ พลานํ.\csromanspacer{} สตฺตนฺนํ โพชฺฌงฺคานํ.\csromanspacer{} อริยสฺส อฏฺฐงฺคิกสฺส มคฺคสฺส ภาวนานุโยคมนุยุตฺตา วิหรนฺติ, เอวรูปาปิ ภิกฺขเว สนฺติ ภิกฺขู อิมสฺมิํ ภิกฺขุสํเฆ.\csromanspacer{} สนฺติ ภิกฺขเว ภิกฺขู อิมสฺมิํ ภิกฺขุสํเฆ เมตฺตาภาวนานุโยคมนุยุตฺตา วิหรนฺติ.\csromanspacer{} กรุณาภาวนานุโยคมนุยุตฺตา วิหรนฺติ.\csromanspacer{} มุทิตาภาวนานุโยคมนุยุตฺตา วิหรนฺติ.\csromanspacer{} อุเปกฺขาภาวนานุโยคมนุยุตฺตา วิหรนฺติ.\csromanspacer{} อสุภภาวนานุโยคมนุยุตฺตา วิหรนฺติ.\csromanspacer{} อนิจฺจสญฺญาภาวนานุโยคมนุยุตฺตา วิหรนฺติ, เอวรูปาปิ ภิกฺขเว สนฺติ ภิกฺขู อิมสฺมิํ ภิกฺขุสํเฆ.\csromanspacer{} สนฺติ ภิกฺขเว ภิกฺขู อิมสฺมิํ ภิกฺขุสํเฆ อานาปานสฺสติภาวนานุโยคมนุยุตฺตา วิหรนฺติ, อานาปานสฺสติ ภิกฺขเว ภาวิตา พหุลีกตา มหปฺผลา โหติ มหานิสํสา.\csromanspacer{} อานาปานสฺสติ ภิกฺขเว ภาวิตา พหุลีกตา จตฺตาโร สติปฏฺฐาเน ปริปูเรติ.\csromanspacer{} จตฺตาโร สติปฏฺฐานา ภาวิตา พหุลีกตา สตฺต โพชฺฌงฺเค ปริปูเรนฺติ.\csromanspacer{} สตฺต โพชฺฌงฺคา ภาวิตา พหุลีกตา วิชฺชาวิมุตฺติํ ปริปูเรนฺติ.}

\proseitem{148}{กถํ ภาวิตา จ ภิกฺขเว อานาปานสฺสติ กถํ พหุลีกตา มหปฺผลา โหติ มหานิสํสา.\csromanspacer{} อิธ ภิกฺขเว ภิกฺขุ อรญฺญคโต วา รุกฺขมูลคโต วา สุญฺญาคารคโต วา นิสีทติ ปลฺลงฺกํ อาภุชิตฺวา}

\vspace{\tipitakaparskip}
\csromancenter{อานาปานสฺสติสุตฺต (118) อุชุํ กายํ ปณิธาย ปริมุขํ สติํ อุปฏฺฐเปตฺวา.\csromanspacer{} โส สโตว อสฺสสติ, สโตว\footnote{สโต (สี, สฺยา, กํ, อิ)} ปสฺสสติ.}
\prose{\csromanfolio{125}ทีฆํ วา อสฺสสนฺโต “ทีฆํ อสฺสสามี”ติ ปชานาติ, ทีฆํ วา ปสฺสสนฺโต “ทีฆํ ปสฺสสามี”ติ ปชานาติ.\csromanspacer{} รสฺสํ วา อสฺสสนฺโต “รสฺสํ อสฺสสามี”ติ ปชานาติ, รสฺสํ วา ปสฺสสนฺโต “รสฺสํ ปสฺสสามี”ติ ปชานาติ. “สพฺพกายปฏิสํเวที อสฺสสิสฺสามี”ติ สิกฺขติ, “สพฺพกายปฏิสํเวที ปสฺสสิสฺสามี”ติ สิกฺขติ. “ปสฺสมฺภยํ กายสงฺขารํ อสฺสสิสฺสามี”ติ สิกฺขติ, “ปสฺสมฺภยํ กายสงฺขารํ ปสฺสสิสฺสามี”ติ สิกฺขติ.}

\prose{“ปีติปฏิสํเวที อสฺสสิสฺสามี”ติ สิกฺขติ, “ปีติปฏิสํเวที ปสฺสสิสฺสามี”ติ สิกฺขติ. “สุขปฏิสํเวที อสฺสสิสฺสามี”ติ สิกฺขติ, “สุขปฏิสํเวที ปสฺสสิสฺสามี”ติ สิกฺขติ. “จิตฺตสงฺขารปฏิสํเวที อสฺสสิสฺสามี”ติ สิกฺขติ, “จิตฺตสงฺขารปฏิสํเวที ปสฺสสิสฺสามี”ติ สิกฺขติ. “ปสฺสมฺภยํ จิตฺตสงฺขารํ อสฺสสิสฺสามี”ติ สิกฺขติ, “ปสฺสมฺภยํ จิตฺตสงฺขารํ ปสฺสสิสฺสามี”ติ สิกฺขติ.}

\prose{“จิตฺตปฏิสํเวที อสฺสสิสฺสามี”ติ สิกฺขติ, “จิตฺตปฏิสํเวที ปสฺสสิสฺสามี”ติ สิกฺขติ. “อภิปฺปโมทยํ จิตฺตํ อสฺสสิสฺสามี”ติ สิกฺขติ, อภิปฺปโมทยํ จิตฺตํ ปสฺสสิสฺสามี”ติ สิกฺขติ. “สมาทหํ จิตฺตํ อสฺสสิสฺสามี”ติ สิกฺขติ, “สมาทหํ จิตฺตํ ปสฺสสิสฺสามี”ติ สิกฺขติ. “วิโมจยํ จิตฺตํ อสฺสสิสฺสามี”ติ สิกฺขติ, “วิโมจยํ จิตฺตํ ปสฺสสิสฺสามี”ติ สิกฺขติ.}

\prose{“อนิจฺจานุปสฺสี อสฺสสิสฺสามี”ติ สิกฺขติ, “อนิจฺจานุปสฺสี ปสฺสสิสฺสามี”ติ สิกฺขติ. “วิราคานุปสฺสี อสฺสสิสฺสามี”ติ สิกฺขติ, “วิราคานุปสฺสี ปสฺสสิสฺสามี”ติ สิกฺขติ. “นิโรธานุปสฺสี อสฺสสิสฺสามี”ติ สิกฺขติ, “นิโรธานุปสฺสี ปสฺสสิสฺสามี”ติ สิกฺขติ. “ปฏินิสฺสคฺคานุปสฺสี อสฺสสิสฺสามี”ติ สิกฺขติ, “ปฏินิสฺสคฺคานุปสฺสี ปสฺสสิสฺสามี”ติ สิกฺขติ.\csromanspacer{} เอวํ ภาวิตา โข ภิกฺขเว อานาปานสฺสติ เอวํ พหุลีกตา มหปฺผลา โหติ มหานิสํสา.}

\proseitem{149}{\csromanfolio{126}กถํ ภาวิตา จ ภิกฺขเว อานาปานสฺสติ กถํ พหุลีกตา จตฺตาโร สติปฏฺฐาเน ปริปูเรติ.\csromanspacer{} ยสฺมิํ สมเย ภิกฺขเว ภิกฺขุ ทีฆํ วา อสฺสสนฺโต “ทีฆํ ปสฺสสามี”ติ ปชานาติ, ทีฆํ วา ปสฺสสนฺโต “ทีฆํ อสฺสสามี”ติ ปชานาติ.\csromanspacer{} รสฺสํ วา อสฺสสนฺโต “รสฺสํ อสฺสสามี”ติ ปชานาติ, รสฺสํ วา ปสฺสสนฺโต “รสฺสํ ปสฺสสามี”ติ ปชานาติ. “สพฺพกายปฏิสํเวที อสฺสสิสฺสามี”ติ สิกฺขติ, “สพฺพกายปฏิสํเวที ปสฺสสิสฺสามี”ติ สิกฺขติ. “ปสฺสมฺภยํ กายสงฺขารํ อสฺสสิสฺสามี”ติ สิกฺขติ, “ปสฺสมฺภยํ กายสงฺขารํ ปสฺสสิสฺสามี”ติ สิกฺขติ.\csromanspacer{} กาเย กายานุปสฺสี ภิกฺขเว ตสฺมิํ สมเย ภิกฺขุ วิหรติ อาตาปี สมฺปชาโน สติมา วิเนยฺย โลเก อภิชฺฌาโทมนสฺสํ.\csromanspacer{} กาเยสุ กายญฺญตราหํ ภิกฺขเว เอวํ วทามิ\csromandash{} ยทิทํ อสฺสาสปสฺสาสา.\csromanspacer{} ตสฺมาติห ภิกฺขเว กาเย กายานุปสฺสี ตสฺมิํ สมเย ภิกฺขุ วิหรติ อาตาปี สมฺปชาโน สติมา วิเนยฺย โลเก อภิชฺฌาโทมนสฺสํ. (1)}

\prose{ยสฺมิํ สมเย ภิกฺขเว ภิกฺขุ “ปีติปฏิสํเวที อสฺสสิสฺสามี”ติ สิกฺขติ, “ปีติปฏิสํเวที ปสฺสสิสฺสามี”ติ สิกฺขติ. “สุขปฏิสํเวที อสฺสสิสฺสามี”ติ สิกฺขติ, “สุขปฏิสํเวที ปสฺสสิสฺสามี”ติ สิกฺขติ. “จิตฺตสงฺขารปฏิสํเวที อสฺสสิสฺสามี”ติ สิกฺขติ, “จิตฺตสงฺขารปฏิสํเวที ปสฺสสิสฺสามี”ติ สิกฺขติ. “ปสฺสมฺภยํ จิตฺตสงฺขารํ อสฺสสิสฺสามี”ติ สิกฺขติ, “ปสฺสมฺภยํ จิตฺตสงฺขารํ อสฺสสิสฺสามี”ติ สิกฺขติ.\csromanspacer{} เวทนาสุ เวทนานุปสฺสี ภิกฺขเว ตสฺมิํ สมเย ภิกฺขุ วิหรติ อาตาปี สมฺปชาโน สติมา วิเนยฺย โลเก อภิชฺฌาโทมนสฺสํ.\csromanspacer{} เวทนาสุ เวทนาญฺญตราหํ ภิกฺขเว เอวํ วทามิ\csromandash{}ยทิทํ อสฺสาสปสฺสาสานํ สาธุกํ มนสิการํ, ตสฺมาติห ภิกฺขเว เวทนาสุ เวทนานุปสฺสี ตสฺมิํ สมเย ภิกฺขุ วิหรติ อาตาปี สมฺปชาโน สติมา วิเนยฺย โลเก อภิชฺฌาโทมนสฺสํ. (2)}

\prose{ยสฺมิํ สมเย ภิกฺขเว ภิกฺขุ “จิตฺตปฏิสํเวที อสฺสสิสฺสามี”ติ สิกฺขติ, “จิตฺตปฏิสํเวที ปสฺสสิสฺสามี”ติ สิกฺขติ. “อภิปฺปโมทยํ จิตฺตํ อสฺสสิสฺสามี”ติ สิกฺขติ, “อภิปฺปโมทยํ จิตฺตํ ปสฺสสิสฺสามี”ติ สิกฺขติ. “สมาทหํ จิตฺตํ อสฺสสิสฺสามี”ติ สิกฺขติ, “สมาทหํ จิตฺตํ ปสฺสสิสฺสามี”ติ สิกฺขติ. “วิโมจยํ จิตฺตํ อสฺสสิสฺสามี”ติ สิกฺขติ,}

\vspace{\tipitakaparskip}
\csromancenter{อานาปานสฺสติสุตฺต (118) “วิโมจยํ จิตฺตํ ปสฺสสิสฺสามี”ติ สิกฺขติ.\csromanspacer{} จิตฺเต จิตฺตานุปสฺสี ภิกฺขเว ตสฺมิํ สมเย ภิกฺขุ วิหรติ อาตาปี สมฺปชาโน สติมา วิเนยฺย โลเก อภิชฺฌาโทมนสฺสํ.\csromanspacer{} นาหํ ภิกฺขเว มุฏฺฐสฺสติสฺส อสมฺปชานสฺส อานาปานสฺสติํ วทามิ, ตสฺมาติห ภิกฺขเว จิตฺเต จิตฺตานุปสฺสี ตสฺมิํ สมเย ภิกฺขุ วิหรติ อาตาปี สมฺปชาโน สติมา วิเนยฺย โลเก อภิชฺฌาโทมนสฺสํ. (3)}
\prose{\csromanfolio{127}ยสฺมิํ สมเย ภิกฺขเว ภิกฺขุ “อนิจฺจานุปสฺสี อสฺสสิสฺสามี”ติ สิกฺขติ, “อนิจฺจานุปสฺสี ปสฺสสิสฺสามี”ติ สิกฺขติ. “วิราคานุปสฺสี อสฺสสิสฺสามี”ติ สิกฺขติ, “วิราคานุปสฺสี ปสฺสสิสฺสามี”ติ สิกฺขติ. “นิโรธานุปสฺสี อสฺสสิสฺสามี”ติ สิกฺขติ, “นิโรธานุปสฺสี ปสฺสสิสฺสามี”ติ สิกฺขติ. “ปฏินิสฺสคฺคานุปสฺสี อสฺสสิสฺสามี”ติ สิกฺขติ, “ปฏินิสฺสคฺคานุปสฺสี ปสฺสสิสฺสามี”ติ สิกฺขติ.\csromanspacer{} ธมฺเมสุ ธมฺมานุปสฺสี ภิกฺขเว ตสฺมิํ สมเย ภิกฺขุ วิหรติ อาตาปี สมฺปชาโน สติมา วิเนยฺย โลเก อภิชฺฌาโทมนสฺสํ.\csromanspacer{} โส ยํ ตํ อภิชฺฌาโทมนสฺสานํ ปหานํ, ตํ ปญฺญาย ทิสฺวา สาธุกํ อชฺฌุเปกฺขิตา โหติ.\csromanspacer{} ตสฺมาติห ภิกฺขเว ธมฺเมสุ ธมฺมานุปสฺสี ตสฺมิํ สมเย ภิกฺขุ วิหรติ อาตาปี สมฺปชาโน สติมา วิเนยฺย โลเก อภิชฺฌาโทมนสฺสํ. (4)}

\prose{เอวํ ภาวิตา โข ภิกฺขเว อานาปานสฺสติ เอวํ พหุลีกตา จตฺตาโร สติปฏฺฐาเน ปริปูเรติ.}

\proseitem{150}{กถํ ภาวิตา จ ภิกฺขเว จตฺตาโร สติปฏฺฐานา กถํ พหุลีกตา สตฺต โพชฺฌงฺเค ปริปูเรนฺติ.\csromanspacer{} ยสฺมิํ สมเย ภิกฺขเว ภิกฺขุ กาเย กายานุปสฺสี วิหรติ อาตาปี สมฺปชาโน สติมา วิเนยฺย โลเก อภิชฺฌาโทมนสฺสํ.\csromanspacer{} อุปฏฺฐิตาสฺส ตสฺมิํ สมเย สติ โหติ อสมฺมุฏฺฐา\footnote{อปฺปมฺมุฏฺฐา (สฺยา, กํ)}.\csromanspacer{} ยสฺมิํ สมเย ภิกฺขเว ภิกฺขุโน อุปฏฺฐิตา สติ โหติ อสมฺมุฏฺฐา, สติสมฺโพชฺฌงฺโค ตสฺมิํ สมเย ภิกฺขุโน อารทฺโธ โหติ, สติสมฺโพชฺฌงฺคํ ตสฺมิํ สมเย ภิกฺขุ ภาเวติ, สติสมฺโพชฺฌงฺโค ตสฺมิํ สมเย ภิกฺขุโน ภาวนาปาริปูริํ คจฺฉติ. (1)}

\prose{โส ตถาสโต วิหรนฺโต ตํ ธมฺมํ ปญฺญาย ปวิจินติ ปวิจยติ\footnote{ปวิจรติ (สี, สฺยา, กํ, อิ)} ปริวีมํสํ อาปชฺชติ.\csromanspacer{} ยสฺมิํ สมเย ภิกฺขเว ภิกฺขุ ตถาสโต \csromanfolio{128}วิหรนฺโต ตํ ธมฺมํ ปญฺญาย ปวิจินติ ปวิจยติ ปริวีมํสํ อาปชฺชติ, ธมฺมวิจยสมฺโพชฺฌงฺโค ตสฺมิํ สมเย ภิกฺขุโน อารทฺโธ โหติ, ธมฺมวิจยสมฺโพชฺฌงฺคํ ตสฺมิํ สมเย ภิกฺขุ ภาเวติ, ธมฺมวิจยสมฺโพชฺฌงฺโค ตสฺมิํ สมเย ภิกฺขุโน ภาวนาปาริปูริํ คจฺฉติ. (2)}

\prose{ตสฺส ตํ ธมฺมํ ปญฺญาย ปวิจินโต ปวิจยโต ปริวีมํสํ อาปชฺชโต อารทฺธํ โหติ วีริยํ อสลฺลีนํ.\csromanspacer{} ยสฺมิํ สมเย ภิกฺขเว ภิกฺขุโน ตํ ธมฺมํ ปญฺญาย ปวิจินโต ปวิจยโต ปริวีมํสํ อาปชฺชโต อารทฺธํ โหติ วีริยํ อสลฺลีนํ, วีริยสมฺโพชฺฌงฺโค ตสฺมิํ สมเย ภิกฺขุโน อารทฺโธ โหติ, วีริยสมฺโพชฺฌงฺคํ ตสฺมิํ สมเย ภิกฺขุ ภาเวติ, วีริยสมฺโพชฺฌงฺโค ตสฺมิํ สมเย ภิกฺขุโน ภาวนาปาริปูริํ คจฺฉติ. (3)}

\prose{อารทฺธวีริยสฺส อุปฺปชฺชติ ปีติ นิรามิสา.\csromanspacer{} ยสฺมิํ สมเย ภิกฺขเว ภิกฺขุโน อารทฺธวีริยสฺส อุปฺปชฺชติ ปีติ นิรามิสา, ปีติสมฺโพชฺฌงฺโค ตสฺมิํ สมเย ภิกฺขุโน อารทฺโธ โหติ, ปีติสมฺโพชฺฌงฺคํ ตสฺมิํ สมเย ภิกฺขุ ภาเวติ, ปีติสมฺโพชฺฌงฺโค ตสฺมิํ สมเย ภิกฺขุโน ภาวนาปาริปูริํ คจฺฉติ. (4)}

\prose{ปีติมนสฺส กาโยปิ ปสฺสมฺภติ จิตฺตมฺปิ ปสฺสมฺภติ.\csromanspacer{} ยสฺมิํ สมเย ภิกฺขเว ภิกฺขุโน ปีติมนสฺส กาโยปิ ปสฺสมฺภติ จิตฺตมฺปิ ปสฺสมฺภติ, ปสฺสทฺธิสมฺโพชฺฌงฺโค ตสฺมิํ สมเย ภิกฺขุโน อารทฺโธ โหติ, ปสฺสทฺธิสมฺโพชฺฌงฺคํ ตสฺมิํ สมเย ภิกฺขุ ภาเวติ.\csromanspacer{} ปสฺสทฺธิสมฺโพชฺฌงฺโค ตสฺมิํ สมเย ภิกฺขุโน ภาวนาปาริปูริํ คจฺฉติ. (5)}

\prose{ปสฺสทฺธกายสฺส สุขิโน จิตฺตํ สมาธิยติ.\csromanspacer{} ยสฺมิํ สมเย ภิกฺขเว ภิกฺขุโน ปสฺสทฺธกายสฺส สุขิโน จิตฺตํ สมาธิยติ, สธิสมฺโพชฺฌงฺโค ตสฺมิํ สมเย ภิกฺขุโน อารทฺโธ โหติ, สมาธิสมฺโพชฺฌงฺคํ ตสฺมิํ สมเย ภิกฺขุ ภาเวติ, สมาธิสมฺโพชฺฌงฺโค ตสฺมิํ สมเย ภิกฺขุโน ภาวนาปาริปูริํ คจฺฉติ. (6)}

\prose{โส ตถาสมาหิตํ จิตฺตํ สาธุกํ อชฺฌุเปกฺขิตา โหติ.\csromanspacer{} ยสฺมิํ สมเย ภิกฺขเว ภิกฺขุ ตถาสมาหิตํ จิตฺตํ สาธุกํ อชฺฌุเปกฺขิตา โหติ.\csromanspacer{} อุเปกฺขาสมฺโพชฺฌงฺโค ตสฺมิํ สมเย ภิกฺขุโน อารทฺโธ โหติ, อุเปกฺขาสมฺโพชฺฌงฺคํ ตสฺมิํ สมเย ภิกฺขุ ภาเวติ, อุเปกฺขาสมฺโพชฺฌงฺโค ตสฺมิํ สมเย ภิกฺขุโน ภาวนาปาริปูริํ คจฺฉติ. (7)}

\csromanheader{2}{8. อานาปานสฺสติสุตฺต (118)}
\proseitem{151}{\csromanfolio{129}ยสฺมิํ สมเย ภิกฺขเว ภิกฺขุ เวทนาสุ ฯเปฯ จิตฺเต ฯเปฯ ธมฺเมสุ ธมฺมานุปสฺสี วิหรติ อาตาปี สมฺปชาโน สติมา วิเนยฺย โลเก อภิชฺฌาโทมนสฺสํ, อุปฏฺฐิตาสฺส ตสฺมิํ สมเย สติ โหติ อสมฺมุฏฺฐา.\csromanspacer{} ยสฺมิํ สมเย ภิกฺขเว ภิกฺขุโน อุปฏฺฐิตา สติ โหติ อสมฺมุฏฺฐา, สติสมฺโพชฺฌงฺโค ตสฺมิํ สมเย ภิกฺขุโน อารทฺโธ โหติ, สติสมฺโพชฺฌงฺคํ ตสฺมิํ สมเย ภิกฺขุ ภาเวติ, สติสมฺโพชฺฌงฺโค ตสฺมิํ สมเย ภิกฺขุโน ภาวนาปาริปูริํ คจฺฉติ. (1)}

\prose{โส ตถาสโต วิหรนฺโต ตํ ธมฺมํ ปญฺญาย ปวิจินติ ปวิจยติ ปริวีมํสํ อาปชฺชติ.\csromanspacer{} ยสฺมิํ สมเย ภิกฺขเว ภิกฺขุ ตถาสโต วิหรนฺโต ตํ ธมฺมํ ปญฺญาย ปวิจินติ ปวิจยติ ปริวีมํสํ อาปชฺชติ, ธมฺมวิจยสมฺโพชฺฌงฺโค ตสฺมิํ สมเย ภิกฺขุโน อารทฺโธ โหติ, ธมฺมวิจยสมฺโพชฺฌงฺคํ ตสฺมิํ สมเย ภิกฺขุ ภาเวติ, ธมฺมวิจยสมฺโพชฺฌงฺโค ตสฺมิํ สมเย ภิกฺขุโน ภาวนาปาริปูริํ คจฺฉติ. (2)}

\prose{ตสฺส ตํ ธมฺมํ ปญฺญาย ปวิจินโต ปวิจยโต ปริวีมํสํ อาปชฺชโต อารทฺธํ โหติ วีริยํ อสลฺลีนํ.\csromanspacer{} ยสฺมิํ สมเย ภิกฺขเว ภิกฺขุโน ตํ ธมฺมํ ปญฺญาย ปวิจินโต ปวิจยโต ปริวีมํสํ อาปชฺชโต อารทฺธํ โหติ วีริยํ อสลฺลีนํ, วีริยสมฺโพชฺฌงฺโค ตสฺมิํ สมเย ภิกฺขุโน อารทฺโธ โหติ, วีริยสมฺโพชฺฌงฺคํ ตสฺมิํ สมเย ภิกฺขุ ภาเวติ, วีริยสมฺโพชฺฌงฺโค ตสฺมิํ สมเย ภิกฺขุโน ภาวนาปาริปูริํ คจฺฉติ. (3)}

\prose{อารทฺธวีริยสฺส อุปฺปชฺชติ ปีติ นิรามิสา.\csromanspacer{} ยสฺมิํ สมเย ภิกฺขเว ภิกฺขุโน อารทฺธวีริยสฺส อุปฺปชฺชติ ปีติ นิรามิสา, ปีติสมฺโพชฺฌงฺโค ตสฺมิํ สมเย ภิกฺขุโน อารทฺโธ โหติ, ปีติสมฺโพชฺฌงฺคํ ตสฺมิํ สมเย ภิกฺขุ ภาเวติ, ปีติสมฺโพชฺฌงฺโค ตสฺมิํ สมเย ภิกฺขุโน ภาวนาปาริปูริํ คจฺฉติ. (4)}

\prose{ปีติมนสฺส กาโยปิ ปสฺสมฺภติ จิตฺตมฺปิ ปสฺสมฺภติ.\csromanspacer{} ยสฺมิํ สมเย ภิกฺขเว ภิกฺขุโน ปีติมนสฺส กาโยปิ ปสฺสมฺภติ จิตฺตมฺปิ ปสฺสมฺภติ, ปสฺสทฺธิสมฺโพชฺฌงฺโค ตสฺมิํ สมเย ภิกฺขุโน อารทฺโธ โหติ, ปสฺสทฺธิมฺโพชฺฌงฺคํ ตสฺมิํ สมเย ภิกฺขุ ภาเวติ, ปสฺสทฺธิสมฺโพชฺฌงฺโค ตสฺมิํ สมเย ภิกฺขุโน ภาวนาปาริปูริํ คจฺฉติ. (5)}

\prose{\csromanfolio{130}ปสฺสทฺธกายสฺส สุขิโน จิตฺตํ สมาธิยติ.\csromanspacer{} ยสฺมิํ สมเย ภิกฺขเว ภิกฺขุโน ปสฺสทฺธกายสฺส สุขิโน จิตฺตํ สมาธิยติ, สมาธิสมฺโพชฺฌงฺโค ตสฺมิํ สมเย ภิกฺขุโน อารทฺโธ โหติ, สมาธิสมฺโพชฺฌงฺคํ ตสฺมิํ สมเย ภิกฺขุ ภาเวติ, สมาธิสมฺโพชฺฌงฺโค ตสฺมิํ สมเย ภิกฺขุโน ภาวนาปาริปูริํ คจฺฉติ. (6)}

\prose{โส ตถาสมาหิตํ จิตฺตํ สาธุกํ อชฺฌุเปกฺขิตา โหติ.\csromanspacer{} ยสฺมิํ สมเย ภิกฺขเว ภิกฺขุ ตถาสมาหิตํ จิตฺตํ สาธุกํ อชฺฌุเปกฺขิตา โหติ, อุเปกฺขาสมฺโพชฺฌงฺโค ตสฺมิํ สมเย ภิกฺขุโน อารทฺโธ โหติ.\csromanspacer{} อุเปกฺขาสมฺโพชฺฌงฺคํ ตสฺมิํ สมเย ภิกฺขุ ภาเวติ, อุเปกฺขาสมฺโพชฺฌงฺโค ตสฺมิํ สมเย ภิกฺขุโน ภาวนาปาริปูริํ คจฺฉติ.\csromanspacer{} เอวํ ภาวิตา โข ภิกฺขเว จตฺตาโร สติปฏฺฐานา เอวํ พหุลีกตา สตฺต สมฺโพชฺฌงฺเค ปริปูเรนฺติ. (7)}

\proseitem{152}{กถํ ภาวิตา จ ภิกฺขเว สตฺต โพชฺฌงฺคา กถํ พหุลีกตา วิชฺชาวิมุตฺติํ ปริปูเรนฺติ.\csromanspacer{} อิธ ภิกฺขเว ภิกฺขุ สติสมฺโพชฺฌงฺคํ ภาเวติ วิเวกนิสฺสิตํ วิราคนิสฺสิตํ นิโรธนิสฺสิตํ โวสฺสคฺคปริณามิํ.\csromanspacer{} ธมฺมวิจยสมฺโพชฺฌงฺคํ ภาเวติ ฯเปฯ.\csromanspacer{} วีริยสมฺโพชฺฌงฺคํ ภาเวติ.\csromanspacer{} ปีติสมฺโพชฺฌงฺคํ ภาเวติ.\csromanspacer{} ปสฺสทฺธิสมฺโพชฺฌงฺคํ ภาเวติ.\csromanspacer{} สมาธิสมฺโพชฺฌงฺคํ ภาเวติ.\csromanspacer{} อุเปกฺขาสมฺโพชฺฌงฺคํ ภาเวติ วิเวกนิสฺสิตํ วิราคนิสฺสิตํ นิโรธนิสฺสิตํ โวสฺสคฺคปริณามิํ.\csromanspacer{} เอวํ ภาวิตา โข ภิกฺขเว สตฺต โพชฺฌงฺคา เอวํ พหุลีกตา วิชฺชาวิมุตฺติํ ปริปูเรนฺตีติ.}

\prose{อิทมโวจ ภควา.\csromanspacer{} อตฺตมนา เต ภิกฺขู ภควโต ภาสิตํ อภินนฺทุนฺติ.}

\nitthitamruled{อานาปานสฺสติสุตฺตํ นิฏฺฐิตํ อฏฺฐมํ.}
\csromanmatikaanchor{mtk.22}
\csromantocmarkref{subsection}{9. กายคตาสติสุตฺต}{mtk.22}
\bookmark[dest={mtk.22},level=2]{9. กายคตาสติสุตฺต}
\csromanheader{2}{9. กายคตาสติสุตฺต}
\proseitem{153}{เอวํ เม สุตํ\csromandash{}เอกํ สมยํ ภควา สาวตฺถิยํ วิหรติ เชตวเน อนาถปิณฺฑิกสฺส อาราเม.\csromanspacer{} อถ โข สมฺพหุลานํ ภิกฺขูนํ ปจฺฉาภตฺตํ ปิณฺฑปาตปฏิกฺกนฺตานํ อุปฏฺฐานสาลายํ สนฺนิสินฺนานํ สนฺนิปติตานํ}

\vspace{\tipitakaparskip}
\csromancenter{กายคตาสติสุตฺต (119) อยมนฺตรากถา อุทปาทิ “อจฺฉริยํ อาวุโส, อพฺภุตํ อาวุโส, ยาวญฺจิทํ เตน ภควตา ชานตา ปสฺสตา อรหตา สมฺมาสมฺพุทฺเธน กายคตาสติ\footnote{กายคตา สติ (สฺยา, กํ, อิ)} ภาวิตา พหุลีกตา มหปฺผลา วุตฺตา มหานิสํสา”ติ.\csromanspacer{} อยญฺจ หิทํ เตสํ ภิกฺขูนํ อนฺตรากถา วิปฺปกตา โหติ.\csromanspacer{} อถ โข ภควา สายนฺหสมยํ ปฏิสลฺลานา วุฏฺฐิโต เยน อุปฏฺฐานสาลา เตนุปสงฺกมิ, อุปสงฺกมิตฺวา ปญฺญตฺเต อาสเน นิสีทิ.\csromanspacer{} นิสชฺช โข ภควา ภิกฺขู อามนฺเตสิ “กาย นุตฺถ ภิกฺขเว เอตรหิ กถาย สนฺนิสินฺนา, กา จ ปน โว อนฺตารากถา วิปฺปกตา”ติ.\csromanspacer{} อิธ ภนฺเต อมฺหากํ ปจฺฉาภตฺตํ ปิณฺฑปาตปฏิกฺกนฺตานํ อุปฏฺฐานสาลายํ สนฺนิสินฺนานํ สนฺนิปติตานํ อยมนฺตรากถา อุทปาทิ “อจฺฉริยํ อาวุโส, อพฺภุตํ อาวุโส, ยาวญฺจิทํ เตน ภควตา ชานตา ปสฺสตา อรหตา สมฺมาสมฺพุทฺเธน กายคตาสติ ภาวิตา พหุลีกตา มหปฺผลา วุตฺตา มหานิสํสา”ติ.\csromanspacer{} อยํ โข โน ภนฺเต อนฺตรากถา วิปฺปกตา, อถ ภควา อนุปฺปตฺโตติ.}
\proseitem{154}{\csromanfolio{131}กถํ ภาวิตา จ ภิกฺขเว กายคตาสติ กถํ พหุลีกตา มหปฺผลา โหติ มหานิสํสา.\csromanspacer{} อิธ ภิกฺขเว ภิกฺขุ อรญฺญคโต วา รุกฺขมูลคโต วา สุญฺญาคารคโต วา นิสีทติ ปลฺลงฺกํ อาภุชิตฺวา อุชุํ กายํ ปณิธาย ปริมุขํ สติํ อุปฏฺฐเปตฺวา.\csromanspacer{} โส สโตว อสฺสสติ, สโตว ปสฺสสติ.\csromanspacer{} ทีฆํ วา อสฺสสนฺโต “ทีฆํ อสฺสสามี”ติ ปชานาติ, ทีฆํ วา ปสฺสสนฺโต “ทีฆํ ปสฺสสามี”ติ ปชานาติ.\csromanspacer{} รสฺสํ วา อสฺสสนฺโต “รสฺสํ อสฺสสามี”ติ ปชานาติ, รสฺสํ วา ปสฺสสนฺโต “รสฺสํ ปสฺสสามี”ติ ปชานาติ. “สพฺพกายปฏิสํเวที อสฺสสิสฺสามี”ติ สิกฺขติ, “สพฺพกายปฏิสํเวที ปสฺสสิสฺสามี”ติ สิกฺขติ. “ปสฺสมฺภยํ กายสงฺขารํ อสฺสสิสฺสามี”ติ สิกฺขติ, “ปสฺสมฺภยํ กายสงฺขารํ ปสฺสสิสฺสามี”ติ สิกฺขติ.\csromanspacer{} ตสฺส เอวํ อปฺปมตฺตสฺส อาตาปิโน ปหิตตฺตสฺส วิหรโต เย เคหสิตา\footnote{เคหสฺสิตา (ฏีกา)} สรสงฺกปฺปา, เต ปหียนฺติ.\csromanspacer{} เตสํ ปหานา อชฺฌตฺตเมว จิตฺตํ สนฺติฏฺฐติ สนฺนิสีทติ เอโกทิ โหติ\footnote{เอโกที โหติ (สี), เอโกทิโภติ (สฺยา, กํ)} สมาธิยติ.\csromanspacer{} เอวํ ภิกฺขเว ภิกฺขุ กายคตาสติํ\footnote{กายคตํ สติํ (สฺยา, กํ, อิ)} ภาเวติ. (1)}

\prose{\csromanfolio{132}ปุน จปรํ ภิกฺขเว ภิกฺขุ คจฺฉนฺโต วา “คจฺฉามี”ติ ปชานาติ, ฐิโต วา “ฐิโตมฺหี”ติ ปชานาติ, นิสินฺโน วา “นิสินฺโนมฺหี”ติ ปชานาติ, สยาโน วา “สยาโนมฺหี”ติ ปชานาติ, ยถา ยถา วา ปนสฺส กาโย ปณิหิโต โหติ, ตถา ตถา นํ ปชานาติ.\csromanspacer{} ตสฺส เอวํ อปฺปมตฺตสฺส อาตาปิโน ปหิตตฺตสฺส วิหรโต เย เคหสิตา สรสงฺกปฺปา, เต ปหียนฺติ.\csromanspacer{} เตสํ ปหานา อชฺฌตฺตเมว จิตฺตํ สนฺติฏฺฐติ สนฺนิสีทติ เอโกทิ โหติ สมาธิยติ.\csromanspacer{} เอวมฺปิ ภิกฺขเว ภิกฺขุ กายคตาสติํ ภเวติ. (2)}

\prose{ปุน จปรํ ภิกฺขเว ภิกฺขุ อภิกฺกนฺเต ปฏิกฺกนฺเต สมฺปชานการี โหติ, อาโลกิเต วิโลกิเต สมฺปชานการี โหติ, สมิญฺชิเต ปสาริเต สมฺปชานการิ โหติ, สํฆาฏิปตฺตจีวรธารเณ สมฺปชานการี โหติ, อสิเต ปีเต ขายิเต สายิเต สมฺปชานการี โหติ, อุจฺจารปสฺสาวกมฺเม สมฺปชานการี โหติ, คเต ฐิเต นิสินฺเน สุตฺเต ชาคริเต ภาสิเต ตุณฺหีภาเว สมฺปชานการี โหติ.\csromanspacer{} ตสฺส เอวํ อปฺปมตฺตสฺส อาตาปิโน ปหิตตฺตสฺส วิหรโต เย เคหสิตา สรสงฺกปฺปา, เต ปหียนฺติ.\csromanspacer{} เตสํ ปหานา อชฺฌตฺตเมว จิตฺตํ สนฺติฏฺฐติ สนฺนิสีทติ เอโกทิ โหติ สมาธิยติ.\csromanspacer{} เอวมฺปิ ภิกฺขเว ภิกฺขุ กายคตาสติํ ภาเวติ. (3)}

\prose{ปุน จปรํ ภิกฺขเว ภิกฺขุ อิมเมว กายํ อุทฺธํ ปาทตลา อโธ เกสมตฺถกา ตจปริยนฺตํ ปูรํ นานปฺปการสฺส อสุจิโน ปจฺจเวกฺขติ “อตฺถิ อิมสฺมิํ กาเย เกสา โลมา นขา ทนฺตา ตโจ, มํสํ นฺหารุ\footnote{นหารุ (สี, สฺยา, กํ, อิ)} อฏฺฐิ อฏฺฐิมิญฺชํ วกฺกํ, หทยํ ยกนํ กิโลมกํ ปิหกํ ปปฺผาสํ, อนฺตํ อนฺตคุณํ อุทริยํ กรีสํ, ปิตฺตํ เสมฺหํ ปุพฺโพ โลหิตํ เสโท เมโท, อสฺสุ วสา เขโฬ สิงฺฆาณิกา ลสิกา มุตฺตนฺ”ติ.}

\prose{เสยฺยถาปิ ภิกฺขเว อุภโตมุขา ปุโตฬิ\footnote{มูโหฬี (สี, สฺยา, กํ, อิ)} ปูรา นานาวิหิตสฺส ธญฺญสฺส.\csromanspacer{} เสยฺยถิทํ, สาลีนํ วีหีนํ มุคฺคานํ มาสานํ ติลานํ ตณฺฑุลานํ.\csromanspacer{} ตเมนํ จกฺขุมา ปุริโส มุญฺจิตฺวา ปจฺจเวกฺเขยฺย “อิเม สาลี อิเม วีหี อิเม มุคฺคา อิเม มาสา อิเม ติลา อิเม ตณฺฑุลา”ติ.\csromanspacer{} เอวเมว โข ภิกฺขเว ภิกฺขุ อิมเมว กายํ อุทฺธํ ปาทตลา อโธ เกสมตฺถกา ตจปริยนฺตํ}

\vspace{\tipitakaparskip}
\csromancenter{กายคตาสติสุตฺต (119) ปูรํ นานปฺปการสฺส อสุจิโน ปจฺจเวกฺขติ “อตฺถิ อิมสฺมิํ กาเย เกสา โลมา นขา ทนฺตา ตโจ, มํสํ นฺหารุ อฏฺฐิ อฏฺฐิมิญฺชํ วกฺกํ, หทยํ ยกนํ กิโลมกํ ปิหกํ ปปฺผาสํ, อนฺตํ อนฺตคุณํ อุทริยํ กรีสํ, ปิตฺตํ เสมฺหํ ปุพฺโพ โลหิตํ เสโท เมโท, อสฺสุ วสา เขโฬ สิงฺฆาณิกา ลสิกา มุตฺตนฺ”ติ.\csromanspacer{} ตสฺส เอวํ อปฺปมตฺตสฺส อาตาปิโน ปหิตตฺตสฺส วิหรโต เย เคหสิตา สรสงฺกปฺปา เต ปหียนฺติ.\csromanspacer{} เตสํ ปหานา อชฺฌตฺตเมว จิตฺตํ สนฺติฏฺฐติ สนฺนิสีทติ เอโกทิ โหติ สมาธิยติ.\csromanspacer{} เอวมฺปิ ภิกฺขเว ภิกฺขุ กายคตาสติํ ภาเวติ. (4)}
\prose{\csromanfolio{133}ปุน จปรํ ภิกฺขเว ภิกฺขุ อิมเมว กายํ ยถาฐิตํ ยถาปณิหิตํ ธาตุโส ปจฺจเวกฺขติ “อตฺถิ อิมสฺมิํ กาเย ปถวีธาตุ อาโปธาตุ เตโชธาตุ วาโยธาตู”ติ.}

\prose{เสยฺยถาปิ ภิกฺขเว ทกฺโข โคฆาตโก วา โวฆาตกนฺเตวาสี วา คาวิํ วธิตฺวา จตุมหาปเถ\footnote{จาตุมฺมหาปเถ (สี, สฺยา, กํ, อิ)} พิลโส วิภชิตฺวา\footnote{ปฏิวิภชิตฺวา (สี, สฺยา, กํ, อิ)} นิสินฺโน อสฺส.\csromanspacer{} เอวเมว โข ภิกฺขเว ภิกฺขุ อิมเมว กายํ ยถาฐิตํ ยถาปณิหิตํ ธาตุโส ปจฺจเวกฺขติ “อตฺถิ อิมสฺมิํ กาเย ปถวีธาตุ อาโปธาตุ เตโชธาตุ วาโยธาตู”ติ.\csromanspacer{} ตสฺส เอวํ อปฺปมตฺตสฺส อาตาปิโน ปหิตตฺตสฺส วิหรโต เย เคหสิตา สรสงฺกปฺปา, เต ปหียนฺติ.\csromanspacer{} เตสํ ปหานา อชฺฌตฺตเมว จิตฺตํ สนฺติฏฺฐติ สนฺนิสีทติ เอโกทิ โหติ สมาธิยติ.\csromanspacer{} เอวมฺปิ ภิกฺขเว ภิกฺขุ กายคตาสติํ ภาเวติ. (5)}

\prose{ปุน จปรํ ภิกฺขเว ภิกฺขุ เสยฺยถาปิ ปสฺเสยฺย สรีรํ สิวถิกาย\footnote{สีวถิกาย (สี, สฺยา, กํ, อิ)} ฉฑฺฑิตํ เอกาหมตํ วา ทฺวีหมตํ วา ตีหมตํ วา อุทฺธุมาตกํ วินีลกํ วิปุพฺพกชาตํ.\csromanspacer{} โส อิมเมว กายํ อุปสํหรติ “อยมฺปิ โข กาโย เอวํธมฺโม เอวํภาวี เอวํอนตีโต”ติ\footnote{เอตํ อนตีโตติ (สี)}.\csromanspacer{} ตสฺส เอวํ อปฺปมตฺตสฺส อาตาปิโน ปหิตตฺตสฺส วิหรโต เย เคหสิตา สรสงฺกปฺปา, เต ปหียนฺติ.\csromanspacer{} เตสํ ปหานา อชฺฌตฺตเมว จิตฺตํ สนฺติฏฺฐติ สนฺนิสีทติ เอโกทิ โหติ สมาธิยติ.\csromanspacer{} เอวมฺปิ ภิกฺขเว ภิกฺขุ กายคตาสติํ ภาเวติ. (6)}

\prose{\csromanfolio{134}ปุน จปรํ ภิกฺขเว ภิกฺขุ เสยฺยถาปิ ปสฺเสยฺย สรีรํ สิวถิกาย ฉฑฺฑิตํ กาเกหิ วา ขชฺชมานํ กุลเลหิ วา ขชฺชมานํ คิชฺเฌหิ วา ขชฺชมานํ กงฺเกหิ วา ขชฺชมานํ สุนเขหิ วา ขชฺชมานํ พฺยคฺเฆหิ วา ขชฺชมานํ ทีปีหิ วา ขชฺชมานํ สิงฺคาเลหิ วา\footnote{คิชฺเฌหิ วา ขชฺชมานํ สุวาเนหิ วา ขชฺชมานํ สิคาเลหิ วา (สิ, สฺยา, กํ, อิ)} ขชฺชมานํ วิวิเธหิ วา ปาณกชาเตหิ ขชฺชมานํ.\csromanspacer{} โส อิมเมว กายํ อุปสํหรติ “อยมฺปิ โข กาโย เอวํธมฺโม เอวํภาวี เอวํอนตีโต”ติ.\csromanspacer{} ตสฺส เอวํ อปฺปมตฺตสฺส ฯเปฯ.\csromanspacer{} เอวมฺปิ ภิกฺขเว ภิกฺขุ กายคตาสติํ ภาเวติ. (7)}

\prose{ปุน จปรํ ภิกฺขเว ภิกฺขุ เสยฺยถาปิ ปสฺเสยฺย สรีรํ สิวถิกาย ฉฑฺฑิตํ อฏฺฐิกสงฺขลิกํ สมํสโลหิตํ นฺหารุสมฺพนฺธํ ฯเปฯ อฏฺฐิกสงฺขลิกํ นิมฺมํสโลหิตมกฺขิตํ นฺหารุสมฺพนฺธํ ฯเปฯ อฏฺฐิกสงฺขลิกํ อปคตมํสโลหิตํ นฺหารุสมฺพนฺธํ ฯเปฯ อฏฺฐิกานิ อปคตสมฺพนฺธานิ\footnote{อปคตนหารุสมฺพนฺธานิ (สฺยา, กํ)} ทิสาวิทิสาวิกฺขิตฺตานิ\footnote{ทิสาวิทิสาสุ วิกฺขิตฺตานิ (สี, อิ)}, อญฺเญน หตฺถฏฺฐิกํ อญฺเญน ปาทฏฺฐิกํ อญฺเญน โคปฺผกฏฺฐิกํ\footnote{อญฺเญน โคปฺผกฏฺฐิกนฺติ อิทํ สี-สฺยา-กํ-อิ-โปตฺถเกสุ นตฺถิ. 5-5. อญฺเญน กฏฏฺฐิกํ อญฺเญน ปิฏฺฐิกณฺฏกํ อญฺเญน สีสกฏาหํ (สี, สฺยา, กํ, อิ)} อญฺเญน ชงฺฆฏฺฐิกํ อญฺเญน อูรุฏฺฐิกํ อญฺเญน กฏิฏฺฐิกํ5 อญฺเญน ผาสุกฏฺฐิกํ อญฺเญน ปิฏฺฐิฏฺฐิกํ อญฺเญน ขนฺธฏฺฐิกํ อญฺเญน คีวฏฺฐิกํ อญฺเญน หนุกฏฺฐิกํ อญฺเญน ทนฺตฏฺฐิกํ อญฺเญน สีสกฏาหํ5.\csromanspacer{} โส อิมเมว กายํ อุปสํหรติ “อยมฺปิ โข กาโย เอวํธมฺโม เอวํภาวี เอวํอนตีโต”ติ.\csromanspacer{} ตสฺส เอวํ อปฺปมตฺตสฺส ฯเปฯ.\csromanspacer{} เอวมฺปิ ภิกฺขเว ภิกฺขุ กายคตาสติํ ภาเวติ. (8- 11)}

\prose{ปุน จปรํ ภิกฺขเว ภิกฺขุ เสยฺยถาปิ ปสฺเสยฺย สรีรํ สิวถิกาย ฉฑฺฑิตํ อฏฺฐิกานิ เสตานิ สงฺขวณฺณปฏิภาคานิ\footnote{สงฺขวณฺณูปนิภานิ (สี, สฺยา, กํ, อิ)} ฯเปฯ อฏฺฐิกานิ ปุญฺชกิตานิ เตโรวสฺสิกานิ ฯเปฯ อฏฺฐิกานิ ปูตีนิ จุณฺณกชาตานิ.\csromanspacer{} โส อิมเมว กายํ อุปสํหรติ “อยมฺปิ โข กาโย เอวํธมฺโม เอวํภาวี เอวํอนตีโต”ติ.\csromanspacer{} ตสฺส เอวํ อปฺปมตฺตสฺส ฯเปฯ.\csromanspacer{} เอวมฺปิ ภิกฺขเว ภิกฺขุ กายคตาสติํ ภาเวติ. (12-14)}

\proseitem{155}{ปุน จปรํ ภิกฺขเว ภิกฺขุ วิวิจฺเจว กาเมหิ ฯเปฯ ปฐมํ ฌานํ อุปสมฺปชฺช วิหรติ, โส อิมเมว กายํ วิเวกเชน ปีติสุเขน อภิ-}

\vspace{\tipitakaparskip}
\csromancenter{กายคตาสติสุตฺต (119) สนฺเทติ ปริสนฺเทติ ปริปูเรติ ปริปฺผรติ, นาสฺส กิญฺจิ สพฺพาวโต กายสฺส วิเวกเชน ปีติสุเขน อปฺผุฏํ โหติ.\csromanspacer{} เสยฺยถาปิ ภิกฺขเว ทกฺโข นฺหาปโก\footnote{นหาปโก (สี, สฺยา, กํ, อิ)} วา นฺหาปกนฺเตวาสี วา กํสถาเล นฺหานียจุณฺณานิ\footnote{นหานียจุณฺณานิ (สี, สฺยา, กํ, อิ)} อากิริตฺวา อุทเกน ปริปฺโผสกํ ปริปฺโผสกํ สนฺเนยฺย, สายํ นฺหานียปิณฺฑิ\footnote{สาสฺส นหานียปิณฺฑี (สี, สฺยา, กํ, อิ)} สฺเนหานุคตา สฺเนหปเรตา สนฺตรพาหิรา ผุฏา สฺเนเหน, น จ ปคฺฆริณี.\csromanspacer{} เอวเมว โข ภิกฺขเว ภิกฺขุ อิมเมว กายํ วิเวกเชน ปีติสุเขน อภิสนฺเทติ ปริสนฺเธติ ปริปูเรติ ปริปฺผรติ, นาสฺส กิญฺจิ สพฺพาวโต กายสฺส วิเวกเชน ปีติสุเขน อปฺผุฏํ โหติ.\csromanspacer{} ตสฺส เอวํ อปฺปมตฺตสฺส ฯเปฯ.\csromanspacer{} เอวมฺปิ ภิกฺขเว ภิกฺขุ กายคตาสติํ ภาเวติ. (15)}
\prose{\csromanfolio{135}ปุน จปรํ ภิกฺขเว ภิกฺขุ วิตกฺกวิจารานํ วูปสมา ฯเปฯ ทุติยํ ฌานํ อุปสมฺปชฺช วิหรติ, โส อิมเมว กายํ สมาธิเชน ปีติสุเขน อภิสนฺเทติ ปริสนฺเทติ ปริปูเรติ ปริปฺผรติ, นาสฺส กิญฺจิ สพฺพาวโต กายสฺส สมาธิเชน ปีติสุเขน อปฺผุฏํ โหติ.\csromanspacer{} เสยฺยถาปิ ภิกฺขเว อุทกรหโท คมฺภีโร อุพฺภิโททโก\footnote{อุพฺภิโตทโก (สฺยา, กํ, ก)}, ตสฺส เนวสฺส ปุรตฺถิมาย ทิสาย อุทกสฺส อายมุขํ, น ปจฺฉิมาย ทิสาย อุทกสฺส อายมุขํ, น อุตฺตราย ทิสาย อุทกสฺส อายมุขํ, น ทกฺขิณาย ทิสาย อุทกสฺส อายมุขํ, เทโว จ น กาเลน กาลํ สมฺมา ธารํ อนุปฺปเวจฺเฉยฺย.\csromanspacer{} อถ โข ตมฺหาว อุทกรหทา สีตา วาริธารา อุพฺภิชฺชิตฺวา ตเมว อุทกรหทํ สีเตน วารินา อภิสนฺเทยฺย ปริสนฺเทยฺย ปริปูเรยฺย ปริปฺผเรยฺย, นาสฺส กิญฺจิ สพฺพาวโต อุทกรหทสฺส สีเตน วารินา อปฺผุฏํ อสฺส.\csromanspacer{} เอวเมว โข ภิกฺขเว ภิกฺขุ อิมเมว กายํ สมาธิเชน ปีติสุเขน อภิสนฺเทติ ปริสนฺเทติ ปริปูเรติ ปริปฺผรติ, นาสฺส กิญฺจิ สพฺพาวโต กายสฺส สมาธิเชน ปีติสุเขน อปฺผุฏํ โหติ.\csromanspacer{} ตสฺส เอวํ อปฺปมตฺตสฺส ฯเปฯ.\csromanspacer{} เอวมฺปิ ภิกฺขเว ภิกฺขุ กายคตาสติํ ภาเวติ. (16)}

\prose{ปุน จปรํ ภิกฺขเว ภิกฺขุ ปีติยา จ วิราคา ฯเปฯ ตติยํ ฌานํ อุปสมฺปชฺช วิหรติ, โส อิมเมว กายํ นิปฺปีติเกน สุเขน อภิสนฺเทติ ปริสนฺเทติ ปริปูเรติ ปริปฺผรติ, นาสฺส กิญฺจิ สพฺพาวโต กายสฺส นิปฺปีติเกน สุเขน อปฺผุฏํ โหติ.\csromanspacer{} เสยฺยถาปิ ภิกฺขเว อุปฺปลินิยํ วา ปทุมินิยํ วา ปุณฺฑรีกินิยํ วา อปฺเปกจฺจานิ อุปฺปลานิ วา ปทุมานิ วา ปุณฺฑรีกานิ \csromanfolio{136}วา อุทเก ชาตานิ อุทเก สํวฑฺฒานิ อุทกานุคฺคตานิ อนฺโตนิมุคฺคโปสีนิ, ตานิ ยาว จคฺคา ยาว จ มูลา สีเตน วารินา อภิสนฺนานิ ปริสนฺนานิ\footnote{อภิสนฺทานิ ปริสนฺทานิ (ก)} ปริปูรานิ ปริปฺผุฏานิ, นาสฺส\footnote{น เนสํ (?)} กิญฺจิ สพฺพาวตํ อุปฺปลานํ วา ปทุมานํ วา ปุณฺฑรีกานํ วา สีเตน วารินา อปฺผุฏํ อสฺส.\csromanspacer{} เอวเมว โข ภิกฺขเว ภิกฺขุ อิมเมว กายํ นิปฺปีติเกน สุเขน อภิสนฺเทติ ปริสนฺเทติ ปริปูเรติ ปริปฺผรติ, นาสฺส กิญฺจิ สพฺพาวโต กายสฺส นิปฺปีติเกน สุเขน อปฺผุฏํ โหติ.\csromanspacer{} ตสฺส เอวํ อปฺปมตฺตสฺส ฯเปฯ.\csromanspacer{} เอวมฺปิ ภิกฺขเว ภิกฺขุ กายคตาสติํ ภาเวติ. (17)}

\prose{ปุน จปรํ ภิกฺขเว ภิกฺขุ สุขสฺส จ ปหานา ฯเปฯ จตุตฺถํ ฌานํ อุปสมฺปชฺช วิหรติ, โส อิมเมว กายํ ปริสุทฺเธน เจตสา ปริโยทาเตน ผริตฺวา นิสินฺโน โหติ, นาสฺส กิญฺจิ สพฺพาวโต กายสฺส ปริสุทฺเธน เจตสา ปริโยทาเตน อปฺผุฏํ โหติ.\csromanspacer{} เสยฺยถาปิ ภิกฺขเว ปุริโส โอทาเตน วตฺเถน สสีสํ ปารุปิตฺวา นิสินฺโน อสฺส, นาสฺส กิญฺจิ สพฺพาวโต กายสฺส โอทาเตน วตฺเถน อปฺผุฏํ อสฺส.\csromanspacer{} เอวเมว โข ภิกฺขเว ภิกฺขุ อิมเมว กายํ ปริสุทฺเธน เจตสา ปริโยทาเตน ผริตฺวา นิสินฺโน โหติ, นาสฺส กิญฺจิ สพฺพาวโต กายสฺส ปริสุทฺเธน เจตสา ปริโยทาเตน อปฺผุฏํ โหติ.\csromanspacer{} ตสฺส เอวํ อปฺปมตฺตสฺส อาตาปิโน ปหิตตฺตสฺส วิหรโต เย เคหสิตา สรสงฺกปฺปา เต ปหียนฺติ, เตสํ ปหานา อชฺฌตฺตเมว จิตฺตํ สนฺติฏฺฐติ สนฺนิสีทติ เอโกทิ โหติ สมาธิยติ.\csromanspacer{} เอวมฺปิ ภิกฺขเว ภิกฺขุ กายคตาสติํ ภาเวติ. (18)}

\proseitem{156}{ยสฺส กสฺสจิ ภิกฺขเว กายคตาสติ ภาวิตา พหุลีกตา, อนฺโตคธาวาสฺส\footnote{อนฺโตคธา ตสฺส (สี, อิ)} กุสลา ธมฺมา เย เกจิ วิชฺชาภาคิยา.\csromanspacer{} เสยฺยถาปิ ภิกฺขเว ยสฺส กสฺสจิ มหาสมุทฺโท เจตสา ผุโฏ, อนฺโตคธาวาสฺส กุนฺนทิโย ยา กาจิ สมุทฺทงฺคมา.\csromanspacer{} เอวเมว โข ภิกฺขเว ยสฺส กสฺสจิ กายคตาสติ ภาวิตา พหุลีกตา, อนฺโตคธาวาสฺส กุสลา ธมฺมา เย เกจิ วิชฺชาภาคิยา.}

\prose{ยสฺส กสฺสจิ ภิกฺขเว กายคตาสติ อภาวิตา อพหุลีกตา, ลภติ ตสฺส มาโร โอตารํ, ลภติ ตสฺส มาโร อารมฺมณํ\footnote{อารมณํ (?)}.}

\vspace{\tipitakaparskip}
\csromancenter{กายคตาสติสุตฺต (119) เสยฺยถาปิ ภิกฺขเว ปุริโส ครุกํ สิลาคุฬํ อลฺลมตฺติกาปุญฺเช ปกฺขิเปยฺย.\csromanspacer{} ตํ กิํ มญฺญถ ภิกฺขเว, อปิ นุ ตํ ครุกํ สิลาคุฬํ อลฺลมตฺติกาปุญฺเช ลเภถ โอตารนฺติ.\csromanspacer{} เอวํ ภนฺเต.\csromanspacer{} เอวเมว โข ภิกฺขเว ยสฺส กสฺสจิ กายคตาสติ อภาวิตา อพหุลีกตา, ลภติ ตสฺส มาโร โอตารํ, ลภติ ตสฺส มาโร อารมฺมณํ.\csromanspacer{} เสยฺยถาปิ ภิกฺขเว สุกฺขํ กฏฺฐํ โกฬาปํ\footnote{โกฬาปํ อารกา อุทกา ถเล นิกฺขิตํ (ก)}, อถ ปุริโส อาคจฺเฉยฺย อุตฺตรารณิํ อาทาย “อคฺคิํ อภินิพฺพตฺเตสฺสามิ เตโช ปาตุกริสฺสามี”ติ.\csromanspacer{} ตํ กิํ มญฺญถ ภิกฺขเว, อปิ นุ โส ปุริโส อมุํ สุกฺขํ กฏฺฐํ โกฬาปํ อุตฺตรารณิํ อาทาย อภิมนฺเถนฺโต\footnote{อภิมตฺเถนฺโต (สฺยา, กํ, อิ, ก) 3. สสฺเนหํ อารกา อุทกา ถเล นิกฺขิตฺตํ (ก)} อคฺคิํ อภินิพฺพตฺเตยฺย เตโช ปาตุกเรยฺยาติ.\csromanspacer{} เอวํ ภนฺเต.\csromanspacer{} เอวเมว โข ภิกฺขเว ยสฺส กสฺสจิ กายคตาสติ อภาวิตา อพหุลีกตา, ลภติ ตสฺส มาโร โอตารํ, ลภติ ตสฺส มาโร อารมฺมณํ.\csromanspacer{} เสยฺยถาปิ ภิกฺขเว อุทกมณิโก ริตฺโต ตุจฺโฉ อาธาเร ฐปิโต, อถ ปุริโส อาคจฺเฉยฺย อุทกภารํ อาทาย.\csromanspacer{} ตํ กิํ มญฺญถ ภิกฺขเว, อปิ นุ โส ปุริโส ลเภถ อุทกสฺส นิกฺเขปนนฺติ.\csromanspacer{} เอวํ ภนฺเต.\csromanspacer{} เอวเมว โข ภิกฺขเว ยสฺส กสฺสจิ กายคตาสติ อภาวิตา อพหุลีกตา, ลภติ ตสฺส มาโร โอตารํ, ลภติ ตสฺส มาโร อารมฺมณํ.}
\proseitem{157}{\csromanfolio{137}ยสฺส กสฺสจิ ภิกฺขเว กายคตาสติ ภาวิตา พหุลีกตา, น ตสฺส ลภติ มาโร โอตารํ, น ตสฺส ลภติ มาโร อารมฺมณํ.\csromanspacer{} เสยฺยถาปิ ภิกฺขเว ปุริโส ลหุกํ สุตฺตคุฬํ สพฺพสารมเย อคฺคฬผลเก ปกฺขิเปยฺย.\csromanspacer{} ตํ กิํ มญฺญถ ภิกฺขเว, อปิ นุ โส ปุริโส ตํ ลหุกํ สุตฺตคุฬํ สพฺพสารมเย อคฺคฬผลเก ลเภถ โอตารนฺติ.\csromanspacer{} โน เหตํ ภนฺเต.\csromanspacer{} เอวเมว โข ภิกฺขเว ยสฺส กสฺสจิ กายคตาสติ ภาวิตา พหุลีกตา, น ตสฺส ลภติ มาโร โอตารํ, น ตสฺส ลภติ มาโร อารมฺมณํ.\csromanspacer{} เสยฺยถาปิ ภิกฺขเว อลฺลํ กฏฺฐํ สสฺเนหํ3, อถ ปุริโส อาคจฺเฉยฺย อุตฺตรารณิํ อาทาย “อคฺคิํ อภินิพฺพตฺเตสฺสามิ เตโช ปาตุกริสฺสามี”ติ.\csromanspacer{} ตํ กิํ มญฺญถ ภิกฺขเว, อปิ นุ โส ปุริโส อมุํ อลฺลํ กฏฺฐํ สสฺเนหํ อุตฺตรารณิํ อาทาย อภิมนฺเถนฺโต อคฺคิํ อภินิพฺพตฺเตยฺย เตโช ปาตุกเรยฺยาติ.\csromanspacer{} โน เหตํ ภนฺเต.\csromanspacer{} เอวเมว โข \csromanfolio{138}ภิกฺขเว ยสฺส กสฺสจิ กายคตาสติ ภาวิตา พหุลีกตา, น ตสฺส ลภติ มาโร โอตารํ, น ตสฺส ลภติ มาโร อารมฺมณํ.\csromanspacer{} เสยฺยถาปิ ภิกฺขเว อุทกมณิโก ปูโร อุทกสฺส สมติตฺติโก กากเปยฺโย อาธาเร ฐปิโต, อถ ปุริโส อาคจฺเฉยฺย อุทกภารํ อาทาย.\csromanspacer{} ตํ กิํ มญฺญถ ภิกฺขเว, อปิ นุ โส ปุริโส ลเภถ อุทกสฺส นิกฺเขปนนฺติ.\csromanspacer{} โน เหตํ ภนฺเต.\csromanspacer{} เอวเมว โข ภิกฺขเว ยสฺส กสฺสจิ กายคตาสติ ภาวิตา พหุลีกตา, น ตสฺส ลภติ มาโร โอตารํ, น ตสฺส ลภติ มาโร อารมฺมณํ.}

\proseitem{158}{ยสฺส กสฺสจิ ภิกฺขเว กายคตาสติ ภาวิตา พหุลีกตา, โส ยสฺส ยสฺส อภิญฺญาสจฺฉิกรณียสฺส ธมฺมสฺส จิตฺตํ อภินินฺนาเมติ อภิญฺญาสจฺฉิกิริยาย, ตตฺร ตเตฺรว สกฺขิภพฺพตํ ปาปุณาติ สติ สติ- อายตเน.\csromanspacer{} เสยฺยถาปิ ภิกฺขเว อุทกมณิโก ปูโร อุทกสฺส สมติตฺติโก กากเปยฺโย อาธาเร ฐปิโต, ตเมนํ พลวา ปุริโส ยโต ยโต อาวิญฺเฉยฺย, อาคจฺเฉยฺย อุทกนฺติ.\csromanspacer{} เอวํ ภนฺเต.\csromanspacer{} เอวเมว โข ภิกฺขเว ยสฺส กสฺสจิ กายคตาสติ ภาวิตา พหุลีกตา, โส ยสฺส ยสฺส อภิญฺญาสจฺฉิกรณียสฺส ธมฺมสฺส จิตฺตํ อภินินฺนาเมติ อภิญฺญาสจฺฉิกิริยาย, ตตฺร ตเตฺรว สกฺขิภพฺพตํ ปาปุณาติ สติ สติ-อายตเน.\csromanspacer{} เสยฺยถาปิ ภิกฺขเว สเม ภูมิภาเค จตุรสฺสา โปกฺขรณี\footnote{โปกฺขริณี (สี)} อสฺส อาฬิพนฺธา ปูรา อุทกสฺส สมติตฺติกา กากเปยฺยา, ตเมนํ พลวา ปุริโส ยโต ยโต อาฬิํ มุญฺเจยฺย, อาคจฺเฉยฺย อุทกนฺติ.\csromanspacer{} เอวํ ภนฺเต.\csromanspacer{} เอวเมว โข ภิกฺขเว ยสฺส กสฺสจิ กายคตาสติ ภาวิตา พหุลีกตา, โส ยสฺส ยสฺส อภิญฺญาสจฺฉิกรณียสฺส ธมฺมสฺส จิตฺตํ อภินินฺนาเมติ อภิญฺญาสจฺฉิกิริยาย, ตตฺร ตเตฺรว สกฺขิภพฺพตํ ปาปุณาติ สติ สติ-อายตเน.\csromanspacer{} เสยฺยถาปิ ภิกฺขเว สุภูมิยํ จตุมหาปเถ อาชญฺญรโถ ยุตฺโต อสฺส ฐิโต โอธสฺตปโตโท\footnote{โอภสฺตปโตโท (ก), อุภนฺตรปโฏโท (สฺยา, กํ) อว + ธํสุ + ต = โอธสฺต-อิติ ปทวิภาโค.}, ตเมนํ ทกฺโข โยคฺคาจริโย อสฺสทมฺมสารถิ อภิรุหิตฺวา วาเมน หตฺเถน รสฺมิโย คเหตฺวา ทกฺขิเณน หตฺเถน ปโตทํ คเหตฺวา เยนิจฺฉกํ ยทิจฺฉกํ}

\vspace{\tipitakaparskip}
\csromancenter{กายคตาสติสุตฺต (119) สาเรยฺยาปิ ปจฺจาสาเรยฺยาปิ.\csromanspacer{} เอวเมว โข ภิกฺขเว ยสฺส กสฺสจิ กายคตาสติ ภาวิตา พหุลีกตา, โส ยสฺส ยสฺส อภิญฺญาสจฺฉิกรณียสฺส ธมฺมสฺส จิตฺตํ อภินินฺนาเมติ อภิญฺญาสจฺฉิกิริยาย, ตตฺร ตเตฺรว สกฺขิภพฺพตํ ปาปุณาติ สติ สติ-อายตเน.}
\proseitem{159}{\csromanfolio{139}กายคตาย ภิกฺขเว สติยา อาเสวิตาย ภาวิตาย พหุลีกตาย ยานีกตาย วตฺถุกตาย อนุฏฺฐิตาย ปริจิตาย สุสมารทฺธาย ทสานิสํสา ปาฏิกงฺขา\csromandash{}อรติรติสโห โหติ, น จ ตํ อรติ สหติ, อุปฺปนํ อรติํ อภิภุยฺย วิหรติ. (1)}

\prose{ภยเภรวสโห โหติ, น จ ตํ ภยเภรวํ สหติ, อุปฺปนฺนํ ภยเภรวํ อภิภุยฺย วิหรติ. (2)}

\prose{ขโม โหติ สีตสฺส อุณฺหาสฺส ชิฆจฺฉาย ปิปาสาย ฑํสมกสวาตาตปสรีสปสมฺผสฺสานํ ทุรุตฺตานํ ทุราคตานํ วจนปถานํ อุปฺปนฺนานํ สารีริกานํ เวทนานํ ทุกฺขานํ ติพฺพานํ ขรานํ กฏุกานํ อสาตานํ อมนาปานํ ปาณหรานํ อธิวาสกชาติโก โหติ. (3)}

\prose{จตุนฺนํ ฌานานํ อาภิเจตสิกานํ ทิฏฺฐธมฺมสุขวิหารานํ นิกามลาภี โหติ อกิจฺฉลาภี อกสิรลาภี. (4)}

\prose{โส อเนกวิหิตํ อิทฺธิวิธํ ปจฺจนุโภติ, เอโกปิ หุตฺวา พหุธา โหติ, พหุธาปิ หุตฺวา เอโก โหติ, อาวิภาวํ ฯเปฯ ยาว พฺรหฺมโลกาปิ กาเยน วสํ วตฺเตติ. (5)}

\prose{ทิพฺพาย โสตธาตุยา วิสุทฺธาย อติกฺกนฺตมานุสิกาย อุโภ สทฺเท สุณาติ ทิพฺเพ จ มานุเส จ เย ทูเร สนฺติเก จ ฯเปฯ. (6)}

\prose{ปรสตฺตานํ ปรปุคฺคลานํ เจตสา เจโต ปริจฺจ ปชานาติ, สราคํ วา จิตฺตํ สราคํ จิตฺตนฺติ ปชานาติ, วีตราคํ วา จิตฺตํ ฯเปฯ สโทสํ วา จิตฺตํ.\csromanspacer{} วีตโทสํ วา จิตฺตํ.\csromanspacer{} สโมหํ วา จิตฺตํ.\csromanspacer{} วีตโมหํ วา จิตฺตํ.\csromanspacer{} สํขิตฺตํ วา จิตฺตํ.\csromanspacer{} วิกฺขิตฺตํ วา จิตฺตํ.\csromanspacer{} มหคฺคตํ วา จิตฺตํ.\csromanspacer{} อมหคฺคตํ วา จิตฺตํ.\csromanspacer{} ส-อุตฺตรํ วา จิตฺตํ.\csromanspacer{} อนุตฺตรํ วา จิตฺตํ.\csromanspacer{} สมาหิตํ วา จิตฺตํ.\csromanspacer{} อสมาหิตํ วา จิตฺตํ.\csromanspacer{} วิมุตฺตํ วา จิตฺตํ.\csromanspacer{} อวิมุตฺตํ วา จิตฺตํ อวิมุตฺตํ จิตฺตนฺติ ปชานาติ. (7)}

\prose{โส อเนกวิหิตํ ปุพฺเพนิวาสํ อนุสฺสรติ.\csromanspacer{} เสยฺยถิทํ, เอกมฺปิ ชาติํ เทฺวปิ ชาติโย ฯเปฯ อิติ สาการํ ส-อุทฺเทสํ อเนกวิหิตํ ปุพฺเพนิวาสํ อนุสฺสรติ. (8)}

\prose{\csromanfolio{140}ทิพฺเพน จกฺขุนา วิสุทฺเธน อติกฺกนฺตมานุสเกน สตฺเต ปสฺสติ จวมาเน อุปปชฺชมาเน หีเน ปณีเต สุวณฺเณ ทุพฺพณฺเณ สุคเต ทุคฺคเต, ยถากมฺมูปเค สตฺเต ปชานาติ. (9)}

\prose{อาสวานํ ขยา อนาสวํ เจโตวิมุตฺติํ ปญฺญาวิมุตฺติํ ทิฏฺเฐว ธมฺเม สยํ อภิญฺญา สจฺฉิกตฺวา อุปสมฺปชฺช วิหรติ. (10)}

\prose{กายคตาย ภิกฺขเว สติยา อาเสวิตาย ภาวิตาย พหุลีกตาย ยานีกตาย วตฺถุกตาย อนุฏฺฐิตาย ปริจิตาย สุสมารทฺธาย อิเม ทสานิสํสา ปาฏิกงฺขาติ.}

\prose{อิทมโวจ ภควา.\csromanspacer{} อตฺตมนา เต ภิกฺขู ภควโต ภาสิตํ อภินนฺทุนฺติ.}

\nitthitamruled{กายคตาสติสุตฺตํ นิฏฺฐิตํ นวมํ.}
\csromanmatikaanchor{mtk.23}
\csromantocmarkref{subsection}{10. สงฺขารุปปตฺติสุตฺต}{mtk.23}
\bookmark[dest={mtk.23},level=2]{10. สงฺขารุปปตฺติสุตฺต}
\csromanheader{2}{10. สงฺขารุปปตฺติสุตฺต}
\proseitem{160}{เอวํ เม สุตํ\csromandash{}เอกํ สมยํ ภควา สาวตฺถิยํ วิหรติ เชตวเน อนาถปิณฺฑิกสฺส อาราเม.\csromanspacer{} ตตฺร โข ภควา ภิกฺขู อามนฺเตสิ “ภิกฺขโว”ติ. “ภทนฺเต”ติ เต ภิกฺขู ภควโต ปจฺจสฺโสสุํ.\csromanspacer{} ภควา เอตทโวจ “สงฺขารุปปตฺติํ\footnote{สงฺขารูปปตฺติํ (สฺยา, กํ), สงฺขารุปฺปตฺติํ (สี, อิ)} โว ภิกฺขเว เทเสสฺสามิ, ตํ สุณาถ สาธุกํ มนสิ กโรถ, ภาสิสฺสามี”ติ. “เอวํ ภนฺเต”ติ โข เต ภิกฺขู ภควโต ปจฺจสฺโสสุํ.\csromanspacer{} ภควา เอตทโวจ\csromandash{}}

\proseitem{161}{อิธ ภิกฺขเว ภิกฺขุ สทฺธาย สมนฺนาคโต โหติ, สีเลน สมนฺนาคโต โหติ, สุเตน สมนฺนาคโต โหติ, จาเคน สมนฺนาคโต โหติ, ปญฺญาย สมนฺนาคโต โหติ.\csromanspacer{} ตสฺส เอวํ โหติ “อโห วตาหํ กายสฺส เภทา ปรํ มรณา ขตฺติยมหาสาลานํ\footnote{ขตฺติยมหาสาลานํ วา (สฺยา, กํ, อิ)} สหพฺยตํ อุปปชฺเชยฺยนฺ”ติ.\csromanspacer{} โส ตํ จิตฺตํ ทหติ, ตํ จิตฺตํ อธิฏฺฐาติ, ตํ จิตฺตํ ภาเวติ.\csromanspacer{} ตสฺส เต สงฺขารา จ วิหารา\footnote{วิหาโร (สี, อิ)} จ เอวํ ภาวิตา เอวํ}

\vspace{\tipitakaparskip}
\csromancenter{สงฺขารุปปตฺติสุตฺต (120) พหุลีกตา ตตฺรุปปตฺติยา\footnote{ตตฺรูปปตฺติยา (สฺยา, กํ), ตตฺรุปฺปตฺติยา (สี, อิ)} สํวตฺตนฺติ.\csromanspacer{} อยํ ภิกฺขเว มคฺโค อยํ ปฏิปทา ตตฺรุปปตฺติยา สํวตฺตติ.}
\proseitem{162}{\csromanfolio{141}ปุน จปรํ ภิกฺขเว ภิกฺขุ สทฺธาย สมนฺนาคโต โหติ, สีเลน สมนฺนาคโต โหติ, สุเตน สมนฺนาคโต โหติ, จาเคน สมนฺนาคโต โหติ, ปญฺญาย สมนฺนาคโต โหติ.\csromanspacer{} ตสฺส เอวํ โหติ “อโห วตาหํ กายสฺส เภทา ปรํ มรณา พฺราหฺมณมหาสาลานํ ฯเปฯ คหปติมหาสาลานํ\footnote{พฺราหฺมณมหาสาลานํ วา คหปติมหาสาลานํ วา (สฺยา, กํ, อิ)} สหพฺยตํ อุปปชฺเชยฺยนฺ”ติ.\csromanspacer{} โส ตํ จิตฺตํ ทหติ, ตํ จิตฺตํ อธิฏฺฐาติ, ตํ จิตฺตํ ภาเวติ.\csromanspacer{} ตสฺส เต สงฺขารา จ วิหารา จ เอวํ ภาวิตา เอวํ พหุลีกตา ตตฺรุปปตฺติยา สํวตฺตนฺติ.\csromanspacer{} อยํ ภิกฺขเว มคฺโค อยํ ปฏิปทา ตตฺรุปปตฺติยา สํวตฺตติ.}

\proseitem{163}{ปุน จปรํ ภิกฺขเว ภิกฺขุ สทฺธาย สมนฺนาคโต โหติ, สีเลน สมนฺนาคโต โหติ, สุเตน สมนฺนาคโต โหติ, จาเคน สมนฺนาคโต โหติ, ปญฺญาย สมนฺนาคโต โหติ.\csromanspacer{} ตสฺส สุตํ โหติ “จาตุมหาราชิกา\footnote{จาตุมฺมหาราชิกา (สี, สฺยา, กํ, อิ)} เทวา ทีฆายุกา วณฺณวนฺโต สุขพหุลา”ติ.\csromanspacer{} ตสฺส เอวํ โหติ “อโห วตาหํ กายสฺส เภทา ปรํ มรณา จาตุมหาราชิกานํ เทวานํ สหพฺยตํ อุปปชฺเชยฺยนฺ”ติ.\csromanspacer{} โส ตํ จิตฺตํ ทหติ, ตํ จิตฺตํ อธิฏฺฐาติ, ตํ จิตฺตํ ภาเวติ.\csromanspacer{} ตสฺส เต สงฺขารา จ วิหารา จ เอวํ ภาวิตา เอวํ พหุลีกตา ตตฺรุปปตฺติยา สํวตฺตนฺติ.\csromanspacer{} อยํ ภิกฺขเว มคฺโค อยํ ปฏิปทา ตตฺรุปปตฺติยา สํวตฺตติ.}

\proseitem{164}{ปุน จปรํ ภิกฺขเว ภิกฺขุ สทฺธาย สมนฺนาคโต โหติ, สีเลน สมนฺนาคโต โหติ, สุเตน สมนฺนาคโต โหติ, จาเคน สมนฺนาคโต โหติ, ปญฺญาย สมนฺนาคโต โหติ.\csromanspacer{} ตสฺส สุตํ โหติ “ตาวติํสา เทวา ฯเปฯ ยามา เทวา.\csromanspacer{} ตุสิตา เทวา.\csromanspacer{} นิมฺมานรตี เทวา.\csromanspacer{} ปรนิมฺมิตวสวตฺตี เทวา ทีฆายุกา วณฺณวนฺโต สุขพหุลา”ติ.\csromanspacer{} ตสฺส เอวํ โหติ “อโห วตาหํ กายสฺส เภทา ปรํ มรณา ปรนิมฺมิตวสวตฺตีนํ เทวานํ สหพฺยตํ อุปปชฺเชยฺยนฺ”ติ.\csromanspacer{} โส ตํ จิตฺตํ ทหติ, ตํ จิตฺตํ อธิฏฺฐาติ, ตํ จิตฺตํ ภาเวติ.\csromanspacer{} ตสฺส เต สงฺขารา จ วิหารา จ \csromanfolio{142}เอวํ ภาวิตา เอวํ พหุลีกตา ตตฺรุปปตฺติยา สํวตฺตนฺติ.\csromanspacer{} อยํ ภิกฺขเว มคฺโค อยํ ปฏิปทา ตตฺรุปปตฺติยา สํวตฺตติ.}

\proseitem{165}{ปุน จปรํ ภิกฺขเว ภิกฺขุ สทฺธาย สมนฺนาคโต โหติ, สีเลน สมนฺนาคโต โหติ, สุเตน สมนฺนาคโต โหติ, จาเคน สมนฺนาคโต โหติ, ปญฺญาย สมนฺนาคโต โหติ.\csromanspacer{} ตสฺส สุตํ โหติ “สหสฺโส พฺรหฺมา ทีฆายุโก วณฺณวา สุขพหุโล”ติ.\csromanspacer{} สหสฺโส ภิกฺขเว พฺรหฺมา สหสฺสิโลกธาตุํ\footnote{สหสฺสิํ โลกธาตุํ (สี)} ผริตฺวา อธิมุจฺจิตฺวา\footnote{อธิมุญฺจิตฺวา (ก)} วิหรติ.\csromanspacer{} เยปิ ตตฺถ สตฺตา อุปปนฺนา, เตปิ ผริตฺวา อธิมุจฺจิตฺวา วิหรติ.\csromanspacer{} เสยฺยถาปิ ภิกฺขเว จกฺขุมา ปุริโส เอกํ อามณฺฑํ หตฺเถ กริตฺวา ปจฺจเวกฺเขยฺย.\csromanspacer{} เอวเมว โข ภิกฺขเว สหสฺโส พฺรหฺมา สหสฺสิโลกธาตุํ ผริตฺวา อธิมุจฺจิตฺวา วิหรติ.\csromanspacer{} เยปิ ตตฺถ สตฺตา อุปปนฺนา, เตปิ ผริตฺวา อธิมุจฺจิตฺวา วิหรติ.\csromanspacer{} ตสฺส เอวํ โหติ “อโห วตาหํ กายสฺส เภทา ปรํ มรณา สหสฺสสฺส พฺรหฺมุโน สหพฺยตํ อุปปชฺเชยฺยนฺ”ติ.\csromanspacer{} โส ตํ จิตฺตํ ทหติ, ตํ จิตฺตํ อธิฏฺฐาติ, ตํ จิตฺตํ ภาเวติ.\csromanspacer{} ตสฺส เต สงฺขารา จ วิหารา จ เอวํ ภาวิตา เอวํ พหุลีกตา ตตฺรุปปตฺติยา สํวตฺตนฺติ.\csromanspacer{} อยํ ภิกฺขเว มคฺโค อยํ ปฏิปทา ตตฺรุปปตฺติยา สํวตฺตติ.}

\proseitem{166}{ปุน จปรํ ภิกฺขเว ภิกฺขุ สทฺธาย สมนฺนาคโต โหติ, สีเลน สมนฺนาคโต โหติ, สุเตน.\csromanspacer{} จาเคน.\csromanspacer{} ปญฺญาย สมนฺนาคโต โหติ.\csromanspacer{} ตสฺส สุตํ โหติ “ทฺวิสหสฺโส พฺรหฺมา ฯเปฯ ติสหสฺโส พฺรหฺมา.\csromanspacer{} จตุสหสฺโส พฺรหฺมา.\csromanspacer{} ปญฺจสหสฺโส พฺรหฺมา ทีฆายุโก วณฺณวา สุขพหุโล”ติ.\csromanspacer{} ปญฺจสหสฺโส ภิกฺขเว พฺรหฺมา ปญฺจสหสฺสิโลกธาตุํ ผริตฺวา อธิมุจฺจิตฺวา วิหรติ.\csromanspacer{} เยปิ ตตฺถ สตฺตา อุปปนฺนา, เตปิ ผริตฺวา อธิมุจฺจิตฺวา วิหรติ.\csromanspacer{} เสยฺยถาปิ ภิกฺขเว จกฺขุมา ปุริโส ปญฺจ อามณฺฑานิ หตฺเถ กริตฺวา ปจฺจเวกฺเขยฺย.\csromanspacer{} เอวเมว โข ภิกฺขเว ปญฺจสหสฺโส พฺรหฺมา ปญฺจสหสฺสิโลกธาตุํ ผริตฺวา อธิมุจฺจิตฺวา วิหรติ.\csromanspacer{} เยปิ ตตฺถ สตฺตา อุปปนฺนา, เตปิ ผริตฺวา อธิมุจฺจิตฺวา วิหรติ.\csromanspacer{} ตสฺส เอวํ โหติ “อโห วตาหํ กายสฺส เภทา ปรํ มรณา ปญฺจสหสฺสสฺส พฺรหฺมุโน สหพฺยตํ อุปปชฺเชยฺยนฺ”ติ.\csromanspacer{} โส ตํ จิตฺตํ ทหติ, ตํ จิตฺตํ อธิฏฺฐาติ, ตํ จิตฺตํ ภาเวติ.}

\vspace{\tipitakaparskip}
\csromancenter{สงฺขารุปปตฺติสุตฺต (120) ตสฺส เต สงฺขารา จ วิหารา จ เอวํ ภาวิตา เอวํ พหุลีกตา ตตฺรุปปตฺติยา สํวตฺตนฺติ.\csromanspacer{} อยํ ภิกฺขเว มคฺโค อยํ ปฏิปทา ตตฺรุปปตฺติยา สํวตฺตติ.}
\proseitem{167}{\csromanfolio{143}ปุน จปรํ ภิกฺขเว ภิกฺขุ สทฺธาย สมนฺนาคโต โหติ, สีเลน สมนฺนาคโต โหติ, สุเตน.\csromanspacer{} จาเคน.\csromanspacer{} ปญฺญาย สมนฺนาคโต โหติ.\csromanspacer{} ตสฺส สุตํ โหติ “ทสสหสฺโส พฺรหฺมา ทีฆายุโก วณฺณวา สุขพหุโล”ติ.\csromanspacer{} ทสสหสฺโส ภิกฺขเว พฺรหฺมา ทสสหสฺสิโลกธาตุํ ผริตฺวา อธิมุจฺจิตฺวา วิหรติ.\csromanspacer{} เยปิ ตตฺถ สตฺตา อุปปนฺนา, เตปิ ผริตฺวา อธิมุจฺจิตฺวา วิหรติ.\csromanspacer{} เสยฺยถาปิ ภิกฺขเว มณิ เวฬุริโย สุโภ ชาติมา อฏฺฐํโส สุปริกมฺมกโต ปณฺฑุกมฺพเล นิกฺขิตฺโต ภาสเต จ ตปเต จ\footnote{ภาสติ จ ตปติ จ (สี, สฺยา, กํ, อิ)} วิโรจติ จ.\csromanspacer{} เอวเมว โข ภิกฺขเว ทสสหสฺโส พฺรหฺมา ทสสหสฺสิโลกธาตุํ ผริตฺวา อธิมุจฺจิตฺวา วิหรติ.\csromanspacer{} เยปิ ตตฺถ สตฺตา อุปปนฺนา, เตปิ ผริตฺวา อธิมุจฺจิตฺวา วิหรติ.\csromanspacer{} ตสฺส เอวํ โหติ “อโห วตาหํ กายสฺส เภทา ปรํ มรณา ทสสหสฺสสฺส พฺรหฺมุโน สหพฺยตํ อุปปชฺเชยฺยนฺ”ติ.\csromanspacer{} โส ตํ จิตฺตํ ทหติ, ตํ จิตฺตํ อธิฏฺฐาติ, ตํ จิตฺตํ ภาเวติ.\csromanspacer{} ตสฺส เต สงฺขารา จ วิหารา จ เอวํ ภาวิตา เอวํ พหุลีกตา ตตฺรุปปตฺติยา สํวตฺตนฺติ.\csromanspacer{} อยํ ภิกฺขเว มคฺโค อยํ ปฏิปทา ตตฺรุปปตฺติยา สํวตฺตติ.}

\proseitem{168}{ปุน จปรํ ภิกฺขเว ภิกฺขุ สทฺธาย สมนฺนาคโต โหติ, สีเลน.\csromanspacer{} สุเตน.\csromanspacer{} จาเคน.\csromanspacer{} ปญฺญาย สมนฺนาคโต โหติ.\csromanspacer{} ตสฺส สุตํ โหติ “สตสหสฺโส พฺรหฺมา ทีฆายุโก วณฺณวา สุขพหุโล”ติ.\csromanspacer{} สตสหสฺโส ภิกฺขเว พฺรหฺมา สตสหสฺสิโลกธาตุํ ผริตฺวา อธิมุจฺจิตฺวา วิหรติ.\csromanspacer{} เยปิ ตตฺถ สตฺตา อุปปนฺนา, เตปิ ผริตฺวา อธิมุจฺจิตฺวา วิหรติ.\csromanspacer{} เสยฺยถาปิ ภิกฺขเว นิกฺขํ\footnote{เนกฺขํ (สี, สฺยา, กํ, อิ)} ชมฺโพนทํ ทกฺขกมฺมารปุตฺต- อุกฺกามุขสุกุสลสมฺปหฏฺฐํ ปณฺฑุกมฺพเล นิกฺขิตฺตํ ภาสเต จ ตปเต จ วิโรจติ จ.\csromanspacer{} เอวเมว โข ภิกฺขเว สตสหสฺโส พฺรหฺมา สตสหสฺสิโลกธาตุํ ผริตฺวา อธิมุจฺจิตฺวา วิหรติ.\csromanspacer{} เยปิ ตตฺถ สตฺตา อุปปนฺนา, เตปิ ผริตฺวา อธิมุจฺจิตฺวา วิหรติ.\csromanspacer{} ตสฺส เอวํ โหติ “อโห วตาหํ กายสฺส เภทา ปรํ มรณา สตสหสฺสสฺส พฺรหฺมุโน สหพฺยตํ \csromanfolio{144}อุปปชฺเชยฺยนฺ”ติ.\csromanspacer{} โส ตํ จิตฺตํ ทหติ, ตํ จิตฺตํ อธิฏฺฐาติ, ตํ จิตฺตํ ภาเวติ.\csromanspacer{} ตสฺส เต สงฺขารา จ วิหารา จ เอวํ ภาวิตา เอวํ พหุลีกตา ตตฺรุปปตฺติยา สํวตฺตนฺติ.\csromanspacer{} อยํ ภิกฺขเว มคฺโค อยํ ปฏิปทา ตตฺรุปปตฺติยา สํวตฺตติ.}

\proseitem{169}{ปุน จปรํ ภิกฺขเว ภิกฺขุ สทฺธาย สมนฺนาคโต โหติ, สีเลน.\csromanspacer{} สุเตน.\csromanspacer{} จาเคน.\csromanspacer{} ปญฺญาย สมนฺนาคโต โหติ.\csromanspacer{} ตสฺส สุตํ โหติ “อาภา เทวา ฯเปฯ ปริตฺตาภา เทวา.\csromanspacer{} อปฺปมาณาภา เทวา.\csromanspacer{} อาภสฺสรา เทวา ทีฆายุกา วณฺณวนฺโต สุขพหุลา”ติ.\csromanspacer{} ตสฺส เอวํ โหติ “อโห วตาหํ กายสฺส เภทา ปรํ มรณา อาภสฺสรานํ เทวานํ สหพฺยตํ อุปปชฺเชยฺยนฺ”ติ.\csromanspacer{} โส ตํ จิตฺตํ ทหติ, ตํ จิตฺตํ อธิฏฺฐาติ, ตํ จิตฺตํ ภาเวติ.\csromanspacer{} ตสฺส เต สงฺขารา จ วิหารา จ เอวํ ภาวิตา เอวํ พหุลีกตา ตตฺรุปปตฺติยา สํวตฺตนฺติ.\csromanspacer{} อยํ ภิกฺขเว มคฺโค อยํ ปฏิปทา ตตฺรุปปตฺติยา สํวตฺตติ.}

\proseitem{170}{ปุน จปรํ ภิกฺขเว ภิกฺขุ สทฺธาย สมนฺนาคโต โหติ, สีเลน.\csromanspacer{} สุเตน.\csromanspacer{} จาเคน.\csromanspacer{} ปญฺญาย สมนฺนาคโต โหติ.\csromanspacer{} ตสฺส สุตํ โหติ “ปริตฺตสุภา เทวา ฯเปฯ อปฺปมาณสุภา เทวา.\csromanspacer{} สุภกิณฺหา เทวา ทีฆายุกา วณฺณวนฺโต สุขพหุลา”ติ.\csromanspacer{} ตสฺส เอวํ โหติ “อโห วตาหํ กายสฺส เภทา ปรํ มรณา สุภกิณฺหา นํ เทวานํ สหพฺยตํ อุปปชฺเชยฺยนฺ”ติ.\csromanspacer{} โส ตํ จิตฺตํ ทหติ, ตํ จิตฺตํ อธิฏฺฐาติ, ตํ จิตฺตํ ภาเวติ.\csromanspacer{} ตสฺส เต สงฺขารา จ วิหารา จ เอวํ ภาวิตา เอวํ พหุลีกตา ตตฺรุปปตฺติยา สํวตฺตติ.\csromanspacer{} อยํ ภิกฺขเว มคฺโค อยํ ปฏิปทา ตตฺรุปปตฺติยา สํวตฺตติ.}

\proseitem{171}{ปุน จปรํ ภิกฺขเว ภิกฺขุ สทฺธาย สมนฺนาคโต โหติ, สีเลน.\csromanspacer{} สุเตน.\csromanspacer{} จาเคน.\csromanspacer{} ปญฺญาย สมนฺนาคโต โหติ.\csromanspacer{} ตสฺส สุตํ โหติ “เวหปฺผลา เทวา ฯเปฯ อวิหา เทวา.\csromanspacer{} อตปฺปา เทวา.\csromanspacer{} สุทสฺสา เทวา.\csromanspacer{} สุทสฺสี เทวา.\csromanspacer{} อกนิฏฺฐา เทวา ทีฆายุกา วณฺณวนฺโต สุขพหุลา”ติ.\csromanspacer{} ตสฺส เอวํ โหติ “อโห วตาหํ กายสฺส เภทา ปรํ มรณา อกนิฏฺฐานํ เทวานํ สหพฺยตํ อุปปชฺเชยฺยนฺ”ติ.\csromanspacer{} โส ตํ จิตฺตํ ทหติ, ตํ จิตฺตํ ภาเวติ.\csromanspacer{} ตสฺส เต สงฺขารา จ}

\vspace{\tipitakaparskip}
\csromancenter{สงฺขารุปปตฺติสุตฺต (120) วิหารา จ เอวํ ภาวิตา เอวํ พหุลีกตา ตตฺรุปปตฺติยา สํวตฺตนฺติ.\csromanspacer{} อยํ ภิกฺขเว มคฺโค อยํ ปฏิปทา ตตฺรุปปตฺติยา สํวตฺตติ.}
\proseitem{172}{\csromanfolio{145}ปุน จปรํ ภิกฺขเว ภิกฺขุ สทฺธาย สมนฺนาคโต โหติ, สีเลน.\csromanspacer{} สุเตน.\csromanspacer{} จาเคน.\csromanspacer{} ปญฺญาย สมนฺนาคโต โหติ.\csromanspacer{} ตสฺส สุตํ โหติ “อากาสานญฺจายตนูปคา เทวา ทีฆายุกา จิรฏฺฐิติกา สุขพหุลา”ติ.\csromanspacer{} ตสฺส เอวํ โหติ “อโห วตาหํ กายสฺส เภทา ปรํ มรณา อากาสานญฺจายตนูปคานํ เทวานํ สหพฺยตํ อุปปชฺเชยฺยนฺ”ติ.\csromanspacer{} โส ตํ จิตฺตํ ทหติ, ตํ จิตฺตํ อธิฏฺฐาติ, ตํ จิตฺตํ ภาเวติ.\csromanspacer{} ตสฺส เต สงฺขารา จ วิหารา จ เอวํ ภาวิตา เอวํ พหุลีกตา ตตฺรุปปตฺติยา สํวตฺตนฺติ.\csromanspacer{} อยํ ภิกฺขเว มคฺโค อยํ ปฏิปทา ตตฺรุปปตฺติยา สํวตฺตติ.}

\proseitem{173}{ปุน จปรํ ภิกฺขเว ภิกฺขุ สทฺธาย สมนฺนาคโต โหติ, สีเลน.\csromanspacer{} สุเตน.\csromanspacer{} จาเคน.\csromanspacer{} ปญฺญาย สมนฺนาคโต โหติ.\csromanspacer{} ตสฺส สุตํ โหติ “วิญฺญาณญฺจายตนูปคา เทวา ทีฆายุกา จิรฏฺฐิติกา สุขพหุลา”ติ.\csromanspacer{} ตสฺส เอวํ โหติ “อโห วตาหํ กายสฺส เภทา ปรํ มรณา วิญฺญาณญฺจายตนูปคานํ เทวานํ สหพฺยตํ อุปปชฺเชยฺยนฺ”ติ.\csromanspacer{} โส ตํ จิตฺตํ ทหติ, ตํ จิตฺตํ อธิฏฺฐาติ, ตํ จิตฺตํ ภาเวติ.\csromanspacer{} ตสฺส เต สงฺขารา จ วิหารา จ เอวํ ภาวิตา เอวํ พหุลีกตา ตตฺรุปปตฺติยา สํวตฺตนฺติ.\csromanspacer{} อยํ ภิกฺขเว มคฺโค อยํ ปฏิปทา ตตฺรุปปตฺติยา สํวตฺตติ.}

\proseitem{174}{ปุน จปรํ ภิกฺขเว ภิกฺขุ สทฺธาย สมนฺนาคโต โหติ, สีเลน.\csromanspacer{} สุเตน.\csromanspacer{} จาเคน.\csromanspacer{} ปญฺญาย สมนฺนาคโต โหติ.\csromanspacer{} ตสฺส สุตํ โหติ “อากิญฺจญฺญายตนูปคา เทวา ฯเปฯ เนวสญฺญานาสญฺญายตนูปคา เทวา ทีฆายุกา จิรฏฺฐิติกา สุขพหุลา”ติ.\csromanspacer{} ตสฺส เอวํ โหติ “อโห วตาหํ กายสฺส เภทา ปรํ มรณา เนวสญฺญานาสญฺญายตนูปคานํ เทวานํ สหพฺยตํ อุปปชฺเชยฺยนฺ”ติ.\csromanspacer{} โส ตํ จิตฺตํ ทหติ, ตํ จิตฺตํ อธิฏฺฐาติ, ตํ จิตฺตํ ภาเวติ.\csromanspacer{} ตสฺส เต สงฺขารา จ วิหารา จ เอวํ ภาวิตา เอวํ พหุลีกตา ตตฺรุปปตฺติยา สํวตฺตนฺติ.\csromanspacer{} อยํ ภิกฺขเว มคฺโค อยํ ปฏิปทา ตตฺรุปปตฺติยา สํวตฺตติ.}

\csromanmatikaanchor{mtk.24}
\csromantocmarkref{subsection}{อุทฺทานคาถา}{mtk.24}
\bookmark[dest={mtk.24},level=2]{อุทฺทานคาถา}
\proseitem{175}{ปุน จปรํ ภิกฺขเว ภิกฺขุ สทฺธาย สมนฺนาคโต โหติ, สีเลน.\csromanspacer{} สุเตน.\csromanspacer{} จาเคน.\csromanspacer{} ปญฺญาย สมนฺนาคโต โหติ.\csromanspacer{} ตสฺส เอวํ \csromanfolio{146}โหติ “อโห วตาหํ อาสวานํ ขยา อนาสวํ เจโตวิมุตฺติํ ปญฺญาวิมุตฺติํ ทิฏฺเฐว ธมฺเม สยํ อภิญฺญา สจฺฉิกตฺวา อุปสมฺปชฺช วิหเรยฺยนฺ”ติ.\csromanspacer{} โส อาสวานํ ขยา อนาสวํ เจโตวิมุตฺติํ ปญฺญาวิมุตฺติํ ทิฏฺเฐว ธมฺเม สยํ อภิญฺญา สจฺฉิกตฺวา อุปสมฺปชฺช วิหรติ.\csromanspacer{} อยํ ภิกฺขเว ภิกฺขุ น กตฺถจิ อุปปชฺชตีติ\footnote{น กตฺถจิ อุปปชฺชติ, น กุหิญฺจิ อุปปชฺชตีติ (สี, อิ), น กตฺถจิ อุปปชฺชติ, น กุหิญฺจิ อุปสมฺปชฺช วิหรตีติ. (ก)}.}

\prose{อิทมโวจ ภควา.\csromanspacer{} อตฺตมนา เต ภิกฺขู ภควโต ภาสิตํ อภินนฺทุนฺติ.}

\nitthitam{สงฺขารุปปตฺติสุตฺตํ นิฏฺฐิตํ ทสมํ.}
\nitthitam{\textbf{อนุปทวคฺโค นิฏฺฐิโต ทุติโย.}}
\titlehead{\textbf{ตสฺสุทฺทานํ}}
\csromangathameasure{อนุปาทโสธนโปริสธมฺโม, \\ เสวิตพฺพพหุธาตุวิภตฺติ. \\ พุทฺธสฺส กิตฺตินาม จตฺตารีเสน, \\ อานาปาโน กายคโต อุปปตฺติ.}
\csromangathagroup{\csromangathastanza{อนุปาทโสธนโปริสธมฺโม, \\ เสวิตพฺพพหุธาตุวิภตฺติ. \\ พุทฺธสฺส กิตฺตินาม จตฺตารีเสน, \\ อานาปาโน กายคโต อุปปตฺติ\footnote{อิโต ปรํ สฺยา-กํ-ก-โปตฺถเกสุ เอวมฺปิ ทิสฺสติ\csromandash{} จนฺทเก วิมเล ปริสุทฺเธ, ปุณฺณสมฺโมทินิโรธ-อตฺตโน. ทนฺธา พหุชนเสวิตํ ธมฺมวรํ, ยํ อนุปทํ วคฺควรํ ทุติยาติ.}.}}

\csromanmatikaanchor{mtk.25}
\csromantocmarknopage{chapter}{3. สุญฺญตวคฺค}{mtk.25}
\bookmark[dest={mtk.25},level=0]{3. สุญฺญตวคฺค}
\chapterheadpageread{3. สุญฺญตวคฺค}
\csromanmatikaanchor{mtk.26}
\csromantocmarkref{subsection}{1. จูฬสุญฺญตสุตฺต}{mtk.26}
\bookmark[dest={mtk.26},level=2]{1. จูฬสุญฺญตสุตฺต}
\csromanheader{2}{1. จูฬสุญฺญตสุตฺต}
\proseitem{176}{\csromanfolio{147}เอวํ เม สุตํ\csromandash{}เอกํ สมยํ ภควา สาวตฺถิยํ วิหรติ ปุพฺพาราเม มิคารมาตุปาสาเท.\csromanspacer{} อถ โข อายสฺมา อานนฺโท สายนฺหสมยํ ปฏิสลฺลานา วุฏฺฐิโต เยน ภควา เตนุปสงฺกมิ, อุปสงฺกมิตฺวา ภควนฺตํ อภิวาเทตฺวา เอกมนฺตํ นิสีทิ, เอกมนฺตํ นิสินฺโน โข อายสฺมา อานนฺโท ภควนฺตํ เอตทโวจ “เอกมิทํ ภนฺเต สมยํ ภควา สกฺเกสุ วิหรติ นครกํ นาม สกฺยานํ นิคโม, ตตฺถ เม ภนฺเต ภควโต สมฺมุขา สุตํ สมฺมุขา ปฏิคฺคหิตํ ‘สุญฺญตาวิหาเรนาหํ อานนฺท เอตรหิ พหุลํ วิหรามี’ติ, กจฺจิ เมตํ ภนฺเต สุสฺสุตํ สุคฺคหิตํ สุมนสิกตํ สูปธาริตนฺ”ติ.\csromanspacer{} ตคฺฆ เต เอตํ อานนฺท สุสฺสุตํ สุคฺคหิตํ สุมนสิกตํ สูปธาริตํ ปุพฺเพปาหํ\footnote{ปุพฺเพจาหํ (สี, สฺยา, กํ, อิ)} อานนฺท เอตรหิปิ\footnote{เอตรหิ จ (สพฺพตฺถ)} สุญฺญตาวิหาเรน พหุลํ วิหรามิ.\csromanspacer{} เสยฺยถาปิ อานนฺท อยํ มิคารมาตุปาสาโท สุญฺโญ หตฺถิควสฺสวฬเวน, สุญฺโญ ชาตรูปรชเตน, สุญฺโญ อิตฺถิปุริสสนฺนิปาเตน.\csromanspacer{} อตฺถิ เจวิทํ อสุญฺญตํ, ยทิทํ ภิกฺขุสํฆํ ปฏิจฺจ เอกตฺตํ.\csromanspacer{} เอวเมว โข อานนฺท ภิกฺขุ อมนสิกริตฺวา คามสญฺญํ อมนสิกริตฺวา มนุสฺสสญฺญํ อรญฺญสญฺญํ ปฏิจฺจ มนสิ กโรติ เอกตฺตํ, ตสฺส อรญฺญสญฺญาย จิตฺตํ ปกฺขนฺทติ ปสีทติ สนฺติฏฺฐติ อธิมุจฺจติ.\csromanspacer{} โส เอวํ ปชานาติ “เย อสฺสุ ทรถา คามสญฺญํ ปฏิจฺจ, เตธ น สนฺติ.\csromanspacer{} เย อสฺสุ ทรถา มนุสฺสสญฺญํ ปฏิจฺจ, เตธ น สนฺติ.\csromanspacer{} อตฺถิ เจวายํ ทรถมตฺตา, ยทิทํ อรญฺญสญฺญํ ปฏิจฺจ เอกตฺตนฺ”ติ.\csromanspacer{} โส “สุญฺญมิทํ สญฺญาคตํ คามสญฺญายา”ติ ปชานาติ, “สุญฺญมิทํ สญฺญาคตํ มนุสฺสสญฺญายา”ติ ปชานาติ.\csromanspacer{} อตฺถิ เจวิทํ อสุญฺญตํ, ยทิทํ อรญฺญสญฺญํ ปฏิจฺจ เอกตฺตนฺติ.\csromanspacer{} อิติ ยํ หิ โข ตตฺถ น โหติ, เตน ตํ สุญฺญํ สมนุปสฺสติ.\csromanspacer{} ยํ ปน ตตฺถ อวสิฏฺฐํ โหติ, ตํ “สนฺตมิทํ อตฺถี”ติ ปชานาติ.\csromanspacer{} เอวมฺปิสฺส เอสา อานนฺท ยถาภุจฺจา อวิปลฺลตฺถา ปริสุทฺธา สุญฺญตาวกฺกนฺติ ภวติ.}

\proseitem{177}{ปุน จปรํ อานนฺท ภิกฺขุ อมนสิกริตฺวา มนุสฺสสญฺญํ อมนสิกริตฺวา อรญฺญสญฺญํ ปถวีสญฺญํ ปฏิจฺจ มนสิ กโรติ เอกตฺตํ, ตสฺส \csromanfolio{148}ปถวีสญฺญาย จิตฺตํ ปกฺขนฺทติ ปสีทติ สนฺติฏฺฐติ อธิมุจฺจติ.\csromanspacer{} เสยฺยถาปิ อานนฺท อาสภจมฺมํ สงฺกุสเตน สุวิหตํ วิคตวลิกํ.\csromanspacer{} เอวเมว โข อานนฺท ภิกฺขุ ยํ อิมิสฺสา ปถวิยา อุกฺกูลวิกฺกูลํ นทีวิทุคฺคํ ขาณุกณฺฏกฏฺฐานํ ปพฺพตวิสมํ, ตํ สพฺพํ\footnote{สพฺพํ (ก)} อมนสิกริตฺวา ปถวีสญฺญํ ปฏิจฺจ มนสิ กโรติ เอกตฺตํ, ตสฺส ปถวีสญฺญาย จิตฺตํ ปกฺขนฺทติ ปสีทติ สนฺติฏฺฐติ อธิมุจฺจติ.\csromanspacer{} โส เอวํ ปชานาติ “เย อสฺสุ ทรถา มนุสฺสสญฺญํ ปฏิจฺจ, เตธ น สนฺติ.\csromanspacer{} เย อสฺสุ ทรถา อรญฺญสญฺญํ ปฏิจฺจ, เตธ น สนฺติ.\csromanspacer{} อตฺถิ เจวายํ ทรถมตฺตา, ยทิทํ ปถวีสญฺญํ ปฏิจฺจ เอกตฺตนฺ”ติ.\csromanspacer{} โส “สุญฺญมิทํ สญฺญาคตํ มนุสฺสสญฺญายา”ติ ปชานาติ, “สุญฺญมิทํ สญฺญาคตํ อรญฺญาสญฺญายา”ติ ปชานาติ.\csromanspacer{} อตฺถิ เจวิทํ อสุญฺญตํ, ยทิทํ ปถวีสญฺญํ ปฏิจฺจ เอกตฺตนฺติ.\csromanspacer{} อิติ ยํ หิ โข ตตฺถ น โหติ, เตน ตํ สุญฺญํ สมนุปสฺสติ.\csromanspacer{} ยํ ปน ตตฺถ อวสิฏฺฐํ โหติ, ตํ “สนฺตมิทํ อตฺถี”ติ ปชานาติ.\csromanspacer{} เอวมฺปิสฺส เอสา อานนฺท ยถาภุจฺจา อวิปลฺลตฺถา ปริสุทฺธา สุญฺญตาวกฺกนฺติ ภวนฺติ.}

\proseitem{178}{ปุน จปรํ อานนฺท ภิกฺขุ อมนสิกริตฺวา อรญฺญสญฺญํ อมนสิกริตฺวา ปถวีสญฺญํ อากาสานญฺจายตนสญฺญํ ปฏิจฺจ มนสิ กโรติ เอกตฺตํ, ตสฺส อากาสานญฺจายตนสญฺญาย จิตฺตํ ปกฺขนฺทติ ปสีทติ สนฺติฏฺฐติ อธิมุจฺจติ.\csromanspacer{} โส เอวํ ปชานาติ “เย อสฺสุ ทรถา อรญฺญสญฺญํ ปฏิจฺจ, เตธ สนฺติ.\csromanspacer{} เย อสฺสุ ทรถา ปถวีสญฺญํ ปฏิจฺจ, เตธ น สนฺติ.\csromanspacer{} อตฺถิ เจวายํ ทรถมตฺตา, ยทิทํ อากาสานญฺจายตนสญฺญํ ปฏิจฺจ เอกตฺตนฺ”ติ.\csromanspacer{} โส “สุญฺญมิทํ สญฺญาคตํ อรญฺญสญฺญายา”ติ ปชานาติ, “สุญฺญมิทํ สญฺญาคตํ ปถวีสญฺญายา”ติ ปชานาติ.\csromanspacer{} อตฺถิ เจวิทํ อสุญฺญตํ, ยทิทํ อากาสานญฺจายตนสญฺญํ ปฏิจฺจ เอกตฺตนฺติ.\csromanspacer{} อิติ ยํ หิ โข ตตฺถ น โหติ, เตน ตํ สุญฺญํ สมนุปสฺสติ.\csromanspacer{} ยํ ปน ตตฺถ อวสิฏฺฐํ โหติ, ตํ “สนฺตมิทํ อตฺถี”ติ ปชานาติ.\csromanspacer{} เอวมฺปิสฺส เอสา อานนฺท ยถาภุจฺจา อวิปลฺลตฺถา ปริสุทฺธา สุญฺญตาวกฺกนฺติ ภวติ.}

\proseitem{179}{ปุน จปรํ อานนฺท ภิกฺขุ อมนสิกริตฺวา ปถวีสญฺญํ อมนสิกริตฺวา อากาสานญฺจายตนสญฺญํ วิญฺญาณญฺจายตนสญฺญํ ปฏิจฺจ มนสิ กโรติ เอกตฺตํ, ตสฺส วิญฺญาณญฺจายตนสญฺญาย จิตฺตํ ปกฺขนฺทติ ปสีทติ สนฺติฏฺฐติ อธิมุจฺจติ.\csromanspacer{} โส เอวํ ปชานาติ “เย อสฺสุ ทรถา}

\csromanheader{2}{2. จูฬสุญฺญตสุตฺต (121)}
\prose{\csromanfolio{149}ปถวีสญฺญํ ปฏิจฺจ, เตธ น สนฺติ.\csromanspacer{} เย อสฺสุ ทรถา อากาสานญฺจายตนสญฺญํ ปฏิจฺจ, เตธ น สนฺติ.\csromanspacer{} อตฺถิ เจวายํ ทรถมตฺตา, ยทิทํ วิญฺญาณญฺจายตนสญฺญํ ปฏิจฺจ เอกตฺตนฺ”ติ.\csromanspacer{} โส “สุญฺญมิทํ สญฺญาคตํ ปถวีสญฺญายา”ติ ปชานาติ, “สุญฺญมิทํ สญฺญาคตํ อากาสานญฺจายตนสญฺญายา”ติ ปชานาติ.\csromanspacer{} อตฺถิ เจวิทํ อสุญฺญตํ, ยทิทํ วิญฺญาณญฺจายตนสญฺญํ ปฏิจฺจ เอกตฺตนฺติ.\csromanspacer{} อิติ ยํ หิ โข ตตฺถ น โหติ, เตน ตํ สุญฺญํ สมนุปสฺสติ.\csromanspacer{} ยํ ปน ตตฺถ อวสิฏฺฐํ โหติ, ตํ “สนฺตมิทํ อตฺถี”ติ ปชานาติ.\csromanspacer{} เอวมฺปิสฺส เอสา อานนฺท ยถาภุจฺจา อวิปลฺลตฺถา ปริสุทฺธา สุญฺญตาวกฺกนฺติ ภวติ.}

\proseitem{180}{ปุน จปรํ อานนฺท ภิกฺขุ อมนสิกริตฺวา อากาสานญฺจายตนสญฺญํ อมนสิกริตฺวา วิญฺญาณญฺจายตนสญฺญํ อากิญฺจญฺญายตนสญฺญํ ปฏิจฺจ มนสิ กโรติ เอกตฺตํ, ตสฺส อากิญฺจญฺญายตนสญฺญาย จิตฺตํ ปกฺขนฺทติ ปสีทติ สนฺติฏฺฐติ อธิมุจฺจติ.\csromanspacer{} โส เอวํ ปชานาติ “เย อสฺสุ ทรถา อากาสานญฺจายตนสญฺญํ ปฏิจฺจ, เตธ น สนฺติ.\csromanspacer{} เย อสฺสุ ทรถา วิญฺญาณญฺจายตนสญฺญํ ปฏิจฺจ, เตธ น สนฺติ.\csromanspacer{} อตฺถิ เจวายํ ทรถมตฺตา, ยทิทํ อากิญฺจยตนสญฺญํ ปฏิจฺจ เอกตฺตนฺ”ติ.\csromanspacer{} โส “สุญฺญมิทํ สญฺญาคตํ อากาสานญฺจายตนสญฺญายา”ติ ปชานาติ, “สุญฺญมิทํ สญฺญาคตํ วิญฺญาณญฺจายตนสญฺญายา”ติ ปชานาติ.\csromanspacer{} อตฺถิ เจวิทํ อสุญฺญตํ, ยทิทํ อากิญฺจญฺญายตนสญฺญํ ปฏิจฺจ เอกตฺตนฺติ.\csromanspacer{} อิติ ยํ หิ โข ตตฺถ น โหติ, เตน ตํ สุญฺญํ สมนุปสฺสติ.\csromanspacer{} ยํ ปน ตตฺต อวสิฏฺฐํ โหติ, ตํ “สนฺตมิทํ อตฺถี”ติ ปชานาติ.\csromanspacer{} เอวมฺปิสฺส เอสา อานนฺท ยถาภุจฺจา อวิปลฺลตฺถา ปริสุทฺธา สุญฺญตาวกฺกนฺติ ภวติ.}

\proseitem{181}{ปุน จปรํ อานนฺท ภิกฺขุ อมนสิกริตฺวา วิญฺญาณญฺจายตนสญฺญํ อมนสิกริตฺวา อากิญฺจญฺญายตนสญฺญํ เนวสญฺญานาสญฺญายตนสญฺญํ ปฏิจฺจ มนสิ กโรติ เอกตฺตํ, ตสฺส เนวสญฺญานาสญฺญายตสญฺญาย จิตฺตํ ปกฺขนฺทติ ปสีทติ สนฺติฏฺฐติ อธิมุจฺจติ.\csromanspacer{} โส เอวํ ปชานาติ “เย อสฺสุ ทรถา วิญฺญาณญฺจายตนสญฺญํ ปฏิจฺจ, เตธ น สนฺติ.\csromanspacer{} เย อสฺสุ ทรถา อากิญฺจญฺญายตนสญฺญํ ปฏิจฺจ, เตธ น สนฺติ.\csromanspacer{} อตฺถิ เจวายํ ทรถมตฺตา, ยทิทํ เนวสญฺญานาสญฺญายตนสญฺญํ ปฏิจฺจ เอกตฺตนฺ”ติ.\csromanspacer{} โส “สุญฺญมิทํ สญฺญาคตํ วิญฺญาณญฺจายตนสญฺญายา”ติ ปชานาติ, “สุญฺญมิทํ สญฺญาคตํ อากิญฺจญฺญายตนสญฺญายา”ติ ปชานาติ.\csromanspacer{} อตฺถิ เจวิทํ \csromanfolio{150}อสุญฺญตํ, ยทิทํ เนวสญฺญานาสญฺญายตนสญฺญํ ปฏิจฺจ เอกตฺตนฺติ.\csromanspacer{} อิติ ยํ หิ โข ตตฺถ น โหติ, เตน ตํ สุญฺญํ สมนุปสฺสติ.\csromanspacer{} ยํ ปน ตตฺถ อวสิฏฺฐํ โหติ, ตํ “สนฺตมิทํ อตฺถี”ติ ปชานาติ.\csromanspacer{} เอวมฺปิสฺส เอสา อานนฺท ยถาภุจฺจา อวิปลฺลตฺถา ปริสุทฺธา สุญฺญตาวกฺกนฺติ ภวติ.}

\proseitem{182}{ปุน จปรํ อานนฺท ภิกฺขุ อมนสิกริตฺวา อากิญฺจญฺญายตนสญฺญํ อมนสิกริตฺวา เนวสญฺญานาสญฺญายตนสญฺญํ อนิมิตฺตํ เจโตสมาธิํ ปฏิจฺจ มนสิ กโรติ เอกตฺตํ, ตสฺส อนิมิตฺเต เจโตสมาธิมฺหิ จิตฺตํ ปกฺขนฺทติ ปสีทติ สนฺติฏฺฐติ อธิมุจฺจติ.\csromanspacer{} โส เอวํ ปชานาติ “เย อสฺสุ ทรถา อากิญฺจญฺญายตนสญฺญํ ปฏิจฺจ, เตธ น สนฺติ.\csromanspacer{} เย อสฺสุ ทรถา เนวสญฺญานาสญฺญายตนสญฺญํ ปฏิจฺจ, เตธ น สนฺติ.\csromanspacer{} อตฺถิ เจวายํ ทรถามตฺตา, ยทิทํ อิมเมว กายํ ปฏิจฺจ สฬายตนิกํ ชีวิตปจฺจยา”ติ.\csromanspacer{} โส “สุญฺญมิทํ สญฺญาคตํ อากิญฺจญฺญายตนสญฺญายา”ติ ปชานาติ, “สุญฺญมิทํ สญฺญาคตํ เนวสญฺญานาสญฺญายตนสญฺญายา”ติ ปชานาติ.\csromanspacer{} อตฺถิ เจวิทํ อสุญฺญตํ, ยทิทํ อิมเมว กายํ ปฏิจฺจ สฬายตนิกํ ชีวิตปจฺจยาติ.\csromanspacer{} อิติ ยํ หิ โข ตตฺถ น โหติ, เตน ตํ สุญฺญํ สมนุปสฺสติ.\csromanspacer{} ยํ ปน ตตฺถ อวสิฏฺฐํ โหติ, ตํ “สนฺตมิทํ อตฺถี”ติ ปชานาติ.\csromanspacer{} เอวมฺปิสฺส เอสา อานนฺท ยถาภุจฺจา อวิปลฺลตฺถา ปริสุทฺธา สุญฺญตาวกฺกนฺติ ภวติ.}

\proseitem{183}{ปุน จปรํ อานนฺท ภิกฺขุ อมนสิกริตฺวา อากิญฺจญฺญายตนสญฺญํ อมนสิกริตฺวา เนวสญฺญานาสญฺญายตนสญฺญํ อนิมิตฺตํ เจโตสมาธิํ ปฏิจฺจ มนสิ กโรติ เอกตฺตํ, ตสฺส อนิมิตฺเต เจโตสมาธิมฺหิ จิตฺตํ ปกฺขนฺทติ ปสีทติ สนฺติฏฺฐติ อธิมุจฺจติ.\csromanspacer{} โส เอวํ ปชานาติ “อยมฺปิ โข อนิมิตฺโต เจโตสมาธิ อภิสงฺขโต อภิสญฺเจตยิโต, ยํ โข ปน กิญฺจิ อภิสงฺขตํ อภิสญฺเจตยิตํ, ตทนิจฺจํ นิโรธธมฺมนฺ”ติ ปชานาติ, ตสฺส เอวํ ชานโต เอวํ ปสฺสโต กามาสวาปิ จิตฺตํ วิมุจฺจติ, ภวาสวาปิ จิตฺตํ วิมุจฺจติ, อวิชฺชาสวาปิ จิตฺตํ วิมุจฺจติ, วิมุตฺตสฺมิํ “วิมุตฺตมฺ”อิติ ญาณํ โหติ, “ขีณา ชาติ, วุสิตํ พฺรหฺมจริยํ, กตํ กรณียํ, นาปรํ อิตฺถตฺตายา”ติ ปชานาติ.\csromanspacer{} โส เอวํ ปชานาติ “เย อสฺสุ ทรถา กามาสวํ ปฏิจฺจ, เตธ น สนฺติ.\csromanspacer{} เย อสฺสุ ทรถา ภวาสวํ ปฏิจฺจ, เตธ น สนฺติ.\csromanspacer{} เย อสฺสุ ทรถา อวิชฺชาสวํ ปฏิจฺจ, เตธ น สนฺติ.\csromanspacer{} อตฺถิ เจวายํ ทรถมตฺตา, ยทิทํ อิมเมว กายํ ปฏิจฺจ}

\csromanheader{2}{2. มหาสุญฺญตสุตฺต (122)}
\prose{\csromanfolio{151}สฬายตนิกํ ชีวิตปจฺจยา”ติ.\csromanspacer{} โส “สุญฺญมิทํ สญฺญาคตํ กามาสเวนา”ติ ปชานาติ, “สุญฺญมิทํ สญฺญาคตํ ภวาสเวนา”ติ ปชานาติ, “สุญฺญมิทํ สญฺญาคตํ อวิชฺชาสเวนา”ติ ปชานาติ, อตฺถิ เจวิทํ อสุญฺญตํ, ยทิทํ อิมเมว กายํ ปฏิจฺจ สฬายตนิกํ ชีวิตปจฺจยาติ.\csromanspacer{} อิติ ยํ หิ โข ตตฺถ น โหติ, เตน ตํ สุญฺญํ สมนุปสฺสติ.\csromanspacer{} ยํ ปน ตตฺถ อวสิฏฺฐํ โหติ, ตํ “สนฺตมิทํ อตฺถี”ติ ปชานาติ.\csromanspacer{} เอวมฺปิสฺส เอสา อานนฺท ยถาภุจฺจา อวิปลฺลตฺถา ปริสุทฺธา ปรมานุตฺตรา สุญฺญตาวกฺกนฺติ ภวติ.}

\proseitem{184}{เยปิ\footnote{เย (สี, อิ)} หิ เกจิ อานนฺท อตีตมทฺธานํ สมณา วา พฺราหฺมณา วา ปริสุทฺธํ ปรมานุตฺตรํ สุญฺญตํ อุปสมฺปชฺช วิหริํสุ, สพฺเพ เต อิมํเยว ปริสุทฺธํ ปรมานุตฺตรํ สุญฺญตํ อุปสมฺปชฺช วิหริํสุ.\csromanspacer{} เยปิ1 หิ เกจิ อานนฺท อนาคตมทฺธานํ สมณา วา พฺราหฺมณา วา ปริสุทฺธํ ปรมานุตฺตรํ สุญฺญตํ อุปสมฺปชฺช วิหริสฺสนฺติ, สพฺเพ เต อิมํเยว ปริสุทฺธํ ปรมานุตฺตรํ สุญฺญตํ อุปสมฺปชฺช วิหริสฺสนฺติ.\csromanspacer{} เยปิ1 หิ เกจิ อานนฺท เอตรหิ สมณา วา พฺราหฺมณา วา ปริสุทฺธํ ปรมานุตฺตรํ สุญฺญตํ อุปสมฺปชฺช วิหรนฺติ, สพฺเพ เต อิมํเยว ปริสุทฺธํ ปรมานุตฺตรํ สุญฺญตํ อุปสมฺปชฺช วิหรนฺติ.\csromanspacer{} ตสฺมาติห อานนฺท ปริสุทฺธํ ปรมานุตฺตรํ สุญฺญตํ อุปสมฺปชฺช วิหริสฺสามาติ\footnote{วิหริสฺสามีติ (อิ, ก)} เอวํ หิ โว\footnote{เต (ก)} อานนฺท สิกฺขิตพฺพนฺติ.}

\prose{อิทมโวจ ภควา.\csromanspacer{} อตฺตมโน อายสฺมา อานนฺโท ภควโต ภาสิตํ อภินนฺทีติ.}

\nitthitamruled{จูฬสุญฺญตสุตฺตํ นิฏฺฐิตํ ปฐมํ.}
\csromanmatikaanchor{mtk.27}
\csromantocmarkref{subsection}{2. มหาสุญฺญตสุตฺต}{mtk.27}
\bookmark[dest={mtk.27},level=2]{2. มหาสุญฺญตสุตฺต}
\csromanheader{2}{2. มหาสุญฺญตสุตฺต}
\proseitem{185}{เอวํ เม สุตํ\csromandash{}เอกํ สมยํ ภควา สกฺเกสุ วิหรติ กปิลวตฺถุสฺมิํ นิโคฺรธาราเม.\csromanspacer{} อถ โข ภควา ปุพฺพณฺหสมยํ นิวาเสตฺวา ปตฺตจีวรมาทาย กปิลวตฺถุํ ปิณฺฑาย ปาวิสิ.\csromanspacer{} กปิลวตฺถุสฺมิํ ปิณฺฑาย จริตฺวา ปจฺฉาภตฺตํ ปิณฺฑปาตปฏิกฺกนฺโต เยน กาฬเขมกสฺส สกฺกสฺส วิหาโร เตนุปสงฺกมิ ทิวาวิหาราย.\csromanspacer{} เตน โข ปน สมเยน \csromanfolio{152}กาฬเขมกสฺส สกฺกสฺส วิหาเร สมฺพหุลานิ เสนาสนานิ ปญฺญตฺตานิ โหนฺติ.\csromanspacer{} อทฺทสา โข ภควา กาฬเขมกสฺส สกฺกสฺส วิหาเร สมฺพหุลานิ เสนาสนานิ ปญฺญตฺตานิ, ทิสฺวาน ภควโต เอตทโหสิ “สมฺพหุลานิ โข กาฬเขมกสฺส สกฺกสฺส วิหาเร เสนาสนานิ ปญฺญตฺตานิ, สมฺพหุลา นุ โข อิธ ภิกฺขู วิหรนฺตี”ติ.}

\proseitem{186}{เตน โข ปน สมเยน อายสฺมา อานนฺโท สมฺพหุเลหิ ภิกฺขูหิ สทฺธิํ ฆฏาย สกฺกสฺส วิหาเร จีวรกมฺมํ กโรติ.\csromanspacer{} อถ โข ภควา สายนฺหสมยํ ปฏิสลฺลานา วุฏฺฐิโต เยน ฆฏาย สกฺกสฺส วิหาโร เตนุปสงฺกมิ, อุปสงฺกมิตฺวา ปญฺญตฺเต อาสเน นิสีทิ, นิสชฺช โข ภควา อายสฺมนฺตํ อานนฺทํ อามนฺเตสิ “สมฺพหุลานิ โข อานนฺท กาฬเขมกสฺส สกฺกสฺส วิหาเร เสนาสนานิ ปญฺญตฺตานิ, สมฺพหุลา นุ โข เอตฺถ ภิกฺขู วิหรนฺตี”ติ.\csromanspacer{} สมฺพหุลานิ ภนฺเต กาฬเขมกสฺส สกฺกสฺส วิหาเร เสนาสนานิ ปญฺญตฺตานิ, สมฺพหุลา ภิกฺขู เอตฺถ วิหรนฺติ, จีวรการสมโย โน ภนฺเต วตฺตตีติ.}

\prose{น โข อานนฺท ภิกฺขุ โสภติ สงฺคณิการาโม สงฺคณิกรโต สงฺคณิการามตํ อนุยุตฺโต คณาราโม คณรโต คณสมฺมุทิโต.\csromanspacer{} โส วตานนฺท ภิกฺขุ สงฺคณิการาโม สงฺคณิกรโต สงฺคณิการามตํ อนุยุตฺโต คณาราโม คณรโต คณสมฺมุทิโต, ยํ ตํ เนกฺขมฺมสุขํ ปวิเวกสุขํ อุปสมสุขํ สมฺโพธิสุขํ\footnote{สมฺโพธสุขํ (สี, อิ), สมฺโพธสุขํ จิตฺเตกคฺคตาสุขํ (ก) อุปริ อรณวิภงฺคสุตฺเต ปน สมฺโพธิสุขนฺเตฺวว ทิสฺสติ.}, ตสฺส สุขสฺส นิกามลาภี ภวิสฺสติ อกิจฺฉลาภี อกสิรลาภีติ เนตํ ฐานํ วิชฺชติ.\csromanspacer{} โย จ โข โส อานนฺท ภิกฺขุ เอโก คณสฺมา วูปกฏฺโฐ วิหรติ, ตสฺเสตํ ภิกฺขุโน ปาฏิกงฺขํ.\csromanspacer{} ยํ ตํ เนกฺขมฺมสุขํ ปวิเวกสุขํ อุปสมสุขํ สมฺโพธิสุขํ, ตสฺส สุขสฺส นิกามลาภี ภวิสฺสติ อกิจฺฉลาภี อกสิรลาภีติ ฐานเมตํ วิชฺชติ.}

\prose{โส วตานนฺท ภิกฺขุ สงฺคณิการาโม สงฺคณิกรโต สงฺคณิการามตํ อนุยุตฺโต คณาราโม คณรโต คณสมฺมุทิโต สามายิกํ วา กนฺตํ เจโตวิมุตฺติํ อุปสมฺปชฺช วิหริสฺสติ, อสามายิกํ วา อกุปฺปนฺติ เนตํ ฐานํ วิชฺชติ.\csromanspacer{} โย จ โข โส อานนฺท ภิกฺขุ เอโก คณสฺมา}

\csromanheader{2}{2. มหาสุญฺญตสุตฺต (122)}
\prose{\csromanfolio{153}วูปกฏฺโฐ วิหรติ, ตสฺเสตํ ภิกฺขุโน ปาฏิกงฺขํ สามายิกํ วา กนฺตํ เจโตวิมุตฺติํ อุปสมฺปชฺช วิหริสฺสติ, อสามายิกํ วา อกุปฺปนฺติ ฐานเมตํ วิชฺชติ.}

\prose{นาหํ อานนฺท เอกํ รูปมฺปิ\footnote{เอกรูปมฺปิ (สี)} สมนุปสฺสามิ, ยตฺถ รตฺตสฺส ยถาภิรตสฺส รูปสฺส วิปริณามญฺญถาภาวา น อุปฺปชฺเชยฺยุํ โสกปริเทวทุกฺขโทมนสฺสูปายาสา.}

\proseitem{187}{อยํ โข ปนานนฺท วิหาโร ตถาคเตน อภิสมฺพุทฺโธ, ยทิทํ สพฺพนิมิตฺตานํ อมนสิการา อชฺฌตฺตํ สุญฺญตํ อุปสมฺปชฺช วิหริตุํ.\footnote{วิหรตํ (ก-สี), วิหรติ (สฺยา, กํ, ก)} ตตฺร เจ อานนฺท ตถาคตํ อิมินา วิหาเรน วิหรนฺตํ ภวนฺติ\footnote{ภควนฺตํ (สี, สฺยา, กํ, ก)} อุปสงฺกมิตาโร ภิกฺขู ภิกฺขุนิโย อุปาสกา อุปาสิกาโย ราชาโน ราชมหามตฺตา ติตฺถิยา ติตฺถิยสาวกา.\csromanspacer{} ตตฺรานนฺท ตถาคโต วิเวกนินฺเนเนว จิตฺเตน วิเวกโปเณน วิเวกปพฺภาเรน วูปกฏฺเฐน เนกฺขมฺมาภิรเตน พฺยนฺตีภูเตน สพฺพโส อาสวฏฺฐานีเยหิ ธมฺเมหิ อญฺญทตฺถุ อุยฺโยชนิกปฏิสํยุตฺตํเยว กถํ กตฺตา โหติ.\csromanspacer{} ตสฺมาติหานนฺท ภิกฺขุ เจปิ อากงฺเขยฺย “อชฺฌตฺตํ สุญฺญตํ อุปสมฺปชฺช วิหเรยฺยนฺ”ติ, เตนานนฺท ภิกฺขุนา อชฺฌตฺตเมว จิตฺตํ สณฺฐเปตพฺพํ สนฺนิสาเทตพฺพํ เอโกทิ กาตพฺพํ สมาทหาตพฺพํ.}

\proseitem{188}{กถญฺจานนฺท ภิกฺขุ อชฺฌตฺตเมว จิตฺตํ สณฺฐเปติ สนฺนิสาเทติ เอโกทิํ กโรติ\footnote{เอโกทิกโรติ (สี, สฺยา, กํ, อิ)} สมาทหติ.\csromanspacer{} อิธานนฺท ภิกฺขุ วิวิจฺเจว กาเมหิ วิวิจฺจ อกุสเลหิ ธมฺเมหิ ฯเปฯ ปฐนํ ฌานํ อุปสมฺปชฺช วิหรติ ฯเปฯ ทุติยํ ฌานํ.\csromanspacer{} ตติยํ ฌานํ.\csromanspacer{} จตุตฺถํ ฌานํ อุปสมฺปชฺช วิหรติ.\csromanspacer{} เอวํ โข อานนฺท ภิกฺขุ อชฺฌตฺตเมว จิตฺตํ สณฺฐเปติ สนฺนิสาเทติ เอโกทิํ กโรติ สมาทหติ.\csromanspacer{} โส อชฺฌตฺตํ สุญฺญตํ มนสิ กโรติ, ตสฺส อชฺฌตฺตํ สุญฺญตํ มนสิกโรโต สุญฺญตาย จิตฺตํ น ปกฺขนฺทติ นปฺปสีทติ น สนฺติฏฺฐติ น วิมุจฺจติ.\csromanspacer{} เอวํ สนฺตเมตํ อานนฺท ภิกฺขุ เอวํ ปชานาติ “อชฺฌตฺตํ สุญฺญตํ โข เม มนสิกโรโต อชฺฌตฺตํ สุญฺญตาย จิตฺตํ น ปกฺขนฺทติ นปฺปสีทติ น สนฺติฏฺฐติ น วิมุจฺจตี”ติ, อิติห ตตฺถ สมฺปชาโน โหติ.\csromanspacer{} โส พหิทฺธา สุญฺญตํ มนสิ กโรติ ฯเปฯ.\csromanspacer{} โส อชฺฌตฺตพหิทฺธา สุญฺญตํ มนสิ \csromanfolio{154}กโรติ ฯเปฯ.\csromanspacer{} โส อาเนญฺชํ มนสิ กโรติ, ตสฺส อาเนญฺชํ มนสิกโรโต อาเนญฺชาย จิตฺตํ น ปกฺขนฺทติ นปฺปสีทติ น สนฺติฏฺฐติ น วิมุจฺจติ.\csromanspacer{} เอวํ สนฺตเมตํ อานนฺท ภิกฺขุ เอวํ ปชานาติ “อาเนญฺชํ โข เม มนสิกโรโต อาเนญฺชาย จิตฺตํ น ปกฺขนฺทติ นปฺปสีทติ น สนฺติฏฺฐติ น วิมุจฺจตี”ติ, อิติห ตตฺถ สมฺปชาโน โหติ.}

\prose{เตนานนฺท ภิกฺขุนา ตสฺมิํ เยว ปุริมสฺมิํ สมาธินิมิตฺเต อชฺฌตฺตเมว จิตฺตํ สณฺฐเปตพฺพํ สนฺนิสาเทตพฺพํ เอโกทิ กาตพฺพํ สมาทหาตพฺพํ.\csromanspacer{} โส อชฺฌตฺตํ สุญฺญตํ มนสิ กโรติ, ตสฺส อชฺฌตฺตํ สุญฺญตํ มนสิกโรโต อชฺฌตฺตํ สุญฺญตาย จิตฺตํ ปกฺขนฺทติ ปสีทติ สนฺติฏฺฐติ วิมุจฺจติ.\csromanspacer{} เอวํ สนฺตเมตํ อานนฺท ภิกฺขุ เอวํ ปชานาติ “อชฺฌตฺตํ สุญฺญตํ โข เม มนสิกโรโต อชฺฌตฺตํ สุญฺญตาย จิตฺตํ ปกฺขนฺทติ ปสีทติ สนฺติฏฺฐติ วิมุจฺจตี”ติ, อิติห ตตฺถ สมฺปชาโน โหติ.\csromanspacer{} โส พหิทฺธา สุญฺญตํ มนสิ กโรติ ฯเปฯ.\csromanspacer{} โส อชฺฌตฺตพหิทฺธา สุญฺญตํ มนสิ กโรติ ฯเปฯ.\csromanspacer{} โส อาเนญฺชํ มนสิ กโรติ, ตสฺส อาเนญฺชํ มนสิกโรโต อาเนญฺชาย จิตฺตํ ปกฺขนฺทติ ปสีทติ สนฺติฏฺฐติ วิมุจฺจติ.\csromanspacer{} เอวํ สนฺตเมตํ อานนฺท ภิกฺขุ เอวํ ปชานาติ “อาเนญฺชํ โข เม มนสิกโรโต อาเนญฺชาย จิตฺตํ ปกฺขนฺทติ ปสีทติ สนฺติฏฺฐติ วิมุจฺจตี”ติ, อิติห ตตฺถ สมฺปชาโน โหติ.}

\proseitem{189}{ตสฺส เจ อานนฺท ภิกฺขุโน อิมินา วิหาเรน วิหรโต จงฺกมาย จิตฺตํ นมติ, โส จงฺกมติ “เอวํ มํ จงฺกมนฺตํ นาภิชฺฌาโทมนสฺสา ปาปกา อกุสลา ธมฺมา อนฺวาสฺสวิสฺสนฺตี”ติ, อิติห ตตฺถ สมฺปชาโน โหติ.\csromanspacer{} ตสฺส เจ อานนฺท ภิกฺขุโน อิมินา วิหาเรน วิหรโต ฐานาย จิตฺตํ นมติ, โส ติฏฺฐติ “เอวํ มํ ฐิตํ นาภิชฺฌาโทมนสฺสา ปาปกา อกุสลา ธมฺมา อนฺวาสฺสวิสฺสนฺตี”ติ, อิติห ตตฺถ สมฺปชาโน โหติ.\csromanspacer{} ตสฺส เจ อานนฺท ภิกฺขุโน อิมินา วิหาเรน วิหรโต นิสชฺชาย จิตฺตํ นมติ, โส นิสีทติ “เอวํ มํ นิสินฺนํ นาภิชฺฌาโทมนสฺสา ปาปกา อกุสลา ธมฺมา อนฺวาสฺสวิสฺสนฺตี”ติ, อิติห ตตฺถ สมฺปชาโน โหติ.\csromanspacer{} ตสฺส เจ อานนฺท ภิกฺขุโน อิมินา วิหาเรน วิหรโต สยนาย จิตฺตํ นมติ, โส สยติ “เอวํ มํ สยนฺตํ นาภิชฺฌาโทมนสฺสา ปาปกา อกุสลา ธมฺมา อนฺวาสฺสวิสฺสนฺตี”ติ, อิติห ตตฺถ สมฺปชาโน โหติ.}

\csromanheader{2}{2. มหาสุญฺญตสุตฺต (122)}
\prose{\csromanfolio{155}ตสฺส เจ อานนฺท ภิกฺขุโน อิมินา วิหาเรน วิหรโต กถาย\footnote{ภสฺสาย (สี), ภาสาย (สฺยา, กํ, อิ)} จิตฺตํ นมติ “โส ยายํ กถา หีนา คมฺมา โปถุชฺชนิกา อนริยา อนตฺถสํหิตา น นิพฺพิทาย น วิราคาย น นิโรธาย น อุปสมาย น อภิญฺญาย น สมฺโพธาย น นิพฺพานาย สํวตฺตติ.\csromanspacer{} เสยฺยถิทํ, ราชกถา โจรกถา มหามตฺตกถา เสนากถา ภยกถา ยุทฺธกถา อนฺนกถา ปานกถา วตฺถกถา สยนกถา มาลากถา คนฺธกถา ญาติกถา ยานกถา คามกถา นิคมกถา นครกถา ชนปทกถา อิตฺถิกถา สุรากถา วิสิขากถา กุมฺภฏฺฐานกถา ปุพฺพเปตกถา นานตฺตกถา โลกกฺขายิกา สมุทฺทกฺขายิกา อิติภวาภวกถา อิติ วา.\csromanspacer{} อิติ เอวรูปิํ กถํ น กเถสฺสามี”ติ, อิติห ตตฺถ สมฺปชาโน โหติ.\csromanspacer{} ยา จ โข อยํ อานนฺท กถา อภิสลฺเลขิกา เจโตวินีวรณสปฺปายา\footnote{เจโตวิจารณสปฺปายา (สี, สฺยา, กํ), เจโตวิวรณสปฺปายา (อิ)} เอกนฺตนิพฺพิทาย วิราคาย นิโรธาย อุปสมาย อภิญฺญาย สมฺโพธาย นิพฺพานาย สํวตฺตติ.\csromanspacer{} เสยฺยถิทํ, อปฺปิจฺฉากถา สนฺตุฏฺฐิกถา ปวิเวกกถา อสํสคฺคกถา วีริยารมฺภกถา สีลกถา สมาธิกถา ปญฺญากถา วิมุตฺติกถา วิมุตฺติญาณทสฺสนกถา, อิติ เอวรูปิํ กถํ กเถสฺสามีติ, อิติห ตตฺถ สมฺปชาโน โหติ.}

\prose{ตสฺส เจ อานนฺท ภิกฺขุโน อิมินา วิหาเรน วิหรโต วิตกฺกาย จิตฺตํ นมติ “โส เย เต วิตกฺกา หีนา คมฺมา โปถุชฺชนิกา อนริยา อนตฺถสํหิตา น นิพฺพิทาย น วิราคาย น นิโรธาย น อุปสมาย น อภิญฺญาย น สมฺโพธาย น นิพฺพานาย สํวตฺตนฺติ.\csromanspacer{} เสยฺยถิทํ, กามวิตกฺโก พฺยาปาทวิตกฺโก วิหิํสาวิตกฺโก.\csromanspacer{} อิติ เอวรูเป วิตกฺเก\footnote{เอวรูเปน วิตกฺเกน (สี, สฺยา, กํ, ก)} น วิตกฺเกสฺสามี”ติ, อิติห ตตฺถ สมฺปชาโน โหติ.\csromanspacer{} เย จ โข อิเม อานนฺท วิตกฺกา อริยา นิยฺยานิกา นิยฺยนฺติ ตกฺกรสฺส สมฺมาทุกฺขกฺขยาย.\csromanspacer{} เสยฺยถิทํ, เนกฺขมฺมวิตกฺโก อพฺยาปาทวิตกฺโก อวิหิํสาวิตกฺโก.\csromanspacer{} อิติ เอวรูเป วิตกฺเก\footnote{เอวรูเปน วิตกฺเกน (ก)} วิตกฺเกสฺสามีติ, อิติห ตตฺถ สมฺปชาโน โหติ.}

\proseitem{190}{ปญฺจ โข อิเม อานนฺท กามคุณา.\csromanspacer{} กตเม ปญฺจ, จกฺขุวิญฺเญยฺยา รูปา อิฏฺฐา กนฺตา มนาปา ปิยรูปา กามูปสํหิตา รชนียา. \csromanfolio{156}โสตวิญฺเญยฺยา สทฺทา.\csromanspacer{} ฆานวิญฺเญยฺยา คนฺธา.\csromanspacer{} ชิวฺหาวิญฺเญยฺยา รสา.\csromanspacer{} กายวิญฺเญยฺยา โผฏฺฐพฺพา อิฏฺฐา กนฺตา มนาปา ปิยรูปา กามูปสํหิตา รชนียา.\csromanspacer{} อิเม โข อานนฺท ปญฺจ กามคุณา.\csromanspacer{} ยตฺถ ภิกฺขุนา อภิกฺขณํ สกํ จิตฺตํ ปจฺจเวกฺขิตพฺพํ “อตฺถิ นุ โข อิเมสุ ปญฺจสุ กามคุเณสุ อญฺญตรสฺมิํ วา อญฺญตรสฺมิํ วา อายตเน อุปฺปชฺชติ เจตโส สมุทาจาโร”ติ.\csromanspacer{} สเจ อานนฺท ภิกฺขุ ปจฺจเวกฺขมาโน เอวํ ปชานาติ “อตฺถิ โข เม อิเมสุ ปญฺจสุ กามคุเณสุ อญฺญตรสฺมิํ วา อญฺญตรสฺมิํ วา อายตเน อุปฺปชฺชติ เจตโส สมุทาจาโร”ติ.\csromanspacer{} เอวํ สนฺตเมตํ\footnote{เอวํ สนฺตํ (ฏฺฐ)} อานนฺท ภิกฺขุ เอวํ ปชานาติ “โย โข อิเมสุ ปญฺจสุ กามคุเณสุ ฉนฺทราโค, โส เม นปฺปหีโน”ติ, อิติห ตตฺถ สมฺปชาโน โหติ.\csromanspacer{} สเจ ปนานนฺท ภิกฺขุ ปจฺจเวกฺขมาโน เอวํ ปชานาติ “นตฺถิ โข เม อิเมสุ ปญฺจสุ กามคุเณสุ อญฺญตรสฺมิํ วา อญฺญตรสฺมิํ วา อายตเน อุปฺปชฺชติ เจตโส สมุทาจาโร”ติ.\csromanspacer{} เอวํ สนฺตเมตํ อานนฺท ภิกฺขุ เอวํ ปชานาติ “โย โข อิเมสุ ปญฺจสุ กามคุเณสุ ฉนฺทราโค, โส เม ปหีโน”ติ, อิติห ตตฺถ สมฺปชาโน โหติ.}

\proseitem{191}{ปญฺจ โข อิเม อานนฺท อุปาทานกฺขนฺธา, ยตฺถ ภิกฺขุนา อุทยพฺพยานุปสฺสินา วิหาตพฺพํ “อิติ รูปํ, อิติ รูปสฺส สมุทโย, อิติ รูปสฺส อตฺถงฺคโม.\csromanspacer{} อิติ เวทนา.\csromanspacer{} อิติ สญฺญา.\csromanspacer{} อิติ สงฺขารา.\csromanspacer{} อิติ วิญฺญาณํ, อิติ วิญฺญาณสฺส สมุทโย, อิติ วิญฺญาณสฺส อตฺถงฺคโม”ติ.\csromanspacer{} ตสฺส อิเมสุ ปญฺจสุ อุปาทานกฺขนฺเธสุ อุทยพฺพยานุปสฺสิโน วิหรโต โย ปญฺจสุ อุปาทานกฺขนฺเธสุ อสฺมิมาโน, โส ปหียติ.\csromanspacer{} เอวํ สนฺตเมตํ อานนฺท ภิกฺขุ เอวํ ปชานาติ “โย โข อิเมสุ ปญฺจสุ อุปาทานกฺขนฺเธสุ อสฺมิมาโน, โส เม ปหีโน”ติ, อิติห ตตฺถ สมฺปชาโน โหติ.\csromanspacer{} อิเม โข เต อานนฺท ธมฺมา เอกนฺตกุสลา กุสลายาติกา\footnote{ธมฺมา เอกนฺตกุสลายติกา (สพฺพตฺถ) อฏฺฐกถาฏีกา โอโลเกตพฺพา.} อริยา โลกุตฺตรา อนวกฺกนฺตา ปาปิมตา.\csromanspacer{} ตํ กิํ มญฺญสิ อานนฺท, กํ อตฺถวสํ สมฺปสฺสมาโน อรหติ สาวโก สตฺถารํ อนุพนฺธิตุํ อปิ ปณุชฺชมาโนติ\footnote{อปิ ปนุชฺชมาโนปีติ (ก-สี), อปิ ปยุชฺชมาโนติ (สฺยา, กํ, อิ)}.\csromanspacer{} ภควํมูลกา โน ภนฺเต ธมฺมา ภควํเนตฺติกา}

\csromanheader{2}{2. มหาสุญฺญตสุตฺต (122)}
\prose{\csromanfolio{157}ภควํปฏิสรณา, สาธุ วต ภนฺเต ภควนฺตํเยว ปฏิภาตุ, เอตสฺส ภาสิตสฺส อตฺโถ, ภควโต สุตฺวา ภิกฺขู ธาเรสฺสนฺตีติ.}

\proseitem{192}{น โข อานนฺท อรหติ สาวโก สตฺถารํ อนุพนฺธิตุํ, ยทิทํ สุตฺตํ เคยฺยํ เวยฺยากรณํ ตสฺส เหตุ\footnote{เวยฺยากรณสฺส เหตุ (ก)}.\csromanspacer{} ตํ กิสฺส เหตุ, ทีฆรตฺตสฺส\footnote{ทีฆรตฺตํ + อสฺสาติ ปทจฺเฉโท.} หิ เต อานนฺท ธมฺมา สุตา ธาตา วจสา ปริจิตา มนสานุเปกฺขิตา ทิฏฺฐิยา สุปฺปฏิวิทฺธา.\csromanspacer{} ยา จ โข อยํ อานนฺท กถา อภิสลฺเลขิกา เจโตวินีวรณสปฺปายา เอกนฺตนิพฺพิทาย วิราคาย นิโรธาย อุปสมาย อภิญฺญาย สมฺโพธาย นิพฺพานาย สํวตฺตติ.\csromanspacer{} เสยฺยถิทํ, อปฺปิจฺฉกถา สนฺตุฏฺฐิกถา ปวิเวกกถา อสํสคฺคกถา วีริยารมฺภกถา สีลกถา สมาธิกถา ปญฺญากถา วิมุตฺติกถา วิมุตฺติญาณทสฺสนกถา.\csromanspacer{} เอวรูปิยา โข อานนฺท กถาย เหตุ อรหติ สาวโก สตฺถารํ อนุพนฺธิตุํ อปิ ปณุชฺชมาโน.}

\prose{เอวํ สนฺเต โข อานนฺท อาจริยูปทฺทโว โหติ.\csromanspacer{} เอวํ สนฺเต อนฺเตวาสูปทฺทโว โหติ, เอวํ สนฺเต พฺรหฺมจารูปทฺทโว โหติ.}

\proseitem{193}{กถญฺจานนฺท อาจริยูปทฺทโว โหติ, อิธานนฺท เอกจฺโจ สตฺถา วิวิตฺตํ เสนาสนํ ภชติ อรญฺญํ รุกฺขมูลํ ปพฺพตํ กนฺทรํ คิริคุหํ สุสานํ วนปตฺถํ อพฺโภกาสํ ปลาลปุญฺชํ.\csromanspacer{} ตสฺส ตถาวูปกฏฺฐสฺส วิหรโต อนฺวาวตฺตนฺติ\footnote{อนฺวาวฏฺฏนฺติ (สี, สฺยา, กํ, อิ)} พฺราหฺมณคหปติกา เนคมา เจว ชานปทา จ.\csromanspacer{} โส อนฺวาวตฺตนฺเตสุ พฺราหฺมณคหติเกสุ เนคเมสุ เจว ชานปเทสุ จ มุจฺฉํ นิกามยติ\footnote{มุจฺฉติ กามยติ (สี, อิ) อฏฺฐกถายํ ปน น ตถา ทิสฺสติ.}, เคธํ อาปชฺชติ, อาวตฺตติ พาหุลฺลาย.\csromanspacer{} อยํ วุจฺจตานนฺท อุปทฺทโว\footnote{อุปทฺทุโต (สี, อิ)} อาจริโย อาจริยูปทฺทเวน อวธิํสุ นํ ปาปกา อกุสลา ธมฺมา สํกิเลสิกา โปโนพฺภวิกา\footnote{โปโนภวิกา (สี, อิ)} สทรา ทุกฺขวิปากา อายติํ ชาติชรามรณิยา.\csromanspacer{} เอวํ โข อานนฺท อาจริยูปทฺทโว โหติ.}

\proseitem{194}{กถญฺจานนฺท อนฺเตวาสูปทฺทโว โหติ, ตสฺเสว โข ปนานนฺท สตฺถุ สาวโก ตสฺส สตฺถุ วิเวกมนุพฺรูหยมาโน วิวิตฺตํ เสนาสนํ \csromanfolio{158}ภชติ อรญฺญํ รุกฺขมูลํ ปพฺพตํ กนฺทรํ คิริคุหํ สุสานํ วนปตฺถํ อพฺโภกาสํ ปลาลปุญฺชํ.\csromanspacer{} ตสฺส ตถาวูปกฏฺฐสฺส วิหรโต อนฺวาวตฺตนฺติ พฺราหฺมณคหปติกา เนคมา เจว ชานปทา จ.\csromanspacer{} โส อนฺวาวตฺตนฺเตสุ พฺราหฺมณคหติเกสุ เนคเมสุ เจว ชานปเทสุ จ มุจฺฉํ นิกามยติ, เคธํ อาปชฺชติ, อาวตฺตติ พาหุลฺลาย.\csromanspacer{} อยํ วุจฺจตานนฺท อุปทฺทโว อนฺเตวาสี อนฺเตวาสูปทฺทเวน อวธิํสุ นํ ปาปกา อกุสลา ธมฺมา สํกิเลสิกา โปโนพฺภวิกา สทรา ทุกฺขวิปากา อายติํ ชาติชรามรณิยา.\csromanspacer{} เอวํ โข อานนฺท อนฺเตวาสูปทฺทโว โหติ.}

\proseitem{195}{กถญฺจานนฺท พฺรหฺมจารูปทฺทโว โหติ, อิธานนฺท ตถาคโต โลเก อุปฺปชฺชติ อรหํ สมฺมาสมฺพุทฺโธ วิชฺชาจรณสมฺปนฺโน สุคโต โลกวิทู อนุตฺตโร ปุริสทมฺมสารถิ สตฺถา เทวมนุสฺสานํ พุทฺโธ ภควา.\csromanspacer{} โส วิวิตฺตํ เสนาสนํ ภชติ อรญฺญํ รุกฺขมูลํ ปพฺพตํ กนฺทรํ คิริคุหํ สุสานํ วนปตฺถํ อพฺโภกาสํ ปลาลปุญฺชํ.\csromanspacer{} ตสฺส ตถาวูปกฏฺฐสฺส วิหรโต อนฺวาวตฺตนฺติ พฺราหฺมณคหปติกา เนคมา เจว ชานาปทา จ.\csromanspacer{} โส อนฺวาวตฺตนฺเตสุ พฺราหฺมณคหปติเกสุ เนคเมสุ เจว ชานปเทสุ จ น มุจฺฉํ นิกามยติ, น เคธํ อาปชฺชติ, น อาวตฺตติ พาหุลฺลาย.\csromanspacer{} ตสฺเสว โข ปนานนฺท สตฺถุ สาวโก ตสฺส สตฺถุ วิเวกมนุพฺรูหยมาโน วิวิตฺตํ เสนาสนํ ภชติ อรญฺญํ รุกฺขมูลํ ปพฺพตํ กนฺทรํ คิริคุหํ สุสานํ วนปตฺถํ อพฺโภกาสํ ปลาลปุญฺชํ.\csromanspacer{} ตสฺส ตถาวูปกฏฺฐสฺส วิหรโต อนฺวาวตฺตนฺติ พฺราหฺมณคหปติกา เนคมา เจว ชานปทา จ.\csromanspacer{} โส อนฺวาวตฺตนฺเตสุ พฺราหฺมณคหปติเกสุ เนคเมสุ เจว ชานปเทสุ จ มุจฺฉํ นิกามยติ, เคธํ อาปชฺชติ, อาวตฺตติ พาหุลฺลาย.\csromanspacer{} อยํ วุจฺจตานนฺท อุปทฺทโว พฺรหฺมจารี พฺรหฺมจารูปทฺทเวน อวธิํสุ นํ ปาปกา อกุสลา ธมฺมา สํกิเลสิกา โปโนพฺภวิกา สทรา ทุกฺขวิปากา อายติํ ชาติชรามรณิยา.\csromanspacer{} เอวํ โข อานนฺท พฺรหฺมจารูปทฺทโว โหติ.}

\prose{ตตฺรานนฺท โย เจวายํ อาจริยูปทฺทโว โย จ อนฺเตวาสูปทฺทโว, อยํ เตหิ พฺรหฺมจารูปทฺทโว ทุกฺขวิปากตโร เจว กฏุกวิปากตโร จ.\csromanspacer{} อปิ จ วินิปาตาย สํวตฺตติ.}

\proseitem{196}{ตสฺมาติห มํ อานนฺท มิตฺตวตาย สมุทาจรถ มา สปตฺตวตาย, ตํ โว ภวิสฺสติ ทีฆรตฺตํ หิตาย สุขาย.}

\csromanheader{2}{3. อจฺฉริย-อพฺภุตสุตฺต (123)}
\prose{\csromanfolio{159}กถญฺจานนฺท สตฺถารํ สาวกา สปตฺตวตาย สมุทาจรนฺติ โน มิตฺตวตาย.\csromanspacer{} อิธานนฺท สตฺถา สาวกานํ ธมฺมํ เทเสติ อนุกมฺปโก หิเตสี อนุกมฺปํ อุปาทาย “อิทํ โว หิตาย อิทํ โว สุขายา”ติ.\csromanspacer{} ตสฺส สาวกา น สุสฺสูสนฺติ, น โสตํ โอทหนฺติ, น อญฺญา จิตฺตํ อุปฏฺฐเปนฺติ, โวกฺกมฺม จ สตฺถุสาสนา วตฺตนฺติ.\csromanspacer{} เอวํ โข อานนฺท สตฺถารํ สาวกา สปตฺตวตาย สมุทาจรนฺติ โน มิตฺตวตาย.}

\prose{กถญฺจานนฺท สตฺถารํ สาวกา มิตฺตวตาย สมุทาจรนฺติ โน สปตฺตวตาย.\csromanspacer{} อิธานนฺท สตฺถา สาวกานํ ธมฺมํ เทเสติ อนุกมฺปโก หิเตสี อนุกมฺปํ อุปาทาย”อิทํ โว หิตาย อิทํ โว สุขายา”ติ.\csromanspacer{} ตสฺส สาวกา สุสฺสูสนฺติ, โสตํ โอทหนฺติ, อญฺญา จิตฺตํ อุปฏฺฐเปนฺติ, น จ โวกฺกม สตฺถุสาสนา วตฺตนฺติ.\csromanspacer{} เอวํ โข อานนฺท สตฺถารํ สาวกา มิตฺตวตาย สมุทาจรนฺติ โน สปตฺตวตาย.}

\prose{ตสฺมาติห มํ อานนฺท มิตฺตวตาย สมุทาจรถ มา สปตฺตวตาย, ตํ โว ภวิสฺสติ ทีฆรตฺตํ หิตาย สุขาย.\csromanspacer{} น โว อหํ อานนฺท ตถา ปรกฺกมิสฺสามิ, ยถา กุมฺภกาโร อามเก อามกมตฺเต.\csromanspacer{} นิคฺคยฺห นิคฺคยฺหาหํ อานนฺท วกฺขามิ, ปวยฺห ปวยฺห อานนฺท วกฺขามิ\footnote{ปวยฺห ปวยฺห (สี, อิ), ปคฺคยฺห ปคฺคยฺห อานนฺท วกฺขามิ (ก)}, โย สาโร โส ฐสฺสตีติ.}

\prose{อิทมโวจ ภควา.\csromanspacer{} อตฺตมโน อายสฺมา อานนฺโท ภควโต ภาสิตํ อภินนฺทีติ.}

\nitthitamruled{มหาสุญฺญตสุตฺตํ นิฏฺฐิตํ ทุติยํ.}
\csromanmatikaanchor{mtk.28}
\csromantocmarkref{subsection}{3. อจฺฉริย-อพฺภุตสุตฺต}{mtk.28}
\bookmark[dest={mtk.28},level=2]{3. อจฺฉริย-อพฺภุตสุตฺต}
\csromanheader{2}{3. อจฺฉริย-อพฺภุตสุตฺต}
\proseitem{197}{เอวํ เม สุตํ\csromandash{}เอกํ สมยํ ภควา สาวตฺถิยํ วิหรติ เชตวเน อนาถปิณฺฑิกสฺส อาราเม.\csromanspacer{} อถ โข สมฺพหุลานํ ภิกฺขูนํ ปจฺฉาภตฺตํ ปิณฺฑปาตปฏิกฺกนฺตานํ อุปฏฺฐานสาลายํ สนฺนิสินฺนานํ สนฺนิปติตานํ อยมนฺตรากถา อุทปาทิ “อจฺฉริยํ อาวุโส, อพฺภุตํ อาวุโส ตถาคตสฺส มหิทฺธิกตา มหานุภาวตา.\csromanspacer{} ยตฺร หิ นาม \csromanfolio{160}ตถาคโต อตีเต พุทฺเธ ปรินิพฺพุเต ฉินฺนปปญฺเจ ฉินฺนวฏุเม ปริยาทินฺนวฏฺเฏ สพฺพทุกฺขวีติวตฺเต ชานิสฺสติ\footnote{อนุสฺสริสฺสติ ชานิสฺสติ (ก)}. ‘เอวํชจฺจา เต ภควนฺโต อเหสุํ อิติปิ, เอวํนามา เต ภควนฺโต อเหสุํ อิติปิ, เอวํโคตฺตา เต ภควนฺโต อเหสุํ อิติปิ, เอวํสีลา เต ภควนฺโต อเหสุํ อิติปิ, เอวํธมฺมา เต ภควนฺโต อเหสุํ อิติปิ, เอวํปญฺญา เต ภควนฺโต อเหสุํ อิติปิ, เอวํวิหารี เต ภควนฺโต อเหสุํ อิติปิ, เอวํวิมุตฺตา เต ภควนฺโต อเหสุํ อิติปี’ติ”.\csromanspacer{} เอวํ วุตฺเต อายสฺมา อานนฺโท เต ภิกฺขู เอตทโวจ “อจฺฉริยา เจว อาวุโส ตถาคตา อจฺฉริยธมฺมสมนฺนาคตา จ, อพฺภุตา เจว อาวุโส ตถาคตา อพฺภุตธมฺมสมนฺนาคตา จา”ติ.\csromanspacer{} อยญฺจ หิทํ เตสํ ภิกฺขูนํ อนฺตรากถา วิปฺปกตา โหติ.}

\proseitem{198}{อถ โข ภควา สายนฺหสมยํ ปฏิสลฺลานา วุฏฺฐิโต เยนุปฏฺฐานสาลา เตนุปสงฺกมิ, อุปสงฺกมิตฺวา ปญฺญตฺเต อาสเน นิสีทิ, นิสชฺช โข ภควา ภิกฺขู อามนฺเตสิ “กาย นุตฺถ ภิกฺขเว เอตรหิ กาถาย สนฺนิสินฺนา, กา จ ปน โว อนฺตรากถา วิปฺปกตา”ติ.\csromanspacer{} อิธ ภนฺเต อมฺหากํ ปจฺฉา ภตฺตํ ปิณฺฑปาตปฏิกฺกนฺตานํ อุปฏฺฐานสาลายํ สนฺนิสินฺนานํ สนฺนิปติตานํ อยมนฺตรากถา อุทปาทิ “อจฺฉริยํ อาวุโส, อพฺภุตํ อาวุโส ตถาคตสฺส มหิทฺธิกตา มหานุภาวตา, ยตฺร หิ นาม ตถาคโต อตีเต พุทฺเธ ปรินิพฺพุเต ฉินฺนปปญฺเจ ฉินฺนวฏุเม ปริยาทินฺนวฏฺเฏ สพฺพทุกฺขวีติวตฺเต ชานิสฺสติ ‘เอวํชจฺจา เต ภควนฺโต อเหสุํ อิติปิ, เอวํนามา.\csromanspacer{} เอวํโคตฺตา.\csromanspacer{} เอวํสีลา.\csromanspacer{} เอวํธมฺมา.\csromanspacer{} เอวํปญฺญา.\csromanspacer{} เอวํวิหารี.\csromanspacer{} เอวํวิมุตฺตา เต ภควนฺโต อเหสุํ อิติปี’ติ”.\csromanspacer{} เอวํ วุตฺเต ภนฺเต อายสฺมา อานนฺโท อเมฺห เอตทโวจ “อจฺฉริยา เจว อาวุโส ตถาคตา อจฺฉริยธมฺมสมนฺนาคตา จ, อพฺภุตา เจว อาวุโส ตถาคตา อพฺภุตธมฺมสมนฺนาคตา จา”ติ.\csromanspacer{} อยํ โข โน ภนฺเต อนฺตรากถา วิปฺปกตา, อถ ภควา อนุปฺปตฺโตติ.}

\proseitem{199}{อถ โข ภควา อายสฺมนฺตํ อานนฺทํ อามนฺเตสิ “ตสฺมาติห ตํ อานนฺท ภิยฺโยโส มตฺตาย ปฏิภนฺตุ ตถาคตสฺส อจฺฉริยา อพฺภุตธมฺมา”ติ\footnote{อพฺภุตา ธมฺมาติ (?)}.}

\csromanheader{2}{3. อจฺฉริย-อพฺภุตสุตฺต (123)}
\prose{\csromanfolio{161}สมฺมุขา เมตํ ภนฺเต ภควโต สุตํ, สมฺมุขา ปฏิคฺคหิตํ “สโต สมฺปชาโน อานนฺท โพธิสตฺโต ตุสิตํ กายํ อุปปชฺชี”ติ.\csromanspacer{} ยมฺปิ ภนฺเต สโต สมฺปชาโน โพธิสตฺโต ตุสิตํ กายํ อุปปชฺชิ, อิทํปาหํ ภนฺเต ภควโต อจฺฉริยํ อพฺภุตธมฺมํ ธาเรมิ. (1)}

\prose{สมฺมุขา เมตํ ภนฺเต ภควโต สุตํ, สมฺมุขา ปฏิคฺคหิตํ “สโต สมฺปชาโน อานนฺท โพธิสตฺโต ตุสิเต กาเย อฏฺฐาสี”ติ.\csromanspacer{} ยมฺปิ ภนฺเต สโต สมฺปชาโน โพธิสตฺโต ตุสิเต กาเย อฏฺฐาสิ, อิทํปาหํ\footnote{อิทํปหํ (สี, สฺยา, กํ, อิ)} ภนฺเต ภควโต อจฺฉริยํ อพฺภุตธมฺมํ ธาเรมิ. (2)}

\proseitem{200}{สมฺมุขา เมตํ ภนฺเต ภควโต สุตํ, สมฺมุขา ปฏิคฺคหิตํ “ยาวตายุกํ อานนฺท โพธิสตฺโต ตุสิเต กาเย อฏฺฐาสี”ติ.\csromanspacer{} ยมฺปิ ภนฺเต ยาวตายุกํ โพธิสตฺโต ตุสิเต กาเย อฏฺฐาสิ, อิทํปาหํ ภนฺเต ภควโต อจฺฉริยํ อพฺภุตธมฺมํ ธาเรมิ. (3)}

\prose{สมฺมุขา เมตํ ภนฺเต ภควโต สุตํ, สมฺมุขา ปฏิคฺคหิตํ “สโต สมฺปชาโน อานนฺท โพธิสตฺโต ตุสิตา กายา จวิตฺวา มาตุกุจฺฉิํ โอกฺกมี”ติ.\csromanspacer{} ยมฺปิ ภนฺเต สโต สมฺปชาโน โพธิสตฺโต ตุสิตา กายา จวิตฺวา มาตุกุจฺฉิํ โอกฺกมิ, อิทํปาหํ ภนฺเต ภควโต อจฺฉริยํ อพฺภุตธมฺมํ ธาเรมิ. (4)}

\proseitem{201}{สมฺมุขา เมตํ ภนฺเต ภควโต สุตํ, สมฺมุขา ปฏิคฺคหิตํ “ยทา อานนฺท โพธิสตฺโต ตุสิตา กายา จวิตฺวา มาตุกุจฺฉิํ โอกฺกมติ, อถ สเทวเก โลเก สมารเก สพฺรหฺมเก สสฺสมณพฺราหฺมณิยา ปชาย สเทวมนุสฺสาย อปฺปมาโณ อุฬาโร โอภาโส โลเก ปาตุภวติ อติกฺกมฺเมว เทวานํ เทวานุภาวํ.\csromanspacer{} ยาปิ ตา โลกนฺตริกา อฆา อสํวุตา อนฺธการา อนฺธการติมิสา, ยตฺถปิเม จนฺทิมสูริยา เอวํมหิทฺธิกา เอวํมหานุภาวา อาภาย นานุโภนฺติ, ตตฺถปิ อปฺปมาโณ อุฬาโร โอภาโส โลเก ปาตุภวติ อติกฺกมฺเมว เทวานํ เทวานุภาวํ.\csromanspacer{} เยปิ ตตฺถ สตฺตา อุปปนฺนา, เตปิ เตโนภาเสน อญฺญมญฺญํ สญฺชานนฺติ ‘อญฺเญปิ กิร โภ สนฺติ สตฺตา อิธูปปนฺนา’ติ.\csromanspacer{} อยญฺจ ทสสหสฺสี โลกธาตุ สงฺกมฺปติ สมฺปกมฺปติ \csromanfolio{162}สมฺปเวธติ, อปฺปมาโณ จ อุฬาโร โอภาโส โลเก ปาตุภวติ อติกฺกมฺเมว เทวานํ เทวานุภาวนฺ”ติ.\csromanspacer{} ยมฺปิ ภนฺเต ฯเปฯ อิทํปาหํ ภนฺเต ภควโต อจฺฉริยํ อพฺภุตธมฺมํ ธาเรมิ. (5)}

\proseitem{202}{สมฺมุขา เมตํ ภนฺเต ภควโต สุตํ, สมฺมุขา ปฏิคฺคหิตํ “ยทา อานนฺท โพธิสตฺโต มาตุกุจฺฉิํ โอกฺกนฺโต โหติ, จตฺตาโร เทวปุตฺตา จตุทฺทิสํ อารกฺขาย อุปคจฺฉนฺติ ‘มา นํ โพธิสตฺตํ วา โพธิสตฺตมาตรํ วา มนุสฺโส วา อมนุสฺโส วา โกจิ วา วิเหเฐสี’ติ”.\csromanspacer{} ยมฺปิ ภนฺเต ฯเปฯ อิทํปาหํ ภนฺเต ภควโต อจฺฉริยํ อพฺภุตธมฺมํ ธาเรมิ. (6)}

\proseitem{203}{สมฺมุขา เมตํ ภนฺเต ภควโต สุตํ, สมฺมุขาปฏิคฺคหิตํ “ยทา อานนฺท โพธิสตฺโต มาตุกุจฺฉิํ โอกฺกนฺโต โหติ, ปกติยา สีลวตี โพธิสตฺตมาตา โหติ, วิรตา ปาณาติปาตา วิรตา อทินฺนาทานา วิรตา กาเมสุมิจฺฉาจารา วิรตา มุสาวาทา วิรตา สุราเมรยมชฺชปมาทฏฺฐานา”ติ.\csromanspacer{} ยมฺปิ ภนฺเต ฯเปฯ อิทํปาหํ ภนฺเต ภควโต อจฺฉริยํ อพฺภุตธมฺมํ ธาเรมิ. (7)}

\prose{สมฺมุขา เมตํ ภนฺเต ภควโต สุตํ, สมฺมุขา ปฏิคฺคหิตํ “ยทา อานนฺท โพธิสตฺโต มาตุกุจฺฉิํ โอกฺกนฺโต โหติ, น โพธิสตฺตมาตุ ปุริเสสุ มานสํ อุปฺปชฺชติ กามคุณูปสํหิตํ, อนติกฺกมนียา จ โพธิสตฺตมาตา โหติ เกนจิ ปุริเสน รตฺตจิตฺเตนา”ติ.\csromanspacer{} ยมฺปิ ภนฺเต ฯเปฯ อิทํปาหํ ภนฺเต ภควโต อจฺฉริยํ อพฺภุตธมฺมํ ธาเรมิ. (8)}

\prose{สมฺมุขา เมตํ ภนฺเต ภควโต สุตํ, สมฺมุขา ปฏิคฺคหิตํ “ยทา อานนฺท โพธิสตฺโต มาตุกุจฺฉิํ โอกฺกนฺโต โหติ, ลาภินี โพธิสตฺตมาตา โหติ ปญฺจนฺนํ กามคุณานํ, สา ปญฺจหิ กามคุเณหิ สมปฺปิตา สมงฺคีภูตา ปริจาเรตี”ติ.\csromanspacer{} ยมฺปิ ภนฺเต ฯเปฯ อิทํปาหํ ภนฺเต ภควโต อจฺฉริยํ อพฺภุตธมฺมํ ธาเรมิ. (9)}

\proseitem{204}{สมฺมุขา เมตํ ภนฺเต ภควโต สุตํ, สมฺมุขา ปฏิคฺคหิตํ “ยทา อานนฺท โพธิสตฺโต มาตุกุจฺฉิํ โอกฺกนฺโต โหติ, น โพธิสตฺตมาตุ โกจิเทว อาพาโธ อุปฺปชฺชติ, สุขินี โพธิสตฺตมาตา โหติ อกิลนฺตกายา, โพธิสตฺตญฺจ โพธิสตฺตมาตา}

\csromanheader{2}{3. อจฺฉริย-อพฺภุตสุตฺต (123)}
\prose{\csromanfolio{163}ติโรกุจฺฉิคตํ ปสฺสติ สพฺพงฺคปจฺจงฺคํ อหีนินฺทฺริยํ, เสยฺยถาปิ อานนฺท มณิ เวฬุริโย สุโภ ชาติมา อฏฺฐํโส สุปริกมฺมกโต, ตตฺราสฺส สุตฺตํ อาวุตํ นีลํ วา ปีตํ วา โลหิตํ วา โอทาตํ วา ปณฺฑุสุตฺตํ วา, ตเมนํ จกฺขุมา ปุริโส หตฺเถ กริตฺวา ปจฺจเวกฺขยฺย ‘อยํ โข มณิ เวฬุริโย สุโภ ชาติมา อฏฺฐํโส สุปริกมฺมกโต, ตตฺริทํ สุตฺตํ อาวุตํ นีลํ วา ปีตํ วา โลหิตํ วา โอทาตํ วา ปณฺฑุสุตฺตํ วา’ติ.\csromanspacer{} เอวเมว โข อานนฺท ยทา โพธิสตฺโต มาตุกุจฺฉิํ โอกฺกนฺโต โหติ, น โพธิสตฺตมาตุ โกจิเทว อาพาโธ อุปฺปชฺชติ, สุขินี โพธิสตฺตมาตา โหติ อกิลนฺตกายา.\csromanspacer{} โพธิสตฺตญฺจ โพธิสตฺตมาตา ติโรกุจฺฉิคตํ ปสฺสติ สพฺพงฺคปจฺจงฺคํ อหีนินฺทฺริยนฺ”ติ.\csromanspacer{} ยมฺปิ ภนฺเต ฯเปฯ อิทํปาหํ ภนฺเต ภควโต อจฺฉริยํ อพฺภุตธมฺมํ ธาเรมิ. (10)}

\proseitem{205}{สมฺมุขา เมตํ ภนฺเต ภควโต สุตํ, สมฺมุขา ปฏิคฺคหิตํ “สตฺตาหชาเต อานนฺท โพธิสตฺเต โพธิสตฺตมาตา กาลํ กโรติ, ตุสิตํ กายํ อุปปชฺชตี”ติ.\csromanspacer{} ยมฺปิ ภนฺเต ฯเปฯ อิทํปาหํ ภนฺเต ภควโต อจฺฉริยํ อพฺภุตธมฺมํ ธาเรมิ. (11)}

\prose{สมฺมุขา เมตํ ภนฺเต ภควโต สุตํ, สมฺมุขา ปฏิคฺคหิตํ “ยถา โข ปนานนฺท อญฺญา อิตฺถิกา นว วา ทส วา มาเส คพฺภํ กุจฺฉินา ปริหริตฺวา วิชายนฺติ.\csromanspacer{} น เหวํ โพธิสตฺตํ โพธิสตฺตมาตา วิชายติ, ทเสว มาสานิ โพธิสตฺตํ โพธิสตฺตมาตา กุจฺฉินา ปริหริตฺวา วิชายตี”ติ.\csromanspacer{} ยมฺปิ ภนฺเต ฯเปฯ อิทํปาหํ ภนฺเต ภควโต อจฺฉริยํ อพฺภุตธมฺมํ ธาเรมิ. (12)}

\prose{สมฺมุขา เมตํ ภนฺเต ภควโต สุตํ, สมฺมุขา ปฏิคฺคหิตํ “ยถา โข ปนานนฺท อญฺญา อิตฺถิกา นิสินฺนา วา นิปนฺนา วา วิชายนฺติ.\csromanspacer{} น เหวํ โพธิสตฺตํ โพธิสตฺตมาตา วิชายติ, ฐิตาว โพธิสตฺตํ โพธิสตฺตมาตา วิชายตี”ติ.\csromanspacer{} ยมฺปิ ภนฺเต ฯเปฯ อิทํปาหํ ภนฺเต ภควโต อจฺฉริยํ อพฺภุตธมฺมํ ธาเรมิ. (13)}

\prose{สมฺมุขา เมตํ ภนฺเต ภควโต สุตํ, สมฺมุขา ปฏิคฺคหิตํ “ยทา อานนฺท โพธิสตฺโต มาตุกุจฺฉิมฺหา นิกฺขมติ, เทวา นํ ปฐมํ ปฏิคฺคณฺหนฺติ, ปจฺฉา มนุสฺสา”ติ.\csromanspacer{} ยมฺปิ ภนฺเต ฯเปฯ อิทํปาหํ ภนฺเต ภควโต อจฺฉริยํ อพฺภุตธมฺมํ ธาเรมิ. (14)}

\proseitem{206}{\csromanfolio{164}สมฺมุขา เมตํ ภนฺเต ภควโต สุตํ, สมฺมุขา ปฏิคฺคหิตํ “ยทา อานนฺท โพธิสตฺโต มาตุกุจฺฉิมฺหา นิกฺขมติ, อปฺปตฺโตว โพธิสตฺโต ปถวิํ โหติ.\csromanspacer{} จตฺตาโร นํ เทวปุตฺตา ปฏิคฺคเหตฺวา มาตุ ปุรโต ฐเปนฺติ ‘อตฺตมนา เทวิ โหหิ, มเหสกฺโข เต ปุตฺโต อุปฺปนฺโน’ติ”.\csromanspacer{} ยมฺปิ ภนฺเต ฯเปฯ อิทํปาหํ ภนฺเต ภควโต อจฺฉริยํ อพฺภุตธมฺมํ ธาเรมิ. (15)}

\prose{สมฺมุขา เมตํ ภนฺเต ภควโต สุตํ, สมฺมุขา ปฏิคฺคหิตํ “ยทา อานนฺท โพธิสตฺโต มาตุกุจฺฉิมฺหา นิกฺขมติ, วิสโทว นิกฺขมติ อมกฺขิโต อุเทน\footnote{อุทฺเทน (สี, สฺยา, กํ, อิ)} อมกฺขิโต เสเมฺหน อมกฺขิโต รุหิเรน อมกฺขิโต เกนจิ อสุจินา, สุทฺโธ วิสโท\footnote{วิสุทฺโธ (สฺยา)}.\csromanspacer{} เสยฺยถาปิ อานนฺท มณิรตนํ กาสิเก วตฺเถ นิกฺขิตฺตํ เนว มณิรตนํ กาสิกํ วตฺถํ มกฺเขติ, นาปิ กาสิกํ วตฺถํ มณิรตนํ มกฺเขติ.\csromanspacer{} ตํ กิสฺส เหตุ, อุภินฺนํ สุทฺธตฺตา.\csromanspacer{} เอวเมว โข อานนฺท ยทา โพธิสตฺโต มาตุกุจฺฉิมฺหา นิกฺขมติ, วิสโทว นิกฺขมติ อมกฺขิโต อุเทน อมกฺขิโต เสเมฺหน อมกฺขิโต รุหิเรน อมกฺขิโต เกนจิ อสุจินา, สุทฺโธ วิสโท”ติ.\csromanspacer{} ยมฺปิ ภนฺเต ฯเปฯ อิทํปาหํ ภนฺเต ภควโต อจฺฉริยํ อพฺภุตธมฺมํ ธาเรมิ. (16)}

\prose{สมฺมุขา เมตํ ภนฺเต ภควโต สุตํ, สมฺมุขา ปฏิคฺคหิตํ “ยทา อานนฺท โพธิสตฺโต มาตุกุจฺฉิมฺหา นิกฺขมติ, เทฺว อุทกสฺส ธารา อนฺตลิกฺขา ปาตุภวนฺติ, เอกา สีตสฺส เอกา อุณฺหสฺส.\csromanspacer{} เยน โพธิสตฺตสฺส อุทกกิจฺจํ กโรนฺติ มาตุ จา”ติ.\csromanspacer{} ยมฺปิ ภนฺเต ฯเปฯ อิทํปาหํ ภนฺเต ภควโต อจฺฉริยํ อพฺภุตธมฺมํ ธาเรมิ. (17)}

\proseitem{207}{สมฺมุขา เมตํ ภนฺเต ภควโต สุตํ, สมฺมุขา ปฏิคฺคหิตํ “ยมฺปติชาโต อานนฺท โพธิสตฺโต สเมหิ ปาเทหิ ปถวิยํ ปติฏฺฐหิตฺวา อุตฺตราภิมุโข สตฺตปทวีติหาเรน คจฺฉติ, เสตมฺหิ ฉตฺเต อนุธาริยมาเน สพฺโพ จ ทิสา วิโลเกติ, อาสภิญฺจ วาจํ ภาสติ ‘อคฺโคหมสฺมิ โลกสฺส, เชฏฺโฐหมสฺมิ โลกสฺส, เสฏฺโฐหมสฺมิ โลกสฺส, อยมนฺติมา ชาติ, นตฺถิ ทานิ ปุนพฺภโว’ติ”.\csromanspacer{} ยมฺปิ ภนฺเต ฯเปฯ อิทํปาหํ ภนฺเต ภควโต อจฺฉริยํ อพฺภุตธมฺมํ ธาเรมิ. (18)}

\csromanmatikaanchor{mtk.29}
\csromantocmarkref{subsection}{4. พากุลสุตฺต}{mtk.29}
\bookmark[dest={mtk.29},level=2]{4. พากุลสุตฺต}
\csromanheader{2}{4. พากุลสุตฺต (124)}
\prose{\csromanfolio{165}สมฺมุขา เมตํ ภนฺเต ภควโต สุตํ, สมฺมุขา ปฏิคฺคหิตํ “ยทา อานนฺท โพธิสตฺโต มาตุกุจฺฉิมฺหา นิกฺขมติ, อถ สเทวเก โลเก สมารเก สพฺรหฺมเก สสฺสมณพฺราหฺมณิยา ปชาย สเทวมนุสฺสาย อปฺปมาโณ อุฬาโร โอภาโส โลเก ปาตุภวติ อติกฺกมฺเมว เทวานํ เทวานุภาวํ.\csromanspacer{} ยาปิ ตา โลกนฺตริกา อฆา อสํวุตา อนฺธการา อนฺธการติมิสา.\csromanspacer{} ยตฺถปิเม จนฺทิมสูริยา เอวํมหิทฺธิกา เอวํมหานุภาวา อาภาย นานุโภนฺติ, ตตฺถปิ อปฺปมาโณ อุฬาโร โอภาโส โลเก ปาตุภวติ อติกฺกมฺเมว เทวานํ เทวานุภาวํ.\csromanspacer{} เยปิ ตตฺถ สตฺตา อุปปนฺนา, เตปิ เตโนภาเสน อญฺญมญฺญํ สญฺชานนฺติ ‘อญฺเญปิ กิร โภ สนฺติ สตฺตา อิธูปปนฺนา’ติ.\csromanspacer{} อยญฺจ ทสสหสฺสี โลกธาตุ สงฺกมฺปติ สมฺปกมฺปติ สมฺปเวธติ, อปฺปมาโณ จ อุฬาโร โอภาโส โลเก ปาตุภวติ อติกฺกมฺเมว เทวานํ เทวานุภาวนฺ”ติ.\csromanspacer{} ยมฺปิ ภนฺเต ฯเปฯ อิทํปาหํ ภนฺเต ภควโต อจฺฉริยํ อพฺภุตธมฺมํ ธาเรมีติ. (19)}

\proseitem{208}{ตสฺมาติห ตฺวํ อานนฺท อิทมฺปิ ตถาคตสฺส อจฺฉริยํ อพฺภุตธมฺมํ ธาเรหิ.\csromanspacer{} อิธานนฺท ตถาคตสฺส วิทิตา เวทนา อุปฺปชฺชนฺติ, วิทิตา อุปฏฺฐหนฺติ, วิทิตา อพฺภตฺถํ คจฺฉนฺติ, วิทิตา สญฺญา อุปฺปชฺชนฺติ ฯเปฯ วิทิตา วิตกฺกา อุปฺปชฺชนฺติ, วิทิตา อุปฏฺฐหนฺติ, วิทิตา อพฺภตฺถํ คจฺฉนฺติ.\csromanspacer{} อิทมฺปิ โข ตฺวํ อานนฺท ตถาคตสฺส อจฺฉริยํ อพฺภุตธมฺมํ ธาเรหีติ.\csromanspacer{} ยมฺปิ ภนฺเต ภควโต วิทิตา เวทนา อุปฺปชฺชนฺติ, วิทิตา อุปฏฺฐหนฺติ, วิทิตา อพฺภตฺถํ คจฺฉนฺติ, วิทิตา สญฺญา ฯเปฯ วิทิตา วิตกฺกา อุปฺปชฺชนฺติ, วิทิตา อุปฏฺฐหนฺติ, วิทิตา อพฺภตฺถํ คจฺฉนฺติ.\csromanspacer{} อิทํปาหํ ภนฺเต ภควโต อจฺฉริยํ อพฺภุตธมฺมํ ธาเรมีติ. (20)}

\prose{อิทมโวจ อายสฺมา อานนฺโท.\csromanspacer{} สมนุญฺโญ สตฺถา อโหสิ.\csromanspacer{} อตฺตมนา จ เต ภิกฺขู อายสฺมโต อานนฺทสฺส ภาสิตํ อภินนฺทุนฺติ.\csromanspacer{} อจฺฉริย-อพฺภุตสุตฺตํ นิฏฺฐิตํ ตติยํ.\textbf{4.}\csromanspacer{}\textbf{ พากุลสุตฺต}}
\csromansectionrule

\proseitem{209}{เอวํ เม สุตํ\csromandash{}เอกํ สมยํ อายสฺมา พากุโล\footnote{พกฺกุโล (สี, สฺยา, กํ, อิ)} ราชคเห วิหรติ เวฬุวเน กลนฺทกนิวาเป.\csromanspacer{} อถ โข อเจลกสฺสโป \csromanfolio{166}อายสฺมโต พากุลสฺส ปุราณคิหิสหาโย เยนายสฺมา พากุโล เตนุปสงฺกมิ, อุปสงฺกมิตฺวา อายสฺมตา พากุเลน สทฺธิํ สมฺโมทิ, สมฺโมทนียํ กถํ สารณียํ วีติสาเรตฺวา เอกมนฺตํ นิสีทิ, เอกมนฺตํ นิสินฺโน โข อเจลกสฺสโป อายสฺมนฺตํ พากุลํ เอตทโวจ\csromandash{}}

\prose{“กีวจิรํ ปพฺพชิโตสิ อาวุโส พากุลา”ติ.\csromanspacer{} อสีติ เม อาวุโส วสฺสานิ ปพฺพชิตสฺสาติ.\csromanspacer{} อิเมหิ ปน เต อาวุโส พากุล อสีติยา วสฺเสหิ กติกฺขตฺตุํ เมถุโน ธมฺโม ปฏิเสวิโตติ.\csromanspacer{} น โข มํ อาวุโส กสฺสป เอวํ ปุจฺฉิตพฺพํ “อิเมหิ ปน เต อาวุโส พากุล อสีติยา วสฺเสหิ กติกฺขตฺตุํ เมถุโน ธมฺโม ปฏิเสวิโต”ติ.\csromanspacer{} เอวญฺจ โข มํ อาวุโส กสฺสป ปุจฺฉิตพฺพํ “อิเมหิ ปน เต อาวุโส พากุล อสีติยา วสฺเสติ กติกฺขตฺตุํ กามสญฺญา อุปฺปนฺนปุพฺพา”ติ. ( )\footnote{(อิเมหิ ปน เต อาวุโส พกฺกุล อสีติยา วสฺเสหิ กติกฺขตฺตุํ กามสญฺญา อุปฺปนฺนปุพฺพาติ.) (สี, อิ)}}

\proseitem{210}{อสีติ เม อาวุโส วสฺสานิ ปพฺพชิตสฺส นาภิชานามิ กามสญฺญํ อุปฺปนฺนปุพฺพํ.\csromanspacer{} ยํปายสฺมา พากุโล อสีติยา วสฺเสหิ นาภิชานาติ กามสญฺญํ อุปฺปนฺนปุพฺพํ, อิทมฺปิ มยํ อายสฺมโต พากุลสฺส อจฺฉริยํ อพฺภุตธมฺมํ ธาเรม. (1)}

\prose{อสีติ เม อาวุโส วสฺสานิ ปพฺพชิตสฺส นาภิชานามิ พฺยาปาทสญฺญํ ฯเปฯ วิหิํสาสญฺญํ อุปฺปนฺนปุพฺพํ.\csromanspacer{} ยํปายสฺมา พากุโล อสีติยา วสฺเสหิ นาภิชานาติ วิหิํสาสญฺญํ อุปฺปนฺนปุพฺพํ, อิทมฺปิ มยํ อายสฺมโต พากุลสฺส อจฺฉริยํ อพฺภุตธมฺมํ ธาเรม. (2-3)}

\prose{อสีติ เม อาวุโส วสฺสานิ ปพฺพชิตสฺส นาภิชานามิ กามวิตกฺกํ อุปฺปนฺนปุพฺพํ.\csromanspacer{} ยํปายสฺมา พากุโล อสีติยา วสฺเสหิ นาภิชานาติ กามวิตกฺกํ อุปฺปนฺนปุพฺพํ, อิทมฺปิ มยํ อายสฺมโต พากุลสฺส อจฺฉริยํ อพฺภุตธมฺมํ ธาเรม. (4)}

\prose{อสีติ เม อาวุโส วสฺสานิ ปพฺพชิตสฺส นาภิชานามิ พฺยาปาทวิตกฺกํ ฯเปฯ วิหิํสาวิตกฺกํ อุปฺปนฺนปุพฺพํ.\csromanspacer{} ยํปายสฺมา พากุโล อสีติยา}

\csromanheader{2}{4. พากุลสุตฺต (124)}
\prose{\csromanfolio{167}วสฺเสหิ นาภิชานาติ วิหิํสาวิตกฺกํ อุปฺปนฺนปุพฺพํ, อิทมฺปิ มยํ อายสฺมโต พากุลสฺส อจฺฉริยํ อพฺภุตธมฺมํ ธาเรม. (5-6)}

\proseitem{211}{อสีติ เม อาวุโส วสฺสานิ ปพฺพชิตสฺส นาภิชานามิ คหปติจีวรํ สาทิตา.\csromanspacer{} ยํปายสฺมา พากุโล อสีติยา วสฺเสหิ นาภิชานาติ คหปติจีวรํ สาทิตา, อิทมฺปิ มยํ อายสฺมโต พากุลสฺส อจฺฉริยํ อพฺภุตธมฺมํ ธาเรม. (7)}

\prose{อสีติ เม อาวุโส วสฺสานิ ปพฺพชิตสฺส นาภิชานามิ สตฺเถน จีวรํ ฉินฺทิตา.\csromanspacer{} ยํปายสฺมา พากุโล อสีติยา วสฺเสหิ นาภิชานาติ สตฺเถน จีวรํ ฉินฺทิตา ฯเปฯ ธาเรม. (8)}

\prose{อสีติ เม อาวุโส วสฺสานิ ปพฺพชิตสฺส นาภิชานามิ สูจิยา จีวรํ สิพฺพิตา ฯเปฯ นาภิชานามิ รชเนน จีวรํ รชิตา.\csromanspacer{} นาภิชานามิ กถิเน\footnote{กฐิเน (สี, สฺยา, กํ, อิ)} จีวรํ สิพฺพิตา.\csromanspacer{} นาภิชานามิ สพฺรหฺมจารีนํ จีวรกมฺเม วิจาริตา\footnote{สพฺรหฺมจารี จีวรกมฺเม พฺยาปาริตา (สี, อิ)}.\csromanspacer{} นาภิชานามิ นิมนฺตนํ สาทิตา.\csromanspacer{} นาภิชานามิ เอวรูปํ จิตฺตํ อุปฺปนฺนปุพฺพํ “อโห วต มํ โกจิ นิมนฺเตยฺยา”ติ.\csromanspacer{} นาภิชานามิ อนฺตรฆเร นิสีทิตา.\csromanspacer{} นาภิชานามิ อนฺตรฆเร ภุญฺชิตา.\csromanspacer{} นาภิชานามิ มาตุคามสฺส อนุพฺยญฺชนโส นิมิตฺตํ คเหตา.\csromanspacer{} นาภิชานามิ มาตุคามสฺส ธมฺมํ เทสิตา อนฺตมโส จตุปฺปทมฺปิ คาถํ.\csromanspacer{} นาภิชานามิ ภิกฺขุนุปสฺสยํ อุปสงฺกมิตา.\csromanspacer{} นาภิชานามิ ภิกฺขุนิยา ธมฺมํ เทสิตา.\csromanspacer{} นาภิชานามิ สิกฺขมานาย ธมฺมํ เทสิตา.\csromanspacer{} นาภิชานามิ สามเณริยา ธมฺมํ เทสิตา.\csromanspacer{} นาภิชานามิ ปพฺพาเชตา.\csromanspacer{} นาภิชานามิ อุปสมฺปาเทตา.\csromanspacer{} นาภิชานามิ นิสฺสยํ ทาตา.\csromanspacer{} นาภิชานามิ สามเณรํ อุปฏฺฐาเปตา.\csromanspacer{} นาภิชานามิ ชนฺตาฆเร นฺหายิตา.\csromanspacer{} นาภิชานามิ จุณฺเณน นฺหายิตา.\csromanspacer{} นาภิชานามิ สพฺรหฺมจารี คตฺตปริกมฺเม วิจาริตา\footnote{พฺยาปาริตา (สี, อิ)}.\csromanspacer{} นาภิชานามิ อาพาธํ อุปฺปนฺนปุพฺพํ อนฺตมโส คทฺทูหนมตฺตมฺปิ.\csromanspacer{} นาภิชานามิ เภสชฺชํ อุปหริตา อนฺตมโส หริตกิขณฺฑมฺปิ.\csromanspacer{} นาภิชานามิ อปสฺเสนกํ อปสฺสยิตา.\csromanspacer{} นาภิชานามิ เสยฺยํ กปฺเปตา.\csromanspacer{} ยํปายสฺมา ฯเปฯ ธาเรม. (9-33)}

\prose{อสีติ เม อาวุโส วสฺสานิ ปพฺพชิตสฺส นาภิชานามิ คามนฺตเสนาสเน วสฺสํ อุปคนฺตา.\csromanspacer{} ยํปาสฺมา พากุโล อสีติยา \csromanfolio{168}วสฺเสหิ นาภิชานาติ คามนฺตเสนาสเน วสฺสํ อุปคนฺตา, อิทมฺปิ มยํ อายสฺมโต พากุลสฺส อจฺฉริยํ อพฺภุตธมฺมํ ธาเรม. (34)}

\prose{สตฺตาหเมว โข อหํ อาวุโส สรโณ รฏฺฐปิณฺฑํ ภุญฺชิํ, อถ อฏฺฐมิยํ อญฺญา อุทปาทิ.\csromanspacer{} ยํปายสฺมา พากุโล สตฺตาหเมว สรโณ รฏฺฐปิณฺฑํ ภุญฺชิ, อถ อฏฺฐมิยํ อญฺญา อุทปาทิ, อิทมฺปิ มยํ อายสฺมโต พากุลสฺส อจฺฉริยํ อพฺภุตธมฺมํ ธาเรม. (35)}

\proseitem{212}{ลเภยฺยาหํ อาวุโส พากุล อิมสฺมิํ ธมฺมวินเย ปพฺพชฺชํ, ลเภยฺยํ อุปสมฺปทนฺติ.\csromanspacer{} อลตฺถ โข อเจลกสฺสโป อิมสฺมิํ ธมฺมวินเย ปพฺพชฺชํ, อลตฺถ อุปสมฺปทํ.\csromanspacer{} อจิรูปสมฺปนฺโน ปนายสฺมา กสฺสโป เอโก วูปกฏฺโฐ อปฺปมตฺโต อาตาปี ปหิตตฺโต วิหรนฺโต นจิรสฺเสว, ยสฺสตฺถาย กุลปุตฺตา สมฺมเทว อคารสฺมา อนคาริยํ ปพฺพชนฺติ, ตทนุตฺตรํ พฺรหฺมจริยปริโยสานํ ทิฏฺเฐว ธมฺเม สยํ อภิญฺญา สจฺฉิกตฺวา อุปสมฺปชฺช วิหาสิ, “ขีณา ชาติ, วุสิตํ พฺรหฺมจริยํ, กตํ กรณียํ, นาปรํ อิตฺถตฺตายา”ติ อพฺภญฺญาสิ.\csromanspacer{} อญฺญตโร โข ปนายสฺมา กสฺสโป อรหตํ อโหสิ.}

\prose{อถ โข อายสฺมา พากุโล อปเรน สมเยน อวาปุรณํ\footnote{อปาปุรณํ (สี, สฺยา, กํ, อิ)} อาทาย วิหาเรน วิหารํ อุปสงฺกมิตฺวา เอวมาห “อภิกฺกมถายสฺมนฺโต, อภิกฺกมถายสฺมนฺโต, อชฺช เม ปรินิพฺพานํ ภวิสฺสตี”ติ.\csromanspacer{} ยํปายสฺมา พากุโล อวาปุรณํ อาทาย วิหาเรน วิหารํ อุปสงฺกมิตฺวา เอวมาห “อภิกฺกมถายสฺมนฺโต, อภิกฺกมถายสฺมนฺโต, อชฺช เม ปรินิพฺพานํ ภวิสฺสตี”ติ, อิทมฺปิ มยํ อายสฺมโต พากุลสฺส อจฺฉริยํ อพฺภุตธมฺมํ ธาเรม. (36)}

\prose{อายสฺมา พากุโล มชฺเฌ ภิกฺขุสํฆสฺส นิสินฺนโกว ปรินิพฺพายิ.\csromanspacer{} ยํปายสฺมา พากุโล มชฺเฌ ภิกฺขุสํฆสฺส นิสินฺนโกว ปรินิพฺพายิ, อิทมฺปิ มยํ อายสฺมโต พากุลสฺส อจฺฉริยํ อพฺภุตธมฺมํ ธาเรมาติ. (37)}

\nitthitam{พากุลสุตฺตํ นิฏฺฐิตํ จตุตฺถํ.}
\csromanheader{2}{5. ทนฺตภูมิสุตฺต (125)}
\csromanmatikaanchor{mtk.30}
\csromantocmarkref{subsection}{5. ทนฺตภูมิสุตฺต}{mtk.30}
\bookmark[dest={mtk.30},level=2]{5. ทนฺตภูมิสุตฺต}
\csromanheader{2}{5. ทนฺตภูมิสุตฺต}
\proseitem{213}{\csromanfolio{169}เอวํ เม สุตํ\csromandash{}เอกํ สมยํ ภควา ราชคเห วิหรติ เวฬุวเน กลนฺทกนิวาเป.\csromanspacer{} เตน โข ปน สมเยน อจิรวโต สมณุทฺเทโส อรญฺญกุฏิกายํ วิหรติ.\csromanspacer{} อถ โข ชยเสโน ราชกุมาโร ชงฺฆาวิหารํ อนุจงฺกมมาโน อนุวิจรมาโน เยน อจิรวโต สมณุทฺเทโส เตนุปสงฺกมิ, อุปสงฺกมิตฺวา อจิรวเตน สมณุทฺเทเสน สทฺธิํ สมฺโมทิ, สมฺโมทนียํ กถํ สารณียํ วีติสาเรตฺวา เอกมนฺตํ นิสีทิ, เอกมนฺตํ นิสินฺโน โข ชยเสโน ราชกุมาโร อจิรวตํ สมณุทฺเทสํ เอตทโวจ\csromandash{}}

\prose{สุตํ เมตํ โภ อคฺคิเวสฺสน “อิธ ภิกฺขุ อปฺปมตฺโต อาตาปี ปหิตตฺโต วิหรนฺโต ผุเสยฺย จิตฺตสฺส เอกคฺคตนฺ”ติ.\csromanspacer{} เอวเมตํ ราชกุมาร, เอวเมตํ ราชกุมาร, อิธ ภิกฺขุ อปฺปมตฺโต อาตาปี ปหิตตฺโต วิหรนฺโต ผุเสยฺย จิตฺตสฺส เอกคฺคตนฺติ.\csromanspacer{} สาธุ เม ภวํ อคฺคิเวสฺสโน ยถาสุตํ ยถาปริยตฺตํ ธมฺมํ เทเสตูติ.\csromanspacer{} น โข เต อหํ ราชกุมาร สกฺโกมิ ยถาสุตํ ยถาปริยตฺตํ ธมฺมํ เทเสตุํ, อหญฺจ หิ เต ราชกุมาร ยถาสุตํ ยถาปริยตฺตํ ธมฺมํ เทเสยฺยํ, ตฺวญฺจ เม ภาสิตสฺส อตฺถํ น อาชาเนยฺยาสิ.\csromanspacer{} โส มมสฺส กิลมโถ สา มมสฺส วิเหสาติ.\csromanspacer{} เทเสตุ เม ภวํ อคฺคิเวสฺสโน ยถาสุตํ ยถาปริยตฺตํ ธมฺมํ, อปฺเปวนามาหํ โภโต อคฺคิเวสฺสนสฺส ภาสิตสฺส อตฺถํ อาชาเนยฺยนฺติ.\csromanspacer{} เทเสยฺยํ โข เต อหํ ราชกุมาร ยถาสุตํ ยถาปริยตฺตํ ธมฺมํ, สเจ เม ตฺวํ ภาสิตสฺส อตฺถํ อาชาเนยฺยาสิ, อิจฺเจตํ กุสลํ.\csromanspacer{} โน เจ เม ตฺวํ ภาสิตสฺส อตฺถํ อาชาเนยฺยาสิ, ยถาสเก ติฏฺเฐยฺยาสิ, น มํ ตตฺถ อุตฺตริํ ปฏิปุจฺเฉยฺยาสีติ.\csromanspacer{} เทเสตุ เม ภวํ อคฺคิเวสฺสโน ยถาสุตํ ยถาปริยตฺตํ ธมฺมํ, สเจ อหํ โภโต อคฺคิเวสนสฺส ภาสิตสฺส อตฺถํ อาชานิสฺสามิ\footnote{อาชาเนยฺยามิ (ก)}, อิจฺเจตํ กุสลํ.\csromanspacer{} โน เจ อหํ โภโต อคฺคิเวสฺสนสฺส ภาสิตสฺส อตฺถํ อาชานิสฺสามิ, ยถาสเก ติฏฺฐิสฺสามิ\footnote{ติฏฺเฐยฺยามิ (ก)}, นาหํ ตตฺถ ภวนฺตํ อคฺคิเวสฺสนํ อุตฺตริํ ปฏิปุจฺฉิสฺสามีติ.}

\proseitem{214}{\csromanfolio{170}อถ โข อจิรวโต สมณุทฺเทโส ชยเสนสฺส ราชกุมารสฺส ยถาสุตํ ยถาปริยตฺตํ ธมฺมํ เทเสสิ.\csromanspacer{} เอวํ วุตฺเต ชยเสโน ราชกุมาโร อจิรวตํ สมณุทฺเทสํ เอตทโวจ “อฏฺฐานเมตํ โภ อคฺคิเวสฺสน อนวกาโส ยํ ภิกฺขุ อปฺปมตฺโต อาตาปี ปหิตตฺโต วิหรนฺโต ผุเสยฺย จิตฺตสฺส เอกคฺคตนฺ”ติ.\csromanspacer{} อถ โข ชยเสโน ราชกุมาโร อจิรวตสฺส สมณุทฺเทสสฺส อฏฺฐานตญฺจ อนวกาสตญฺจ ปเวเทตฺวา อุฏฺฐายาสนา ปกฺกามิ.}

\prose{อถ โข อจิรวโต สมณุทฺเทโส อจิรปกฺกนฺเต ชยเสเน ราชกุมาเร เยน ภควา เตนุปสงฺกมิ, อุปสงฺกมิตฺวา ภควนฺตํ อภิวาเทตฺวา เอกมนฺตํ นิสีทิ, เอกมนฺตํ นิสินฺโน โข อจิรวโต สมณุทฺเทโส ยาวตโก อโหสิ ชยเสเนน ราชกุมาเรน สทฺธิํ กถาสลฺลาโป, ตํ สพฺพํ ภควโต อาโรเจสิ.}

\prose{เอวํ วุตฺเต ภควา อจิรวตํ สมณุทฺเทสํ เอตทโวจ\csromandash{}ตํ กุเตตฺถ อคฺคิเวสฺสน ลพฺภา, ยํ ตํ เนกฺขมฺเมน ญาตพฺพํ เนกฺขมฺเมน ทฏฺฐพฺพํ เนกฺขมฺเมน ปตฺตพฺพํ เนกฺขมฺเมน สจฺฉิกาตพฺพํ, ตํ วต ชยเสโน ราชกุมาโร กามมชฺเฌ วสนฺโต กาเม ปริภุญฺชนฺโต กามวิตกฺเกหิ ขชฺชมาโน กามปริฬาเหน ปริฑยฺหมาโน กามปริเยสนาย อุสฺสุโก\footnote{อุสฺสุกฺโก (สพฺพตฺถ)} ญสฺสติ วา ทกฺขติ วา สจฺฉิ วา กริสฺสตีติ เนตํ ฐานํ วิชฺชติ.}

\proseitem{215}{เสยฺยถาปิสฺสุ อคฺคิเวสฺสน เทฺว หตฺถิทมฺมา วา อสฺสทมฺมา วา โคทมฺมา วา สุทนฺตา สุวินีตา, เทฺว หตฺถิทมฺมา วา อสฺสทมฺมา วา โคทมฺมา วา อทนฺตา อวินีตา.\csromanspacer{} ตํ กิํ มญฺญสิ อคฺคิเวสฺสน, เย เต เทฺว หตฺถิทมฺมา วา อสฺสทมฺมาวา โคทมฺมา วา สุทนฺตา สุวินีตา, อปิ นุ เต ทนฺตาว ทนฺตการณํ คจฺเฉยฺยุํ, ทนฺตาว ทนฺตภูมิํ สมฺปาปุเณยฺยุนฺติ.\csromanspacer{} เอวํ ภนฺเต.\csromanspacer{} เย ปน เต เทฺว หตฺถิทมฺมา วา อสฺสทมฺมา วา โคทมฺมา วา อทนฺตา อวินีตา, อปิ นุ เต อทนฺตาว ทนฺตการณํ คจฺเฉยฺยุํ, อทนฺตาว ทนฺตภูมิํ สมฺปาปุเณยฺยุํ.\csromanspacer{} เสยฺยถาปิ เต เทฺว หตฺถิทมฺมา วา อสฺสทมฺมา วา โคทมฺมา วา สุทนฺตา สุวินีตาติ.\csromanspacer{} โน เหตํ ภนฺเต.\csromanspacer{} เอวเมว โข อคฺคิเวสฺสน ยํ ตํ เนกฺขมฺเมน ญาตพฺพํ เนกฺขมฺเมน ทฏฺฐพฺพํ เนกฺขมฺเมน ปตฺตพฺพํ เนกฺขมฺเมน สจฺฉิกาตพฺพํ, ตํ วต ชยเสโน ราชกุมาโร กามมชฺเฌ วสนฺโต}

\csromanheader{2}{5. ทนฺตภูมิสุตฺต (125)}
\prose{\csromanfolio{171}กาเม ปริภุญฺชนฺโต กามวิตกฺเกหิ ขชฺชมาโน กามปริฬาเหน ปริฑยฺหมาโน กามปริเยสนาย อุสฺสุโก ญสฺสติ วา ทกฺขติ วา สจฺฉิ วา กริสฺสตีติ เนตํ ฐานํ วิชฺชติ.}

\proseitem{216}{เสยฺยถาปิ อคฺคิเวสฺสน คามสฺส วา นิคมสฺส วา อวิทูเร มหาปพฺพโต, ตเมนํ เทฺว สหายกา ตมฺหา คามา วา นิคมา วา นิกฺขมิตฺวา หตฺถวิลงฺฆเกน เยน โส ปพฺพโต เตนุปสงฺกเมยฺยุํ, อุปสงฺกมิตฺวา เอโก สหายโก เหฏฺฐา ปพฺพตปาเท ติฏฺเฐยฺย, เอโก สหายโก อุปริปพฺพตํ อาโรเหยฺย, ตเมนํ เหฏฺฐา ปพฺพตปาเท ฐิโต สหายโก อุปริปพฺพเต ฐิตํ สหายกํ เอวํ วเทยฺย “ยํ สมฺม กิํ ตฺวํ ปสฺสสิ อุปริปพฺพเต ฐิโต”ติ.\csromanspacer{} โส เอวํ วเทยฺย “ปสฺสามิ โข อหํ สมฺม อุปริปพฺพเต ฐิโต อารามรามเณยฺยกํ วนรามเณยฺยกํ ภูมิรามเณยฺยกํ โปกฺขรณีรามเณยฺยกนฺ”ติ.}

\prose{โส เอวํ วเทยฺย “อฏฺฐานํ โข เอตํ สมฺม อนวกาโส, ยํ ตฺวํ อุปริปพฺพเต ฐิโต ปสฺเสยฺยาสิ อารามรามเณยฺยกํ วนรามเณยฺยกํ ภูมิรามเณยฺยกํ โปกฺขรณีรามเณยฺยกนฺ”ติ.\csromanspacer{} ตเมนํ อุปริปพฺพเต ฐิโต สหายโก เหฏฺฐิมปพฺพตปาทํ โอโรหิตฺวา ตํ สหายกํ พาหายํ คเหตฺวา อุปริปพฺพตํ อาโรเปตฺวา มุหุตฺตํ อสฺสาเสตฺวา เอวํ วเทยฺย “ยํ สมฺม กิํ ตฺวํ ปสฺสสิ อุปริปพฺพเต ฐิโต”ติ.\csromanspacer{} โส เอวํ วเทยฺย “ปสฺสามิ โข อหํ สมฺม อุปริปพฺพเต ฐิโต อารามรามเณยฺยกํ วนรามเณยฺยกํ ภูมิรามเณยฺยกํ โปกฺขรณีรามเณยฺยกนฺ”ติ.}

\prose{โส เอวํ วเทยฺย “อิทาเนว โข เต สมฺม ภาสิตํ มยํ เอวํ อาชานาม, อฏฺฐานํ โข เอตํ สมฺม อนวกาโส, ยํ ตฺวํ อุปริปพฺพเต ฐิโต ปสฺเสยฺยาสิ อารามรามเณยฺยกํ วนรามเณยฺยกํ ภูมิรามเณยฺยกํ โปกฺขรณีรามเณยฺยกนฺ”ติ.\csromanspacer{} อิทาเนว จ ปน เต ภาสิตํ มยํ เอวํ อาชานาม “ปสฺสามิ โข อหํ สมฺม อุปริปพฺพเต ฐิโต อารามรามเณยฺยกํ วนรามเณยฺยกํ ภูมิรามเณยฺยกํ โปกฺขรณีรามเณยฺยกนฺ”ติ.\csromanspacer{} โส เอวํ วเทยฺย “ตถา หิ ปนาหํ สมฺม อิมินา มหตา ปพฺพเตน อาวุโต\footnote{อาวโฏ (สี-ฏฺฐ, อิ), อาวุโฏ (สฺยา, กํ, ก)} ทฏฺเฐยฺยํ นาทฺทสนฺ”ติ.}

\prose{\csromanfolio{172}อโต มหนฺตตเรน อคฺคิเวสฺสน อวิชฺชาขนฺเธน ชยเสโน ราชกุมาโร อาวุโต นิวุโต\footnote{นิวุโฏ (สฺยา, กํ, อิ, ก)} โอผุโฏ\footnote{โอวุโต (สี), โอวุโฏ (สฺยา, กํ, อิ)} ปริโยนทฺโธ, โส วต ยํ ตํ เนกฺขมฺเมน ญาตพฺพํ เนกฺขมฺเมน ทฏฺฐพฺพํ เนกฺขมฺเมน ปตฺตพฺพํ เนกฺขมฺเมน สจฺฉิกาตพฺพํ, ตํ วต ชยเสโน ราชกุมาโร กามมชฺเฌ วสนฺโต กาเม ปริภุญฺชนฺโต กามวิตกฺเกหิ ขชฺชมาโน กามปริฬาเหน ปริฑยฺหมาโน กามปริเยสนาย อุสฺสุโก ญสฺสติ วา ทกฺขติ วา สจฺฉิ วา กริสฺสตีติ เนตํ ฐานํ วิชฺชติ.\csromanspacer{} สเจ โข ตํ อคฺคิเวสฺสน ชยเสนสฺส ราชกุมารสฺส อิมา เทฺว อุปมา ปฏิภาเยยฺยุํ\footnote{ปฏิภาเสยฺยุํ (สี, สฺยา, กํ, อิ)}, อนจฺฉริยํ เต ชยเสโน ราชกุมาโร ปสีเทยฺย, ปสนฺโน จ เต ปสนฺนาการํ กเรยฺยาติ, กุโต ปน มํ ภนฺเต ชยเสนสฺส ราชกุมารสฺส อิมา เทฺว อุปมา ปฏิภายิสฺสนฺติ\footnote{ปฏิภาสิสฺสนฺติ (สี, สฺยา, กํ, อิ)} อนจฺฉริยา ปุพฺเพ อสฺสุตปุพฺพา เสยฺยถาปิ ภควนฺตนฺติ.}

\proseitem{217}{เสยฺยถาปิ อคฺคิเวสฺสน ราชา ขตฺติโย มุทฺธาวสิตฺโต นาควนิกํ อามนฺเตติ “เอหิ ตฺวํ สมฺม นาควนิก รญฺโญ นาคํ อภิรุหิตฺวา นาควนํ ปวิสิตฺวา อารญฺญกํ นาคํ อติปสฺสิตฺวา รญฺโญ นาคสฺส คีวายํ อุปนิพนฺธาหี”ติ. “เอวํ เทวา”ติ โข อคฺคิเวสฺสน นาควนิโก รญฺโญ ขตฺติยสฺส มุทฺธาวสิตฺตสฺส ปฏิสฺสุตฺวา รญฺโญ นาคํ อภิรุหิตฺวา นาควนํ ปวิสิตฺวา อารญฺญกํ นาคํ อติปสฺสิตฺวา รญฺโญ นาคสฺส คีวายํ อุปนิพนฺธติ, ตเมนํ รญฺโญ นาโค อพฺโภกาสํ นีหรติ.\csromanspacer{} เอตฺตาวตา โข อคฺคิเวสฺสน อารญฺญโก นาโค อพฺโภกาสํ คโต โหติ.\csromanspacer{} เอตฺถเคธา\footnote{เอตเคธา (สี, อิ)} หิ อคฺคิเวสฺสน อารญฺญกา นาคา ยทิทํ นาควนํ.\csromanspacer{} ตเมนํ นาควนิโก รญฺโญ ขตฺติยสฺส มุทฺธาวสิตฺตสฺส อาโรเจสิ “อพฺโภกาสคโต โข\footnote{โข เต (สฺยา, กํ, ก)} เทว อารญฺญโก นาโค”ติ.\csromanspacer{} อถ โข อคฺคิเวสฺสน ตเมนํ ราชา ขตฺติโย มุทฺธาวสิตฺโต หตฺถิทมกํ อามนฺเตสิ “เอหิ ตฺวํ สมฺม หตฺถิทมก อารญฺญกํ นาคํ ทมยาหิ, อารญฺญกานญฺเจว สีลานํ อภินิมฺมทนาย อารญฺญกานญฺเจว สรสงฺกปฺปานํ อภินิมฺมทนาย อารญฺญกานญฺเจว ทรถกิลมถปริฬาหานํ อภินิมฺมทนาย คามนฺเต อภิรมาปนาย มนุสฺสกนฺเตสุ สีเลสุ สมาทปนายา”ติ\footnote{สมาทาปนายาติ (?)}.}

\csromanheader{2}{5. ทนฺตภูมิสุตฺต (125)}
\prose{\csromanfolio{173}“เอวํ เทวา”ติ โข อคฺคิเวสฺสน หตฺถิทมโก รญฺโญ ขตฺติยสฺส มุทฺธาวสิตฺตสฺส ปฏิสฺสุตฺวา มหนฺตํ ถมฺภํ ปถวิยํ นิขณิตฺวา อารญฺญกสฺส นาคสฺส คีวายํ อุปนิพนฺธติ อารญฺญกานญฺเจว สีลานํ อภินิมฺมทนาย อารญฺญกานญฺเจว สรสงฺกปฺปานํ อภินิมฺมทนาย อารญฺญกานญฺเจว ทรถกิลมถปริฬาหานํ อภินิมฺมทนาย คามนฺเต อภิรมาปนาย มนุสฺสกนฺเตสุ สีเลสุ สมาทปนาย.\csromanspacer{} ตเมนํ หตฺถิทมโก ยา สา วาจา เนลา กณฺณสุขา เปมนียา หทยงฺคมา โปรี พหุชนกนฺตา พหุชนมนาปา, ตถารูปาหิ วาจาหิ สมุทาจรติ.\csromanspacer{} ยโต โข อคฺคิเวสฺสน อารญฺญโก นาโค หตฺถิทมกสฺส ยา สา วาจา เนลา กณฺณสุขา เปมนียา หทยงฺคมา โปรี พหุชนกนฺตา พหุชนมนาปา, ตถารูปาหิ วาจาหิ สมุทาจริยมาโน สุสฺสูสติ, โสตํ โอทหติ, อญฺญา จิตฺตํ อุปฏฺฐาเปติ.\csromanspacer{} ตเมนํ หตฺถิทมโก อุตฺตริ ติณฆาโสทกํ อนุปฺปเวจฺฉติ.}

\prose{ยโต โข อคฺคิเวสฺสน อารญฺญโก นาโค หตฺถิทมกสฺส ติณฆาโสทกํ ปฏิคฺคณฺหาติ, ตตฺร หตฺถิทมกสฺส เอวํ โหติ “ชีวิสฺสติ โข\footnote{นุ โข (สี, ก)} ทานิ อารญฺญโก\footnote{รญฺโญ (สี, อิ)} นาโค”ติ.\csromanspacer{} ตเมนํ หตฺถิทมโก อุตฺตริ การณํ กาเรติ “อาทิย โภ นิกฺขิป โภ”ติ.\csromanspacer{} ยโต โข อคฺคิเวสฺสน อารญฺญโก นาโค หตฺถิทมกสฺส อาทานนิกฺเขเป วจนกโร โหติ โอวาทปฺปฏิกโร, ตเมนํ หตฺถิทมโก อุตฺตริ การณํ กาเรติ “อภิกฺกม โภ ปฏิกฺกม โภ”ติ.\csromanspacer{} ยโต โข อคฺคิเวสฺสน อารญฺญโก นาโค หตฺถิทมกสฺส อภิกฺกมปฏิกฺกมวจนกโร โหติ โอวาทปฺปฏิกโร, ตเมนํ หตฺถิทมโก อุตฺตริ การณํ กาเรติ “อุฏฺฐห โภ นิสีท โภ”ติ.\csromanspacer{} ยโต โข อคฺคิเวสฺสน อารญฺญโก นาโค หตฺถิทมกสฺส อุฏฺฐานนิสชฺชาย วจนกโร โหติ โอวาทปฺปฏิกโร, ตเมนํ หตฺถิทมโก อุตฺตริ อาเนญฺชํ นาม การณํ กาเรติ, มหนฺตสฺส ผลกํ โสณฺฑาย อุปนิพนฺธติ.\csromanspacer{} โตมรหตฺโถ จ ปุริโส อุปริคีวาย นิสินฺโน โหติ.\csromanspacer{} สมนฺตโต จ โตมรหตฺถา ปุริสา ปริวาเรตฺวา ฐิตา โหนฺติ.\csromanspacer{} หตฺถิทมโก จ ทีฆโตมรยฏฺฐิํ คเหตฺวา ปุรโต ฐิโต โหติ.\csromanspacer{} โส อาเนญฺชํ การณํ การิยมาโน เนว ปุริเม ปาเท โจเปติ, น ปจฺฉิเม ปาเท โจเปติ, น ปุริมกายํ โจเปติ, น ปจฺฉิมกายํ โจเปติ, น สีสํ โจเปติ, น กณฺเณ โจเปติ, น ทนฺเต}

\prose{\csromanfolio{174}โจเปติ, น นงฺคุฏฺฐํ โจเปติ, น โสณฺฑํ โจเปติ.\csromanspacer{} โส โหติ อารญฺญโก นาโค ขโม สตฺติปฺปหารานํ อสิปฺปหารานํ อุสุปฺปหารานํ สรปตฺตปฺปหารานํ\footnote{ปรสตฺถปฺปหารานํ (สี), ปรสตฺตุปฺปหารานํ (สฺยา, กํ, อิ)} เภริปณววํสสงฺขฑิณฺฑิมนินฺนาทสทฺทานํ\footnote{เภริปณวสงฺข ติณวนินฺนาทสทฺทานํ (อิ)} สพฺพวงฺกโทสนิหิตนินฺนีตกสาโว ราชารโห ราชโภคฺโค รญฺโญ องฺคนฺเตว สงฺขํ คจฺฉติ.}

\proseitem{218}{เอวเมว โข อคฺคิเวสฺสน อิธ ตถาคโต โลเก อุปฺปชฺชติ อรหํ สมฺมาสมฺพุทฺโธ วิชฺชาจรณสมฺปนฺโน สุคโต โลกวิทู อนุตฺตโร ปุริสทมฺมสารถิ สตฺถา เทวมนุสฺสานํ พุทฺโธ ภควา, โส อิมํ โลกํ สเทวกํ สมารกํ สพฺรหฺมกํ สสฺสมณพฺราหฺมณิํ ปชํ สเทวมนุสฺสํ สยํ อภิญฺญา สจฺฉิกตฺวา ปเวเทติ, โส ธมฺมํ เทเสติ อาทิกลฺยาณํ มชฺเฌกลฺยาณํ ปริโยสานกลฺยาณํ สาตฺถํ สพฺยญฺชนํ เกวลปริปุณฺณํ ปริสุทฺธํ พฺรหฺมจริยํ ปกาเสติ, ตํ ธมฺมํ สุณาติ คหปติ วา คหปติปุตฺโต วา อญฺญตรสฺมิํ วา กุเล ปจฺจาชาโต, โส ตํ ธมฺมํ สุตฺวา ตถาคเต สทฺธํ ปฏิลภติ, โส เตน สทฺธาปฏิลาเภน สมนฺนาคโต อิติ ปฏิสญฺจิกฺขติ “สมฺพาโธ ฆราวาโส รชาปโถ, อพฺโภกาโส ปพฺพชฺชา, นยิทํ สุกรํ อคารํ อชฺฌาวสตา เอกนฺตปริปุณฺณํ เอกนฺตปริสุทฺธํ สงฺขลิขิตํ พฺรหฺมจริยํ จริตุํ, ยํนูนาหํ เกสมสฺสุํ โอหาเรตฺวา กาสายานิ วตฺถานิ อจฺฉาเทตฺวา อคารสฺมา อนคาริยํ ปพฺพเชยฺยนฺ”ติ.}

\prose{โส อปเรน สมเยน อปฺปํ วา โภคกฺขนฺธํ ปหาย มหนฺตํ วา โภคกฺขนฺธํ ปหาย อปฺปํ วา ญาติปริวฏฺฏํ ปหาย มหนฺตํ วา ญาติปริวฏฺฏํ ปหาย เกสมสฺสุํ โอหาเรตฺวา กาสายานิ วตฺถานิ อจฺฉาเทตฺวา อคารสฺมา อนคาริยํ ปพฺพชติ.\csromanspacer{} เอตฺตาวตา โข อคฺคิเวสฺสน อริยสาวโก อพฺโภกาสคโต โหติ.\csromanspacer{} เอตฺถเคธา หิ อคฺคิเวสฺสน เทวมนุสฺสา ยทิทํ ปญฺจ กามคุณา, ตเมนํ ตถาคโต อุตฺตริํ วิเนติ “เอหิ ตฺวํ ภิกฺขุ สีลวา โหติ, ปาติโมกฺขสํวรสํวุโต วิหราหิ อาจารโคจรสมฺปนฺโน, อณุมตฺเตสุ วชฺเชสุ ภยทสฺสาวี สมาทาย สิกฺขสฺสุ สิกฺขาปเทสู”ติ.}

\prose{ยโต โข อคฺคิเวสฺสน อริยสาวโก สีลวา โหติ, ปาติโมกฺขสํวรสํวุโต วิหรติ อาจารโคจรสมฺปนฺโน, อณุมตฺเตสุ}

\csromanheader{2}{5. ทนฺตภูมิสุตฺต (125)}
\prose{\csromanfolio{175}วชฺเชสุ ภยทสฺสาวี สมาทาย สิกฺขติ สิกฺขาปเทสุ.\csromanspacer{} ตเมนํ ตถาคโต อุตฺตริํ วิเนติ “เอหิ ตฺวํ ภิกฺขุ อินฺทฺริเยสุ คุตฺตทฺวาโร โหติ, จกฺขุนา รูปํ ทิสฺวา มา นิมิตฺตคฺคาหี ฯเปฯ. (ยถา คณกโมคฺคลฺลานสุตฺตนฺเต, เอวํ วิตฺถาเรตพฺพานิ.)}

\proseitem{219}{โส อิเม ปญฺจ นีวรเณ ปหาย เจตโส อุปกฺกิเลเส ปญฺญาย ทุพฺพลีกรเณ กาเย กายานุปสฺสี วิหรติ อาตาปี สมฺปชาโน สติมา วิเนยฺย โลเก อภิชฺฌาโทมนสฺสํ.\csromanspacer{} เวทนาสุ ฯเปฯ จิตฺเต.\csromanspacer{} ธมฺเมสุ ธมฺมานุปสฺสี วิหรติ อาตาปี สมฺปชาโน สติมา วิเนยฺย โลเก อภิชฺฌาโทมนสฺสํ.\csromanspacer{} เสยฺยถาปิ อคฺคิเวสฺสน หตฺถิทมโก มหนฺตํ ถมฺภํ ปถวิยํ นิขณิตฺวา อารญฺญกสฺส นาคสฺส คีวายํ อุปนิพนฺธติ, อารญฺญกานญฺเจว สีลานํ อภินิมฺมทนาย อารญฺญกานญฺเจว สรสงฺกปฺปานํ อภินิมฺมทนาย อารญฺญกานญฺเจว ทรถกิลมถปริฬาหานํ อภินิมฺมทนาย คามนฺเต อภิรมาปนาย มนุสฺสกนฺเตสุ สีเลสุ สมาทปนาย.\csromanspacer{} เอวเมว โข อคฺคิเวสฺสน อริยสาวกสฺส อิเม จตฺตาโร สติปฏฺฐานา เจตโส อุปนิพนฺธนา โหนฺติ เคหสิตานญฺเจว สีลานํ อภินิมฺมทนาย เคหสิตานญฺเจว สรสงฺกปฺปานํ อภินิมฺมทนาย เคหสิตานญฺเจว ทรถกิลมถปริฬาหานํ อภินิมฺมทนาย ญายสฺส อธิคมาย นิพฺพานสฺส สจฺฉิกิริยาย.}

\proseitem{220}{ตเมนํ ตถาคโต อุตฺตริํ วิเนติ “เอหิ ตฺวํ ภิกฺขุ กาเย กายานุปสฺสี วิหราหิ, มา จ กามูปสํหิตํ วิตกฺกํ วิตกฺเกสิ.\csromanspacer{} เวทนาสุ.\csromanspacer{} จิตฺเต.\csromanspacer{} ธมฺเมสุ ธมฺมานุปสฺสี วิหราหิ, มา จ กามูปสํหิตํ วิตกฺกํ วิตกฺเกสี”ติ.}

\prose{โส วิตกฺกวิจารานํ วูปสมา อชฺฌตฺตํ สมฺปสาทนํ เจตโส เอโกทิภาวํ อวิตกฺกํ อวิจารํ สมาธิชํ ปีติสุขํ ทุติยํ ฌานํ.\csromanspacer{} ตติยํ ฌานํ.\csromanspacer{} จตุตฺถํ ฌานํ อุปสมฺปชฺช วิหรติ.\csromanspacer{} โส เอวํ สมาหิเต จิตฺเต ปริสุทฺเธ ปริโยทาเต อนงฺคเณ วิคตูปกฺกิเลเส มุทุภูเต กมฺมนิเย ฐิเต อาเนญฺชปฺปตฺเต ปุพฺเพนิวาสานุสฺสติญาณาย จิตฺตํ อภินินฺนาเมติ, โส อเนกวิหิตํ ปุพฺเพนิวาสํ อนุสฺสรติ.\csromanspacer{} เสยฺยถิทํ, เอกมฺปิ ชาติํ เทฺวปิ ชาติโย ฯเปฯ อิติ สาการํ ส-อุทฺเทสํ อเนกวิหิตํ ปุพฺเพนิวาสํ อนุสฺสรติ.}

\proseitem{221}{\csromanfolio{176}โส เอวํ สมาหิเต จิตฺเต ปริสุทฺเธ ปริโยทาเต อนงฺคเณ วิคตูปกฺกิเลเส มุทุภูเต กมฺมนิเย ฐิเต อาเนญฺชปฺปตฺเต สตฺตานํ จุตูปปาตญาณาย จิตฺตํ อภินินฺนาเมติ, โส ทิพฺเพน จกฺขุนา วิสุทฺเธน อติกฺกนฺตมานุสเกน สตฺเต ปสฺสติ จวมาเน อุปปชฺชมาเน หีเน ปณีเต สุวณฺเณ ทุพฺพณฺเณ สุคเต ทุคฺคเต ฯเปฯ ยถากมฺมูปเค สตฺเต ปชานาติ.}

\prose{โส เอวํ สมาหิเต จิตฺเต ปริสุทฺเธ ปริโยทาเต อนงฺคเณ วิคตูปกฺกิเลเส มุทุภูเต กมฺมนิเย ฐิเต อาเนญฺชปฺปตฺเต อาสวานํ ขยญาณาย จิตฺตํ อภินินฺนาเมติ, โส อิทํ ทุกฺขนฺติ ยถาภูตํ ปชานาติ, อยํ ทุกฺขสมุทโยติ ยถาภูตํ ปชานาติ, อยํ ทุกฺขนิโรโธติ ยถาภูตํ ปชานาติ, อยํ ทุกฺขนิโรธคามินี ปฏิปทาติ ยถาภูตํ ปชานาติ, อิเม อาสวาติ ยถาภูตํ ปชานาติ, อยํ อาสวสมุทโยติ ยถาภูตํ ปชานาติ, อยํ อาสวนิโรโธติ ยถาภูตํ ปชานาติ, อยํ อาสวนิโรธคามินี ปฏิปทาติ ยถาภูตํ ปชานาติ.\csromanspacer{} ตสฺส เอวํ ชานโต เอวํ ปสฺสโต กามาสวาปิ จิตฺตํ วิมุจฺจติ, ภวาสวาปิ จิตฺตํ วิมุจฺจติ, อวิชฺชาสวาปิ จิตฺตํ วิมุจฺจติ, วิมุตฺตสฺมิํ “วิมุตฺตมฺ”อิติ ญาณํ โหติ “ขีณา ชาติ, วุสิตํ พฺรหฺมจริยํ, กตํ กรณียํ, นาปรํ อิตฺถตฺตายา”ติ ปชานาติ.}

\prose{โส โหติ ภิกฺขุ ขโม สีตสฺส อุณฺหสฺส ชิฆจฺฉาย ปิปาสาย ฑํสมกสวาตาตปสรีสปสมฺผสฺสานํ ทุรุตฺตานํ ทุราคตานํ วจนปถานํ อุปฺปนฺนานํ สารีริกานํ เวทนานํ ทุกฺขานํ ติพฺพานํ ขรานํ กฏุกานํ อสาตานํ อมนาปานํ ปาณหรานํ อธิวาสกชาติโก โหติ สพฺพราคโทสโมหนิหิตนินฺนีตกสาโว อาหุเนยฺโย ปาหุเนยฺโย ทกฺขิเณยฺโย อญฺชลิกรณีโย อนุตฺตรํ ปุญฺญกฺเขตฺตํ โลกสฺส.}

\proseitem{222}{มหลฺลโก เจปิ อคฺคิเวสฺสน รญฺโญ นาโค อทนฺโต อวินีโต กาลํ กโรติ, อทนฺตมรณํ\footnote{อทนฺตํ มรณํ (ก)} มหลฺลโก รญฺโญ นาโค กาลงฺกโตเตฺวว สงฺขํ คจฺฉติ.\csromanspacer{} มชฺฌิโม เจปิ อคฺคิเวสฺสน รญฺโญ นาโค.\csromanspacer{} ทหโร เจปิ อคฺคิเวสฺสน รญฺโญ นาโค อทนฺโต อวินีโต}

\csromanheader{2}{6. ภูมิชสุตฺต (126)}
\prose{\csromanfolio{177}กาลํ กโรติ, อทนฺตมรณํ ทหโร รญฺโญ นาโค กาลงฺกโตเตฺวว สงฺขํ คจฺฉติ.\csromanspacer{} เอวเมว โข อคฺคิเวสฺสน เถโร เจปิ ภิกฺขุ อขีณาสโว กาลํ กโรติ, อทนฺตมรณํ เถโร ภิกฺขุ กาลงฺกโตเตฺวว สงฺขํ คจฺฉติ.\csromanspacer{} มชฺฌิโม เจปิ อคฺคิเวสฺสน ภิกฺขุ.\csromanspacer{} นโว เจปิ อคฺคิเวสฺสน ภิกฺขุ อขีณาสโว กาลํ กโรติ, อทนฺตมรณํ นโว ภิกฺขุ กาลงฺกโตเตฺวว สงฺขํ คจฺฉติ.}

\prose{มหลฺลโก เจปิ อคฺคิเวสฺสน รญฺโญ นาโค สุทนฺโต สุวินีโต กาลํ กโรติ, ทนฺตมรณํ มหลฺลโก รญฺโญ นาโค กาลงฺกโตเตฺวว สงฺขํ คจฺฉติ.\csromanspacer{} มชฺฌิโม เจปิ อคฺคิเวสฺสน รญฺโญ นาโค.\csromanspacer{} ทหโร เจปิ อคฺคิเวสฺสน รญฺโญ นาโค สุทนฺโต สุวินีโต กาลํ กโรติ, ทนฺตมรณํ ทหโร รญฺโญ นาโค กาลงฺกโตเตฺวว สงฺขํ คจฺฉติ.\csromanspacer{} เอวเมว โข อคฺคิเวสฺสน เถโร เจปิ ภิกฺขุ ขีณาสโว กาลํ กโรติ, ทนฺตมรณํ เถโร ภิกฺขุ กาลงฺกโตเตฺวว สงฺขํ คจฺฉติ.\csromanspacer{} มชฺฌิโม เจปิ อคฺคิเวสฺสน ภิกฺขุ.\csromanspacer{} นโว เจปิ อคฺคิเวสฺสน ภิกฺขุ ขีณาสโว กาลํ กโรติ, ทนฺตมรณํ นโว ภิกฺขุ กาลงฺกโตเตฺวว สงฺขํ คจฺฉตีติ.}

\prose{อิทมโวจ ภควา.\csromanspacer{} อตฺตมโน อจิรวโต สมณุทฺเทโส ภควโต ภาสิตํ อภินนฺทีติ.}

\nitthitamruled{ทนฺตภูมิสุตฺตํ นิฏฺฐิตํ ปญฺจมํ.}
\csromanmatikaanchor{mtk.31}
\csromantocmarkref{subsection}{6. ภูมิชสุตฺต}{mtk.31}
\bookmark[dest={mtk.31},level=2]{6. ภูมิชสุตฺต}
\csromanheader{2}{6. ภูมิชสุตฺต}
\proseitem{223}{เอวํ เม สุตํ\csromandash{}เอกํ สมยํ ภควา ราชคเห วิหรติ เวฬุวเน กลนฺทกนิวาเป.\csromanspacer{} อถ โข อายสฺมา ภูมิโช ปุพฺพณฺหสมยํ นิวาเสตฺวา ปตฺตจีวรมาทาย เยน ชยเสนสฺส ราชกุมารสฺส นิเวสนํ เตนุปสงฺกมิ, อุปสงฺกมิตฺวา ปญฺญตฺเต อาสเน นิสีทิ.\csromanspacer{} อถ โข ชยเสโน ราชกุมาโร เยนายสฺมา ภูมิโช เตนุปสงฺกมิ, อุปสงฺกมิตฺวา อายสฺมตา ภูมิเชน สทฺธิํ สมฺโมทิ, สมฺโมทนียํ กถํ สารณียํ วีติสาเรตฺวา เอกมนฺตํ นิสีทิ, เอกมนฺตํ นิสินฺโน โข ชยเสโน ราชกุมาโร อายสฺมนฺตํ ภูมิชํ เอตทโวจ “สนฺติ โภ ภูมิช เอเก \csromanfolio{178}สมณพฺราหฺมณา เอวํวาทิโน เอวํทิฏฺฐิโน ‘อาสญฺเจปิ กริตฺวา พฺรหฺมจริยํ จรนฺติ, อภพฺพา\footnote{จรติ, อภพฺโพ (สี, อิ) เอวมุปริปิ เอกวจเนเนว ทิสฺสติ.} ผลสฺส อธิคมาย.\csromanspacer{} อนาสญฺเจปิ กริตฺวา พฺรหฺมจริยํ จรินฺติ, อภพฺพา ผลสฺส อธิคมาย.\csromanspacer{} อาสญฺจ อนาสญฺเจปิ\footnote{อาสญฺจ อนาสญฺจ เจปิ (ฏฺฐ)} กริตฺวา พฺรหฺมจริยํ จรนฺติ, อภพฺพา ผลสฺส อธิคมาย.\csromanspacer{} เนวาสํ นานาสญฺเจปิ กริตฺวา พฺรหฺมจริยํ จรนฺติ, อภพฺพา ผลสฺส อธิคมายา’ติ.\csromanspacer{} อิธ โภโต ภูมิชสฺส สตฺถา กิํ วาที\footnote{กิํ วาที กิํทิฏฺฐี (สฺยา, กํ, ก)} กิมกฺขายี”ติ.\csromanspacer{} น โข เมตํ ราชกุมาร ภควโต สมฺมุขา สุตํ, สมฺมุขา ปฏิคฺคหิตํ.\csromanspacer{} ฐานญฺจ โข เอตํ วิชฺชติ, ยํ ภควา เอวํ พฺยากเรยฺย “อาสญฺเจปิ กริตฺวา อโยนิโส พฺรหฺมจริยํ จรนฺติ, อภพฺพา ผลสฺส อธิคมาย.\csromanspacer{} อนาสญฺเจปิ กริตฺวา อโยนิโส พฺรหฺมจริยํ จรนฺติ, อภพฺพา ผลสฺส อธิคมาย.\csromanspacer{} อาสญฺจ อนาสญฺเจปิ กริตฺวา อโยนิโส พฺรหฺมจริยํ จรนฺติ, อภพฺพา ผลสฺส อธิคมาย.\csromanspacer{} เนวาสํ นานาสญฺเจปิ กริตฺวา อโยนิโส พฺรหฺมจริยํ จรนฺติ, อภพฺพา ผลสฺส อธิคมาย.\csromanspacer{} อาสญฺเจปิ กริตฺวา โยนิโส พฺรหฺมจริยํ จรินฺติ, ภพฺพา ผลสฺส อธิคมาย.\csromanspacer{} อนาสญฺเจปิ กริตฺวา โยนิโส พฺรหฺมจริยํ จรนฺติ, ภพฺพา ผลสฺส อธิคมาย.\csromanspacer{} อาสญฺจ อนาสญฺเจปิ กริตฺวา โยนิโส พฺรหฺมจริยํ จรินฺติ, ภพฺพา ผลสฺส อธิคมาย.\csromanspacer{} เนวาสํ นานาสญฺเจปิ กริตฺวา โยนิโส พฺรหฺมจริยํ จรนฺติ, ภพฺพา ผลสฺส อธิคมายา”ติ.\csromanspacer{} น โข เม ตํ ราชกุมาร ภควโต สมฺมุขา สุตํ, สมฺมุขา ปฏิคฺคหิตํ.\csromanspacer{} ฐานญฺจ โข เอตํ วิชฺชติ, ยํ ภควา เอวํ พฺยากเรยฺยาติ.\csromanspacer{} สเจ โข โภโต ภูมิชสฺส สตฺถา เอวํวาที\footnote{เอวํวาที เอวํทิฏฺฐี (สฺยา, กํ, ก)} เอวมกฺขายี, อทฺธา โภโต ภูมิชสฺส สตฺถา สพฺเพสํเยว ปุถุสมณพฺราหฺมณานํ มุทฺธานํ\footnote{พุทฺธานํ (ก) มุทฺธานนฺติมุทฺธํ, มตฺถกนฺติ อตฺโถ.} มญฺเญ อาหจฺจ ติฏฺฐตีติ.\csromanspacer{} อถ โข ชยเสโน ราชกุมาโร อายสฺมนฺตํ ภูมิชํ สเกเนว ถาลิปาเกน ปริวิสิ.}

\proseitem{224}{อถ โข อายสฺมา ภูมิโช ปจฺฉาภตฺตํ ปิณฺฑปาตปฏิกฺกนฺโต เยน ภควา เตนุปสงฺกมิ, อุปสงฺกมิตฺวา ภควนฺตํ อภิวาเทตฺวา เอกมนฺตํ นิสีทิ, เอกมนฺตํ นิสินฺโน โข อายสฺมา ภูมิโช ภควนฺตํ เอตทโวจ\csromandash{}อิธาหํ ภนฺเต ปุพฺพณฺหสมยํ นิวาเสตฺวา ปตฺตจีวรมาทาย เยน ชยเสนสฺส ราชกุมารสฺส นิเวสนํ เตนุปสงฺกมิํ, อุปสงฺกมิตฺวา}

\csromanheader{2}{6. ภูมิชสุตฺต (126)}
\prose{\csromanfolio{179}ปญฺญตฺเต อาสเน นิสีทิํ.\csromanspacer{} อถ โข ภนฺเต ชยเสโน ราชกุมาโร เยนาหํ เตนุปสงฺกมิ, อุปสงฺกมิตฺวา มยา สทฺธิํ สมฺโมทิ, สมฺโมทนียํ กถํ สารณียํ วิติสาเรตฺวา เอกมนฺตํ นิสีทิ, เอกมนฺตํ นิสินฺโน โข ภนฺเต ชยเสโน ราชกุมาโร มํ เอตทโวจ “สนฺติ โภ ภูมิช เอเก สมณพฺราหฺมณา เอวํวาทิโน เอวํทิฏฺฐิโน ‘อาสญฺเจปิ กริตฺวา พฺรหฺมจริยํ จรนฺติ, อภพฺพา ผลสฺส อธิคมาย.\csromanspacer{} อนาสญฺเจปิ กริตฺวา ฯเปฯ.\csromanspacer{} อาสญฺจ อนาสญฺเจปิ กริตฺวา พฺรหฺมจริยํ จรนฺติ, อภพฺพา ผลสฺส อธิคมาย.\csromanspacer{} เนวาสํ นานาสญฺเจปิ กริตฺวา พฺรหฺมจริยํ จรนฺติ, อภพฺพา ผลสฺส อธิคมายา’ติ.\csromanspacer{} อิธ โภโต ภูมิชสฺส สตฺถา กิํวาที กิมกฺขายี”ติ.\csromanspacer{} เอวํ วุตฺเต อหํ ภนฺเต ชยเสนํ ราชกุมารํ เอตทโวจํ “น โข เม ตํ ราชกุมาร ภควโต สมฺมุขา สุตํ, สมฺมุขา ปฏิคฺคหิตํ.\csromanspacer{} ฐานญฺจ โข เอตํ วิชฺชติ, ยํ ภควา เอวํ พฺยากเรยฺย ‘อาสญฺเจปิ กริตฺวา อโยนิโส พฺรหฺมจริยํ จรนฺติ, อภพฺพา ผลสฺส อธิคมาย.\csromanspacer{} อนาสญฺเจปิ กริตฺวา อโยนิโส พฺรหฺมจริยํ จรนฺติ, อภพฺพา ผลสฺส อธิคมาย.\csromanspacer{} อาสญฺจ อนาสญฺเจปิ กริตฺวา อโยนิโส พฺรหฺมจริยํ จรนฺติ, อภพฺพา ผลสฺส อธิคมาย.\csromanspacer{} เนวาสํ นานาสญฺเจปิ กริตฺวา อโยนิโส พฺรหฺมจริยํ จรนฺติ, อภพฺพา ผลสฺส อธิคมาย.\csromanspacer{} อาสญฺเจปิ กริตฺวา โยนิโส พฺรหฺมจริยํ จรนฺติ, ภพฺพา ผลสฺส อธิคมาย.\csromanspacer{} อนาสญฺเจปิ กริตฺวา ฯเปฯ.\csromanspacer{} อาสญฺจ อนาสญฺเจปิ กริตฺวา ฯเปฯ.\csromanspacer{} เนวาสํ นานาสญฺเจปิ กริตฺวา โยนิโส พฺรหฺมจริยํ จรนฺติ, ภพฺพา ผลสฺส อธิคมายา’ติ, น โข เม ตํ ราชกุมาร ภควโต สมฺมุขา สุตํ, สมฺมุขา ปฏิคฺคหิตํ.\csromanspacer{} ฐานญฺจ โข เอตํ วิชฺชติ, ยํ ภควา เอวํ พฺยากเรยฺยาติ.\csromanspacer{} สเจ โภโต ภูมิชสฺส สตฺถา เอวํวาที เอวมกฺขายี, อทฺธา โภโต ภูมิชสฺส สตฺถา สพฺเพสํเยว ปุถุสมณพฺราหฺมณานํ มุทฺธานํ มญฺเญ อาหจฺจ ติฏฺฐตี”ติ.\csromanspacer{} กจฺจาหํ ภนฺเต เอวํ ปุฏฺโฐ เอวํ พฺยากรมาโน วุตฺตวาที เจว ภควโต โหมิ, น จ ภควนฺตํ อภูเตน อพฺภาจิกฺขามิ, ธมฺมสฺส จานุธมฺมํ พฺยากโรมิ, น จ โกจิ สหธมฺมิโก วาทานุวาโท คารยฺหํ ฐานํ อาคจฺฉตีติ.}

\prose{ตคฺฆ ตฺวํ ภูมิช เอวํ ปุฏฺโฐ เอวํ พฺยากรมาโน วุตฺตวาที เจว เม โหสิ, น จ มํ อภูเตน อพฺภาจิกฺขสิ, ธมฺมสฺส จานุธมฺมํ พฺยากโรสิ, น จ โกจิ สหธมฺมิโก วาทานุวาโท คารยฺหํ ฐานํ อาคจฺฉติ.\csromanspacer{} เย หิ เกจิ ภูมิช สมณา วา พฺราหฺมณา วา มิจฺฉาทิฏฺฐิโน มิจฺฉาสงฺกปฺปา มิจฺฉาวาจา}

\prose{\csromanfolio{180}มิจฺฉากมฺมนฺตา มิจฺฉา-อาชีวา มิจฺฉาวายามา มิจฺฉาสตี มิจฺฉาสมาธิโน, เต อาสญฺเจปิ กริตฺวา พฺรหฺมจริยํ จรนฺติ, อภพฺพา ผลสฺส อธิคมาย.\csromanspacer{} อนาสญฺเจปิ กริตฺวา พฺรหฺมจริยํ จรนฺติ, อภพฺพา ผลสฺส อธิคมาย.\csromanspacer{} อาสญฺจ อนาสญฺเจปิ กริตฺวา พฺรหฺมจริยํ จรนฺติ, อภพฺพา ผลสฺส อธิคมาย.\csromanspacer{} เนวาสํ นานาสญฺเจปิ กริตฺวา พฺรหฺมจริยํ จรนฺติ, อภพฺพา ผลสฺส อธิคมาย.\csromanspacer{} ตํ กิสฺส เหตุ, อโยนิ เหสา ภูมิช ผลสฺส อธิคมาย.}

\proseitem{225}{เสยฺยถาปิ ภูมิช ปุริโส เตลตฺถิโก เตลคเวสี เตลปริเยสนํ จรมาโน วาลิกํ โทณิยา อากิริตฺวา อุทเกน ปริปฺโผสกํ ปริปฺโผสกํ ปีเฬยฺย, อาสญฺเจปิ กริตฺวา วาลิกํ โทณิยา อากิริตฺวา อุทเกน ปริปฺโผสกํ ปปฺโผสกํ ปีเฬยฺย, อภพฺโพ เตลสฺส อธิคมาย.\csromanspacer{} อนาสญฺเจปิ กริตฺวา วาลิกํ โทณิยา อากิริตฺวา อุทเกน ปริปฺโผสกํ ปริปฺโผสกํ ปีเฬยฺย, อภพฺโพ เตลสฺส อธิคมาย.\csromanspacer{} อาสญฺจ อนาสญฺเจโป กริตฺวา วาลิกํ โทณิยา อากิริตฺวา อุทเกน ปริปฺโผสกํ ปริปฺโผสกํ ปีเฬยฺย, อภพฺโพ เตลสฺส อธิคมาย.\csromanspacer{} เนวาสํ นานาสญฺเจปิ กริตฺวา วาลิกํ โทณิยา อากิริตฺวา อุทเกน ปริปฺโผสกํ ปริปฺโผสกํ ปีเฬยฺย, อภพฺโพ เตลสฺส อธิคมาย.\csromanspacer{} ตํ กิสฺส เหตุ, อโยนิ เหสา ภูมิช เตลสฺส อธิคมาย.\csromanspacer{} เอวเมว โข ภูมิช เย หิ เกจิ สมณา วา พฺราหฺมณา วา มิจฺฉาทิฏฺฐิโน มิจฺฉาสงฺกปฺปา มิจฺฉาวาจา มิจฺฉากมฺมนฺตา มิจฺฉา-อาชีวา มิจฺฉาวายามา มิจฺฉาสตี มิจฺฉาสมาธิโน, เต อาสญฺเจปิ กริตฺวา พฺรหฺมจริยํ จรนฺติ, อภพฺพา ผลสฺส อธิคมาย.\csromanspacer{} อนาสญฺเจปิ กริตฺวา พฺรหฺมจริยํ จรนฺติ, อภพฺพา ผลสฺส อธิคมาย.\csromanspacer{} อาสญฺจ อนาสญฺเจปิ กริตฺวา พฺรหฺมจริยํ จรนฺติ, อภพฺพา ผลสฺส อธิคมาย.\csromanspacer{} เนวาสํ นานาสญฺเจปิ กริตฺวา พฺรหฺมจริยํ จรนฺติ, อภพฺพา ผลสฺส อธิคมาย.\csromanspacer{} ตํ กิสฺส เหตุ, อโยนิ เหสา ภูมิช ผลสฺส อธิคมาย.}

\prose{เสยฺยถาปิ ภูมิช ปุริโส ขีรตฺถิโก ขีรคเวสี ขีรปริเยสนํ จรมาโน คาวิํ ตรุณวจฺฉํ วิสาณโต อาวิญฺเฉยฺย\footnote{อาวิญฺเชยฺย (สี, สฺยา, กํ, อิ)}, อาสญฺเจปิ กริตฺวา คาวิํ ตรุณวจฺฉํ วิสาณโต อาวิญฺฉยฺย, อภพฺโพ ขีรสฺส อธิคมาย.\csromanspacer{} อนาสญฺเจปิ กริตฺวา ฯเปฯ.\csromanspacer{} อาสญฺจ อนาสญฺเจปิ กริตฺวา ฯเปฯ.}

\csromanheader{2}{6. ภูมิชสุตฺต (126)}
\prose{\csromanfolio{181}เนวาสํ นานาสญฺเจปิ กริตฺวา คาวิํ ตรุณวจฺฉํ วิสาณโต อาวิญฺเฉยฺย, อภพฺโพ ขีรสฺส อธิคมาย.\csromanspacer{} ตํ กิสฺส เหตุ, อโยนิ เหสา ภูมิช ขีรสฺส อธิคมาย.\csromanspacer{} เอวเมว โข ภูมิช เย หิ เกจิ สมณา วา พฺราหฺมณา วา มิจฺฉาทิฏฺฐิโน ฯเปฯ มิจฺฉาสมาธิโน.\csromanspacer{} เต อาสญฺเจปิ กริตฺวา พฺรหฺมจริยํ จรนฺติ, อภพฺพา ผลสฺส อธิคมาย.\csromanspacer{} อนาสญฺเจปิ กริตฺวา ฯเปฯ.\csromanspacer{} อาสญฺจ อนาสญฺเจปิ กริตฺวา ฯเปฯ.\csromanspacer{} เนวาสํ นานาสญฺเจปิ กริตฺวา พฺรหฺมจริยํ จรนฺติ, อภพฺพา ผลสฺส อธิคมาย.\csromanspacer{} ตํ กิสฺส เหตุ, อโยนิ เหสา ภูมิช ผลสฺส อธิคมาย.}

\proseitem{226}{เสยฺยถาปิ ภูมิช ปุริโส นวนีตตฺถิโก นวนีตคเวสี นวนีตปริเยสนํ จรมาโน อุทกํ กลเส อาสิญฺจิตฺวา มตฺเถน\footnote{มนฺเถน (สี), มตฺเตน (ก)} อาวิญฺเฉยฺย, อาสญฺเจปิ กริตฺวา อุทกํ กลเส อาสิญฺจิตฺวา มตฺเถน อาวิญฺเฉยฺย, อภพฺโพ นวนีตสฺส อธิคมาย.\csromanspacer{} อานาสญฺเจปิ กริตฺวา ฯเปฯ.\csromanspacer{} อาสญฺจ อนาสญฺเจปิ กริตฺวา ฯเปฯ.\csromanspacer{} เนวาสํ นานาสญฺเจปิ กริตฺวา อุทกํ กลเส อาสิญฺจิตฺวา มตฺเถน อาวิญฺเฉยฺย, อภพฺโพ นวนีตสฺส อธิคมาย.\csromanspacer{} ตํ กิสฺส เหตุ, อโยนิ เหสา ภูมิช นวนีตสฺส อธิคมาย.\csromanspacer{} เอวเมว โข ภูมิช เย หิ เกจิ สมณา วา พฺราหฺมณา วา มิจฺฉาทิฏฺฐิโน ฯเปฯ มิจฺฉาสมาธิโน.\csromanspacer{} เต อาสญฺเจปิ กริตฺวา พฺรหฺมจริยํ จรนฺติ, อภพฺพา ผลสฺส อธิคมาย.\csromanspacer{} อนาสญฺเจปิ กริตฺวา ฯเปฯ.\csromanspacer{} อาสญฺจ อนาสญฺเจปิ กริตฺวา ฯเปฯ.\csromanspacer{} เนวาสํ นานาสญฺเจปิ กริตฺวา พฺรหฺมจริยํ จรนฺติ, อภพฺพา ผลสฺส อธิคมาย.\csromanspacer{} ตํ กิสฺส เหตุ, อโยนิ เหสา ภูมิช ผลสฺส อธิคมาย.}

\prose{เสยฺยถาปิ ภูมิช ปุริโส อคฺคิตฺถิโก\footnote{อคฺคตฺถิโก (สี)} อคฺคิคเวสี อคฺคิปริเยสนํ จรมาโน อลฺลํ กฏฺฐํ สสฺเนหํ อุตฺตรารณิํ อาทาย อภิมนฺเถยฺย\footnote{อภิมตฺเถยฺย (สฺยา, กํ, อิ, ก)}, อาสญฺเจปิ กริตฺวา อลฺลํ กฏฺฐํ สสฺเนหํ อุตฺตรารณิํ อาทาย อภิมนฺเถยฺย, อภพฺโพ อคฺคิสฺส อธิคมาย.\csromanspacer{} อนาสญฺเจปิ กริตฺวา ฯเปฯ.\csromanspacer{} อาสญฺจ อนาสญฺเจปิ กริตฺวา ฯเปฯ.\csromanspacer{} เนวาสํ นานาสญฺเจปิ กริตฺวา อลฺลํ กฏฺฐํ สสฺเนหํ อุตฺตรารณิํ อาทาย อภิมนฺเถยฺย, อภพฺโพ อคฺคิสฺส อธิคมาย.\csromanspacer{} ตํ กิสฺส เหตุ, อโยนิ เหสา ภูมิช อคฺคิสฺส อธิคมาย.\csromanspacer{} เอวเมว โข ภูมิช เย หิ เกจิ สมณา วา พฺราหฺมณา \csromanfolio{182}วา มิจฺฉาทิฏฺฐิโน ฯเปฯ มิจฺฉาสมาธิโน.\csromanspacer{} เต อาสญฺเจปิ กริตฺวา พฺรหฺมจริยํ จรนฺติ, อภพฺพา ผลสฺส อธิคมาย.\csromanspacer{} อนาสญฺเจปิ กริตฺวา ฯเปฯ.\csromanspacer{} อาสญฺจ อนาสญฺเจปิ กริตฺวา ฯเปฯ.\csromanspacer{} เนวาสํ นานาสญฺเจปิ กริตฺวา พฺรหฺมจริยํ จรนฺติ, อภพฺพา ผลสฺส อธิคมาย.\csromanspacer{} ตํ กิสฺส เหตุ, อโยนิ เหสา ภูมิช ผลสฺส อธิคมาย.\csromanspacer{} เย หิ เกจิ ภูมิช สมณา วา พฺราหฺมณา วา สมฺมาทิฏฺฐิโน สมฺมาสงฺกปฺปา สมฺมาวาจา สมฺมากมฺมนฺตา สมฺมา-อาชีวา สมฺมาวายามา สมฺมาสตี สมฺมาสมาธิโน, เต อาสญฺเจปิ กริตฺวา พฺรหฺมจริยํ จรนฺติ, ภพฺพา ผลสฺส อธิคมาย.\csromanspacer{} อนาสญฺเจปิ กริตฺวา พฺรหฺมจริยํ จรนฺติ, ภพฺพา ผลสฺส อธิคมาย.\csromanspacer{} อาสญฺจ อนาสญฺเจปิ กริตฺวา พฺรหฺมจริยํ จรนฺติ, ภพฺพา ผลสฺส อธิคมาย.\csromanspacer{} เนวาสํ นานาสญฺเจปิ กริตฺวา พฺรหฺมจริยํ จรนฺติ, ภพฺพา ผลสฺส อธิคมาย.\csromanspacer{} ตํ กิสฺส เหตุ, โยนิ เหสา ภูมิช ผลสฺส อธิคมาย.}

\proseitem{227}{เสยฺยถาปิ ภูมิช ปุริโส เตลตฺถิโก เตลคเวสี เตลปริเยสนํ จรมาโน ติลปิฏฺฐํ โทณิยา อากิริตฺวา อุทเกน ปริปฺโผสกํ ปริปฺโผสกํ ปีเฬยฺย, อาสญฺเจปิ กริตฺวา ติลปิฏฺฐํ โทณิยา อากิริตฺวา อุทเกน ปริปฺโผสกํ ปริปฺโผสกํ ปีเฬยฺย, ภพฺโพ เตลสฺส อธิคมาย.\csromanspacer{} อนาสญฺเจปิ กริตฺวา ฯเปฯ.\csromanspacer{} อาสญฺจ อนาสญฺเจปิ กริตฺวา ฯเปฯ.\csromanspacer{} เนวาสํ นานาสญฺเจปิ กริตฺวา ติลปิฏฺฐํ โทณิยา อากิริตฺวา อุทเกน ปริปฺโผสกํ ปริปฺโผสกํ ปีเฬยฺย, ภพฺโพ เตลสฺส อธิคมาย.\csromanspacer{} ตํ กิสฺส เหตุ, โยนิ เหสา ภูมิช เตลสฺส อธิคมาย.\csromanspacer{} เอวเมว โข ภูมิช เย หิ เกจิ สมณา วา พฺราหฺมณา วา สมฺมาทิฏฺฐิโน ฯเปฯ สมฺมาสมาธิโน, เต อาสญฺเจปิ กริตฺวา พฺรหฺมจริยํ จรนฺติ, ภพฺพา ผลสฺส อธิคมาย.\csromanspacer{} อนาสญฺเจปิ กริตฺวา ฯเปฯ.\csromanspacer{} อาสญฺจ อนาสญฺเจปิ กริตฺวา ฯเปฯ.\csromanspacer{} เนวาสํ นานาสญฺเจปิ กริตฺวา พฺรหฺมจริยํ จรนฺติ, ภพฺพา ผลสฺส อธิคมาย.\csromanspacer{} ตํ กิสฺส เหตุ, โยนิ เหสา ภูมิช ผลสฺส อธิคมาย.}

\prose{เสยฺยถาปิ ภูมิช ปุริโส ขีรตฺถิโก ขีรคเวสี ขีรปริเยสนํ จรมาโน คาวิํ ตรุณวจฺฉํ ถนโต อาวิญฺเฉยฺย, อาสญฺเจปิ กริตฺวา คาวิํ ตรุณวจฺฉํ ถนโต อาวิญฺเฉยฺย, ภพฺโพ ขีรสฺส อธิคมาย.\csromanspacer{} อนาสญฺเจปิ กริตฺวา ฯเปฯ.\csromanspacer{} อาสญฺจ อนาสญฺเจปิ กริตฺวา ฯเปฯ.\csromanspacer{} เนวาสํ นานาสญฺเจปิ กริตฺวา คาวิํ ตรุณวจฺฉํ ถนโต อาวิญฺเฉยฺย, ภพฺโพ ขีรสฺส อธิคมาย.\csromanspacer{} ตํ กิสฺส เหตุ, โยนิ เหสา ภูมิช ขีรสฺส อธิคมาย.}

\csromanheader{2}{6. ภูมิชสุตฺต (126)}
\prose{\csromanfolio{183}เอวเมว โข ภูมิช เย หิ เกจิ สมณา วา พฺราหฺมณา วา สมฺมาทิฏฺฐิโน ฯเปฯ สมฺมาสมาธิโน, เต อาสญฺเจปิ กริตฺวา ฯเปฯ.\csromanspacer{} อนาสญฺเจปิ กริตฺวา ฯเปฯ.\csromanspacer{} อาสญฺจ อนาสญฺเจปิ กริตฺวา ฯเปฯ.\csromanspacer{} เนวาสํ นานาสญฺเจปิ กริตฺวา พฺรหฺมจริยํ จรนฺติ, ภพฺพา ผลสฺส อธิคมาย.\csromanspacer{} ตํ กิสฺส เหตุ, โยนิ เหสา ภูมิช ผลสฺส อธิคมาย.}

\proseitem{228}{เสยฺยถาปิ ภูมิช ปุริโส นวนีตตฺถิโก นวนีตคเวสี นวนีตปริเยสนํ จรมาโน ทธิํ กลเส อาสิญฺจิตฺวา มตฺเถน อาวิญฺเฉยฺย, อาสญฺเจปิ กริตฺวา ทธิํ กลเส อาสิญฺจิตฺวา มตฺเถน อาวิญฺเฉยฺย, ภพฺโพ นวนีตสฺส อธิคมาย.\csromanspacer{} อนาสญฺเจปิ กริตฺวา.\csromanspacer{} อาสญฺจ อนาสญฺเจปิ กริตฺวา.\csromanspacer{} เนวาสํ นานาสญฺเจปิ กริตฺวา ทธิํ กลเส อาสิญฺจิตฺวา มตฺเถน อาวิญฺเฉยฺย, ภพฺโพ นวนีตสฺส อธิคมาย.\csromanspacer{} ตํ กิสฺส เหตุ, โยนิ เหสา ภูมิช นวนีตสฺส อธิคมาย.\csromanspacer{} เอวเมว โข ภูมิช เย หิ เกจิ สมณา วา พฺราหฺมณา วา สมฺมาทิฏฺฐิโน ฯเปฯ สมฺมาสมาธิโน.\csromanspacer{} เต อาสญฺเจปิ กริตฺวา พฺรหฺมจริยํ จรนฺติ, ภพฺพา ผลสฺส อธิคมาย.\csromanspacer{} อนาสญฺเจปิ กริตฺวา.\csromanspacer{} อาสญฺจ อนาสญฺเจปิ กริตฺวา.\csromanspacer{} เนวาสํ นานาสญฺเจปิ กริตฺวา พฺรหฺมจริยํ จรนฺติ, ภพฺพา ผลสฺส อธิคมาย.\csromanspacer{} ตํ กิสฺส เหตุ, โยนิ เหสา ภูมิช ผลสฺส อธิคมาย.}

\prose{เสยฺยถาปิ ภูมิช ปุริโส อคฺคิตฺถิโก อคฺคิคเวสี อคฺคิปริเยสนํ จรมาโน สุกฺขํ กฏฺฐํ โกฬาปํ อุตฺตรารณิํ อาทาย อภิมนฺเถยฺย, ( )\footnote{(ภพฺโพ อคฺคิสฺส อธิคมาย) (สพฺพตฺถ)} อาสญฺเจปิ กริตฺวา.\csromanspacer{} อนาสญฺเจปิ กริตฺวา.\csromanspacer{} อาสญฺจ อนาสญฺเจปิ กริตฺวา.\csromanspacer{} เนวาสํ นานาสญฺเจปิ กริตฺวา สุกฺขํ กฏฺฐํ โกฬาปํ อุตฺตารารณิํ อาทาย อภิมนฺเถยฺย, ภพฺโพ อคฺคิสฺส อธิคมาย.\csromanspacer{} ตํ กิสฺส เหตุ, โยนิ เหสา ภูมิช อคฺคิสฺส อธิคมาย.\csromanspacer{} เอวเมว โข ภูมิช เย หิ เกจิ สมณา วา พฺราหฺมณา วา สมฺมาทิฏฺฐิโน ฯเปฯ สมฺมาสมาธิโน.\csromanspacer{} เต อาสญฺเจปิ กริตฺวา พฺรหฺมจริยํ จรนฺติ, ภพฺพา ผลสฺส อธิคมาย.\csromanspacer{} อนาสญฺเจปิ กริตฺวา พฺรหฺมจริยํ จรนฺติ, ภพฺพา ผลสฺส อธิคมาย.\csromanspacer{} อาสญฺจ อนาสญฺเจปิ กริตฺวา พฺรหฺมจริยํ จรนฺติ, ภพฺพา ผลสฺส อธิคมาย.\csromanspacer{} เนวาสํ นานาสญฺเจปิ กริตฺวา พฺรหฺมจริยํ จรนฺติ, ภพฺพา ผลสฺส อธิคมาย.\csromanspacer{} ตํ กิสฺส เหตุ, โยนิ เหสา ภูมิช ผลสฺส อธิคมาย.}

\prose{\csromanfolio{184}สเจ โข ตํ ภูมิช ชยเสนสฺส ราชกุมารสฺส อิมา จตสฺโส อุปมา ปฏิภาเยยฺยุํ, อนจฺฉริยํ เต ชยเสโน ราชกุมาโร ปสีเทยฺย, ปสนฺโน จ เต ปสนฺนาการํ กเรยฺยาติ.\csromanspacer{} กุโต ปน มํ ภนฺเต ชยเสนสฺส ราชกุมารสฺส อิมา จตสฺโส อุปมา ปฏิภายิสฺสนฺติ, อนจฺฉริยา ปุพฺเพ อสฺสุตปุพฺพา, เสยฺยถาปิ ภควนฺตนฺติ.}

\prose{อิทมโวจ ภควา.\csromanspacer{} อตฺตมโน อายสฺมา ภูมิโช ภควโต ภาสิตํ อภินนฺทีติ.}

\nitthitamruled{ภูมิชสุตฺตํ นิฏฺฐิตํ ฉฏฺฐํ.}
\csromanmatikaanchor{mtk.32}
\csromantocmarkref{subsection}{7. อนุรุทฺธสุตฺต}{mtk.32}
\bookmark[dest={mtk.32},level=2]{7. อนุรุทฺธสุตฺต}
\csromanheader{2}{7. อนุรุทฺธสุตฺต}
\proseitem{229}{เอวํ เม สุตํ\csromandash{}เอกํ สมยํ ภควา สาวตฺถิยํ วิหรติ เชตวเน อนาถปิณฺฑิกสฺส อาราเม.\csromanspacer{} อถ โข ปญฺจกงฺโค ถปติ อญฺญตรํ ปุริสํ อามนฺเตสิ “เอหิ ตฺวํ อมฺโภ ปุริส เยนายสฺมา อนุรุทฺโธ เตนุปสงฺกม, อุปสงฺกมิตฺวา มม วจเนน อายสฺมโต อนุรุทฺธสฺส ปาเท สิรสา วนฺทาหิ\footnote{วนฺทาหิ, เอวญฺจ วเทหิ (สี, อิ)} ‘ปญฺจกงฺโค ภนฺเต ถปติ อายสฺมโต อนุรุทฺธสฺส ปาเท สิรสา วนฺทตี’ติ, เอวญฺจ วเทหิ\footnote{เอวญฺจ วเทติ (สี, อิ)} ‘อธิวาเสตุ กิร ภนฺเต อายสฺมา อนุรุทฺโธ ปญฺจกงฺคสฺส ถปติสฺส สฺวาตนาย อตฺตจตุตฺโถ ภตฺตํ, เยน จ กิร ภนฺเต อายสฺมา อนุรุทฺโธ ปเควตรํ อาคจฺเฉยฺย, ปญฺจกงฺโค ภนฺเต ถปติ\footnote{ปญฺจกงฺโค ถปหิ (สี, อิ)} พหุกิจฺโจ พหุกรณีโย ราชกรณีเยนา’ติ”. “เอวํ ภนฺเต”ติ โข โส ปุริโส ปญฺจกงฺคสฺส ถปติสฺส ปฏิสฺสุตฺวา เยนายสฺมา อนุรุทฺโธ เตนุปสงฺกมิ, อุปสงฺกมิตฺวา อายสฺมนฺตํ อนุรุทฺธํ อภิวาเทตฺวา เอกมนฺตํ นิสีทิ, เอกมนฺตํ นิสินฺโน โข โส ปุริโส อายสฺมนฺตํ อนุรุทฺธํ เอตทโวจ “ปญฺจกงฺโค ภนฺเต ถปติ อายสฺมโต อนุรุทฺธสฺส ปาเท สิรสา วนฺทติ, เอวญฺจ วเทติ ‘อธิวาเสตุ กิร ภนฺเต อายสฺมา อนุรุทฺโธ ปญฺจกงฺคสฺส ถปติสฺส สฺวาตนาย อตฺตจตุตฺโถ ภตฺตํ, เยน จ กิร ภนฺเต อายสฺมา อนุรุทฺโธ ปเควตรํ อาคจฺเฉยฺย, ปญฺจกงฺโค ภนฺเต ถปติ พหุกิจฺโจ พหุกรณีโย ราชกรณีเยนา’ติ”.\csromanspacer{} อธิวาเสสิ โข อายสฺมา อนุรุทฺโธ ตุณฺหีภาเวน.}

\csromanheader{2}{7. อนุรุทฺธสุตฺต (127)}
\proseitem{230}{\csromanfolio{185}อถ โข อายสฺมา อนุรุทฺโธ ตสฺสา รตฺติยา อจฺจเยน ปุพฺพณฺหสมยํ นิวาเสตฺวา ปตฺตจีวรมาทาย เยน ปญฺจกงฺคสฺส ถปติสฺส นิเวสนํ เตนุปสงฺกมิ, อุปสงฺกมิตฺวา ปญฺญตฺเต อาสเน นิสีทิ.\csromanspacer{} อถ โข ปญฺจกงฺโค ถปติ อายสฺมนฺตํ อนุรุทฺธํ ปณีเตน ขาทนีเยน โภชนีเยน สหตฺถา สนฺตปฺเปสิ สมฺปวาเรสิ.\csromanspacer{} อถ โข ปญฺจกงฺโค ถปติ อายสฺมนฺตํ อนุรุทฺธํ ภุตฺตาวิํ โอนีตปตฺตปาณิํ อญฺญตรํ นีจํ อาสนํ คเหตฺวา เอกมนฺตํ นิสีทิ, เอกมนฺตํ นิสินฺโน โข ปญฺจกงฺโค ถปติ อายสฺมนฺตํ อนุรุทฺธํ เอตทโวจ\csromandash{}}

\prose{อิธ มํ ภนฺเต เถรา ภิกฺขู อุปสงฺกมิตฺวา เอวมาหํสุ “อปฺปมาณํ คหปติ เจโตวิมุตฺติํ ภาเวหี”ติ\footnote{อปฺปมาณา คหปติ เจโตวิมุตฺติ ภาเวตพฺพาติ (ก)}.\csromanspacer{} เอกจฺเจ เถรา เอวมาหํสุ “มหคฺคตํ คหปติ เจโตวิมุตฺติํ ภาเวหี”ติ.\csromanspacer{} ยา จายํ ภนฺเต อปฺปมาณา เจโตวิมุตฺติ ยา จ มหคฺคตา เจโตวิมุตฺติ, อิเม ธมฺมา นานตฺถา เจว นานาพฺยญฺชนา จ, อุทาหุ เอกตฺถา พฺยญฺชนเมว นานนฺติ? เตน หิ คหปติ ตํเยเวตฺถ ปฏิภาตุ อปณฺณกนฺเต อิโต ภวิสฺสตีติ.\csromanspacer{} มยฺหํ โข ภนฺเต เอวํ โหติ “ยา จายํ อปฺปมาณา เจโตวิมุตฺติ ยา จ มหคฺคตา เจโตวิมุตฺติ, อิเม ธมฺมา เอกตฺถา พฺยญฺชเมว นานนฺ”ติ.\csromanspacer{} ยา จายํ คหปติ อปฺปมาณา เจโตวิมุตฺติ ยา จ มหคฺคตา เจโตวิมุตฺติ, อิเม ธมฺมา นานตฺถา เจว นานาพฺยญฺชนา จ, ตทมินาเปตํ คหปติ ปริยาเยน เวทิตพฺพํ, ยถา อิเม ธมฺมา นานตฺถา เจว นานาพฺยญฺชนา จ.}

\prose{กตมา จ คหปติ อปฺปมาณา เจโตวิมุตฺติ, อิธ คหปติ ภิกฺขุ เมตฺตาสหคเตน เจตสา เอกํ ทิสํ ผริตฺวา วิหรติ.\csromanspacer{} ตถา ทุติยํ.\csromanspacer{} ตถา ตติยํ.\csromanspacer{} ตถา จตุตฺถํ.\csromanspacer{} อิติ อุทฺธมโธ ติริยํ สพฺพธิ สพฺพตฺตตาย สพฺพาวนฺตํ โลกํ เมตฺตาสหคเตน เจตสา วิปุเลน มหคฺคเตน อปฺปมาเณน อเวเรน อพฺยาพชฺเฌน ผริตฺวา วิหรติ.\csromanspacer{} กรุณาสหคเตน เจตสา.\csromanspacer{} มุทิตาสหคเตน เจตสา.\csromanspacer{} อุเปกฺขาสหคเตน เจตสา เอกํ ทิสํ ผริตฺวา วิหรติ.\csromanspacer{} ตถา ทุติยํ.\csromanspacer{} ตถา ตติยํ.\csromanspacer{} ตถา จตุตฺถํ.\csromanspacer{} อิติ อุทฺเธมโธ ติริยํ สพฺพธิ สพฺพตฺตตาย สพฺพาวนฺตํ โลกํ อุเปกฺขาสหคเตน เจตสา วิปุเลน มหคฺคเตน อปฺปมาเณน อเวเรน อพฺยาพชฺเฌน ผริตฺวา วิหรติ, อยํ วุจฺจติ คหปติ อปฺปมาณา เจโตวิมุตฺติ.}

\proseitem{231}{\csromanfolio{186}กตมา จ คหปติ มหคฺคตา เจโตวิมุตฺติ, อิธ คหปติ ภิกฺขุ ยาวตา เอกํ รุกฺขมูลํ “มหคฺคตนฺ”ติ ผริตฺวา อธิมุจฺจิตฺวา วิหรติ, อยํ วุจฺจติ คหปติ มหคฺคตา เจโตวิมุตฺติ.\csromanspacer{} อิธ ปน คหปติ ภิกฺขุ ยาวตา เทฺว วา ตีณิ วา รุกฺขมูลานิ “มหคฺคตนฺ”ติ ผริตฺวา อธิมุจฺจิตฺวา วิหรติ, อยมฺปิ\footnote{อยํ (สฺยา, กํ, ก)} วุจฺจติ คหปติ มหคฺคตา เจโตวิมุตฺติ.\csromanspacer{} อิธ ปน คหปติ ภิกฺขุ ยาวตา เอกํ คามกฺเขตฺตํ “มหคฺคตนฺ”ติ ผริตฺวา อธิมุจฺจิตฺวา วิหรติ, อยมฺปิ วุจฺจติ มหคฺคตา เจโตวิมุตฺติ.\csromanspacer{} อิธ ปน คหปติ ภิกฺขุ ยาวตา เทฺว วา ตีณิ วา คามกฺเขตฺตานิ “มหคฺคตนฺ”ติ ผริตฺวา อธิมุจฺจิตฺวา วิหรติ, อยมฺปิ วุจฺจติ คหปติ มหคฺคตา เจโตวิมุตฺติ.\csromanspacer{} อิธ ปน คหปติ ภิกฺขุ ยาวตา เอกํ มหารชฺชํ “มหคฺคตนฺ”ติ ผริตฺวา อธิมุจฺจิตฺวา วิหรติ, อยมฺปิ วุจฺจติ คหปติ มหคฺคตา เจโตวิมุตฺติ.\csromanspacer{} อิธ ปน คหปติ ภิกฺขุ ยาวตา เทฺว วา ตีณิ วา มหารชฺชานิ “มหคฺคตนฺ”ติ ผริตฺวา อธิมุจฺจิตฺวา วิหรติ, อยมฺปิ วุจฺจติ คหปติ มหคฺคตา เจโตวิมุตฺติ.\csromanspacer{} อิธ ปน คหปติ ภิกฺขุ ยาวตา สมุทฺทปริยนฺตํ ปถวิํ “มหคฺคตนฺ”ติ ผริตฺวา อธิมุจฺจิตฺวา วิหรติ, อยมฺปิ วุจฺจติ คหปติ มหคฺคตา เจโตวิมุตฺติ.\csromanspacer{} อิมินา โข เอตํ คหปติ ปริยาเยน เวทิตพฺพํ, ยถา อิเม ธมฺมา นานตฺถา เจว นานาพฺยญฺชนา จ.}

\proseitem{232}{จตสฺโส โข อิมา คหปติ ภวูปปตฺติโย.\csromanspacer{} กตมา จตสฺโส, อิธ คหปติ เอกจฺโจ “ปริตฺตาภา”ติ ผริตฺวา อธิมุจฺจิตฺวา วิหรติ, โส กายสฺส เภทา ปรํ มรณา ปริตฺตาภานํ เทวานํ สหพฺยตํ อุปปชฺชติ.\csromanspacer{} อิธ ปน คหปติ เอกจฺโจ “อปฺปมาณาภา”ติ ผริตฺวา อธิมุจฺจิตฺวา วิหรติ, โส กายสฺส เภทา ปรํ มรณา อปฺปมาณาภานํ เทวานํ สหพฺยตํ อุปปชฺชติ.\csromanspacer{} อิธ ปน คหปติ เอกจฺโจ “สํกิลิฏฺฐาภา”ติ ผริตฺวา อธิมุจฺจิตฺวา วิหรติ, โส กายสฺส เภทา ปรํ มรณา สํกิลิฏฺฐาภานํ เทวานํ สหพฺยตํ อุปปชฺชติ.\csromanspacer{} อิธ ปน คหปติ เอกจฺโจ “ปริสุทฺธาภา”ติ ผริตฺวา วิหรติ, โส กายสฺส เภทา ปรํ มรณา ปริสุทฺธาภานํ เทวานํ สหพฺยตํ อุปปชฺชติ.\csromanspacer{} อิมา โข คหปติ จตสฺโส ภวูปปตฺติโย.}

\prose{โหติ โข โส คหปติ สมโย, ยา ตา เทวตา เอกชฺฌํ สนฺนิปตนฺติ, ตาสํ เอกชฺฌํ สนฺนิปติตานํ วณฺณนานตฺตํ หิ โข ปญฺญายติ,}

\csromanheader{2}{7. อนุรุทฺธสุตฺต (127)}
\prose{\csromanfolio{187}โน จ อาภานานตฺตํ.\csromanspacer{} เสยฺยถาปิ คหปติ ปุริโส สมฺพหุลานิ เตลปฺปทีปานิ เอกํ ฆรํ ปเวเสยฺย, เตสํ เอกํ ฆรํ ปเวสิตานํ อจฺจินานตฺตํ หิ โข ปญฺญาเยถ, โน จ อาภานานตฺตํ.\csromanspacer{} เอวเมว โข คหปติ โหติ โข โส สมโย, ยา ตา เทวตา เอกชฺฌํ สนฺนิปตนฺติ, ตาสํ เอกชฺฌํ สนฺนิปติตานํ วณฺณนานตฺตํ หิ โข ปญฺญายติ, โน จ อาภานานตฺตํ.}

\prose{โหติ โข โส คหปติ สมโย, ยา ตา เทวตา ตโต วิปกฺกมนฺติ, ตาสํ ตโต วิปกฺกมนฺตีนํ วณฺณนานตฺตญฺเจว ปญฺญายติ อาภานานตฺตญฺจ.\csromanspacer{} เสยฺยถาปิ คหปติ ปุริโส ตานิ สมฺพหุลานิ เตลปฺปทีปานิ ตมฺหา ฆรา นีหเรยฺย, เตสํ ตโต นีหตานํ\footnote{นีหรนฺตานํ (สี, สฺยา, กํ, อิ)} อจฺจินานตฺตญฺเจว ปญฺญาเยถ อาภานานตฺตญฺจ.\csromanspacer{} เอวเมว โข คหปติ โหติ โข โส สมโย, ยา ตา เทวตา ตโต วิปกฺกมนฺติ, ตาสํ ตโต วิปกฺกมนฺตีนํ วณฺณนานตฺตญฺเจว ปญฺญายติ อาภานานตฺตญฺจ.}

\prose{น โข คหปติ ตาสํ เทวตานํ เอวํ โหติ “อิทํ อมฺหากํ นิจฺจนฺติ วา ธุวนฺติ วา สสฺสตนฺติ วา”, อปิ จ ยตฺถ ยตฺเถว ตา\footnote{ยา (ก)} เทวตา อภินิวิสนฺติ, ตตฺถ ตตฺเถว ตา เทวตา อภิรมนฺติ.\csromanspacer{} เสยฺยถาปิ คหปติ มกฺขิกานํ กาเชน วา ปิฏเกน วา หรียมานานํ น เอวํ โหติ “อิทํ อมฺหากํ นิจฺจนฺติ วา ธุวนฺติ วา สสฺสตนฺติ วา”, อปิ จ ยตฺถ ยตฺเถว ตา2 มกฺขิกา อภินิวิสนฺติ, ตตฺถ ตตฺเถว ตา มกฺขิกา อภิรมนฺติ.\csromanspacer{} เอวเมว โข คหปติ ตาสํ เทวตานํ น เอวํ โหติ “อิทํ อมฺหากํ นิจฺจนฺติ วา ธุวนฺติ วา สสฺสตนฺติ วา”, อปิ จ ยตฺถ ยตฺเถว ตา เทวตา อภินิวิสนฺติ, ตตฺถ ตตฺเถว ตา เทวตา อภิรมนฺตีติ.}

\proseitem{233}{เอวํ วุตฺเต อายสฺมา สภิโย กจฺจาโน\footnote{กจฺจายโน (สี)} อายสฺมนฺตํ อนุรุทฺธํ เอตทโวจ “สาธุ ภนฺเต อนุรุทฺธ, อตฺถิ จ เม เอตฺถ อุตฺตริํ ปฏิปุจฺฉิตพฺพํ ‘ยา ตา ภนฺเต เทวตา อาภา, สพฺพา ตา ปริตฺตาภา, อุทาหุ สนฺเตตฺถ เอกจฺจา เทวตา อปฺปมาณาภา’ติ”.\csromanspacer{} ตทงฺเคน โข อาวุโส กจฺจาน สนฺเตตฺถ เอกจฺจา เทวตา ปริตฺตาภา, สนฺติ ปเนตฺถ เอกจฺจา เทวตา อปฺปมาณาภาติ.\csromanspacer{} โก นุ โข ภนฺเต อนุรุทฺธ เหตุ โก ปจฺจโย, เยน ตาสํ เทวตานํ เอกํ เทวนิกายํ \csromanfolio{188}อุปปนฺนานํ สนฺเตตฺถ เอกจฺจา เทวตา ปริตฺตาภา, สนฺติ ปเนตฺถ เอกจฺจา เทวตา อปฺปมาณาภาติ.}

\prose{เตน หาวุโส กจฺจาน ตํเยเวตฺถ ปฏิปุจฺฉิสฺสามิ.\csromanspacer{} ยถา เต ขเมยฺย, ตถา นํ พฺยากเรยฺยาสิ.\csromanspacer{} ตํ กิํ มญฺญสิ อาวุโส กจฺจาน, ยฺวายํ ภิกฺขุ ยาวตา เอกํ รุกฺขมูลํ “มหคฺคตนฺ”ติ ผริตฺวา อธิมุจฺจิตฺวา วิหรติ, โยจายํ\footnote{โยปายํ (ก)} ภิกฺขุ ยาวตา เทฺว วา ตีณิ วา รุกฺขมูลานิ “มหคฺคตนฺ”ติ ผริตฺวา อธิมุจฺจิตฺวา วิหรติ.\csromanspacer{} อิมาสํ อุภินฺนํ จิตฺตภาวนานํ กตมา จิตฺตภาวนา มหคฺคตตราติ.\csromanspacer{} ยฺวายํ ภนฺเต ภิกฺขุ ยาวตา เทฺว วา ตีณิ วา รุกฺขมูลานิ “มหคฺคตนฺ”ติ ผริตฺวา อธิมุจฺจิตฺวา วิหรติ, อยํ อิมาสํ อุภินฺนํ จิตฺตภาวนานํ มหคฺคตตราติ.}

\prose{ตํ กิํ มญฺญสิ อาวุโส กจฺจาน, ยฺวายํ ภิกฺขุ ยาวตา เทฺว วา ตีณิ วา รุกฺขมูลานิ “มหคฺคตนฺ”ติ ผริตฺวา อธิมุจฺจิตฺวา วิหรติ, โยจายํ ภิกฺขุ ยาวตา เอกํ คามกฺเขตฺตํ “มหคฺคตนฺ”ติ ผริตฺวา อธิมุจฺจิตฺวา วิหรติ.\csromanspacer{} อิมาสํ อุภินฺนํ จิตฺตภาวนานํ กตมา จิตฺตภาวนา มหคฺคตตราติ.\csromanspacer{} ยฺวายํ ภนฺเต ภิกฺขุ ยาวตา เอกํ คามกฺเขตฺตํ “มหคฺคตนฺ”ติ ผริตฺวา อธิมุจฺจิตฺวา วิหรติ, อยํ อิมาสํ อุภินฺนํ จิตฺตภาวนานํ มหคฺคตตราติ.}

\prose{ตํ กิํ มญฺญสิ อาวุโส กจฺจาน, ยฺวายํ ภิกฺขุ ยาวตา เอกํ คามกฺเขตฺตํ “มหคฺคตนฺ”ติ ผริตฺวา อธิมุจฺจิตฺวา วิหรติ, โยจายํ ภิกฺขุ ยาวตา เทฺว วา ตีณิ วา คามกฺเขตฺตานิ “มหคฺคตนฺ”ติ ผริตฺวา อธิมุจฺจิตฺวา วิหรติ.\csromanspacer{} อิมาสํ อุภินฺนํ จิตฺตภาวนานํ กตมา จิตฺตภาวนา มหคฺคตตราติ.\csromanspacer{} ยฺวายํ ภนฺเต ภิกฺขุ ยาวตา เทฺว วา ตีณิ วา คามกฺเขตฺตานิ “มหคฺคตนฺ”ติ ผริตฺวา อธิมุจฺจิตฺวา วิหรติ, อยํ อิมาสํ อุภินฺนํ จิตฺตภาวนานํ มหคฺคตตราติ.}

\prose{ตํ กิํ มญฺญสิ อาวุโส กจฺจาน, ยฺวายํ ภิกฺขุ ยาวตา เทฺว วา ตีณิ วา คามกฺเขตฺตานิ “มหคฺคตนฺ”ติ ผริตฺวา อธิมุจฺจิตฺวา วิหรติ, โยจายํ ภิกฺขุ ยาวตา เอกํ มหารชฺชํ “มหคฺคตนฺ”ติ ผริตฺวา อธิมุจฺจิตฺวา วิหรติ.\csromanspacer{} อิมาสํ อุภินฺนํ จิตฺตภาวนานํ กตมา จิตฺตภาวนา}

\csromanheader{2}{7. อนุรุทฺธสุตฺต (127)}
\prose{\csromanfolio{189}มหคฺคตตราติ.\csromanspacer{} ยฺวายํ ภนฺเต ภิกฺขุ ยาวตา เอกํ มหารชฺชํ “มหคฺคตนฺ”ติ ผริตฺวา อธิมุจฺจิตฺวา วิหรติ, อยํ อิมาสํ อุภินฺนํ จิตฺตภาวนานํ มหคฺคตตราติ.}

\prose{ตํ กิํ มญฺญสิ อาวุโส กจฺจาน, ยฺวายํ ภิกฺขุ ยาวตา เอกํ มหารชฺชํ “มหคฺคตนฺ”ติ ผริตฺวา อธิมุจฺจิตฺวา วิหรติ, โยจายํ ภิกฺขุ ยาวตา เทฺว วา ตีณิ วา มหารชฺชานิ “มหคฺคตนฺ”ติ ผริตฺวา อธิมุจฺจิตฺวา วิหรติ.\csromanspacer{} อิมาสํ อุภินฺนํ จิตฺตภาวนานํ กตมา จิตฺตภาวนา มหคฺคตตราติ.\csromanspacer{} ยฺวายํ ภนฺเต ภิกฺขุ ยาวตา เทฺว วา ตีณิ วา มหารชฺชานิ “มหคฺคตนฺ”ติ ผริตฺวา อธิมุจฺจิตฺวา วิหรติ, อยํ อิมาสํ อุภินฺนํ จิตฺตภาวนานํ มหคฺคตตราติ.}

\prose{ตํ กิํ มญฺญสิ อาวุโส กจฺจาน, ยฺวายํ ภิกฺขุ ยาวตา เทฺว วา ตีณิ วา มหารชฺชานิ “มหคฺคตนฺ”ติ ผริตฺวา อธิมุจฺจิตฺวา วิหรติ, โยจายํ ภิกฺขุ ยาวตา สมุทฺทปริยนฺตํ ปถวิํ “มหคฺคตนฺ”ติ ผริตฺวา อธิมุจฺจิตฺวา วิหรติ.\csromanspacer{} อิมาสํ อุภินฺนํ จิตฺตภาวนานํ กตมา จิตฺตภาวนา มหคฺคตตราติ.\csromanspacer{} ยฺวายํ ภนฺเต ภิกฺขุ ยาวตา สมุทฺทปริยนฺตํ ปถวิํ “มหคฺคตนฺ”ติ ผริตฺวา อธิมุจฺจิตฺวา วิหรติ, อยํ อิมาสํ อุภินฺนํ จิตฺตภาวนานํ มหคฺคตตราติ.\csromanspacer{} อยํ โข อาวุโส กจฺจาน เหตุ อยํ ปจฺจโย, เยน ตาสํ เทวตานํ เอกํ เทวนิกายํ อุปปนฺนานํ สนฺเตตฺถ เอกจฺจา เทวตา ปริตฺตาภา, สนฺติ ปเนตฺถ เอกจฺจา เทวตา อปฺปมาณาภาติ.}

\proseitem{234}{สาธุ ภนฺเต อนุรุทฺธ, อตฺถิ จ เม เอตฺถ อุตฺตริํ ปฏิปุจฺฉิตพฺพํ “ยาวตา\footnote{ยา ตา (ก)} ภนฺเต เทวตา อาภา, สพฺพา ตา สํกิลิฏฺฐาภา, อุทาหุ สนฺเตตฺถ เอกจฺจา เทวตา ปริสุทฺธาภา”ติ.\csromanspacer{} ตทงฺเคน โข อาวุโส กจฺจาน สนฺเตตฺถ เอกจฺจา เทวตา สํกิลิฏฺฐาภา, สนฺติ ปเนตฺถ เอกจฺจา เทวตา ปริสุทฺธาภาติ.\csromanspacer{} โก นุ โข ภนฺเต อนุรุทฺธ เหตุ โก ปจฺจโย, เยน ตาสํ เทวตานํ เอกํ เทวนิกายํ อุปปนฺนานํ สนฺเตตฺถ เอกจฺจา เทวตา สํกิลิฏฺฐาภา, สนฺติ ปเนตฺถ เอกจฺจา เทวตา ปริสุทฺธาภาติ.}

\prose{\csromanfolio{190}เตน หาวุโส กจฺจาน อุปมํ เต กริสฺสามิ, อุปมายปิเธกจฺเจ\footnote{อุปมายมิเธกจฺเจ (ก)} วิญฺญู ปุริสา ภาสิตสฺส อตฺถํ อาชานนฺติ.\csromanspacer{} เสยฺยถาปิ อาวุโส กจฺจาน เตลปฺปทีปสฺส ฌายโต เตลมฺปิ อปริสุทฺธํ วฏฺฏิปิ อปริสุทฺธา, โส เตลสฺสปิ อปริสุทฺธตฺตา วฏฺฏิยาปิ อปริสุทฺธตฺตา อนฺธนฺธํ วิย ฌายติ.\csromanspacer{} เอวเมว โข อาวุโส กจฺจาน อิเธกจฺโจ ภิกฺขุ “สํกิลิฏฺฐาภา”ติ ผริตฺวา อธิมุจฺจิตฺวา วิหรติ, ตสฺส กายทุฏฺฐุลฺลมฺปิ น สุปฺปฏิปฺปสฺสทฺธํ โหติ, ถินมิทฺธมฺปิ น สุสมูหตํ โหติ, อุทฺธจฺจกุกฺกุจฺจมฺปิ น สุปฺปฏิวินีตํ โหติ, โส กายทุฏฺฐุลฺลสฺสปิ น สุปฺปฏิปฺปสฺสทฺธตฺตา ถินมิทฺธสฺสปิ น สุสมูหตตฺตา อุทฺธจฺจกุกฺกุจฺจสฺสปิ น สุปฺปฏิวินีตตฺตา อนฺธนฺธํ วิย ฌายติ, โส กายสฺส เภทา ปรํ มรณา สํกิลิฏฺฐาภานํ เทวานํ สหพฺยตํ อุปปชฺชติ.\csromanspacer{} เสยฺยถาปิ อาวุโส กจฺจาน เตลปฺปทีปสฺส ฌายโต เตลมฺปิ ปริสุทฺธํ วฏฺฏิปิ ปริสุทฺธา, โส เตลสฺสปิ ปริสุทฺธตฺตา วฏฺฏิยาปิ ปริสุทฺธตฺตา น อนฺธนฺธํ วิย ฌายติ.\csromanspacer{} เอวเมว โข อาวุโส กจฺจาน อิเธกจฺโจ ภิกฺขุ “ปริสุทฺธาภา”ติ ผริตฺวา อธิมุจฺจิตฺวา วิหรติ, ตสฺส กายทุฏฺฐุลฺลมฺปิ สุปฺปฏิปฺปสฺสทฺธํ โหติ, ถินมิทฺธมฺปิ สุสมูหตํ โหติ, อุทฺธจฺจกุกฺกุจฺจมฺปิ สุปฺปฏิวินีตํ โหติ, โส กายทุฏฺฐุลฺลสฺสปิ สุปฺปฏิปฺปสฺสทฺธตฺตา ถินมิทฺธสฺสปิ สุสมูหตตฺตา อุทฺธจฺจกุกฺกุจฺจสฺสปิ สุปฺปฏิวินีตตฺตา น อนฺธนฺธํ วิย ฌายติ, โส กายสฺส เภทา ปรํ มรณา ปริสุทฺธาภานํ เทวานํ สหพฺยตํ อุปปชฺชติ.\csromanspacer{} อยํ โข อาวุโส กจฺจาน เหตุ อยํ ปจฺจโย, เยน ตาสํ เทวตานํ เอกํ เทวนิกายํ อุปปนฺนานํ สนฺเตตฺถ เอกจฺจา เทวตา สํกิลิฏฺฐาภา, สนฺติ ปเนตฺถ เอกจฺจา เทวตา ปริสุทฺธาภาติ.}

\proseitem{235}{เอวํ วุตฺเต อายสฺมา สภิโย กจฺจาโน อายสฺมนฺตํ อนุรุทฺธํ เอตทโวจ “สาธุ ภนฺเต อนุรุทฺธ, น ภนฺเต อายสฺมา อนุรุทฺโธ เอวมาห ‘เอวํ เม สุตนฺ’ติ วา, ‘เอวํ อรหติ ภวิตุนฺ’ติ วา, อถ จ ปน ภนฺเต อายสฺมา อนุรุทฺโธ เอวมฺปิ ตา เทวตา อิติปิ ตา เทวตาเตฺวว ภาสติ.\csromanspacer{} ตสฺส มยฺหํ ภนฺเต เอวํ โหติ ‘อทฺธา อายสฺมตา อนุรุทฺเธน ตาหิ เทวตาหิ สทฺธิํ สนฺนิวุตฺถปุพฺพญฺเจว สลฺลปิตปุพฺพญฺจ, สากจฺฉา จ สมาปชฺชิตปุพฺพา’ติ”.\csromanspacer{} อทฺธา โข อยํ อาวุโส กจฺจาน อาสชฺช อุปนีย วาจา ภาสิตา, อปิ จ เต อหํ พฺยากริสฺสามิ, ทีฆรตฺตํ โข เม}

\csromanheader{2}{8. อุปกฺกิเลสสุตฺต (128)}
\prose{\csromanfolio{191}อาวุโส กจฺจาน ตาหิ เทวตาหิ สทฺธิํ สนฺนิวุตฺถปุพฺพญฺเจว สลฺลปิตปุพฺพญฺจ, สากจฺฉา จ สมาปชฺชิตปุพฺพาติ.}

\prose{เอวํ วุตฺเต อายสฺมา สภิโย กจฺจาโน ปญฺจกงฺค ถปติํ เอตทโวจ “ลาภา เต คหปติ, สุลทฺธํ เต คหปติ, ยํ ตฺวํ เจว ตํ กงฺขาธมฺมํ ปหาสิ\footnote{ปชหสิ (ก)}, มยญฺจิมํ\footnote{ยมฺปิมํ (สี, สฺยา, กํ, อิ)} ธมฺมปริยายํ อลตฺถมฺหา สวนายา”ติ.}

\nitthitamruled{อนุรุทฺธสุตฺตํ นิฏฺฐิตํ สตฺตมํ.}
\csromanmatikaanchor{mtk.33}
\csromantocmarkref{subsection}{8. อุปกฺกิเลสสุตฺต}{mtk.33}
\bookmark[dest={mtk.33},level=2]{8. อุปกฺกิเลสสุตฺต}
\csromanheader{2}{8. อุปกฺกิเลสสุตฺต}
\proseitem{236}{เอวํ เม สุตํ\csromandash{}เอกํ สมยํ ภควา โกสมฺพิยํ วิหรติ โฆสิตาราเม.\csromanspacer{} เตน โข ปน สมเยน โกสมฺพิยํ ภิกฺขู ภณฺฑนชาตา กลหชาตา วิวาทาปนฺนา อญฺญมญฺญํ มุขสตฺตีหิ วิตุทนฺตา วิหรนฺติ.\csromanspacer{} อถ โข อญฺญตโร ภิกฺขุ เยน ภควา เตนุปสงฺกมิ, อุปสงฺกมิตฺวา ภควนฺตํ อภิวาเทตฺวา เอกมนฺตํ อฏฺฐาสิ, เอกมนฺตํ ฐิโต โข โส ภิกฺขุ ภควนฺตํ เอตทโวจ “อิธ ภนฺเต โกสมฺพิยํ ภิกฺขู ภณฺฑนชาตา กลหชาตา วิวาทาปนฺนา อญฺญมญฺญํ มุขสตฺตีหิ วิตุทนฺตา วิหรนฺติ, สาธุ ภนฺเต ภควา เยน เต ภิกฺขู เตนุปสงฺกมตุ อนุกมฺปํ อุปาทายา”ติ.\csromanspacer{} อธิวาเสสิ ภควา ตุณฺหีภาเวน.\csromanspacer{} อถ โข ภควา เยน เต ภิกฺขู เตนุปสงฺกมิ, อุปสงฺกมิตฺวา เต ภิกฺขู เอตทโวจ “อลํ ภิกฺขเว มา ภณฺฑนํ มา กลหํ มา วิคฺคหํ มา วิวาทนฺ”ติ.}

\prose{เอวํ วุตฺเต อญฺญตโร ภิกฺขุ ภควนฺตํ เอตทโวจ “อาคเมตุ ภนฺเต ภควา ธมฺมสฺสามี, อปฺโปสฺสุกฺโก ภนฺเต ภควา ทิฏฺฐธมฺมสุขวิหารํ อนุยุตฺโต วิหรตุ, มยเมเตน ภณฺฑเนน กลเหน วิคฺคเหน วิวาเทน ปญฺญายิสฺสามา”ติ.\csromanspacer{} ทุติยมฺปิ โข ภควา เต ภิกฺขู เอตทโวจ “อลํ ภิกฺขเว มา ภณฺฑนํ มา กลหํ มา วิคฺคหํ มา วิวาทนฺ”ติ.\csromanspacer{} ทุติยมฺปิ โข โส ภิกฺขุ ภควนฺตํ เอตทโวจ “อาคเมตุ ภนฺเต ภควา ธมฺมสฺสามี, อปฺโปสฺสุกฺโก ภนฺเต ภควา ทิฏฺฐธมฺมสุขวิหารํ อนุยุตฺโต วิหรตุ, มยเมเตน ภณฺฑเนน กลเหน วิคฺคเหน วิวาเทน ปญฺญายิสฺสามา”ติ. \csromanfolio{192}ตติยมฺปิ โข ภควา เต ภิกฺขู เอตทโวจ “อลํ ภิกฺขเว มา ภณฺฑนํ มา กลหํ มา วิคฺคหํ มา วิทาทนฺ”ติ.\csromanspacer{} ตติยมฺปิ โข โส ภิกฺขุ ภควนฺตํ เอตทโวจ “อาคเมตุ ภนฺเต ภควา ธมฺมสฺสามี, อปฺโปสฺสุกฺโก ภนฺเต ภควา ทิฏฺฐธมฺมสุขวิหารํ อนุยุตฺโต วิหรตุ, มยเมเตน ภณฺฑเนน กลเหน วิคฺคเหน วิวาเทน ปญฺญายิสฺสามา”ติ.}

\prose{อถ โข ภควา ปุพฺพณฺหาสมยํ นิวาเสตฺวา ปตฺตจีวรมาทาย โกสมฺพิํ ปิณฺฑาย ปาวิสิ, โกสมฺพิยํ ปิณฺฑาย จริตฺวา ปจฺฉาภตฺตํ ปิณฺฑปาตปฏิกฺกนฺโต เสนาสนํ สํสาเมตฺวา ปตฺตจีวรมาทาย ฐิตโกว อิมา คาถา อภาสิ\csromandash{}}

\csromanhangingbegin
\hangingheaditem{237}{ปุถุสทฺโท สมชโน, น พาโล โกจิ มญฺญถ.}
\hangingbody{สํฆสฺมิํ ภิชฺชมานสฺมิํ, นาญฺญํ ภิยฺโย อมญฺญรุํ.}
\hangingbody{ปริมุฏฺฐา ปณฺฑิตาภาสา, วาจาโคจรภาณิโน.}
\hangingbody{ยาวิจฺฉนฺติ มุขายามํ, เยน นีตา น ตํ วิทู.}
\hangingbody{อกฺโกจฺฉิ มํ อวธิ มํ, อชินิ มํ อหาสิ เม.}
\hangingbody{เย จ ตํ อุปนยฺหนฺติ, เวรํ เตสํ น สมฺมติ.}
\hangingbody{อกฺโกจฺฉิ มํ อวธิ มํ, อชินิ มํ อหาสิ เม.}
\hangingbody{เย จ ตํ นุปนยฺหนฺติ, เวรํ เตสูปสมฺมติ.}
\hangingbody{น หิ เวเรน เวรานิ, สมฺมนฺตีธ กุทาจนํ.}
\hangingbody{อเวเรน จ สมฺมนฺติ, เอส ธมฺโม สนนฺตโน.}
\hangingbody{ปเร จ น วิชานนฺติ, มยเมตฺถ ยมามเส.}
\hangingbody{เย จ ตตฺถ วิชานนฺติ, ตโต สมฺมนฺติ เมธคา.}
\hangingbody{อฏฺฐิจฺฉินฺนา ปาณหรา, ควสฺสธนหาริโน.}
\hangingbody{รฏฺฐํ วิลุมฺปมานานํ, เตสมฺปิ โหติ สงฺคติ.}
\hangingbody{กสฺมา ตุมฺหาก โน สิยา.}
\hangingbody{สเจ ลเภถ นิปกํ สหายํ,}
\hangingbody{สทฺธิํ จรํ สาธุวิหาริ ธีรํ.}
\hangingbody{อภิภุยฺย สพฺพานิ ปริสฺสยานิ,}
\hangingbody{จเรยฺย เตนตฺตมโน สตีมา.}
\csromanhangingend

\csromanheader{2}{8. อุปกฺกิเลสสุตฺต (128)}
\csromangathameasure{โน เจ ลเภถ นิปกํ สหายํ, \\ สทฺธิํ จรํ สาธุวิหาริ ธีรํ. \\ ราชาว รฏฺฐํ วิชิตํ ปหาย, \\ เอโก จเร มาตงฺครญฺเญว นาโค. \\ เอกสฺส จริตํ เสยฺโย, \\ นตฺถิ พาเล สหายตา. \\ เอโก จเร น จ ปาปานิ กยิรา, \\ อปฺโปสฺสุกฺโก มาตงฺครญฺเญว นาโคติ.}
\csromangathagroup{\csromangathastanza{\csromanfolio{193}โน เจ ลเภถ นิปกํ สหายํ, \\ สทฺธิํ จรํ สาธุวิหาริ ธีรํ. \\ ราชาว รฏฺฐํ วิชิตํ ปหาย, \\ เอโก จเร มาตงฺครญฺเญว นาโค.}\csromangathastanza{เอกสฺส จริตํ เสยฺโย, \\ นตฺถิ พาเล สหายตา. \\ เอโก จเร น จ ปาปานิ กยิรา, \\ อปฺโปสฺสุกฺโก มาตงฺครญฺเญว นาโคติ.}}

\proseitem{238}{อถ โข ภควา ฐิตโกว อิมา คาถา ภาสิตฺวา เยน พาลกโลณการคาโม\footnote{พาลกโลณกคาโม (ก), ตถา วินเยปิ.} เตนุปสงฺกมิ.\csromanspacer{} เตน โข ปน สมเยน อายสฺมา ภคุ พาลกโลณการคาเม วิหรติ.\csromanspacer{} อทฺทสา โข อายสฺมา ภคุ ภควนฺตํ ทูรโตว อาคจฺฉนฺตํ, ทิสฺวาน อาสนํ ปญฺญเปสิ, อุทกญฺจ ปาทานํ โธวนํ\footnote{อุทกญฺจ ปาทานํ (สี, สฺยา, กํ, อิ)}.\csromanspacer{} นิสีทิ ภควา ปญฺญตฺเต อาสเน, นิสชฺช ปาเท ปกฺขาเลสิ.\csromanspacer{} อายสฺมาปิ โข ภคุ ภควนฺตํ อภิวาเทตฺวา เอกมนฺตํ นิสีทิ, เอกมนฺตํ นิสินฺนํ โข อายสฺมนฺตํ ภคุํ ภควา เอตทโวจ “กจฺจิ ภิกฺขุ ขมนียํ, กจฺจิ ยาปนียํ, กจฺจิ ปิณฺฑเกน น กิลมสี”ติ.\csromanspacer{} ขมนียํ ภควา, ยาปนียํ ภควา, น จาหํ ภนฺเต ปิณฺฑเกน กิลมามีติ.\csromanspacer{} อถ โข ภควา อายสฺมนฺตํ ภคุํ ธมฺมิยา กถาย สนฺทสฺเสตฺวา สมาทเปตฺวา สมุตฺเตเชตฺวา สมฺปหํเสตฺวา อุฏฺฐายาสนา เยน ปาจีนวํสทาโย เตนุปสงฺกมิ.}

\prose{เตน โข ปน สมเยน อายสฺมา จ อนุรุทฺโธ อายสฺมา จ นนฺทิโย\footnote{ภทฺทิโย (ม 2 นฬกปาเน 125)} อายสฺมา จ กิมิโล\footnote{กิมฺพิโล (สี, สฺยา, กํ, อิ)} ปาจีนวํสทาเย วิหรนฺติ.\csromanspacer{} อทฺทสา โข ทายปาโล ภควนฺตํ ทูรโตว อาคจฺฉนฺตํ, ทิสฺวาน ภควนฺตํ เอตทโวจ “มา มหาสมณ เอตํ ทายํ ปาวิสิ, สนฺเตตฺถ ตโย กุลปุตฺตา อตฺตกามรูปา วิหรนฺติ, มา เตสํ อผาสุมกาสี”ติ.\csromanspacer{} อสฺโสสิ โข อายสฺมา อนุรุทฺโธ ทายปาลสฺส ภควตา สทฺธิํ มนฺตยมานสฺส, สุตฺวาน ทายปาลํ เอตทโวจ “มา อาวุโส ทายปาล ภควนฺตํ วาเรสิ, สตฺถา โน ภควา อนุปฺปตฺโต”ติ.}

\proseitem{239}{\csromanfolio{194}อถ โข อายสฺมา อนุรุทฺโธ เยนายสฺมา จ นนฺทิโย เยนายสฺมา จ กิมิโล เตนุปสงฺกมิ, อุปสงฺกมิตฺวา อายสฺมนฺตํ จ นนฺทิยํ อายสฺมนฺตํ จ กิมิลํ เอตทโวจ “อภิกฺกมถายสฺมนฺโต อภิกฺกมถายสฺมนฺโต, สตฺถา โน ภควา อนุปฺปตฺโต”ติ.\csromanspacer{} อถ โข อายสฺมา จ อนุรุทฺโธ อายสฺมา จ นนฺทิโย อายสฺมา จ กิมิโล ภควนฺตํ ปจฺจุคฺคนฺตฺวา เอโก ภควโต ปตฺตจีวรํ ปฏิคฺคเหสิ, เอโก อาสนํ ปญฺญเปสิ, เอโก ปาโททกํ อุปฏฺฐเปสิ.\csromanspacer{} นิสีทิ ภควา ปญฺญตฺเต อาสเน, นิสชฺช ปาเท ปกฺขาเลสิ.\csromanspacer{} เตปิ โข อายสฺมนฺโต ภควนฺตํ อภิวาเทตฺวา เอกมนฺตํ นิสีทิํสุ, เอกมนฺตํ นิสินฺนํ โข อายสฺมนฺตํ อนุรุทฺธํ ภควา เอตทโวจ “กจฺจิ โว อนุรุทฺธา ขมนียํ, กจฺจิ ยาปนียํ, กจฺจิ ปิณฺฑเกน น กิลมถา”ติ.\csromanspacer{} ขมนียํ ภควา, ยาปนียํ ภควา, น จ มยํ ภนฺเต ปิณฺฑเกน กิลมามาติ.\csromanspacer{} กจฺจิ ปน โว อนุรุทฺธา สมคฺคา สมฺโมทมานา อวิวทมานา ขีโรทกีภูตา อญฺญมญฺญํ ปิยจกฺขูหิ สมฺปสฺสนฺตา วิหรถาติ.\csromanspacer{} ตคฺฆ มยํ ภนฺเต สมคฺคา สมฺโมทมานา อวิวทมานา ขีโรทกีภูตา อญฺญมญฺญํ ปิยจกฺขูหิ สมฺปสฺสนฺตา วิหรามาติ.\csromanspacer{} ยถา กถํ ปน ตุเมฺห อนุรุทฺธา สมคฺคา สมฺโมทมานา อวิวทมานา ขีโรทกีภูตา อญฺญมญฺญํ ปิยจกฺขูหิ สมฺปสฺสนฺตา วิหรถาติ.\csromanspacer{} อิธ มยฺหํ ภนฺเต เอวํ โหติ “ลาภา วต เม, สุลทฺธํ วต เม, โยหํ เอวรูเปหิ สพฺรหฺมจารีหิ สทฺธิํ วิหรามี”ติ.\csromanspacer{} ตสฺส มยฺหํ ภนฺเต อิเมสุ อายสฺมนฺเตสุ เมตฺตํ กายกมฺมํ ปจฺจุปฏฺฐิตํ อาวิ เจว รโห จ, เมตฺตํ วจีกมฺมํ ปจฺจุปฏฺฐิตํ อาวิ เจว รโห จ, เมตฺตํ มโนกมฺมํ ปจฺจุปฏฺฐิตํ อาวิ เจว รโห จ.\csromanspacer{} ตสฺส มยฺหํ ภนฺเต เอวํ โหติ “ยํนูนาหํ สกํ จิตฺตํ นิกฺขิปิตฺวา อิเมสํเยว อายสฺมนฺตานํ จิตฺตสฺส วเสน วตฺเตยฺยนฺ”ติ.\csromanspacer{} โส โข อหํ ภนฺเต สกํ จิตฺตํ นิกฺขิปิตฺวา อิเมสํเยว อายสฺมนฺตานํ จิตฺตสฺส วเสน วตฺตามิ.\csromanspacer{} นานา หิ โข โน ภนฺเต กายา, เอกญฺจ ปน มญฺเญ จิตฺตนฺติ.}

\prose{อายสฺมาปิ โข นนฺทิโย ฯเปฯ.\csromanspacer{} อายสฺมาปิ โข กิมิโล ภควนฺตํ เอตทโวจ “มยฺหมฺปิ โข ภนฺเต เอวํ โหติ ‘ลาภา วต เม, สุลทฺธํ วต เม, โยหํ เอวรูเปหิ สพฺรหฺมจารีหิ สทฺธิํ วิหรามี’ติ.\csromanspacer{} ตสฺส มยฺหํ ภนฺเต อิเมสุ อายสฺมนฺเตสุ เมตฺตํ กายกมฺมํ ปจฺจุปฏฺฐิตํ อาวิ เจว รโห จ, เมตฺตํ วจีกมฺมํ ปจฺจุปฏฺฐิตํ อาวิ เจว รโห จ, เมตฺตํ มโนกมฺมํ}

\csromanheader{2}{8. อุปกฺกิเลสสุตฺต (128)}
\prose{\csromanfolio{195}ปจฺจุปฏฺฐิตํ อาวิ เจว รโห จ.\csromanspacer{} ตสฺส มยฺหํ ภนฺเต เอวํ โหติ ‘ยํนูนาหํ สกํ จิตฺตํ นิกฺขิปิตฺวา อิเมสํเยว อายสฺมนฺตานํ จิตฺตสฺส วเสน วตฺเตยฺยนฺ’ติ.\csromanspacer{} โส โข อหํ ภนฺเต สกํ จิตฺตํ นิกฺขิปิตฺวา อิเมสํเยว อายสฺมนฺตานํ จิตฺตสฺส วเสน วตฺตามิ.\csromanspacer{} นานา หิ โข โน ภนฺเต กายา, เอกญฺจ ปน มญฺเญ จิตฺตนฺติ.\csromanspacer{} เอวํ โข มยํ ภนฺเต สมคฺคา สมฺโมทมานา อวิวทมานา ขีโรทกีภูตา อญฺญมญฺญํ ปิยจกฺขูหิ สมฺปสฺสนฺตา วิหรามา”ติ.}

\proseitem{240}{สาธุ สาธุ อนุรุทฺธา, กจฺจิ ปน โว อนุรุทฺธา อปฺปมตฺตา อาตาปิโน ปหิตตฺตา วิหรถาติ.\csromanspacer{} ตคฺฆ มยํ ภนฺเต อปฺปมตฺตา อาตาปิโน ปหิตตฺตา วิหรามาติ.\csromanspacer{} ยถา กถํ ปน ตุเมฺห อนุรุทฺธา อปฺปมตฺตา อาตาปิโน ปหิตตฺตา วิหรถาติ.\csromanspacer{} อิธ ภนฺเต อมฺหากํ โย ปฐมํ คามโต ปิณฺฑาย ปฏิกฺกมติ, โส อาสนานิ ปญฺญเปติ, ปานียํ ปริโภชนียํ อุปฏฺฐาเปติ, อวกฺการปาติํ อุปฏฺฐาเปติ.\csromanspacer{} โย ปจฺฉา คามโต ปิณฺฑาย ปฏิกฺกมติ, สเจ โหติ ภุตฺตาวเสโส, สเจ อากงฺขติ ภุญฺชติ, โน เจ อากงฺขติ อปฺปหริเต วา ฉฑฺเฑติ, อปาณเก วา อุทเก โอปิลาเปติ, โส อาสนานิ ปฏิสาเมติ, ปานียํ ปริโภชนียํ ปฏิสาเมติ, อวกฺการปาติํ โธวิตฺวา ปฏิสาเมติ, ภตฺตคฺคํ สมฺมชฺชติ.\csromanspacer{} โย ปสฺสติ ปานียฆฏํ วา ปริโภชนียฆฏํ วา วจฺจฆฏํ วา ริตฺตํ ตุจฺฉํ, โส อุปฏฺฐาเปติ.\csromanspacer{} สจสฺส โหติ อวิสยฺหํ, หตฺถวิกาเรน ทุติยํ อามนฺเตตฺวา หตฺถวิลงฺฆเกน อุปฏฺฐาเปม\footnote{อุปฏฺฐเปติ (สี)}, น เตฺวว มยํ ภนฺเต ตปฺปจฺจยา วาจํ ภินฺทาม.\csromanspacer{} ปญฺจาหิกํ โข ปน มยํ ภนฺเต สพฺพรตฺติํ ธมฺมิยา กถาย สนฺนิสีทาม.\csromanspacer{} เอวํ โข มยํ ภนฺเต อปฺปมตฺตา อาตาปิโน ปหิตตฺตา วิหรามาติ.}

\proseitem{241}{สาธุ สาธุ อนุรุทฺธา, อตฺถิ ปน โว อนุรุทฺธา เอวํ อปฺปมตฺตานํ อาตาปีนํ ปหิตตฺตานํ วิหรตํ อุตฺตริ มนุสฺสธมฺมา อลมริยญาณทสฺสนวิเสโส อธิคโต ผาสุวิหาโรติ.\csromanspacer{} อิธ มยํ ภนฺเต อปฺปมตฺตา อาตาปิโน ปหิตตฺตา วิหรนฺตา โอภาสญฺเจว สญฺชานาม ทสฺสนญฺจ รูปานํ.\csromanspacer{} โส โข ปน โน โอภาโส นจิรสฺเสว อนฺตรธายติ ทสฺสนญฺจ รูปานํ, ตญฺจ นิมิตฺตํ นปฺปฏิวิชฺฌามาติ.}

\prose{\csromanfolio{196}ตํ โข ปน โว อนุรุทฺธา นิมิตฺตํ ปฏิวิชฺฌิตพฺพํ.\csromanspacer{} อหมฺปิ สุทํ อนุรุทฺธา ปุพฺเพว สมฺโพธา อนภิสมฺพุทฺโธ โพธิสตฺโตว สมาโน โอภาสญฺเจว สญฺชานามิ ทสฺสนญฺจ รูปานํ.\csromanspacer{} โส โข ปน เม โอภาโส นจิรสฺเสว อนฺตรธายติ ทสฺสนญฺจ รูปานํ.\csromanspacer{} ตสฺส มยฺหํ อนุรุทฺธา เอตทโหสิ “โก นุ โข เหตุ โก ปจฺจโย, เยน เม โอภาโส อนฺตรธายติ ทสฺสนญฺจ รูปานนฺ”ติ.\csromanspacer{} ตสฺส มยฺหํ อนุรุทฺธา เอตทโหสิ “วิจิกิจฺฉา โข เม อุทปาทิ, วิจิกิจฺฉาธิกรณญฺจ ปน เม สมาธิ จวิ, สมาธิมฺหิ จุเต โอภาโส อนฺตรธายติ ทสฺสนญฺจ รูปานํ, โสหํ ตถา กริสฺสามิ ยถา เม ปุน น วิจิกิจฺฉา อุปฺปชฺชิสฺสตี”ติ. (1)}

\prose{โส โข อหํ อนุรุทฺธา อปฺปมตฺโต อาตาปี ปหิตตฺโต วิหรนฺโต โอภาสญฺเจว สญฺชานามิ ทสฺสนญฺจ รูปานํ.\csromanspacer{} โส โข ปน เม โอภาโส นจิรสฺเสว อนฺตรธายติ ทสฺสนญฺจ รูปานํ.\csromanspacer{} ตสฺส มยฺหํ อนุรุทฺธา เอตทโหสิ “โก นุ โข เหตุ โก ปจฺจโย, เยน เม โอภาโส อนฺตรธายติ ทสฺสนญฺจ รูปานนฺ”ติ.\csromanspacer{} ตสฺส มยฺหํ อนุรุทฺธา เอตทโหสิ “อมนสิกาโร โข เม อุทปาทิ, อมนสิการาธิกรณญฺจ ปน เม สมาธิ จวิ, สมาธิมฺหิ จุเต โอภาโส อนฺตรธายติ ทสฺสนญฺจ รูปานํ, โสหํ ตถา กริสฺสามิ ยถา เม ปุน น วิจิกิจฺฉา อุปฺปชฺชิสฺสติ, น อมนสิกาโร”ติ. (2)}

\prose{โส โข อหํ อนุรุทฺธา ฯเปฯ.\csromanspacer{} ตสฺส มยฺหํ อนุรุทฺธา เอตทโหสิ “ถินมิทฺธํ โข เม อุทปาทิ, ถินมิทฺธาธิกรณญฺจ ปน เม สมาธิ จวิ, สมาธิมฺหิ จุเต โอภาโส อนฺตรธายติ ทสฺสนญฺจ รูปานํ, โสหํ ตถา กริสฺสามิ ยถา เม ปุน น วิจิกิจฺฉา อุปฺปชฺชิสฺสติ, น อมนสิกาโร, น ถินมิทฺธนฺ”ติ. (3)}

\prose{โส โข อหํ อนุรุทฺธา ฯเปฯ.\csromanspacer{} ตสฺส มยฺหํ อนุรุทฺธา เอตทโหสิ “ฉมฺภิตตฺตํ โข เม อุทปาทิ, ฉมฺภิตตฺตาธิกรณญฺจ ปน เม สมาธิ จวิ, สมาธิมฺหิ จุเต โอภาโส อนฺตรธายติ ทสฺสนญฺจ รูปานํ.\csromanspacer{} เสยฺยถาปิ อนุรุทฺธา ปุริโส อทฺธานมคฺคปฺปฏิปนฺโน, ตสฺส อุภโตปสฺเส วฏฺฏกา\footnote{วธกา (สี, สฺยา, กํ, อิ)} อุปฺปเตยฺยุํ, ตสฺส ตโตนิทานํ ฉมฺภิตตฺตํ อุปฺปชฺเชยฺย.\csromanspacer{} เอวเมว โข เม อนุรุทฺธา ฉมฺภิตตฺตํ อุทปาทิ, ฉมฺภิตตฺตาธิกรณญฺจ ปน เม สมาธิ จวิ,}

\csromanheader{2}{8. อุปกฺกิเลสสุตฺต (128)}
\prose{\csromanfolio{197}สมาธิมฺหิ จุเต โอภาโส อนฺตรธายติ ทสฺสนญฺจ รูปานํ, โสหํ ตถา กริสฺสามิ ยถา เม ปุน น วิจิกิจฺฉา อุปฺปชฺชิสฺสติ, น อมนสิกาโร, น ถินมิทฺธํ, น ฉมฺภิตตฺตนฺ”ติ. (4)}

\prose{โส โข อหํ อนุรุทฺธา ฯเปฯ.\csromanspacer{} ตสฺส มยฺหํ อนุรุทฺธา เอตทโหสิ “อุปฺปิลํ\footnote{อุพฺพิลฺลํ (สี, อิ), อุพฺพิลํ (สฺยา, กํ)} โข เม อุทปาทิ, อุปฺปิลาธิกรณญฺจ ปน เม สมาธิ จวิ, สมาธิมฺหิ จุเต โอภาโส อนฺตรธายติ ทสฺสนญฺจ รูปานํ.\csromanspacer{} เสยฺยถาปิ อนุรุทฺธา ปุริโส เอกํ นิธิมุขํ คเวสนฺโต สกิเทว ปญฺจนิธิมุขานิ อธิคจฺเฉยฺย, ตสฺส ตโตนิทานํ อุปฺปิลํ อุปฺปชฺเชยฺย.\csromanspacer{} เอวเมว โข เม อนุรุทฺธา อุปฺปิลํ อุทปาทิ, อุปฺปิลาธิกรณญฺจ ปน เม สมาธิ จวิ, สมาธิมฺหิ จุเต โอภาโส อนฺตรธายติ ทสฺสนญฺจ รูปานํ, โสหํ ตถา กริสฺสามิ ยถา เม ปุน น วิจิกิจฺฉา อุปฺปชฺชิสฺสติ, น อมนสิกาโร, น ถินมิทฺธํ, น ฉมฺภิตตฺตํ, น อุปฺปิลนฺ”ติ. (5)}

\prose{โส โข อหํ อนุรุทฺธา ฯเปฯ.\csromanspacer{} ตสฺส มยฺหํ อนุรุทฺธา เอตทโหสิ “ทุฏฺฐุลฺลํ โข เม อุทปาทิ, ทุฏฺฐุลฺลาธิกรณญฺจ ปน เม สมาธิ จวิ, สมาธิมฺหิ จุเต โอภาโส อนฺตรธายติ ทสฺสนญฺจ รูปานํ, โสหํ ตถา กริสฺสามิ ยถา เม ปุน น วิจิกิจฺฉา อุปฺปชฺชิสฺสติ, น อมนสิกาโร, น ถินมิทฺธํ, น ฉมฺภิตตฺตํ, น อุปฺปิลํ, น ทุฏฺฐุลฺลนฺ”ติ. (6)}

\prose{โส โข อหํ อนุรุทฺธา ฯเปฯ.\csromanspacer{} ตสฺส มยฺหํ อนุรุทฺธา เอตทโหสิ “อจฺจารทฺธวีริยํ โข เม อุทปาทิ, อจฺจารทฺธวีริยาธิกรณญฺจ ปน เม สมาธิ จวิ, สมาธิมฺหิ จุเต โอภาโส อนฺตรธายติ ทสฺสนญฺจ รูปานํ.\csromanspacer{} เสยฺยถาปิ อนุรุทฺธา ปุริโส อุโภหิ หตฺเถหิ วฏฺฏกํ คาฬฺหํ คเณฺหยฺย, โส ตตฺเถว ปตเมยฺย\footnote{มตเมยฺย (พหูสุ) ป + ตม + เอยฺย = ปตเมยฺย-อิติ ปทวิภาโค.}.\csromanspacer{} เอวเมว โข เม อนุรุทฺธา อจฺจารทฺธํวีริยํ อุทปาทิ, อจฺจารทฺธวีริยาธิกรณญฺจ ปน เม สมาธิ จวิ, สมาธิมฺหิ จุเต โอภาโส อนฺตรธายติ ทสฺสนญฺจ รูปานํ, โสหํ ตถา กริสฺสามิ ยถา เม ปุน น วิจิกิจฺฉา อุปฺปชฺชิสฺสติ, น อมนสิกาโร, น ถินมิทฺธํ, น ฉมฺภิตตฺตํ, น อุปฺปิลํ, น ทุฏฺฐุลฺลํ, น อจฺจารทฺธวีริยนฺ”ติ. (7)}

\prose{โส โข อหํ อนุรุทฺธา ฯเปฯ.\csromanspacer{} ตสฺส มยฺหํ อนุรุทฺธา เอตทโหสิ “อติลีนวีริยํ โข เม อุทปาทิ, อติลีนวีริยาธิกรณญฺจ ปน}

\prose{\csromanfolio{198}เม สมาธิ จวิ, สมาธิมฺหิ จุเต โอภาโส อนฺตรธายติ ทสฺสนญฺจ รูปานํ.\csromanspacer{} เสยฺยถาปิ อนุรุทฺธา ปุริโส วฏฺฏกํ สิถิลํ คเณฺหยฺย, โส ตสฺส หตฺถโต อุปฺปเตยฺย.\csromanspacer{} เอวเมว โข เม อนุรุทฺธา อติลีนวีริยํ อุทปาทิ, อติลีนวีริยาธิกรณญฺจ ปน เม สมาธิ จวิ, สมาธิมฺหิ จุเต โอภาโส อนฺตรธายติ ทสฺสนญฺจ รูปานํ, โสหํ ตถา กริสฺสามิ ยถา เม ปุน น วิจิกิจฺฉา อุปฺปชฺชิสฺสติ, น อมนสิกาโร, น ถินมิทฺธํ, น ฉมฺภิตตฺตํ, น อุปฺปิลํ, น ทุฏฺฐุลฺลํ, น อจฺฉารทฺธวีริยํ, น อติลีนวีริยนฺ”ติ. (8)}

\prose{โส โข อหํ อนุรุทฺธา ฯเปฯ.\csromanspacer{} ตสฺส มยฺหํ อนุรุทฺธา เอตทโหสิ “อภิชปฺปา โข เม อุทปาทิ, อภิชปฺปาธิกรณญฺจ ปน เม สมาธิ จวิ, สมาธิมฺหิ จุเต โอภาโส อนฺตรธายติ ทสฺสนญฺจ รูปานํ.\csromanspacer{} โสหํ ตถา กริสฺสามิ, ยถา เม ปุน น วิจิกิจฺฉา อุปฺปชฺชิสฺสติ, น อมนสิกาโร, น ถินมิทฺธํ, น ฉมฺภิตตฺตํ, น อุปฺปิลํ, น ทุฏฺฐุลฺลํ, น อจฺจารทฺธวีริยํ, น อติลีนวีริยํ, น อภิชปฺปา”ติ. (9)}

\prose{โส โข อหํ อนุรุทฺธา ฯเปฯ.\csromanspacer{} ตสฺส มยฺหํ อนุรุทฺธา เอตทโหสิ “นานตฺตสญฺญา โข เม อุทปาทิ, นานตฺตสญฺญาธิกรณญฺจ ปน เม สมาธิ จวิ, สมาธิมฺหิ จุเต โอภาโส อนฺตรธายติ ทสฺสนญฺจ รูปานํ.\csromanspacer{} โสหํ ตถา กริสฺสามิ, ยถา เม ปุน น วิจิกิจฺฉา อุปฺปชฺชิสฺสติ, น อมนสิกาโร, น ถินมิทฺธํ, น ฉมฺภิตตฺตํ, น อุปฺปิลํ, น ทุฏฺฐุลฺลํ, น อจฺจารทฺธวีริยํ, น อติลีนวีริยํ, น อภิชปฺปา, น นานตฺตสญฺญา”ติ. (10)}

\prose{โส โข อหํ อนุรุทฺธา อปฺปมตฺโต อาตาปี ปหิตตฺโต วิหรนฺโต โอภาสญฺเจว สญฺชานามิ ทสฺสนญฺจ รูปานํ.\csromanspacer{} โส โข ปน เม โอภาโส นจิรสฺเสว อนฺตรธายติ ทสฺสนญฺจ รูปานํ.\csromanspacer{} ตสฺส มยฺหํ อนุรุทฺธา เอตทโหสิ “โก นุ โข เหตุ โก ปจฺจโย, เยน เม โอภาโส อนฺตรธายติ ทสฺสนญฺจ รูปานนฺ”ติ.\csromanspacer{} ตสฺส มยฺหํ อนุรุทฺธา เอตทโหสิ “อตินิชฺฌายิตตฺตํ โข เม รูปานํ อุทปาทิ, อตินิชฺฌายิตตฺตาธิกรณญฺจ ปน เม รูปานํ สมาธิ จวิ, สมาธิมฺหิ จุเต โอภาโส อนฺตรธายติ ทสฺสนญฺจ รูปานํ.\csromanspacer{} โสหํ ตถา กริสฺสามิ, ยถา เม ปุน น วิจิกิจฺฉา อุปฺปชฺชิสฺสติ, น อมนสิกาโร, น ถินมิทฺธํ, น ฉมฺภิตตฺตํ, น อุปฺปิลํ, น ทุฏฺฐุลฺลํ, น อจฺจารทฺธวีริยํ, น อติลีนวีริยํ, น อภิชปฺปา, น นานตฺตสญฺญา, น อตินิชฺฌายิตตฺตํ รูปานนฺ”ติ. (11)}

\csromanheader{2}{8. อุปกฺกิเลสสุตฺต (128)}
\proseitem{242}{\csromanfolio{199}โส โข อหํ อนุรุทฺธา “วิจิกิจฺฉา จิตฺตสฺส อุปกฺกิเลโส”ติ อิติ วิทิตฺวา วิจิกิจฺฉํ จิตฺตสฺส อุปกฺกิเลสํ ปชหิํ. “อมนสิกาโร จิตฺตสฺส อุปกฺกิเลโส”ติ อิติ วิทิตฺวา อมนสิการํ จิตฺตสฺส อุปกฺกิเลสํ ปชหิํ. “ถินมิทฺธํ จิตฺตสฺส อุปกฺกิเลโส”ติ อิติ วิทิตฺวา ถินมิทฺธํ จิตฺตสฺส อุปกฺกิเลสํ ปชหิํ. “ฉมฺภิตตฺตํ จิตฺตสฺส อุปกฺกิเลโส”ติ อิติ วิทิตฺวา ฉมฺภิตตฺตํ จิตฺตสฺส อุปกฺกิเลสํ ปชหิํ. “อุปฺปิลํ จิตฺตสฺส อุปกฺกิเลโส”ติ อิติ วิทิตฺวา อุปฺปิลํ จิตฺตสฺส อุปกฺกิเลสํ ปชหิํ. “ทุฏฺฐุลฺลํ จิตฺตสฺส อุปกฺกิเลโส”ติ อิติ วิทิตฺวา ทุฏฺฐุลฺลํ จิตฺตสฺส อุปกฺกิเลสํ ปชหิํ. “อจฺจารทฺธวีริยํ จิตฺตสฺส อุปกฺกิเลโส”ติ อิติ วิทิตฺวา อจฺจารทฺธวีริยํ จิตฺตสฺส อุปกฺกิเลสํ ปชหิํ. “อติลีนวีริยํ จิตฺตสฺส อุปกฺกิเลโส”ติ อิติ วิทิตฺวา อติลีนวีริยํ จิตฺตสฺส อุปกฺกิเลสํ ปชหิํ. “อภิชปฺปา จิตฺตสฺส อุปกฺกิเลโส”ติ อิติ วิทิตฺวา อภิชปฺปํ จิตฺตสฺส อุปกฺกิเลสํ ปชหิํ. “นานตฺตสญฺญา จิตฺตสฺส อุปกฺกิเลโส”ติ อิติ วิทิตฺวา นานตฺตสญฺญํ จิตฺตสฺส อุปกฺกิเลสํ ปชหิํ. “อตินิชฺฌายิตตฺตํ รูปานํ จิตฺตสฺส อุปกฺกิเลโส”ติ อิติ วิทิตฺวา อตินิชฺฌายิตตฺตํ รูปานํ จิตฺตสฺส อุปกฺกิเลสํ ปชหิํ.}

\proseitem{243}{โส โข อหํ อนุรุทฺธา อปฺปมตฺโต อาตาปี ปหิตตฺโต วิหรนฺโต โอภาสํ หิ โข สญฺชานามิ, น จ รูปานิ ปสฺสามิ.\csromanspacer{} รูปานิ หิ โข ปสฺสามิ, น จ โอภาสํ สญฺชานามิ เกวลมฺปิ รตฺติํ เกวลมฺปิ ทิวํ\footnote{ทิวสํ (สี, สฺยา, กํ, อิ)} เกวลมฺปิ รตฺตินฺทิวํ\footnote{รตฺติทิวํ (ก)}.\csromanspacer{} ตสฺส มยฺหํ อนุรุทฺธา เอตทโหสิ “โก นุ โข เหตุ โก ปจฺจโย, ยฺวาหํ โอภาสํ หิ โข สญฺชานามิ, น จ รูปานิ ปสฺสามิ.\csromanspacer{} รูปานิ หิ โข\footnote{โข ตสฺมิํ สมเย (สี, ก)} ปสฺสามิ, น จ โอภาสํ สญฺชานามิ เกวลมฺปิ รตฺติํ เกวลมฺปิ ทิวํ เกวลมฺปิ รตฺตินฺทิวนฺ”ติ.\csromanspacer{} ตสฺส มยฺหํ อนุรุทฺธา เอตทโหสิ “ยสฺมิํ หิ โข อหํ สมเย รูปนิมิตฺตํ อมนสิกริตฺวา โอภาสนิมิตฺตํ มนสิ กโรมิ, โอภาสํ หิ โข ตสฺมิํ สมเย สญฺชานามิ, น จ รูปานิ ปสฺสามิ.\csromanspacer{} ยสฺมิํ ปนาหํ สมเย โอภาสนิมิตฺตํ อมนสิกริตฺวา รูปนิมิตฺตํ มนสิ กโรมิ, รูปานิ หิ โข ตสฺมิํ สมเย ปสฺสามิ, น จ โอภาสํ สญฺชานามิ เกวลมฺปิ รตฺติํ เกวลมฺปิ ทิวํ เกวลมฺปิ รตฺตินฺทิวนฺ”ติ.}

\prose{\csromanfolio{200}โส โข อหํ อนุรุทฺธา อปฺปมตฺโต อาตาปี ปหิตตฺโต วิหรนฺโต ปริตฺตํ เจว โอภาสํ สญฺชานามิ, ปริตฺตานิ จ รูปานิ ปสฺสามิ.\csromanspacer{} อปฺปมาณํ เจว โอภาสํ สญฺชานามิ, อปฺปมาณานิ จ รูปานิ ปสฺสามิ เกวลมฺปิ รตฺติํ เกวลมฺปิ ทิวํ เกวลมฺปิ รตฺตินฺทิวํ.\csromanspacer{} ตสฺส มยฺหํ อนุรุทฺธา เอตทโหสิ “โก นุ โข เหตุ โก ปจฺจโย, ยฺวาหํ ปริตฺตญฺเจว โอภาสํ สญฺชานามิ, ปริตฺตานิ จ รูปานิ ปสฺสามิ.\csromanspacer{} อปฺปมาณญฺเจว โอภาสํ สญฺชานามิ, อปฺปมาณานิ จ รูปานิ ปสฺสามิ เกวลมฺปิ รตฺติํ เกวลมฺปิ ทิวํ เกวลมฺปิ รตฺตินฺทิวนฺ”ติ.\csromanspacer{} ตสฺส มยฺหํ อนุรุทฺธา เอตทโหสิ “ยสฺมิํ โข เม สมเย ปริตฺโต สมาธิ โหติ, ปริตฺตํ เม ตสฺมิํ สมเย จกฺขุ โหติ.\csromanspacer{} โสหํ ปริตฺเตน จกฺขุนา ปริตฺตญฺเจว โอภาสํ สญฺชานามิ, ปริตฺตานิ จ รูปานิ ปสฺสามิ.\csromanspacer{} ยสฺมิํ ปน เม สมเย อปฺปมาโณ สมาธิ โหติ, อปฺปมาณํ เม ตสฺมิํ สมเย จกฺขุ โหติ.\csromanspacer{} โสหํ อปฺปมาเณน จกฺขุนา อปฺปมาณญฺเจว โอภาสํ สญฺชานามิ, อปฺปมาณานิ จ รูปานิ ปสฺสามิ เกวลมฺปิ รตฺติํ เกวลมฺปิ ทิวํ เกวลมฺปิ รตฺตินฺทิวนฺ”ติ.}

\proseitem{244}{ยโต โข เม อนุรุทฺธา “วิจิกิจฺฉา จิตฺตสฺส อุปกฺกิเลโส”ติ อิติ วิทิตฺวา วิจิกิจฺฉา จิตฺตสฺส อุปกฺกิเลโส ปหีโน อโหสิ. “อมนสิกาโร จิตฺตสฺส อุปกฺกิเลโส”ติ อิติ วิทิตฺวา อมนสิกาโร จิตฺตสฺส อุปกฺกิเลโส ปหีโน อโหสิ. “ถินมิทฺธํ จิตฺตสฺส อุปกฺกิเลโส”ติ อิติ วิทิตฺวา ถินมิทฺธํ จิตฺตสฺส อุปกฺกิเลโส ปหีโน อโหสิ. “ฉมฺภิตตฺตํ จิตฺตสฺส อุปกฺกิเลโส”ติ อิติ วิทิตฺวา ฉมฺภิตตฺตํ จิตฺตสฺส อุปกฺกิเลโส ปหีโน อโหสิ. “อุปฺปิลํ จิตฺตสฺส อุปกฺกิเลโส”ติ วิทิตฺวา อุปฺปิลํ จิตฺตสฺส อุปกฺกิเลโส ปหีโน อโหสิ. “ทุฏฺฐุลฺลํ จิตฺตสฺส อุปกฺกิเลโส”ติ อิติ วิทิตฺวา ทุฏฺฐุลฺลํ จิตฺตสฺส อุปกฺกิเลโส ปหีโน อโหสิ. “อจฺจารทฺธวีริยํ จิตฺตสฺส อุปกฺกิเลโส”ติ อิติ วิทิตฺวา อจฺจารทฺธวีริยํ จิตฺตสฺส อุปกฺกิเลโส ปหีโน อโหสิ. “อติลีนวีริยํ จิตฺตสฺส อุปกฺกิเลโส”ติ อิติ วิทิตฺวา อติลีนวีริยํ จิตฺตสฺส อุปกฺกิเลโส ปหีโน อโหสิ. “อภิชปฺปา จิตฺตสฺส อุปกฺกิเลโส”ติ อิติ วิทิตฺวา อภิชปฺปา จิตฺตสฺส อุปกฺกิเลโส ปหีโน อโหสิ. “นานตฺตสญฺญา จิตฺตสฺส อุปกฺกิเลโส”ติ อิติ วิทิตฺวา นานตฺตสญฺญา จิตฺตสฺส อุปกฺกิเลโส ปหีโน อโหสิ. “อตินิชฺฌายิตตฺตํ รูปานํ จิตฺตสฺส อุปกฺกิเลโส”ติ อิติ}

\csromanheader{2}{9. พาลปณฺฑิตสุตฺต (129)}
\prose{\csromanfolio{201}วิทิตฺวา อตินิชฺฌายิตตฺตํ รูปานํ จิตฺตสฺส อุปกฺกิเลโส ปหีโน อโหสิ.}

\proseitem{245}{ตสฺส มยฺหํ อนุรุทฺธา เอตทโหสิ “เย โข เม จิตฺตสฺส อุปกฺกิเลสา, เต เม ปหีนา.\csromanspacer{} หนฺท ทานาหํ ติวิเธน สมาธิํ ภาเวมี”ติ\footnote{ภาเวสินฺติ (สี, สฺยา, กํ)}.\csromanspacer{} โส โข อหํ อนุรุทฺธา สวิตกฺกมฺปิ สวิจารํ สมาธิํ ภาเวสิํ\footnote{ภาเวมิ (ก)}, อวิตกฺกมฺปิ วิจารมตฺตํ สมาธิํ ภาเวสิํ, อวิตกฺกมฺปิ อวิจารํ สมาธิํ ภาเวสิํ, สปฺปีติกมฺปิ สมาธิํ ภาเวสิํ, นิปฺปีติกมฺปิ สมาธิํ ภาเวสิํ, สาตสหคตมฺปิ สมาธิํ ภาเวสิํ, อุเปกฺขาสหคตมฺปิ สมาธิํ ภาเวสิํ.\csromanspacer{} ยโต โข เม อนุรุทฺธา สวิตกฺโกปิ สวิจาโร สมาธิ ภาวิโต อโหสิ, อวิตกฺโกปิ วิจารมตฺโต สมาธิ ภาวิโต อโหสิ, อวิตกฺโกปิ อวิจาโร สมาธิ ภาวิโต อโหสิ, สปฺปีติโกปิ สมาธิ ภาวิโต อโหสิ, นิปฺปีติโกปิ สมาธิ ภาวิโต อโหสิ, สาตสหคโตปิ สมาธิ ภาวิโต อโหสิ, อุเปกฺขาสหคโตปิ สมาธิ ภาวิโต อโหสิ, ญาณญฺจ ปน เม ทสฺสนํ อุทปาทิ, อกุปฺปา เม เจโตวิมุตฺติ, อยมนฺติมา ชาติ, นตฺถิ ทานิ ปุนพฺภโวติ.}

\prose{อิทมโวจ ภควา.\csromanspacer{} อตฺตมโน อายสฺมา อนุรุทฺโธ ภควโต ภาสิตํ อภินนฺทีติ.}

\nitthitamruled{อุปกฺกิเลสสุตฺตํ นิฏฺฐิตํ อฏฺฐมํ.}
\csromanmatikaanchor{mtk.34}
\csromantocmarkref{subsection}{9. พาลปณฺฑิตสุตฺต}{mtk.34}
\bookmark[dest={mtk.34},level=2]{9. พาลปณฺฑิตสุตฺต}
\csromanheader{2}{9. พาลปณฺฑิตสุตฺต}
\proseitem{246}{เอวํ เม สุตํ\csromandash{}เอกํ สมยํ ภควา สาวตฺถิยํ วิหรติ เชตวเน อนาถปิณฺฑิกสฺส อาราเม.\csromanspacer{} ตตฺร โข ภควา ภิกฺขู อามนฺเตสิ “ภิกฺขโว”ติ. “ภทนฺเต”ติ เต ภิกฺขู ภควโต ปจฺจสฺโสสุํ.\csromanspacer{} ภควา เอตทโวจ\csromandash{}}

\prose{ตีณิมานิ ภิกฺขเว พาลสฺส พาลลกฺขณานิ พาลนิมิตฺตานิ พาลาปทานานิ.\csromanspacer{} กตมานิ ตีณิ, อิธ ภิกฺขเว พาโล ทุจฺจินฺติตจินฺตี จ โหติ ทุพฺภาสิตภาสี จ ทุกฺกฏกมฺมการี จ.\csromanspacer{} โน เจตํ\footnote{โน เจทํ (สํ 2. 25)} ภิกฺขเว \csromanfolio{202}พาโล ทุจฺจินฺติตจินฺตี จ อภวิสฺส ทุพฺภาสิตภาสี จ ทุกฺกฏกมฺมการี จ, เกน นํ\footnote{น เตน นํ (ก), น นํ (?)} ปณฺฑิตา ชาเนยฺยุํ “พาโล อยํ ภวํ อสปฺปุริโส”ติ.\csromanspacer{} ยสฺมา จ โข ภิกฺขเว พาโล ทุจฺจินฺติตจินฺตี จ โหติ ทุพฺภาสิตภาสี จ ทุกฺกฏกมฺมการี จ, ตสฺมา นํ ปณฺฑิตา ชานนฺติ “พาโล อยํ ภวํ อสปฺปุริโส”ติ.\csromanspacer{} ส โข โส ภิกฺขเว พาโล ติวิธํ ทิฏฺเฐว ธมฺเม ทุกฺขํ โทมนสฺสํ ปฏิสํเวเทติ.\csromanspacer{} สเจ ภิกฺขเว พาโล สภายํ วา นิสินฺโน โหติ, รถิกาย\footnote{รถิยาย (พหูสุ)} วา นิสินฺโน โหติ, สิงฺฆาฏเก วา นิสินฺโน โหติ.\csromanspacer{} ตตฺร เจ ชโน ตชฺชํ ตสฺสารุปฺปํ กถํ มนฺเตติ.\csromanspacer{} สเจ ภิกฺขเว พาโล ปาณาติปาตี โหติ, อทินฺนาทายี โหติ, กาเมสุมิจฺฉาจารี โหติ, มุสาวาที โหติ, สุราเมรยมชฺชปมาทฏฺฐายี โหติ.\csromanspacer{} ตตฺร ภิกฺขเว พาลสฺส เอวํ โหติ “ยํ โข ชโน ตชฺชํ ตสฺสารุปฺปํ กถํ มนฺเตติ, สํวิชฺชนฺเตว เต\footnote{สํวิชฺชนฺเต เต จ (สี, สฺยา, กํ, อิ)} ธมฺมา มยิ, อหญฺจ เตสุ ธมฺเมสุ สนฺทิสฺสามี”ติ.\csromanspacer{} อิทํ ภิกฺขเว พาโล ปฐมํ ทิฏฺเฐว ธมฺเม ทุกฺขํ โทมนสฺสํ ปฏิสํเวเทติ.}

\proseitem{247}{ปุน จปรํ ภิกฺขเว พาโล ปสฺสติ ราชาโน โจรํ อาคุจาริํ คเหตฺวา วิวิธา กมฺมการณา กาเรนฺเต กสาหิปิ ตาเฬนฺเต เวตฺเตหิปิ ตาเฬนฺเต อทฺธทณฺฑเกหิปิ ตาเฬนฺเต หตฺถมฺปิ ฉินฺทนฺเต ปาทมฺปิ ฉินฺทนฺเต หตฺถปาทมฺปิ ฉินฺทนฺเต กณฺณมฺปิ ฉินฺทนฺเต นาสมฺปิ ฉินฺทนฺเต กณฺณนาสมฺปิ ฉินฺทนฺเต พิลงฺคถาลิกมฺปิ กโรนฺเต สงฺขมุณฺฑิกมฺปิ กโรนฺเต ราหุมุขมฺปิ กโรนฺเต โชติมาลิกมฺปิ กโรนฺเต หตฺถปชฺโชติกมฺปิ กโรนฺเต เอรกวตฺติกมฺปิ กโรนฺเต จีรกวาสิกมฺปิ กโรนฺเต เอเณยฺยกมฺปิ กโรนฺเต พฬิสมํสิกมฺปิ กโรนฺเต กหาปณิกมฺปิ กโรนฺเต ขาราปตจฺฉิกมฺปิ\footnote{ขาราปฏิจฺฉกมฺปิ (ก)} กโรนฺเต ปลิฆปริวตฺติกมฺปิ กโรนฺเต ปลาลปีฐกมฺปิ\footnote{ปลาลปิฏฺฐกมฺปิ (อิ)} กโรนฺเต ตตฺเตนปิ เตเลน โอสิญฺจนฺเต สุนเขหิปิ ขาทาเปนฺเต ชีวนฺตมฺปิ สูเล อุตฺตาเสนฺเต อสินาปิ สีสํ ฉินฺทนฺเต.\csromanspacer{} ตตฺร ภิกฺขเว พาลสฺส เอวํ โหติ “ยถารูปานํ โข ปาปกานํ กมฺมานํ เหตุ ราชาโน โจรํ อาคุจาริํ คเหตฺวา วิวิธา กมฺมการณา กาเรนฺติ, กสาหิปิ ตาเฬนฺติ ฯเปฯ อสินาปิ สีสํ ฉินฺทนฺติ.\csromanspacer{} สํวิชฺชนฺเตว เต ธมฺมา มยิ, อหญฺจ เตสุ ธมฺเมสุ สนฺทิสฺสามิ, มํ เจปิ ราชาโน\footnote{สเจ มมฺปิ (ก)} ชาเนยฺยุํ, มมฺปิ ราชาโน}

\csromanheader{2}{9. พาลปณฺฑิตสุตฺต (129)}
\prose{\csromanfolio{203}คเหตฺวา วิวิธา กมฺมการณา กาเรยฺยุํ, กสาหิปิ ตาเฬยฺยุํ ฯเปฯ ชีวนฺตมฺปิ สูเล อุตฺตาเสยฺยุํ, อสินาปิ สีสํ ฉินฺเทยฺยุนฺ”ติ.\csromanspacer{} อิทมฺปิ ภิกฺขเว พาโล ทุติยํ ทิฏฺเฐว ธมฺเม ทุกฺขํ โทมนสฺสํ ปฏิสํเวเทติ.}

\proseitem{248}{ปุน จปรํ ภิกฺขเว พาลํ ปีฐสมารูฬฺหํ วา มญฺจสมารูฬฺหํ วา ฉมายํ\footnote{ฉมาย (สี, อิ)} วา เสมานํ ยานิสฺส ปุพฺเพ ปาปกานิ กมฺมานิ กตานิ กาเยน ทุจฺจริตานิ วาจาย ทุจฺจริตานิ มนสา ทุจฺจริตานิ, ตานิสฺส ตมฺหิ สมเย โอลมฺพนฺติ อชฺโฌลมฺพนฺติ อภิปฺปลมฺพนฺติ.\csromanspacer{} เสยฺยถาปิ ภิกฺขเว มหตํ ปพฺพตกูฏานํ ฉายา สายนฺหสมยํ ปถวิยา โอลมฺพนฺติ อชฺโฌลมฺพนฺติ อภิปฺปลมฺพนฺติ.\csromanspacer{} เอวเมว โข ภิกฺขเว พาลํ ปีฐสมารูฬฺหํ วา มญฺจสมารูฬฺหํ วา ฉมายํ วา เสมานํ ยานิสฺส ปุพฺเพ ปาปกานิ กมฺมานิ กตานิ กาเยน ทุจฺจริตานิ วาจาย ทุจฺจริตานิ มนสา ทุจฺจริตานิ, ตานิสฺส ตมฺหิ สมเย โอลมฺพนฺติ อชฺโฌลมฺพนฺติ อภิปฺปลมฺพนฺติ.\csromanspacer{} ตตฺร ภิกฺขเว พาลสฺส เอวํ โหติ “อกตํ วต เม กลฺยาณํ, อกตํ กุสลํ, อกตํ ภีรุตฺตาณํ.\csromanspacer{} กตํ ปาปํ, กตํ ลุทฺทํ, กตํ กิพฺพิสํ, ยาวตา โภ อกตลฺยาณานํ อกตกุสลานํ อกตภีรุตฺตาณานํ กตปาปานํ กตลุทฺทานํ กตกิพฺพิสานํ คติ, ตํ คติํ เปจฺจ คจฺฉามี”ติ.\csromanspacer{} โส โสจติ กิลมติ ปริเทวติ อุรตฺตาฬิํ กนฺทติ สมฺโมหํ อาปชฺชติ.\csromanspacer{} อิทมฺปิ ภิกฺขเว พาโล ตติยํ ทิฏฺเฐว ธมฺเม ทุกฺขํ โทมนสฺสํ ปฏิสํเวเทติ.}

\prose{ส โข โส ภิกฺขเว พาโล กาเยน ทุจฺจริตํ จริตฺวา วาจาย ทุจฺจริตํ จริตฺวา มนสา ทุจฺจริตํ จริตฺวา กายสฺส เภทา ปรํ มรณา อปายํ ทุคฺคติํ วินิปาตํ นิรยํ อุปปชฺชติ.\csromanspacer{} ยํ โข ตํ ภิกฺขเว สมฺมา วทมาโน วเทยฺย “เอกนฺตํ อนิฏฺฐํ เอกนฺตํ อกนฺตํ เอกนฺตํ อมนาปนฺ”ติ.\csromanspacer{} นิรยเมว ตํ สมฺมา วทมาโน วเทยฺย “เอกนฺตํ อนิฏฺฐํ เอกนฺตํ อกนฺตํ เอกนฺตํ อมนาปนฺ”ติ.\csromanspacer{} ยาวญฺจิทํ ภิกฺขเว อุปมาปิ\footnote{อุปมาหิปิ (สี)} น สุกรา, ยาว ทุกฺขา นิรยาติ.}

\proseitem{249}{เอวํ วุตฺเต อญฺญตโร ภิกฺขุ ภควนฺตํ เอตทโวจ “สกฺกา ปน ภนฺเต อุปมํ กาตุนฺ”ติ. “สกฺกา ภิกฺขู”ติ ภควา อโวจ, เสยฺยถาปิ ภิกฺขุ โจรํ อาคุจาริํ คเหตฺวา รญฺโญ ทสฺเสยฺยุํ “อยํ โข เทว โจโร อาคุจารี, อิมสฺส ยํ อิจฺฉสิ, ตํ ทณฺฑํ ปเณหี”ติ.\csromanspacer{} ตเมนํ ราชา เอวํ วเทยฺย “คจฺฉถ โภ อิมํ ปุริสํ ปุพฺพณฺหสมยํ สตฺติสเตน \csromanfolio{204}หนถา”ติ.\csromanspacer{} ตเมนํ ปุพฺพณฺหสมยํ สตฺติสเตน หเนยฺยุํ, อถ ราชา มชฺฌนฺหิกสมยํ\footnote{มชฺฌนฺติกสมยํ (สี, สฺยา, กํ, ก), มชฺฌนฺติกํ สมยํ (อิ)} เอวํ วเทยฺย “อมฺโภ กถํ โส ปุริโส”ติ.\csromanspacer{} ตเถว เทว ชีวตีติ.\csromanspacer{} ตเมนํ ราชา เอวํ วเทยฺย “คจฺฉถ โภ ตํ ปุริสํ มชฺฌนฺหิกสมยํ สตฺติสเตน หนถา”ติ.\csromanspacer{} ตเมนํ มชฺฌนฺหิกสมยํ สตฺติสเตน หเนยฺยุํ.\csromanspacer{} อถ ราชา สายนฺหสมยํ เอวํ วเทยฺย “อมฺโภ กถํ โส ปุริโส”ติ.\csromanspacer{} ตเถว เทว ชีวตีติ.\csromanspacer{} ตเมนํ ราชา เอวํ วเทยฺย “คจฺฉถ โภ ตํ ปุริสํ สายนฺหสมยํ สตฺติสเตน หนถา”ติ.\csromanspacer{} ตเมนํ สายนฺหสมยํ สตฺติสเตน หเนยฺยุํ.\csromanspacer{} ตํ กิํ มญฺญถ ภิกฺขเว, อปิ นุ โส ปุริโส ตีหิ สตฺติสเตหิ หญฺญมาโน ตโตนิทานํ ทุกฺขํ โทมนสฺสํ ปฏิสํเวทิเยถาติ.\csromanspacer{} เอกิสฺสาปิ ภนฺเต สตฺติยา หญฺญมาโน โส ปุริโส ตโตนิทานํ ทุกฺขํ โทมนสฺสํ ปฏิสํเวทิเยถ, โก ปน วาโท ตีติ สตฺติสเตหีติ.}

\proseitem{250}{อถ โข ภควา ปริตฺตํ ปาณิมตฺตํ ปาสาณํ คเหตฺวา ภิกฺขู อามนฺเตสิ “ตํ กิํ มญฺญถ ภิกฺขเว, กตโม นุ โข มหนฺตตโร, โย จายํ มยา ปริตฺโต ปาณิมตฺโต ปาสาโณ คหิโต, โย จ หิมวา ปพฺพตราชา”ติ.\csromanspacer{} อปฺปมตฺตโก อยํ ภนฺเต ภควตา ปริตฺโต ปาณิมตฺโต ปาสาโณ คหิโต, หิมวนฺตํ ปพฺพตราชานํ อุปนิธาย สงฺขมฺปิ น อุเปติ, กลภาคมฺปิ น อุเปติ, อุปนิธมฺปิ\footnote{อุปนิธิมฺปิ (สี,อิ)} น อุเปติ.\csromanspacer{} เอวเมว โข ภิกฺขเว ยํ โส ปุริโส ตีหิ สตฺติสเตหิ หญฺญมาโน ตโตนิทานํ ทุกฺขํ โทมนสฺสํ ปฏิสํเวเทติ, ตํ นิรยกสฺส ทุกฺขสฺส อุปนิธาย สงฺขมฺปิ น อุเปติ, กลภาคมฺปิ น อุเปติ, อุปนิธมฺปิ น อุเปติ.}

\prose{ตเมนํ ภิกฺขเว นิรยปาลา ปญฺจวิธพนฺธนํ นาม กมฺมการณํ กโรนฺติ, ตตฺตํ อโยขิลํ\footnote{อโยขีลํ (สี, สฺยา, กํ, อิ)} หตฺเถ คเมนฺติ, ตตฺตํ อโยขิลํ ทุติเย หตฺเถ คเมนฺติ, ตตฺตํ อโยขิลํ ปาเท คเมนฺติ, ตตฺตํ อโยขิลํ ทุติเย ปาเท คเมนฺติ, ตตฺตํ อโยขิลํ มชฺเฌ อุรสฺมิํ คเมนฺติ.\csromanspacer{} โส ตตฺถ ทุกฺขา ติพฺพา ขารา กฏุกา เวทนา เวเทติ.\csromanspacer{} น จ ตาว กาลํ กโรติ, ยาว น ตํ ปาปกมฺมํ พฺยนฺตีโหติ\footnote{พฺยนฺติโหติ (อิ, ก)}.\csromanspacer{} ตเมนํ ภิกฺขเว นิรยปาลา สํเวเสตฺวา กุฐารีหิ\footnote{กุธารีหิ (ก)} ตจฺฉนฺติ.\csromanspacer{} โส ตตฺถ ทุกฺขา ติพฺพา ฯเปฯ พฺยนฺตีโหติ.\csromanspacer{} ตเมนํ ภิกฺขเว นิรยปาลา อุทฺธํปาทํ อโธสิรํ คเหตฺวา วาสีหิ ตจฺฉนฺติ.}

\csromanheader{2}{9. พาลปณฺฑิตสุตฺต (129)}
\prose{\csromanfolio{205}โส ตตฺถ ทุกฺขา ติพฺพา ฯเปฯ พฺยนฺตีโหติ.\csromanspacer{} ตเมนํ ภิกฺขเว นิรยปาลา รเถ โยเชตฺวา อาทิตฺตาย ปถวิยา สมฺปชฺชลิตาย สโชติภูตาย\footnote{สญฺโชติภูตาย (สฺยา, กํ, อิ)} สาเรนฺติปิ ปจฺจาสาเรนฺติปิ.\csromanspacer{} โส ตตฺถ ทุกฺขา ติพฺพา ฯเปฯ พฺยนฺตีโหติ.\csromanspacer{} ตเมนํ ภิกฺขเว นิรยปาลา มหนฺตํ องฺคารปพฺพตํ อาทิตฺตํ สมฺปชฺชลิตํ สโชติภูตํ อาโรเปนฺติปิ โอโรเปนฺติปิ.\csromanspacer{} โส ตตฺถ ทุกฺขา ติพฺพา ขรา กฏุกา เวทนา เวเทติ.\csromanspacer{} น จ ตาว กาลํ กโรติ, ยาว น ตํ ปาปกมฺมํ พฺยนฺตีโหติ.\csromanspacer{} ตเมนํ ภิกฺขเว นิรยปาลา อุทฺธํปาทํ อโธสิรํ คเหตฺวา ตตฺตาย โลหกุมฺภิยา ปกฺขิปนฺติ อาทิตฺตาย สมฺปชฺชลิตาย สโชติภูตาย.\csromanspacer{} โส ตตฺถ เผณุทฺเทหกํ ปจฺจติ, โส ตตฺถ เผณุทฺเทหกํ ปจฺจมาโน สกิมฺปิ อุทฺธํ คจฺฉติ, สกิมฺปิ อโธ คจฺฉติ, สกิมฺปิ ติริยํ คจฺฉติ.\csromanspacer{} โส ตตฺถ ทุกฺขา ติพฺพา ขรา กฏุกา เวทนา เวเทติ.\csromanspacer{} น จ ตาว กาลํ กโรติ, ยาว น ตํ ปาปกมฺมํ พฺยนฺตีโหติ.\csromanspacer{} ตเมนํ ภิกฺขเว นิรยปาลา\footnote{นิรยปาลา ปุนปฺปุนํ (ก)} มหานิรเย ปกฺขิปนฺติ.\csromanspacer{} โส โข ปน ภิกฺขเว มหานิรโย\csromandash{}}

\csromangathameasure{จตุกฺกณฺโณ จตุทฺวาโร, วิภตฺโต ภาคโส มิโต. \\ อโยปาการปริยนฺโต, อยสา ปฏิกุชฺชิโต. \\ ตสฺส อโยมยา ภูมิ, ชลิตา เตชสา ยุตา. \\ สมนฺตา โยชนสตํ, ผริตฺวา ติฏฺฐติ สพฺพทา.}
\csromangathagroup{\csromangathastanza{จตุกฺกณฺโณ จตุทฺวาโร, วิภตฺโต ภาคโส มิโต. \\ อโยปาการปริยนฺโต, อยสา ปฏิกุชฺชิโต.}\csromangathastanza{ตสฺส อโยมยา ภูมิ, ชลิตา เตชสา ยุตา. \\ สมนฺตา โยชนสตํ, ผริตฺวา ติฏฺฐติ สพฺพทา.}}

\prose{อเนกปริยาเยนปิ โข อหํ ภิกฺขเว นิรยกถํ กเถยฺยํ, ยาวญฺจิทํ ภิกฺขเว น สุกรา อกฺขาเนน ปาปุณิตุํ ยาว ทุกฺขา นิรยา.}

\proseitem{251}{สนฺติ ภิกฺขเว ติรจฺฉานคตา ปาณา ติณภกฺขา, เต อลฺลานิปิ ติณานิ สุกฺขานิปิ ติณานิ ทนฺตุลฺเลหกํ ขาทนฺติ.\csromanspacer{} กตเม จ ภิกฺขเว ติรจฺฉานคตา ปาณา ติณภกฺขา, หตฺถี อสฺสา โคณา คทฺรภา อชา มิคา เย วา ปนญฺเญปิ เกจิ ติรจฺฉานคตา ปาณา ติณภกฺขา.\csromanspacer{} ส โข โส ภิกฺขเว พาโล อิธ ปุพฺเพ รสาโท อิธ ปาปานิ กมฺมานิ กริตฺวา กายสฺส เภทา ปรํ มรณา เตสํ สตฺตานํ สหพฺยตํ อุปปชฺชติ, เย เต สตฺตา ติณภกฺขา.}

\prose{สนฺติ ภิกฺขเว ติรจฺฉานคตา ปาณา คูถภกฺขา, เต ทูรโตว คูถคนฺธํ ฆายิตฺวา ธาวนฺติ “เอตฺถ ภุญฺชิสฺสาม เอตฺถ ภุญฺชิสฺสามา”ติ. \csromanfolio{206}เสยฺยถาปิ นาม พฺราหฺมณา อาหุติคนฺเธน ธาวนฺติ “เอตฺถ ภุญฺชิสฺสาม เอตฺถ ภุญฺชิสฺสามา”ติ.\csromanspacer{} เอวเมว โข ภิกฺขเว สนฺติ ติรจฺฉานคตา ปาณา คูถภกฺขา, เต ทูรโตว คูถคนฺธํ ฆายิตฺวา ธาวนฺติ “เอตฺถ ภุญฺชิสฺสาม เอตฺถ ภุญฺชิสฺสามา”ติ.\csromanspacer{} กตเม จ ภิกฺขเว ติรจฺฉานคตา ปาณา คูถภกฺขา, กุกฺกุฏา สูกรา โสณา สิงฺคาลา เย วา ปนญฺเญปิ เกจิ ติรจฺฉานคตา ปาณา คูถภกฺขา.\csromanspacer{} ส โข โส ภิกฺขเว พาโล อิธ ปุพฺเพ รสาโท อิธ ปาปานิ กมฺมานิ กริตฺวา กายสฺส เภทา ปรํ มรณา เตสํ สตฺตานํ สหพฺยตํ อุปปชฺชติ, เย เต สตฺตา คูถภกฺขา.}

\prose{สนฺติ ภิกฺขเว ติรจฺฉานคตา ปาณา อนฺธกาเร ชายนฺติ, อนฺธกาเร ชียนฺติ\footnote{ชิยฺยนฺติ (ก)}, อนฺธกาเร มียนฺติ\footnote{มิยฺยนฺติ (ก)}.\csromanspacer{} กตเม จ ภิกฺขเว ติรจฺฉานคตา ปาณา อนฺธกาเร ชายนฺติ, อนฺธกาเร ชียนฺติ, อนฺธกาเร มียนฺติ.\csromanspacer{} กีฏา ปุฬวา\footnote{ปฏงฺคา (สฺยา, กํ, ก)} คณฺฑุปฺปาทา เย วา ปนญฺเญปิ เกจิ ติรจฺฉานคตา ปาณา อนฺธกาเร ชายนฺติ, อนฺธกาเร ชียนฺติ, อนฺธกาเร มียนฺติ.\csromanspacer{} ส โข โส ภิกฺขเว พาโล อิธ ปุพฺเพ รสาโท อิธ ปาปานิ กมฺมานิ กริตฺวา กายสฺส เภทา ปรํ มรณา เตสํ สตฺตานํ สหพฺยตํ อุปปชฺชติ, เย เต สตฺตา อนฺธกาเร ชายนฺติ, อนฺธกาเร ชียนฺติ, อนฺธกาเร มียนฺติ.}

\prose{สนฺติ ภิกฺขเว ติรจฺฉานคตา ปาณา อุทกสฺมิํ ชายนฺติ, อุทกสฺมิํ ชียนฺติ, อุทกสฺมิํ มียนฺติ.\csromanspacer{} กตเม จ ภิกฺขเว ติรจฺฉานคตา ปาณา อุทกสฺมิํ ชายนฺติ, อุทกสฺมิํ ชียนฺติ, อุทกสฺมิํ มียนฺติ.\csromanspacer{} มจฺฉา กจฺฉปา สุสุมารา เย วา ปนญฺเญปิ เกจิ ติรจฺฉานคตา ปาณา อุทกสฺมิํ ชายนฺติ, อุทกสฺมิํ ชียนฺติ, อุทกสฺมิํ มียนฺติ.\csromanspacer{} ส โข โส ภิกฺขเว พาโล อิธ ปุพฺเพ รสาโท อิธ ปาปานิ กมฺมานิ กริตฺวา กายสฺส เภทา ปรํ มรณา เตสํ สตฺตานํ สหพฺยตํ อุปปชฺชติ, เย เต สตฺตา อุทกสฺมิํ ชายนฺติ, อุทกสฺมิํ ชียนฺติ, อุทกสฺมิํ มียนฺติ.}

\prose{สนฺติ ภิกฺขเว ติรจฺฉานคตา ปาณา อสุจิสฺมิํ ชายนฺติ, อสุจิสฺมิํ ชียนฺติ, อสุจิสฺมิํ มียนฺติ.\csromanspacer{} กตเม จ ภิกฺขเว ติรจฺฉานคตา ปาณา อสุจิสฺมิํ ชายนฺติ, อสุจิสฺมิํ ชียนฺติ, อสุจิสฺมิํ มียนฺติ.\csromanspacer{} เย เต ภิกฺขเว สตฺตา ปูติมจฺเฉ วา ชายนฺติ, ปูติมจฺเฉ วา ชียนฺติ, ปูติมจฺเฉ วา มียนฺติ.\csromanspacer{} ปูติกุณเป วา ฯเปฯ.\csromanspacer{} ปูติกุมฺมาเส วา.\csromanspacer{} จนฺทนิกาย วา.}

\csromanheader{2}{9. พาลปณฺฑิตสุตฺต (129)}
\prose{\csromanfolio{207}โอลิคลฺเล วา ชายนฺติ. (เย วา ปนญฺเญปิ เกจิ ติรจฺฉานคตา ปาณา อสุจิสฺมิํ ชายนฺติ, อสุจิสฺมิํ ชียนฺติ, อสุจิสฺมิํ มียนฺติ.)\footnote{( ) นตฺถิ สี-สฺยา-กํ-อิ-โปตฺถเกสุ. 2. วสฺสสตสฺส วสฺสสหสฺสสฺส วสฺสสตสหสฺสสฺส (สี), วสฺสสตสฺส (สฺยา, กํ, อิ) 3. ( ) นตฺถิ สี-อิ-โปตฺถเกสุ.} ส โข โส ภิกฺขเว พาโล อิธ ปุพฺเพ รสาโท อิธ ปาปานิ กมฺมานิ กริตฺวา กายสฺส เภทา ปรํ มรณา เตสํ สตฺตานํ สหพฺยตํ อุปปชฺชติ, เย เต สตฺตา อสุจิสฺมิํ ชายนฺติ, อสุจิสฺมิํ ชียนฺติ, อสุจิสฺมิํ มียนฺติ.}

\prose{อเนกปริยาเยนปิ โข อหํ ภิกฺขเว ติรจฺฉานโยนิกถํ กเถยฺยํ, ยาวญฺจิทํ ภิกฺขเว น สุกรํ อกฺขาเนน ปาปุณิตุํ, ยาว ทุกฺขา ติรจฺฉานโยนิ.}

\proseitem{252}{เสยฺยถาปิ ภิกฺขเว ปุริโส เอกจฺฉิคฺคลํ ยุคํ มหาสมุทฺเท ปกฺขิเปยฺย, ตเมนํ ปุรตฺถิโม วาโต ปจฺฉิเมน สํหเรยฺย, ปจฺฉิโม วาโต ปุรตฺถิเมน สํหเรยฺย, อุตฺตโร วาโต ทกฺขิเณน สํหเรยฺย, ทกฺขิโณ วาโต อุตฺตเรนํ สํหเรยฺย.\csromanspacer{} ตตฺราสฺส กาโณ กจฺฉโป, โส วสฺสสตสฺส วสฺสสตสฺส2 อจฺจเยน สกิํ อุมฺมุชฺเชยฺย.\csromanspacer{} ตํ กิํ มญฺญถ ภิกฺขเว, อปิ นุ โส กาโณ กจฺฉโป อมุสฺมิํ เอกจฺฉิคฺคเล ยุเค คีวํ ปเวเสยฺยาติ. (โน เหตํ ภนฺเต.)3 ยทิ ปน\footnote{ยทิ นูน (สี, สฺยา, กํ, อิ)} ภนฺเต กทาจิ กรหจิ ทีฆสฺส อทฺธุโน อจฺจเยนาติ.\csromanspacer{} ขิปฺปตรํ โข โส ภิกฺขเว กาโณ กจฺฉโป อมุสฺมิํ เอกจฺฉิคฺคเล ยุเค คีวํ ปเวเสยฺย.\csromanspacer{} อโต ทุลฺลภตราหํ ภิกฺขเว มนุสฺสตฺตํ วทามิ สกิํ วินิปาตคเตน พาเลน.\csromanspacer{} ตํ กิสฺส เหตุ, น เหตฺถ ภิกฺขเว อตฺถิ ธมฺมจริยา สมจริยา กุสลกิริยา ปุญฺญกิริยา, อญฺญมญฺญขาทิกา เอตฺถ ภิกฺขเว วตฺตติ ทุพฺพลขาทิกา.}

\prose{ส โข โส ภิกฺขเว พาโล สเจ กทาจิ กรหจิ ทีฆสฺส อทฺธุโน อจฺจเยน มนุสฺสตฺตํ อาคจฺฉติ, ยานิ ตานิ นีจกุลานิ จณฺฑาลกุลํ วา เนสาทกุลํ วา เวนกุลํ\footnote{เวณกุลํ (สี, อิ)} วา รถการกุลํ วา ปุกฺกุสกุลํ วา, ตถารูเป กุเล ปจฺจาชายติ ทลิทฺเท อปฺปนฺนปานโภชเน กสิรวุตฺติเก, ยตฺถ กสิเรน ฆาสจฺฉาโท ลพฺภติ, โส จ โหติ ทุพฺพณฺโณ \csromanfolio{208}ทุทฺทสิโก โอโกฏิมโก พวฺหาพาโธ\footnote{พหฺวาพาโธ (ก)} กาโณ วา กุณี วา ขุชฺโช วา ปกฺขหโต วา น ลาภี อนฺนสฺส ปานสฺส วตฺถสฺส ยานสฺส มาลาคนฺธวิเลปนสฺส เสยฺยาวสถปทีเปยฺยสฺส.\csromanspacer{} โส กาเยน ทุจฺจริตํ จรติ, วาจาย ทุจฺจริตํ จรติ, มนสา ทุจฺจริตํ จรติ, โส กาเยน ทุจฺจริตํ จริตฺวา วาจาย ทุจฺจริตํ จริตฺวา มนสา ทุจฺจริตํ จริตฺวา กายสฺส เภทา ปรํ มรณา อปายํ ทุคฺคติํ วินิปาตํ นิรยํ อุปปชฺชติ.}

\prose{เสยฺยถาปิ ภิกฺขเว อกฺขธุตฺโต ปฐเมเนว กลิคฺคเหน ปุตฺตมฺปิ ชีเยถ, ทารมฺปิ ชีเยถ, สพฺพํ สาปเตยฺยมฺปิ ชีเยถ, อุตฺตริปิ อธิพนฺธํ\footnote{อนุพนฺธํ (สี, อิ), อทฺธุพนฺธํ (สฺยา, กํ)} นิคจฺเฉยฺย.\csromanspacer{} อปฺปมตฺตโก โส ภิกฺขเว กลิคฺคโห ยํ โส อกฺขธุตฺโต ปฐเมเนว กลิคฺคเหน ปุตฺตมฺปิ ชีเยถ, ทารมฺปิ ชีเยถ, สพฺพํ สาปเตยฺยมฺปิ ชีเยถ, อุตฺตริปิ อธิพนฺธํ นิคจฺเฉยฺย, อถ โข อยเมว ตโต มหนฺตตโร กลิคฺคโห, ยํ โส พาโล กาเยน ทุจฺจริตํ จริตฺวา วาจาย ทุจฺจริตํ จริตฺวา มนสา ทุจฺจริตํ จริตฺวา กายสฺส เภทา ปรํ มรณา อปายํ ทุคฺคติํ วินิปาตํ นิรยํ อุปปชฺชติ.\csromanspacer{} อยํ ภิกฺขเว เกวลา ปริปูรา\footnote{เกวลปริปูรา (สี, อิ) ม 1. 191 ปิฏฺเฐ ปาฬิยา สํสนฺเทตพฺพา.} พาลภูมีติ.}

\proseitem{253}{ตีณิมานิ ภิกฺขเว ปณฺฑิตสฺส ปณฺฑิตลกฺขณานิ ปณฺฑิตนิมิตฺตานิ ปณฺฑิตาปทานานิ.\csromanspacer{} กตมานิ ตีณิ, อิธ ภิกฺขเว ปณฺฑิโต สุจินฺติตจินฺตี จ โหติ สุภาสิตภาสี จ สุกตกมฺมการี จ.\csromanspacer{} โน เจตํ ภิกฺขเว ปณฺฑิโต สุจินฺติตจินฺตี จ อภวิสฺส สุภาสิตภาสี จ สุกตกมฺมการี จ, เกน นํ\footnote{น เตน นํ (ก), น นํ (?)} ปณฺฑิตา ชาเนยฺยุํ “ปณฺฑิโต อยํ ภวํ สปฺปุริโส”ติ.\csromanspacer{} ยสฺมา จ โข ภิกฺขเว ปณฺฑิโต สุจินฺติตจินฺตี จ โหติ สุภาสิตภาสี จ สุกตกมฺมการี จ, ตสฺมา นํ ปณฺฑิตา ชานนฺติ “ปณฺฑิโต อยํ ภวํ สปฺปุริโส”ติ.\csromanspacer{} ส โข โส ภิกฺขเว ปณฺฑิโต ติวิธํ ทิฏฺเฐว ธมฺเม สุขํ โสมนสฺสํ ปฏิสํเวเทติ.\csromanspacer{} สเจ ภิกฺขเว ปณฺฑิโต สภายํ วา นิสินฺโน โหติ รถิกาย วา นิสินฺโน โหติ, สิงฺฆาฏเก วา นิสินฺโน โหติ, ตตฺร เจ ชโน ตชฺชํ ตสฺสารุปฺปํ กถํ}

\csromanheader{2}{9. พาลปณฺฑิตสุตฺต (129)}
\prose{\csromanfolio{209}มนฺเตติ.\csromanspacer{} สเจ ภิกฺขเว ปณฺฑิโต ปาณาติปาตา ปฏิวิรโต โหติ, อทินฺนาทานา ปฏิวิรโต โหติ, กาเมสุมิจฺฉาจารา ปฏิวิรโต โหติ, มุสาวาทา ปฏิวิรโต โหติ, สุราเมรยมชฺชปฺปมาทฏฺฐานา ปฏิวิรโต โหติ.\csromanspacer{} ตตฺร ภิกฺขเว ปณฺฑิตสฺส เอวํ โหติ “ยํ โข ชโน ตชฺชํ ตสฺสารุปฺปํ กถํ มนฺเตติ.\csromanspacer{} สํวิชฺชนฺเตว เต ธมฺมา มยิ, อหญฺจ เตสุ ธมฺเมสุ สนฺทิสฺสามี”ติ.\csromanspacer{} อิทํ ภิกฺขเว ปณฺฑิโต ปฐมํ ทิฏฺเฐว ธมฺเม สุขํ โสมนสฺสํ ปฏิสํเวเทติ.}

\proseitem{254}{ปุน จปรํ ภิกฺขเว ปณฺฑิโต ปสฺสติ ราชาโน โจรํ อาคุจาริํ คเหตฺวา วิวิธา กมฺมการณา กาเรนฺเต กสาหิปิ ตาเฬนฺเต เวตฺเตหิปิ ตาเฬนฺเต อทฺธทณฺฑเกหิปิ ตาเฬนฺเต หตฺถมฺปิ ฉินฺเทนฺเต ปาทมฺปิ ฉินฺเทนฺเต หตฺถปาทมฺปิ ฉินฺทนฺเต กณฺณมฺปิ ฉินฺทนฺเต นาสมฺปิ ฉินฺทนฺเต กณฺณนาสมฺปิ ฉินฺทนฺเต พิลงฺคถาลิกมฺปิ กโรนฺเต สงฺขมุณฺฑิกมฺปิ กโรนฺเต ราหุมุขมฺปิ กโรนฺเต โชติมาลิกมฺปิ กโรนฺเต หตฺถปชฺโชติกมฺปิ กโรนฺเต เอรกวตฺติกมฺปิ กโรนฺเต จีรกวาสิกมฺปิ กโรนฺเต เอเณยฺยกมฺปิ กโรนฺเต พลิสมํสิกมฺปิ กโรนฺเต กหาปณิกมฺปิ กโรนฺเต ขาราปตจฺฉิกมฺปิ กโรนฺเต ปลิฆปริวตฺติกมฺปิ กโรนฺเต ปลาลปีฐกมฺปิ กโรนฺเต ตตฺเตนปิ เตเลน โอสิญฺจนฺเต สุนเขหิปิ ขาทาเปนฺเต ชีวนฺตมฺปิ สูเล อุตฺตาเสนฺเต อสินาปิ สีสํ ฉินฺทนฺเต.\csromanspacer{} ตตฺร ภิกฺขเว ปณฺฑิตสฺส เอวํ โหติ “ยถารูปานํ โข ปาปกานํ กมฺมานํ เหตุ ราชาโน โจรํ อาคุจาริํ คเหตฺวา วิวิธา กมฺมการณา กาเรนฺติ, กสาหิปิ ตาเฬนฺติ, เวตฺเตหิปิ ตาเฬนฺติ, อทฺธทณฺฑเกหิปิ ตาเฬนฺติ, หตฺถมฺปิ ฉินฺทนฺติ, ปาทมฺปิ ฉินฺทนฺติ, หตฺถปาทมฺปิ ฉินฺทนฺติ, กณฺณมฺปิ ฉินฺทนฺติ, นาสมฺปิ ฉินฺทนฺติ, กณฺณนาสมฺปิ ฉินฺทนฺติ, พิลงฺคถาลิกมฺปิ กโรนฺติ, สงฺขมุณฺฑิกมฺปิ กโรนฺติ, ราหุมุขมฺปิ กโรนฺติ, โชติมาลิกมฺปิ กโรนฺติ, หตฺถปชฺโชติกมฺปิ กโรนฺติ, เอรกวตฺติกมฺปิ กโรนฺติ, จีรกวาสิกมฺปิ กโรนฺติ, เอเณยฺยกมฺปิ กโรนฺติ, พลิสมํสิกมฺปิ กโรนฺติ, กหาปณิกมฺปิ กโรนฺติ, ขาราปตจฺฉิกมฺปิ กโรนฺติ, ปลิฆปริวตฺติกมฺปิ กโรนฺติ, ปลาลปีฐกมฺปิ กโรนฺติ, ตตฺเตนปิ เตเลน โอสิญฺจนฺติ, สุนเขหิปิ ขาทาเปนฺติ, ชีวนฺตมฺปิ สูเล อุตฺตาเสนฺติ, อสินาปิ สีสํ ฉินฺทนฺติ, น เต ธมฺมา มยิ สํวิชฺชนฺติ, อหญฺจ น เตสุ ธมฺเมสุ สนฺทิสฺสามี”ติ.\csromanspacer{} อิทมฺปิ ภิกฺขเว ปณฺฑิโต ทุติยํ ทิฏฺเฐว ธมฺเม สุขํ โสมนสฺสํ ปฏิสํเวเทติ.}

\proseitem{225}{\csromanfolio{210}ปุน จปรํ ภิกฺขเว ปณฺฑิตํ ปีฐสมารูฬฺหํ วา มญฺจสมารูฬฺหํ วา ฉมายํ วา เสมานํ ยานิสฺส ปุพฺเพ กลฺยาณานิ กมฺมานิ กตานิ กาเยน สุจริตานิ วาจาย สุจริตานิ มนสา สุจริตานิ, ตานิสฺส ตมฺหิ สมเย โอลมฺพนฺติ ฯเปฯ.\csromanspacer{} เสยฺยถาปิ ภิกฺขเว มหตํ ปพฺพตกูฏานํ ฉายา สายนฺหสมยํ ปถวิยา โอลมฺพนฺติ อชฺโฌลมฺพนฺติ อภิปฺปลมฺพนฺติ.\csromanspacer{} เอวเมว โข ภิกฺขเว ปณฺฑิตํ ปีฐสมารูฬฺหํ วา มญฺจสมารูฬํ วา ฉมายํ วา เสมานํ ยานิสฺส ปุพฺเพ กลฺยาณานิ กมฺมานิ กตานิ กาเยน สุจริตานิ วาจาย สุจริตานิ มนสา สุจริตานิ, ตานิสฺส ตมฺหิ สมเย โอลมฺพนฺติ อชฺโฌลมฺพนฺติ อภิปฺปลมฺพนฺติ.\csromanspacer{} ตตฺร ภิกฺขเว ปณฺฑิตสฺส เอวํ โหติ “อกตํ วต เม ปาปํ, อกตํ ลุทฺทํ, อกตํ กิพฺพิสํ, กตํ กลฺยาณํ, กตํ กุสลํ, กตํ ภีรุตฺตาณํ, ยาวตา โภ อกตปาปานํ อกตลุทฺทานํ อกตกิพฺพิสานํ กตกลฺยาณานํ กตกุสลานํ กตภีรุตฺตาณานํ คติ, ตํ คติํ เปจฺจ คจฺฉามี”ติ.\csromanspacer{} โส น โสจติ น กิลมติ น ปริเทวติ น อุรตฺตาฬิํ กนฺทติ น สมฺโมหํ อาปชฺชติ.\csromanspacer{} อิทมฺปิ ภิกฺขเว ปณฺฑิโต ตติยํ ทิฏฺเฐว ธมฺเม สุขํ โสมนสฺสํ ปฏิสํเวเทติ.}

\prose{ส โข โส ภิกฺขเว ปณฺฑิโต กาเยน สุจริตํ จริตฺวา วาจาย สุจริตํ จริตฺวา มนสา สุจริตํ จริตฺวา กายสฺส เภทา ปรํ มรณา สุคติํ สคฺคํ โลกํ อุปปชฺชติ.\csromanspacer{} ยํ โข ตํ ภิกฺขเว สมฺมา วทมาโน วเทยฺย “เอกนฺตํ อิฏฺฐํ เอกนฺตํ กนฺตํ เอกนฺตํ มนาปนฺ”ติ.\csromanspacer{} สคฺคเมว ตํ สมฺมา วทมาโน วเทยฺย “เอกนฺตํ อิฏฺฐํ เอกนฺตํ กนฺตํ เอกนฺตํ มนาปนฺ”ติ.\csromanspacer{} ยาวญฺจิทํ ภิกฺขเว อุปมาปิ น สุกรา, ยาว สุขา สคฺคาติ.}

\proseitem{256}{เอวํ วุตฺเต อญฺญตโร ภิกฺขุ ภควนฺตํ เอตทโวจ “สกฺกา ปน ภนฺเต อุปมํ กาตุนฺ”ติ. “สกฺกา ภิกฺขู”ติ ภควา อโวจ, เสยฺยถาปิ ภิกฺขเว ราชา จกฺกวตฺตี สตฺตหิ รตเนหิ สมนฺนาคโต จตูหิ จ อิทฺธีหิ, ตโตนิทานํ สุขํ โสมนสฺสํ ปฏิสํเวเทติ.\csromanspacer{} กตเมหิ สตฺตหิ, อิธ ภิกฺขเว รญฺโญ ขตฺติยสฺส มุทฺธาวสิตฺตสฺส ตทหุโปสเถ ปนฺนรเส สีสํนฺหาตสฺส อุโปสถิกสฺส อุปริปาสาทวรคตสฺส ทิพฺพํ จกฺกรตนํ ปาตุภวติ สหสฺสารํ สเนมิกํ สนาภิกํ สพฺพาการปริปูรํ, ตํ ทิสฺวาน รญฺโญ ขตฺติยสฺส มุทฺธาวสิตฺตสฺส เอวํ โหติ\footnote{เอตทโหสิ (สฺยา, กํ, ก)} “สุตํ โข ปน}

\csromanheader{2}{9. พาลปณฺฑิตสุตฺต (129)}
\prose{\csromanfolio{211}เมตํ ‘ยสฺส รญฺโญ ขตฺติยสฺส มุทฺธาวสิตฺตสฺส ตทหุโปสเถ ปนฺนรเส สีสํนฺหาตสฺส อุโปสถิกสฺส อุปริปาสาทวรคตสฺส ทิพฺพํ จกฺกรตนํ ปาตุภวติ สหสฺสารํ สเนมิกํ สนาภิกํ สพฺพาการปริปูรํ, โส โหติ ราชา จกฺกวตฺตี’ติ, อสฺสํ นุ โข อหํ ราชา จกฺกวตฺตี”ติ.}

\prose{อถ โข ภิกฺขเว ราชา ขตฺติโย มุทฺธาวสิตฺโต วาเมน หตฺเถน ภิงฺการํ คเหตฺวา ทกฺขิเณน หตฺเถน จกฺกรตนํ อพฺภุกฺกิรติ “ปวตฺตตุ ภวํ จกฺกรตนํ, อภิวิชินาตุ ภวํ จกฺกรตนนฺ”ติ.\csromanspacer{} อถ โข ตํ ภิกฺขเว จกฺกรตนํ ปุรตฺถิมํ ทิสํ ปวตฺตติ อนฺวเทว ราชา จกฺกวตฺตี สทฺธิํ จตุรงฺคินิยา เสนาย.\csromanspacer{} ยสฺมิํ โข ปน ภิกฺขเว ปเทเส จกฺกรตนํ ปติฏฺฐาติ, ตตฺถ ราชา จกฺกวตฺตี วาสํ อุเปติ สทฺธิํ จตุรงฺคินิยา เสนาย.\csromanspacer{} เย โข ปน ภิกฺขเว ปุรตฺถิมาย ทิสาย ปฏิราชาโน, เต ราชานํ จกฺกวตฺติํ อุปสงฺกมิตฺวา เอวมาหํสุ “เอหิ โข มหาราช, สฺวาคตํ เต มหารช\footnote{สฺวาคตํ มหาราช (สี, สฺยา, กํ, อิ)}, สกํ เต มหาราช, อนุสาส มหาราชา”ติ.\csromanspacer{} ราชา จกฺกวตฺตี เอวมาห “ปาโณ น หนฺตพฺโพ, อทินฺนํ นาทาตพฺพํ, กาเมสุมิจฺฉา น จริตพฺพา, มุสา น ภาสิตพฺพา, มชฺชํ น ปาตพฺพํ, ยถาภุตฺตญฺจ ภุญฺชถา”ติ.\csromanspacer{} เย โข ปน ภิกฺขเว ปุรตฺถิมาย ทิสาย ปฏิราชาโน, เต รญฺโญ จกฺกวตฺติสฺส อนุยนฺตา\footnote{อนุยุตฺตา (สี, สฺยา, กํ, อิ)} ภวนฺติ\footnote{อเหสุํ (สฺยา, กํ, ก)}.}

\proseitem{257}{อถ โข ตํ ภิกฺขเว จกฺกรตนํ ปุรตฺถิมํ สมุทฺทํ อชฺโฌคาเหตฺวา\footnote{อชฺโฌคเหตฺวา (สี, สฺยา, กํ, อิ)} ปจฺจุตฺตริตฺวา ทกฺขิณํ ทิสํ ปวตฺตติ ฯเปฯ ทกฺขิณํ สมุทฺทํ อชฺโฌคาเหตฺวา ปจฺจุตฺตริตฺวา ปจฺฉิมํ ทิสํ ปวตฺตติ ฯเปฯ ปจฺฉิมํ สมุทฺทํ อชฺโฌคาเหตฺวา ปจฺจุตฺตริตฺวา อุตฺตรํ ทิสํ ปวตฺตติ อนฺวเทว ราชา จกฺกวตฺตี สทฺธิํ จตุรงฺคินิยา เสนาย.\csromanspacer{} ยสฺมิํ โข ปน ภิกฺขเว ปเทเส จกฺกรตนํ ปติฏฺฐาติ, ตตฺถ ราชา จกฺกวตฺตี วาสํ อุเปติ สทฺธิํ จตุรงฺคินิยา เสนาย.}

\prose{เย โข ปน ภิกฺขเว อุตฺตราย ทิสาย ปฏิราชาโน, เต ราชานํ จกฺกวตฺติํ อุปสงฺกมิตฺวา เอวมาหํสุ “เอหิ โข มหาราช, สฺวาคตํ เต มหาราช, สกํ เต มหาราช, อนุสาส มหาราชา”ติ.\csromanspacer{} ราชา จกฺกวตฺตี เอวมาห “ปาโณ น หนฺตพฺโพ, อทินฺนํ นาทาตพฺพํ, กาเมสุมิจฺฉา \csromanfolio{212}น จริตพฺพา, มุสา น ภาสิตพฺพา, มชฺชํ น ปาตพฺพํ, ยถาภุตฺตญฺจ ภุญฺชถา”ติ.\csromanspacer{} เย โข ปน ภิกฺขเว อุตฺตราย ทิสาย ปฏิราชาโน, เต รญฺโญ จกฺกวตฺติสฺส อนุยนฺตา ภวนฺติ.}

\prose{อถ โข ตํ ภิกฺขเว จกฺกรตนํ สมุทฺทปริยนฺตํ ปถวิํ อภิวิชินิตฺวา ตเมว ราชธานิํ ปจฺจาคนฺตฺวา รญฺโญ จกฺกวตฺติสฺส อนฺเตปุรทฺวาเร อกฺขาหตํ มญฺเญ ติฏฺฐติ รญฺโญ จกฺกวตฺติสฺส อนฺเตปุรทฺวารํ อุปโสภยมานํ.\csromanspacer{} รญฺโญ ภิกฺขเว จกฺกวตฺติสฺส เอวรูปํ จกฺกรตนํ ปาตุภวติ. (1)}

\proseitem{258}{ปุน จปรํ ภิกฺขเว รญฺโญ จกฺกวตฺติสฺส หตฺถิรตนํ ปาตุภวติ, สพฺพเสโต สตฺตปฺปติฏฺโฐ อิทฺธิมา เวหาสงฺคโม อุโปสโถ นาม นาคราชา, ตํ ทิสฺวาน รญฺโญ จกฺกวตฺติสฺส จิตฺตํ ปสีทติ “ภทฺทกํ วต โภ หตฺถิยานํ, สเจ ทมถํ อุเปยฺยา”ติ.\csromanspacer{} อถ โข ตํ ภิกฺขเว หตฺถิรตนํ เสยฺยถาปิ นาม ภทฺโท หตฺถาชานีโย ทีฆรตฺตํ สุปริทนฺโต เอวเมว ทมถํ อุเปติ.\csromanspacer{} ภูตปุพฺพํ ภิกฺขเว ราชา จกฺกวตฺตี ตเมว หตฺถิรตนํ วีมํสมาโน ปุพฺพณฺหสมยํ อภิรุหิตฺวา สมุทฺทปริยนฺตํ ปถวิํ อนุสํยายิตฺวา ตเมว ราชธานิํ ปจฺจาคนฺตฺวา ปาตราสมกาสิ.\csromanspacer{} รญฺโญ ภิกฺขเว จกฺกวตฺติสฺส เอวรูปํ หตฺถิรตนํ ปาตุภวติ.(2)}

\prose{ปุน จปรํ ภิกฺขเว รญฺโญ จกฺกวตฺติสฺส อสฺสรตนํ ปาตุภวติ, สพฺพเสโต กาฬสีโส มุญฺชเกโส อิทฺธิมา เวหาสงฺคโม วลหโก นาม อสฺสราชา, ตํ ทิสฺวาน รญฺโญ จกฺกวตฺติสฺส จิตฺตํ ปสีทติ “ภทฺทกํ วต โภ อสฺสยานํ, สเจ ทมถํ อุเปยฺยา”ติ.\csromanspacer{} อถ โข ตํ ภิกฺขเว อสฺสรตนํ เสยฺยถาปิ นาม ภทฺโท อสฺสาชานีโย ทีฆรตฺตํ สุปริทนฺโต เอวเมว ทมถํ อุเปติ.\csromanspacer{} ภูตปุพฺพํ ภิกฺขเว ราชา จกฺกวตฺตี ตเมว อสฺสรตนํ วีมํสมาโน ปุพฺพณฺหสมยํ อภิรุหิตฺวา สมุทฺทปริยนฺตํ ปถวิํ อนุสํยายิตฺวา ตเมว ราชธานิํ ปจฺจาคนฺตฺวา ปาตราสมกาสิ.\csromanspacer{} รญฺโญ ภิกฺขเว จกฺกวตฺติสฺส เอวรูปํ อสฺสรตนํ ปาตุภวติ. (3)}

\prose{ปุน จปรํ ภิกฺขเว รญฺโญ จกฺกวตฺติสฺส มณิรตนํ ปาตุภวติ, โส โหติ มณิ เวฬุริโย สุโภ ชาติมา อฏฺฐํ โส สุปริกมฺมกโต.\csromanspacer{} ตสฺส โข ปน ภิกฺขเว มณิรตนสฺส อาภา สมนฺตา โยชนํ ผุฏา โหติ.\csromanspacer{} ภูตปุพฺพํ ภิกฺขเว ราชา จกฺกวตฺตี ตเมว มณิรตนํ วีมํสมาโน}

\csromanheader{2}{9. พาลปณฺฑิตสุตฺต (129)}
\prose{\csromanfolio{213}จตุรงฺคินิํ เสนํ สนฺนยฺหิตฺวา มณิํ ธชคฺคํ อาโรเปตฺวา รตฺตนฺธการติมิสาย ปายาสิ.\csromanspacer{} เย โข ปน ภิกฺขเว สมนฺตา คามา อเหสุํ, เต เตโนภาเสน กมฺมนฺเต ปโยเชสุํ ทิวาติ มญฺญมานา.\csromanspacer{} รญฺโญ ภิกฺขเว จกฺกวตฺติสฺส เอวรูปํ มณิรตนํ ปาตุภวติ. (4)}

\prose{ปุน จปรํ ภิกฺขเว รญฺโญ จกฺกวตฺติสฺส อิตฺถิรตนํ ปาตุภวติ, สา อภิรูปา ทสฺสนียา ปาสาทิกา ปรมาย วณฺณโปกฺขรตาย สมนฺนาคตา นาติทีฆา นาติรสฺสา นาติกิสา นาติถูลา นาติกาฬิกา\footnote{นาติกาฬี (สี, อิ)} นาจฺโจทาตา อติกฺกนฺตา มานุสํ วณฺณํ อปฺปตฺตา ทิพฺพํ วณฺณํ.\csromanspacer{} ตสฺส โข ปน ภิกฺขเว อิตฺถิรตนสฺส เอวรูโป กายสมฺผสฺโส โหติ เสยฺยถาปิ นาม ตูลปิจุโน วา กปฺปาสปิจุโน วา.\csromanspacer{} ตสฺส โข ปน ภิกฺขเว อิตฺถิรตนสฺส สีเต อุณฺหานิ คตฺตานิ โหนฺติ, อุเณฺห สีตานิ คตฺตานิ โหนฺติ.\csromanspacer{} ตสฺส โข ปน ภิกฺขเว อิตฺถิรตนสฺส กายโต จนฺทนคนฺโธ วายติ, มุขโต อุปฺปลคนฺโธ วายติ.\csromanspacer{} ตํ โข ปน ภิกฺขเว อิตฺถิรตนํ รญฺโญ จกฺกวตฺติสฺส ปุพฺพุฏฺฐายินี โหติ ปจฺฉานิปาตินี กิํการปฏิสฺสาวินี มนาปจารินี ปิยวาทินี.\csromanspacer{} ตํ โข ปน ภิกฺขเว อิตฺถิรตนํ ราชานํ จกฺกวตฺติํ มนสาปิ โน อติจรติ, กุโต ปน กาเยน.\csromanspacer{} รญฺโญ ภิกฺขเว จกฺกวตฺติสฺส เอวรูปํ อิตฺถิรตนํ ปาตุภวติ. (5)}

\prose{ปุน จปรํ ภิกฺขเว รญฺโญ จกฺกวตฺติสฺส คหปติรตนํ ปาตุภวติ, ตสฺส กมฺมวิปากชํ ทิพฺพจกฺขุ ปาตุภวติ, เยน นิธิํ ปสฺสติ สสฺสามิกมฺปิ อสฺสามิกมฺปิ.\csromanspacer{} โส ราชานํ จกฺกวตฺติํ อุปสงฺกมิตฺวา เอวมาห “อปฺโปสฺสุกฺโก ตฺวํ เทว โหหิ, อหํ เต ธเนน ธนกรณียํ\footnote{ธเนน กรณียํ (ก)} กริสฺสามี”ติ.\csromanspacer{} ภูตปุพฺพํ ภิกฺขเว ราชา จกฺกวตฺตี ตเมว คหปติรตนํ วีมํสมาโน นาวํ อภิรุหิตฺวา มชฺเฌ คงฺคาย นทิยา โสตํ โอคาหิตฺวา\footnote{โอคเหตฺวา (สี, อิ)} คหปติรตนํ เอตทโวจ “อตฺโถ เม คหปติ หิรญฺญสุวณฺเณนา”ติ.\csromanspacer{} เตน หิ มหาราช เอกํ ตีรํ นาวา อุเปตูติ.\csromanspacer{} อิเธว เม คหปติ อตฺโถ หิรญฺญสุวณฺเณนาติ.\csromanspacer{} อถ โข ตํ ภิกฺขเว คหปติรตนํ อุโภหิ หตฺเถหิ อุทเก โอมสิตฺวา ปูรํ หิรญฺญสุวณฺณสฺส กุมฺภิํ อุทฺธริตฺวา ราชานํ จกฺกวตฺติํ เอตทโวจ “อลเมตฺตาวตา มหาราช, กตเมตฺตาวตา มหาราช, ปูชิตเมตฺตาวตา มหาราชา”ติ.\csromanspacer{} ราชา จกฺกวตฺตี เอวมาห “อลเมตฺตาวตา คหปติ, กตเมตฺตาวตา คหปติ, ปูชิตเมตฺตาวตา}

\prose{\csromanfolio{214}คหปตี”ติ.\csromanspacer{} รญฺโญ ภิกฺขเว จกฺกวตฺติสฺส เอวรูปํ คหปติรตนํ ปาตุภวติ. (6)}

\prose{ปุน จปรํ ภิกฺขเว รญฺโญ จกฺกวตฺติสฺส ปริณายกรตนํ ปาตุภวติ ปณฺฑิโต พฺยตฺโต เมธาวี ปฏิพโล ราชานํ จกฺกวตฺติํ อุปยาเปตพฺพํ อุปยาเปตุํ\footnote{อุปฏฺฐเปตพฺพํ อุปฏฺฐเปตุํ (สี, สฺยา, กํ, อิ)} อปยาเปตพฺพํ อปยาเปตุํ ฐเปตพฺพํ ฐเปตุํ.\csromanspacer{} โส ราชานํ จกฺกวตฺติํ อุปสงฺกมิตฺวา เอวมาห “อปฺโปสฺสุกฺโก ตฺวํ เทว โหหิ, อหมนุสาสิสฺสามี”ติ.\csromanspacer{} รญฺโญ ภิกฺขเว จกฺกวตฺติสฺส เอวรูปํ ปริณายกรตนํ ปาตุภวติ.\csromanspacer{} ราชา ภิกฺขเว จกฺกวตฺตี อิเมหิ สตฺตหิ รตเนหิ สมนฺนาคโต โหติ. (7)}

\proseitem{259}{กตมาหิ จตูหิ อิทฺธีหิ.\csromanspacer{} อิธ ภิกฺขเว ราชา จกฺกวตฺตี อภิรูโป โหติ ทสฺสนีโย ปาสาทิโก ปรมาย วณฺณโปกฺขรตาย สมนฺนาคโต อติวิย อญฺเญหิ มนุสฺเสหิ.\csromanspacer{} ราชา ภิกฺขเว จกฺกวตฺตี อิมาย ปฐมาย อิทฺธิยา สมนฺนาคโต โหติ.}

\prose{ปุน จปรํ ภิกฺขเว ราชา จกฺกวตฺตี ทีฆายุโก โหติ จิรฏฺฐิติโก อติวิย อญฺเญหิ มนุสฺเสหิ.\csromanspacer{} ราชา ภิกฺขเว จกฺกวตฺตี อิมาย ทุติยาย อิทฺธิยา สมนฺนาคโต โหติ.}

\prose{ปุน จปรํ ภิกฺขเว ราชา จกฺกวตฺตี อปฺปาพาโธ โหติ อปฺปาตงฺโก สมเวปากินิยา คหณิยา สมนฺนาคโต นาติสีตาย นาจฺจุณฺหาย อติวิย อญฺเญหิ มนุสฺเสหิ.\csromanspacer{} ราชา ภิกฺขเว จกฺกวตฺตี อิมาย ตติยาย อิทฺธิยา สมนฺนาคโต โหติ.}

\prose{ปุน จปรํ ภิกฺขเว ราชา จกฺกวตฺตี พฺราหฺมณคหปติกานํ ปิโย โหติ มนาโป.\csromanspacer{} เสยฺยถาปิ ภิกฺขเว ปิตา ปุตฺตานํ ปิโย โหติ มนาโป.\csromanspacer{} เอวเมว โข ภิกฺขเว ราชา จกฺกวตฺตี พฺราหฺมณคหปติกานํ ปิโย โหติ มนาโป.\csromanspacer{} รญฺโญปิ ภิกฺขเว จกฺกวตฺติสฺส พฺราหฺมณคหปติกา ปิยา โหนฺติ มนาปา.\csromanspacer{} เสยฺยถาปิ ภิกฺขเว ปิตุ ปุตฺตา ปิยา โหนฺติ มนาปา.\csromanspacer{} เอวเมว โข ภิกฺขเว รญฺโญปิ จกฺกวตฺติสฺส พฺราหฺมณคหปติกา ปิยา โหนฺติ มนาปา.}

\csromanheader{2}{9. พาลปณฺฑิตสุตฺต (129)}
\prose{\csromanfolio{215}ภูตปุพฺพํ ภิกฺขเว ราชา จกฺกวตฺตี จตุรงฺคินิยา เสนาย อุยฺยานภูมิํ นิยฺยาสิ.\csromanspacer{} อถ โข ภิกฺขเว พฺราหฺมณคหปติกา ราชานํ จกฺกวตฺติํ อุปสงฺกมิตฺวา เอวมาหํสุ “อตรมาโน เทว ยาหิ, ยถา ตํ มยํ จิรตรํ ปสฺเสยฺยามา”ติ.\csromanspacer{} ราชาปิ ภิกฺขเว จกฺกวตฺตี สารถิํ อามนฺเตสิ “อตรมาโน สารถิ เปเสหิ, ยถา มํ พฺราหฺมณคหปติกา จิรตรํ ปสฺเสยฺยุนฺ”ติ.\csromanspacer{} ราชา ภิกฺขเว จกฺกวตฺตี อิมาย จตุตฺถาย อิทฺธิยา สมนฺนาคโต โหติ.\csromanspacer{} ราชา ภิกฺขเว จกฺกวตฺตี อิมาหิ จตูหิ อิทฺธีหิ สมนฺนาคโต โหติ.}

\prose{ตํ กิํ มญฺญถ ภิกฺขเว, อปิ นุ โข ราชา จกฺกวตฺตี อิเมหิ สตฺตหิ รตเนหิ สมนฺนาคโต อิมาหิ จตูหิ จ อิทฺธีหิ ตโตนิทานํ สุขํ โสมนสฺสํ ปฏิสํเวทิเยถาติ.\csromanspacer{} เอกเมเกนปิ ภนฺเต รตเนน\footnote{เตน รตเนน (สี)} สมนฺนาคโต ราชา จกฺกวตฺตี ตโตนิทานํ สุขํ โสมนสฺสํ ปฏิสํเวทิเยถ, โก ปน วาโท สตฺตหิ รตเนหิ จตูหิ จ อิทฺธีหีติ.}

\proseitem{260}{อถ โข ภควา ปริตฺตํ ปาณิมตฺตํ ปาสาณํ คเหตฺวา ภิกฺขู อามนฺเตสิ “ตํ กิํ มญฺญถ ภิกฺขเว, กตโม นุ โข มหนฺตตโร, โย จายํ มยา ปริตฺโต ปาณิมตฺโต ปาสาโณ คหิโต, โย จ หิมวา ปพฺพตราชา”ติ.\csromanspacer{} อปฺปมตฺตโก อยํ ภนฺเต ภควตา ปริตฺโต ปาณิมตฺโต ปาสาโณ คหิโต, หิมวนฺตํ ปพฺพตราชานํ อุปนิธาย สงฺขามฺปิ น อุเปติ, กลภาคมฺปิ น อุเปติ, อุปนิธมฺปิ น อุเปตีติ.\csromanspacer{} เอวเมว โข ภิกฺขเว ยํ ราชา จกฺกวตฺตี สตฺตหิ รตเนหิ สมนฺนาคโต จตูหิ จ อิทฺธีหิ ตโตนิทานํ สุขํ โสมนสฺสํ ปฏิสํเวเทติ.\csromanspacer{} ตํ ทิพฺพสฺส สุขสฺส อุปนิธาย สงฺขามฺปิ น อุเปติ, กลภาคมฺปิ น อุเปติ, อุปนิธมฺปิ น อุเปติ.}

\prose{ส โข โส ภิกฺขเว ปณฺฑิโต สเจ กทาจิ กรหจิ ทีฆสฺส อทฺธุโน อจฺจเยน มนุสฺสตฺตํ อาคจฺฉติ, ยานิ ตานิ อุจฺจากุลานิ ขตฺติยมหาสาลกุลํ วา พฺราหฺมณมหาสาลกุลํ วา คหปติมหาสาลกุลํ วา ตถารูเป กุเล ปจฺจาชายติ อฑฺเฒ มหทฺธเน มหาโภเค ปหูตชาตรูปรชเต ปหูตวิตฺตูปกรเณ ปหูตธนธญฺเญ.\csromanspacer{} โส จ โหติ อภิรูโป ทสฺสนีโย ปาสาทิโก \csromanfolio{216}ปรมาย วณฺณโปกฺขรตาย สมนฺนาคโต ลาภี อนฺนสฺส ปานสฺส วตฺถสฺส ยานสฺส มาลาคนฺธวิเลปนสฺส เสยฺยาวสถปทีเปยฺยสฺส.\csromanspacer{} โส กาเยน สุจริตํ จรติ, วาจาย สุจริตํ จรติ, มนสา สุจริตํ จรติ.\csromanspacer{} โส กาเยน สุจริตํ จริตฺวา วาจาย สุจริตํ จริตฺวา มนสา สุจริตํ จริตฺวา กายสฺส เภทา ปรํ มรณา สุคติํ สคฺคํ โลกํ อุปปชฺชติ.\csromanspacer{} เสยฺยถาปิ ภิกฺขเว อกฺขธุตฺโต ปฐเมเนว กฏคฺคเหน มหนฺตํ โภคกฺขนฺธํ อธิคจฺเฉยฺย, อปฺปมตฺตโก โส ภิกฺขเว กฏคฺคโห, ยํ โส อกฺขธุตฺโต ปฐเมเนว กฏคฺคเหน มหนฺตํ โภคกฺขนฺทํ อธิคจฺเฉยฺย, อถ โข อยเมว ตโต มหนฺตตโร กฏคฺคโห.\csromanspacer{} ยํ โส ปณฺฑิโต กาเยน สุจริตํ จริตฺวา วาจาย สุจริตํ จริตฺวา มนสา สุจริตํ จริตฺวา กายสฺส เภทา ปรํ มรณา สุคติํ สคฺคํ อุปปชฺชติ.\csromanspacer{} อยํ ภิกฺขเว เกวลา ปริปูรา ปณฺฑิตภูมีติ.}

\prose{อิทมโวจ ภควา.\csromanspacer{} อตฺตมนา เต ภิกฺขู ภควโต ภาสิตํ อภินนฺทุนฺติ.}

\nitthitamruled{พาลปณฺฑิตสุตฺตํ นิฏฺฐิตํ นวมํ.}
\csromanmatikaanchor{mtk.35}
\csromantocmarkref{subsection}{10. เทวทูตสุตฺต}{mtk.35}
\bookmark[dest={mtk.35},level=2]{10. เทวทูตสุตฺต}
\csromanheader{2}{10. เทวทูตสุตฺต}
\proseitem{261}{เอวํ เม สุตํ\csromandash{}เอกํ สมยํ ภควา สาวตฺถิยํ วิหรติ เชตวเน อนาถปิณฺฑิกสฺส อาราเม.\csromanspacer{} ตตฺร โข ภควา ภิกฺขู อามนฺเตสิ “ภิกฺขโว”ติ. “ภทนฺเต”ติ เต ภิกฺขู ภควโต ปจฺจสฺโสสุํ.\csromanspacer{} ภควา เอตทโวจ\csromandash{}}

\prose{เสยฺยถาปิ ภิกฺขเว เทฺว อคารา สทฺวารา\footnote{สนฺธิทฺวารา (ก)}, ตตฺถ จกฺขุมา ปุริโส มชฺเฌ ฐิโต ปสฺเสยฺย มนุสฺเส เคหํ ปวิสนฺเตปิ นิกฺขมนฺเตปิ อนุจงฺกมนฺเตปิ อนุวิจรนฺเตปิ.\csromanspacer{} เอวเมว โข อหํ ภิกฺขเว ทิพฺเพน จกฺขุนา วิสุทฺเธน อติกฺกนฺตมานุสเกน สตฺเต ปสฺสามิ จวมาเน อุปปชฺชมาเน หีเน ปณีเต สุวณฺเณ ทุพฺพณฺเณ สุคเต ทุคฺคเต, ยถากมฺมูปเค สตฺเต ปชานามิ “อิเม วต โภนฺโต สตฺตา กายสุจริเตน สมนฺนาคตา วจีสุจริเตน สมนฺนาคตา มโนสุจริเตน สมนฺนาคตา อริยานํ}

\csromanheader{2}{10. เทวทูตสุตฺต (130)}
\prose{\csromanfolio{217}อนุปวาทกา สมฺมาทิฏฺฐิกา สมฺมาทิฏฺฐิกมฺมสมาทานา, เต กายสฺส เภทา ปรํ มรณา สุคติํ สคฺคํ โลกํ อุปปนฺนา.\csromanspacer{} อิเม วา ปน โภนฺโต สตฺตา กายสุจริเตน สมนฺนาคตา วจีสุจริเตน สมนฺนาคตา มโนสุจริเตน สมนฺนาคตา อริยานํ อนุปวาทกา สมฺมาทิฏฺฐิกา สมฺมาทิฏฺฐิกมฺมสมาทานา, เต กายสฺส เภทา ปรํ มรณา มนุสฺเสสุ อุปปนฺนา.\csromanspacer{} อิเม วต โภนฺโต สตฺตา กายทุจฺจริเตน สมนฺนาคตา วจีทุจฺจริเตน สมนฺนาคตา มโนทุจฺจริเตน สมนฺนาคตา อริยานํ อุปวาทกา มิจฺฉาทิฏฺฐิกา มิจฺฉาทิฏฺฐิกมฺมสมาทานา, เต กายสฺส เภทา ปรํ มรณา เปตฺติวิสยํ อุปปนฺนา.\csromanspacer{} อิเม วา ปน โภนฺโต สตฺตา กายทุจฺจริเตน สมนฺนาคตา วจีทุจฺจริเตน สมนฺนาคตา มโนทุจฺจริเตน สมนฺนาคตา อริยานํ อุปวาทกา มิจฺฉาทิฏฺฐิกา มิจฺฉาทิฏฺฐิกมฺมสมาทานา, เต กายสฺส เภทา ปรํ มรณา ติรจฺฉานโยนิํ อุปปนฺนา.\csromanspacer{} อิเม วา ปน โภนฺโต สตฺตา กายทุจฺจริตา สมนฺนาคตา วจีทุจฺจริเตน สมนฺนาคตา มโนทุจฺจริเตน สมนฺนาคตา อริยานํ อุปวาทกา มิจฺฉาทิฏฺฐิกา มิจฺฉาทิฏฺฐิกมฺมสมาทานา, เต กายสฺส เภทา ปรํ มรณา อปายํ ทุคฺคติํ วินิปาตํ นิรยํ อุปปนฺนา”ติ.}

\proseitem{262}{ตเมนํ ภิกฺขเว นิรยปาลา นานาพาหาสุ คเหตฺวา ยมสฺส รญฺโญ ทสฺเสนฺติ. “อยํ เทว ปุริโส อมตฺเตยฺโย อเปตฺเตยฺโย อสามญฺโญ อพฺราหฺมญฺโญ น กุเล เชฏฺฐาปจายี, อิมสฺส เทโว ทณฺฑํ ปเณตู”ติ.\csromanspacer{} ตเมนํ ภิกฺขเว ยโม ราชา ปฐมํ เทวทูตํ สมนุยุญฺชติ สมนุคาหติ สมนุภาสติ “อมฺโภ ปุริส น ตฺวํ อทฺทส มนุสฺเสสุ ปฐมํ เทวทูตํ ปาตุภูตนฺ”ติ? โส เอวมาห “นาทฺทสํ ภนฺเต”ติ.}

\prose{ตเมนํ ภิกฺขเว ยโม ราชา เอวมาห “อมฺโภ ปุริส น ตฺวํ อทฺทส มนุสฺเสสุ ทหรํ กุมารํ มนฺทํ อุตฺตานเสยฺยกํ สเก มุตฺตกรีเส ปลิปนฺนํ เสมานนฺ”ติ? โส เอวมาห “อทฺทสํ ภนฺเต”ติ.}

\prose{ตเมนํ ภิกฺขเว ยโม ราชา เอวมาห “อมฺโภ ปุริส ตสฺส เต วิญฺญุสฺส สโต มหลฺลกสฺส น เอตทโหสิ ‘อหมฺปิ โขมฺหิ ชาติธมฺโม, ชาติํ อนตีโต, หนฺทาหํ กลฺยาณํ กโรมิ กาเยน วาจาย มนสา’ติ”? โส เอวมาห “นาสกฺขิสฺสํ ภนฺเต ปมาทสฺสํ ภนฺเต”ติ.}

\prose{\csromanfolio{218}ตเมนํ ภิกฺขเว ยโม ราชา เอวมาห “อมฺโภ ปุริส ปมาทวตาย น กลฺยาณมกาสิ กาเยน วาจาย มนสา, ตคฺฆ ตฺวํ อมฺโภ ปุริส ตถา กริสฺสนฺติ, ยถา ตํ ปมตฺตํ.\csromanspacer{} ตํ โข ปน เต เอตํ ปาปกมฺมํ\footnote{ปาปํ กมฺมํ (สี, อิ)} เนว มาตรา กตํ, น ปิตรา กตํ, น ภาตรา กตํ, น ภคินิยา กตํ, น มิตฺตามจฺเจหิ กตํ, น ญาติสาโลหิเตหิ กตํ, น สมณพฺราหฺมเณหิ กตํ, น เทวตาหิ กตํ, ตยาเวตํ ปาปกมฺมํ1 กตํ, ตฺวญฺเญเวตสฺส วิปากํ ปฏิสํเวทิสฺสสี”ติ. (1)}

\proseitem{263}{ตเมนํ ภิกฺขเว ยโม ราชา ปฐมํ เทวทูตํ สมนุยุญฺชิตฺวา สมนุคาหิตฺวา สมนุภาสิตฺวา ทุติยํ เทวทูตํ สมนุยุญฺชติ สมนุคาหติ สมนุภาสติ “อมฺโภ ปุริส น ตฺวํ อทฺทส มนุสฺเสสุ ทุติยํ เทวทูตํ ปาตุภูตนฺ”ติ? โส เอวมาห “นาทฺทสํ ภนฺเต”ติ.}

\prose{ตเมนํ ภิกฺขเว ยโม ราชา เอวมาห “อมฺโภ ปุริส น ตฺวํ อทฺทส มนุสฺเสสุ อิตฺถิํ วา ปุริสํ วา ( )\footnote{(อาสีติกํ วา นาวุติกํ วา วสฺสสติกํ วา ชาติยา) (ก-สี, สฺยา, กํ, อิ) ติกงฺคุตฺตเรปิ.} ชิณฺณํ โคปานสิวงฺกํ โภคฺคํ ทณฺฑปรายนํ ปเวธมานํ คจฺฉนฺตํ อาตุรํ คตโยพฺพนํ ขณฺฑทนฺตํ ปลิตเกสํ วิลูนํ ขลิตสิรํ\footnote{ขลิตํ สิโร (สี), ขลิตํสิรํ (สฺยา, กํ, อิ)} วลินํ ติลกาหตคตฺตนฺ”ติ? โส เอวมาห “อทฺทสํ ภนฺเต”ติ.}

\prose{ตเมนํ ภิกฺขเว ยโม ราชา เอวมาห “อมฺโภ ปุริส ตสฺส เต วิญฺญุสฺส สโต มหลฺลกสฺส น เอตทโหสิ ‘อหมฺปิ โขมฺหิ ชราธมฺโม, ชรํ อนตีโต, หนฺทาหํ กลฺยาณํ กโรมิ กาเยน วาจาย มนสา’ติ”? โส เอวมาห “นาสกฺขิสฺสํ ภนฺเต, ปมาทสฺสํ ภนฺเต”ติ.}

\prose{ตเมนํ ภิกฺขเว ยโม ราชา เอวมาห “อมฺโภ ปุริส ปมาทวตาย น กลฺยาณมกาสิ กาเยน วาจาย มนสา, ตคฺฆ ตฺวํ อมฺโภ ปุริส ตถา กริสฺสนฺติ, ยถา ตํ ปมตฺตํ.\csromanspacer{} ตํ โข ปน เต เอตํ ปาปกมฺมํ เนว มาตรา กตํ, น ปิตรา กตํ, น ภาตรา กตํ, น ภคินิยา กตํ, น มิตฺตามจฺเจหิ กตํ, น ญาติสาโลหิเตหิ กตํ,}

\csromanheader{2}{10. เทวทูตสุตฺต (130)}
\prose{\csromanfolio{219}น สมณพฺราหฺมเณหิ กตํ, น เทวตาหิ กตํ, ตยาเวตํ ปาปกมฺมํ กตํ, ตฺวญฺเญเวตสฺส วิปากํ ปฏิสํเวทิสฺสสี”ติ. (2)}

\proseitem{264}{ตเมนํ ภิกฺขเว ยโม ราชา ทุติยํ เทวทูตํ สมนุยุญฺชิตฺวา สมนุคาหิตฺวา สมนุภาสิตฺวา ตติยํ เทวทูตํ สมนุยุญฺชติ สมนุคาหติ สมนุภาสติ “อมฺโภ ปุริส น ตฺวํ อทฺทส มนุสฺเสสุ ตติยํ เทวทูตํ ปาตุภูตนฺ”ติ? โส เอวมาห “นาทฺทสํ ภนฺเต”ติ.}

\prose{ตเมนํ ภิกฺขเว ยโม ราชา เอวมาห “อมฺโภ ปุริส น ตฺวํ อทฺทส มนุสฺเสสุ อิตฺถิํ วา ปุริสํ วา อาพาธิกํ ทุกฺขิตํ พาฬฺหคิลานํ สเก มุตฺตกรีเส ปลิปนฺนํ เสมานํ อญฺเญหิ วุฏฺฐาปิยมานํ อญฺเญหิ สํเวสิยมานนฺ”ติ? โส เอวมาห “อทฺทสํ ภนฺเต”ติ.}

\prose{ตเมนํ ภิกฺขเว ยโม ราชา เอวมาห “อมฺโภ ปุริส ตสฺส เต วิญฺญุสฺส สโต มหลฺลกสฺส น เอตทโหสิ ‘อหมฺปิ โขมฺหิ พฺยาธิธมฺโม, พฺยาธิํ อนตีโต, หนฺทาหํ กลฺยาณํ กโรมิ กาเยน วาจาย มนสา’ติ”? โส เอวมาห “นาสกฺขิสฺสํ ภนฺเต, ปมาทสฺสํ ภนฺเต”ติ.}

\prose{ตเมนํ ภิกฺขเว ยโม ราชา เอวมาห “อมฺโภ ปุริส ปมาทวตาย น กลฺยาณมกาสิ กาเยน วาจาย มนสา, ตคฺฆ ตฺวํ อมฺโภ ปุริส ตถา กริสฺสนฺติ, ยถา ตํ ปมตฺตํ.\csromanspacer{} ตํ โข ปน เต เอตํ ปาปกมฺมํ เนว มาตรา กตํ, น ปิตรา กตํ, น ภาตรา กตํ, น ภคินิยา กตํ, น มิตฺตามจฺเจหิ กตํ, น ญาติสาโลหิเตหิ กตํ, น สมณพฺราหฺมเณหิ กตํ, น เทวตาหิ กตํ, ตยาเวตํ ปาปกมฺมํ กตํ, ตฺวญฺเญเวตสฺส วิปากํ ปฏิสํเวทิสฺสสี”ติ. (3)}

\proseitem{265}{ตเมนํ ภิกฺขเว ยโม ราชา ตติยํ เทวทูตํ สมนุยุญฺชิตฺวา สมนุคาหิตฺวา สมนุภาสิตฺวา จตุตฺถํ เทวทูตํ สมนุยุญฺชติ สมนุคาหติ สมนุภาสติ “อมฺโภ ปุริส น ตฺวํ อทฺทส มนุสฺเสสุ จตุตฺถํ เทวทูตํ ปาตุภูตนฺ”ติ? โส เอวมาห “นาทฺทสํ ภนฺเต”ติ.}

\prose{ตเมนํ ภิกฺขเว ยโม ราชา เอวมาห “อมฺโภ ปุริส น ตฺวํ อทฺทส มนุสฺเสสุ ราชาโน โจรํ อาคุจริํ คเหตฺวา วิวิธา กมฺมการณา กาเรนฺเต กสาหิปิ ตาเฬนฺเต เวตฺเตหิปิ ตาเฬนฺเต อทฺธทณฺฑเกหิปิ ตาเฬนฺเต หตฺถมฺปิ ฉินฺทนฺเต ปาทมฺปิ ฉินฺทนฺเต หตฺถปาทมฺปิ ฉินฺทนฺเต กณฺณมฺปิ \csromanfolio{220}ฉินฺทนฺเต นาสมฺปิ ฉินฺทนฺเต กณฺณนาสมฺปิ ฉินฺทนฺเต พิลงฺคถาลิกมฺปิ กโรนฺเต สงฺขมุณฺฑิกมฺปิ กโรนฺเต ราหุมุขมฺปิ กโรนฺเต โชติมาลิกมฺปิ กโรนฺเต หตฺถปชฺโชติกมฺปิ กโรนฺเต เอรกวตฺติกมฺปิ กโรนฺเต จีรกวาสิกมฺปิ กโรนฺเต เอเณยฺยกมฺปิ กโรนฺเต พฬิสมํสิกมฺปิ กโรนฺเต กหาปณิกมฺปิ กโรนฺเต ขาราปตจฺฉิกมฺปิ กโรนฺเต ปลิฆปริวตฺติกมฺปิ กโรนฺเต ปลาลปีฐกมฺปิ กโรนฺเต ตตฺเตนปิ เตเลน โอสิญฺจนฺเต สุนเขหิปิ ขาทาเปนฺเต ชีวนฺตมฺปิ สูเล อุตฺตาเสนฺเต อสินาปิ สีสํ ฉินฺทนฺเต”ติ? โส เอวมาห “อทฺทสํ ภนฺเต”ติ.}

\prose{ตเมนํ ภิกฺขเว ยโม ราชา เอวมาห “อมฺโภ ปุริส ตสฺส เต วิญฺญุสฺส สโต มหลฺลกสฺส น เอตทโหสิ ‘เย กิร โภ ปาปกานิ กมฺมานิ กโรนฺติ, เต ทิฏฺเฐว เอวรูปา วิวิธา กมฺมการณา กรียนฺติ, กิมงฺคํ\footnote{กิมงฺค (สี, อิ)} ปน ปรตฺถ.\csromanspacer{} หนฺทาหํ กลฺยาณํ กโรมิ กาเยน วาจาย มนสา’ติ”? โส เอวมาห “นาสกฺขิสฺสํ ภนฺเต, ปมาทสฺสํ ภนฺเต”ติ.}

\prose{ตเมนํ ภิกฺขเว ยโม ราชา เอวมาห “อมฺโภ ปุริส ปมาทวตาย น กลฺยาณมกาสิ กาเยน วาจาย มนสา, ตคฺฆ ตฺวํ อมฺโภ ปุริส ตถา กริสฺสนฺติ, ยถา ตํ ปมตฺตํ.\csromanspacer{} ตํ โข ปน เต เอตํ ปาปกมฺมํ เนว มาตรา กตํ, น ปิตรา กตํ, น ภาตรา กตํ, น ภคินิยา กตํ, น มิตฺตามจฺเจหิ กตํ, น ญาติสาโลหิเตหิ กตํ, น สมณพฺราหฺมเณหิ กตํ, น เทวตาหิ กตํ, ตยาเวตํ ปาปกมฺมํ กตํ, ตฺวญฺเญเวตสฺส วิปากํ ปฏิสํเวทิสฺสสี”ติ. (4)}

\proseitem{266}{ตเมนํ ภิกฺขเว ยโม ราชา จตุตฺถํ เทวทูตํ สมนุยุญฺชิตฺวา สมนุคาหิตฺวา สมนุภาสิตฺวา ปญฺจมํ เทวทูตํ สมนุยุญฺชติ สมนุคาหติ สมนุภาสติ “อมฺโภ ปุริส น ตฺวํ อทฺทส มนุสฺเสสุ ปญฺจมํ เทวทูตํ ปาตุภูตนฺ”ติ? โส เอวมาห “นาทฺทสํ ภนฺเต”ติ.}

\prose{ตเมนํ ภิกฺขเว ยโม ราชา เอวมาห “อมฺโภ ปุริส น ตฺวํ อทฺทส มนุสฺเสสุ อิตฺถิํ วา ปุริสํ วา เอกาหมตํ วา ทฺวีหมตํ วา ตีหมตํ วา อุทฺธุมาตกํ วินีลกํ วิปุพฺพกชาตนฺ”ติ? โส เอวมาห “อทฺทสํ ภนฺเต”ติ.}

\csromanheader{2}{10. เทวทูตสุตฺต (130)}
\prose{\csromanfolio{221}ตเมนํ ภิกฺขเว ยโม ราชา เอวมาห “อมฺโภ ปุริส ตสฺส เต วิญฺญุสฺส สโต มหลฺลกสฺส น เอตทโหสิ ‘อหมฺปิ โขมฺหิ มรณธมฺโม, มรณํ อนตีโต, หนฺทาหํ กลฺยาณํ กโรมิ กาเยน วาจาย มนสา’ติ”? โส เอวมาห “นาสกฺขิสฺสํ ภนฺเต, ปมาทสฺสํ ภนฺเต”ติ.}

\prose{ตเมนํ ภิกฺขเว ยโม ราชา เอวมาห “อมฺโภ ปุริส ปมาทวตาย น กลฺยาณมกาสิ กาเยน วาจาย มนสา, ตคฺฆ ตฺวํ อมฺโภ ปุริส ตถา กริสฺสนฺติ, ยถา ตํ ปมตฺตํ.\csromanspacer{} ตํ โข ปน เต เอตํ ปาปกมฺมํ เนว มาตรา กตํ, น ปิตรา กตํ, น ภาตรา กตํ, น ภคินิยา กตํ, น มิตฺตามจฺเจหิ กตํ, น ญาติสาโลหิเตหิ กตํ, น สมณพฺราหฺมเณหิ กตํ, น เทวตาหิ กตํ, ตยาเวตํ ปาปกมฺมํ กตํ, ตฺวญฺเญเวตสฺส วิปากํ ปฏิสํเวทิสฺสสี”ติ. (5)}

\proseitem{267}{ตเมนํ ภิกฺขเว ยโม ราชา ปญฺจมํ เทวทูตํ สมนุยุญฺชิตฺวา สมนุคาหิตฺวา สมนุภาสิตฺวา ตุณฺหี โหติ.\csromanspacer{} ตเมนํ ภิกฺขเว นิรยปาลา ปญฺจวิธพนฺธนํ นาม กมฺมการณํ กโรนฺติ, ตตฺตํ อโยขิลํ หตฺเถ คเมนฺติ, ตตฺตํ อโยขิลํ ทุติเย หตฺเถ คเมนฺติ, ตตฺตํ อโยขิลํ ปาเท คเมนฺติ, ตตฺตํ อโยขิลํ ทุติเย ปาเท คเมนฺติ, ตตฺตํ อโยขิลํ มชฺเฌ อุรสฺมิํ คเมนฺติ.\csromanspacer{} โส ตตฺถ ทุกฺขา ติพฺพา ขรา กฏุกา เวทนา เวเทติ, น จ ตาว กาลํ กโรติ, ยาว น ตํ ปาปกมฺมํ พฺยนฺตีโหติ.\csromanspacer{} ตเมนํ ภิกฺขเว นิรยปาลา อุทฺธํปาทํ อโธสิรํ คเหตฺวา วาสีหิ ตจฺฉนฺติ ฯเปฯ.\csromanspacer{} ตเมนํ ภิกฺขเว นิรยปาลา รเถ โยเชตฺวา อาทิตฺตาย ปถวิยา สมฺปชฺชลิตาย สโชติภูตาย สาเรนฺติปิ ปจฺจาสาเรนฺติปิ ฯเปฯ.\csromanspacer{} ตเมนํ ภิกฺขเว นิรยปาลา มหนฺตํ องฺคารปพฺพตํ อาทิตฺตํ สมฺปชฺชลิตํ สโชติภูตํ อาโรเปนฺติปิ โอโรเปนฺติปิ ฯเปฯ.\csromanspacer{} ตเมนํ ภิกฺขเว นิรยปาลา อุทฺธํปาทํ อโธสิรํ คเหตฺวา ตตฺตาย โลหกุมฺภิยา ปกฺขิปนฺติ อาทิตฺตาย สมฺปชฺชลิตาย สโชติภูตาย.\csromanspacer{} โส ตตฺถ เผณุทฺเทสกํ ปจฺจติ, โส ตตฺถ เผณุทฺเทหกํ ปจฺจมาโน สกิมฺปิ อุทฺธํ คจฺฉติ, สกิมฺปิ อโธ คจฺฉติ, สกิมฺปิ กิริยํ คจฺฉติ.\csromanspacer{} โส ตตฺถ ทุกฺขา ติพฺพา ขรา กฏุกา เวทนา เวเทติ, น จ ตาว กาลํ กโรติ, ยาว น ปาปกมฺมํ พฺยนฺตีโหติ.\csromanspacer{} ตเมนํ ภิกฺขเว นิรยปาลา มหานิรเย ปกฺขิปนฺติ.\csromanspacer{} โส โข ปน ภิกฺขเว มหานิรโย\csromandash{}}

\csromangathameasure{จตุกฺกณฺโณ จตุทฺวาโร, วิภตฺโต ภาคโส มิโต. \\ อโยปาการปริยนฺโต, อยสา ปฏิกุชฺชิโต. \\ ตสฺส อโยมยา ภูมิ, ชลิตา เตชสายุตา. \\ สมนฺตา โยชนสตํ, ผริตฺวา ติฏฺฐติ สพฺพทา.}
\csromangathagroup{\csromangathastanza{\csromanfolio{222}จตุกฺกณฺโณ จตุทฺวาโร, วิภตฺโต ภาคโส มิโต. \\ อโยปาการปริยนฺโต, อยสา ปฏิกุชฺชิโต.}\csromangathastanza{ตสฺส อโยมยา ภูมิ, ชลิตา เตชสายุตา. \\ สมนฺตา โยชนสตํ, ผริตฺวา ติฏฺฐติ สพฺพทา.}}

\proseitem{268}{ตสฺส โข ปน ภิกฺขเว มหานิรยสฺส ปุรตฺถิมาย ภิตฺติยา อจฺจิ อุฏฺฐหิตฺวา ปจฺฉิมาย ภิตฺติยา ปฏิหญฺญติ.\csromanspacer{} ปจฺฉิมาย ภิตฺติยา อจฺจิ อุฏฺฐหิตฺวา ปุรตฺถิมาย ภิตฺติยา ปฏิหญฺญติ.\csromanspacer{} อุตฺตราย ภิตฺติยา อจฺจิ อุฏฺฐหิตฺวา ทกฺขิณาย ภิตฺติยา ปฏิหญฺญติ.\csromanspacer{} ทกฺขิณาย ภิตฺติยา อจฺจิ อุฏฺฐหิตฺวา อุตฺตราย ภิตฺติยา ปฏิหญฺญติ.\csromanspacer{} เหฏฺฐา อจฺจิ อุฏฺฐหิตฺวา อุปริ ปฏิหญฺญติ.\csromanspacer{} อุปริโต อจฺจิ อุฏฺฐหิตฺวา เหฏฺฐา ปฏิหญฺญติ.\csromanspacer{} โส ตตฺถ ทุกฺขา ติพฺพา ขรา กฏุกา เวทนา เวเทติ, น จ ตาว กาลํ กโรติ, ยาว น ตํ ปาปกมฺมํ พฺยนฺตีโหติ.}

\prose{โหติ โข โส ภิกฺขเว สมโย, ยํ กทาจิ กรหจิ ทีฆสฺส อทฺธุโน อจฺจเยน ตสฺส มหานิรยสฺส ปุรตฺถิมํ ทฺวารํ อปาปุรียติ\footnote{อวาปุรียติ (สี)}.\csromanspacer{} โส ตตฺถ สีเฆน ชเวน ธาวติ, ตสฺส สีเฆน ชเวน ธาวโต ฉวิมฺปิ ฑยฺหติ, จมฺมมฺปิ ฑยฺหติ, มํสมฺปิ ฑยฺหติ, นฺหารุมฺปิ ฑยฺหติ, อฏฺฐีนิปิ สมฺปธูปายนฺติ, อุพฺภตํ ตาทิสเมว โหติ.\csromanspacer{} ยโต จ โข โส ภิกฺขเว พหุสมฺปตฺโต โหติ, อถ ตํ ทฺวารํ ปิธียติ\footnote{ปิถียติ (สี, สฺยา, กํ, อิ)}.\csromanspacer{} โส ตตฺถ ทุกฺขา ติพฺพา ขรา กฏุกา เวทนา เวเทติ, น จ ตาว กาลํ กโรติ, ยาว น ตํ ปาปกมฺมํ พฺยนฺตีโหติ.}

\prose{โหติ โข โส ภิกฺขเว สมโย, ยํ กทาจิ กรหจิ ทีฆสฺส อทฺธุโน อจฺจเยน ตสฺส มหานิรยสฺส ปจฺฉิมํ ทฺวารํ อปาปุรียติ ฯเปฯ อุตฺตรํ ทฺวารํ อปาปุรียติ ฯเปฯ ทกฺขิณํ ทฺวารํ อปาปุรียติ.\csromanspacer{} โส ตตฺถ สีเฆน ชเวน ธาวติ, ตสฺส สีเฆน ชเวน ธาวโต ฉวิมฺปิ ฑยฺหติ, จมฺมมฺปิ ฑยฺหติ, มํสมฺปิ ฑยฺหติ, นฺหารุมฺปิ ฑยฺหติ, อฏฺฐีนิปิ สมฺปทูปายนฺติ, อุพฺภตํ ตาทิสเมว โหติ.\csromanspacer{} ยโต จ โข โส ภิกฺขเว พหุสมฺปตฺโต โหติ, อถ ตํ ทฺวารํ ปิธียติ.\csromanspacer{} โส ตตฺถ ทุกฺขา ติพฺพา ขรา กฏุกา เวทนา เวเทติ, น จ ตาว กาลํ กโรติ, ยาว น ตํ ปาปกมฺมํ พฺยนฺตีโหติ.}

\csromanheader{2}{10. เทวทูตสุตฺต (130)}
\prose{\csromanfolio{223}โหติ โข โส ภิกฺขเว สมโย, ยํ กทาจิ กรหจิ ทีฆสฺส อทฺธุโน อจฺจเยน ตสฺส มหานิรยสฺส ปุรตฺถิมํ ทฺวารํ อปาปุรียติ.\csromanspacer{} โส ตตฺถ สีเฆน ชเวน ธาวติ, ตสฺส สีเฆน ชเวน ธาวโต ฉวิมฺปิ ฑยฺหติ, จมฺมมฺปิ ฑยฺหติ, มํสมฺปิ ฑยฺหติ, นฺหารุมฺปิ ฑยฺหติ, อฏฺฐีนิปิ สมฺปธูปายนฺติ, อุพฺภตํ ตาทิสเมว โหติ.\csromanspacer{} โส เตน ทฺวาเรน นิกฺขมติ.}

\proseitem{269}{ตสฺส โข ปน ภิกฺขเว มหานิรยสฺส สมนนฺตรา สหิตเมว มหนฺโต คูถนิรโย, โส ตตฺถ ปตติ.\csromanspacer{} ตสฺมิํ โข ปน ภิกฺขเว คูถนิรเย สูจิมุขา ปาณา ฉวิํ ฉินฺทนฺติ, ฉวิํ เฉตฺวา จมฺมํ ฉินฺทนฺติ, จมฺมํ เฉตฺวา มํสํ ฉินฺทนฺติ, มํสํ เฉตฺวา นฺหารุํ ฉินฺทนฺติ, นฺหารุํ เฉตฺวา อฏฺฐิํ ฉินฺทนฺติ, อฏฺฐิํ เฉตฺวา อฏฺฐิมิญฺชํ ขาทนฺติ.\csromanspacer{} โส ตตฺถ ทุกฺขา ติพฺพา ขรา กฏุกา เวทนา เวเทติ, น จ ตาว กาลํ กโรติ, ยาว น ตํ ปาปกมฺมํ พฺยนฺตีโหติ. (1)}

\prose{ตสฺส โข ปน ภิกฺขเว คูถนิรยสฺส สมนนฺตรา สหิตเมว มหนฺโต กุกฺกุลนิรโย, โส ตตฺถ ปตติ.\csromanspacer{} โส ตตฺถ ทุกฺขา ติพฺพา ขรา กฏุกา เวทนา เวเทติ, น จ ตาว กาลํ กโรติ, ยาว น ตํ ปาปกมฺมํ พฺยนฺตีโหติ. (2)}

\prose{ตสฺส โข ปน ภิกฺขเว กุกฺกุลนิรยสฺส สมนนฺตรา สหิตเมว มหนฺตํ สิมฺพลิวนํ อุทฺธํ\footnote{อุจฺจํ (สฺยา, กํ), อุพฺภโต (ก)} โยชนมุคฺคตํ โสฬสงฺคุลกณฺฏกํ\footnote{โสฬสงฺคุลกณฺฑกํ (สี)} อาทิตฺตํ สมฺปชฺชลิตํ สโชติภูตํ, ตตฺถ อาโรเปนฺติปิ โอโรเปนฺติปิ.\csromanspacer{} โส ตตฺถ ทุกฺขา ติพฺพา ขรา กฏุกา เวทนา เวเทติ, น จ ตาว กาลํ กโรติ, ยาว น ตํ ปาปกมฺมํ พฺยนฺตีโหติ. (3)}

\prose{ตสฺส โข ปน ภิกฺขเว สิมฺพลิวนสฺส สมนนฺตรา สหิตเมว มหนฺตํ อสิปตฺตวนํ, โส ตตฺถ ปวิสติ.\csromanspacer{} ตสฺส วาเตริตานิ ปตฺตานิ ปติตานิ หตฺถมฺปิ ฉินฺทนฺติ, ปาทมฺปิ ฉินฺทนฺติ, หตฺถปาทมฺปิ ฉินฺทนฺติ, กณฺณมฺปิ ฉินฺทนฺติ, นาสมฺปิ ฉินฺทนฺติ, กณฺณนาสมฺปิ ฉินฺทนฺติ.\csromanspacer{} โส ตตฺถ ทุกฺขา ติพฺพา ขรา กฏุกา เวทนา เวเทติ, น จ ตาว กาลํ กโรติ, ยาว น ตํ ปาปกมฺมํ พฺยนฺตีโหติ. (4)}

\prose{ตสฺส โข ปน ภิกฺขเว อสิปตฺตวนสฺส สมนนฺตรา สหิตเมว มหตี ขาโรทกา นที\footnote{ขาโรทิกา นที (สี)}, โส ตตฺถ ปตติ.\csromanspacer{} โส ตตฺถ อนุโสตมฺปิ \csromanfolio{224}วุยฺหติ, ปฏิโสตมฺปิ วุยฺหติ, อนุโสตปฏิโสตมฺปิ วุยฺหติ.\csromanspacer{} โส ตตฺถ ทุกฺขา ติพฺพา ขรา กฏุกา เวทนา เวเทติ, น จ ตาว กาลํ กโรติ, ยาว น ตํ ปาปกมฺมํ พฺยนฺตีโหติ.(5)}

\proseitem{270}{ตเมนํ ภิกฺขเว นิรยปาลา พลิเสน อุทฺธริตฺวา ถเล ปติฏฺฐาเปตฺวา เอวมาหํสุ “อมฺโภ ปุริส กิํ อิจฺฉสี”ติ.\csromanspacer{} โส เอวมาห “ชิฆจฺฉิโตสฺมิ ภนฺเต”ติ.\csromanspacer{} ตเมนํ ภิกฺขเว นิรยปาลา ตตฺเตน อโยสงฺกุนา มุขํ วิวริตฺวา อาทิตฺเตน สมฺปชฺชลิเตน สโชติภูเตน ตตฺตํ โลหคุฬํ มุเข ปกฺขิปนฺติ อาทิตฺตํ สมฺปชฺชลิตํ สโชติภูตํ.\csromanspacer{} โส ตสฺส\footnote{ตํ ตสฺส (ก), ตสฺส (สี, อิ)} โอฏฺฐมฺปิ ทหติ\footnote{qอยฺหติ (สี, สฺยา, กํ, อิ)}, มุขมฺปิ ทหติ, กณฺฐมฺปิ ทหติ, อุรมฺปิ\footnote{อุทรมฺปิ (สี, สฺยา, กํ)} ทหติ, อนฺตมฺปิ อนฺตคุณมฺปิ อาทาย อโธภาคา นิกฺขมติ.\csromanspacer{} โส ตตฺถ ทุกฺขา ติพฺพา ขรา กฏุกา เวทนา เวเทติ, น จ ตาว กาลํ กโรติ, ยาว น ตํ ปาปกมฺมํ พฺยนฺตีโหติ.\csromanspacer{} ตเมนํ ภิกฺขเว นิรยปาลา เอวมาหํสุ “อมฺโภ ปุริส กิํ อิจฺฉสี”ติ.\csromanspacer{} โส เอวมาห “ปิปาสิโตสฺมิ ภนฺเต”ติ.\csromanspacer{} ตเมนํ ภิกฺขเว นิรยปาลา ตตฺเตน อโยสงฺกุนา มุขํ วิวริตฺวา อาทิตฺเตน สมฺปชฺชลิเตน สโชติภูเตน ตตฺตํ ตมฺพโลหํ มุเข อาสิญฺจนฺติ อาทิตฺตํ สมฺปชฺชลิตํ สโชติภูตํ.\csromanspacer{} ตํ ตสฺส\footnote{เอตฺถ ปน ปาฐเภโท นตฺถิ.} โอฏฺฐมฺปิ ทหติ, มุขมฺปิ ทหติ, กณฺฐมฺปิ ทหติ, อุรมฺปิ ทหติ, อนฺตมฺปิ อนฺตคุณมฺปิ อาทาย อโธภาคา นิกฺขมติ.\csromanspacer{} โส ตตฺถ ทุกฺขา ติพฺพา ขรา กฏุกา เวทนา เวเทติ, น จ ตาว กาลํ กโรติ, ยาว น ตํ ปาปกมฺมํ พฺยนฺตีโหติ.\csromanspacer{} ตเมนํ ภิกฺขเว นิรยปาลา ปุน มหานิรเย ปกฺขิปนฺติ.}

\prose{ภูตปุพฺพํ ภิกฺขเว ยมสฺส รญฺโญ เอตทโหสิ “เย กิร โภ โลเก ปาปกานิ อกุสลานิ กมฺมานิ กโรนฺติ, เต เอวรูปา วิวิธา กมฺมการณา กรียนฺติ.\csromanspacer{} อโห วตาหํ มนุสฺสตฺตํ ลเภยฺยํ, ตถาคโต จ โลเก อุปฺปชฺเชยฺย อรหํ สมฺมาสมฺพุทฺโธ, ตญฺจาหํ ภควนฺตํ ปยิรุปาเสยฺยํ, โส จ เม ภควา ธมฺมํ เทเสยฺย, ตสฺส จาหํ ภควโต ธมฺมํ อาชาเนยฺยนฺ”ติ.\csromanspacer{} ตํ โข ปนาหํ ภิกฺขเว นาญฺญสฺส สมณสฺส วา พฺราหฺมณสฺส วา สุตฺวา วทามิ.\csromanspacer{} อปิ จ ยเทว สามํ ญาตํ สามํ ทิฏฺฐํ สามํ วิทิตํ, ตเทวาหํ วทามีติ.}

\csromanheader{2}{10. เทวทูตสุตฺต (130)}
\csromanmatikaanchor{mtk.36}
\csromantocmarkref{subsection}{อุทฺทานคาถา}{mtk.36}
\bookmark[dest={mtk.36},level=2]{อุทฺทานคาถา}
\proseitem{271}{\csromanfolio{225}อิทมโวจ ภควา.\csromanspacer{} อิทํ วตฺวาน\footnote{อิทํ วตฺวา (สี, อิ) เอวมีทิเสสุ ฐาเนสุ.} สุคโต, อถาปรํ เอตทโวจ สตฺถา\csromandash{}}

\csromangathameasure{“โจทิตา เทวทูเตหิ, เย ปมชฺชนฺติ มาณวา. \\ เต ทีฆรตฺตํ โสจนฺติ, หีนกายูปคา นรา. \\ เย จ โข เทวทูเตหิ, สนฺโต สปฺปุริสา อิธ. \\ โจทิตา นปฺปมชฺชนฺติ, อริยธมฺเม กุทาจนํ. \\ อุปาทาเน ภยํ ทิสฺวา, ชาติมรณสมฺภเว. \\ อนุปาทา วิมุจฺจนฺติ, ชาติมรณสงฺขเย. \\ เต เขมปฺปตฺตา สุขิโน, ทิฏฺฐธมฺมาภินิพฺพุตา. \\ สพฺพเวรภยาตีตา, สพฺพทุกฺขํ อุปจฺจคุนฺ”ติ.}
\csromangathagroup{\csromangathastanza{“โจทิตา เทวทูเตหิ, เย ปมชฺชนฺติ มาณวา. \\ เต ทีฆรตฺตํ โสจนฺติ, หีนกายูปคา นรา.}\csromangathastanza{เย จ โข เทวทูเตหิ, สนฺโต สปฺปุริสา อิธ. \\ โจทิตา นปฺปมชฺชนฺติ, อริยธมฺเม กุทาจนํ.}\csromangathastanza{อุปาทาเน ภยํ ทิสฺวา, ชาติมรณสมฺภเว. \\ อนุปาทา วิมุจฺจนฺติ, ชาติมรณสงฺขเย.}\csromangathastanza{เต เขมปฺปตฺตา สุขิโน, ทิฏฺฐธมฺมาภินิพฺพุตา. \\ สพฺพเวรภยาตีตา, สพฺพทุกฺขํ\footnote{สพฺพทุกฺขา (ก)} อุปจฺจคุนฺ”ติ.}}

\nitthitam{เทวทูตสุตฺตํ นิฏฺฐิตํ ทสมํ.}
\nitthitam{\textbf{สุญฺญตวคฺโค นิฏฺฐิโต ตติโย.}}
\titlehead{\textbf{ตสฺสุทฺทานํ}}
\csromangathameasure{ทฺวิธาว สุญฺญตา โหติ, อพฺภุตธมฺมพากุลํ. \\ อจิรวตภูมิชนาโม, อนุรุทฺธุปกฺกิเลสํ. \\ พาลปณฺฑิโต เทวทูตญฺจ เต ทสาติ.}
\csromangathagroup{\csromangathastanza{ทฺวิธาว สุญฺญตา โหติ, อพฺภุตธมฺมพากุลํ. \\ อจิรวตภูมิชนาโม, อนุรุทฺธุปกฺกิเลสํ. \\ พาลปณฺฑิโต เทวทูตญฺจ เต ทสาติ.}}

\csromanmatikaanchor{mtk.37}
\csromantocmarknopage{chapter}{4. วิภงฺควคฺค}{mtk.37}
\bookmark[dest={mtk.37},level=0]{4. วิภงฺควคฺค}
\chapterheadpageread{4. วิภงฺควคฺค}
\csromanmatikaanchor{mtk.38}
\csromantocmarkref{subsection}{1. ภทฺเทกรตฺตสุตฺต}{mtk.38}
\bookmark[dest={mtk.38},level=2]{1. ภทฺเทกรตฺตสุตฺต}
\csromanheader{2}{1. ภทฺเทกรตฺตสุตฺต}
\proseitem{272}{\csromanfolio{226}เอวํ เม สุตํ\csromandash{}เอกํ สมยํ ภควา สาวตฺถิยํ วิหรติ เชตวเน อนาถปิณฺฑิกสฺส อาราเม.\csromanspacer{} ตตฺร โข ภควา ภิกฺขู อามนฺเตสิ “ภิกฺขโว”ติ. “ภทนฺเต”ติ เต ภิกฺขู ภควโต ปจฺจสฺโสสุํ.\csromanspacer{} ภควา เอตทโวจ “ภทฺเทกรตฺตสฺส โว ภิกฺขเว อุทฺเทสญฺจ วิภงฺคญฺจ เทเสสฺสามิ, ตํ สุณาถ สาธุกํ มนสิ กโรถ, ภาสิสฺสามี”ติ. “เอวํ ภนฺเต”ติ โข เต ภิกฺขู ภควโต ปจฺจสฺโสสุํ.\csromanspacer{} ภควา เอตทโวจ\csromandash{}}

\csromangathameasure{อตีตํ นานฺวาคเมยฺย, นปฺปฏิกงฺเข อนาคตํ. \\ ยทตีตํ ปหีนํ ตํ, อปฺปตฺตญฺจ อนาคตํ. \\ ปจฺจุปฺปนฺนญฺจ โย ธมฺมํ, ตตฺถ ตตฺถ วิปสฺสติ. \\ อสํหีรํ อสํกุปฺปํ, ตํ วิทฺวา มนุพฺรูหเย. \\ อชฺเชว กิจฺจมาตปฺปํ, โก ชญฺญา มรณํ สุเว. \\ น หิ โน สงฺครํ เตน, มหาเสเนน มจฺจุนา. \\ เอวํ วิหาริํ อาตาปิํ, อโหรตฺตมตนฺทิตํ. \\ ตํ เว “ภทฺเทกรตฺโต”ติ, สนฺโต อาจิกฺขเต มุนิ.}
\csromangathagroup{\csromangathastanza{อตีตํ นานฺวาคเมยฺย, นปฺปฏิกงฺเข อนาคตํ. \\ ยทตีตํ ปหีนํ ตํ, อปฺปตฺตญฺจ อนาคตํ.}\csromangathastanza{ปจฺจุปฺปนฺนญฺจ โย\footnote{ยํ (เนตฺติปาฬิ)} ธมฺมํ, ตตฺถ ตตฺถ วิปสฺสติ. \\ อสํหีรํ\footnote{ยํ (เนตฺติปาฬิ)} อสํกุปฺปํ, ตํ วิทฺวา มนุพฺรูหเย.}\csromangathastanza{อชฺเชว กิจฺจมาตปฺปํ\footnote{กิจฺจํ อาตปฺปํ (สี, ก)}, โก ชญฺญา มรณํ สุเว. \\ น หิ โน สงฺครํ เตน, มหาเสเนน มจฺจุนา.}\csromangathastanza{เอวํ วิหาริํ อาตาปิํ, อโหรตฺตมตนฺทิตํ. \\ ตํ เว “ภทฺเทกรตฺโต”ติ, สนฺโต อาจิกฺขเต มุนิ\footnote{มุนีติ (สี, สฺยา, กํ, อิ)}.}}

\proseitem{273}{กถญฺจ ภิกฺขเว อตีตํ อนฺวาคเมติ. “เอวํรูโป อโหสิํ อตีตมทฺธานนฺ”ติ ตตฺถ นนฺทิํ สมนฺวาเนติ. “เอวํเวทโน อโหสิํ อตีตมทฺธานนฺ”ติ ตตฺถ นนฺทิํ สมนฺวาเนติ. “เอวํสญฺโญ อโหสิํ อตีตมทฺโธนนฺ”ติ ตตฺถ นนฺทิํ สมนฺวาเนติ. “เอวํสงฺขาโร อโหสิํ อตีตมทฺธานนฺ”ติ ตตฺถ นนฺทิํ สมนฺวาเนติ. “เอวํวิญฺญาโณ อโหสิํ อตีตมทฺธานนฺ”ติ ตตฺถ นนฺทิํ สมนฺวาเนติ.\csromanspacer{} เอวํ โข ภิกฺขเว อตีตํ อนฺวาคเมติ.}

\prose{กถญฺจ ภิกฺขเว อตีตํ นานฺวาคเมติ. “เอวํรูโป อโหสิํ อตีตมทฺธานนฺ”ติ ตตฺถ นนฺทิํ น สมนฺวาเนติ. “เอวํเวทโน อโหสิํ อตีตมทฺธานนฺ”ติ ตตฺถ นนฺทิํ น สมนฺวาเนติ. “เอวํสญฺโญ อโหสิํ}

\csromanheader{2}{1. ภทฺเทกรตฺตสุตฺต (131)}
\prose{\csromanfolio{227}อตีตมทฺธานนฺ”ติ ตตฺถ นนฺทิํ น สมนฺวาเนติ. “เอวํสงฺขาโร อโหสิํ อตีตมทฺธานนฺ”ติ ตตฺถ นนฺทิํ น สมนฺวาเนติ. “เอวํวิญฺญาโณ อโหสิํ อตีตมทฺธานนฺ”ติ ตตฺถ นนฺทิํ น สมนฺวาเนติ.\csromanspacer{} เอวํ โข ภิกฺขเว อตีตํ นานฺวาคเมติ.}

\proseitem{274}{กถญฺจ ภิกฺขเว อนาคตํ ปฏิกงฺขติ. “เอวํรูโป สิยํ อนาคตมทฺธานนฺ”ติ ตตฺถ นนฺทิํ สมนฺวาเนติ.\csromanspacer{} เอวํเวทโน สิยํ ฯเปฯ.\csromanspacer{} เอวํสญฺโญ สิยํ.\csromanspacer{} เอวํสงฺขาโร สิยํ. “เอวํวิญฺญาโณ สิยํ อนาคตมทฺธานนฺ”ติ ตตฺถ นนฺทิํ สมนฺวาเนติ.\csromanspacer{} เอวํ โข ภิกฺขเว อนาคตํ ปฏิกงฺขติ.}

\prose{กถญฺจ ภิกฺขเว อนาคตํ นปฺปฏิกงฺขติ. “เอวํรูโป สิยํ อนาคตมทฺธานนฺ”ติ ตตฺถ นนฺทิํ น สมนฺวาเนติ.\csromanspacer{} เอวํเวทโน สิยํ.\csromanspacer{} เอวํสญฺโญ สิยํ.\csromanspacer{} เอวํสงฺขาโร สิยํ. “เอวํวิญฺญาโณ สิยํ อนาคตมทฺธานนฺ”ติ ตตฺถ นนฺทิํ น สมนฺวาเนติ.\csromanspacer{} เอวํ โข ภิกฺขเว อนาคตํ นปฺปฏิกงฺขติ.}

\proseitem{275}{กถญฺจ ภิกฺขเว ปจฺจุปฺปนฺเนสุ ธมฺเมสุ สํหีรติ.\csromanspacer{} อิธ ภิกฺขเว อสฺสุตวา ปุถุชฺชโน อริยานํ อทสฺสาวี อริยธมฺมสฺส อโกวิโท อริยธมฺเม อวินีโต สปฺปุริสานํ อทสฺสาวี สปฺปุริสธมฺมสฺส อโกวิโท สปฺปุริสธมฺเม อวินีโต รูปํ อตฺตโต สมนุปสฺสติ, รูปวนฺตํ วา อตฺตานํ, อตฺตนิ วา รูปํ, รูสฺมิํ วา อตฺตานํ.\csromanspacer{} เวทนํ ฯเปฯ.\csromanspacer{} สญฺญํ.\csromanspacer{} สงฺขาเร.\csromanspacer{} วิญฺญาณํ อตฺตโต สมนุปสฺสติ, วิญฺญาณวนฺตํ วา อตฺตานํ, อตฺตนิ วา วิญฺญาณํ, วิญฺญาณสฺมิํ วา อตฺตานํ.\csromanspacer{} เอวํ โข ภิกฺขเว ปจฺจุนฺเนสุ ธมฺเมสุ สํหีรติ.}

\prose{กถญฺจ ภิกฺขเว ปจฺจุปฺปนฺเนสุ ธมฺเมสุ น สํหีรติ.\csromanspacer{} อิธ ภิกฺขเว สุตวา อริยสาวโก อริยานํ ทสฺสาวี อริยธมฺมสฺส โกวิโท อริยธมฺเม สุวินีโต สปฺปุริสานํ ทสฺสาวี สปฺปุริสธมฺมสฺส โกวิโท สปฺปุริสธมฺเม สุวินีโต น รูปํ อตฺตโต สมนุปสฺสติ, น รูปวนฺตํ วา อตฺตานํ, น อตฺตนิ วา รูปํ, น รูปสฺมิํ วา อตฺตานํ.\csromanspacer{} น เวทนํ.\csromanspacer{} น สญฺญํ.\csromanspacer{} น สงฺขาเร.\csromanspacer{} น วิญฺญาณํ อตฺตโต สมนุปสฺสติ, น วิญฺญาณวนฺตํ วา อตฺตานํ, น อตฺตนิ วา วิญฺญาณํ, น วิญฺญาณสฺมิํ วา อตฺตานํ.\csromanspacer{} เอวํ โข ภิกฺขเว ปจฺจุปฺปนฺเนสุ ธมฺเมสุ น สํหีรติ.}

\csromangathameasure{อตีตํ นานฺวาคเมยฺย, นปฺปฏิกงฺเข อนาคตํ. \\ ยทตีตํ ปหีนํ ตํ, อปฺปตฺตญฺจ อนาคตํ. \\ ปจฺจุปฺปนฺนญฺจ โย ธมฺมํ, ตตฺถ ตตฺถ วิปสฺสติ. \\ อสํหีรํ อสํกุปฺปํ, ตํ วิทฺวา มนุพฺรูหเย. \\ อชฺเชว กิจฺจมาตปฺปํ, โก ชญฺญา มรณํ สุเว. \\ น หิ โน สงฺครํ เตน, มหาเสเนน มจฺจุนา. \\ เอวํ วิหาริํ อาตาปิํ, อโหรตฺตมตนฺทิตํ. \\ ตํ เว “ภทฺเทกรตฺโต”ติ, สนฺโต อาจิกฺขเต มุนีติ.}
\csromangathagroup{\csromangathastanza{อตีตํ นานฺวาคเมยฺย, นปฺปฏิกงฺเข อนาคตํ. \\ ยทตีตํ ปหีนํ ตํ, อปฺปตฺตญฺจ อนาคตํ.}\csromangathastanza{\csromanfolio{228}ปจฺจุปฺปนฺนญฺจ โย ธมฺมํ, ตตฺถ ตตฺถ วิปสฺสติ. \\ อสํหีรํ อสํกุปฺปํ, ตํ วิทฺวา มนุพฺรูหเย.}\csromangathastanza{อชฺเชว กิจฺจมาตปฺปํ, โก ชญฺญา มรณํ สุเว. \\ น หิ โน สงฺครํ เตน, มหาเสเนน มจฺจุนา.}\csromangathastanza{เอวํ วิหาริํ อาตาปิํ, อโหรตฺตมตนฺทิตํ. \\ ตํ เว “ภทฺเทกรตฺโต”ติ, สนฺโต อาจิกฺขเต มุนีติ.}}

\prose{“ภทฺเทกรตฺตสฺส โว ภิกฺขเว อุทฺเทสญฺจ วิภงฺคญฺจ เทเสสฺสามี”ติ อิติ ยํ ตํ วุตฺตํ, อิทเมตํ ปฏิจฺจ วุตฺตนฺติ.}

\prose{อิทมโวจ ภควา.\csromanspacer{} อตฺตมนา เต ภิกฺขู ภควโต ภาสิตํ อภินนฺทุนฺติ.}

\nitthitamruled{ภทฺเทกรตฺตสุตฺตํ นิฏฺฐิตํ ปฐมํ.}
\csromanmatikaanchor{mtk.39}
\csromantocmarkref{subsection}{2. อานนฺทภทฺเทกรตฺตสุตฺต}{mtk.39}
\bookmark[dest={mtk.39},level=2]{2. อานนฺทภทฺเทกรตฺตสุตฺต}
\csromanheader{2}{2. อานนฺทภทฺเทกรตฺตสุตฺต}
\proseitem{276}{เอวํ เม สุตํ\csromandash{}เอกํ สมยํ ภควา สาวตฺถิยํ วิหรติ เชตวเน อนาถปิณฺฑิกสฺส อาราเม.\csromanspacer{} เตน โข ปน สมเยน อายสฺมา อานนฺโท อุปฏฺฐานสาลายํ ภิกฺขูนํ ธมฺมิยา กถาย สนฺทสฺเสติ สมาทเปติ สมุตฺเตเชติ สมฺปหํเสติ, ภทฺเทกรตฺตสฺส อุทฺเทสญฺจ วิภงฺคญฺจ ภาสติ.}

\prose{อถ โข ภควา สายนฺหสมยํ ปฏิสลฺลานา วุฏฺฐิโต เยนุปฏฺฐานสาลา เตนุปสงฺกมิ, อุปสงฺกมิตฺวา ปญฺญตฺเต อาสเน นิสีทิ, นิสชฺช โข ภควา ภิกฺขู อามนฺเตสิ “โก นุ โข ภิกฺขเว อุปฏฺฐานสาลายํ ภิกฺขูนํ ธมฺมิยา กถาย สนฺทสฺเสสิ สมาทเปสิ สมุตฺเตเชสิ สมฺปหํเสสิ, ภทฺเทกรตฺตสฺส อุทฺเทสญฺจ วิภงฺคญฺจ อภาสี”ติ.\csromanspacer{} อายสฺมา ภนฺเต อานนฺโท อุปฏฺฐานสาลายํ ภิกฺขูนํ ธมฺมิยา กถาย สนฺทสฺเสสิ สมาทเปสิ สมุตฺเตเชสิ สมฺปหํเสสิ, ภทฺเทกรตฺตสฺส อุทฺเทสญฺจ วิภงฺคญฺจ อภาสี”ติ.}

\prose{อถ โข ภควา อายสฺมนฺตํ อานนฺทํ อามนฺเตสิ “ยถา กถํ ปน ตฺวํ อานนฺท ภิกฺขูนํ ธมฺมิยา กถาย สนฺทสฺเสสิ สมาทเปสิ สมุตฺเตเชสิ}

\csromanheader{2}{2. อานนฺทภทฺเทกรตฺตสุตฺต (132)}
\prose{\csromanfolio{229}สมฺปหํเสสิ, ภทฺเทกรตฺตสฺส อุทฺเทสญฺจ วิภงฺคญฺจ อภาสี”ติ.\csromanspacer{} เอวํ โข อหํ ภนฺเต ภิกฺขูนํ ธมฺมิยา กถาย สนฺทสฺเสสิํ สมาทเปสิํ สมุตฺเตเชสิํ สมฺปหํเสสิํ, ภทฺเทกรตฺตสฺส อุทฺเทสญฺจ วิภงฺคญฺจ อภาสิํ\csromandash{}}

\csromangathameasure{อตีตํ นานฺวาคเมยฺย, นปฺปฏิกงฺเข อนาคตํ. \\ ยทตีตํ ปหีนํ ตํ, อปฺปตฺตญฺจ อนาคตํ. \\ ปจฺจุปฺปนฺนญฺจ โย ธมฺมํ, ตตฺถ ตตฺถ วิปสฺสติ. \\ อสํหีรํ อสํกุปฺปํ, ตํ วิทฺวา มนุพฺรูหเย. \\ อชฺเชว กิจฺจมาตปฺปํ, โก ชญฺญา มรณํ สุเว. \\ น หิ โน สงฺครํ เตน, มหาเสเนน มจฺจุนา. \\ เอวํ วิหาริํ อาตาปิํ, อโหรตฺตมตนฺทิตํ. \\ ตํ เว “ภทฺเทกรตฺโต”ติ, สนฺโต อาจิกฺขเต มุนิ.}
\csromangathagroup{\csromangathastanza{อตีตํ นานฺวาคเมยฺย, นปฺปฏิกงฺเข อนาคตํ. \\ ยทตีตํ ปหีนํ ตํ, อปฺปตฺตญฺจ อนาคตํ.}\csromangathastanza{ปจฺจุปฺปนฺนญฺจ โย ธมฺมํ, ตตฺถ ตตฺถ วิปสฺสติ. \\ อสํหีรํ อสํกุปฺปํ, ตํ วิทฺวา มนุพฺรูหเย.}\csromangathastanza{อชฺเชว กิจฺจมาตปฺปํ, โก ชญฺญา มรณํ สุเว. \\ น หิ โน สงฺครํ เตน, มหาเสเนน มจฺจุนา.}\csromangathastanza{เอวํ วิหาริํ อาตาปิํ, อโหรตฺตมตนฺทิตํ. \\ ตํ เว “ภทฺเทกรตฺโต”ติ, สนฺโต อาจิกฺขเต มุนิ.}}

\proseitem{277}{กถญฺจ อาวุโส อตีตํ อนฺวาคเมติ. “เอวํรูโป อโหสิํ อตีตมทฺธานนฺ”ติ ตตฺถ นนฺทิํ สมนฺวาเนติ. “เอวํเวทโน อโหสิํ อตีตมทฺธานนฺ”ติ ตตฺถ นนฺทิํ สมนฺวาเนติ. “เอวํสญฺโญ อโหสิํ อตีตมทฺธานนฺ”ติ ตตฺถ นนฺทิํ สมนฺวาเนติ. “เอวํสงฺขาโร อโหสิํ อตีตมทฺธานนฺ”ติ ตตฺถ นนฺทิํ สมนฺวาเนติ. “เอวํวิญฺญาโณ อโหสิํ อตีตมทฺธานนฺ”ติ ตตฺถ นนฺทิํ สมนฺวาเนติ.\csromanspacer{} เอวํ โข อาวุโส อตีตํ อนฺวาคเมติ.}

\prose{กถญฺจ อาวุโส อตีตํ นานฺวาคเมติ. “เอวํรูโป อโหสิํ อตีตมทฺธานนฺ”ติ ตตฺถ นนฺทิํ น สมนฺวาเนติ. “เอวํเวทโน อโหสิํ อตีตมทฺธานนฺ”ติ ตตฺถ นนฺทิํ น สมนฺวาเนติ. “เอวํสญฺโญ อโหสิํ อตีตมทฺธานนฺ”ติ ตตฺถ นนฺทิํ น สมนฺวาเนติ. “เอวํสงฺขาโร อโหสิํ อตีตมทฺธานนฺ”ติ ตตฺถ นนฺทิํ น สมนฺวาเนติ. “เอวํวิญฺญาโณ อโหสิํ อตีตมทฺธานนฺ”ติ ตตฺถ นนฺทิํ น สมนฺวาเนติ.\csromanspacer{} เอวํ โข อาวุโส อตีตํ นานฺวาคเมติ.}

\prose{กถญฺจ อาวุโส อนาคตํ ปฏิกงฺขติ. “เอวํรูโป สิยํ อนาคตมทฺธานนฺ”ติ ตตฺถ นนฺทิํ สมนฺวาเนติ.\csromanspacer{} เอวํเวทโน สิยํ ฯเปฯ.\csromanspacer{} เอวํสญฺโญ สิยํ. \csromanfolio{230}เอวํสงฺขาโร สิยํ. “เอวํวิญฺญาโณ สิยํ อนาคตมทฺธานนฺ”ติ ตตฺถ นนฺทิํ สมนฺวาเนติ.\csromanspacer{} เอวํ โข อาวุโส อนาคตํ ปฏิกงฺขติ.}

\prose{กถญฺจ อาวุโส อนาคตํ นปฺปฏิกงฺขติ. “เอวํรูโป สิยํ อนาคตมทฺธานนฺ”ติ ตตฺถ นนฺทิํ น สมนฺวาเนติ.\csromanspacer{} เอวํเวทโน สิยํ ฯเปฯ.\csromanspacer{} เอวํ สญฺโญ สิยํ.\csromanspacer{} เอวํสงฺขาโร สิยํ. “เอวํวิญฺญาโณ สิยํ อนาคตมทฺธานนฺ”ติ ตตฺถ นนฺทิํ น สมนฺวาเนติ.\csromanspacer{} เอวํ โข อาวุโส อนาคตํ นปฺปฏิกงฺขติ.}

\prose{กถญฺจ อาวุโส ปจฺจุปฺปนฺเนสุ ธมฺเมสุ สํหีรติ.\csromanspacer{} อิธ อาวุโส อสฺสุตวา ปุถุชฺชโน อริยานํ อทสฺสาวี อริยธมฺมสฺส อโกวิโท อริยธมฺเม อวินีโต สปฺปุริสานํ อทสฺสาวี สปฺปุริสธมฺมสฺส อโกวิโท สปฺปุริสธมฺเม อวินีโต รูปํ อตฺตโต สมนุปสฺสติ, รูปวนฺตํ วา อตฺตานํ, อตฺตนิ วา รูปํ, รูปสฺมิํ วา อตฺตานํ.\csromanspacer{} เวทนํ.\csromanspacer{} สญฺญํ.\csromanspacer{} สงฺขาเร.\csromanspacer{} วิญฺญาณํ อตฺตโต สมนุปสฺสติ, วิญฺญาณวนฺตํ วา อตฺตานํ, อตฺตนิ วา วิญฺญาณํ, วิญฺญาณสฺมิํ วา อตฺตานํ.\csromanspacer{} เอวํ โข อาวุโส ปจฺจุปฺปนฺเนสุ ธมฺเมสุ สํหีรติ.}

\prose{กถญฺจ อาวุโส ปจฺจุปฺปนฺเนสุ ธมฺเมสุ น สํหีรติ.\csromanspacer{} อิธ อาวุโส สุตวา อริยสาวโก อริยานํ ทสฺสาวี อริยธมฺมสฺส โกวิโท อริยธมฺเม สุวินีโต สปฺปุริสานํ ทสฺสาวี สปฺปุริสธมฺมสฺส โกวิโท สปฺปุริสธมฺเม สุวินีโต น รูปํ อตฺตโต สมนุปสฺสติ, น รูปวนฺตํ วา อตฺตานํ, น อตฺตนิ วา รูปํ, น รูปสฺมิํ วา อตฺตานํ.\csromanspacer{} น เวทนํ.\csromanspacer{} น สญฺญํ.\csromanspacer{} น สงฺขาเร.\csromanspacer{} น วิญฺญาณํ อตฺตโต สมนุปสฺสติ, น วิญฺญาณวนฺตํ วา อตฺตานํ, น อตฺตนิ วา วิญฺญาณํ, น วิญฺญาณสฺมิํ วา อตฺตานํ.\csromanspacer{} เอวํ โข อาวุโส ปจฺจุปฺปนฺเนสุ ธมฺเมสุ น สํหีรติ.}

\csromangathameasure{อตีตํ นานฺวาคเมยฺย, นปฺปฏิกงฺเข อนาคตํ. \\ ยทตีตํ ปหีนํ ตํ, อปฺปตฺตญฺจ อนาคตํ. \\ ปจฺจุปฺปนฺนญฺจ โย ธมฺมํ, ตตฺถ ตตฺถ วิปสฺสติ. \\ อสํหีรํ อสํกุปฺปํ, ตํ วิทฺวา มนุพฺรูหเย. \\ อชฺเชว กิจฺจมาตปฺปํ, โก ชญฺญา มรณํ สุเว. \\ น หิ โน สงฺครํ เตน, มหาเสเนน มจฺจุนา. \\ เอวํ วิหาริํ อาตาปิํ, อโหรตฺตมตนฺทิตํ. \\ ตํ เว “ภทฺเทกรตฺโต”ติ, สนฺโต อาจิกฺขเต มุนีติ.}
\csromangathagroup{\csromangathastanza{อตีตํ นานฺวาคเมยฺย, นปฺปฏิกงฺเข อนาคตํ. \\ ยทตีตํ ปหีนํ ตํ, อปฺปตฺตญฺจ อนาคตํ.}\csromangathastanza{ปจฺจุปฺปนฺนญฺจ โย ธมฺมํ, ตตฺถ ตตฺถ วิปสฺสติ. \\ อสํหีรํ อสํกุปฺปํ, ตํ วิทฺวา มนุพฺรูหเย.}\csromangathastanza{อชฺเชว กิจฺจมาตปฺปํ, โก ชญฺญา มรณํ สุเว. \\ น หิ โน สงฺครํ เตน, มหาเสเนน มจฺจุนา.}\csromangathastanza{เอวํ วิหาริํ อาตาปิํ, อโหรตฺตมตนฺทิตํ. \\ ตํ เว “ภทฺเทกรตฺโต”ติ, สนฺโต อาจิกฺขเต มุนีติ.}}

\csromanheader{2}{3. มหากจฺจานภทฺเทกรตฺตสุตฺต (133)}
\prose{\csromanfolio{231}เอวํ โข อหํ ภนฺเต ภิกฺขูนํ ธมฺมิยา กถาย สนฺทสฺเสสิํ สมาทเปสิํ สมุตฺเตเชสิํ สมฺปหํเสสิํ, ภทฺเทกรตฺตสฺส อุทฺเทสญฺจ วิภงฺคญฺจ อภาสินฺติ.}

\proseitem{278}{สาธุ สาธุ อานนฺท, สาธุ โข ตฺวํ อานนฺท ภิกฺขูนํ ธมฺมิยา กถาย สนฺเทสฺเสสิ สมาทเปสิ สมุตฺเตเชสิ สมฺปหํเสสิ, ภทฺเทกรตฺตสฺส อุทฺเทสญฺจ วิภงฺคญฺจ อภาสิ\csromandash{}}

\prose{“อตีตํ นานฺวาคเมยฺย ฯเปฯ.}

\csromangathameasure{ตํ เว ‘ภทฺเทกรตฺโต’ติ, สนฺโต อาจิกฺขเต มุนี”ติ.}
\csromangathagroup{\csromangathastanza{ตํ เว ‘ภทฺเทกรตฺโต’ติ, สนฺโต อาจิกฺขเต มุนี”ติ.}}

\prose{กถญฺจ อานนฺท อตีตํ อนฺวาคเมติ ฯเปฯ.\csromanspacer{} เอวํ โข อานนฺท อตีตํ อนฺวาคเมติ.\csromanspacer{} กถญฺจ อานนฺท อตีตํ นานฺวาคเมติ ฯเปฯ.\csromanspacer{} เอวํ โข อานนฺท อตีตํ นานฺวาคเมติ.\csromanspacer{} กถญฺจ อานนฺท อนาคตํ ปฏิกงฺขติ ฯเปฯ.\csromanspacer{} เอวํ โข อานนฺท อนาคตํ ปฏิกงฺขติ.\csromanspacer{} กถญฺจ อานนฺท อนาคตํ นปฺปฏิกงฺขติ ฯเปฯ.\csromanspacer{} เอวํ โข อานนฺท อนาคตํ นปฺปฏิกงฺขติ.\csromanspacer{} กถญฺจ อานนฺท ปจฺจุปฺปนฺเนสุ ธมฺเมสุ สํหีรติ ฯเปฯ.\csromanspacer{} เอวํ โข อานนฺท ปจฺจุปฺปนฺเนสุ ธมฺเมสุ สํหีรติ.\csromanspacer{} กถญฺจ อานนฺท ปจฺจุปฺปนฺเนสุ ธมฺเมสุ น สํหีรติ ฯเปฯ.\csromanspacer{} เอวํ โข อานนฺท ปจฺจุนฺเนสุ ธมฺเมสุ น สํหีรติ.}

\prose{อตีตํ นานฺวาคเมยฺย ฯเปฯ.}

\csromangathameasure{ตํ เว “ภทฺเทกรตฺโต”ติ, สนฺโต อาจิกฺขเต มุนีติ.}
\csromangathagroup{\csromangathastanza{ตํ เว “ภทฺเทกรตฺโต”ติ, สนฺโต อาจิกฺขเต มุนีติ.}}

\prose{อิทมโวจ ภควา.\csromanspacer{} อตฺตมโน อายสฺมา อานนฺโท ภควโต ภาสิตํ อภินนฺทีติ.}

\nitthitamruled{อานนฺทภทฺเทกรตฺตสุตฺตํ นิฏฺฐิตํ ทุติยํ.}
\csromanmatikaanchor{mtk.40}
\csromantocmarkref{subsection}{3. มหากจฺจานภทฺเทกรตฺตสุตฺต}{mtk.40}
\bookmark[dest={mtk.40},level=2]{3. มหากจฺจานภทฺเทกรตฺตสุตฺต}
\csromanheader{2}{3. มหากจฺจานภทฺเทกรตฺตสุตฺต}
\proseitem{279}{เอวํ เม สุตํ\csromandash{}เอกํ สมยํ ภควา ราชคเห วิหรติ ตโปทาราเม.\csromanspacer{} อถ โข อายสฺมา สมิทฺธิ รตฺติยา ปจฺจูสสมยํ ปจฺจุฏฺฐาย เยน ตโปโท\footnote{ตโปทา (สี)} เตนุปสงฺกมิ คตฺตานิ ปริสิญฺจิตุํ.\csromanspacer{} ตโปเท \csromanfolio{232}คตฺตานิ ปริสิญฺจิตฺวา ปจฺจุตฺตริตฺวา เอกจีวโร อฏฺฐาสิ คตฺตานิ ปุพฺพาปยมาโน\footnote{สุกฺขาปยมาโน (ก)}.\csromanspacer{} อถ โข อญฺญตรา เทวตา อภิกฺกนฺตาย รตฺติยา อภิกฺกนฺตวณฺณา เกวลกปฺปํ ตโปทํ โอภาเสตฺวา เยนายสฺมา สมิทฺธิ เตนุปสงฺกมิ, อุปสงฺกมิตฺวา เอกมนฺตํ อฏฺฐาสิ, เอกมนฺตํ ฐิตา โข สา เทวตา อายสฺมนฺตํ สมิทฺธิํ เอตทโวจ “ธาเรสิ ตฺวํ ภิกฺขุ ภทฺเทกรตฺตสฺส อุทฺเทสญฺจ วิภงฺคญฺจา”ติ, น โข อหํ อาวุโส ธาเรมิ ภทฺเทกรตฺตสฺส อุทฺเทสญฺจ วิภงฺคญฺจ, ตฺวํ ปนาวุโส ธาเรสิ ภทฺเทกรตฺตสฺส อุทฺเทสญฺจ วิภงฺคญฺจาติ.\csromanspacer{} อหมฺปิ โข ภิกฺขุ น ธาเรมิ ภทฺเทกรตฺตสฺส อุทฺเทสญฺจ วิภงฺคญฺจ.\csromanspacer{} ธาเรสิ ปน ตฺวํ ภิกฺขุ ภทฺเทกรตฺติโย คาถาติ.\csromanspacer{} น โข อหํ อาวุโส ธาเรมิ ภทฺเทกรตฺติโย คาถาติ.\csromanspacer{} ตฺวํ ปนาวุโส ธาเรสิ ภทฺเทกตฺติโย คาถาติ.\csromanspacer{} อหมฺปิ โข ภิกฺขุ น ธาเรมิ ภทฺเทกรตฺติโย คาถาติ.\csromanspacer{} อุคฺคณฺหาหิ ตฺวํ ภิกฺขุ ภทฺเทกรตฺตสฺส อุทฺเทสญฺจ วิภงฺคญฺจ, ปริยาปุณาหิ ตฺวํ ภิกฺขุ ภทฺเทกรตฺตสฺส อุทฺเทสญฺจ วิภงฺคญฺจ, ธาเรหิ ตฺวํ ภิกฺขุ ภทฺเทกรตฺตสฺส อุทฺเทสญฺจ วิภงฺคญฺจ, อตฺถสํหิโต ภิกฺขุ ภทฺเทกรตฺตสฺส อุทฺเทโส จ วิภงฺโค จ อาทิพฺรหฺมจริยโกติ.\csromanspacer{} อิทมโวจ สา เทวตา, อิทํ วตฺวา ตตฺเถวนฺตรธายิ.}

\proseitem{280}{อถ โข อายสฺมา สมิทฺธิ ตสฺสา รตฺติยา อจฺจเยน เยน ภควา เตนุปสงฺกมิ, อุปสงฺกมิตฺวา ภควนฺตํ อภิวาเทตฺวา เอกมนฺตํ นิสีทิ, เอกมนฺตํ นิสินฺโน โข อายสฺมา สมิทฺธิ ภควนฺตํ เอตทโวจ\csromandash{}}

\prose{อิธาหํ ภนฺเต รตฺติยา ปจฺจูสสมยํ ปจฺจุฏฺฐาย เยน ตโปโท เตนุปสงฺกมิํ คตฺตานิ ปริสิญฺจิตุํ.\csromanspacer{} ตโปเท คตฺตานิ ปริสิญฺจตฺวา ปจฺจุตฺตริตฺวา เอกจีวโร อฏฺฐาสิํ คตฺตานิ ปุพฺพาปยมาโน.\csromanspacer{} อถ โข ภนฺเต อญฺญตรา เทวตา อภิกฺกนฺตาย รตฺติยา อภิกฺกนฺตวณฺณา เกวลกปฺปํ ตโปทํ โอภาเสตฺวา เยนาหํ เตนุปสงฺกมิ, อุปสงฺกมิตฺวา เอกมนฺตํ อฏฺฐาสิ, เอกมนฺตํ ฐิตา โข สา เทวตา มํ เอตทโวจ “ธาเรสิ ตฺวํ ภิกฺขุ ภทฺเทกรตฺตสฺส อุทฺเทสญฺจ วิภงฺคญฺจา”ติ.}

\prose{เอวํ วุตฺเต อหํ ภนฺเต ตํ เทวตํ เอตทโวจํ “น โข อหํ อาวุโส ธาเรมิ ภทฺเทกรตฺตสฺส อุทฺเทสญฺจ วิภงฺคญฺจ, ตฺวํ ปนาวุโส}

\csromanheader{2}{3. มหากจฺจานภทฺเทกรตฺตสุตฺต (133)}
\prose{\csromanfolio{233}ธาเรสิ ภทฺเทกรตฺตสฺส อุทฺเทสญฺจ วิภงฺคญฺจา”ติ.\csromanspacer{} อหมฺปิ โข ภิกฺขุ น ธาเรมิ ภทฺเทกรตฺตสฺส อุทฺเทสญฺจ วิภงฺคญฺจ, ธาเรสิ ปน ตฺวํ ภิกฺขุ ภทฺเทกรตฺติโย คาถาติ.\csromanspacer{} น โข อหํ อาวุโส ธาเรมิ ภทฺเทกรตฺติโย คาถาติ, ตฺวํ ปนาวุโส ธาเรสิ ภทฺเทกรตฺติโย คาถาติ.\csromanspacer{} อหมฺปิ โข ภิกฺขุ น ธาเรมิ ภทฺเทกรตฺติโย คาถาติ, อุคฺคณฺหาหิ ตฺวํ ภิกฺขุ ภทฺเทกรตฺตสฺส อุทฺเทสญฺจ วิภงฺคญฺจ, ปริยาปุณาหิ ตฺวํ ภิกฺขุ ภทฺเทกรตฺตสฺส อุทฺเทสญฺจ วิภงฺคญฺจ, ธาเรหิ ตฺวํ ภิกฺขุ ภทฺเทกรตฺตสฺส อุทฺเทสญฺจ วิภงฺคญฺจ, อตฺถสํหิโต ภิกฺขุ ภทฺเทกรตฺตสฺส อุทฺเทโส จ วิภงฺโค จ อาทิพฺรจริยโกติ.\csromanspacer{} อิทมโวจ ภนฺเต สา เทวตา, อิทํ วตฺวา ตตฺเถวนฺตรธายิ.\csromanspacer{} สาธุ เม ภนฺเต ภควา ภทฺเทกรตฺตสฺส อุทฺเทสญฺจ วิภงฺคญฺจ เทเสตูติ.\csromanspacer{} เตน หิ ภิกฺขู สุณาหิ สาธุกํ มนสิ กโรหิ, ภาสิสฺสามีติ. “เอวํ ภนฺเต”ติ โข อายสฺมา สมิทฺธิ ภควโต ปจฺจสฺโสสิ.\csromanspacer{} ภควา เอตทโวจ\csromandash{}}

\csromangathameasure{“อตีตํ นานฺวาคเมยฺย, นปฺปฏิกงฺข อนาคตํ. \\ ยทตีตํ ปหีนํ ตํ, อปฺปตฺตญฺจ อนาคตํ. \\ ปจฺจุปฺปนฺนญฺจ โย ธมฺมํ, ตตฺถ ตตฺถ วิปสฺสติ. \\ อสํหีรํ อสํกุปฺปํ, ตํ วิทฺวา มนุพฺรูหเย. \\ อชฺเชว กิจฺจมาตปฺปํ, \\ โก ชญฺญา มรณํ สุเว. \\ น หิ โน สงฺครํ เตน, มหาเสเนน มจฺจุนา, \\ เอวํวิหาริํ อาตาปิํ, \\ อโหรตฺตมตนฺทิตํ. ตํ เว ‘ภทฺเทกรตฺโต’ติ, \\ สนฺโต อาจิกฺขเต มุนี”ติ.}
\csromangathagroup{\csromangathastanza{“อตีตํ นานฺวาคเมยฺย, นปฺปฏิกงฺข อนาคตํ. \\ ยทตีตํ ปหีนํ ตํ, อปฺปตฺตญฺจ อนาคตํ.}\csromangathastanza{ปจฺจุปฺปนฺนญฺจ โย ธมฺมํ, ตตฺถ ตตฺถ วิปสฺสติ. \\ อสํหีรํ อสํกุปฺปํ, ตํ วิทฺวา มนุพฺรูหเย.}\csromangathastanza{อชฺเชว กิจฺจมาตปฺปํ, \\ โก ชญฺญา มรณํ สุเว. \\ น หิ โน สงฺครํ เตน, มหาเสเนน มจฺจุนา, \\ เอวํวิหาริํ อาตาปิํ, \\ อโหรตฺตมตนฺทิตํ. ตํ เว ‘ภทฺเทกรตฺโต’ติ, \\ สนฺโต อาจิกฺขเต มุนี”ติ.}}

\prose{อิทมโวจ ภควา, อิทํ วตฺวาน สุคโต อุฏฺฐายาสนา วิหารํ ปาวิสิ.\csromanspacer{} อถ โข เตสํ ภิกฺขูนํ อจิรปกฺกนฺตสฺส ภควโต เอตทโหสิ “อิทํ โข โน อาวุโส ภควา สํขิตฺเตน อุทฺเทสํ อุทฺทิสิตฺวา วิตฺถาเรน อตฺถํ อวิภชิตฺวา อุฏฺฐายาสนา วิหารํ ปวิฏฺโฐ}

\csromangathameasure{‘อตีตํ นานฺวาคเมยฺย, นปฺปฏิกงฺเข อนาคตํ. \\ ยทตีตํ ปหีนํ ตํ, อปฺปตฺตญฺจ อนาคตํ. \\ ปจฺจุปฺปนฺนญฺจ โย ธมฺมํ, ตตฺถ ตตฺถ วิปสฺสติ. \\ อสํหีรํ อสํกุปฺปํ, ตํ วิทฺวา มนุพฺรูหเย. \\ อชฺเชว กิจฺจมาตปฺปํ, โก ชญฺญา มรณํ สุเว. \\ น หิ โน สงฺครํ เตน, มหาเสเนน มจฺจุนา. \\ เอวํ วิหาริํ อาตาปิํ, อโหรตฺตมตนฺทิตํ. \\ ตํ เว ภทฺเทกรตฺโตติ, สนฺโต อาจิกฺขเต มุนี’ติ.}
\csromangathagroup{\csromangathastanza{‘อตีตํ นานฺวาคเมยฺย, นปฺปฏิกงฺเข อนาคตํ. \\ ยทตีตํ ปหีนํ ตํ, อปฺปตฺตญฺจ อนาคตํ.}\csromangathastanza{ปจฺจุปฺปนฺนญฺจ โย ธมฺมํ, ตตฺถ ตตฺถ วิปสฺสติ. \\ อสํหีรํ อสํกุปฺปํ, ตํ วิทฺวา มนุพฺรูหเย.}\csromangathastanza{\csromanfolio{234}อชฺเชว กิจฺจมาตปฺปํ, โก ชญฺญา มรณํ สุเว. \\ น หิ โน สงฺครํ เตน, มหาเสเนน มจฺจุนา.}\csromangathastanza{เอวํ วิหาริํ อาตาปิํ, อโหรตฺตมตนฺทิตํ. \\ ตํ เว ภทฺเทกรตฺโตติ, สนฺโต อาจิกฺขเต มุนี’ติ.}}

\prose{โก นุ โข อิมสฺส ภควตา สํขิตฺเตน อุทฺเทสสฺส อุทฺทิฏฺฐสฺส วิตฺถาเรน อตฺถํ อวิภตฺตสฺส วิตฺถาเรน อตฺถํ วิภเชยฺยา”ติ.}

\prose{อถ โข เตสํ ภิกฺขูนํ เอตทโหสิ “อยํ โข อายสฺมา มหากจฺจาโน สตฺถุ เจว สํวณฺณิโต, สมฺภาวิโต จ วิญฺญูนํ สพฺรหฺมจารีนํ, ปโหติ จายสฺมา มหากจฺจาโน อิมสฺส ภควตา สํขิตฺเตน อุทฺเทสสฺส อุทฺทิฏฺฐสฺส วิตฺถาเรน อตฺถํ อวิภตฺตสฺส วิตฺถาเรน อตฺถํ วิภชิตุํ, ยํนูน มยํ เยนายสฺมา มหากจฺจาโน เตนุปสงฺกเมยฺยาม, อุปสงฺกมิตฺวา อายสฺมนฺตํ มหากจฺจานํ เอตมตฺถํ ปฏิปุจฺเฉยฺยามา”ติ.}

\proseitem{281}{อถ โข เต ภิกฺขู เยนายสฺมา มหากจฺจาโน เตนุปสงฺกมิํสุ, อุปสงฺกมิตฺวา อายสฺมตา มหากจฺจาเนน สทฺธิํ สมฺโมทิํสุ, สมฺโมทนียํ กถํ สารณียํ วีติสาเรตฺวา เอกมนฺตํ นิสีทิํสุ, เอกมนฺตํ นิสินฺนา โข เต ภิกฺขู อายสฺมนฺตํ มหากจฺจานํ เอตทโวจุํ\csromandash{}อิทํ โข โน อาวุโส กจฺจาน ภควา สํขิตฺเตน อุทฺเทสํ อุทฺทิสิตฺวา วิตฺถาเรน อตฺถํ อวิภชิตฺวา อุฏฺฐายาสนา วิหารํ ปวิฏฺโฐ}

\prose{“อตีตํ นานฺวาคเมยฺย ฯเปฯ.}

\csromangathameasure{ตํ เว ภทฺเทกรตฺโตติ, สนฺโต อาจิกฺขเต มุนี”ติ.}
\csromangathagroup{\csromangathastanza{ตํ เว ภทฺเทกรตฺโตติ, สนฺโต อาจิกฺขเต มุนี”ติ.}}

\prose{เตสํ โน อาวุโส กจฺจาน อมฺหากํ อจิรปกฺกนฺตสฺส ภควโต เอตทโหสิ “อิทํ โข โน อาวุโส ภควา สํขิตฺเตน อุทฺเทสํ อุทฺทิสิตฺวา วิตฺถาเรน อตฺถํ อวิภชิตฺวา อุฏฺฐายาสนา วิหารํ ปวิฏฺโฐ}

\prose{‘อตีตํ นานฺวาคเมยฺย ฯเปฯ.}

\csromangathameasure{ตํ เว ภทฺเทกรตฺโตติ, สนฺโต อาจิกฺขเต มุนี’ติ.}
\csromangathagroup{\csromangathastanza{ตํ เว ภทฺเทกรตฺโตติ, สนฺโต อาจิกฺขเต มุนี’ติ.}}

\prose{โก นุ โข อิมสฺส ภควตา สํขิตฺเตน อุทฺเทสสฺส อุทฺทิฏฺฐสฺส วิตฺถาเรน อตฺถํ อวิภตฺตสฺส วิตฺถาเรน อตฺถํ วิภเชยฺยา”ติ.\csromanspacer{} เตสํ}

\csromanheader{2}{3. มหากจฺจานภทฺเทกรตฺตสุตฺต (133)}
\prose{\csromanfolio{235}โน อาวุโส กจฺจาน อมฺหากํ เอตทโหสิ “อยํ โข อายสฺมา มหากจฺจาโน สตฺถุ เจว สํวณฺณิโต, สมฺภาวิโต จ วิญฺญูนํ สพฺรหฺมจารีนํ, ปโหติ จายสฺมา มหากจฺจาโน อิมสฺส ภควตา สํขิตฺเตน อุทฺเทสสฺส อุทฺทิฏฺฐสฺส วิตฺถาเรน อตฺถํ อวิภตฺตสฺส วิตฺถาเรน อตฺถํ วิภชิตุํ, ยํนูน มยํ เยนายสฺมา มหากจฺจาโน เตนุปสงฺกเมยฺยาม, อุปสงฺกมิตฺวา อายสฺมนฺตํ มหากจฺจานํ เอตมตฺถํ ปฏิปุจฺเฉยฺยามา”ติ.\csromanspacer{} วิภชตายสฺมา มหากจฺจาโนติ.}

\prose{เสยฺยถาปิ อาวุโส ปุริโส สารตฺถิโก สารคเวสี สารปริเยสนํ จรมาโน มหโต รุกฺขสฺส ติฏฺฐโต สารวโต อติกฺกมฺเมว มูลํ อติกฺกมฺม ขนฺธํ สาขาปลาเส สารํ ปริเยสิตพฺพํ มญฺเญยฺย.\csromanspacer{} เอวํสมฺปทมิทํ อายสฺมนฺตานํ, สตฺถริ สมฺมุขีภูเต ตํ ภควนฺตํ อติสิตฺวา อเมฺห เอตมตฺถํ ปฏิปุจฺฉิตพฺพํ มญฺญถ\footnote{มญฺเญถ (อิ)}, โส หาวุโส ภควา ชานํ ชานาติ, ปสฺสํ ปสฺสติ, จกฺขุภูโต ญาณภูโต ธมฺมภูโต พฺรหฺมภูโต วตฺตา ปวตฺตา อตฺถสฺส นินฺเนตา อมตสฺส ทาตา ธมฺมสฺสามี ตถาคโต, โส เจว ปเนตสฺส กาโล อโหสิ, ยํ ภควนฺตํเยว เอตมตฺถํ ปฏิปุจฺเฉยฺยาถ, ยถา โว ภควา พฺยากเรยฺย, ตถา นํ ธาเรยฺยาถาติ.}

\prose{อทฺธาวุโส กจฺจาน ภควา ชานํ ชานาติ, ปสฺสํ ปสฺสติ, จกฺขุภูโต ญาณภูโต ธมฺมภูโต พฺรหฺมภูโต วตฺตา ปวตฺตา อตฺถสฺส นินฺเนตา อมตสฺส ทาตา ธมฺมสฺสามี ตถาคโต, โส เจว ปเนตสฺส กาโล อโหสิ, ยํ ภควนฺตํเยว เอตมตฺถํ ปฏิปุจฺเฉยฺยาม, ยถา โน ภควา พฺยากเรยฺย, ตถา นํ ธาเรยฺยาม, อปิ จายสฺมา มหากจฺจาโน สตฺถุ เจว สํวณฺณิโต, สมฺโภวิโต จ วิญฺญูนํ สพฺรหฺมจารีนํ, ปโหติ จายสฺมา มหากจฺจาโน อิมสฺส ภควตา สํขิตฺเตน อุทฺเทสสฺส อุทฺทิฏฺฐสฺส วิตฺถาเรน อตฺถํ อวิภตฺตสฺส วิตฺถาเรน อตฺถํ วิภชิตุํ, วิภชตายสฺมา มหากจฺจาโน อครุํ กริตฺวาติ\footnote{อครุกริตฺวา (สี, สฺยา, กํ, อิ)}.}

\prose{เตน หาวุโส สุณาถ สาธุกํ มนสิ กโรถ, ภาสิสฺสามีติ. “เอวมาวุโส”ติ โข เต ภิกฺขู อายสฺมโต มหากจฺจานสฺส ปจฺจสฺโสสุํ.\csromanspacer{} อายสฺมา มหากจฺจาโน เอตทโวจ\csromandash{}}

\prose{\csromanfolio{236}ยํ โข โน อาวุโส ภควา สํขิตฺเตน อุทฺเทสํ อุทฺทิสิตฺวา วิตฺถาเรน อตฺถํ อวิภชิตฺวา อุฏฺฐายาสนา วิหารํ ปวิฏฺโฐ}

\prose{“อตีตํ นานฺวาคเมยฺย ฯเปฯ.}

\csromangathameasure{ตํ เว ภทฺเทกรตฺโตติ, สนฺโต อาจิกฺขเต มุนี”ติ.}
\csromangathagroup{\csromangathastanza{ตํ เว ภทฺเทกรตฺโตติ, สนฺโต อาจิกฺขเต มุนี”ติ.}}

\prose{อิมสฺส โข อหํ อาวุโส ภควตา สํขิตฺเตน อุทฺเทสสฺส อุทฺทิฏฺฐสฺส วิตฺถาเรน อตฺถํ อวิภตฺตสฺส เอวํ วิตฺถาเรน อตฺถํ อาชานามิ.}

\proseitem{282}{กถญฺจ อาวุโส อตีตํ อนฺวาคเมติ.\csromanspacer{} อิติ เม จกฺขุ อโหสิ อตีตมทฺธานํ, อิติ รูปาติ, ตตฺถ ฉนฺทราคปฺปฏิพทฺธํ\footnote{ฉนฺทราคปฺปฏิพนฺธํ (ก)} โหติ วิญฺญาณํ, ฉนฺทราคปฺปฏิพทฺธตฺตา วิญฺญาณสฺส ตทภินนฺทติ, ตทภินนฺทนฺโต อตีตํ อนฺวาคเมติ.\csromanspacer{} อิติ เม โสตํ อโหสิ อตีตมทฺธานํ, อิติ สทฺทาติ ฯเปฯ.\csromanspacer{} อิติ เม ฆานํ อโหสิ อตีตมทฺธานํ, อิติ คนฺธาติ.\csromanspacer{} อิติ เม ชิวฺหา อโหสิ อตีตมทฺธานํ, อิติ รสาติ.\csromanspacer{} อิติ เม กาโย อโหสิ อตีตมทฺธานํ, อิติ โผฏฺฐพฺพาติ.\csromanspacer{} อิติ เม มโน อโหสิ อตีตมทฺธานํ, อิติ ธมฺมาติ, ตตฺถ ฉนฺทราคปฺปฏิพทฺธํ โหติ วิญฺญาณํ, ฉนฺทราคปฺปฏิพทฺธตฺตา วิญฺญาณสฺส ตทภินนฺทติ, ตทภินนฺทนฺโต อตีตํ อนฺวาคเมติ.\csromanspacer{} เอวํ โข อาวุโส อตีตํ อนฺวาคเมติ.}

\prose{กถญฺจ อาวุโส อตีตํ นานฺวาคเมติ.\csromanspacer{} อิติ เม จกฺขุ อโหสิ อตีตมทฺธานํ, อิติ รูปาติ, ตตฺถ น ฉนฺทราคปฺปฏิพทฺธํ โหติ วิญฺญาณํ, น ฉนฺทราคปฺปฏิพทฺธตฺตา วิญฺญาณสฺส น ตทภินนฺทติ, น ตทภินนฺโท อตีตํ นานฺวาคเมติ.\csromanspacer{} อิติ เม โสตํ อโหสิ อตีตมทฺธานํ, อิติ สทฺทาติ ฯเปฯ.\csromanspacer{} อิติ เม ฆานํ อโหสิ อตีตมทฺธานํ, อิติ คนฺธาติ.\csromanspacer{} อิติ เม ชิวฺหา อโหสิ อตีตมทฺธานํ, อิติ รสาติ.\csromanspacer{} อิติ เม กาโย อโหสิ อตีตมทฺธานํ, อิติ โผฏฺฐพฺพาติ.\csromanspacer{} อิติ เม มโน อโหสิ อตีตมทฺธานํ, อิติ ธมฺมาติ, ตตฺถ น ฉนฺทราคปฺปฏิพทฺธํ โหติ วิญฺญาณํ, น ฉนฺทราคปฺปฏิพทฺธตฺตา วิญฺญาณสฺส น ตทภินนฺทติ, น ตทภินนฺทนฺโต อตีตํ นานฺวาคเมติ.\csromanspacer{} เอวํ โข อาวุโส อตีตํ นานฺวาคเมติ.}

\csromanheader{2}{3. มหากจฺจานภทฺเทกรตฺตสุตฺต (133)}
\proseitem{283}{\csromanfolio{237}กถญฺจ อาวุโส อนาคตํ ปฏิกงฺขติ.\csromanspacer{} อิติ เม จกฺขุ สิยา อนาคตมทฺธานํ, อิติ รูปาติ อปฺปฏิลทฺธสฺส ปฏิลาภาย จิตฺตํ ปณิทหติ, เจตโส ปณิธานปจฺจยา ตทภินนฺทติ, ตทภินนฺทนฺโต อนาคตํ ปฏิกงฺขติ.\csromanspacer{} อิติ เม โสตํ สิยา อนาคตมทฺธานํ, อิติ สทฺธาติ ฯเปฯ.\csromanspacer{} อิติ เม ฆานํ สิยา อนาคตมทฺธานํ, อิติ คนฺธาติ.\csromanspacer{} อิติ เม ชิวฺหา สิยา อนาคตมทฺธานํ, อิติ รสาติ.\csromanspacer{} อิติ เม กาโย สิยา อนาคตมทฺธานํ, อิติ โผฏฺฐพฺพาติ.\csromanspacer{} อิติ เม มโน สิยา อนาคตมทฺธานํ, อิติ ธมฺมาติ อปฺปฏิลทฺธสฺส ปฏิลาภาย จิตฺตํ ปณิทหติ, เจตโส ปณิธานปจฺจยา ตทภินนฺทติ, ตทภินนฺทนฺโต อนาคตํ ปฏิกงฺขติ.\csromanspacer{} เอวํ โข อาวุโส อนาคตํ ปฏิกงฺขติ.}

\prose{กถญฺจ อาวุโส อนาคตํ นปฺปฏิกงฺขติ.\csromanspacer{} อิติ เม จกฺขุ สิยา อนาคตมทฺธานํ, อิติ รูปาติ อปฺปฏิลทฺธสฺส ปฏิลาภาย จิตฺตํ นปฺปณิทหติ, เจตโส อปฺปณิธานปจฺจยา น ตทภินนฺทติ, น ตทภินนฺทนฺโต อนาคตํ นปฺปฏิกงฺขติ.\csromanspacer{} อิติ เม โสตํ สิยา อนาคตมทฺธานํ, อิติ สทฺทาติ ฯเปฯ.\csromanspacer{} อิติ เม ฆานํ สิยา อนาคตมทฺธานํ, อิติ คนฺธาติ.\csromanspacer{} อิติ เม ชิวฺหา สิยา อนาคตมทฺธานํ, อิติ รสาติ.\csromanspacer{} อิติ เม กาโย สิยา อนาคตมทฺธานํ, อิติ โผฏฺฐพฺพาติ.\csromanspacer{} อิติ เม มโน สิยา อนาคตมทฺธานํ, อิติ ธมฺมาติ อปฺปฏิลทฺธสฺส ปฏิลาภาย จิตฺตํ นปฺปณิทหติ, เจตโส อปฺปณิธานปจฺจยา น ตทภินนฺทติ, น ตทภินนฺทนฺโต อนาคตํ นปฺปฏิกงฺขติ.\csromanspacer{} เอวํ โข อาวุโส อนาคตํ นปฺปฏิกงฺขติ.}

\proseitem{284}{กถญฺจ อาวุโส ปจฺจุปฺปนฺเนสุ ธมฺเมสุ สํหีรติ.\csromanspacer{} ยญฺจาวุโส จกฺขุ เย จ รูปา, อุภยเมตํ ปจฺจุปฺปนฺนํ, ตสฺมิํ เจ ปจฺจุปฺปนฺเน ฉนฺทราคปฺปฏิพทฺธํ โหติ วิญฺญาณํ, ฉนฺทราคปฺปฏิพทฺธตฺตา วิญฺญาณสฺส ตทภินนฺทติ, ตทภินนฺทนฺโต ปจฺจุปฺปนฺเนสุ ธมฺเมสุ สํหีรติ.\csromanspacer{} ยญฺจาวุโส โสตํ เย จ สทฺทา ฯเปฯ.\csromanspacer{} ยญฺจาวุโส ฆานํ เย จ คนฺธา.\csromanspacer{} ยา จาวุโส ชิวฺหา เย จ รสา.\csromanspacer{} โย จาวุโส กาโย เย จ โผฏฺฐพฺพา.\csromanspacer{} โย จาวุโส มโน เย จ ธมฺมา, อุภยเมตํ ปจฺจุปฺปนฺนํ, ตสฺมิํ เจ ปจฺจุปฺปนฺเน ฉนฺทราคปฺปฏิพทฺธํ โหติ วิญฺญาณํ, ฉนฺทราคปฺปฏิพทฺธตฺตา วิญฺญาณสฺส ตทภินนฺทติ, ตทภินนฺทนฺโต ปจฺจุปฺปนฺเนสุ ธมฺเมสุ สํหีรติ.\csromanspacer{} เอวํ โข อาวุโส ปจฺจุปฺปนฺเนสุ ธมฺเมสุ สํหีรติ.}

\prose{\csromanfolio{238}กถญฺจ อาวุโส ปจฺจุปฺปนฺเนสุ ธมฺเมสุ น สํหีรติ.\csromanspacer{} ยญฺจาวุโส จกฺขุ เย จ รูปา, อุภยเมตํ ปจฺจุปฺปนฺนํ, ตสฺมิํ เจ ปจฺจุปฺปนฺเน น ฉนฺทราคปฺปฏิพทฺธํ โหติ วิญฺญาณํ, น ฉนฺทราคปฺปฏิพทฺธตฺตา วิญฺญาณสฺส น ตทภินนฺทติ, น ตทภินนฺทนฺโต ปจฺจุปฺปนฺเนสุ ธมฺเมสุ น สํหีรติ.\csromanspacer{} ยญฺจาวุโส โสตํ เย จ สทฺทา ฯเปฯ.\csromanspacer{} ยญฺจาวุโส ฆานํ เย จ คนฺธา.\csromanspacer{} ยา จาวุโส ชิวฺหา เย จ รสา.\csromanspacer{} โย จาวุโส กาโย เย จ โผฏฺฐพฺพา.\csromanspacer{} โย จาวุโส มโน เย จ ธมฺมา, อุภยเมตํ ปจฺจุปฺปนฺนํ, ตสฺมิํ เจ ปจฺจุปฺปนฺเน น ฉนฺทราคปฺปฏิพทฺธํ โหติ วิญฺญาณํ, น ฉนฺทราคปฺปฏิพทฺธตฺตา วิญฺญาณสฺส น ตทภินนฺทติ, น ตทภินนฺทนฺโต ปจฺจุปฺปนฺเนสุ ธมฺเมสุ น สํหีรติ.\csromanspacer{} เอวํ โข อาวุโส ปจฺจุปฺปนฺเนสุ ธมฺเมสุ น สํหีรติ.}

\proseitem{285}{ยํ โข โน อาวุโส ภควา สํขิตฺเตน อุทฺเทสํ อุทฺทิสิตฺวา วิตฺถาเรน อตฺถํ อวิภชิตฺวา อุฏฺฐายาสนา วิหารํ ปวิฏฺโฐ}

\prose{“อตีตํ นานฺวาคเมยฺย ฯเปฯ.}

\csromangathameasure{ตํ เว ภทฺเทกรตฺโตติ, สนฺโต อาจิกฺขเต มุนี”ติ.}
\csromangathagroup{\csromangathastanza{ตํ เว ภทฺเทกรตฺโตติ, สนฺโต อาจิกฺขเต มุนี”ติ.}}

\prose{อิมสฺส โข อหํ อาวุโส ภควตา สํขิตฺเตน อุทฺเทสสฺส อุทฺทิฏฺฐสฺส วิตฺถาเรน อตฺถํ อวิภตฺตสฺส เอวํ วิตฺถาเรน อตฺถํ อาชานามิ, อากงฺขมานา จ ปน ตุเมฺห อายสฺมนฺโต ภควนฺตํ เยว อุปสงฺกมิตฺวา เอตมตฺถํ ปฏิปุจฺเฉยฺยาถ, ยถา โว ภควา พฺยากโรติ ตถา นํ ธาเรยฺยาถาติ.}

\prose{อถ โข เต ภิกฺขู อายสฺมโต มหากจฺจานสฺส ภาสิตํ อภินนฺทิตฺวา อนุโมทิตฺวา อุฏฺฐายาสนา เยน ภควา เตนุปสงฺกมิํสุ, อุปสงฺกมิตฺวา ภควนฺตํ อภิวาเทตฺวา เอกมนฺตํ นิสีทิํสุ, เอกมนฺตํ นิสินฺนา โข เต ภิกฺขู ภควนฺตํ เอตทโวจุํ\csromandash{}ยํ โข โน ภนฺเต ภควา สํขิตฺเตน อุทฺเทสํ อุทฺทิสิตฺวา วิตฺถาเรน อตฺถํ อวิภชิตฺวา อุฏฺฐายาสนา วิหารํ ปวิฏฺโฐ}

\prose{“อตีตํ นานฺวาคเมยฺย ฯเปฯ.}

\csromangathameasure{ตํ เว ภทฺเทกรตฺโตติ, สนฺโต อาจิกฺขเต มุนี”ติ.}
\csromangathagroup{\csromangathastanza{ตํ เว ภทฺเทกรตฺโตติ, สนฺโต อาจิกฺขเต มุนี”ติ.}}

\prose{เตสํ โน ภนฺเต อมฺหากํ อจิรปกฺกนฺตสฺส ภควโต เอตทโหสิ “อิทํ โข โน อาวุโส ภควา สํขิตฺเตน อุทฺเทสํ อุทฺทิสิตฺวา วิตฺถาเรน อตฺถํ อวิภชิตฺวา อุฏฺฐายาสนา วิหารํ ปวิฏฺโฐ}

\csromanheader{2}{3. มหากจฺจานภทฺเทกรตฺตสุตฺต (133)}
\csromangathameasure{‘อตีตํ นานฺวาคเมยฺย, นปฺปฏิกงฺเข อนาคตํ, \\ ยทตีตํ ปหีนํ ตํ, \\ อปฺปตฺตญฺจ อนาคตํ. \\ ปจฺจุปฺปนฺนญฺจ โย ธมฺมํ, \\ ตตฺถ ตตฺถ วิปสฺสติ. อสํหีรํ อสํกุปฺปํ, \\ ตํ วิทฺวา มนุพฺรูหเย. อชฺเชว กิจฺจมาตปฺปํ, \\ โก ชญฺญา มรณํ สุเว. น หิ โน สงฺครํ เตน, \\ มหาเสเนน มจฺจุนา. เอวํวิหาริํ อาตาปิํ, \\ อโหรตฺตมตนฺทิตํ. ตํ เว ภทฺเทกรตฺโตติ, \\ สนฺโต อาจิกฺขเต มุนี’ติ.}
\csromangathagroup{\csromangathastanza{\csromanfolio{239}‘อตีตํ นานฺวาคเมยฺย, นปฺปฏิกงฺเข อนาคตํ, \\ ยทตีตํ ปหีนํ ตํ, \\ อปฺปตฺตญฺจ อนาคตํ. \\ ปจฺจุปฺปนฺนญฺจ โย ธมฺมํ, \\ ตตฺถ ตตฺถ วิปสฺสติ. อสํหีรํ อสํกุปฺปํ, \\ ตํ วิทฺวา มนุพฺรูหเย. อชฺเชว กิจฺจมาตปฺปํ,}\csromangathastanza{โก ชญฺญา มรณํ สุเว. น หิ โน สงฺครํ เตน, \\ มหาเสเนน มจฺจุนา. เอวํวิหาริํ อาตาปิํ,}\csromangathastanza{อโหรตฺตมตนฺทิตํ. ตํ เว ภทฺเทกรตฺโตติ, \\ สนฺโต อาจิกฺขเต มุนี’ติ.}}

\prose{โก นุ โข อิมสฺส ภควตา สํขิตฺเตน อุทฺเทสสฺส อุทฺทิฏฺฐสฺส วิตฺถาเรน อตฺถํ อวิภตฺตสฺส วิตฺถาเรน อตฺถํ วิภเชยฺยา”ติ.\csromanspacer{} เตสํ โน ภนฺเต อมฺหากํ เอตทโหสิ “อยํ โข อายสฺมา มหากจฺจาโน สตฺถุ เจว สํวณฺณิโต, สมฺภาวิโต จ วิญฺญูนํ สพฺรหฺมจารีนํ, ปโหติ จายสฺมา มหากจฺจาโน อิมสฺส ภควโต สํขิตฺเตน อุทฺเทสสฺส อุทฺทิฏฺฐสฺส วิตฺถาเรน อตฺถํ อวิภตฺตสฺส วิตฺถาเรน อตฺถํ วิภชิตุํ.\csromanspacer{} ยํนูน มยํ เยนายสฺมา มหากจฺจาโน เตนุปสงฺกเมยฺยาม, อุปสงฺกมิตฺวา อายสฺมนฺตํ มหากจฺจานํ เอตมตฺถํ ปฏิปุจฺเฉยฺยามา”ติ.\csromanspacer{} อถ โข มยํ ภนฺเต เยนายสฺมา มหากจฺจาโน เตนุปสงฺกมิมฺห, อุปสงฺกมิตฺวา อายสฺมนฺตํ มหากจฺจานํ เอตมตฺถํ ปฏิปุจฺฉิมฺห, เตสํ โน ภนฺเต อายสฺมตา มหากจฺจาเนน อิเมหิ อากาเรหิ อิเมหิ ปเทหิ อิเมหิ พฺยญฺชเนหิ อตฺโถ วิภตฺโตติ.}

\prose{ปณฺฑิโต ภิกฺขเว มหากจฺจาโน, มหาปญฺโญ ภิกฺขเว มหากจฺจาโน, มํ เจปิ ตุเมฺห ภิกฺขเว เอตมตฺถํ ปฏิปุจฺเฉยฺยาถ, อหมฺปิ ตํ เอวเมวํ พฺยากเรยฺยํ.\csromanspacer{} ยถา ตํ มหากจฺจาเนน พฺยากตํ, เอโส เจเวตสฺส อตฺโถ, เอวญฺจ นํ ธาเรถาติ.}

\prose{อิทมโวจ ภควา.\csromanspacer{} อตฺตมนา เต ภิกฺขู ภควโต ภาสิตํ อภินนฺทุนฺติ.}

\nitthitam{มหากจฺจานภทฺเทกรตฺตสุตฺตํ นิฏฺฐิตํ ตติยํ.}
\csromanmatikaanchor{mtk.41}
\csromantocmarkref{subsection}{4. โลมสกงฺคิยภทฺเทกรตฺตสุตฺต}{mtk.41}
\bookmark[dest={mtk.41},level=2]{4. โลมสกงฺคิยภทฺเทกรตฺตสุตฺต}
\csromanheader{2}{4. โลมสกงฺคิยภทฺเทกรตฺตสุตฺต}
\proseitem{286}{\csromanfolio{240}เอวํ เม สุตํ\csromandash{}เอกํ สมยํ ภควา สาวตฺถิยํ วิหรติ เชตวเน อนาถปิณฺฑิกสฺส อาราเม.\csromanspacer{} เตน โข ปน สมเยน อายสฺมา โลมสกงฺคิโย\footnote{โลมสกกงฺคิโย (ฏีกา)} สกฺเกสุ วิหรติ กปิลวตฺถุสฺมิํ นิโคฺรธาราเม.\csromanspacer{} อถ โข จนฺทโน เทวปุตฺโต อภิกฺกนฺตาย รตฺติยา อภิกฺกนฺตวณฺโณ เกวลกปฺปํ นิโคฺรธารามํ โอภาเสตฺวา เยนายสฺมา โลมสกงฺคิโย เตนุปสงฺกมิ, อุปสงฺกมิตฺวา เอกมนฺตํ อฏฺฐาสิ, เอกมนฺตํ ฐิโต โข จนฺทโน เทวปุตฺโต อายสฺมนฺตํ โลมสกงฺคิยํ เอตทโวจ “ธาเรสิ ตฺวํ ภิกฺขุ ภทฺเทกรตฺตสฺส อุทฺเทสญฺจ วิภงฺคญฺจา”ติ.\csromanspacer{} น โข อหํ อาวุโส ธาเรมิ ภทฺเทกรตฺตสฺส อุทฺเทสญฺจ วิภงฺคญฺจ, ตฺวํ ปนาวุโส ธาเรสิ ภทฺเทกรตฺตสฺส อุทฺเทสญฺจ วิภงฺคญฺจาติ.\csromanspacer{} อหมฺปิ โข ภิกฺขุ น ธาเรมิ ภทฺเทกรตฺตสฺส อุทฺเทสญฺจ วิภงฺคญฺจ, ธาเรสิ ปน ตฺวํ ภิกฺขุ ภทฺเทกรตฺติโย คาถาติ.\csromanspacer{} น โข อหํ อาวุโส ธาเรมิ ภทฺเทกรตฺติโย คาถา, ตฺวํ ปนาวุโส ธาเรสิ ภทฺเทกรตฺติโย คาถาติ.\csromanspacer{} ธาเรมิ โข อหํ ภิกฺขุ ภทฺเทกรตฺติโย คาถาติ.\csromanspacer{} ยถา กถํ ปน ตฺวํ อาวุโส ธาเรสิ ภทฺเทกรตฺติโย คาถาติ.\csromanspacer{} เอกมิทํ ภิกฺขุ สมยํ ภควา เทเวสุ ตาวติํเสสุ วิหรติ ปาริจฺฉตฺตกมูเล ปณฺฑุกมฺพลสิลายํ, ตตฺร ภควา เทวานํ ตาวติํสานํ ภทฺเทกรตฺตสฺส อุทฺเทสญฺจ วิภงฺคญฺจ อภาสิ\csromandash{}}

\csromangathameasure{“อตีตํ นานฺวาคเมยฺย, นปฺปฏิกงฺเข อนาคตํ. \\ ยทตีตํ ปหีนํ ตํ, อปฺปตฺตญฺจ อนาคตํ. \\ ปจฺจุปฺปนฺนญฺจ โย ธมฺมํ, ตตฺถ ตตฺถ วิปสฺสติ. \\ อสํหีรํ อสํกุปฺปํ, ตํ วิทฺวา มนุพฺรูหเย. \\ อชฺเชว กิจฺจมาตปฺปํ, โก ชญฺญา มรณํ สุเว. \\ น หิ โน สงฺครํ เตน, มหาเสเนน มจฺจุนา. \\ เอวํวิหาริํ อาตาปิํ, อโหรตฺตมตนฺทิตํ. \\ ตํ เว ‘ภทฺเทกรตฺโต’ติ, สนฺโต อาจิกฺขเต มุนี”ติ.}
\csromangathagroup{\csromangathastanza{“อตีตํ นานฺวาคเมยฺย, นปฺปฏิกงฺเข อนาคตํ. \\ ยทตีตํ ปหีนํ ตํ, อปฺปตฺตญฺจ อนาคตํ.}\csromangathastanza{ปจฺจุปฺปนฺนญฺจ โย ธมฺมํ, ตตฺถ ตตฺถ วิปสฺสติ. \\ อสํหีรํ อสํกุปฺปํ, ตํ วิทฺวา มนุพฺรูหเย.}\csromangathastanza{อชฺเชว กิจฺจมาตปฺปํ, โก ชญฺญา มรณํ สุเว. \\ น หิ โน สงฺครํ เตน, มหาเสเนน มจฺจุนา.}\csromangathastanza{เอวํวิหาริํ อาตาปิํ, อโหรตฺตมตนฺทิตํ. \\ ตํ เว ‘ภทฺเทกรตฺโต’ติ, สนฺโต อาจิกฺขเต มุนี”ติ.}}

\prose{เอวํ โข อหํ ภิกฺขุ ธาเรมิ ภทฺเทกรตฺติโย คาถา.\csromanspacer{} อุคณฺหาหิ ตฺวํ ภิกฺขุ ภทฺเทกรตฺตสฺส อุทฺเทสญฺจ วิภงฺคญฺจ, ปริยาปุณาหิ}

\csromanheader{2}{4. โลมสกงฺคิยภทฺเทกรตฺตสุตฺต (134)}
\prose{\csromanfolio{241}ตฺวํ ภิกฺขุ ภทฺเทกรตฺตสฺส อุทฺเทสญฺจ วิภงฺคญฺจ, ธาเรหิ ตฺวํ ภิกฺขุ ภทฺเทกรตฺตสฺส อุทฺเทสญฺจ วิภงฺคญฺจ, อตฺถสํหิโต ภิกฺขุ ภทฺเทกรตฺตสฺส อุทฺเทโส จ วิภงฺโค จ อาทิพฺรหฺมจริยโกติ.\csromanspacer{} อิทมโวจ จนฺทโน เทวปุตฺโต, อิทํ วตฺวา ตตฺเถวนฺตรธายิ.}

\proseitem{287}{อถ โข อายสฺมา โลมสกงฺคิโย ตสฺสา รตฺติยา อจฺจเยน เสนาสนํ สํสาเมตฺวา ปตฺตจีวรมาทาย เยน สาวตฺถิ เตน จาริกํ ปกฺกามิ, อนุปุพฺเพน จาริกํ จรมาโน เยน สาวตฺถิ เชตวนํ อนาถปิณฺฑิกสฺส อาราโม เยน ภควา เตนุปสงฺกมิ, อุปสงฺกมิตฺวา ภควนฺตํ อภิวาเทตฺวา เอกมนฺตํ นิสีทิ, เอกมนฺตํ นิสินฺโน โข อายสฺมา โลมสกงฺคิโย ภควนฺตํ เอตทโวจ\csromandash{}}

\prose{เอกมิทาหํ ภนฺเต สมยํ สกฺเกสุ วิหรามิ กปิลวตฺถุสฺมิํ นิโคฺรธาราเม.\csromanspacer{} อถ โข ภนฺเต อญฺญตโร เทวปุตฺโต อภิกฺกนฺตาย รตฺติยา อภิกฺกนฺตวณฺโณ เกวลกปฺปํ นิโคฺรธารามํ โอภาเสตฺวา เยนาหํ เตนุปสงฺกมิ, อุปสงฺกมิตฺวา เอกมนฺตํ อฏฺฐาสิ, เอกมนฺตํ ฐิโต โข ภนฺเต โส เทวปุตฺโต มํ เอตทโวจ “ธาเรสิ ตฺวํ ภิกฺขุ ภทฺเทกรตฺตสฺส อุทฺเทสญฺจ วิภงฺคญฺจา”ติ.\csromanspacer{} เอวํ วุตฺเต อหํ ภนฺเต ตํ เทวปุตฺตํ เอตทโวจํ “น โข อหํ อาวุโส ธาเรมิ ภทฺเทกรตฺตสฺส อุทฺเทสญฺจ วิภงฺคญฺจ, ตฺวํ ปนาวุโส ธาเรสิ ภทฺเทกรตฺตสฺส อุทฺเทสญฺจ วิภงฺคญฺจา”ติ.\csromanspacer{} อหมฺปิ โข ภิกฺขุ น ธาเรมิ ภทฺเทกรตฺตสฺส อุทฺเทสญฺจ วิภงฺคญฺจ, ธาเรสิ ปน ตฺวํ ภิกฺขุ ภทฺเทกรตฺติโย คาถาติ.\csromanspacer{} น โข อหํ อาวุโส ธาเรมิ ภทฺเทกรตฺติโย คาถา, ตฺวํ ปนาวุโส ธาเรสิ ภทฺเทกรตฺติโย คาถาติ.\csromanspacer{} ธาเรมิ โข อหํ ภิกฺขุ ภทฺเทกรตฺติโย คาถาติ.\csromanspacer{} ยถา กถํ ปน ตฺวํ อาวุโส ธาเรสิ ภทฺเทกรตฺติโย คาถาติ.\csromanspacer{} เอกมิทํ ภิกฺขุ สมยํ ภควา เทเวสุ ตาวติํเสสุ วิหรติ ปาริจฺฉตฺตกมูเล ปณฺฑุกมฺพลสิลายํ, ตตฺร โข ภควา เทวานํ ตาวติํสานํ ภทฺเทกรตฺตสฺส อุทฺเทสญฺจ วิภงฺคญฺจ อภาสิ\csromandash{}}

\prose{“อตีตํ นานฺวาคเมยฺย ฯเปฯ.}

\csromangathameasure{ตํ เว ภทฺเทกรตฺโตติ, สนฺโต อาจิกฺขเต มุนี”ติ.}
\csromangathagroup{\csromangathastanza{ตํ เว ภทฺเทกรตฺโตติ, สนฺโต อาจิกฺขเต มุนี”ติ.}}

\prose{เอวํ โข อหํ ภิกฺขุ ธาเรมิ ภทฺเทกรตฺติโย คาถา, อุคฺคณฺหาหิ ตฺวํ ภิกฺขุ ภทฺเทกรตฺตสฺส อุทฺเทสญฺจ วิภงฺคญฺจ, ปริยาปุณาหิ ตฺวํ ภิกฺขุ \csromanfolio{242}ภทฺเทกรตฺตสฺส อุทฺเทสญฺจ วิภงฺคญฺจ, ธาเรหิ ตฺวํ ภิกฺขุ ภทฺเทกรตฺตสฺส อุทฺเทสญฺจ วิภงฺคญฺจ, อตฺถสํหิโต ภิกฺขุ ภทฺเทกรตฺตสฺส อุทฺเทโส จ วิภงฺโค จ อาทิพฺรหฺมจริยโกติ.\csromanspacer{} อิทมโวจ ภนฺเต โส เทวปุตฺโต, อิทํ วตฺวา ตตฺเถวนฺตรธายิ.\csromanspacer{} สาธุ เม ภนฺเต ภควา ภทฺเทกรตฺตสฺส อุทฺเทสญฺจ วิภงฺคญฺจ เทเสตูติ.}

\proseitem{288}{ชานาสิ ปน ตฺวํ ภิกฺขุ ตํ เทวปุตฺตนฺติ.\csromanspacer{} น โข อหํ ภนฺเต ชานามิ ตํ เทวปุตฺตนฺติ.\csromanspacer{} จนฺทโน นาม โส ภิกฺขุ เทวปุตฺโต, จนฺทโน ภิกฺขุ เทวปุตฺโต อฏฺฐิํ กตฺวา\footnote{อฏฺฐิกตฺวา (สี, สฺยา, กํ, อิ)} มนสิกตฺวา สพฺพเจตสา\footnote{สพฺพํ เจตโส (สี, สฺยา, กํ, อิ), สพฺพํ เจตสา (ก)} สมนฺนาหริตฺวา โอหิตโสโต ธมฺมํ สุณาติ.\csromanspacer{} เตน หิ ภิกฺขุ สุณาหิ สาธุกํ มนสิ กโรหิ, ภาสิสฺสามีติ. “เอวํ ภนฺเต”ติ โข อายสฺมา โลมสกงฺคิโย ภควโต ปจฺจสฺโสสิ.\csromanspacer{} ภควา เอตทโวจ\csromandash{}}

\csromangathameasure{อตีตํ นานฺวาคเมยฺย, นปฺปฏิกงฺเข อนาคตํ. \\ ยทตีตํ ปหีนํ ตํ, อปฺปตฺตญฺจ อนาคตํ. \\ ปจฺจุปฺปนฺนญฺจ โย ธมฺมํ, ตตฺถ ตตฺถ วิปสฺสติ. \\ อสํหีรํ อสํกุปฺปํ, ตํ วิทฺวา มนุพฺรูหเย. \\ อชฺเชว กิจฺจมาตปฺปํ, โก ชญฺญา มรณํ สุเว. \\ น หิ โน สงฺครํ เตน, มหาเสเนน มจฺจุนา. \\ เอวํวิหาริํ อาตาปิํ, อโหรตฺตมตนฺทิตํ. \\ ตํ เว “ภทฺเทกรตฺโต”ติ, สนฺโต อาจิกฺขเต มุนิ.}
\csromangathagroup{\csromangathastanza{อตีตํ นานฺวาคเมยฺย, นปฺปฏิกงฺเข อนาคตํ. \\ ยทตีตํ ปหีนํ ตํ, อปฺปตฺตญฺจ อนาคตํ.}\csromangathastanza{ปจฺจุปฺปนฺนญฺจ โย ธมฺมํ, ตตฺถ ตตฺถ วิปสฺสติ. \\ อสํหีรํ อสํกุปฺปํ, ตํ วิทฺวา มนุพฺรูหเย.}\csromangathastanza{อชฺเชว กิจฺจมาตปฺปํ, โก ชญฺญา มรณํ สุเว. \\ น หิ โน สงฺครํ เตน, มหาเสเนน มจฺจุนา.}\csromangathastanza{เอวํวิหาริํ อาตาปิํ, อโหรตฺตมตนฺทิตํ. \\ ตํ เว “ภทฺเทกรตฺโต”ติ, สนฺโต อาจิกฺขเต มุนิ.}}

\prose{กถญฺจ ภิกฺขุ อตีตํ อนฺวาคเมติ ฯเปฯ.\csromanspacer{} เอวํ โข ภิกฺขุ อตีตํ อนฺวาคเมติ.\csromanspacer{} กถญฺจ ภิกฺขุ อตีตํ นานฺวาคเมติ ฯเปฯ.\csromanspacer{} เอวํ โข ภิกฺขุ อตีตํ นานฺวาคเมติ.\csromanspacer{} กถญฺจ ภิกฺขุ อนาคตํ ปฏิกงฺขติ ฯเปฯ.\csromanspacer{} เอวํ โข ภิกฺขุ อนาคตํ ปฏิกงฺขติ.\csromanspacer{} กถญฺจ ภิกฺขุ อนาคตํ นปฺปฏิกงฺขติ ฯเปฯ.\csromanspacer{} เอวํ โข ภิกฺขุ อนาคตํ นปฺปฏิกงฺขติ.\csromanspacer{} กถญฺจ ภิกฺขุ ปจฺจุปฺปนฺเนสุ ธมฺเมสุ สํหีรติ ฯเปฯ.\csromanspacer{} เอวํ โข ภิกฺขุ ปจฺจุปฺปนฺเนสุ ธมฺเมสุ สํหีรติ.\csromanspacer{} กถญฺจ ภิกฺขุ ปจฺจุปฺปนฺเนสุ ธมฺเมสุ น สํหีรติ ฯเปฯ.\csromanspacer{} เอวํ โข ภิกฺขุ ปจฺจุปฺปนฺเนสุ ธมฺเมสุ น สํหีรติ.}

\csromangathameasure{อตีตํ นานฺวาคเมยฺย, นปฺปฏิกงฺเข อนาคตํ. \\ ยทตีตํ ปหีนํ ตํ, อปฺปตฺตญฺจ อนาคตํ.}
\csromangathagroup{\csromangathastanza{อตีตํ นานฺวาคเมยฺย, นปฺปฏิกงฺเข อนาคตํ. \\ ยทตีตํ ปหีนํ ตํ, อปฺปตฺตญฺจ อนาคตํ.}}

\csromanheader{2}{5. จูฬกมฺมวิภงฺคสุตฺต (\textbf{1}35)}
\csromangathameasure{ปจฺจุปฺปนฺนญฺจ โย ธมฺมํ, ตตฺถ ตตฺถ วิปสฺสติ. \\ อสํหีรํ อสํกุปฺปํ, ตํ วิทฺวา มนุพฺรูหเย. \\ อชฺเชว กิจฺจมาตปฺปํ, โก ชญฺญา มรณํ สุเว. \\ น หิ โน สงฺครํ เตน, มหาเสเนน มจฺจุนา. \\ เอวํวิหาริํ อาตาปิํ, อโหรตฺตมตนฺทิตํ. \\ ตํ เว “ภทฺเทกรตฺโต”ติ, สนฺโต อาจิกฺขเต มุนีติ.}
\csromangathagroup{\csromangathastanza{\csromanfolio{243}ปจฺจุปฺปนฺนญฺจ โย ธมฺมํ, ตตฺถ ตตฺถ วิปสฺสติ. \\ อสํหีรํ อสํกุปฺปํ, ตํ วิทฺวา มนุพฺรูหเย.}\csromangathastanza{อชฺเชว กิจฺจมาตปฺปํ, โก ชญฺญา มรณํ สุเว. \\ น หิ โน สงฺครํ เตน, มหาเสเนน มจฺจุนา.}\csromangathastanza{เอวํวิหาริํ อาตาปิํ, อโหรตฺตมตนฺทิตํ. \\ ตํ เว “ภทฺเทกรตฺโต”ติ, สนฺโต อาจิกฺขเต มุนีติ.}}

\prose{อิทมโวจ ภควา.\csromanspacer{} อตฺตมโน อายสฺมา โลมสกงฺคิโย ภควโต ภาสิตํ อภินนฺทีติ.}

\nitthitamruled{โลมสกงฺคิยภทฺเทกรตฺตสุตฺตํ นิฏฺฐิตํ จตุตฺถํ.}
\csromanmatikaanchor{mtk.42}
\csromantocmarkref{subsection}{5. จูฬกมฺมวิภงฺคสุตฺต}{mtk.42}
\bookmark[dest={mtk.42},level=2]{5. จูฬกมฺมวิภงฺคสุตฺต}
\csromanheader{2}{5. จูฬกมฺมวิภงฺคสุตฺต\footnote{สุภสุตฺตนฺติปิ วุจฺจติ.}}
\proseitem{289}{เอวํ เม สุตํ\csromandash{}เอกํ สมยํ ภควา สาวตฺถิยํ วิหรติ เชตวเน อนาถปิณฺฑิกสฺส อาราเม.\csromanspacer{} อถ โข สุโภ มาณโว โตเทยฺยปุตฺโต เยน ภควา เตนุปสงฺกมิ, อุปสงฺกมิตฺวา ภควตา สทฺธิํ สมฺโมทิ, สมฺโมทนียํ กถํ สารณียํ วีติสาเรตฺวา เอกมนฺตํ นิสีทิ, เอกมนฺตํ นิสินฺโน โข สุโภ มาณโว โตเทยฺยปุตฺโต ภควนฺตํ เอตทโวจ\csromandash{}}

\prose{โก นุ โข โภ โคตม เหตุ โก ปจฺจโย, เยน มนุสฺสานํเยว สตํ มนุสฺสภูตานํ ทิสฺสนฺติ หีนปฺปณีตตา.\csromanspacer{} ทิสฺสนฺติ หิ โภ โคตม มนุสฺสา อปฺปายุกา, ทิสฺสนฺติ ทีฆายุกา.\csromanspacer{} ทิสฺสนฺติ พวฺหาพาธา\footnote{พหฺวาพาธา (สฺยา, กํ, ก)}, ทิสฺสนฺติ อปฺปาพาธา.\csromanspacer{} ทิสฺสนฺติ ทุพฺพณฺณา, ทิสฺสนฺติ วณฺณวนฺโต.\csromanspacer{} ทิสฺสนฺติ อปฺเปสกฺขา, ทิสฺสนฺติ มเหสกฺขา.\csromanspacer{} ทิสฺสนฺติ อปฺปโภคา, ทิสฺสนฺติ มหาโภคา.\csromanspacer{} ทิสฺสนฺติ นีจกุลีนา, ทิสฺสนฺติ อุจฺจากุลีนา.\csromanspacer{} ทิสฺสนฺติ ทุปฺปญฺญา, ทิสฺสนฺติ ปญฺญวนฺโต\footnote{ปญฺญาวนฺโต (สี, อิ)}.\csromanspacer{} โก นุ โข โภ โคตม เหตุ โก ปจฺจโย, เยน มนุสฺสานํเยว สตํ มนุสฺสภูตานํ ทิสฺสนฺติ หีนปฺปณีตตาติ.}

\prose{\csromanfolio{244}กมฺมสฺสกา มาณว สตฺตา กมฺมทายาทา กมฺมโยนี กมฺมพนฺธู\footnote{กมฺมโยนิ กมฺมพนฺธุ (สี)} กมฺมปฺปฏิสรณา, กมฺมํ สตฺเต วิภชติ ยทิทํ หีนปฺปณีตตายาติ.\csromanspacer{} น โข อหํ อิมสฺส โภโต โคตมสฺส สํขิตฺเตน ภาสิตสฺส วิตฺถาเรน อตฺถํ อวิภตฺตสฺส วิตฺถาเรน อตฺถํ อาชานามิ.\csromanspacer{} สาธุ เม ภวํ โคตโม ตถา ธมฺมํ เทเสตุ, ยถา อหํ อิมสฺส โภโต โคตมสฺส สํขิตฺเตน ภาสิตสฺส วิตฺถาเรน อตฺถํ อวิภตฺตสฺส วิตฺถาเรน อตฺถํ อาชาเนยฺยนฺติ.}

\proseitem{290}{เตน หิ มาณว สุณาหิ สาธุกํ มนสิ กโรหิ, ภาสิสฺสามีติ. “เอวํ โภ”ติ โข สุโภ มาณโว โตเทยฺยปุตฺโต ภควโต ปจฺจสฺโสสิ.\csromanspacer{} ภควา เอตทโวจ\csromandash{}}

\prose{อิธ มาณว เอกจฺโจ อิตฺถี วา ปุริโส วา ปาณาติปาตี โหติ, ลุทฺโท โลหิตปาณิ หตปหเต นิวิฏฺโฐ อทยาปนฺโน ปาณภูเตสุ\footnote{สพฺพปาณภูเตสุ (สี, ก)}, โส เตน กมฺเมน เอวํ สมตฺเตน เอวํ สมาทินฺเนน\footnote{สมาทิณฺเณน (อิ, ก)} กายสฺส เภทา ปรํ มรณา อปายํ ทุคฺคติํ วินิปาตํ นิรยํ อุปปชฺชติ.\csromanspacer{} โน เจ กายสฺส เภทา ปรํ มรณา อปายํ ทุคฺคติํ วินิปาตํ นิรยํ อุปปชฺชติ, สเจ มนุสฺสตฺตํ อาคจฺฉติ, ยตฺถ ยตฺถ ปจฺจาชายติ, อปฺปายุโก โหติ.\csromanspacer{} อปฺปายุกสํวตฺตนิกา เอสา มาณว ปฏิปทา ยทิทํ ปาณาติปาตี โหติ, ลุทฺโท โลหิตปาณิ หตปหเต นิวิฏฺโฐ อทยาปนฺโน ปาณภูเตสุ. (1)}

\prose{อิธ ปน มาณว เอกจฺโจ อิตฺถี วา ปุริโส วา ปาณาติปาตํ ปหาย ปาณาติปาตา ปฏิวิรโต โหติ, นิหิตทณฺโฑ นิหิตสตฺโถ ลชฺชี ทยาปนฺโน สพฺพปาณภูตหิตานุกมฺปี วิหรติ, โส เตน กมฺเมน เอวํ สมตฺเตน เอวํ สมาทินฺเนน กายสฺส เภทา ปรํ มรณา สุคติํ สคฺคํ โลกํ อุปปชฺชติ.\csromanspacer{} โน เจ กายสฺส เภทา ปรํ มรณา สุคติํ สคฺคํ โลกํ อุปปชฺชติ, สเจ มนุสฺสตฺตํ อาคจฺฉติ, ยตฺถ ยตฺถ ปจฺจาชายติ, ทีฆายุโก โหติ.\csromanspacer{} ทีฆายุกสํวตฺตนิกา เอสา มาณว ปฏิปทา ยทิทํ ปาณาติปาตํ ปหาย ปาณาติปาตา ปฏิวิรโต โหติ, นิหิตทณฺโฑ นิหิตสตฺโถ ลชฺชี ทยาปนฺโน สพฺพปาณภูตหิตานุกมฺปี วิหรติ. (1)}

\csromanheader{2}{5. จูฬกมฺมวิภงฺคสุตฺต (135)}
\proseitem{291}{\csromanfolio{245}อิธ มาณว เอกจฺโจ อิตฺถี วา ปุริโส วา สตฺตานํ วิเหฐกชาติโก โหติ ปาณินา วา เลฑฺฑุนา วา ทณฺเฑน วา สตฺเถน วา, โส เตน กมฺเมน เอวํ สมตฺเตน เอวํ สมาทินฺเนน กายสฺส เภทา ปรํ มรณา อปายํ ทุคฺคติํ วินิปาตํ นิรยํ อุปปชฺชติ.\csromanspacer{} โน เจ กายสฺส เภทา ปรํ มรณา อปายํ ทุคฺคติํ วินิปาตํ นิรยํ อุปปชฺชติ, สเจ มนุสฺสตฺตํ อาคจฺฉติ, ยตฺถ ยตฺถ ปจฺจาชายติ, พวฺหาพาโธ โหติ.\csromanspacer{} พวฺหาพาธสํวตฺตนิกา เอสา มาณว ปฏิปทา ยทิทํ สตฺตานํ วิเหฐกชาติโก โหติ ปาณินา วา เลฑฺฑุนา วา ทณฺเฑนวา สตฺเถน วา. (2)}

\prose{อิธ ปน มาณว เอกจฺโจ อิตฺถี วา ปุริโส วา สตฺตานํ อวิเหฐกชาติโก โหติ ปาณินา วา เลฑฺฑุนา วา ทณฺเฑน วา สตฺเถน วา, โส เตน กมฺเมน เอวํ สมตฺเตน เอวํ สมาทินฺเนน กายสฺส เภทา ปรํ มรณา สุคติํ สคฺคํ โลกํ อุปปชฺชติ.\csromanspacer{} โน เจ กายสฺส เภทา ปรํ มรณา สุคติํ สคฺคํ โลกํ อุปปชฺชติ, สเจ มนุสฺสตฺตํ อาคจฺฉติ, ยตฺถ ยตฺถ ปจฺจาชายติ, อปฺปาพาโธ โหติ.\csromanspacer{} อปฺปาพาธสํวตฺตนิกา เอสา มาณว ปฏิปทา ยทิทํ สตฺตานํ อวิเหฐกชาติโก โหติ ปาณินา วา เลฑฺฑุนา วา ทณฺเฑน วา สตฺเถน วา. (2)}

\proseitem{292}{อิธ มาณว เอกจฺโจ อิตฺถี วา ปุริโส วา โกธโน โหติ อุปายาสพหุโล, อปฺปมฺปิ วุตฺโต สมาโน อภิสชฺชติ กุปฺปติ พฺยาปชฺชติ ปติฏฺฐียติ, โกปญฺจ โทสญฺจ อปฺปจฺจยญฺจ ปาตุกโรติ, โส เตน กมฺเมน เอวํ สมตฺเตน เอวํ สมาทินฺเนน กายสฺส เภทา ปรํ มรณา อปายํ ทุคฺคติํ วินิปาตํ นิรยํ อุปปชฺชติ.\csromanspacer{} โน เจ กายสฺส เภทา ปรํ มรณา อปายํ ทุคฺคติํ วินิปาตํ นิรยํ อุปปชฺชติ, สเจ มนุสฺสตฺตํ อาคจฺฉติ, ยตฺถ ยตฺถ ปจฺจาชายติ, ทุพฺพณฺโณ โหติ.\csromanspacer{} ทุพฺพณฺณสํวตฺตนิกา เอสา มาณว ปฏิปทา ยทิทํ โกธโน โหติ อุปายาสพหุโล, อปฺปมฺปิ วุตฺโต สมาโน อภิสชฺชติ กุปฺปติ พฺยาปชฺชติ ปติฏฺฐียติ, โกปญฺจ โทสญฺจ อปฺปจฺจยญฺจ ปาตุกโรติ. (3)}

\prose{อิธ ปน มาณว เอกจฺโจ อิตฺถี วา ปุริโส วา อกฺโกธโน โหติ อนุปายาสพหุโล, พหุมฺปิ วุตฺโต สมาโน นาภิสชฺชติ น กุปฺปติ น พฺยาปชฺชติ น ปติฏฺฐียติ, น โกปญฺจ โทสญฺจ อปฺปจฺจยญฺจ ปาตุกโรติ, โส เตน กมฺเมน เอวํ สมตฺเตน เอวํ สมาทินฺเนน กายสฺส เภทา ปรํ มรณา สุคติํ สคฺคํ โลกํ อุปปชฺชติ.\csromanspacer{} โน เจ กายสฺส เภทา ปรํ \csromanfolio{246}มรณา สุคติํ สคฺคํ โลกํ อุปปชฺชติ, สเจ มนุสฺสตฺตํ อาคจฺฉติ, ยตฺถ ยตฺถ ปจฺจาชายติ, ปาสาทิโก โหติ.\csromanspacer{} ปาสาทิกสํวตฺตนิกา เอสา มาณว ปฏิปทา ยทิทํ อกฺโกธโน โหติ อนุปายาสพหุโล, พหุมฺปิ วุตฺโต สมาโน นาภิสชฺชติ น กุปฺปติ น พฺยาปชฺชติ น ปติฏฺฐียติ, น โกปญฺจ โทสญฺจ อปฺปจฺจยญฺจ ปาตุกโรติ. (3)}

\proseitem{293}{อิธ มาณว เอกจฺโจ อิตฺถี วา ปุริโส วา อิสฺสามนโก โหติ, ปรลาภสกฺการครุการมานนวนฺทนปูชนาสุ อิสฺสติ อุปทุสฺสติ อิสฺสํ พนฺธติ, โส เตน กมฺเมน เอวํ สมตฺเตน เอวํ สมาทินฺเนน กายสฺส เภทา ปรํ มรณา อปายํ ทุคฺคติํ วินิปาตํ นิรยํ อุปปชฺชติ.\csromanspacer{} โน เจ กายสฺส เภทา ปรํ มรณา อปายํ ทุคฺคติํ วินิปาตํ นิรยํ อุปปชฺชติ, สเจ มนุสฺสตฺตํ อาคจฺฉติ, ยตฺถ ยตฺถ ปจฺจาชายติ, อปฺเปสกฺโข โหติ.\csromanspacer{} อปฺเปสกฺขสํวตฺตนิกา เอสา มาณว ปฏิปทา ยทิทํ อิสฺสามนโก โหติ, ปรลาภสกฺการครุการมานนวนฺทนปูชนาสุ อิสฺสติ อุปทุสฺสติ อิสฺสํ พนฺธติ. (4)}

\prose{อิธ ปน มาณว เอกจฺโจ อิตฺถี วา ปุริโส วา อนิสฺสามนโก โหติ, ปรลาภสกฺการครุการมานนวนฺทนปูชนาสุ น อิสฺสติ น อุปทุสฺสติ น อิสฺสํ พนฺทติ, โส เตน กมฺเมน เอวํ สมตฺเตน เอวํ สมาทินฺเนน กายสฺส เภทา ปรํ มรณา สุคติํ สคฺคํ โลกํ อุปปชฺชติ.\csromanspacer{} โน เจ กายสฺส เภทา ปรํ มรณา สุคติํ สคฺคํ โลกํ อุปปชฺชติ, สเจ มนุสฺสตฺตํ อาคจฺฉติ, ยตฺถ ยตฺถ ปจฺจาชายติ, มเหสกฺโข โหติ.\csromanspacer{} มเหสกฺขสํวตฺตนิกา เอสา มาณว ปฏิปทา ยทิทํ อนิสฺสามนโก โหติ, ปรลาภสกฺการครุการมานนวนฺทนปูชนาสุ น อิสฺสติ น อุปทุสฺสติ น อิสฺสํ พนฺธติ. (4)}

\proseitem{294}{อิธ มาณว เอกจฺโจ อิตฺถี วา ปุริโส วา น ทาตา โหติ สมณสฺส วา พฺราหฺมณสฺส วา อนฺนํ ปานํ วตฺถํ ยานํ มาลาคนฺธวิเลปนํ เสยฺยาวสถปทีเปยฺยํ, โส เตน กมฺเมน เอวํ สมตฺเตน เอวํ สมาทินฺเนน กายสฺส เภทา ปรํ มรณา อปายํ ทุคฺคติํ วินิปาตํ นิรยํ อุปปชฺชติ.\csromanspacer{} โน เจ กายสฺส เภทา ปรํ มรณา อปายํ ทุคฺคติํ วินิปาตํ นิรยํ อุปปชฺชติ, สเจ มนุสฺสตฺตํ อาคจฺฉติ, ยตฺถ ยตฺถ ปจฺจาชายติ, อปฺปโภโค โหติ.\csromanspacer{} อปฺปโภคสํวตฺตนิกา เอสา มาณว ปฏิปทา}

\csromanheader{2}{5. จูฬกมฺมวิภงฺคสุตฺต (135)}
\prose{\csromanfolio{247}ยทิทํ น ทาตา โหติ สมณสฺส วา พฺราหฺมณสฺส วา อนฺนํ ปานํ วตฺถํ ยานํ มาลาคนฺธวิเลปนํ เสยฺยาวสถปทีเปยฺยํ. (5)}

\prose{อิธ ปน มาณว เอกจฺโจ อิตฺถี วา ปุริโส วา ทาตา โหติ สมณสฺส วา พฺราหฺมณสฺส วา อนฺนํ ปานํ วตฺถํ ยานํ มาลาคนฺธวิเลปนํ เสยฺยาวสถปทีเปยฺยํ, โส เตน กมฺเมน เอวํ สมตฺเตน เอวํ สมาทินฺเนน กายสฺส เภทา ปรํ มรณา สุคติํ สคฺคํ โลกํ อุปปชฺชติ.\csromanspacer{} โน เจ กายสฺส เภทา ปรํ มรณา สุคติํ สคฺคํ โลกํ อุปปชฺชติ, สเจ มนุสฺสตฺตํ อาคจฺฉติ, ยตฺถ ยตฺถ ปจฺจาชายติ, มหาโภโค โหติ.\csromanspacer{} มหาโภคสํวตฺตนิกา เอสา มาณว ปฏิปทา ยทิทํ ทาตา โหติ สมณสฺส วา พฺราหฺมณสฺส วา อนฺนํ ปานํ วตฺถํ ยานํ มาลาคนฺธวิเลปนํ เสยฺยาวสถปทีเปยฺยํ. (5)}

\proseitem{295}{อิธ มาณว เอกจฺโจ อิตฺถี วา ปุริโส วา ถทฺโธ โหติ อติมานี\csromandash{} อภิวาเทตพฺพํ น อภิวาเทติ, ปจฺจุฏฺฐาตพฺพํ น ปจฺจุฏฺเฐติ, อาสนารหสฺส น อาสนํ เทติ, มคฺคารหสฺส น มคฺคํ เทติ, สกฺกาตพฺพํ น สกฺกโรติ, ครุกาตพฺพํ น ครุกโรติ, มาเนตพฺพํ น มาเนติ, ปูเชตพฺพํ น ปูเชติ, โส เตน กมฺเมน เอวํ สมตฺเตน เอวํ สมาทินฺเนน กายสฺส เภทา ปรํ มรณา อปายํ ทุคฺคติํ วินิปาตํ นิรยํ อุปปชฺชติ.\csromanspacer{} โน เจ กายสฺส เภทา ปรํ มรณา อปายํ ทุคฺคติํ วินิปาตํ นิรยํ อุปปชฺชติ, สเจ มนุสฺสตฺตํ อาคจฺฉติ, ยตฺถ ยตฺถ ปจฺจาชายติ, นีจกุลีโน โหติ.\csromanspacer{} นีจกุลีนสํวตฺตนิกา เอสา มาณว ปฏิปทา ยทิทํ ถทฺโธ โหติ อติมานี\csromandash{}อภิวาเทตพฺพํ น อภิวาเทติ, ปจฺจุฏฺฐาตพฺพํ น ปจฺจุฏฺเฐติ, อาสนารหสฺส น อาสนํ เทติ, มคฺคารหสฺส น มคฺคํ เทติ, สกฺกาตพฺพํ น สกฺกโรติ, ครุกาตพฺพํ น ครุกโรติ, มาเนตพฺพํ น มาเนติ, ปูเชตพฺพํ น ปูเชติ. (6)}

\prose{อิธ ปน มาณว เอกจฺโจ อิตฺถี วา ปุริโส วา อตฺถทฺโธ โหติ อนติมานี\csromandash{} อภิวาเทตพฺพํ อภิวาเทติ, ปจฺจุฏฺฐาตพฺพํ ปจฺจุฏฺเฐติ, อาสนารหสฺส อาสนํ เทติ, มคฺคารหสฺส มคฺคํ เทติ, สกฺกาตพฺพํ สกฺกโรติ, ครุกาตพฺพํ ครุกโรติ, มาเนตพฺพํ มาเนติ, ปูเชตพฺพํ ปูเชติ, โส เตน กมฺเมน เอวํ สมตฺเตน เอวํ สมาทินฺเนน กายสฺส เภทา ปรํ มรณา สุคติํ สคฺคํ โลกํ อุปปชฺชติ.\csromanspacer{} โน เจ กายสฺส เภทา ปรํ มรณา สุคติํ สคฺคํ อุปปชฺชติ, สเจ \csromanfolio{248}มนุสฺสตฺตํ อาคจฺฉติ, ยตฺถ ยตฺถ ปจฺจาชายติ, อุจฺจากุลีโน โหติ.\csromanspacer{} อุจฺจากุลีนสํวตฺตนิกา เอสา มาณว ปฏิปทา ยทิทํ อตฺถทฺโธ โหติ อนติมานี\csromandash{}อภิวาเทตพฺพํ อภิวาเทติ, ปจฺจุฏฺฐาตพฺพํ ปจฺจุฏฺเฐติ, อาสนารหสฺส อาสนํ เทติ, มคฺคารหสฺส มคฺคํ เทติ, สกฺกาตพฺพํ สกฺกโรติ, ครุกาตพฺพํ ครุกโรติ, มาเนตพฺพํ มาเนติ, ปูเชตพฺพํ ปูเชติ. (6)}

\proseitem{296}{อิธ มาณว เอกจฺโจ อิตฺถี วา ปุริโส วา สมณํ วา พฺราหฺมณํ วา อุปสงฺกมิตฺวา น ปริปุจฺฉิตา โหติ “กิํ ภนฺเต กุสลํ, กิํ อกุสลํ.\csromanspacer{} กิํ สาวชฺชํ, กิํ อนวชฺชํ.\csromanspacer{} กิํ เสวิตพฺพํ, กิํ น เสวิตพฺพํ.\csromanspacer{} กิํ เม กรียมานํ ทีฆรตฺตํ อหิตาย ทุกฺขาย โหติ, กิํ วา ปน เม กรียมานํ ทีฆรตฺตํ หิตาย สุขาย โหตี”ติ, โส เตน กมฺเมน เอวํ สมตฺเตน เอวํ สมาทินฺเนน กายสฺส เภทา ปรํ มรณา อปายํ ทุคฺคติํ วินิปาตํ นิรยํ อุปปชฺชติ.\csromanspacer{} โน เจ กายสฺส เภทา ปรํ มรณา อปายํ ทุคฺคติํ วินิปาตํ นิรยํ อุปปชฺชติ, สเจ มนุสฺสตฺตํ อาคจฺฉติ, ยตฺถ ยตฺถ ปจฺจาชายติ, ทุปฺปญฺโญ โหติ.\csromanspacer{} ทุปฺปญฺญสํวตฺตนิกา เอสา มาณว ปฏิปทา ยทิทํ สมณํ วา พฺราหฺมณํ วา อุปสงฺกมิตฺวา น ปริปุจฺฉิตา โหติ “กิํ ภนฺเต กุสลํ, กิํ อกุสลํ.\csromanspacer{} กิํ สาวชฺชํ, กิํ อนวชฺชํ.\csromanspacer{} กิํ เสวิตพฺพํ, กิํ น เสวิตพฺพํ.\csromanspacer{} กิํ เม กรียมานํ ทีฆรตฺตํ อหิตาย ทุกฺขาย โหติ, กิํ วา ปน เม กรียมานํ ทีฆรตฺตํ หิตาย สุขาย โหตี”ติ. (7)}

\prose{อิธ ปน มาณว เอกจฺโจ อิตฺถี วา ปุริโส วา สมณํ วา พฺราหฺมณํ วา อุปสงฺกมิตฺวา ปริปุจฺฉิตา โหติ “กิํ ภนฺเต กุสลํ, กิํ อกุสลํ.\csromanspacer{} กิํ สาวชฺชํ, กิํ อนวชฺชํ.\csromanspacer{} กิํ เสวิตพฺพํ, กิํ น เสวิตพฺพํ.\csromanspacer{} กิํ เม กรียมานํ ทีฆรตฺตํ อหิตาย ทุกฺขาย โหติ, กิํ วา ปน เม กรียมานํ ทีฆรตฺตํ หิตาย สุขาย โหตี”ติ.\csromanspacer{} โส เตน กมฺเมน เอวํ สมตฺเตน เอวํ สมาทินฺเนน กายสฺส เภทา ปรํ มรณา สุคติํ สคฺคํ โลกํ อุปปชฺชติ.\csromanspacer{} โน เจ กายสฺส เภทา ปรํ มรณา สุคติํ สคฺคํ โลกํ อุปปชฺชติ, สเจ มนุสฺสตฺตํ อาคจฺฉติ, ยตฺถ ยตฺถ ปจฺจาชายติ, มหาปญฺโญ โหติ.\csromanspacer{} มหาปญฺญสํวตฺตนิกา เอสา มาณว ปฏิปทา ยทิทํ สมณํ วา พฺราหฺมณํ วา อุปสงฺกมิตฺวา ปริปุจฺฉิตา โหติ “กิํ ภนฺเต กุสลํ, กิํ อกุสลํ.\csromanspacer{} กิํ สาวชฺชํ, กิํ อนวชฺชํ.\csromanspacer{} กิํ}

\csromanheader{2}{6. มหากมฺมวิภงฺคสุตฺต (136)}
\prose{\csromanfolio{249}เสวิตพฺพํ, กิํ น เสวิตพฺพํ.\csromanspacer{} กิํ เม กรียมานํ ทีฆรตฺตํ อหิตาย ทุกฺขาย โหติ, กิํ วา ปน เม กรียมานํ ทีฆรตฺตํ หิตาย สุขาย โหตี”ติ. (7)}

\proseitem{297}{อิติ โข มาณว อปฺปายุกสํวตฺตนิกา ปฏิปทา อปฺปายุกตฺตํ อุปเนติ, ทีฆายุกสํวตฺตนิกา ปฏิปทา ทีฆายุกตฺตํ อุปเนติ.\csromanspacer{} พวฺหาพาธสํวตฺตนิกา ปฏิปทา พวฺหาพาธตฺตํ อุปเนติ, อปฺปาพาธสํวตฺตนิกา ปฏิปทา อปฺปาพาธตฺตํ อุปเนติ.\csromanspacer{} ทุพฺพณฺณสํวตฺตนิกา ปฏิปทา ทุพฺพณฺณตฺตํ อุปเนติ, ปาสาทิกสํวตฺตนิกา ปฏิปทา ปาสาทิกตฺตํ อุปเนติ.\csromanspacer{} อปฺเปสกฺขสํวตฺตนิกา ปฏิปทา อปฺเปสกฺขตฺตํ อุปเนติ, มเหสกฺขสํวตฺตนิกา ปฏิปทา มเหสกฺขตฺตํ อุปเนติ.\csromanspacer{} อปฺปโภคสํวตฺตนิกา ปฏิปทา อปฺปโภคตฺตํ อุปเนติ, มหาโภคสํวตฺตนิกา ปฏิปทา มหาโภคตฺตํ อุปเนติ.\csromanspacer{} นีจกุลีนสํวตฺตนิกา ปฏิปทา นีจกุลีนตฺตํ อุปเนติ, อุจฺจากุลีนสํวตฺตนิกา ปฏิปทา อุจฺจากุลีนตฺตํ อุปเนติ.\csromanspacer{} ทุปฺปญฺญสํวตฺตนิกา ปฏิปทา ทุปฺปญฺญตฺตํ อุปเนติ, มหาปญฺญสํวตฺตนิกา ปฏิปทา มหาปญฺญตฺตํ อุปเนติ.\csromanspacer{} กมฺมสฺสกา มาณว สตฺตา กมฺมทายาทา กมฺมโยนี กมฺมพนฺธู กมฺมปฺปฏิสรณา, กมฺมํ สตฺเต วิภชติ ยทิทํ หีนปฺปณีตตายาติ.}

\prose{เอวํ วุตฺเต สุโภ มาณโว โตเทยฺยปุตฺโต ภควนฺตํ เอตทโวจ “อภิกฺกนฺตํ โภ โคตม, อภิกฺกนฺตํ โภ โคตม, เสยฺยถาปิ โภ โคตม นิกฺขุชฺชิตํ วา อุกฺกุชฺเชยฺย, ปฏิจฺฉนฺนํ วา วิวเรยฺย, มูฬฺหสฺส วา มคฺคํ อาจิกฺเขยฺย, อนฺธกาเร วา เตลปชฺโชตํ ธาเรยฺย ‘จกฺขุมนฺโต รูปานิ ทกฺขนฺตี’ติ.\csromanspacer{} เอวเมวํ โภตา โคตเมน อเนกปริยาเยน ธมฺโม ปกาสิโต, เอสาหํ ภวนฺตํ โคตมํ สรณํ คจฺฉามิ ธมฺมญฺจ ภิกฺขุสํฆญฺจ, อุปาสกํ มํ ภวํ โคตโม ธาเรตุ อชฺชตคฺเค ปาณุเปตํ สรณํ คตนฺ”ติ.}

\nitthitamruled{จูฬกมฺมวิภงฺคสุตฺตํ นิฏฺฐิตํ ปญฺจมํ.}
\csromanmatikaanchor{mtk.43}
\csromantocmarkref{subsection}{6. มหากมฺมวิภงฺคสุตฺต}{mtk.43}
\bookmark[dest={mtk.43},level=2]{6. มหากมฺมวิภงฺคสุตฺต}
\csromanheader{2}{6. มหากมฺมวิภงฺคสุตฺต}
\proseitem{298}{เอวํ เม สุตํ\csromandash{}เอกํ สมยํ ภควา ราชคเห วิหรติ เวฬุวเน กลนฺทกนิวาเป.\csromanspacer{} เตน โข ปน สมเยน อายสฺมา สมิทฺธิ \csromanfolio{250}อรญฺญกุฏิกายํ วิหรติ.\csromanspacer{} อถ โข โปตลิปุตฺโต ปริพฺพาชโก ชงฺฆาวิหารํ อนุจงฺกมมาโน อนุวิจรมาโน เยนายสฺมา สมิทฺธิ เตนุปสงฺกมิ, อุปสงฺกมิตฺวา อายสฺมตา สมิทฺธินา สทฺธิํ สมฺโมทิ, สมฺโมทนียํ กถํ สารณียํ วีติสาเรตฺวา เอกมนฺตํ นิสีทิ, เอกมนฺตํ นิสินฺโน โข โปตลิปุตฺโต ปริพฺพาชโก อายสฺมนฺตํ สมิทฺธิํ เอตทโวจ “สมฺมุขา เมตํ อาวุโส สมิทฺธิ สมณสฺส โคตมสฺส สุตํ สมฺมุขา ปฏิคฺคหิตํ ‘โมฆํ กายกมฺมํ โมฆํ วจีกมฺมํ, มโนกมฺมเมว สจฺจนฺ’ติ.\csromanspacer{} อตฺถิ จ สา\footnote{อตฺถิเจสา (สี, ก)} สมาปตฺติ, ยํ สมาปตฺติํ สมาปนฺโน น กิญฺจิ เวทิยตี”ติ.\csromanspacer{} มา เหวํ อาวุโส โปตลิปุตฺต อวจ, (มา เหวํ อาวุโส โปตลิปุตฺต อวจ,)\footnote{( ) สฺยา-กํ-โปตฺถเกสุ นตฺถิ.} มา ภควนฺตํ อพฺภาจิกฺขิ, น หิ สาธุ ภควโต อพฺภกฺขานํ, น หิ ภควา เอวํ วเทยฺย “โมฆํ กายกมฺมํ โมฆํ วจีกมฺมํ, มโนกมฺมเมว สจฺจนฺ”ติ.\csromanspacer{} อตฺถิ จ โข\footnote{อตฺถิ เจว โข (สี, ก)} สา อาวุโส สมาปตฺติ, ยํ สมาปตฺติํ สมาปนฺโน น กิญฺจิ เวทิยตีติ.\csromanspacer{} กีวจิรํ ปพฺพชิโตสิ อาวุโส สมิทฺธีติ.\csromanspacer{} นจิรํ อาวุโส ตีณิ วสฺสานีติ.\csromanspacer{} เอตฺถ ทานิ มยํ เถเร ภิกฺขู กิํ วกฺขาม, ยตฺร หิ นาม เอวํนโว ภิกฺขุ\footnote{นวเกน ภิกฺขุนา (ก)} สตฺถารํ ปริรกฺขิตพฺพํ มญฺญิสฺสติ.\csromanspacer{} สญฺเจตนิกํ อาวุโส สมิทฺธิ กมฺมํ กตฺวา กาเยน วาจาย มนสา กิํ โส เวทิยตีติ.\csromanspacer{} สญฺเจตนิกํ อาวุโส โปตลิปุตฺต กมฺมํ กตฺวา กาเยน วาจาย มนสา ทุกฺขํ โส เวทิยตีติ.\csromanspacer{} อถ โข โปตลิปุตฺโต ปริพฺพาชโก อายสฺมโต สมิทฺธิสฺส ภาสิตํ เนว อภินนฺทิ นปฺปฏิกฺโกสิ, อนภินนฺทิตฺวา อปฺปฏิกฺโกสิตฺวา อุฏฺฐายาสนา ปกฺกามิ.}

\proseitem{299}{อถ โข อายสฺมา สมิทฺธิ อจิรปกฺกนฺเต โปตลิปุตฺเต ปริพฺพาชเก เยนายสฺมา อานนฺโท เตนุปสงฺกมิ, อุปสงฺกมิตฺวา อายสฺมตา อานนฺเทน สทฺธิํ สมฺโมทิ, สมฺโมทนียํ กถํ สารณียํ วีติสาเรตฺวา เอกมนฺตํ นิสีทิ, เอกมนฺตํ นิสินฺโน โข อายสฺมา สมิทฺธิ ยาวตโก อโหสิ โปตลิปุตฺเตน ปริพฺพาชเกน สทฺธิํ กถาสลฺลาโป, ตํ สพฺพํ อายสฺมโต อานนฺทสฺส อาโรเจสิ.}

\prose{เอวํ วุตฺเต อายสฺมา อานนฺโท อายสฺมนฺตํ สมิทฺธิํ เอตทโวจ “อตฺถิ โข อิทํ อาวุโส สมิทฺธิ กถาปาภตํ ภควนฺตํ ทสฺสนาย,}

\csromanheader{2}{6. มหากมฺมวิภงฺคสุตฺต (136)}
\prose{\csromanfolio{251}อายามาวุโส สมิทฺธิ เยน ภควา เตนุปสงฺกมิสฺสาม, อุปสงฺกมิตฺวา เอตมตฺถํ ภควโต อาโรเจสฺสาม, ยถา โน ภควา พฺยากริสฺสติ, ตถา นํ ธาเรสฺสามา”ติ. “เอวมาวุโส”ติ โข อายสฺมา สมิทฺธิ อายสฺมโต อานนฺทสฺส ปจฺจสฺโสสิ.}

\prose{อถ โข อายสฺมา จ อานนฺโท อายสฺมา จ สมิทฺธิ เยน ภควา เตนุปสงฺกมิํสุ, อุปสงฺกมิตฺวา ภควนฺตํ อภิวาเทตฺวา เอกมนฺตํ นิสีทิํสุ, เอกมนฺตํ นิสินฺโน โข อายสฺมา อานนฺโท ยาวตโก อโหสิ อายสฺมโต สมิทฺธิสฺส โปตลิปุตฺเตน ปริพฺพาชเกน สทฺธิํ กถาสลฺลาโป, ตํ สพฺพํ ภควโต อาโรเจสิ.\csromanspacer{} เอวํ วุตฺเต ภควา อายสฺมนฺตํ อานนฺทํ เอตทโวจ “ทสฺสนมฺปิ โข อหํ อานนฺท โปตลิปุตฺตสฺส ปริพฺพาชกสฺส นาภิชานามิ, กุโต ปเนวรูปํ กถาสลฺลาปํ.\csromanspacer{} อิมินา จ อานนฺท สมิทฺธินา โมฆปุริเสน โปตลิปุตฺตสฺส ปริพฺพาชกสฺส วิภชฺชพฺยากรณีโย ปโญฺห เอกํเสน พฺยากโต”ติ.\csromanspacer{} เอวํ วุตฺเต อายสฺมา อุทายี ภควนฺตํ เอตทโวจ “สเจ ปน\footnote{กิํ ปน (ก)} ภนฺเต อายสฺมตา สมิทฺธินา อิทํ สนฺธาย ภาสิตํ ‘ยํ กิญฺจิ เวทยิตํ, ตํ ทุกฺขสฺมินฺ’ติ”.}

\proseitem{300}{อถ โข\footnote{เอวํ วุตฺเต (สฺยา, กํ)} ภควา อายสฺมนฺตํ อานนฺทํ อามนฺเตสิ “ปสฺสสิ โน ตฺวํ อานนฺท อิมสฺส อุทายิสฺส โมฆปุริสสฺส อุมฺมงฺคํ\footnote{อุมฺมคฺคํ (สี, สฺยา, กํ, อิ), อุมงฺคํ (ก)}, อญฺญาสิํ โข อหํ อานนฺท อิทาเนวายํ อุทายี โมฆปุริโส อุมฺมุชฺชมาโน อโยนิโส อุมฺมุชฺชิสฺสตี”ติ.\csromanspacer{} อาทิํเยว\footnote{อาทิโสว (สี, อิ), อาทิเยว (ก)} อานนฺท โปตลิปุตฺเตน ปริพฺพาชเกน ติสฺโส เวทนา ปุจฺฉิตา.\csromanspacer{} สจายํ อานนฺท สมิทฺธิ โมฆปุริโส โปตลิปุตฺตสฺส ปริพฺพาชกสฺส เอวํ ปุฏฺโฐ เอวํ พฺยากเรยฺย “สญฺเจตนิกํ อาวุโส โปตลิปุตฺต กมฺมํ กตฺวา กาเยน วาจาย มนสา สุขเวทนียํ, สุขํ โส เวทยติ, สญฺเจตนิกํ อาวุโส โปตลิปุตฺต กมฺมํ กตฺวา กาเยน วาจาย มนสา ทุกฺขเวทนียํ, ทุกฺขํ โส เวทยติ, สญฺเจตนิกํ อาวุโส โปตลิปุตฺต กมฺมํ กตฺวา กาเยน วาจาย มนสา อทุกฺขมสุขเวทนียํ, อทุกฺขมสุขํ โส เวทยตี”ติ.\csromanspacer{} เอวํ พฺยากรมาโน โข อานนฺท สมิทฺธิ โมฆปุริโส โปตลิปุตฺตสฺส ปริพฺพาชกสฺส สมฺมา (พฺยากรมาโน)\footnote{( ) นตฺถิ (สี, สฺยา, กํ, อิ)} พฺยากเรยฺย. \csromanfolio{252}อปิ จ อานนฺท เก จ\footnote{เกจิ (ก)} อญฺญติตฺถิยา ปริพฺพาชกา พาลา อพฺยตฺตา, เก จ1 ตถาคตสฺส มหากมฺมวิภงฺคํ ชานิสฺสนฺติ.\csromanspacer{} สเจ ตุเมฺห อานนฺท สุเณยฺยาถ ตถาคตสฺส มหากมฺมวิภงฺคํ วิภชนฺตสฺสาติ.}

\prose{เอตสฺส ภควา กาโล, เอตสฺส สุคต กาโล, ยํ ภควา มหากมฺมวิภงฺคํ วิภเชยฺย, ภควโต สุตฺวา ภิกฺขู ธาเรสฺสนฺตีติ.\csromanspacer{} เตน หานนฺท สุณาหิ สาธุกํ มนสิ กโรหิ, ภาสิสฺสามีติ. “เอวํ ภนฺเต”ติ โข อายสฺมา อานนฺโท ภควโต ปจฺจสฺโสสิ.\csromanspacer{} ภควา เอตทโวจ\csromandash{}}

\prose{จตฺตาโรเม อานนฺท ปุคฺคลา สนฺโต สํวิชฺชมานา โลกสฺมิํ.\csromanspacer{} กตเม จตฺตาโร, อิธานนฺท เอกจฺโจ ปุคฺคโล อิธ ปาณาติปาตี โหติ, อทินฺนาทายี โหติ, กาเมสุมิจฺฉาจารี โหติ, มุสาวาที โหติ, ปิสุณวาโจ โหติ, ผรุสวาโจ โหติ, สมฺผปฺปลาปี โหติ, อภิชฺฌาลุ โหติ, พฺยาปนฺนจิตฺโต โหติ, มิจฺฉาทิฏฺฐิ โหติ, โส กายสฺส เภทา ปรํ มรณา อปายํ ทุคฺคติํ นิรยํ อุปปชฺชติ. (1)}

\prose{อิธ ปนานนฺท เอกจฺโจ ปุคฺคโล อิธ ปาณาติปาตี โหติ, อทินฺนาทายี โหติ, กาเมสุมิจฺฉาจารี โหติ, มุสาวาที โหติ, ปิสุณวาโจ โหติ, ผรุสวาโจ โหติ, สมฺผปฺปลาปี โหติ, อภิชฺฌาลุ โหติ, พฺยาปนฺนจิตฺโต โหติ, มิจฺฉาทิฏฺฐิ โหติ, โส กายสฺส เภทา ปรํ มรณา สุคติํ สคฺคํ โลกํ อุปปชฺชติ. (2)}

\prose{อิธานนฺท เอกจฺโจ ปุคฺคโล อิธ ปาณาติปาตา ปฏิวิรโต โหติ, อทินฺนาทานา ปฏิวิรโต โหติ, กาเมสุมิจฺฉาจารา ปฏิวิรโต โหติ, มุสาวาทา ปฏิวิรโต โหติ, ปิสุณาย วาจาย ปฏิวิรโต โหติ, ผรุสาย วาจาย ปฏิวิรโต โหติ, สมฺผปฺปลาปา ปฏิวิรโต โหติ, อนภิชฺฌาลุ โหติ, อพฺยาปนฺนจิตฺโต โหติ, สมฺมาทิฏฺฐิ โหติ, โส กายสฺส เภทา ปรํ มรณา สุคติํ สคฺคํ โลกํ อุปปชฺชติ. (3)}

\prose{อิธ ปนานนฺท เอกจฺโจ ปุคฺคโล อิธ ปาณาติปาตา ปฏิวิรโต โหติ, อทินฺนาทานา ปฏิวิรโต โหติ, กาเมสุมิจฺฉาจารา ปฏิวิรโต โหติ, มุสาวาทา ปฏิวิรโต โหติ, ปิสุณาย}

\csromanheader{2}{6. มหากมฺมวิภงฺคสุตฺต (136)}
\prose{\csromanfolio{253}วาจาย ปฏิวิรโต โหติ, ผรุสาย วาจาย ปฏิวิรโต โหติ, สมฺผปฺปลาปา ปฏิวิรโต โหติ, อนภิชฺฌาลุ โหติ, อพฺยาปนฺนจิตฺโต โหติ, สมฺมาทิฏฺฐิ โหติ, โส กายสฺส เภทา ปรํ มรณา อปายํ ทุคฺคติํ วินิปาตํ นิรยํ อุปปชฺชติ. (4)}

\proseitem{301}{อิธานนฺท เอกจฺโจ สมโณ วา พฺราหฺมโณ วา อาตปฺปมนฺวาย ปธานมนฺวาย อนุโยคมนฺวาย อปฺปมาทมนฺวาย สมฺมามนสิการมนฺวาย ตถารูปํ เจโตสมาธิํ ผุสติ, ยถาสมาหิเต จิตฺเต ทิพฺเพน จกฺขุนา วิสุทฺเธน อติกฺกนฺตมานุสเกน อมุํ ปุคฺคลํ ปสฺสติ อิธ ปาณาติปาติํ อทินฺนาทายิํ กาเมสุมิจฺฉาจาริํ มุสาวาทิํ ปิสุณวาจํ ผรุสวาจํ สมฺผปฺปลาปิํ อภิชฺฌาลุํ พฺยาปนฺนจิตฺตํ มิจฺฉาทิฏฺฐิํ, กายสฺส เภทา ปรํ มรณา ปสฺสติ อปายํ ทุคฺคติํ วินิปาตํ นิรยํ อุปปนฺนํ.\csromanspacer{} โส เอวมาห “อตฺถิ กิร โภ ปาปกานิ กมฺมานิ, อตฺถิ ทุจฺจริตสฺส วิปาโก, อมาหํ\footnote{อปาหํ (สี, อิ, ก) อมุํ + อหํ = อมาหํ-อิติ ปทวิภาโค.} ปุคฺคลํ อทฺทสํ อิธ ปาณาติปาติํ อทินฺนาทายิํ ฯเปฯ มิจฺฉาทิฏฺฐิํ, กายสฺส เภทา ปรํ มรณา ปสฺสามิ อปายํ ทุคฺคติํ วินิปาตํ นิรยํ อุปปนฺนนฺ”ติ.\csromanspacer{} โส เอวมาห “โย กิร โภ ปาณาติปาตี อทินฺนาทายี ฯเปฯ มิจฺฉาทิฏฺฐิ, สพฺโพ โส กายสฺส เภทา ปรํ มรณา อปายํ ทุคฺคติํ วินิปาตํ นิรยํ อุปปชฺชติ.\csromanspacer{} เย เอวํ ชานนฺติ, เต สมฺมา ชานนฺติ.\csromanspacer{} เย อญฺญถา ชานนฺติ, มิจฺฉา เตสํ ญาณนฺ”ติ\footnote{มิจฺฉา เต สญฺชานนฺติ (ก)}.\csromanspacer{} อิติ โส ยเทว ตสฺส สามํ ญาตํ สามํ ทิฏฺฐํ สามํ วิทิตํ, ตเทว ตตฺถ ถามสา ปรามาสา\footnote{ปรามสฺส (สี, อิ)} อภินิวิสฺส โวหรติ “อิทเมว สจฺจํ, โมฆมญฺญนฺ”ติ. (1)}

\prose{อิธ ปนานนฺท เอกจฺโจ สมโณ วา พฺราหฺมโณวา อาตปฺปมนฺวาย ปธานมนฺวาย อนุโยคมนฺวาย อปฺปมาทมนฺวาย สมฺมามนสิการมนฺวาย ตถารูปํ เจโตสมาธิํ ผุสติ, ยถาสมาหิเต จิตฺเต ทิพฺเพน จกฺขุนา วิสุทฺเธน อติกฺกนฺตมานุสเกน อมุํ ปุคฺคลํ ปสฺสติ อิธ ปาณาติปาติํ อทินฺนาทายิํ ฯเปฯ มิจฺฉาทิฏฺฐิํ, กายสฺส เภทา ปรํ มรณา ปสฺสติ สุคติํ สคฺคํ โลกํ อุปปนฺนํ.\csromanspacer{} โส เอวมาห “นตฺถิ กิร โภ ปาปกานิ กมฺมานิ, นตฺถิ ทุจฺจริตสฺส วิปาโก, อมาหํ ปุคฺคลํ อทฺทสํ อิธ ปาณาติปาติํ อทินฺนาทายิํ ฯเปฯ มิจฺฉาทิฏฺฐิํ, กายสฺส เภทา ปรํ มรณา ปสฺสามิ สุคติํ สคฺคํ โลกํ อุปปนฺนนฺ”ติ.\csromanspacer{} โส เอวมาห \csromanfolio{254}“โย กิร โภ ปาณาติปาตี อทินฺนาทายี ฯเปฯ มิจฺฉาทิฏฺฐิ, สพฺโพ โส กายสฺส เภทา ปรํ มรณา สุคติํ สคฺคํ โลกํ อุปปชฺชติ.\csromanspacer{} เย เอวํ ชานนฺติ, เต สมฺมา ชานนฺติ.\csromanspacer{} เย อญฺญถา ชานนฺติ, มิจฺฉา เตสํ ญาณนฺ”ติ.\csromanspacer{} อิติ โส ยเทว ตสฺส สามํ ญาตํ สามํ ทิฏฺฐํ สามํ วิทิตํ, ตเทว ตตฺถ ถามสา ปรามาสา อภินิวิสฺส โวหรติ “อิทเมว สจฺจํ โมฆมญฺญนฺ”ติ. (2)}

\prose{อิธานนฺท เอกจฺโจ สมโณ วา พฺราหฺมโณ วา อาตปฺปมนฺวาย ปธานมนฺวาย อนุโยคมนฺวาย อปฺปมาทมนฺวาย สมฺมามนสิการมนฺวาย ตถารูปํ เจโตสมาธิํ ผุสติ, ยถาสมาหิเต จิตฺเต ทิพฺเพน จกฺขุนา วิสุทฺเธน อติกฺกนฺตมานุสเกน อมุํ ปุคฺคลํ ปสฺสติ อิธ ปาณาติปาตา ปฏิวิรตํ อทินฺนาทานา ปฏิวิรตํ กาเมสุมิจฺฉาจารา ปฏิวิรตํ มุสาวาทา ปฏิวิรตํ ปิสุณาย วาจาย ปฏิวิรตํ ผรุสาย วาจาย ปฏิวิรตํ สมฺผปฺปลาปา ปฏิวิรตํ อนภิชฺฌาลุํ อพฺยาปนฺนจิตฺตํ สมฺมาทิฏฺฐิํ, กายสฺส เภทา ปรํ มรณา ปสฺสติ สุคติํ สคฺคํ โลกํ อุปปนฺนํ.\csromanspacer{} โส เอวมาห “อตฺถิ กิร โภ กลฺยาณานิ กมฺมานิ, อตฺถิ สุจริตสฺส วิปาโก, อมาหํ ปุคฺคลํ อทฺทสํ อิธ ปาณาติปาตา ปฏิวิรตํ อทินฺนาทานา ปฏิวิรตํ ฯเปฯ สมฺมาทิฏฺฐิํ, กายสฺส เภทา ปรํ มรณา ปสฺสามิ สุคติํ สคฺคํ โลกํ อุปปนฺนนฺ”ติ.\csromanspacer{} โส เอวมาห “โย กิร โภ ปาณาติปาตา ปฏิวิรโต อทินฺนาทานา ปฏิวิรโต ฯเปฯ สมฺมาทิฏฺฐิ, สพฺโพ โส กายสฺส เภทา ปรํ มรณา สุคติํ สคฺคํ โลกํ อุปปชฺชติ.\csromanspacer{} เย เอวํ ชานนฺติ, เต สมฺมา ชานนฺติ.\csromanspacer{} เย อญฺญาถา ชานนฺติ, มิจฺฉา เตสํ ญาณนฺ”ติ.\csromanspacer{} อิติ โส ยเทว ตสฺส สามํ ญาตํ สามํ ทิฏฺฐํ สามํ วิทิตํ, ตเทว ตตฺถ ถามสา ปรามาสา อภินิวิสฺส โวหรติ “อิทเมว สจฺจํ โมฆมญฺญนฺ”ติ. (3)}

\prose{อิธ ปนานนฺท เอกจฺจา สมโณ วา พฺราหฺมโณ วา อาตปฺปมนฺวาย ปธานมนฺวาย อนุโยคมนฺวาย อปฺปมาทมนฺวาย สมฺมามนสิการมนฺวาย ตถารูปํ เจโตสมาธิํ ผุสติ, ยถาสมาหิเต จิตฺเต ทิพฺเพน จกฺขุนา วิสุทฺเธน อติกฺกนฺตมานุสเกน อมุํ ปุคฺคลํ ปสฺสติ อิธ ปาณาติปาตา ปฏิวิรตํ ฯเปฯ สมฺมาทิฏฺฐิํ, กายสฺส เภทา ปรํ มรณา ปสฺสติ อปายํ ทุคฺคติํ วินิปาตํ นิรยํ อุปปนฺนํ.\csromanspacer{} โส เอวมาห “นตฺถิ กิร โภ กลฺยาณานิ กมฺมานิ, นตฺถิ สุจริตสฺส วิปาโก,}

\csromanheader{2}{6. มหากมฺมวิภงฺคสุตฺต (136)}
\prose{\csromanfolio{255}อมาหํ ปุคฺคลํ อทฺทสํ อิธ ปาณาติปาตา ปฏิวิรตํ อทินฺนาทานา ปฏิวิรตํ ฯเปฯ สมฺมาทิฏฺฐิํ, กายสฺส เภทา ปรํ มรณา ปสฺสามิ อปายํ ทุคฺคติํ วินิปาตํ นิรยํ อุปปนฺนนฺ”ติ.\csromanspacer{} โส เอวมาห “โย กิร โภ ปาณาติปาตา ปฏิวิรโต อทินฺนาทานา ปฏิวิรโต ฯเปฯ สมฺมาทิฏฺฐิ, สพฺโพ โส กายสฺส เภทา ปรํ มรณา อปายํ ทุคฺคติํ วินิปาตํ นิรยํ อุปปชฺชติ.\csromanspacer{} เย เอวํ ชานนฺติ, เต สมฺมา ชานนฺติ.\csromanspacer{} เย อญฺญถา ชานนฺติ, มิจฺฉา เตสํ ญาณนฺ”ติ.\csromanspacer{} อิติ โส ยเทว ตสฺส สามํ ญาตํ สามํ ทิฏฺฐํ สามํ วิทิตํ, ตเทว ตตฺถ ถามสา ปรามาสา อภินิวิสฺส โวหรติ “อิทเมว สจฺจํ โมฆมญฺญนฺ”ติ. (4)}

\proseitem{302}{ตตฺรานนฺท ยฺวายํ สมโณ วา พฺราหฺมโณ วา เอวมาห “อตฺถิ กิร โภ ปาปกานิ กมฺมานิ, อตฺถิ ทุจฺจริตสฺส วิปาโก”ติ, อิทมสฺส อนุชานามิ.\csromanspacer{} ยมฺปิ โส เอวมาห “อมาหํ ปุคฺคลํ อทฺทสํ อิธ ปาณาติปาติํ อทินฺนาทายิํ ฯเปฯ มิจฺฉาทิฏฺฐิํ, กายสฺส เภทา ปรํ มรณา ปสฺสามิ อปายํ ทุคฺคติํ วินิปาตํ นิรยํ อุปปนฺนนฺ”ติ, อิทมฺปิสฺส อนุชานามิ.\csromanspacer{} ยญฺจ โข โส เอวมาห “โย กิร โภ ปาณาติปาตี อทินฺนาทายี ฯเปฯ มิจฺฉาทิฏฺฐิ, สพฺโพ โส กายสฺส เภทา ปรํ มรณา อปายํ ทุคฺคติํ วินิปาตํ นิรยํ อุปปชฺชตี”ติ, อิทมสฺส นานุชานามิ.\csromanspacer{} ยมฺปิ โส เอวมาห “เย เอวํ ชานนฺติ, เต สมฺมา ชานนฺติ.\csromanspacer{} เย อญฺญถา ชานนฺติ, มิจฺฉา เตสํ ญาณนฺ”ติ, อิทมฺปิสฺส นานุชานามิ.\csromanspacer{} ยมฺปิ โส ยเทว ตสฺส สามํ ญาตํ สามํ ทิฏฺฐํ สามํ วิทิตํ, ตเทว ตตฺถ ถามสา ปรามาสา อภินิวิสฺส โวหรติ “อิทเมว สจฺจํ โมฆมญฺญนฺ”ติ, อิทมฺปิสฺส นานุชานามิ.\csromanspacer{} ตํ กิสฺส เหตุ, อญฺญถา หิ อานนฺท ตถาคตสฺส มหากมฺมวิภงฺเค ญาณํ โหติ. (1)}

\prose{ตตฺรานนฺท ยฺวายํ สมโณ วา พฺราหฺมโณ วา เอวมาห “นตฺถิ กิร โภ ปาปกานิ กมฺมานิ, นตฺถิ ทุจฺจริตสฺส วิปาโก”ติ, อิทมสฺส นานุชานามิ.\csromanspacer{} ยญฺจ โข โส เอวมาห “อมาหํ ปุคฺคลํ อทฺทสํ อิธ ปาณาติปาติํ อทินฺนาทายิํ ฯเปฯ มิจฺฉาทิฏฺฐิํ, กายสฺส เภทา ปรํ มรณา ปสฺสามิ สุคติํ สคฺคํ โลกํ อุปปนฺนนฺ”ติ, อิทมสฺส อนุชานามิ.\csromanspacer{} ยญฺจ โข โส เอวมาห “โย กิร โภ ปาณาติปาตี อทินฺนาทายี ฯเปฯ มิจฺฉาทิฏฺฐิ, สพฺโพ โส กายสฺส เภทา ปรํ มรณา สุคติํ สคฺคํ โลกํ อุปปชฺชตี”ติ, อิทมสฺส นานุชานามิ.\csromanspacer{} ยมฺปิ โส เอวมาห “เย \csromanfolio{256}เอวํ ชานนฺติ, เต สมฺมา ชานนฺติ.\csromanspacer{} เย อญฺญถา ชานนฺติ, มิจฺฉา เตสํ ญาณนฺ”ติ, อิทมฺปิสฺส นานุชานามิ.\csromanspacer{} ยมฺปิ โส ยเทว ตสฺส สามํ ญาตํ สามํ ทิฏฺฐํ สามํ วิทิตํ, ตเทว ตตฺถ ถามสา ปรามาสา อภินิวิสฺส โวหรติ “อิทเมว สจฺจํ โมฆมญฺญนฺ”ติ, อิทมฺปิสฺส นานุชานามิ.\csromanspacer{} ตํ กิสฺส เหตุ, อญฺญถา หิ อานนฺท ตถาคตสฺส มหากมฺมวิภงฺเค ญาณํ โหติ. (2)}

\prose{ตตฺรานนฺท ยฺวายํ สมโณ วา พฺราหฺมโณ วา เอวมาห “อตฺถิ กิร โภ กลฺยาณานิ กมฺมานิ, อตฺถิ สุจริตสฺส วิปาโก”ติ, อิทมสฺส อนุชานามิ.\csromanspacer{} ยมฺปิ โส เอวมาห “อมาหํ ปุคฺคลํ อทฺทสํ อิธ ปาณาติปาตา ปฏิวิรตํ อทินฺนาทานา ปฏิวิรตํ ฯเปฯ สมฺมาทิฏฺฐิํ, กายสฺส เภทา ปรํ มรณา ปสฺสามิ สุคติํ สคฺคํ โลกํ อุปปนฺนนฺ”ติ, อิทมฺปิสฺส อนุชานามิ.\csromanspacer{} ยญฺจ โข โส เอวมาห “โย กิร โภ ปาณาติปาตา ปฏิวิรโต อทินฺนาทานา ปฏิวิรโต ฯเปฯ สมฺมาทิฏฺฐิ, สพฺโพ โส กายสฺส เภทา ปรํ มรณา สุคติํ สคฺคํ โลกํ อุปปชฺชตี”ติ, อิทมสฺส นานุชานามิ.\csromanspacer{} ยมฺปิ โส เอวมาห “เย เอวํ ชานนฺติ, เต สมฺมา ชานนฺติ.\csromanspacer{} เย อญฺญถา ชานนฺติ, มิจฺฉา เตสํ ญาณนฺ”ติ, อิทมฺปิสฺส นานุชานามิ.\csromanspacer{} ยมฺปิ โส ยเทว ตสฺส สามํ ญาตํ สามํ ทิฏฺฐํ สามํ วิทิตํ, ตเทว ตตฺถ ถามสา ปรามาสา อภินิวิสฺส โวหรติ “อิทเมว สจฺจํ โมฆมญฺญนฺ”ติ, อิทมฺปิสฺส นานุชานามิ.\csromanspacer{} ตํ กิสฺส เหตุ, อญฺญถา หิ อานนฺท ตถาคตสฺส มหากมฺมวิภงฺเค ญาณํ โหติ. (3)}

\prose{ตตฺรานนฺท ยฺวายํ สมโณ วา พฺราหฺมโณ วา เอวมาห “นตฺถิ กิร โภ กลฺยาณานิ กมฺมานิ, นตฺถิ สุจริตสฺส วิปาโก”ติ, อิทมสฺส นานุชานามิ.\csromanspacer{} ยญฺจ โข โส เอวมาห “อมาหํ ปุคฺคลํ อทฺทสํ อิธ ปาณาติปาตา ปฏิวิรตํ อทินฺนาทานา ปฏิวิรตํ ฯเปฯ สมฺมาทิฏฺฐิํ, กายสฺส เภทา ปรํ มรณา ปสฺสามิ อปายํ ทุคฺคติํ วินิปาตํ นิรยํ อุปปนฺนนฺ”ติ, อิทมสฺส อนุชานามิ.\csromanspacer{} ยญฺจ โข โส เอวมาห “โย กิร โภ ปาณาติปาตา ปฏิวิรโต อทินฺนาทานา ปฏิวิรโต ฯเปฯ สมฺมาทิฏฺฐิ, สพฺโพ โส กายสฺส เภทา ปรํ มรณา อปายํ ทุคฺคติํ วินิปายํ นิรยํ อุปปชฺชตี”ติ, อิทมสฺส นานุชานามิ.\csromanspacer{} ยญฺจ โข โส เอวมาห “เย เอวํ ชานนฺติ, เต สมฺมา ชานนฺติ.\csromanspacer{} เย อญฺญถา ชานนฺติ, มิจฺฉา เตสํ ญาณนฺ”ติ, อิทมฺปิสฺส นานุชานามิ.\csromanspacer{} ยมฺปิ โส ยเทว ตสฺส สามํ}

\csromanheader{2}{6. มหากมฺมวิภงฺคสุตฺต (136)}
\prose{\csromanfolio{257}ญาตํ สามํ ทิฏฺฐํ สามํ วิทิตํ, ตเทว ตตฺถ ถามสา ปรามาสา อภินิวิสฺส โวหรติ “อิทเมว สจฺจํ โมฆมญฺญนฺ”ติ, อิทมฺปิสฺส นานุชานามิ.\csromanspacer{} ตํ กิสฺส เหตุ, อญฺญถา หิ อานนฺท ตถาคตสฺส มหากมฺมวิภงฺเค ญาณํ โหติ. (4)}

\proseitem{303}{ตตฺรานนฺท ยฺวายํ ปุคฺคโล อิธ ปาณาติปาตี อทินฺนาทายี ฯเปฯ มิจฺฉาทิฏฺฐิ กายสฺส เภทา ปรํ มรณา อปายํ ทุคฺคติํ วินิปาตํ นิรยํ อุปปชฺชติ.\csromanspacer{} ปุพฺเพ วาสฺส ตํ กตํ โหติ ปาปกมฺมํ ทุกฺขเวทนียํ, ปจฺฉา วาสฺส ตํ กตํ โหติ ปาปกมฺมํ ทุกฺขเวทนียํ, มรณกาเล วาสฺส โหติ มิจฺฉาทิฏฺฐิ สมตฺตา สมาทินฺนา, เตน โส กายสฺส เภทา ปรํ มรณา อปายํ ทุคฺคติํ วินิปาตํ นิรยํ อุปปชฺชติ.\csromanspacer{} ยญฺจ โข โส อิธ ปาณาติปาตี โหติ, อทินฺนาทายี โหติ ฯเปฯ มิจฺฉาทิฏฺฐิ โหติ, ตสฺส ทิฏฺเฐว ธมฺเม วิปากํ ปฏิสํเวเทติ อุปปชฺช วา\footnote{อุปปชฺชํ วา (สี, อิ), อุปปชฺเช วา (สฺยา, กํ, ก) อุปปชฺชิตฺวาติ สํวณฺณนายสํสนฺเทตพฺพา.} อปเร วา ปริยาเย. (1)}

\prose{ตตฺรานนฺท ยฺวายํ ปุคฺคโล อิธ ปาณาติปาตี อทินฺนาทายี ฯเปฯ มิจฺฉาทิฏฺฐิ กายสฺส เภทา ปรํ มรณา สุคติํ สคฺคํ โลกํ อุปปชฺชติ.\csromanspacer{} ปุพฺเพ วาสฺส ตํ กตํ โหติ กลฺยาณกมฺมํ สุขเวทนียํ, ปจฺฉา วาสฺส ตํ กตํ โหติ กลฺยาณกมฺมํ สุขเวทนียํ, มรณกาเล วาสฺส โหติ สมฺมาทิฏฺฐิ สมตฺตา สมาทินฺนา.\csromanspacer{} เตน โส กายสฺส เภทา ปรํ มรณา สุคติํ สคฺคํ โลกํ อุปปชฺชติ.\csromanspacer{} ยญฺจ โข โส อิธ ปาณาติปาตี โหติ, อทินฺนาทายี โหติ ฯเปฯ มิจฺฉาทิฏฺฐิ โหติ, ตสฺส ทิฏฺเฐว ธมฺเม วิปากํ ปฏิสํเวเทติ อุปปชฺช วา อปเร วา ปริยาเย. (2)}

\prose{ตตฺรานนฺท ยฺวายํ ปุคฺคโล อิธ ปาณาติปาตา ปฏิวิรโต อทินฺนาทานา ปฏิวิรโต ฯเปฯ สมฺมาทิฏฺฐิ กายสฺส เภทา ปรํ มรณา สุคติํ สคฺคํ โลกํ อุปปชฺชติ.\csromanspacer{} ปุพฺเพ วาสฺส ตํ กตํ โหติ กลฺยาณกมฺมํ สุขเวทนียํ, ปจฺฉา วาสฺส ตํ กตํ โหติ กลฺยาณกมฺมํ สุขเวทนียํ, มรณกาเล วาสฺส โหติ สมฺมาทิฏฺฐิ สมตฺตา สมาทินฺนา.\csromanspacer{} เตน โส อิธ ปาณาติปาตา ปฏิวิรโต โหติ, อทินฺนาทานา ปฏิวิรโต โหติ ฯเปฯ สมฺมาทิฏฺฐิ โหติ, ตสฺส ทิฏฺเฐว ธมฺเม วิปากํ ปฏิสํเวเทติ อุปปชฺช วา อปเร วา ปริยาเย. (3)}

\prose{\csromanfolio{258}ตตฺรานนฺท ยฺวายํ ปุคฺคโล อิธ ปาณาติปาตา ปฏิวิรโต อทินฺนาทานา ปฏิวิรโต ฯเปฯ สมฺมาทิฏฺฐิ กายสฺส เภทา ปรํ มรณา อปายํ ทุคฺคติํ วินิปาตํ นิรยํ อุปปชฺชติ.\csromanspacer{} ปุพฺเพ วาสฺส ตํ กตํ โหติ ปาปกมฺมํ ทุกฺขเวทนียํ, ปจฺฉา วาสฺส ตํ กตํ โหติ ปาปกมฺมํ ทุกฺขเวทนียํ, มรณกาเล วาสฺส โหติ มิจฺฉาทิฏฺฐิ สมตฺตา สมาทินฺนา.\csromanspacer{} เตน โส กายสฺส เภทา ปรํ มรณา อปายํ ทุคฺคติํ วินิปาตํ นิรยํ อุปปชฺชติ.\csromanspacer{} ยญฺจ โข โส อิธ ปาณาติปาตา ปฏิวิรโต โหติ, อทินฺนาทานา ปฏิวิรโต โหติ ฯเปฯ สมฺมาทิฏฺฐิ โหติ, ตสฺส ทิฏฺเฐว ธมฺเม วิปากํ ปฏิสํเวเทติ อุปปชฺช วา อปเร วา ปริยาเย. (4)}

\prose{อิติ โข อานนฺท อตฺถิ กมฺมํ อภพฺพํ อภพฺพาภาสํ, อตฺถิ กมฺมํ อภพฺพํ ภพฺพาภาสํ, อตฺถิ กมฺมํ ภพฺพญฺเจว ภพฺพาภาสญฺจ, อตฺถิ กมฺมํ ภพฺพํ อภพฺพาภาสนฺติ.}

\prose{อิทมโวจ ภควา.\csromanspacer{} อตฺตมโน อายสฺมา อานนฺโท ภควโต ภาสิตํ อภินนฺทีติ.}

\nitthitamruled{มหากมฺมวิภงฺคสุตฺตํ นิฏฺฐิตํ ฉฏฺฐํ.}
\csromanmatikaanchor{mtk.44}
\csromantocmarkref{subsection}{7. สฬายตนวิภงฺคสุตฺต}{mtk.44}
\bookmark[dest={mtk.44},level=2]{7. สฬายตนวิภงฺคสุตฺต}
\csromanheader{2}{7. สฬายตนวิภงฺคสุตฺต}
\proseitem{304}{เอวํ เม สุตํ\csromandash{}เอกํ สมยํ ภควา สาวตฺถิยํ วิหรติ เชตวเน อนาถปิณฺฑิกสฺส อาราเม.\csromanspacer{} ตตฺร โข ภควา ภิกฺขู อามนฺเตสิ “ภิกฺขโว”ติ. “ภทนฺเต”ติ เต ภิกฺขู ภควโต ปจฺจสฺโสสุํ.\csromanspacer{} ภควา เอตทโวจ\csromandash{}สฬายตนวิภงฺคํ โว ภิกฺขเว เทเสสฺสามิ, ตํ สุณาถ สาธุกํ มนสิ กโรถ, ภาสิสฺสามีติ. “เอวํ ภนฺเต”ติ โข เต ภิกฺขู ภควโต ปจฺจสฺโสสุํ.\csromanspacer{} ภควา เอตทโวจ\csromandash{}}

\prose{ฉ อชฺฌตฺติกานิ อายตนานิ เวทิตพฺพานิ, ฉ พาหิรานิ อายตนานิ เวทิตพฺพานิ, ฉ วิญฺญาณกายา เวทิตพฺพา, ฉ ผสฺสกายา เวทิตพฺพา, อฏฺฐารส มโนปวิจารา เวทิตพฺพา, ฉตฺติํส สตฺตปทา เวทิตพฺพา.\csromanspacer{} ตตฺร อิทํ นิสฺสาย อิทํ ปชหถ, ตโย สติปฏฺฐานา ยทริโย เสวติ,}

\csromanheader{2}{7. สฬายตนวิภงฺคสุตฺต (137)}
\prose{\csromanfolio{259}ยทริโย เสวมาโน สตฺถา คณมนุสาสิตุมรหติ, โส วุจฺจติ โยคฺคาจริยานํ\footnote{โย คาจริยานํ (ก)} อนุตฺตโร ปุริสทมฺมสารถีติ.\csromanspacer{} อยมุทฺเทโส สฬายตนวิภงฺคสฺส.}

\proseitem{305}{“ฉ อชฺฌตฺติกานิ อายตนานิ เวทิตพฺพานี”ติ อิติ โข ปเนตํ วุตฺตํ, กิญฺเจตํ ปฏิจฺจ วุตฺตํ.\csromanspacer{} จกฺขายตนํ โสตายตนํ ฆานายตนํ ชิวฺหายตนํ กายายตนํ มนายตนํ. “ฉ อชฺฌตฺติกานิ อายตนานิ เวทิตพฺพานี”ติ อิติ ยํ ตํ วุตฺตํ, อิธเมตํ ปฏิจฺจ วุตฺตํ. (1)}

\prose{“ฉ พาหิรานิ อายตนานิ เวทิตพฺพานี”ติ อิติ โข ปเนตํ วุตฺตํ, กิญฺเจตํ ปฏิจฺจ วุตฺตํ.\csromanspacer{} รูปายตนํ สทฺทายตนํ คนฺธายตนํ รสายตนํ โผฏฺฐพฺพายตนํ ธมฺมายตนํ. “ฉ พาหิรานิ อายตนานิ เวทิตพฺพานี”ติ อิติ ยํ ตํ วุตฺตํ, อิทเมตํ ปฏิจฺจ วุตฺตํ. (2)}

\prose{“ฉ วิญฺญาณกายา เวทิตพฺพา”ติ อิติ โข ปเนตํ วุตฺตํ, กิญฺเจตํ ปฏิจฺจ วุตฺตํ.\csromanspacer{} จกฺขุวิญฺญาณํ โสตวิญฺญาณํ ฆานวิญฺญาณํ ชิวฺหาวิญฺญาณํ กายวิญฺญาณํ มโนวิญฺญาณํ. “ฉ วิญฺญาณกายา เวทิตพฺพา”ติ อิติ ยํ ตํ วุตฺตํ, อิทเมตํ ปฏิจฺจ วุตฺตํ. (3)}

\prose{“ฉ ผสฺสกายา เวทิตพฺพา”ติ อิติ โข ปเนตํ วุตฺตํ, กิญฺเจตํ ปฏิจฺจ วุตฺตํ.\csromanspacer{} จกฺขุสมฺผสฺโส โสตสมฺผสฺโส ฆานสมฺผสฺโส ชิวฺหาสมฺผสฺโส กายสมฺผสฺโส มโนสมฺผสฺโส. “ฉ ผสฺสกายา เวทิตพฺพา”ติ อิติ ยํ ตํ วุตฺตํ, อิทเมตํ ปฏิจฺจ วุตฺตํ. (4)}

\prose{“อฏฺฐารส มโนปวิจารา เวทิตพฺพา”ติ อิติ โข ปเนตํ วุตฺตํ, กิญฺเจตํ ปฏิจฺจ วุตฺตํ.\csromanspacer{} จกฺขุนา รูปํ ทิสฺวา โสมนสฺสฏฺฐานียํ รูปํ อุปวิจรติ, โทมนสฺสฏฺฐานียํ รูปํ อุปวิจรติ, อุเปกฺขาฏฺฐานียํ รูปํ อุปวิจรติ.\csromanspacer{} โสเตน สทฺทํ สุตฺวา ฯเปฯ.\csromanspacer{} ฆาเนน คนฺธํ ฆายิตฺวา.\csromanspacer{} ชิวฺหาย รสํ สายิตฺวา.\csromanspacer{} กาเยน โผฏฺฐพฺพํ ผุสิตฺวา.\csromanspacer{} มนสา ธมฺมํ วิญฺญาย โสมนสฺสฏฺฐานียํ ธมฺมํ อุปวิจรติ, โทมนสฺสฏฺฐานียํ ธมฺมํ อุปวิจรติ, อุเปกฺขาฏฺฐานียํ ธมฺมํ อุปวิจรติ.\csromanspacer{} อิติ ฉ โสมนสฺสูปวิจารา ฉ โทมนสฺสูปวิจารา ฉ อุเปกฺขูปวิจารา. “อฏฺฐารส มโนปวิจารา เวทิตพฺพา”ติ อิติ ยํ ตํ วุตฺตํ, อิทเมตํ ปฏิจฺจ วุตฺตํ. (5)}

\proseitem{306}{\csromanfolio{260}“ฉตฺติํส สตฺตปทา เวทิตพฺพา”ติ อิติ โข ปน ปเนตํ วุตฺตํ, กิญฺเจตํ ปฏิจฺจ วุตฺตํ.\csromanspacer{} ฉ เคหสิตานิ\footnote{เคหสฺสิตานิ (?)} โสมนสฺสานิ ฉ เนกฺขมฺมสิตานิ\footnote{เนกฺขมฺมสฺสิตานิ (ฏีกา)} โสมนสฺสานิ ฉ เคหสิตานิ โทมนสฺสานิ ฉ เนกฺขมฺมสิตานิ โทมนสฺสานิ ฉ เคหสิตา อุเปกฺขา ฉ เนกฺขมฺมสิตา อุเปกฺขา.\csromanspacer{} ตตฺถ กตมานิ ฉ เคหสิตานิ โสมนสฺสานิ.\csromanspacer{} จกฺขุวิญฺเญยฺยานํ รูปานํ อิฏฺฐานํ กนฺตานํ มนาปานํ มโนรมานํ โลกามิสปฏิสํยุตฺตานํ ปฏิลาภํ วา ปฏิลาภโต สมนุปสฺสโต ปุพฺเพ วา ปฏิลทฺธปุพฺพํ อตีตํ นิรุทฺธํ วิปริณตํ สมนุสฺสรโต อุปฺปชฺชติ โสมนสฺสํ.\csromanspacer{} ยํ เอวรูปํ โสมนสฺสํ, อิทํ วุจฺจติ เคหสิตํ โสมนสฺสํ.\csromanspacer{} โสตวิญฺเญยฺยานํ สทฺทานํ.\csromanspacer{} ฆานวิญฺเญยฺยานํ คนฺธานํ.\csromanspacer{} ชิวฺหาวิญฺเญยฺยานํ รสานํ.\csromanspacer{} กายวิญฺเญยฺยานํ โผฏฺฐพฺพานํ.\csromanspacer{} มโนวิญฺเญยฺยานํ ธมฺมานํ อิฏฺฐานํ กนฺตานํ มนาปานํ ฯเปฯ โสมนสฺสํ.\csromanspacer{} ยํ เอวรูปํ โสมนสฺสํ, อิทํ วุจฺจติ เคหสิตํ โสมนสฺสํ.\csromanspacer{} อิมานิ ฉ เคหสิตานิ โสมนสฺสานิ.}

\prose{ตตฺถ กตมานิ ฉ เนกฺขมฺมสิตานิ โสมนสฺสานิ.\csromanspacer{} รูปานํเตฺวว อนจฺจตํ วิทิตฺวา วิปริณามวิราคนิโรธํ\footnote{วิปริณามํ วิราคํ นิโรธํ (ก)} “ปุพฺเพ เจว รูปา เอตรหิ จ, สพฺเพ เต รูปา อนิจฺจา ทุกฺขา วิปริณามธมฺมา”ติ เอวเมตํ ยถาภูตํ สมฺมปฺปญฺญาย ปสฺสโต อุปฺปชฺชติ โสมนสฺสํ.\csromanspacer{} ยํ เอวรูปํ โสมนสฺสํ, อิทํ วุจฺจติ เนกฺขมฺมสิตํ โสมนสฺสํ.\csromanspacer{} สทฺทานํเตฺวว.\csromanspacer{} คนฺธานํเตฺวว.\csromanspacer{} รสานํเตฺวว.\csromanspacer{} โผฏฺฐพฺพานํเตฺวว.\csromanspacer{} ธมฺมานํเตฺวว อนิจฺจตํ วิทิตฺวา วิปริณามวิราคนิโรธํ “ปุพฺเพ เจว ธมฺมา เอตรหิ จ, สพฺเพ เต ธมฺมา อนิจฺจา ทุกฺขา วิปริณามธมฺมา”ติ เอวเมตํ ยถาภูตํ สมฺมปฺปญฺญาย ปสฺสโต อุปฺปชฺชติ โสมนสฺสํ.\csromanspacer{} ยํ เอวรูปํ โสมนสฺสํ, อิทํ วุจฺจติ เนกฺขมฺมสิตํ โสมนสฺสํ.\csromanspacer{} อิมานิ ฉ เนกฺขมฺมสิตานิ โสมนสฺสานิ.}

\proseitem{307}{ตตฺถ กตมานิ ฉ เคหสิตานิ โทมนสฺสานิ.\csromanspacer{} จกฺขุวิญฺเญยฺยานํ รูปานํ ฯเปฯ.\csromanspacer{} โสตวิญฺเญยฺยานํ สทฺทานํ.\csromanspacer{} ฆานวิญฺเญยฺยานํ คนฺธานํ.\csromanspacer{} ชิวฺหาวิญฺเญยฺยานํ รสานํ.\csromanspacer{} กายวิญฺเญยฺยานํ โผฏฺฐพฺพานํ.\csromanspacer{} มโนวิญฺเญยฺยานํ ธมฺมานํ อิฏฺฐานํ กนฺตานํ มนาปานํ มโนรมานํ โลกามิสปฏิสํยุตฺตานํ อปฺปฏิลาภํ วา อปฺปฏิลาภโต สมนุปสฺสโต ปุพฺเพ วา อปฺปฏิลทฺธปุพฺพํ อตีตํ นิรุทฺธํ วิปริณตํ สมนุสฺสรโต อุปฺปชฺชติ โทมนสฺสํ.\csromanspacer{} ยํ เอวรูปํ}

\csromanheader{2}{7. สฬายตนวิภงฺคสุตฺต (137)}
\prose{\csromanfolio{261}โทมนสฺสํ, อิทํ วุจฺจติ เคหสิตํ โทมนสฺสํ.\csromanspacer{} อิมานิ ฉ เคหสิตานิ โทมนสฺสานิ.}

\prose{ตตฺถ กตมานิ ฉ เนกฺขมฺมสิตานิ โทมนสฺสานิ.\csromanspacer{} รูปานํเตฺวว อนิจฺจตํ วิทิตฺวา วิปริณามวิราคนิโรธํ “ปุพฺเพ เจว รูปา เอตรหิ จ, สพฺเพ เต รูปา อนิจฺจา ทุกฺขา วิปริณามธมฺมา”ติ เอวเมตํ ยถาภูตํ สมฺมปฺปญฺญาย ทิสฺวา อนุตฺตเรสุ วิโมกฺเขสุ ปิหํ อุปฏฺฐาเปติ “กุทาสฺสุ\footnote{กทาสฺสุ (สฺยา, กํ, อิ)} นามาหํ ตทายตนํ อุปสมฺปชฺช วิหริสฺสามิ, ยทริยา เอตรหิ อายตนํ อุปสมฺปชฺช วิหรนฺตี”ติ.\csromanspacer{} อิติ อนุตฺตเรสุ วิโมกฺเขสุ ปิหํ อุปฏฺฐาปยโต อุปฺปชฺชติ ปิหปจฺจยา โทมนสฺสํ.\csromanspacer{} ยํ เอวรูปํ โทมนสฺสํ, อิทํ วุจฺจติ เนกฺขมฺมสิตํ โทมนสฺสํ สทฺทานํเตฺวว ฯเปฯ.\csromanspacer{} คนฺธานํเตฺวว.\csromanspacer{} รสานํเตฺวว.\csromanspacer{} โผฏฺฐพฺพานํเตฺวว.\csromanspacer{} ธมฺมานํเตฺวว อนิจฺจตํ วิทิตฺวา วิปริณามวิราคนิโรธํ “ปุพฺเพ เจว ธมฺมา เอตรหิ จ, สพฺเพ เต ธมฺมา อนิจฺจา ทุกฺขา วิปริณามธมฺมา”ติ เอวเมตํ ยถาภูตํ สมฺมปฺปญฺญาย ทิสฺวา อนุตฺตเรสุ วิโมกฺเขสุ ปิหํ อุปฏฺฐาเปติ “กุทาสฺสุ นามาหํ ตทายตนํ อุปสมฺปชฺช วิหริสฺสามิ, ยทริยา เอตรหิ อายตนํ อุปสมฺปชฺช วิหรนฺตี”ติ.\csromanspacer{} อิติ อนุตฺตเรสุ วิโมกฺเขสุ ปิหํ อุปฏฺฐาปยโต อุปฺปชฺชติ ปิหปจฺจยา โทมนสฺสํ.\csromanspacer{} ยํ เอวรูปํ โทมนสฺสํ, อิทํ วุจฺจติ เนกฺขมฺมสิตํ โทมนสฺสํ.\csromanspacer{} อิมานิ ฉ เนกฺขมฺมสิตานิ โทมนสฺสานิ.}

\proseitem{308}{ตตฺถ กตมา ฉ เคหสิตา อุเปกฺขา.\csromanspacer{} จกฺขุนา รูปํ ทิสฺวา อุปฺปชฺชติ อุเปกฺขา พาลสฺส มูฬฺหสฺส ( )\footnote{(มนฺทสฺส) (ก)} ปุถุชฺชนสฺส อโนธิชินสฺส อวิปากชินสฺส อนาทีนวทสฺสวิโน อสฺสุตวโต ปุถุชฺชนสฺส.\csromanspacer{} ยา เอวรูปา อุเปกฺขา รูปํ สา นาติวตฺตติ, ตสฺมา สา\footnote{สายํ (ก)} อุเปกฺขา เคหสิตาติ วุจฺจติ.\csromanspacer{} โสเตน สทฺทํ สุตฺวา.\csromanspacer{} ฆาเนน คนฺธํ ฆายิตฺวา.\csromanspacer{} ชิวฺหาย รสํ สายิตฺวา.\csromanspacer{} กาเยน โผฏฺฐพฺพํ ผุสิตฺวา.\csromanspacer{} มนสา ธมฺมํ วิญฺญาย อุปฺปชฺชติ อุเปกฺขา พาลสฺส มูฬฺหสฺส ปุถุชฺชนสฺส อโนธิชินสฺส อวิปากชินสฺส อนาทีนวทสฺสาวิโน อสฺสุตวโต ปุถุชฺชนสฺส.\csromanspacer{} ยา เอวรูปา อุเปกฺขา, ธมฺมํ สา นาติวตฺตติ.\csromanspacer{} ตสฺมา สา อุเปกฺขา เคหสิตาติ วุจฺจติ.\csromanspacer{} อิมา ฉ เคหสิตา อุเปกฺขา.}

\prose{ตตฺถ กตมา ฉ เนกฺขมฺมสิตา อุเปกฺขา.\csromanspacer{} รูปานํเตฺวว อนิจฺจตํ วิทิตฺวา วิปริณามวิราคนิโรธํ “ปุพฺเพ เจว รูปา เอตรหิ จ, สพฺเพ เต \csromanfolio{262}รูปา อนิจฺจา ทุกฺขา วิปริณามธมฺมา”ติ เอวเมตํ ยถาภูตํ สมฺมปฺปญฺญาย ปสฺสโต อุปฺปชฺชติ อุเปกฺขา, ยา เอวรูปา อุเปกฺขา, รูปํ สา อติวตฺตติ.\csromanspacer{} ตสฺมา สา อุเปกฺขา เนกฺขมฺมสิตาติ วุจฺจติ.\csromanspacer{} สทฺทานํ เตฺวว.\csromanspacer{} คนฺธานํเตฺวว.\csromanspacer{} รสานํเตฺวว.\csromanspacer{} โผฏฺฐพฺพานํเตฺวว.\csromanspacer{} ธมฺมานํเตฺวว อนิจฺจตํ วิทิตฺวา วิปริณามวิราคนิโรธํ “ปุพฺเพ เจว ธมฺมา เอตรหิ จ, สพฺเพ เต ธมฺมา อนิจฺจา ทุกฺขา วิปริณามธมฺมา”ติ เอวเมตํ ยถาภูตํ สมฺมปฺปญฺญาย ปสฺสโต อุปฺปชฺชติ อุเปกฺขา.\csromanspacer{} ยา เอวรูปา อุเปกฺขา, ธมฺมํ สา อติวตฺตติ.\csromanspacer{} ตสฺมา สา อุเปกฺขา เนกฺขมฺมสิตาติ วุจฺจติ.\csromanspacer{} อิมา ฉ เนกฺขมฺมสิตา อุเปกฺขา. “ฉตฺติํส สตฺตปทา เวทิตพฺพา”ติ อิติ ยํ ตํ วุตฺตํ, อิทเมตํ ปฏิจฺจ วุตฺตํ.}

\proseitem{309}{“ตตฺร อิทํ นิสฺสาย อิทํ ปชหถา”ติ อิติ โข ปเนตํ วุตฺตํ, กิญฺเจตํ ปฏิจฺจ วุตฺตํ.\csromanspacer{} ตตฺร ภิกฺขเว ยานิ ฉ เนกฺขมฺมสิตานิ โสมนสฺสานิ, ตานิ นิสฺสาย ตานิ อาคมฺม, ยานิ ฉ เคหสิตานิ โสมนสฺสานิ, ตานิ ปชหถ ตานิ สมติกฺกมถ.\csromanspacer{} เอวเมเตสํ ปหานํ โหติ, เอวเมเตสํ สมติกฺกโม โหติ.}

\prose{ตตฺร ภิกฺขเว ยานิ ฉ เนกฺขมฺมสิตานิ โทมนสฺสานิ, ตานิ นิสฺสาย ตานิ อาคมฺม, ยานิ ฉ เคหสิตานิ โทมนสฺสานิ, ตานิ ปชหถ ตานิ สม- ติกฺกมถ.\csromanspacer{} เอวเมเตสํ ปหานํ โหติ, เอวเมเตสํ สมติกฺกโม โหติ.}

\prose{ตตฺร ภิกฺขเว ยา ฉ เนกฺขมฺมสิตา อุเปกฺขา, ตา นิสฺสาย ตา อาคมฺม, ยา ฉ เคหสิตา อุเปกฺขา, ตา ปชหถ ตา สมติกฺกมถ.\csromanspacer{} เอวเมตาสํ ปหานํ โหติ, เอวเมตาสํ สมติกฺกโม โหติ.}

\prose{ตตฺร ภิกฺขเว ยานิ ฉ เนกฺขมฺมสิตานิ โสมนสฺสานิ, ตานิ นิสฺสาย ตานิ อาคมฺม.\csromanspacer{} ยานิ ฉ เนกฺขมฺมสิตานิ โทมนสฺสานิ, ตานิ ปชหถ ตานิ สมติกฺกมถ.\csromanspacer{} เอวเมเตสํ ปหานํ โหติ, เอวเมเตสํ สมติกฺกโม โหติ.}

\prose{ตตฺร ภิกฺขเว ยา ฉ เนกฺขมฺมสิตา อุเปกฺขา, ตา นิสฺสาย ตา อาคมฺม, ยานิ ฉ เนกฺขมฺมสิตานิ โสมนสฺสานิ, ตานิ ปชหถ ตานิ สมติกฺกมถ.\csromanspacer{} เอวเมเตสํ ปหานํ โหติ, เอวเมเตสํ สมติกฺกโม โหติ.}

\csromanheader{2}{7. สฬายตนวิภงฺคสุตฺต (137)}
\proseitem{310}{\csromanfolio{263}อตฺถิ ภิกฺขเว อุเปกฺขา นานตฺตา นานตฺตสิตา, อตฺถิ อุเปกฺขา เอกตฺตา เอกตฺตสิตา.\csromanspacer{} กตมา จ ภิกฺขเว อุเปกฺขา นานตฺตา นานตฺตสิตา.\csromanspacer{} อตฺถิ ภิกฺขเว อุเปกฺขา รูเปสุ, อตฺถิ สทฺเทสุ, อตฺถิ คนฺเธสุ, อตฺถิ รเสสุ, อตฺถิ โผฏฺฐพฺเพสุ.\csromanspacer{} อยํ ภิกฺขเว อุเปกฺขา นานตฺตา นานตฺตสิตา.\csromanspacer{} กตมา จ ภิกฺขเว อุเปกฺขา เอกตฺตา เอกตฺตสิตา.\csromanspacer{} อตฺถิ ภิกฺขเว อุเปกฺขา อากาสานญฺจายตนนิสฺสิตา, อตฺถิ วิญฺญาณญฺจายตนนิสฺสิตา, อตฺถิ อากิญฺจญฺญายตนนิสฺสิตา, อตฺถิ เนวสญฺญานาสญฺญายตนนิสฺสิตา.\csromanspacer{} อยํ ภิกฺขเว อุเปกฺขา เอกตฺตา เอกตฺตสิตา.}

\prose{ตตฺร ภิกฺขเว ยายํ อุเปกฺขา เอกตฺตา เอกตฺตสิตา, ตํ นิสฺสาย ตํ อาคมฺม, ยายํ อุเปกฺขา นานตฺตา นานตฺตสิตา, ตํ ปชหถ ตํ สมติกฺกมถ.\csromanspacer{} เอวเมติสฺสา ปหานํ โหติ, เอวเมติสฺสา สมติกฺกโม โหติ.}

\prose{อตมฺมยตํ ภิกฺขเว นิสฺสาย อตมฺมยตํ อาคมฺม ยายํ อุเปกฺขา เอกตฺตา เอกตฺตสิตา, ตํ ปชหถ ตํ สมติกฺกมถ.\csromanspacer{} เอวเมติสฺสา ปหานํ โหติ, เอวเมติสฺสา สมติกฺกโม โหติ. “ตตฺร อิทํ นิสฺสาย อิทํ ปชหถา”ติ อิติ ยํ ตํ วุตฺตํ, อิทเมตํ ปฏิจฺจ วุตฺตํ.}

\proseitem{311}{“ตโย สติปฏฺฐานา ยทริโย เสวติ, ยทริโย เสวมาโน สตฺถา คณมนุสาสิตุมรหตี”ติ อิติ โข ปเนตํ วุตฺตํ, กิญฺเจตํ ปฏิจฺจ วุตฺตํ.\csromanspacer{} อิธ ภิกฺขเว สตฺถา สาวกานํ ธมฺมํ เทเสติ อนุกมฺปโก หิเตสี อนุกมฺปํ อุปาทาย “อิทํ โว หิตาย อิทํ โว สุขายา”ติ.\csromanspacer{} ตสฺส สาวกา น สุสฺสูสนฺติ, น โสตํ โอทหนฺติ, น อญฺญา จิตฺตํ อุปฏฺฐเปนฺติ, โวกฺกมฺม จ สตฺถุสาสนา วตฺตนฺติ.\csromanspacer{} ตตฺร ภิกฺขเว ตถาคโต น เจว อนตฺตมโน โหติ, น จ อนตฺตมนตํ ปฏิสํเวเทติ, อนวสฺสุโต จ วิหรติ สโต สมฺปชาโน.\csromanspacer{} อิทํ ภิกฺขเว ปฐมํ สติปฏฺฐานํ.\csromanspacer{} ยทริโย เสวติ, ยทริโย เสวมาโน สตฺถา คณมนุสาสิตุมรหติ.}

\prose{ปุน จปรํ ภิกฺขเว สตฺถา สาวกานํ ธมฺมํ เทเสติ อนุกมฺปโก หิเตสี อนุกมฺปํ อุปาทาย “อิทํ โว หิตาย อิทํ โว สุขายา”ติ.\csromanspacer{} ตสฺส เอกจฺเจ สาวกา น สุสฺสูสนฺติ, น โสตํ โอทหนฺติ, น อญฺญา จิตฺตํ อุปฏฺฐเปนฺติ, โวกฺกมฺม จ สตฺถุสาสนา วตฺตนฺติ.\csromanspacer{} เอกจฺเจ สาวกา สุสฺสูสนฺติ, โสตํ โอทหนฺติ, อญฺญา จิตฺตํ อุปฏฺฐเปนฺติ, น จ โวกฺกมฺม สตฺถุสาสนา \csromanfolio{264}วตฺตนฺติ.\csromanspacer{} ตตฺร ภิกฺขเว ตถาคโต น เจว อนตฺตมโน โหติ, น จ อนตฺตมนตํ ปฏิสํเวเทติ, น จ อตฺตมโน โหติ, น จ อตฺตมนตํ ปฏิสํเวเทติ.\csromanspacer{} อนตฺตมนตา จ อตฺตมนตา จ ตทุภยํ อภินิวชฺเชตฺวา อุเปกฺขโก วิหรติ สโต สมฺปชาโน.\csromanspacer{} อิทํ วุจฺจติ ภิกฺขเว ทุติยํ สติปฏฺฐานํ.\csromanspacer{} ยทริโย เสวติ, ยทริโย เสวมาโน สตฺถา คณมนุสาสิตุมรหติ.}

\prose{ปุน จปรํ ภิกฺขเว สตฺถา สาวกานํ ธมฺมํ เทเสติ อนุกมฺปโก หิเตสี อนุกมฺปํ อุปาทาย “อิทํ โว หิตาย อิทํ โว สุขายา”ติ.\csromanspacer{} ตสฺส สาวกา สุสฺสูสนฺติ, โสตํ โอทหนฺติ, อญฺญาจิตฺตํ อุปฏฺฐเปนฺติ, น จ โวกฺกมฺม สตฺถุสาสนา วตฺตนฺติ.\csromanspacer{} ตตฺร ภิกฺขเว ตถาคโต อตฺตมโน เจว โหติ, อตฺตมนตญฺจ ปฏิสํเวเทติ, อนวสฺสุโต จ วิหรติ สโต สมฺปชาโน.\csromanspacer{} อิทํ วุจฺจติ ภิกฺขเว ตติยํ สติปฏฺฐานํ.\csromanspacer{} ยทริโย เสวติ, ยทริโย เสวมาโน สตฺถา คณมนุสาสิตุมรหติ. “ตโย สติปฏฺฐานา, ยทริโย เสวติ, ยทริโย เสวมาโน สตฺถา คณมนุสาสิตุมรหตี”ติ อิติ ยํ ตํ วุตฺตํ, อิทเมตํ ปฏิจฺจ วุตฺตํ.}

\proseitem{312}{“โส วุจฺจติ โยคฺคาจริยานํ อนุตฺตโร ปุริสทมฺมสารถี”ติ อิติ โข ปเนตํ วุตฺตํ, กิญฺเจตํ ปฏิจฺจ วุตฺตํ.\csromanspacer{} หตฺถิทมเกน ภิกฺขเว หตฺถิทมฺโม สาริโต เอกํเยว ทิสํ ธาวติ ปุรตฺถิมํ วา ปจฺฉิมํ วา อุตฺตรํ วา ทกฺขิณํ วา.\csromanspacer{} อสฺสทมเกน ภิกฺขเว อสฺสทมฺโม สาริโต เอกํเยว ทิสํ ธาวติ ปุรตฺถิมํ วา ปจฺฉิมํ วา อุตฺตรํ วา ทกฺขิณํ วา.\csromanspacer{} โคทมเกน ภิกฺขเว โคทมฺโม สาริโต เอกํเยว ทิสํ ธาวติ ปุรตฺถิมํ วา ปจฺฉิมํ วา อุตฺตรํ วา ทกฺขิณํ วา.\csromanspacer{} ตถาคเตน หิ ภิกฺขเว อรหตา สมฺมาสมฺพุทฺเธน ปุริสทมฺโม สาริโต อฏฺฐ ทิสา วิธาวติ.\csromanspacer{} รูปี รูปานิ ปสฺสติ, อยํ เอกา ทิสา.\csromanspacer{} อชฺฌตฺตํ อรูปสญฺญี พหิทฺธา รูปานิ ปสฺสติ, อยํ ทุติยา ทิสา.\csromanspacer{} สุภนฺเตฺวว อธิมุตฺโต โหติ, อยํ ตติยา ทิสา.\csromanspacer{} สพฺพโส รูปสญฺญานํ สมติกฺกมา ปฏิฆสญฺญานํ อตฺถงฺคมา นานตฺตสญฺญานํ อมนสิการา “อนนฺโต อากาโส”ติ อากาสานญฺจายตนํ อุปสมฺปชฺช วิหรติ, อยํ จตุตฺถี ทิสา.\csromanspacer{} สพฺพโส อากาสานญฺจายตนํ สมติกฺกมฺม “อนนฺตํ วิญฺญาณนฺ”ติ วิญฺญาณญฺจายตนํ อุปสมฺปชฺช วิหรติ, อยํ ปญฺจมี ทิสา.\csromanspacer{} สพฺพโส วิญฺญาณญฺจายตนํ สมติกฺกมฺม “นตฺถิ กิญฺจี”ติ อากิญฺจญฺญายตนํ}

\csromanheader{2}{8. อุทฺเทสวิภงฺคสุตฺต (138)}
\prose{\csromanfolio{265}อุปสมฺปชฺช วิหรติ, อยํ ฉฏฺฐี ทิสา.\csromanspacer{} สพฺพโส อากิญฺจญฺญายตนํ สมติกฺกมฺม เนวสญฺญานาสญฺญายตนํ อุปสมฺปชฺช วิหรติ, อยํ สตฺตมี ทิสา.\csromanspacer{} สพฺพโส เนวสญฺญานาสญฺญายตนํ สมติกฺกมฺม สญฺญาเวทยิตนิโรธํ อุปสมฺปชฺช วิหรติ, อยํ อฏฺฐมี ทิสา.\csromanspacer{} ตถาคเตน ภิกฺขเว อรหตา สมฺมาสมฺพุทฺเธน ปุริสทมฺโม สาริโต.\csromanspacer{} อิมา อฏฺฐ ทิสา วิธาวติ, โส วุจฺจติ “โยคฺคาจริยานํ อนุตฺตโร ปุริสทมฺมสารถี”ติ อิติ ยํ ตํ วุตฺตํ, อิทเมตํ ปฏิจฺจ วุตฺตนฺติ.}

\prose{อิทมโวจ ภควา.\csromanspacer{} อตฺตมนา เต ภิกฺขู ภควโต ภาสิตํ อภินนฺทุนฺติ.}

\nitthitamruled{สฬายตนวิภงฺคสุตฺตํ นิฏฺฐิตํ สตฺตมํ.}
\csromanmatikaanchor{mtk.45}
\csromantocmarkref{subsection}{8. อุทฺเทสวิภงฺคสุตฺต}{mtk.45}
\bookmark[dest={mtk.45},level=2]{8. อุทฺเทสวิภงฺคสุตฺต}
\csromanheader{2}{8. อุทฺเทสวิภงฺคสุตฺต}
\proseitem{313}{เอวํ เม สุตํ\csromandash{}เอกํ สมยํ ภควา สาวตฺถิยํ วิหรติ เชตวเน อนาถปิณฺฑิกสฺส อาราเม.\csromanspacer{} ตตฺร โข ภควา ภิกฺขู อามนฺเตสิ “ภิกฺขโว”ติ. “ภทนฺเต”ติ เต ภิกฺขู ภควโต ปจฺจสฺโสสุํ.\csromanspacer{} ภควา เอตทโวจ\csromandash{}อุทฺเทสวิภงฺคํ โว ภิกฺขเว เทเสสฺสามิ, ตํ สุณาถ สาธุกํ มนสิ กโรถ, ภาสิสฺสามีติ. “เอวํ ภนฺเต”ติ โข เต ภิกฺขู ภควโต ปจฺจสฺโสสุํ.\csromanspacer{} ภควา เอตทโวจ\csromandash{}}

\prose{“ตถา ตถา ภิกฺขเว ภิกฺขุ อุปปริกฺเขยฺย, ยถา ยถา\footnote{ยถา ยถาสฺส (สี, สฺยา, กํ, อิ)} อุปปริกฺขโต พหิทฺธา จสฺส วิญฺญาณํ อวิกฺขิตฺตํ อวิสฏํ อชฺฌตฺตํ อสณฺฐิตํ อนุปาทาย น ปริตสฺเสยฺย, พหิทฺธา ภิกฺขเว วิญฺญาเณ อวิกฺขิตฺเต อวิสเฏ สติ อชฺฌตฺตํ อสณฺฐิเต อนุปาทาย อปริตสฺสโต อายติํ ชาติชรามรณทุกฺขสมุทยสมฺภโว น โหตี”ติ.\csromanspacer{} อิทมโวจ ภควา, อิทํ วตฺวาน สุคโต อุฏฺฐายาสนา วิหารํ ปาวิสิ.}

\proseitem{314}{อถ โข เตสํ ภิกฺขูนํ อจิรปกฺกนฺตสฺส ภควโต เอตทโหสิ “อิทํ โข โน อาวุโส ภควา สํขิตฺเตน อุทฺเทสํ อุทฺทิสิตฺวา วิตฺถาเรน อตฺถํ อวิภชิตฺวา อุฏฺฐายาสนา วิหารํ ปวิฏฺโฐ ‘ตถา ตถา \csromanfolio{266}ภิกฺขเว ภิกฺขุ อุปปริกฺเขยฺย, ยถา ยถา อุปปริกฺขโต พหิทฺธา จสฺส วิญฺญาณํ อวิกฺขิตฺตํ อวิสฏํ อชฺฌตฺตํ อสณฺฐิตํ อนุปาทาย น ปริตสฺเสยฺย, พหิทฺธา ภิกฺขเว วิญฺญาเณ อวิกฺขิตฺเต อวิสเฏ สติ อชฺฌตฺตํ อสณฺฐิเต อนุปาทาย อปริตสฺสโต อายติํ ชาติชรามรณทุกฺขสมุทยสมฺภโว น โหตี’ติ, โก นุ โข อิมสฺส ภควตา สํขิตฺเตน อุทฺเทสสฺส อุทฺทิฏฺฐสฺส วิตฺถาเรน อตฺถํ อวิภตฺตสฺส วิตฺถาเรน อตฺถํ วิภเชยฺยา”ติ.\csromanspacer{} อถ โข เตสํ ภิกฺขูนํ เอตทโหสิ “อยํ โข อายสฺมา มหากจฺจาโน สตฺถุ เจว สํวณฺณิโต, สมฺภาวิโต จ วิญฺญูนํ สพฺรหฺมจารีนํ, ปโหติ จายสฺมา มหากจฺจาโน อิมสฺส ภควตา สํขิตฺเตน อุทฺเทสสฺส อุทฺทิฏฺฐสฺส วิตฺถาเรน อตฺถํ อวิภตฺตสฺส วิตฺถาเรน อตฺถํ วิภชิตุํ, ยํนูน มยํ เยนายสฺมา มหากจฺจาโน เตนุปสงฺกเมยฺยาม, อุปสงฺกมิตฺวา อายสฺมนฺตํ มหากจฺจานํ เอตมตฺถํ ปฏิปุจฺเฉยฺยามา”ติ.}

\prose{อถ โข เต ภิกฺขู เยนายสฺมา มหากจฺจาโน เตนุปสงฺกมิํสุ, อุปสงฺกมิตฺวา อายสฺมตา มหากจฺจาเนน สทฺธิํ สมฺโมทิํสุ, สมฺโมทนียํ กถํ สารณียํ วีติสาเรตฺวา เอกมนฺตํ นิสีทิํสุ, เอกมนฺตํ นิสินฺนา โข เต ภิกฺขู อายสฺมนฺตํ มหากจฺจานํ เอตทโวจุํ\csromandash{}}

\prose{อิทํ โข โน อาวุโส กจฺจาน ภควา สํขิตฺเตน อุทฺเทสํ อุทฺทิสิตฺวา วิตฺถาเรน อตฺถํ อวิภชิตฺวา อุฏฺฐายาสนา วิหารํ ปวิฏฺโฐ “ตถา ตถา ภิกฺขเว ภิกฺขุ อุปปริกฺเขยฺย, ยถา ยถา อุปปริกฺขโต พหิทฺธา จสฺส วิญฺญาณํ อวิกฺขิตฺตํ อวิสฏํ อชฺฌตฺตํ อสณฺฐิตํ อนุปาทาย น ปริตสฺเสยฺย, พหิทฺธา ภิกฺขเว วิญฺญาเณ อวิกฺขิตฺเต อวิสเฏ สติ อชฺฌตฺตํ อสณฺฐิเต อนุปาทาย อปริตสฺสโต อายติํ ชาติชรามรณทุกฺขสมุทยสมฺภโว น โหตี”ติ.\csromanspacer{} เตสํ โน อาวุโส กจฺจาน อมฺหากํ อจิรปกฺกนฺตสฺส ภควโต เอตทโหสิ “อิทํ โข โน อาวุโส ภควา สํขิตฺเตน อุทฺเทสํ อุทฺทิสิตฺวา วิตฺถาเรน อตฺถํ อวิภชิตฺวา อุฏฺฐายาสนา วิหารํ ปวิฏฺโฐ ‘ตถา ตถา ภิกฺขเว ภิกฺขุ อุปปริกฺเขยฺย, ยถา ยถา อุปปริกฺขโต พหิทฺธา จสฺส วิญฺญาณํ อวิกฺขิตฺตํ อวิสฏํ อชฺฌตฺตํ อสณฺฐิตํ อนุปาทาย น ปริตสฺเสยฺย, พหิทฺธา ภิกฺขเว วิญฺญาเณ อวิกฺขิตฺเต อวิสเฏ สติ อชฺฌตฺตํ อสณฺฐิเต อนุปาทาย อปริตสฺสโต อายติํ ชาติชรามรณทุกฺขสมุทยสมฺภโว น โหตี’ติ.\csromanspacer{} โก นุ โข อิมสฺส ภควตา สํขิตฺเตน อุทฺเทสสฺส อุทฺทิฏฺฐสฺส วิตฺถาเรน}

\csromanheader{2}{8. อุทฺเทสวิภงฺคสุตฺต (138)}
\prose{\csromanfolio{267}อตฺถํ อวิภตฺตสฺส วิตฺถาเรน อตฺถํ วิภเชยฺยา”ติ.\csromanspacer{} เตสํ โน อาวุโส กจฺจาน อมฺหากํ เอตทโหสิ “อยํ โข อายสฺมา มหากจฺจาโน สตฺถุ เจว สํวณฺณิโต, สมฺภาวิโต จ วิญฺญูนํ สพฺรหฺมจารีนํ, ปโหติ จายสฺมา มหากจฺจาโน อิมสฺส ภควตา สํขิตฺเตน อุทฺเทสสฺส อุทฺทิฏฺฐสฺส วิตฺถาเรน อตฺถํ อวิภตฺตสฺส วิตฺถาเรน อตฺถํ วิภชิตุํ, ยํนูน มยํ เยนายสฺมา มหากจฺจาโน เตนุปสงฺกเมยฺยาม, อุปสงฺกมิตฺวา อายสฺมนฺตํ มหากจฺจานํ เอตมตฺถํ ปฏิปุจฺเฉยฺยามา”ติ.\csromanspacer{} วิภชตายสฺมา มหากจฺจาโนติ.}

\proseitem{315}{เสยฺยถาปิ อาวุโส ปุริโส สารตฺถิโก สารคเวสี สารปริเยสนํ จรมาโน มหโต รุกฺขสฺส ติฏฺฐโต สารวโต อติกฺกมฺเมว มูลํ อติกฺกมฺม ขนฺธํ สาขาปลาเส สารํ ปริเยสิตพฺพํ มญฺเญยฺย.\csromanspacer{} เอวํ สมฺปทมิทํ, อายสฺมนฺตานํ สตฺถริ สมฺมุขีภูเต ตํ ภควนฺตํ อติสิตฺวา อเมฺห เอตมตฺถํ ปฏิปุจฺฉิตพฺพํ มญฺญถ.\csromanspacer{} โส หาวุโส ภควา ชานํ ชานาติ ปสฺสํ ปสฺสติ จกฺขุภูโต ญาณภูโต ธมฺมภูโต พฺรหฺมภูโต วตฺตา ปวตฺตา อตฺถสฺส นินฺเนตา อมตสฺส ทาตา ธมฺมสฺสามี ตถาคโต, โส เจว ปเนตสฺส กาโล อโหสิ, ยํ ภควนฺตํเยว เอตมตฺถํ ปฏิปุจฺเฉยฺยาถ, ยถา โว ภควา พฺยากเรยฺย, ตถา นํ ธาเรยฺยาถาติ.\csromanspacer{} อทฺธาวุโส กจฺจาน ภควา ชานํ ชานาติ ปสฺสํ ปสฺสติ จกฺขุภูโต ญาณภูโต ธมฺมภูโต พฺรหฺมภูโต วตฺตา ปวตฺตา อตฺถสฺส นินฺเนตา อมตสฺส ทาตา ธมฺมสฺสามี ตถาคโต, โส เจว ปเนตสฺส กาโล อโหสิ, ยํ ภควนฺตํเยว เอตมตฺถํ ปฏิปุจฺเฉยฺยาม, ยถา โน ภควา พฺยากเรยฺย, ตถา นํ ธาเรยฺยาม.\csromanspacer{} อปิ จายสฺมา มหากจฺจาโน สตฺถุ เจว สํวณฺณิโต, สมฺภาวิโต จ วิญฺญูนํ สพฺรหฺมจารีนํ, ปโหติ จายสฺมา มหากจฺจาโน อิมสฺส ภควโต สํขิตฺเตน อุทฺเทสสฺส อุทฺทิฏฺฐสฺส วิตฺถาเรน อตฺถํ อวิภตฺตสฺส วิตฺถาเรน อตฺถํ วิภชิตุํ, วิภชตายสฺมา มหากจฺจาโน อครุํ กริตฺวาติ.\csromanspacer{} เตนหาวุโส สุณาถ สาธุกํ มนสิ กโรถ, ภาสิสฺสามีติ. “เอวมาวุโส”ติ โข เต ภิกฺขู อายสฺมโต มหากจฺจานสฺส ปจฺจสฺโสสุํ.\csromanspacer{} อายสฺมา มหากจฺจาโน เอตทโวจ\csromandash{}}

\prose{ยํ โข โน อาวุโส ภควา สํขิตฺเตน อุทฺเทสํ อุทฺทิสิตฺวา วิตฺถาเรน อตฺถํ อวิภชิตฺวา อุฏฺฐายาสนา วิหารํ ปวิฏฺโฐ “ตถา ตถา ภิกฺขเว \csromanfolio{268}ภิกฺขุ อุปปริกฺเขยฺย, ยถา ยถา อุปปริกฺขโต พหิทฺธา จสฺส วิญฺญาณํ อวิกฺขิตฺตํ อวิสฏํ อชฺฌตฺตํ อสณฺฐิตํ อนุปาทาย น ปริตสฺเสยฺย, พหิทฺธา ภิกฺขเว วิญฺญาเณ อวิกฺขิตฺเต อวิสเฏ สติ อชฺฌตฺตํ อสณฺฐิเต อนุปาทาย อปริตสฺสโต อายติํ ชาติชรามรณทุกฺขสมุทยสมฺภโว น โหตี”ติ.\csromanspacer{} อิมสฺส โข อหํ อาวุโส ภควตา สํขิตฺเตน อุทฺเทสสฺส อุทฺทิฏฺฐสฺส วิตฺถาเรน อตฺถํ อวิภตฺตสฺส เอวํ วิตฺถาเรน อตฺถํ อาชานามิ.}

\proseitem{316}{กถญฺจาวุโส พหิทฺธา วิญฺญาณํ วิกฺขิตฺตํ วิสฏนฺติ วุจฺจติ.\csromanspacer{} อิธาวุโส ภิกฺขุโน จกฺขุนา รูปํ ทิสฺวา รูปนิมิตฺตานุสาริ วิญฺญาณํ โหติ รูปนิมิตฺตสฺสาทคธิตํ\footnote{...คถิตํ (สี, อิ)} รูปนิมิตฺตสฺสาทวินิพนฺธํ\footnote{...วินิพทฺธํ (สี, อิ)} รูปนิมิตฺตสฺสาทสํโยชนสํยุตฺตํ พหิทฺธา วิญฺญาณํ วิกฺขิตฺตํ วิสฏนฺติ วุจฺจติ.\csromanspacer{} โสเตน สทฺทํ สุตฺวา ฯเปฯ.\csromanspacer{} ฆาเนน คนฺธํ ฆายิตฺวา.\csromanspacer{} ชิวฺหาย รสํ สายิตฺวา.\csromanspacer{} กาเยน โผฏฺฐพฺพํ ผุสิตฺวา.\csromanspacer{} มนสา ธมฺมํ วิญฺญาย ธมฺมนิมิตฺตานุสาริ วิญฺญาณํ โหติ ธมฺมนิมิตฺตสฺสาทคธิตํ ธมฺมนิมิตฺตสฺสาทวินิพนฺธํ ธมฺมนิมิตฺตสฺสาทสํโยชนสํยุตฺตํ พหิทฺธา วิญฺญาณํ วิกฺขิตฺตํ วิสฏนฺติ วุจฺจติ.\csromanspacer{} เอวํ โข อาวุโส พหิทฺธา วิญฺญาณํ วิกฺขิตฺตํ วิสฏนฺติ วุจฺจติ.}

\proseitem{317}{กถญฺจาวุโส พหิทฺธา วิญฺญาณํ อวิกฺขิตฺตํ อวิสฏนฺติ วุจฺจติ.\csromanspacer{} อิธาวุโส ภิกฺขุโน จกฺขุนา รูปํ ทิสฺวา น รูปนิมิตฺตานุสาริ วิญฺญาณํ โหติ, น รูปนิมิตฺตสฺสาทคธิตํ น รูปนิมิตฺตสฺสาทวินิพนฺธํ น รูปนิมิตฺตสฺสาทสํโยชนสํยุตฺตํ พหิทฺธา วิญฺญาณํ อวิกฺขิตฺตํ อวิสฏนฺติ วุจฺจติ.\csromanspacer{} โสเตน สทฺทํ สุตฺวา ฯเปฯ.\csromanspacer{} ฆาเนน คนฺธํ ฆายิตฺวา.\csromanspacer{} ชิวฺหาย รสํ สายิตฺวา.\csromanspacer{} กาเยน โผฏฺฐพฺพํ ผุสิตฺวา.\csromanspacer{} มนสา ธมฺมํ วิญฺญาย น ธมฺมนิมิตฺตานุสาริ วิญฺญาณํ โหติ, น ธมฺมนิมิตฺตสฺสาทคธิตํ น ธมฺมนิมิตฺตสฺสาทวินิพนฺธํ น ธมฺมนิมิตฺตสฺสาทสํโยชนสํยุตฺตํ พหิทฺธา วิญฺญาณํ อวิกฺขิตฺตํ อวิสฏนฺติ วุจฺจติ.\csromanspacer{} เอวํ โข อาวุโส พหิทฺธา วิญฺญาณํ อวิกฺขิตฺตํ อวิสฏนฺติ วุจฺจติ.}

\proseitem{318}{กถญฺจาวุโส อชฺฌตฺตํ\footnote{อชฺฌตฺตํ จิตฺตํ (สี, สฺยา, กํ, อิ)} สณฺฐิตนฺติ วุจฺจติ.\csromanspacer{} อิธาวุโส ภิกฺขุ วิวิจฺเจว กาเมหิ วิวิจฺจ อกุสเลหิ ธมฺเมหิ สวิตกฺกํ สวิจารํ วิเวกชํ ปีติสุขํ ปฐมํ ฌานํ อุปสมฺปชฺช วิหรติ, ตสฺส วิเวกชปีติสุขานุสาริ วิญฺญาณํ โหติ, วิเวกชปีติสุขสฺสาทคธิตํ}

\csromanheader{2}{8. อุทฺเทสวิภงฺคสุตฺต (138)}
\prose{\csromanfolio{269}วิเวกชปีติสุขสฺสาทวินิพนฺธํ วิเวกชปีติสุขสฺสาทสํโยชนสํยุตฺตํ อชฺฌตฺตํ สณฺฐิตนฺติ วุจฺจติ.}

\prose{ปุน จปรํ อาวุโส ภิกฺขุ วิตกฺกวิจารานํ วูปสมา อชฺฌตฺตํ สมฺปสาทนํ เจตโส เอโกทิภาวํ อวิตกฺกํ อวิจารํ สมาธิชํ ปีติสุขํ ทุติยํ ฌานํ อุปสมฺปชฺช วิหรติ, ตสฺส สมาธิชปีติสุขานุสาริ วิญฺญาณํ โหติ, สมาธิชปีติสุขสฺสาทคธิตํ สมาธิชปีติสุขสฺสาทวินิพนฺธํ สมาธิชปีติสุขสฺสาทสํโยชนสํยุตฺตํ อชฺฌตฺตํ จิตฺตํ สณฺฐิตนฺติ วุจฺจติ.}

\prose{ปุน จปรํ อาวุโส ภิกฺขุ ปีติยา จ วิราคา อุเปกฺขโก จ วิหรติ สโต จ สมฺปชาโน, สุขญฺจ กาเยน ปฏิสํเวเทติ, ยํ ตํ อริยา อาจิกฺขนฺติ “อุเปกฺขโก สติมา สุขวิหารี”ติ.\csromanspacer{} ตติยํ ฌานํ อุปสมฺปชฺช วิหรติ, ตสฺส อุเปกฺขานุสาริ วิญฺญาณํ โหติ อุเปกฺขาสุขสฺสาทคธิตํ อุเปกฺขาสุขสฺสาทวินิพนฺธํ อุเปกฺขาสุขสฺสาทสํโยชนสํยุตฺตํ อชฺฌตฺตํ จิตฺตํ สณฺฐิตนฺติ วุจฺจติ.}

\prose{ปุน จปรํ อาวุโส ภิกฺขุ สุขสฺส จ ปหานา ทุกฺขสฺส จ ปหานา ปุพฺเพว โสมนสฺสโทมนสฺสานํ อตฺถงฺคมา อทุกฺขมสุขํ อุเปกฺขาสติปาริสุทฺธิํ จตุตฺถํ ฌานํ อุปสมฺปชฺช วิหรติ, ตสฺส อทุกฺขมสุขานุสาริ วิญฺญาณํ โหติ อทุกฺขมสุขสฺสาทคธิตํ อทุกฺขมสุขสฺสาทวินิพนฺธํ อทุกฺขมสุขสฺสาทสํโยชนสํยุตฺตํ อชฺฌตฺตํ จิตฺตํ สณฺฐิตนฺติ วุจฺจติ.\csromanspacer{} เอวํ โข อาวุโส อชฺฌตฺตํ\footnote{อชฺฌตฺตํ จิตฺตํ (สี, สฺยา, กํ, อิ)} สณฺฐิตนฺติ วุจฺจติ.}

\proseitem{319}{กถญฺจาวุโส อชฺฌตฺตํ1 อสณฺฐิตนฺติ วุจฺจติ.\csromanspacer{} อิธาวุโส ภิกฺขุ วิวิจฺเจว กาเมหิ วิวิจฺจ อกุสเลหิ ธมฺเมหิ ฯเปฯ ปฐมํ ฌานํ อุปสมฺปชฺช วิหรติ, ตสฺส น วิเวกชปีติสุขานุสาริ วิญฺญาณํ โหติ, น วิเวกชปีติสุขสฺสาทคธิตํ น วิเวกชปีติสุขสฺสาทวินิพนฺธํ น วิเวกชปีติสุขสฺสาทสํโยชนสํยุตฺตํ อชฺฌตฺตํ จิตฺตํ อสณฺฐิตนฺติ วุจฺจติ.}

\prose{ปุน จปรํ อาวุโส ภิกฺขุ วิตกฺกวิจารานํ วูปสมา ฯเปฯ ทุติยํ ฌานํ อุปสมฺปชฺช วิหรติ, ตสฺส น สมาธิชปีติสุขานุสาริ วิญฺญาณํ โหติ น สมาธิชปีติสุขสฺสาทคธิตํ น สมาธิชปีติสุขสฺสาทวินิพนฺธํ น สมาธิชปีติสุขสฺสาทสํโยชนสํยุตฺตํ อชฺฌตฺตํ จิตฺตํ อสณฺฐิตนฺติ วุจฺจติ.}

\prose{\csromanfolio{270}ปุน จปรํ อาวุโส ภิกฺขุ ปีติยา จ วิราคา ฯเปฯ ตติยํ ฌานํ อุปสมฺปชฺช วิหรติ, ตสฺส น อุเปกฺขานุสาริ วิญฺญาณํ โหติ น อุเปกฺขาสุขสฺสาทคธิตํ น อุเปกฺขาสุขสฺสาทวินิพนฺธํ น อุเปกฺขาสุขสฺสาทสํโยชนสํยุตฺตํ อชฺฌตฺตํ จิตฺตํ อสณฺฐิตนฺติ วุจฺจติ.}

\prose{ปุน จปรํ อาวุโส ภิกฺขุ สุขสฺส จ ปหานา ทุกฺขสฺส จ ปหานา ปุพฺเพว โสมนสฺสโทมนสฺสานํ อตฺถงฺคมา อทุกฺขมสุขํ อุเปกฺขาสติปาริสุทฺธิํ จตุตฺถํ ฌานํ อุปสมฺปชฺช วิหรติ, ตสฺส น อทุกฺขมสุขานุสาริ วิญฺญาณํ โหติ น อทุกฺขมสุขสฺสาทคธิตํ น อทุกฺขมสุขสฺสาทวินิพนฺธํ น อทุกฺขมสุขสฺสาทสํโยชน- สํยุตฺตํ อชฺฌตฺตํ จิตฺตํ อสณฺฐิตนฺติ วุจฺจติ.\csromanspacer{} เอวํ โข อาวุโส อชฺฌตฺตํ\footnote{อชฺฌตฺตํ จิตฺตํ (สี, สฺยา, กํ, อิ)} อสณฺฐิตนฺติ วุจฺจติ.}

\proseitem{320}{กถญฺจาวุโส อนุปาทา ปริตสฺสนา โหติ.\csromanspacer{} อิธาวุโส อสฺสุตวา ปุถุชฺชโน อริยานํ อทสฺสาวี อริยธมฺมสฺส อโกวิโท อริยธมฺเม อวินีโต สปฺปุริสานํ อทสฺสาวี สปฺปุริสธมฺมสฺส อโกวิโท สปฺปุริสธมฺเม อวินีโต รูปํ อตฺตโต สมนุปสฺสติ, รูปวนฺตํ วา อตฺตานํ, อตฺตนิ วา รูปํ, รูปสฺมิํ วา อตฺตานํ.\csromanspacer{} ตสฺส ตํ รูปํ วิปริณมติ อญฺญถา โหติ, ตสฺส รูปวิปริณามญฺญถาภาวา รูปวิปริณามานุปริวตฺติ วิญฺญาณํ โหติ, ตสฺส รูปวิหริณามานุปริวตฺตชา ปริตสฺสนา ธมฺมสมุปฺปาทา จิตฺตํ ปริยาทาย ติฏฺฐนฺติ, เจตโส ปริยาทานา อุตฺตาสวา จ โหติ วิฆาตวา จ อเปกฺขวา จ อนุปาทาย จ ปริตสฺสติ.\csromanspacer{} เวทนํ ฯเปฯ.\csromanspacer{} สญฺญํ.\csromanspacer{} สงฺขาเร.\csromanspacer{} วิญฺญาณํ อตฺตโต สมนุปสฺสติ, วิญฺญาณวนฺตํ วา อตฺตานํ, อตฺตนิ วา วิญฺญาณํ, วิญฺญาณสฺมิํ วา อตฺตานํ.\csromanspacer{} ตสฺส ตํ วิญฺญาณํ วิปริณมติ อญฺญถา โหติ, ตสฺส วิญฺญาณวิปริณามญฺญถาภาวา วิญฺญาณวิปริณามานุปริวตฺติ วิญฺญาณํ โหติ.\csromanspacer{} ตสฺส วิญฺญาณวิปริณามานุปริวตฺตชา ปริตสฺสนา ธมฺมสมุปฺปาทา จิตฺตํ ปริยาทาย ติฏฺฐนฺติ, เจตโส ปริยาทานา อุตฺตาสวา จ โหติ วิฆาตวา จ อเปกฺขวา จ อนุปาทาย จ ปริตสฺสติ.\csromanspacer{} เอวํ โข อาวุโส อนุปาทา ปริตสฺสนา โหติ.}

\proseitem{321}{กถญฺจาวุโส อนุปาทานา อปริตสฺสนา โหติ.\csromanspacer{} อิธาวุโส สุตวา อริยสาวโก อริยานํ ทสฺสาวี อริยธมฺมสฺส โกวิโท}

\csromanheader{2}{8. อุทฺเทสวิภงฺคสุตฺต (138)}
\prose{\csromanfolio{271}อริยธมฺเม สุวินีโต สปฺปุริสานํ ทสฺสาวี สปฺปุริสธมฺมสฺส โกวิโท สปฺปุริสธมฺเม สุวินีโต น รูปํ อตฺตโต สมนุปสฺสติ, น รูปวนฺตํ วา อตฺตานํ, น อตฺตนิ วา รูปํ, น รูปสฺมิํ วา อตฺตานํ.\csromanspacer{} ตสฺส ตํ รูปํ วิปริณมติ อญฺญถา โหติ, ตสฺส รูปวิปริณามญฺญถาภาวา น จ รูปวิปริณามานุปริวตฺติ วิญฺญาณํ โหติ, ตสฺส น รูปวิปริณามานุปริวตฺตชา ปริตสฺสนา ธมฺมสมุปฺปาทา จิตฺตํ ปริยาทาย ติฏฺฐนฺติ, เจตโส ปริยาทานา น เจวุตฺตาสวา\footnote{น จ อุตฺตาสวา (สี)} โหติ น จ วิฆาตวา น จ อเปกฺขวา อนุปาทาย จ น ปริตสฺสติ.\csromanspacer{} น เวทนํ.\csromanspacer{} น สญฺญํ.\csromanspacer{} น สงฺขาเร.\csromanspacer{} น วิญฺญาณํ อตฺตโต สมนุปสฺสติ, น วิญฺญาณวนฺตํ วา อตฺตานํ น อตฺตนิ วา วิญฺญาณํ, น วิญฺญาณสฺมิํ วา อตฺตานํ.\csromanspacer{} ตสฺส ตํ วิญฺญาณํ วิปริณมติ อญฺญถา โหติ, ตสฺส วิญฺญาณวิปริณามญฺญถาภาวา น จ วิญฺญาณวิปริณามานุปริวตฺติ วิญฺญาณํ โหติ, ตสฺส น วิญฺญาณวิปริณามานุปริวตฺตชา ปริตสฺสนา ธมฺมสมุปฺปาทา จิตฺตํ ปริยาทาย ติฏฺฐนฺติ, เจตโส ปริยาทานา น เจวุตฺตาสวา โหติ น วิฆาตวา น จ อเปกฺขวา อนุปาทาย จ น ปริตสฺสติ.\csromanspacer{} เอวํ โข อาวุโส อนุปาทา อปริตสฺสนา โหติ.}

\prose{ยํ โข โน อาวุโส ภควา สํขิตฺเตน อุทฺเทสํ อุทฺทิสิตฺวา วิตฺถาเรน อตฺถํ อวิภชิตฺวา อุฏฺฐายาสนา วิหารํ ปวิฏฺโฐ “ตถา ตถา ภิกฺขเว ภิกฺขุ อุปปริกฺเขยฺย, ยถา ยถา อุปปริกฺขโต พหิทฺธา จสฺส วิญฺญาณํ อวิกฺขิตฺตํ อวิสฏํ อชฺฌตฺตํ อสณฺฐิตํ อนุปาทาย น ปริตสฺเสยฺย, พหิทฺธา ภิกฺขเว วิญฺญาเณ อวิกฺขิตฺเต อวิสเฏ สติ อชฺฌตฺตํ อสณฺฐิเต อนุปาทาย อปริตสฺสโต อายติํ ชาติชรามรณทุกฺขสมุทยสมฺภโว น โหตี”ติ.\csromanspacer{} อิมสฺส โข อหํ อาวุโส ภควตา สํขิตฺเตน อุทฺเทสสฺส อุทฺทิฏฺฐสฺส วิตฺถาเรน อตฺถํ อวิภตฺตสฺส เอวํ วิตฺถาเรน อตฺถํ อาชานามิ, อากงฺขมานา จ ปน ตุเมฺห อายสฺมนฺโต ภควนฺตํเยว อุปสงฺกมิตฺวา เอตมตฺถํ ปฏิปุจฺเฉยฺยาถ.\csromanspacer{} ยถา โว ภควา พฺยากโรติ, ตถา นํ ธาเรยฺยาถาติ.}

\proseitem{322}{อถ โข เต ภิกฺขู อายสฺมโต มหากจฺจานสฺส ภาสิตํ อภินนฺทิตฺวา อนุโมทิตฺวา อุฏฺฐายาสนา เยน ภควา เตนุปสงฺกมิํสุ, อุปสงฺกมิตฺวา ภควนฺตํ อภิวาเทตฺวา เอกมนฺตํ นิสีทิํสุ, เอกมนฺตํ นิสินฺนา โข เต ภิกฺขู ภควนฺตํ เอตทโวจุํ\csromandash{}}

\prose{\csromanfolio{272}ยํ โข โน ภนฺเต ภควา สํขิตฺเตน อุทฺเทสํ อุทฺทิสิตฺวา วิตฺถาเรน อตฺถํ อวิภชิตฺวา อุฏฺฐายาสนา วิหารํ ปวิฏฺโฐ “ตถา ตถา ภิกฺขเว ภิกฺขุ อุปปริกฺเขยฺย, ยถา ยถา อุปปริกฺขโต พหิทฺธา จสฺส วิญฺญาณํ อวิกฺขิตฺตํ อวิสฏํ อชฺฌตฺตํ อสณฺฐิตํ อนุปาทาย น ปริตสฺเสยฺย, พหิทฺธา ภิกฺขเว วิญฺญาเณ อวิกฺขิตฺเต อวิสเฏ สติ อชฺฌตฺตํ อสณฺฐิเต อนุปาทาย อปริตสฺสโต อายติํ ชาติชรามรณทุกฺขสมุทยสมฺภโว น โหตี”ติ.}

\prose{เตสํ โน ภนฺเต อมฺหากํ อจิรปกฺกนฺตสฺส ภควโต เอตทโหสิ “อิทํ โข โน อาวุโส ภควา สํขิตฺเตน อุทฺเทสํ อุทฺทิสิตฺวา วิตฺถาเรน อตฺถํ อวิภชิตฺวา อุฏฺฐายาสนา วิหารํ ปวิฏฺโฐ ‘ตถา ตถา ภิกฺขเว ภิกฺขุ อุปปริกฺเขยฺย, ยถา ยถา อุปปริกฺขโต พหิทฺธา จสฺส วิญฺญาณํ อวิกฺขิตฺตํ อวิสฏํ อชฺฌตฺตํ อสณฺฐิตํ อนุปาทาย น ปริตสฺเสยฺย, พหิทฺธา ภิกฺขเว วิญฺญาเณ อวิกฺขิตฺเต อวิสเฏ สติ อชฺฌตฺตํ อสณฺฐิเต อนุปาทาย อปริตสฺสโต อายติํ ชาติชรามรณทุกฺขสมุทยสมฺภโว น โหตี’ติ, โก นุ โข อิมสฺส ภควตา สํขิตฺเตน อุทฺเทสสฺส อุทฺทิฏฺฐสฺส วิตฺถาเรน อตฺถํ อวิภตฺตสฺส วิตฺถาเรน อตฺถํ วิภเชยฺยา”ติ.\csromanspacer{} เตสํ โน ภนฺเต อมฺหากํ เอตทโหสิ “อยํ โข อายสฺมา มหากจฺจาโน สตฺถุ เจว สํวณฺณิโต, สมฺภาวิโต จ วิญฺญูนํ สพฺรหฺมจารีนํ, ปโหติ จายสฺมา มหากจฺจาโน อิมสฺส ภควตา สํขิตฺเตน อุทฺเทสสฺส อุทฺทิฏฺฐสฺส วิตฺถาเรน อตฺถํ อวิภตฺตสฺส วิตฺถาเรน อตฺถํ วิภชิตุํ, ยํนูน มยํ เยนายสฺมา มหากจฺจาโน เตนุปสงฺกเมยฺยาม, อุปสงฺกมิตฺวา อายสฺมนฺตํ มหากจฺจานํ เอตมตฺถํ ปฏิปุจฺเฉยฺยามา”ติ.}

\prose{อถ โข มยํ ภนฺเต เยนายสฺมา มหากจฺจาโน เตนุปสงฺกมิมฺห, อุปสงฺกมิตฺวา อายสฺมนฺตํ มหากจฺจานํ เอตมตฺถํ ปฏิปุจฺฉิมฺห.\csromanspacer{} เตสํ โน ภนฺเต อายสฺมตา มหากจฺจาเนน อิเมหิ อากาเรหิ อิเมหิ ปเทหิ อิเมหิ พฺยญฺชเนหิ อตฺโถ วิภตฺโตติ.}

\prose{ปณฺฑิโต ภิกฺขเว มหากจฺจาโน, มหาปญฺโญ ภิกฺขเว มหากจฺจาโน.\csromanspacer{} มํ เจปิ ตุเมฺห ภิกฺขเว เอตมตฺถํ ปฏิปุจฺเฉยฺยาถ, อหมฺปิ เอวเมวํ}

\csromanheader{2}{9. อรณวิภงฺคสุตฺต (139)}
\prose{\csromanfolio{273}พฺยากเรยฺยํ, ยถา ตํ มหากจฺจาเนน พฺยากตํ, เอโส เจเวตสฺส\footnote{เอโส เจตสฺส (สี, อิ), เอโส เจว ตสฺส (สฺยา, กํ), เอโสเยว ตสฺส (ก)} อตฺโถ.\csromanspacer{} เอวญฺจ นํ ธาเรยฺยาถาติ.}

\prose{อิทมโวจ ภควา.\csromanspacer{} อตฺตมนา เต ภิกฺขู ภควโต ภาสิตํ อภินนฺทุนฺติ.}

\nitthitamruled{อุทฺเทสวิภงฺคสุตฺตํ นิฏฺฐิตํ อฏฺฐมํ.}
\csromanmatikaanchor{mtk.46}
\csromantocmarkref{subsection}{9. อรณวิภงฺคสุตฺต}{mtk.46}
\bookmark[dest={mtk.46},level=2]{9. อรณวิภงฺคสุตฺต}
\csromanheader{2}{9. อรณวิภงฺคสุตฺต}
\proseitem{323}{เอวํ เม สุตํ\csromandash{}เอกํ สมยํ ภควา สาวตฺถิยํ วิหรติ เชตวเน อนาถปิณฺฑิกสฺส อาราเม.\csromanspacer{} ตตฺร โข ภควา ภิกฺขู อามนฺเตสิ “ภิกฺขโว”ติ. “ภทนฺเต”ติ เต ภิกฺขู ภควโต ปจฺจสฺโสสุํ.\csromanspacer{} ภควา เอตทโวจ “อรณวิภงฺคํ โว ภิกฺขเว เทเสสฺสามิ, ตํ สุณาถ สาธุกํ มนสิ กโรถ, ภาสิสฺสามี”ติ. “เอวํ ภนฺเต”ติ โข เต ภิกฺขู ภควโต ปจฺจสฺโสสุํ.\csromanspacer{} ภควา เอตทโวจ\csromandash{}}

\prose{น กามสุขมนุยุญฺเชยฺย หีนํ คมฺมํ โปถุชฺชนิกํ อนริยํ อนตฺถสํหิตํ, น จ อตฺถกิลมถานุโยคมนุยุญฺเชยฺย ทุกฺขํ อนริยํ อนตฺถสํหิตํ.\csromanspacer{} เอเต โข ภิกฺขเว\footnote{เอเต โข (สี), เอเต เต (สฺยา, กํ, อิ)} อุโภ อนฺเต อนุปคมฺม มชฺฌิมา ปฏิปทา ตถาคเตน อภิสมฺพุทฺธา จกฺขุกรณี ญาณกรณี อุปสมาย อภิญฺญาย สมฺโพธาย นิพฺพานาย สํวตฺตติ.\csromanspacer{} อุสฺสาทนญฺจ ชญฺญา, อปสาทนญฺจ ชญฺญา, อุสฺสาทนญฺจ ญตฺวา อปสาทนญฺจ ญตฺวา เนวุสฺสาเทยฺย น อปสาเทยฺย\footnote{นาปสาเทยฺย (สี)} ธมฺมเมว เทเสยฺย.\csromanspacer{} สุขวินิจฺฉยํ ชญฺญา, สุขวินิจฺฉยํ ญตฺวา อชฺฌตฺตํ สุขมนุยุญฺเชยฺย, รโหวาทํ น ภาเสยฺย, สมฺมุขา น ขีณํ\footnote{นาติขีณํ (สฺยา, กํ, ก)} ภเณ, อตรมาโนว ภาเสยฺย โน ตรมาโน, ชนปทนิรุตฺติํ นาภินิเวเสยฺย, สมญฺญํ นาติธาเวยฺยาติ.\csromanspacer{} อยมุทฺเทโส อรณวิภงฺคสฺส.}

\proseitem{324}{“น กามสุขมนุยุญฺเชยฺย หีนํ คมฺมํ โปถุชฺชนิกํ อนริยํ อนตฺถสํหิตํ, น จ อตฺตกิลมถานุโยคมนุยุญฺเชยฺย ทุกฺขํ อนริยํ \csromanfolio{274}อนตฺถสํหิตนฺ”ติ อิติ โข ปเนตํ วุตฺตํ, กิญฺเจตํ ปฏิจฺจ วุตฺตํ.\csromanspacer{} โย กามปฏิสนฺธิสุขิโน โสมนสฺสานุโยโค หีโน คมฺโม โปถุชฺชนิโก อนริโย อนตฺถสํหิโต, สทุกฺขา เอโส ธมฺโม ส-อุปฆาโต ส-อุปายาโส สปริฬาโห มิจฺฉาปฏิปทา.\csromanspacer{} โย กามปฏิสนฺธิสุขิโน โสมนสฺสานุโยคํ อนนุโยโค หีนํ คมฺมํ โปถุชฺชนิกํ อนริยํ อนตฺถสํหิตํ, อทุกฺโข เอโส ธมฺโม อนุปฆาโต อนุปายาโส อปริฬาโห สมฺมาปฏิปทา.\csromanspacer{} โย อตฺตกิลมถานุโยโว ทุกฺโข อนริโย อนตฺถสํหิโต, สทุกฺโข เอโส ธมฺโม ส-อุปฆาโต ส-อุปายาโส สปริฬาโห มิจฺฉาปฏิปทา.\csromanspacer{} โย อตฺตกิลมถานุโยคํ อนนุโยโค ทุกฺขํ อนริยํ อนตฺถสํหิตํ, อทุกฺโข เอโส ธมฺโม อนุปฆาโต อนุปายาโส อปริฬาโห สมฺมาปฏิปทา. “น กามสุขนุยุญฺเชยฺย หีนํ คมฺมํ, โปถุชฺชนิกํ อนริยํ อนตฺถสํหิตํ, น จ อตฺถกิลมถานุโยคํ อนุยุญฺเชยฺย ทุกฺขํ อนริยํ อนตฺถสํหิตนฺ”ติ อิติ ยํ ตํ วุตฺตํ, อิทเมตํ ปฏิจฺจ วุตฺตํ.}

\proseitem{325}{“เอเต โข อุโภ อนฺเต อนุปคมฺม มชฺฌิมา ปฏิปทา ตถาคเตน อภิสมฺพุทฺธา จกฺขุกรณี ญาณกรณี อุปสมาย อภิญฺญาย สมฺโพธาย นิพฺพานาย สํวตฺตตี”ติ อิติ โข ปเนตํ วุตฺตํ, กิญฺเจตํ ปฏิจฺจ วุตฺตํ.\csromanspacer{} อยเมว อริโย อฏฺฐงฺคิโก มคฺโค.\csromanspacer{} เสยฺยถิทํ, สมฺมาทิฏฺฐิ สมฺมาสงฺกปฺโป สมฺมาวาจา สมฺมากมฺมนฺโต สมฺมา-อาชีโว สมฺมาวายาโม สมฺมาสติ สมฺมาสมาธิ. “เอเต โข อุโภ อนฺเต อนุปคมฺม มชฺฌิมา ปฏิปทา ตถาคเตน อภิสมฺพุทฺธา จกฺขุกรณี ญาณกรณี อุปสมาย อภิญฺญาย สมฺโพธาย นิพฺพานาย สํวตฺตตี”ติ อิติ ยํ ตํ วุตฺตํ, อิทเมตํ ปฏิจฺจ วุตฺตํ.}

\proseitem{326}{“อุสฺสาทนญฺจ ชญฺญา, อปสาทนญฺจ ชญฺญา, อุสฺสาทนญฺจ ญตฺวา อปสาทนญฺจ ญตฺวา เนวุสฺสาเทยฺย น อปสาเทยฺย ธมฺมเมว เทเสยฺยา”ติ อิติ โข ปเนตํ วุตฺตํ, กิญฺเจตํ ปฏิจฺจ วุตฺตํ.\csromanspacer{} กถญฺจ ภิกฺขเว อุสฺสาทนา จ โหติ อปสาทนา จ, โน จ ธมฺมเทสนา.\csromanspacer{} เย กามปฏิสนฺธิสุขิโน โสมนสฺสานุโยคํ อนุยุตฺตา หีนํ คมฺมํ โปถุชฺชนิกํ อนริยํ อนตฺถสํหิตํ, สพฺเพ เต สทุกฺขา ส-อุปฆาตา ส-อุปายาสา สปริฬาหา มิจฺฉาปฏิปนฺนาติ อิติ วทํ\footnote{อิติ ปรํ (ก)} อิตฺเถเก อปสาเทติ.}

\csromanheader{2}{9. อรณวิภงฺคสุตฺต (139)}
\prose{\csromanfolio{275}เย กามปฏิสนฺธิสุขิโน โสมนสฺสานุโยคํ อนนุยุตฺตา หีนํ คมฺมํ โปถุชฺชนิกํ อนริยํ อนตฺถสํหิตํ, สพฺเพ เต อทุกฺขา อนุปฆาตา อนุปายาสา อปริฬาหา สมฺมาปฏิปนฺนาติ อิติ วทํ อิตฺเถเก อุสฺสาเทติ.}

\prose{เย อตฺตกิลมถานุโยคํ อนุยุตฺตา ทุกฺขํ อนริยํ อนตฺถสํหิตํ, สพฺเพ เต สทุกฺขา ส-อุปฆาตา ส-อุปายาสา สปริฬาหา มิจฺฉาปฏิปนฺนาติ อิติ วทํ อิตฺเถเก อปสาเทติ.}

\prose{เย อตฺตกิลมถานุโยคํ อนนุยุตฺตา ทุกฺขํ อนริยํ อนตฺถสํหิตํ, สพฺเพ เต อทุกฺขา อนุปฆาตา อนุปายาสา อปริฬาหา สมฺมาปฏิปนฺนาติ อิติ วทํ อิตฺเถเก อุสฺสาเทติ.}

\prose{เยสํ เกสญฺจิ ภวสํโยชนํ อปฺปหีนํ, สพฺเพ เต สทุกฺขา ส- อุปฆาตา ส-อุปายาสา สปริฬาหา มิจฺฉาปฏิปนฺนาติ อิติ วทํ อิตฺเถเก อปสาเทติ.}

\prose{เยสํ เกสญฺจิ ภวสํโยชนํ ปหีนํ, สพฺเพ เต อทุกฺขา อนุปฆาตา อนุปายาสา อปริฬาหา สมฺมาปฏิปนฺนาติ อิติ วทํ อิตฺเถเก อุสฺสาเทติ.\csromanspacer{} เอวํ โข ภิกฺขเว อุสฺสาทนา จ โหติ อปสาทนา จ, โน จ ธมฺมเทสนา.}

\proseitem{327}{กถญฺจ ภิกฺขเว เนวุสฺสาทนา โหติ น อปสาทนา ธมฺมเทสนา จ\footnote{ธมฺมเทสนาว (สฺย, กํ)}.\csromanspacer{} เย กามปฏิสนฺธิสุขิโน โสมนสฺสานุโยคํ อนุยุตฺตา หีนํ คมฺมํ โปถุชฺชนิกํ อนริยํ อนตฺถสํหิตํ, สพฺเพ เต สทุกฺขา ส-อุปฆาตา ส-อุปายาสา สปริฬาหา มิจฺฉาปฏิปนฺนาติ น เอวมาห, “อนุโยโค จ โข สทุกฺโข เอโส ธมฺโม ส-อุปฆาโต ส-อุปายาโส สปริฬาโห มิจฺฉาปฏิปทา”ติ อิติ วทํ ธมฺมเมว เทเสติ.}

\prose{เย กามปฏิสนฺธิสุขิโน โสมนสฺสานุโยคํ อนนุยุตฺตา หีนํ คมฺมํ โปถุชฺชนิกํ อนริยํ อนตฺถสํหิตํ, สพฺเพ เต อทุกฺขา อนุปฆาตา อนุปายาสา อปริฬาหา สมฺมาปฏิปนฺนาติ น เอวมาห, “อนนุโยโค จ โข อทุกฺโข เอโส ธมฺโม อนุปฆาโต อนุปายาโส อปริฬาโห สมฺมาปฏิปทา”ติ อิติ วทํ ธมฺมเมว เทเสติ.}

\prose{\csromanfolio{276}เย อตฺตกิลมถานุโยคํ อนุยุตฺตา ทุกฺขํ อนริยํ อนตฺถสํหิตํ, สพฺเพ เต สทุกฺขา ส-อุปฆาตา ส-อุปายาสา สปริฬาหา มิจฺฉาปฏิปนฺนาติ น เอวมาห, “อนุโยโค จ โข สทุกฺโข เอโส ธมฺโม ส- อุปฆาโต ส-อุปายาโส สปริฬาโห มิจฺฉาปฏิปทา”ติ อิติ วทํ ธมฺมเมว เทเสติ.}

\prose{เย อตฺตกิลมถานุโยคํ อนนุยุตฺตา ทุกฺขํ อนริยํ อนตฺถสํหิตํ, สพฺเพ เต อทุกฺขา อนุปฆาตา อนุปายาสา อปริฬาหา สมฺมาปฏิปนฺนาติ น เอวมาห, “อนนุโยโค จ โข อทุกฺโข เอโส ธมฺโม อนุปฆาโต อนุปายาโส อปริฬาโห สมฺมาปฏิปทา”ติ อิติ วทํ ธมฺมเมว เทเสติ.}

\prose{เยสํ เกสญฺจิ ภวสํโยชนํ อปฺปหีนํ, สพฺเพ เต สทุกฺขา ส- อุปฆาตา ส-อุปายาสา สปริฬาหา มิจฺฉาปฏิปนฺนาติ น เอวมาห, “ภวสํโยชเน จ โข อปฺปหีเน ภโวปิ อปฺปหีโน โหตี”ติ อิติ วทํ ธมฺมเมว เทเสติ.}

\prose{เยสํ เกสญฺจิ ภวสํโยชนํ ปหีนํ, สพฺเพ เต อทุกฺขา อนุปฆาตา อนุปายาสา อปริฬาหา สมฺมาปฏิปนฺนาติ น เอวมาห, “ภวสํโยชเน จ โข ปหีเน ภโวปิ ปหีโน โหตี”ติ อิติ วทํ ธมฺมเมว เทเสติ.\csromanspacer{} เอวํ โข ภิกฺขเว เนวุสฺสาทนา โหติ น อปสาทนา ธมฺมเทสนา จ. “อุสฺสาทนญฺจ ชญฺญา, อปสาทนญฺจ ชญฺญา, อุสฺสาทนญฺจ ญตฺวา อปสาทนญฺจ ญตฺวา เนวุสฺสาเทยฺย น อปสาเทยฺย ธมฺมเมว เทเสยฺยา”ติ อิติ ยํ ตํ วุตฺตํ, อิทเมตํ ปฏิจฺจ วุตฺตํ.}

\proseitem{328}{“สุขวินิจฺฉยํ ชญฺญา, สุขวินิจฺฉยํ ญตฺวา อชฺฌตฺตํ สุขมนุยุญฺเชยฺยา”ติ อิติ โข ปเนตํ วุตฺตํ, กิญฺเจตํ ปฏิจฺจ วุตฺตํ.\csromanspacer{} ปญฺจิเม ภิกฺขเว กามคุณา.\csromanspacer{} กตเม ปญฺจ, จกฺขุวิญฺเญยฺยา รูปา อิฏฺฐา กนฺตา มนาปา ปิยรูปา กามูปสํหิตา รชนียา.\csromanspacer{} โสตวิญฺเญยฺยา สทฺทา.\csromanspacer{} ฆานวิญฺเญยฺยา คนฺธา.\csromanspacer{} ชิวฺหาวิญฺเญยฺยา รสา.\csromanspacer{} กายวิญฺเญยฺยา โผฏฺฐพฺพา อิฏฺฐา กนฺตา มนาปา ปิยรูปา กามูปสํหิตา รชนียา.\csromanspacer{} อิเม โข ภิกฺขเว ปญฺจ กามคุณา.\csromanspacer{} ยํ โข ภิกฺขเว อิเม ปญฺจ กามคุเณ ปฏิจฺจ อุปฺปชฺชติ สุขํ โสมนสฺสํ.\csromanspacer{} อิทํ วุจฺจติ กามสุขํ มีฬฺหสุขํ ปุถุชฺชนสุขํ อนริยสุขํ, “น อาเสวิตพฺพํ น ภาเวตพฺพํ น พหุลีกาตพฺพํ, ภายิตพฺพํ เอตสฺส สุขสฺสา”ติ วทามิ.}

\csromanheader{2}{9. อรณวิภงฺคสุตฺต (139)}
\prose{\csromanfolio{277}อิธ ภิกฺขเว ภิกฺขุ วิวิจฺเจว กาเมหิ วิวิจฺจ อกุสเลหิ ธมฺเมหิ สวิตกฺกํ สวิจารํ วิเวกชํ ปีติสุขํ ปฐมํ ฌานํ อุปสมฺปชฺช วิหรติ.\csromanspacer{} วิตกฺกวิจารานํ วูปสมา อชฺฌตฺตํ สมฺปสาทนํ เจตโส เอโกทิภาวํ อวิตกฺกํ อวิจารํ สมาธิชํ ปีติสุขํ ทุติยํ ฌานํ อุปสมฺปชฺช วิหรติ.\csromanspacer{} ปีติยา จ วิราคา อุเปกฺขโก จ วิหรติ ฯเปฯ ตติยํ ฌานํ ฯเปฯ จตุตฺถํ ฌานํ อุปสมฺปชฺช วิหรติ.\csromanspacer{} อิทํ วุจฺจติ เนกฺขมฺมสุขํ ปวิเวกสุขํ อุปสมสุขํ สมฺโพธิสุขํ, “อาเสวิตพฺพํ ภาเวตพฺพํ พหุลีกาตพฺพํ, น ภายิตพฺพํ เอตสฺส สุขสฺสา”ติ วทามิ. “สุขวินิจฺฉยํ ชญฺญา, สุขวินิจฺฉยํ ญตฺวา อชฺฌตฺตํ สุขมนุยุญฺเชยฺยา”ติ อิติ ยํ ตํ วุตฺตํ, อิทเมตํ ปฏิจฺจ วุตฺตํ.}

\proseitem{329}{“รโหวาทํ น ภาเสยฺย, สมฺมุขา น ขีณํ ภเณ”ติ อิติ โข ปเนตํ วุตฺตํ, กิญฺเจตํ ปฏิจฺจ วุตฺตํ.\csromanspacer{} ตตฺร ภิกฺขเว ยํ ชญฺญา รโหวาทํ อภูตํ อตจฺฉํ อนตฺถสํหิตํ, สสกฺกํ\footnote{สมฺปตฺตํ (ก)} ตํ รโหวาทํ น ภาเสยฺย.\csromanspacer{} ยมฺปิ ชญฺญา รโหวาทํ ภูตํ ตจฺฉํ อนตฺถสํหิตํ, ตสฺสปิ สิกฺเขยฺย อวจนาย.\csromanspacer{} ยญฺจ โข ชญฺญา รโหวาทํ ภูตํ ตจฺฉํ อตฺถสํหิตํ, ตตฺร กาลญฺญู อสฺส ตสฺส รโหวาทสฺส วจนาย.\csromanspacer{} ตตฺร ภิกฺขเว ยํ ชญฺญา สมฺมุขา ขีณวาทํ อภูตํ อตจฺฉํ อนตฺถสํหิตํ, สสกฺกํ1 ตํ สมฺมุขา ขีณวาทํ น ภาเสยฺย.\csromanspacer{} ยมฺปิ ชญฺญา สมฺมุขา ขีณวาทํ ภูตํ ตจฺฉํ อนตฺถสํหิตํ, ตสฺสปิ สิกฺเขยฺย อวจนาย.\csromanspacer{} ยญฺจ โข ชญฺญา สมฺมุขา ขีณวาทํ ภูตํ ตจฺฉํ อตฺถสํหิตํ, ตตฺร กาลญฺญู อสฺส ตสฺส สมฺมุขา ขีณวาทสฺส วจนาย. “รโหวาทํ น ภาเสยฺย, สมฺมุขา น ขีณํ ภเณ”ติ อิติ ยํ ตํ วุตฺตํ, อิทเมตํ ปฏิจฺจ วุตฺตํ.}

\proseitem{330}{“อตรมาโนว ภาเสยฺย โน ตรมาโน”ติ อิติ โข ปเนตํ วุตฺตํ, กิญฺเจตํ ปฏิจฺจ วุตฺตํ.\csromanspacer{} ตตฺร ภิกฺขเว ตรมานสฺส ภาสโต กาโยปิ กิลมติ, จิตฺตมฺปิ อุปหญฺญติ\footnote{อูหญฺญติ (สี)}, สโรปิ อุปหญฺญติ2, กณฺโฐปิ อาตุรียติ, อวิสฏฺฐมฺปิ โหติ อวิญฺเญยฺยํ ตรมานสฺส ภาสิตํ.\csromanspacer{} ตตฺร ภิกฺขเว อตรมานสฺส ภาสโต กาโยปิ น กิลมติ, จิตฺตมฺปิ น อุปหญฺญติ, สโรปิ น อุปหญฺญติ, กณฺโฐปิ น อาตุริยติ, วิสฏฺฐมฺปิ โหติ วิญฺเญยฺยํ อตรมานสฺส ภาสิตํ. “อตรมาโนว ภาเสยฺย โน ตรมาโน”ติ อิติ ยํ ตํ วุตฺตํ, อิทเมตํ ปฏิจฺจ วุตฺตํ.}

\proseitem{331}{\csromanfolio{278}“ชนปทนิรุตฺติํ นาภินิเวเสยฺย สมญฺญํ นาติธาเวยฺยา”ติ อิติ โข ปเนตํ วุตฺตํ, กิญฺเจตํ ปฏิจฺจ วุตฺตํ.\csromanspacer{} กถญฺจ ภิกฺขเว ชนปทนิรุตฺติยา จ อภินิเวโส โหติ สมญฺญาย จ อติสาโร.\csromanspacer{} อิธ ภิกฺขเว ตเทเวกจฺเจสุ ชนปเทสุ ปาตีติ สญฺชานนฺติ, ปตฺตนฺติ สญฺชานนฺติ, วิตฺตนฺติ\footnote{ปิฏฺฐนฺติ (สฺยา, กํ)} สญฺชานนฺติ, สราวนฺติ สญฺชานนฺติ, ธาโรปนฺติ\footnote{หโรสนฺติ (สฺยา, กํ) 3. ปิสีลนฺติ (สี, อิ), ปิปิลนฺติ (สฺยา, กํ)} สญฺชานนฺติ, โปณนฺติ สญฺชานนฺติ, ปิสีลวนฺติ3 สญฺชานนฺติ.\csromanspacer{} อิติ ยถา ยถา นํ เตสุ เตสุ ชนปเทสุ สญฺชานนฺติ, ตถา ตถา ถามสา ปรามาสา\footnote{ปรามสฺส (สี)} อภินิวิสฺส โวหรติ “อิทเมว สจฺจํ โมฆมญฺญนฺ”ติ.\csromanspacer{} เอวํ โข ภิกฺขเว ชนปทนิรุตฺติยา จ อภินิเวโส โหติ สมญฺญาย จ อติสาโร.}

\proseitem{332}{กถญฺจ ภิกฺขเว ชนปทนิรุตฺติยา จ อนภินิเวโส โหติ สมญฺญาย จ อนติสาโร.\csromanspacer{} อิธ ภิกฺขเว ตเทเวกจฺเจสุ ชนปเทสุ ปาตีติ สญฺชานนฺติ, ปตฺตนฺติ สญฺชานนฺติ, วิตฺตนฺติ สญฺชานนฺติ, สราวนฺติ สญฺชานนฺติ, ธาโรปนฺติ สญฺชานนฺติ, โปณนฺติ สญฺชานนฺติ, ปิสีลวนฺติ สญฺชานนฺติ.\csromanspacer{} อิติ ยถา ยถา นํ เตสุ เตสุ ชนปเทสุ สญฺชานนฺติ, “อิทํ กิร’เม\footnote{อิทํ กิร เต จ (ก)} อายสฺมนฺโต สนฺธาย โวหรนฺตี”ติ ตถา ตถา โวหรติ อปรามสํ.\csromanspacer{} เอวํ โข ภิกฺขเว ชนปทนิรุตฺติยา จ อนภินิเวโส โหติ สมญฺญาย จ อนติสาโร. “ชนปทนิรุตฺติํ นาภินิเวเสยฺย สมญฺญํ นาติธาเวยฺยา”ติ อิติ ยํ ตํ วุตฺตํ, อิทเมตํ ปฏิจฺจ วุตฺตํ.}

\proseitem{333}{ตตฺร ภิกฺขเว โย กามปฏิสนฺธิสุขิโน โสมนสฺสานุโยโค หีโน คมฺโม โปถุชฺชนิโก อนริโย อนตฺถสํหิโต, สทุกฺโข เอโส ธมฺโม ส-อุปฆาโต ส-อุปายาโส สปริฬาโห มิจฺฉาปฏิปทา.\csromanspacer{} ตสฺมา เอโส ธมฺโม สรโณ.\csromanspacer{} ตตฺร ภิกฺขเว โย กามปฏิสนฺธิสุขิโน โสมนสฺสานุโยคํ อนนุโยโค หีนํ คมฺมํ โปถุชฺชนิกํ อนริยํ อนตฺถสํหิตํ, อทุกฺโข เอโส ธมฺโม อนุปฆาโต อนุปายาโส อปริฬาโห สมฺมาปฏิปทา.\csromanspacer{} ตสฺมา เอโส ธมฺโม อรโณ.}

\proseitem{334}{ตตฺร ภิกฺขเว โย อตฺตกิลมถานุโยโค ทุกฺโข อนริโย อนตฺถสํหิโต, สทุกฺโข เอโส ธมฺโม ส-อุปฆาโต}

\csromanheader{2}{9. อรณวิภงฺคสุตฺต (139)}
\prose{\csromanfolio{279}ส-อุปายาโส สปริฬาโห มิจฺฉาปฏิปทา.\csromanspacer{} ตสฺมา เอโส ธมฺโม สรโณ.\csromanspacer{} ตตฺร ภิกฺขเว โย อตฺตกิลมถานุโยคํ อนนุโยโค ทุกฺขํ อนริยํ อนตฺถสํหิตํ, อทุกฺโข เอโส ธมฺโม อนุปฆาโต อนุปายาโส อปริฬาโห สมฺมาปฏิปทา.\csromanspacer{} ตสฺมา เอโส ธมฺโม อรโณ.}

\proseitem{335}{ตตฺร ภิกฺขเว ยายํ มชฺฌิมา ปฏิปทา ตถาคเตน อภิสมฺพุทฺธา จกฺขุกรณี ญาณกรณี อุปสมาย อภิญฺญาย สมฺโพธาย นิพฺพานาย สํวตฺตติ, อทุกฺโข เอโส ธมฺโม อนุปฆาโต อนุปายาโส อปริฬาโห สมฺมาปฏิปทา.\csromanspacer{} ตสฺมา เอโส ธมฺโม อรโณ.}

\proseitem{336}{ตตฺร ภิกฺขเว ยายํ อุสฺสาทนา จ อปสาทนา จ โน จ ธมฺมเทสนา, สทุกฺโข เอโส ธมฺโม ส-อุปฆาโต ส-อุปายาโส สปริฬาโห มิจฺฉาปฏิปทา.\csromanspacer{} ตสฺมา เอโส ธมฺโม สรโณ.\csromanspacer{} ตตฺร ภิกฺขเว ยายํ เนวุสฺสาทนา จ น อปสาทนา จ ธมฺมเทสนา จ, อทุกฺโข เอโส ธมฺโม อนุปฆาโต อนุปายาโส อปริฬาโห สมฺมาปฏิปทา.\csromanspacer{} ตสฺมา เอโส ธมฺโม อรโณ.}

\proseitem{337}{ตตฺร ภิกฺขเว ยมิทํ กามสุขํ มีฬฺหสุขํ โปถุชฺชนสุขํ อนริยสุขํ, สทุกฺโข เอโส ธมฺโม ส-อุปฆาโต ส- อุปายาโส สปริฬาโห มิจฺฉาปฏิปทา.\csromanspacer{} ตสฺมา เอโส ธมฺโม สรโณ.\csromanspacer{} ตตฺร ภิกฺขเว ยมิทํ เนกฺขมฺมสุขํ ปวิเวกสุขํ อุปสมสุขํ สมฺโพธิสุขํ, อทุกฺโข เอโส ธมฺโม อนุปฆาโต อนุปายาโส อปริฬาโห สมฺมาปฏิปทา.\csromanspacer{} ตสฺมา เอโส ธมฺโม อรโณ.}

\proseitem{338}{ตตฺร ภิกฺขเว ยฺวายํ รโหวาโท อภูโต อตจฺโฉ อนตฺถสํหิโต, สทุกฺโข เอโส ธมฺโม ส-อุปฆาโต ส-อุปายาโส สปริฬาโห มิจฺฉาปฏิปทา.\csromanspacer{} ตสฺมา เอโส ธมฺโม สรโณ.\csromanspacer{} ตตฺร ภิกฺขเว ยฺวายํ รโหวาโท ภูโต ตจฺโฉ อนตฺถสํหิโต, สทุกฺโข เอโส ธมฺโม ส-อุปฆาโต ส-อุปายาโส สปริฬาโห มิจฺฉาปฏิปทา.\csromanspacer{} ตสฺมา เอโส ธมฺโม สรโณ.\csromanspacer{} ตตฺร ภิกฺขเว ยฺวายํ รโหวาโท ภูโต ตจฺโฉ อตฺถสํหิโต, อทุกฺโข เอโส ธมฺโม อนุปฆาโต อนุปายาโส อปริฬาโห สมฺมาปฏิปทา.\csromanspacer{} ตสฺมา เอโส ธมฺโม อรโณ.}

\proseitem{339}{\csromanfolio{280}ตตฺร ภิกฺขเว ยฺวายํ สมฺมุขา ขีณวาโท อภูโต อตจฺโฉ อนตฺถสํหิโต, สทุกฺโข เอโส ธมฺโม ส-อุปฆาโต ส-อุปายาโส สปริฬาโห มิจฺฉาปฏิปทา.\csromanspacer{} ตสฺมา เอโส ธมฺโม สรโณ.\csromanspacer{} ตตฺร ภิกฺขเว ยฺวายํ สมฺมุขา ขีณวาโท ภูโต ตจฺโฉ อนตฺถสํหิโต, สทุกฺโข เอโส ธมฺโม ส-อุปฆาโต ส-อุปายาโส สปริฬาโห มิจฺฉาปฏิปทา.\csromanspacer{} ตสฺมา เอโส ธมฺโม สรโณ.\csromanspacer{} ตตฺร ภิกฺขเว ยฺวายํ สมฺมุขา ขีณวาโท ภูโต ตจฺโฉ อตฺถสํหิโต, อทุกฺโข เอโส ธมฺโม อนุปฆาโต อปริฬาโห สมฺมาปฏิปทา.\csromanspacer{} ตสฺมา เอโส ธมฺโม อรโณ.}

\proseitem{340}{ตตฺร ภิกฺขเว ยมิทํ ตรมานสฺส ภาสิตํ, สทุกฺโข เอโส ธมฺโม ส-อุปฆาโต ส-อุปายาโส สปริฬาโห มิจฺฉาปฏิปทา.\csromanspacer{} ตสฺมา เอโส ธมฺโม สรโณ.\csromanspacer{} ตตฺร ภิกฺขเว ยมิทํ อตรมานสฺส ภาสิตํ, อทุกฺโข.\csromanspacer{} ตตฺร ภิกฺขเว ยฺวายํ ชนปทนิรุตฺติยา จ อนภินิเวโส สมญฺญาย จ อนติสาโร, อทุกฺโข เอโส ธมฺโม อนุปฆาโต อนุปายาโส อปริฬาโห สมฺมาปฏิปทา.\csromanspacer{} ตสฺมา เอโส ธมฺโม อรโณ.}

\proseitem{341}{ตตฺร ภิกฺขเว ยฺวายํ ชนปทนิรุตฺติยา จ อภินิเวโส สมญฺญาย จ อติสาโร, สทุกฺโข เอโส ธมฺโม ส-อุปฆาโต ส-อุปายาโส สปริฬาโห มิจฺฉาปฏิปทา.\csromanspacer{} ตสฺมา เอโส ธมฺโม สรโณ.\csromanspacer{} ตตฺร ภิกฺขเว ยฺวายํ ชนปทนิรุตฺติยา จ อนภินิเวโส สมญฺญาย จ อนติสาโร, อทุกฺโข เอโส ธมฺโม อนุปฆาโต อนุปายาโส อปริฬาโห สมฺมาปฏิปทา.\csromanspacer{} ตสฺมา เอโส ธมฺโม อรโณ.}

\prose{ตสฺมาติห ภิกฺขเว สรณญฺจ ธมฺมํ ชานิสฺสาม, อรณญฺจ ธมฺมํ ชานิสฺสาม, สรณญฺจ ธมฺมํ ญตฺวา อรณญฺจ ธมฺมํ ญตฺวา อรณปฏิปทํ ปฏิปชฺชิสฺสามาติ เอวํ หิ โว ภิกฺขเว สิกฺขิตพฺพํ.\csromanspacer{} สุภูติ จ ปน ภิกฺขเว กุลปุตฺโต อรณปฏิปทํ ปฏิปนฺโนติ.}

\prose{อิทมโวจ ภควา.\csromanspacer{} อตฺตมนา เต ภิกฺขู ภควโต ภาสิตํ อภินนฺทุนฺติ.}

\nitthitam{อรณวิภงฺคสุตฺตํ นิฏฺฐิตํ นวมํ.}
\csromanheader{2}{10. ธาตุวิภงฺคสุตฺต (140)}
\csromanmatikaanchor{mtk.47}
\csromantocmarkref{subsection}{10. ธาตุวิภงฺคสุตฺต}{mtk.47}
\bookmark[dest={mtk.47},level=2]{10. ธาตุวิภงฺคสุตฺต}
\csromanheader{2}{10. ธาตุวิภงฺคสุตฺต}
\proseitem{342}{\csromanfolio{281}เอวํ เม สุตํ\csromandash{}เอกํ สมยํ ภควา มคเธสุ จาริกํ จรมาโน เยน ราชคหํ ตทวสริ, เยน ภคฺคโว กุมฺภกาโร เตนุปสงฺกมิ, อุปสงฺกมิตฺวา ภคฺควํ กุมฺภการํ เอตทโวจ “สเจ เต ภคฺคว อครุ, วิหเรมุ อาเวสเน\footnote{วิหรามาเวสเน (สี, อิ), วิหราม นิเวสเน (สฺยา, กํ), วิหเรมุ นิเวสเน (ก)} เอกรตฺตนฺ”ติ.\csromanspacer{} น โข เม ภนฺเต ครุ, อตฺถิ เจตฺถ ปพฺพชิโต ปฐมํ วาสูปคโต, สเจ โส อนุชานาติ, วิหรถ\footnote{วิหร (สี, อิ)} ภนฺเต ยถาสุขนฺติ.}

\prose{เตน โข ปน สมเยน ปุกฺกุสาติ นาม กุลปุตฺโต ภควนฺตํ อุทฺทิสฺส สทฺธาย อคารสฺมา อนคาริยํ ปพฺพชิโต, โส ตสฺมิํ กุมฺภการาเวสเน\footnote{กุมฺภการนิเวสเน (สฺยา, กํ, ก)} ปฐมํ วาสูปคโต โหติ.\csromanspacer{} อถ โข ภควา เยนายสฺมา ปุกฺกุสาติ เตนุปสงฺกมิ, อุปสงฺกมิตฺวา อายสฺมนฺตํ ปุกฺกุสาติํ เอตทโวจ “สเจ เต ภิกฺขุ อครุ, วิหเรมุ อาเวสเน เอกรตฺตนฺ”ติ.\csromanspacer{} อุรุนฺทํ\footnote{อูรุนฺทํ (สี, สฺยา, กํ, อิ), อุรุทฺธํ (ก) ที 2 สกฺกปญฺหสุตฺตฏีกา โอโลเกตพฺพา.} อาวุโส กุมฺภการาเวสนํ, วิหรตายสฺมา ยถาสุขนฺติ.}

\prose{อถ โข ภควา กุมฺภการาเวสนํ ปวิสิตฺวา เอกมนฺตํ ติณสนฺถารกํ\footnote{ติณสนฺถริกํ (สี), ติณสนฺถรกํ (สฺยา, กํ)} ปญฺญาเปตฺวา นิสีทิ ปลฺลงฺกํ อาภุชิตฺวา อุชุํ กายํ ปณิธาย ปริมุขํ สติํ อุปฏฺฐเปตฺวา.\csromanspacer{} อถ โข ภควา พหุเทว รตฺติํ นิสชฺชาย วีตินาเมสิ.\csromanspacer{} อายสฺมาปิ โข ปุกฺกุสาติ พหุเทว รตฺติํ นิสชฺชาย วีตินาเมสิ.}

\prose{อถ โข ภควโต เอตทโหสิ “ปาสาทิกํ โข อยํ กุลปุตฺโต อิริยติ, ยํนูนาหํ ปุจฺเฉยฺยนฺ”ติ.\csromanspacer{} อถ โข ภควา อายสฺมนฺตํ ปุกฺกุสาติํ เอตทโวจ “กํสิ ตฺวํ ภิกฺขุ อุทฺทิสฺส ปพฺพชิโต, โก วา เต สตฺถา, กสฺส วา ตฺวํ ธมฺมํ โรเจสี”ติ.\csromanspacer{} อตฺถาวุโส สมโณ โคตโม สกฺยปุตฺโต สกฺยกุลา ปพฺพชิโต, ตํ โข ปน ภควนฺตํ โคตมํ เอวํ กลฺยาโณ กิตฺติสทฺโท อพฺภุคฺคโต “อิติปิ โส ภควา อรหํ สมฺมาสมฺพุทฺโธ วิชฺชาจรณสมฺปนฺโน สุคโต โลกวิทู อนุตฺตโร ปุริสทมฺมสารถิ สตฺถา เทวมนุสฺสานํ พุทฺโธ \csromanfolio{282}ภควา”ติ.\csromanspacer{} ตาหํ ภควนฺตํ อุทฺทิสฺส ปพฺพชิโต, โส จ เม ภควา สตฺถา, ตสฺส จาหํ ภควโต ธมฺมํ โรเจมีติ.\csromanspacer{} กหํ ปน ภิกฺขุ เอตรหิ โส ภควา วิหรติ อรหํ สมฺมาสมฺพุทฺโธติ.\csromanspacer{} อตฺถาวุโส อุตฺตเรสุ ชนปเทสุ สาวตฺถิ นาม นครํ, ตตฺถ โส ภควา เอตรหิ วิหรติ อรหํ สมฺมาสมฺพุทฺโธติ.\csromanspacer{} ทิฏฺฐปุพฺโพ ปน เต ภิกฺขุ โส ภควา, ทิสฺวา จ ปน ชาเนยฺยาสีติ.\csromanspacer{} น โข เม อาวุโส ทิฏฺฐปุพฺโพ โส ภควา, ทิสฺวา จาหํ น ชาเนยฺยนฺติ.}

\prose{อถ โข ภควโต เอตทโหสิ “มมญฺจ ขฺวายํ\footnote{มํ ตฺวายํ (สี), มมํ ขฺวายํ (สฺยา, กํ), มํ ขฺวายํ (อิ)} กุลปุตฺโต อุทฺทิสฺส ปพฺพชิโต, ยํนูนสฺสาหํ ธมฺมํ เทเสยฺยนฺ”ติ.\csromanspacer{} อถ โข ภควา อายสฺมนฺตํ ปุกฺกุสาติํ อามนฺเตสิ “ธมฺมํ เต ภิกฺขุ เทเสสฺสามิ, ตํ สุณาหิ สาธุกํ มนสิ กโรหิ, ภาสิสฺสามี”ติ. “เอวมาวุโส”ติ โข อายสฺมา ปุกฺกุสาติ ภควโต ปจฺจสฺโสสิ.\csromanspacer{} ภควา เอตทโวจ\csromandash{}}

\proseitem{343}{ฉธาตุโร\footnote{ฉทฺธาตุโร (สี)} อยํ ภิกฺขุ ปุริโส ฉผสฺสายตโน อฏฺฐารสมโนปวิจาโร จตุราธิฏฺฐาโน.\csromanspacer{} ยตฺถ ฐิตํ มญฺญสฺสวา นปฺปวตฺตนฺติ, มญฺญสฺสเว โข ปน นปฺปวตฺตมาเน มุนิ สนฺโตติ วุจฺจติ.\csromanspacer{} ปญฺญํ นปฺปมชฺเชยฺย, สจฺจมนุรกฺเขยฺย, จาคมนุพฺรูเหยฺย, สนฺติเมว โส สิกฺเขยฺยาติ.\csromanspacer{} อยมุทฺเทโส ธาตุวิภงฺคสฺส\footnote{ฉธาตุวิภงฺคสฺส (สี, สฺยา, กํ, อิ)}.}

\proseitem{344}{“ฉธาตุโร อยํ ภิกฺขุ ปุริโส”ติ อิติ โข ปเนตํ วุตฺตํ, กิญฺเจตํ ปฏิจฺจ วุตฺตํ. (ฉยิมา ภิกฺขุ ธาตุโย,)\footnote{( ) นตฺถิ สี-อิ-โปตฺถเกสุ.} ปถวีธาตุ อาโปธาตุ เตโชธาตุ วาโยธาตุ อากาสธาตุ วิญฺญาณธาตุ. “ฉธาตุโร อยํ ภิกฺขุ ปุริโส”ติ อิติ ยํ ตํ วุตฺตํ, อิทเมตํ ปฏิจฺจ วุตฺตํ.}

\proseitem{345}{“ฉผสฺสายตโน อยํ ภิกฺขุ ปุริโส”ติ อิติ โข ปเนตํ วุตฺตํ, กิญฺเจตํ ปฏิจฺจ วุตฺตํ.\csromanspacer{} จกฺขุสมฺผสฺสายตนํ โสตสมฺผสฺสายตนํ ฆานสมฺผสฺสายตนํ ชิวฺหาสมฺผสฺสายตนํ กายสมฺผสฺสายตนํ มโนสมฺผสฺสายตนํ. “ฉผสฺสายตโน อยํ ภิกฺขุ ปุริโส”ติ อิติ ยํ ตํ วุตฺตํ, อิทเมตํ ปฏิจฺจ วุตฺตํ.}

\csromanheader{2}{10. ธาตุวิภงฺคสุตฺต (140)}
\proseitem{346}{\csromanfolio{283}“อฏฺฐารสมโนปวิจาโร อยํ ภิกฺขุ ปุริโส”ติ อิติ โข ปเนตํ วุตฺตํ, กิญฺเจตํ ปฏิจฺจ วุตฺตํ.\csromanspacer{} จกฺขุนา รูปํ ทิสฺวา โสมนสฺสฏฺฐานิยํ รูปํ อุปวิจรติ, โทมนสฺสฏฺฐานิยํ รูปํ อุปวิจรติ, อุเปกฺขาฏฺฐานิยํ รูปํ อุปวิจรติ.\csromanspacer{} โสเตน สทฺทํ สุตฺวา ฯเปฯ.\csromanspacer{} ฆาเนน คนฺธํ ฆายิตฺวา.\csromanspacer{} ชิวฺหาย รสํ สายิตฺวา.\csromanspacer{} กาเยน โผฏฺฐพฺพํ ผุสิตฺวา.\csromanspacer{} มนสา ธมฺมํ วิญฺญาย โสมนสฺสฏฺฐานิยํ ธมฺมํ อุปวิจรติ, โทมนสฺสฏฺฐานิยํ ธมฺมํ อุปวิจรติ, อุเปกฺขาฏฺฐานิยํ ธมฺมํ อุปวิจรติ.\csromanspacer{} อิติ ฉ โสมนสฺสุปวิจารา ฉ โทมนสฺสุปวิจารา ฉ อุเปกฺขุปวิจารา. “อฏฺฐารสมโนปวิจาโร อยํ ภิกฺขุ ปุริโส”ติ อิติ ยํ ตํ วุตฺตํ, อิทเมตํ ปฏิจฺจ วุตฺตํ.}

\proseitem{347}{“จตุราธิฏฺฐาโน อยํ ภิกฺขุ ปุริโส”ติ อิติ โข ปเนตํ วุตฺตํ, กิญฺเจตํ ปฏิจฺจ วุตฺตํ.\csromanspacer{} ปญฺญาธิฏฺฐาโน สจฺจาธิฏฺฐาโน จาคาธิฏฺฐาโน อุปสมาธิฏฺฐาโน. “จตุราธิฏฺฐาโน อยํ ภิกฺขุ ปุริโส”ติ อิติ ยํ ตํ วุตฺตํ, อิทเมตํ ปฏิจฺจ วุตฺตํ.}

\proseitem{348}{“ปญฺญํ นปฺปมชฺเชยฺย, สจฺจมนุรกฺเขยฺย, จาคมนุพฺรูเหยฺย, สนฺติเมว โส สิกฺเขยฺยา”ติ อิติ โข ปเนตํ วุตฺตํ, กิญฺเจตํ ปฏิจฺจ วุตฺตํ.\csromanspacer{} กถญฺจ ภิกฺขุ ปญฺญํ นปฺปมชฺชติ.\csromanspacer{} ฉยิมา ภิกฺขุ ธาตุโย.\csromanspacer{} ปถวีธาตุ อาโปธาตุ เตโชธาตุ วาโยธาตุ อากาสธาตุ วิญฺญาณธาตุ.}

\proseitem{349}{กตมา จ ภิกฺขุ ปถวีธาตุ.\csromanspacer{} ปถวีธาตุ สิยา อชฺฌตฺติกา, สิยา พาหิรา.\csromanspacer{} กตมา จ ภิกฺขุ อชฺฌตฺติกา ปถวีธาตุ.\csromanspacer{} ยํ อชฺฌตฺตํ ปจฺจตฺตํ กกฺขฬํ ขริคตํ อุปาทินฺนํ\footnote{อุปาทิณฺณํ (อิ, ก)}.\csromanspacer{} เสยฺยถิทํ, เกสา โลมา นขา ทนฺตา ตโจ, มํสํ นฺหารุ อฏฺฐิ อฏฺฐิมิญฺชํ\footnote{อฏฺฐิมิญฺชา (สี, อิ)} วกฺกํ, หทยํ ยกนํ กิโลมกํ ปิหกํ ปปฺผาสํ, อนฺตํ อนฺตคุณํ อุทริยํ กรีสํ, ยํ วา ปนญฺญมฺปิ กิญฺจิ อชฺฌตฺตํ ปจฺจตฺตํ กกฺขฬํ ขริคตํ อุปาทินฺนํ.\csromanspacer{} อยํ วุจฺจติ ภิกฺขุ อชฺฌตฺติกา ปถวีธาตุ.\csromanspacer{} ยา เจว โข ปน อชฺฌตฺติกา ปถวีธาตุ, ยา จ พาหิรา ปถวีธาตุ, ปถวีธาตุเรเวสา.\csromanspacer{} ตํ “เนตํ มม, เนโสหมสฺมิ, น เมโส อตฺตา”ติ เอวเมตํ ยถาภูตํ สมฺมปฺปญฺญาย ทฏฺฐพฺพํ.\csromanspacer{} เอวเมตํ ยถาภูตํ สมฺมปฺปญฺญาย ทิสฺวา ปถวีธาตุยา นิพฺพินฺทติ, ปถวีธาตุยา จิตฺตํ วิราเชติ.}

\proseitem{350}{\csromanfolio{284}กตมา จ ภิกฺขุ อาโปธาตุ.\csromanspacer{} อาโปธาตุ สิยา อชฺฌตฺติกา, สิยา พาหิรา.\csromanspacer{} กตมา จ ภิกฺขุ อชฺฌตฺติกา อาโปธาตุ.\csromanspacer{} ยํ อชฺฌตฺตํ ปจฺจตฺตํ อาโป อาโปคตํ อุปาทินฺนํ.\csromanspacer{} เสยฺยถิทํ, ปิตฺตํ เสมฺหํ ปุพฺโพ โลหิตํ เสโท เมโท อสฺสุ วสา เขโฬ สิงฺฆาณิกา ลสิกา มุตฺตํ, ยํ วา ปนญฺญมฺปิ กิญฺจิ อชฺฌตฺตํ ปจฺจตฺตํ อาโป อาโปคตํ อุปาทินฺนํ.\csromanspacer{} อยํ วุจฺจติ ภิกฺขุ อชฺฌตฺติกา อาโปธาตุ.\csromanspacer{} ยา เจว โข ปน อชฺฌตฺติกา อาโปธาตุ, ยา จ พาหิรา อาโปธาตุ, อาโปธาตุเรเวสา.\csromanspacer{} ตํ “เนตํ มม, เนโสหมสฺมิ, น เมโส อตฺตา”ติ เอวเมตํ ยถาภูตํ สมฺมปฺปญฺญาย ทฏฺฐพฺพํ.\csromanspacer{} เอวเมตํ ยถาภูตํ สมฺมปฺปญฺญาย ทิสฺวา อาโปธาตุยา นิพฺพินฺทติ, อาโปธาตุยา จิตฺตํ วิราเชติ.}

\proseitem{351}{กตมา จ ภิกฺขุ เตโชธาตุ.\csromanspacer{} เตโชธาตุ สิยา อชฺฌตฺติกา, สิยา พาหิรา.\csromanspacer{} กตมา จ ภิกฺขุ อชฺฌตฺติกา เตโชธาตุ.\csromanspacer{} ยํ อชฺฌตํ ปจฺจตฺตํ เตโช เตโชคตํ อุปาทินฺนํ.\csromanspacer{} เสยฺยถิทํ, เยน จ สนฺตปฺปติ, เยน จ ชีรียติ, เยน จ ปริฑยฺหติ, เยน จ อสิตปีตขายิตสายิตํ สมฺมา ปริณามํ คจฺฉติ, ยํ วา ปนญฺญมฺปิ กิญฺจิ อชฺฌตฺตํ ปจฺจตฺตํ เตโช เตโชคตํ อุปาทินฺนํ.\csromanspacer{} อยํ วุจฺจติ ภิกฺขุ อชฺฌตฺติกา เตโชธาตุ.\csromanspacer{} ยา เจว โข ปน อชฺฌตฺติกา เตโชธาตุ, ยา จ พาหิรา เตโชธาตุ, เตโชธาตุเรเวสา.\csromanspacer{} ตํ “เนตํ มม, เนโสหมสฺมิ, น เมโส อตฺตา”ติ เอวเมตํ ยถาภูตํ สมฺมปฺปญฺญาย ทฏฺฐพฺพํ.\csromanspacer{} เอวเมตํ ยถาภูตํ สมฺมปฺปญฺญาย ทิสฺวา เตโชธาตุยา นิพฺพินฺทติ, เตโชธาตุยา จิตฺตํ วิราเชติ.}

\proseitem{352}{กตมา จ ภิกฺขู วาโยธาตุ.\csromanspacer{} วาโยธาตุ สิยา อชฺฌตฺติกา, สิยา พาหิรา.\csromanspacer{} กตมา จ ภิกฺขุ อชฺฌตฺติกา วาโยธาตุ.\csromanspacer{} ยํ อชฺฌตฺตํ ปจฺจตฺตํ วาโย วาโยคตํ อุปาทินฺนํ.\csromanspacer{} เสยฺยถิทํ, อุทฺธงฺคมา วาตา, อโธคมา วาตา, กุจฺฉิสยา วาตา, โกฏฺฐาสยา\footnote{โกฏฺฐสยา (สี, สฺยา, กํ, อิ)} วาตา, องฺคมงฺคานุสาริโน วาตา, อสฺสาโส ปสฺสาโส อิติ, ยํ วา ปนญฺญมฺปิ กิญฺจิ อชฺฌตฺตํ ปจฺจตฺตํ วาโย วาโยคตํ อุปาทินฺนํ.\csromanspacer{} อยํ วุจฺจติ ภิกฺขุ อชฺฌตฺติกา วาโยธาตุ.\csromanspacer{} ยา เจว โข ปน อชฺฌตฺติกา วาโยธาตุ, ยา จ พาหิรา วาโยธาตุ, วาโยธาตุเรเวสา.\csromanspacer{} ตํ “เนตํ มม,}

\csromanheader{2}{10. ธาตุวิภงฺคสุตฺต (140)}
\prose{\csromanfolio{285}เนโสหมสฺมิ, น เมโส อตฺตา”ติ เอวเมตํ ยถาภูตํ สมฺมปฺปญฺญาย ทฏฺฐพฺพํ.\csromanspacer{} เอวเมตํ ยถาภูตํ สมฺมปฺปญฺญาย ทิสฺวา วาโยธาตุยา นิพฺพินฺทติ, วาโยธาตุยา จิตฺตํ วิราเชติ.}

\proseitem{353}{กตมา จ ภิกฺขุ อากาสธาตุ.\csromanspacer{} อากาสธาตุ สิยา อชฺฌตฺติกา, สิยา พาหิรา.\csromanspacer{} กตมา จ ภิกฺขุ อชฺฌตฺติกา อากาสธาตุ.\csromanspacer{} ยํ อชฺฌตฺตํ ปจฺจตฺตํ อากาสํ อากาสคตํ อุปาทินฺนํ.\csromanspacer{} เสยฺยถิทํ, กณฺณจฺฉิทฺทํ นาสจฺฉิทฺทํ มุขทฺวารํ, เยน จ อสิตปีตขายิตสายิตํ อชฺโฌหรติ, ยตฺถ จ อสิตปีตขายิตสายิตํ สนฺติฏฺฐติ, เยน จ อสิตปีตขายิตสายิตํ อโธภาคํ\footnote{อโธภาคา (สี, สฺยา, กํ, อิ) เทวทูตสุตฺเตน สเมติ.} นิกฺขมติ, ยํ วา ปนญฺญมฺปิ กิญฺจิ อชฺฌตฺตํ ปจฺจตฺตํ อากาสํ อากาสคตํ อฆํ อฆคตํ วิวรํ วิวรคตํ อสมฺผุฏฺฐํ มํสโลหิเตหิ อุปาทินฺนํ.\csromanspacer{} อยํ วุจฺจติ ภิกฺขุ อชฺฌตฺติกา อากาสธาตุ.\csromanspacer{} ยา เจว โข ปน อชฺฌตฺติกา อากาสธาตุ, ยา จ พาหิรา อากาสธาตุ, อากาสธาตุเรเวสา.\csromanspacer{} ตํ “เนตํ มม, เนโสหมสฺมิ, น เมโส อตฺตา”ติ เอวเมตํ ยถาภูตํ สมฺมปฺปญฺญาย ทฏฺฐพฺพํ.\csromanspacer{} เอวเมตํ ยถาภูตํ สมฺมปฺปญฺญาย ทิสฺวา อากาสธาตุยา นิพฺพินฺทติ, อากาสธาตุยา จิตฺตํ วิราเชติ.}

\proseitem{354}{อถาปรํ วิญฺญาณํเยว อวสิสฺสติ ปริสุทฺธํ ปริโยทาตํ, เตน จ วิญฺญาเณน กิํ\footnote{เตน วิญฺญาเณน กิญฺจ (สี)} วิชานาติ.\csromanspacer{} สุขนฺติปิ วิชานาติ, ทุกฺขนฺติปิ วิชานาติ, อทุกฺขมสุขนฺติปิ วิชานาติ.\csromanspacer{} สุขเวทนิยํ ภิกฺขุ ผสฺสํ ปฏิจฺจ อุปฺปชฺชติ สุขา เวทนา, โส สุขํ เวทนํ เวทยมาโน สุขํ เวทนํ เวทยามีติ ปชานาติ.\csromanspacer{} ตสฺเสว สุขเวทนิยสฺส ผสฺสสฺส นิโรธา ยํ ตชฺชํ เวทยิตํ สุขเวทนิยํ ผสฺสํ ปฏิจฺจ อุปฺปนฺนา สุขา เวทนา, สา นิรุชฺฌติ, สา วูปสมฺมตีติ ปชานาติ.}

\proseitem{355}{ทุกฺขเวทนิยํ ภิกฺขุ ผสฺสํ ปฏิจฺจ อุปฺปชฺชติ ทุกฺขา เวทนา, โส ทุกฺขํ เวทนํ เวทยมาโน ทุกฺขํ เวทนํ เวทยามีติ ปชานาติ.\csromanspacer{} ตสฺเสว ทุกฺขเวทนิยสฺส ผสฺสสฺส นิโรธา ยํ ตชฺชํ เวทยิตํ ทุกฺขเวทนิยํ ผสฺสํ ปฏิจฺจ อุปฺปนฺนา ทุกฺขา เวทนา, สา นิรุชฺฌติ, สา วูปสมฺมตีติ ปชานาติ.}

\proseitem{356}{อทุกฺขมสุขเวทนิยํ ภิกฺขุ ผสฺสํ ปฏิจฺจ อุปฺปชฺชติ อทุกฺขมสุขา เวทนา, โส อทุกฺขมสุขํ เวทนํ เวทยมาโน อทุกฺขมสุขํ เวทนํ \csromanfolio{286}เวทยามีติ ปชานาติ.\csromanspacer{} ตสฺเสว อทุกฺขมสุขเวทนิยสฺส ผสฺสสฺส นิโรธา ยํ ตชฺชํ เวทยิตํ อทุกฺขมสุขเวทนิยํ ผสฺสํ ปฏิจฺจ อุปฺปนฺนา อทุกฺขมสุขา เวทนา, สา นิรุชฺฌติ, สา วูปสมฺมตีติ ปชานาติ.}

\proseitem{357}{เสยฺยถาปิ ภิกฺขุ ทฺวินฺนํ กฏฺฐานํ สํฆฏฺฏา\footnote{สมฺผสฺส (สี, อิ), สงฺฆฏา (สฺยา, กํ)} สโมธานา อุสฺมา ชายติ เตโช อภินิพฺพตฺตติ, เตสํเยว ทฺวินฺนํ กฏฺฐานํ นานาภาวา วินิกฺเขปา ยา ตชฺชา อุสฺมา, สา นิรุชฺฌติ, สา วูปสมฺมติ.\csromanspacer{} เอวเมว โข ภิกฺขุ สุขเวทนิยํ ผสฺสํ ปฏิจฺจ อุปฺปชฺชติ สุขา เวทนา, โส สุขํ เวทนํ เวทยมาโน สุขํ เวทนํ เวทยามีติ ปชานาติ.\csromanspacer{} ตสฺเสว สุขเวทนิยสฺส ผสฺสสฺส นิโรธา ยํ ตชฺชํ เวทยิตํ สุขเวทนิยํ ผสฺสํ ปฏิจฺจ อุปฺปนฺนา สุขา เวทนา, สา นิรุชฺฌติ, สา วูสมฺมตีติ ปชานาติ.}

\proseitem{358}{ทุกฺขเวทนิยํ ภิกฺขุ ผสฺสํ ปฏิจฺจ อุปฺปชฺชติ ทุกฺขา เวทนา, โส ทุกฺขํ เวทนํ เวทยมาโน ทุกฺขํ เวทนํ เวทยามีติ ปชานาติ.\csromanspacer{} ตสฺเสว ทุกฺขเวทนิยสฺส ผสฺสสฺส นิโรธา ยํ ตชฺชํ เวทยิตํ ทุกฺขเวทนิยํ ผสฺสํ ปฏิจฺจ อุปฺปนฺนา ทุกฺขา เวทนา, สา นิรุชฺฌติ, สา วูปสมฺมตีติ ปชานาติ.}

\proseitem{359}{อทุกฺขมสุขเวทนิยํ ภิกฺขุ ผสฺสํ ปฏิจฺจ อุปฺปชฺชติ อทุกฺขมสุขา เวทนา, โส อทุกฺขมสุขํ เวทนํ เวทยมาโน อทุกฺขมสุขํ เวทนํ เวทยามีติ ปชานาติ.\csromanspacer{} ตสฺเสว อทุกฺขมสุขเวทนิยสฺส ผสฺสสฺส นิโรธา ยํ ตชฺชํ เวทยิตํ อทุกฺขมสุขเวทนิยํ ผสฺสํ ปฏิจฺจ อุปฺปนฺนา อทุกฺขมสุขา เวทนา, สา นิรุชฺฌติ, สา วูปสมฺมตีติปชานาติ.}

\proseitem{360}{อถาปรํ อุเปกฺขาเยว อวสิสฺสติ ปริสุทฺธา ปริโยทาตา มุทุ จ กมฺมญฺญา จ ปภสฺสรา จ.\csromanspacer{} เสยฺยถาปิ ภิกฺขุ ทกฺโข สุวณฺณกาโร วา สุวณฺณกาเรนฺเตวาสี วา อุกฺกํ พนฺเธยฺย, อุกฺกํ พนฺธิตฺวา อุกฺกามุขํ อาลิมฺเปยฺย, อุกฺกามุขํ อาลิมฺเปตฺวา สณฺฑาเสน ชาตรูปํ คเหตฺวา อุกฺกามุเข ปกฺขิเปยฺย.\csromanspacer{} ตเมนํ กาเลน กาลํ อภิธเมยฺย, กาเลน กาลํ อุทเกน ปริปฺโผเสยฺย, กาเลน กาลํ อชฺฌุเปกฺเขยฺย.\csromanspacer{} ตํ โหติ ชาตรูปํ\footnote{ชาตรูปํ ธนฺตํ (สี, อิ) 3. นิหตํ (สฺยา, กํ, ก)} สุธนฺตํ นิทฺธนฺตํ นีหฏํ3 นินฺนีตกสาวํ\footnote{นิหตกสาวํ (ก)} มุทุ จ กมฺมญฺญญฺจ ปภสฺสรญฺจ.\csromanspacer{} ยสฺสา ยสฺสา จ ปิฬนฺธนวิกติยา อากงฺขติ, ยทิ ปฏฺฏิกาย\footnote{ปวฏฺฏิกาย (สี, สฺยา)}}

\csromanheader{2}{10. ธาตุวิภงฺคสุตฺต (140)}
\prose{\csromanfolio{287}ยทิ กุณฺฑลาย ยทิ คีเวยฺยกาย ยทิ สุวณฺณมาลาย, ตญฺจสฺส อตฺถํ อนุโภติ.\csromanspacer{} เอวเมว โข ภิกฺขุ อถาปรํ อุเปกฺขาเยว อวสิสฺสติ ปริสุทฺธา ปริโยทาตา มุทุ จ กมฺมญฺญา จ ปภสฺสรา จ.}

\proseitem{361}{โส เอวํ ปชานาติ “อิมญฺเจ อหํ อุเปกฺขํ เอวํ ปริสุทฺธํ เอวํ ปริโยทาตํ อากาสานญฺจายตนํ อุปสํหเรยฺยํ, ตทนุธมฺมญฺจ จิตฺตํ ภาเวยฺยํ, เอวํ เม อยํ อุเปกฺขา ตํนิสฺสิตา ตทุปาทานา จิรํ ทีฆมทฺธานํ ติฏฺเฐยฺย.\csromanspacer{} อิมญฺเจ อหํ อุเปกฺขํ เอวํ ปริสุทฺธํ เอวํ ปริโยทาตํ วิญฺญาณญฺจายตนํ อุปสํหเรยฺยํ, ตทนุธมฺมญฺจ จิตฺตํ ภาเวยฺยํ, เอวํ เม อยํ อุเปกฺขา ตํนิสฺสิตา ตทุปาทานา จิรํ ทีฆมทฺธานํ ติฏฺเฐยฺย.\csromanspacer{} อิมญฺเจ อหํ อุเปกฺขํ เอวํ ปริสุทฺธํ เอวํ ปริโยทาตํ อากิญฺจญฺญายตนํ อุปสํหเรยฺยํ, ตทนุธมฺมญฺจ จิตฺตํ ภาเวยฺยํ, เอวํ เม อยํ อุเปกฺขา ตํนิสฺสิตา ตทุปาทานา จิรํ ทีฆมทฺธานํ ติฏฺเฐยฺย.\csromanspacer{} อิมญฺเจ อหํ อุเปกฺขํ เอวํ ปริสุทฺธํ เอวํ ปริโยทาตํ เนวสญฺญานาสญฺญายตนํ อุปสํหเรยฺยํ, ตทนุธมฺมญฺจ จิตฺตํ ภาเวยฺยํ, เอวํ เม อยํ อุเปกฺขา ตํนิสฺสิตา ตทุปาทานา จิรํ ทีฆมทฺธานํ ติฏฺเฐยฺยา”ติ.}

\proseitem{362}{โส เอวํ ปชานาติ “อิมญฺเจ อหํ อุเปกฺขํ เอวํ ปริสุทฺธํ เอวํ ปริโยทาตํ อากาสานญฺจายตนํ อุปสํหเรยฺยํ, ตทนุธมฺมญฺจ จิตฺตํ ภาเวยฺยํ, สงฺขตเมตํ.\csromanspacer{} อิมญฺเจ อหํ อุเปกฺขํ เอวํ ปริสุทฺธํ เอวํ ปริโยทาตํ วิญฺญาณญฺจายตนํ อุปสํหเรยฺยํ, ตทนุธมฺมญฺจ จิตฺตํ ภาเวยฺยํ, สงฺขตเมตํ.\csromanspacer{} อิมญฺเจ อหํ อุเปกฺขํ เอวํ ปริสุทฺธํ เอวํ ปริโยทาตํ อากิญฺจญฺญายตนํ อุปสํหเรยฺยํ, ตทนุธมฺมญฺจ จิตฺตํ ภาเวยฺยํ, สงฺขตเมตํ.\csromanspacer{} อิมญฺเจ อหํ อุเปกฺขํ เอวํ ปริสุทฺธํ เอวํ ปริโยทาตํ เนวสญฺญานาสญฺญายตนํ อุปสํหเรยฺยํ, ตทนุธมฺมญฺจ จิตฺตํ ภาเวยฺยํ, สงฺขตเมตนฺติ.}

\prose{โส เนว ตํ อภิสงฺขโรติ น อภิสญฺเจตยติ ภวาย วา วิภวาย วา, โส อนภิสงฺขโรนฺโต อนภิสญฺเจตยนฺโต ภวาย วา วิภวาย วา น กิญฺจิ โลเก อุปาทิยติ, อนุปาทิยํ น ปริตสฺสติ อปริตสฺสํ ปจฺจตฺตํเยว ปรินิพฺพายติ, “ขีณา ชาติ, วุสิตํ พฺรหฺมจริยํ, กตํ กรณียํ, นาปรํ อิตฺถตฺตายา”ติ ปชานาติ.}

\proseitem{363}{โส สุขญฺเจ เวทนํ เวเทติ, สา อนิจฺจาติ ปชานาติ, อนชฺโฌสิตาติ ปชานาติ, อนภินนฺทิตาติ ปชานาติ.\csromanspacer{} ทุกฺขญฺเจ เวทนํ \csromanfolio{288}เวเทติ, สา อนิจฺจาติ ปชานาติ, อนชฺโฌสิตาติ ปชานาติ, อนภินนฺทิตาติ ปชานาติ.\csromanspacer{} อทุกฺขมสุขญฺเจ เวทนํ เวเทติ, สา อนิจฺจาติ ปชานาติ, อนชฺโฌสิตาติ ปชานาติ, อนภินนฺทิตาติ ปชานาติ.}

\proseitem{364}{โส สุขญฺเจ เวทนํ เวเทติ, วิสํยุตฺโต นํ เวเทติ.\csromanspacer{} ทุกฺขญฺเจ เวทนํ เวเทติ, วิสํยุตฺโต นํ เวเทติ.\csromanspacer{} อทุกฺขมสุขญฺเจ เวทนํ เวเทติ, วิสํยุตฺโต นํ เวเทติ.\csromanspacer{} โส กายปริยนฺติกํ เวทนํ เวทยมาโน กายปริยนฺติกํ เวทนํ เวทยามีติ ปชานาติ, ชีวิตปริยนฺติกํ เวทนํ เวทยมาโน ชีวิตปริยนฺติกํ เวทนํ เวทยามีติ ปชานาติ.\csromanspacer{} กายสฺส เภทา ปรํ มรณา อุทฺธํ ชีวิตปริยาทานา อิเธว สพฺพเวทยิตานิ อนภินนฺทิตานิ สีตีภวิสฺสนฺตีติ ปชานาติ.}

\proseitem{365}{เสยฺยถาปิ ภิกฺขุ เตลญฺจ ปฏิจฺจ วฏฺฏิญฺจ ปฏิจฺจ เตลปฺปทีโป ฌายติ, ตสฺเสว เตลสฺส จ วฏฺฏิยา จ ปริยาทานา อญฺญสฺส จ อนุปหารา\footnote{อนุปาหารา (สี, อิ), อนุปาทานา (ก)} อนาหาโร นิพฺพายติ.\csromanspacer{} เอวเมว โข ภิกฺขุ กายปริยนฺติกํ เวทนํ เวทยมาโน กายปริยนฺติกํ เวทนํ เวทยามีติ ปชานาติ, ชีวิตปริยนฺติกํ เวทนํ เวทยมาโน ชีวิตปริยนฺติกํ เวทนํ เวทยามีติ ปชานาติ.\csromanspacer{} กายสฺส เภทา ปรํ ปรํ มรณา อุทฺธํ ชีวิตปริยาทานา อิเธว สพฺพเวทยิตานิ อนภินนฺทิตานิ สีตีภวิสฺสนฺตีติ ปชานาติ.\csromanspacer{} ตสฺมา เอวํ สมนฺนาคโต ภิกฺขุ อิมินา ปรเมน ปญฺญาธิฏฺฐาเนน สมนฺนาคโต โหติ.\csromanspacer{} เอสา หิ ภิกฺขุ ปรมา อริยา ปญฺญา, ยทิทํ สพฺพทุกฺขกฺขเย ญาณํ.}

\proseitem{366}{ตสฺส สา วิมุตฺติ สจฺเจ ฐิตา อกุปฺปา โหติ.\csromanspacer{} ตํ หิ ภิกฺขุ มุสา ยํ โมสธมฺมํ, ตํ สจฺจํ ยํ อโมสธมฺมํ นิพฺพานํ.\csromanspacer{} ตสฺมา เอวํ สมนฺนาคโต ภิกฺขุ อิมินา ปรเมน สจฺจาธิฏฺฐาเนน สมนฺนาคโต โหติ.\csromanspacer{} เอตํ หิ ภิกฺขุ ปรมํ อริยสจฺจํ, ยทิทํ อโมสธมฺมํ นิพฺพานํ.}

\proseitem{367}{ตสฺเสว โข ปน ปุพฺเพ อวิทฺทสุโน อุปธี โหนฺติ สมตฺตา สมาทินฺนา.\csromanspacer{} ตฺยาสฺส ปหีนา โหนฺติ อุจฺฉินฺนมูลา ตาลาวตฺถุกตา อนภาวํกตา อายติํ อนุปฺปาทธมฺมา.\csromanspacer{} ตสฺมา เอวํ สมนฺนาคโต}

\csromanheader{2}{10. ธาตุวิภงฺคสุตฺต (140)}
\prose{\csromanfolio{289}ภิกฺขุ อิมินา ปรเมน จาคาธิฏฺฐาเนน สมนฺนาคโต โหติ.\csromanspacer{} เอโส หิ ภิกฺขุ ปรโม อริโย จาโค, ยทิทํ สพฺพูปธิปฏินิสฺสคฺโค.}

\proseitem{368}{ตสฺเสว โข ปน ปุพฺเพ อวิทฺทสุโน อภิชฺฌา โหติ ฉนฺโท สาราโค, สฺวาสฺส ปหีโน โหติ อุจฺฉินฺนมูโล ตาลาวตฺถุกโต อนภาวํกโต อายติํ อนุปฺปาทธมฺโม.\csromanspacer{} ตสฺเสว โข ปน ปุพฺเพ อวิทฺทสุโน อาฆาโต โหติ พฺยาปาโท สมฺปโทโส, สฺวาสฺส ปหีโน โหติ อุจฺฉินฺนมูโล ตาลาวตฺถุกโต อนภาวํกโต อายติํ อนุปฺปาทธมฺโม.\csromanspacer{} ตสฺเสว โข ปน ปุพฺเพ อวิทฺทสุโน อวิชฺชา โหติ สมฺโมโห, สฺวาสฺส ปหีโน โหติ อุจฺฉินฺนมูโล ตาลาวตฺถุกโต อนภาวํกโต อายติํ อนุปฺปาทธมฺโม.\csromanspacer{} ตสฺมา เอวํ สมนฺนาคโต ภิกฺขุ อิมินา ปรเมน อุปสมาธิฏฺฐาเนน สมนฺนาคโต โหติ.\csromanspacer{} เอโส หิ ภิกฺขุ ปรโม อริโย อุปสโม, ยทิทํ ราคโทสโมหานํ อุปสโม. “ปญฺญํ นปฺปมชฺเชยฺย, สจฺจมนุรกฺขยฺย, จาคมนุพฺรูเหยฺย, สนฺติเมว โส สิกฺเขยฺยา”ติ อิติ ยํ ตํ วุตฺตํ อิทเมตํ ปฏิจฺจ วุตฺตํ.}

\proseitem{369}{“ยตฺถ ฐิตํ มญฺญสฺสวา นปฺปวตฺตนฺติ, มญฺญสฺสเว โข ปน นปฺปวตฺตมาเน มุนิ สนฺโตติ วุจฺจตี”ติ อิติ โข ปเนตํ วุตฺตํ, กิญฺเจตํ ปฏิจฺจ วุตฺตํ.\csromanspacer{} อสฺมีติ ภิกฺขุ มญฺญิตเมตํ, อยมหมสฺมีติ มญฺญิตเมตํ.\csromanspacer{} ภวิสฺสนฺติ มญฺญิตเมตํ, น ภวิสฺสนฺติ มญฺญิตเมตํ.\csromanspacer{} รูปี ภวิสฺสนฺติ มญฺญิตเมตํ, อรูปี ภวิสฺสนฺติ มญฺญิตเมตํ, สญฺญี ภวิสฺสนฺติ มญฺญิตเมตํ, อสญฺญีภวิสฺสนฺติ มญฺญิตเมตํ, เนวสญฺญีนาสญฺญี ภวิสฺสนฺติ มญฺญิตเมตํ.\csromanspacer{} มญฺญิตํ ภิกฺขุ โรโค มญฺญิตํ คณฺโฑ มญฺญิตํ สลฺลํ, สพฺพมญฺญิตานํเตฺวว ภิกฺขุ สมติกฺกมา มุนิ สนฺโตติ วุจฺจติ.\csromanspacer{} มุนิ โข ปน ภิกฺขุ สนฺโต น ชายติ น ชียติ น มียติ น กุปฺปติ น ปิเหติ.\csromanspacer{} ตญฺหิสฺส ภิกฺขุ นตฺถิ, เยน ชาเยถ.\csromanspacer{} อชายมาโน กิํ ชียิสฺสติ, อชียมาโน กิํ มียิสฺสติ, อมียมาโน กิํ กุปฺปิสฺสติ, อกุปฺปมาโน กิสฺส\footnote{กิํ (ก)} ปิเหสฺสติ. “ยตฺถ ฐิตํ มญฺญสฺสวา นปฺปวตฺตนฺติ, มญฺญสฺสเว โข ปน นปฺปวตฺตมาเน มุนิ สนฺโตติ วุจฺจตี”ติ อิติ ยํ ตํ วุตฺตํ, อิทเมตํ ปฏิจฺจ วุตฺตํ.\csromanspacer{} อิมํ โข เม ตฺวํ ภิกฺขุ สํขิตฺเตน ฉธาตุวิภงฺคํ ธาเรหีติ.}

\proseitem{370}{\csromanfolio{290}อถ โข อายสฺมา ปุกฺกุสาติ “สตฺถา กิร เม อนุปฺปตฺโต, สุคโต กิร เม อนุปฺปตฺโต, สมฺมาสมฺพุทฺโธ กิร เม อนุปฺปตฺโต”ติ อุฏฺฐายาสนา เอกํสํ จีวรํ กตฺวา ภควโต ปาเทสุ สิรสา นิปติตฺวา ภควนฺตํ เอตทโวจ “อจโย มํ ภนฺเต อจฺจคมา ยถาพาลํ ยถามูฬฺหํ ยถา- อกุสลํ, โยหํ ภควนฺตํ อาวุโสวาเทน สมุทาจริตพฺพํ อมญฺญิสฺสํ, ตสฺส เม ภนฺเต ภควา อจฺจยํ อจฺจยโต ปฏิคฺคณฺหาตุ อายติํ สํวรายา”ติ.\csromanspacer{} ตคฺฆ ตฺวํ ภิกฺขุ อจฺจโย อจฺจคมา ยถาพาลํ ยถามูฬฺหํ ยถา-อกุสลํ, ยํ มํ ตฺวํ อาวุโสวาเทน สมุทาจริตพฺพํ อมญฺญิตฺถ, ยโต จ โข ตฺวํ ภิกฺขุ อจฺจยํ อจฺจยโต ทิสฺวา ยถาธมฺมํ ปฏิกโรสิ, ตํ เต มยํ ปฏิคฺคณฺหาม.\csromanspacer{} วุทฺธิเหสา ภิกฺขุ อริยสฺส วินเย, โย อจฺจยํ อจฺจยโต ทิสฺวา ยถาธมฺมํ ปฏิกโรติ, อายติํ สํวรํ อาปชฺชตีติ.\csromanspacer{} ลเภยฺยาหํ ภนฺเต ภควโต สนฺติเก อุปสมฺปทนฺติ.\csromanspacer{} ปริปุณฺณํ ปน เต ภิกฺขุ ปตฺตจีวรนฺติ.\csromanspacer{} น โข เม ภนฺเต ปริปุณฺณํ ปตฺตจีวรนฺติ.\csromanspacer{} น โข ภิกฺขุ ตถาคตา อปริปุณฺณปตฺตจีวรํ อุปสมฺปาเทนฺตีติ.\csromanspacer{} อถ โข อายสฺมา ปุกฺกุสาติ ภควโต ภาสิตํ อภินนฺทิตฺวา อนุโมทิตฺวา อุฏฺฐายาสนา ภควนฺตํ อภิวาเทตฺวา ปทกฺขิณํ กตฺวา ปตฺตจีวรปริเยสนํ ปกฺกามิ.}

\prose{อถ โข อายสฺมนฺตํ ปุกฺกุสาติํ ปตฺตจีวรปริเยสนํ จรนฺตํ วิพฺภนฺตา คาวี\footnote{ภนฺตคาวี (สี, อิ), คาวี (สฺยา, กํ)} ชีวิตา โวโรเปสิ.\csromanspacer{} อถ โข สมฺพหุลา ภิกฺขู เยน ภควา เตนุปสงฺกมิํสุ, อุปสงฺกมิตฺวา ภควนฺตํ อภิวาเทตฺวา เอกมนฺตํ นิสีทิํสุ.\csromanspacer{} เอกมนฺตํ นิสินฺนา โข เต ภิกฺขู ภควนฺตํ เอตทโวจุํ “โย โส ภนฺเต ปุกฺกุสาติ นาม กุลปุตฺโต ภควตา สํขิตฺเตน โอวาเทน โอวทิโต, โส กาลงฺกโต, ตสฺส กา คติ, โก อภิสมฺปราโย”ติ.\csromanspacer{} ปณฺฑิโต ภิกฺขเว ปุกฺกุสาติ กุลปุตฺโต ปจฺจปาทิ ธมฺมสฺสานุธมฺมํ, น จ มํ ธมฺมาธิกรณํ วิเหเสสิ\footnote{วิเหเฐสิ (สี, สฺยา, กํ) วิเหเสติ (ก)}.\csromanspacer{} ปุกฺกุสาติ ภิกฺขเว กุลปุตฺโต ปญฺจนฺนํ โอรมฺภาคิยานํ สํโยชนานํ ปริกฺขยา โอปปาติโก ตตฺถ ปรินิพฺพายี อนาวตฺติธมฺโม ตสฺมา โลกาติ.}

\prose{อิทมโวจ ภควา.\csromanspacer{} อตฺตมนา เต ภิกฺขู ภควโต ภาสิตํ อภินนฺทุนฺติ.}

\nitthitam{ธาตุวิงฺคสุตฺตํ นิฏฺฐิตํ ทสมํ.}
\csromanheader{2}{11. สจฺจวิภงฺคสุตฺต (141)}
\csromanmatikaanchor{mtk.48}
\csromantocmarkref{subsection}{11. สจฺจวิภงฺคสุตฺต}{mtk.48}
\bookmark[dest={mtk.48},level=2]{11. สจฺจวิภงฺคสุตฺต}
\csromanheader{2}{11. สจฺจวิภงฺคสุตฺต}
\proseitem{371}{\csromanfolio{291}เอวํ เม สุตํ\csromandash{}เอกํ สมยํ ภควา พาราณสิยํ วิหรติ อิสิปตเน มิคทาเย.\csromanspacer{} ตตฺร โข ภควา ภิกฺขู อามนฺเตสิ “ภิกฺขโว”ติ. “ภทนฺเต”ติ เต ภิกฺขู ภควโต ปจฺจสฺโสสุํ.\csromanspacer{} ภควา เอตทโวจ\csromandash{}}

\prose{ตถาคเตน ภิกฺขเว อรหตา สมฺมาสมฺพุทฺเธน พาราณสิยํ อิสิปตเน มิคทาเย อนุตฺตรํ ธมฺมจกฺกํ ปวตฺติตํ อปฺปฏิวตฺติยํ สมเณน วา พฺราหฺมเณน วา เทเวน วา มาเรน วา พฺรหฺมุนา วา เกนจิ วา โลกสฺมิํ, ยทิทํ จตุนฺนํ อริยสจฺจานํ อาจิกฺขนา เทสนา ปญฺญาปนา ปฏฺฐปนา วิวรณา วิภชนา อุตฺตานีกมฺมํ.\csromanspacer{} กตเมสํ จตุนฺนํ, ทุกฺขสฺส อริยสจฺจสฺส อาจิกฺขนา เทสนา ปญฺญาปนา ปฏฺฐปนา วิวรณา วิภชนา อุตฺตานีกมฺมํ.\csromanspacer{} ทุกฺขสมุทยสฺส อริยสจฺจสฺส อาจิกฺขนา เทสนา ปญฺญาปนา ปฏฺฐปนา วิวรณา วิภชนา อุตฺตานีกมฺมํ.\csromanspacer{} ทุกฺขนิโรธสฺส อริยสจฺจสฺส อาจิกฺขนา เทสนา ปญฺญาปนา ปฏฺฐปนา วิวรณา วิภชนา อุตฺตานีกมฺมํ.\csromanspacer{} ทุกฺขนิโรธคามินิยา ปฏิปทาย อริยสจฺจสฺส อาจิกฺขนา เทสนา ปญฺญาปนา ปฏฺฐปนา วิวรณา วิภชนา อุตฺตานีกมฺมํ.\csromanspacer{} ตถาคเตน ภิกฺขเว อรหตา สมฺมาสมฺพุทฺเธน พาราณสิยํ อิสิปตเน มิคทาเย อนุตฺตรํ ธมฺมจกฺกํ ปวตฺติตํ อปฺปฏิวตฺติยํ สมเณน วา พฺราหฺมเณน วา เทเวน วา มาเรน วา พฺรหฺมุนา วา เกนจิ วา โลกสฺมิํ, ยทิทํ อิเมสํ จตุนฺนํ อริยสจฺจานํ อาจิกฺขนา เทสนา ปญฺญาปนา ปฏฺฐปนา วิวรณา วิภชนา อุตฺตานีกมฺมํ.}

\prose{เสวถ ภิกฺขเว สาริปุตฺตโมคฺคลฺลาเน, ภชถ ภิกฺขเว สาริปุตฺตโมคฺคลฺลาเน.\csromanspacer{} ปณฺฑิตา ภิกฺขู อนุคฺคาหกา สพฺรหฺมจารีนํ.\csromanspacer{} เสยฺยถาปิ ภิกฺขเว ชเนตา\footnote{ชเนตฺติ (สี, อิ)} เอวํ สาริปุตฺโต, เสยฺยถาปิ ชาตสฺส อาปาเทตา เอวํ โมคฺคลฺลาโน.\csromanspacer{} สาริปุตฺโต ภิกฺขเว โสตาปตฺติผเล วิเนติ, โมคฺคลฺลาโน อุตฺตมตฺเถ.\csromanspacer{} สาริปุตฺโต ภิกฺขเว ปโหติ จตฺตาริ อริยสจฺจานิ วิตฺถาเรน อาจิกฺขิตุํ เทเสตุํ ปญฺญาเปตุํ ปฏฺฐเปตุํ วิวริตุํ วิภชิตุํ อุตฺตานีกาตุนฺติ.\csromanspacer{} อิทมโวจ ภควา, อิทํ วตฺวาน สุคโต อุฏฺฐายาสนา วิหารํ ปาวิสิ.}

\proseitem{372}{ตตฺร โข อายสฺมา สาริปุตฺโต อจิรปกฺกนฺตสฺส ภควโต ภิกฺขู อามนฺเตสิ “อาวุโส ภิกฺขเว”ติ. “อาวุโส”ติ โข เต \csromanfolio{292}ภิกฺขู อายสฺมโต สาริปุตฺตสฺส ปจฺจสฺโสสุํ.\csromanspacer{} อายสฺมา สาริปุตฺโต เอตทโวจ\csromandash{}}

\prose{ตถาคเตน อาวุโส อรหตา สมฺมาสมฺพุทฺเธน พาราณสิยํ อิสิปตเน มิคทาเย อนุตฺตรํ ธมฺมจกฺกํ ปวตฺติตํ อปฺปฏิวตฺติยํ สมเณน วา พฺราหฺมเณน วา เทเวน วา มาเรน วา พฺรหฺมุนา วา เกนจิ วา โลกสฺมิํ, ยทิทํ จตุนฺนํ อริยสจฺจานํ อาจิกฺขนา เทสนา ปญฺญาปนา ปฏฺฐปนา วิวรณา วิภชนา อุตฺตานีกมฺมํ.\csromanspacer{} กตเมสํ จตุนฺนํ, ทุกฺขสฺส อริยสจฺจสฺส อาจิกฺขนา เทสนา ปญฺญาปนา ปฏฺฐปนา วิวรณา วิภชนา อุตฺตานีกมฺมํ.\csromanspacer{} ทุกฺขสมุทยสฺส อริยสจฺจสฺส อาจิกฺขนา เทสนา ปญฺญาปนา ปฏฺฐปนา วิวรณา วิภชนา อุตฺตานีกมฺมํ.\csromanspacer{} ทุกฺขนิโรธสฺส อริยสจฺจสฺส อาจิกฺขนา เทสนา ปญฺญาปนา ปฏฺฐปนา วิวรณา วิภชนา อุตฺตานีกมฺมํ.\csromanspacer{} ทุกฺขนิโรธคามินิยา ปฏิปทาย อริยสจฺจสฺส อาจิกฺขนา เทสนา ปญฺญาปนา ปฏฺฐปนา วิวรณา วิภชนา อุตฺตานีกมฺมํ.}

\proseitem{373}{กตมญฺจาวุโส ทุกฺขํ อริยสจฺจํ.\csromanspacer{} ชาติปิ ทุกฺขา, ชราปิ ทุกฺขา, มรณมฺปิ ทุกฺขํ, โสกปริเทวทุกฺขโทมนสฺสุปายาสาปิ ทุกฺขา, ยมฺปิจฺฉํ น ลภติ ตมฺปิ ทุกฺขํ, สํขิตฺเตน ปญฺจุปาทานกฺขนฺธา ทุกฺขา.}

\prose{กตมา จาวุโส ชาติ.\csromanspacer{} ยา เตสํ เตสํ สตฺตานํ ตมฺหิ ตมฺหิ สตฺตนิกาเย ชาติ สญฺชาติ โอกฺกนฺติ อภินิพฺพตฺติ ขนฺธานํ ปาตุภาโว อายตนานํ ปฏิลาโภ.\csromanspacer{} อยํ วุจฺจตาวุโส ชาติ.}

\prose{กตมา จาวุโส ชรา.\csromanspacer{} ยา เตสํ เตสํ สตฺตานํ ตมฺหิ ตมฺหิ สตฺตนิกาเย ชรา ชีรณตา ขณฺฑิจฺจํ ปาลิจฺจํ วลิตฺตจตา อายุโน สํหานิ อินฺทฺริยานํ ปริปาโก.\csromanspacer{} อยํ วุจฺจตาวุโส ชรา.}

\prose{กตมญฺจาวุโส มรณํ.\csromanspacer{} ยา เตสํ เตสํ สตฺตานํ ตมฺหา ตมฺหา สตฺตนิกายา จุติ จวนตา เภโท อนฺตรธานํ มจฺจุ มรณํ กาลํกิริยา ขนฺธานํ เภโท กเฬวรสฺส นิกฺเขโป ชีวิตินฺทฺริยสฺสุปจฺเฉโท.\csromanspacer{} อิทํ วุจฺจตาวุโส มรณํ.}

\prose{กตโม จาวุโส โสโก.\csromanspacer{} โย โข อาวุโส อญฺญตรญฺญตเรน พฺยสเนน สมนฺนาคตสฺส อญฺญตรญฺญตเรน ทุกฺขธมฺเมน ผุฏฺฐสฺส โสโก โสจนา โสจิตตฺตํ อนฺโตโสโก อนฺโตปริโสโก.\csromanspacer{} อยํ วุจฺจตาวุโส โสโก.}

\csromanheader{2}{11. สจฺจวิภงฺคสุตฺต (141)}
\prose{\csromanfolio{293}กตโม จาวุโส ปริเทโว.\csromanspacer{} โย โข อาวุโส อญฺญตรญฺญตเรน พฺยสเนน สมนฺนาคตสฺส อญฺญตรญฺญตเรน ทุกฺขธมฺเมน ผุฏฺฐสฺส อาเทโว ปริเทโว อาเทวนา ปริเทวนา อาเทวิตตฺตํ ปริเทวิตตฺตํ.\csromanspacer{} อยํ วุจฺจตาวุโส ปริเทโว.}

\prose{กตมญฺจาวุโส ทุกฺขํ.\csromanspacer{} ยํ โข อาวุโส กายิกํ ทุกฺขํ กายิกํ อสาตํ กายสมฺผสฺสชํ ทุกฺขํ อสาตํ เวทยิตํ.\csromanspacer{} อิทํ วุจฺจตาวุโส ทุกฺขํ.}

\prose{กตมญฺจาวุโส โทมนสฺสํ.\csromanspacer{} ยํ โข อาวุโส เจตสิกํ ทุกฺขํ เจตสิกํ อสาตํ มโนสมฺผสฺสชํ ทุกฺขํ อสาตํ เวทยิตํ.\csromanspacer{} อิทํ วุจฺจตาวุโส โทมนสฺสํ.}

\prose{กตโม จาวุโส อุปายาโส.\csromanspacer{} โย โข อาวุโส อญฺญตรญฺญตเรน พฺยสเนน สมนฺนาคตสฺส อญฺญตรญฺญตเรน ทุกฺขธมฺเมน ผุฏฺฐสฺส อายาโส อุปายาโส อายาสิตตฺตํ อุปายาสิตตฺตํ.\csromanspacer{} อยํ วุจฺจตาวุโส อุปายาโส.}

\prose{กตมญฺจาวุโส ยมฺปิจฺฉํ น ลภติ ตมฺปิ ทุกฺขํ.\csromanspacer{} ชาติธมฺมานํ อาวุโส สตฺตานํ เอวํ อิจฺฉา อุปฺปชฺชติ “อโห วต มยํ น ชาติธมฺมา อสฺสาม, น จ วต โน ชาติ อาคจฺเฉยฺยา”ติ.\csromanspacer{} น โข ปเนตํ อิจฺฉาย ปตฺตพฺพํ, อิทมฺปิ ยมฺปิจฺฉํ น ลภติ ตมฺปิ ทุกฺขํ.\csromanspacer{} ชราธมฺมานํ อาวุโส สตฺตานํ ฯเปฯ.\csromanspacer{} พฺยาธิธมฺมานํ อาวุโส สตฺตานํ.\csromanspacer{} มรณธมฺมานํ อาวุโส สตฺตานํ.\csromanspacer{} โสกปริเทวทุกฺขโทมนสฺสุปา- ยาสธมฺมานํ อาวุโส สตฺตานํ เอวํ อิจฺฉา อุปฺปชฺชติ “อโห วต มยํ น โสกปริเทวทุกฺขโทมนสฺสุปายาสธมฺมา อสฺสาม, น จ วต โน โสกปริเทวทุกฺขโทมนสฺสุปายาสา อาคจฺเฉยฺยุนฺ”ติ.\csromanspacer{} น โข ปเนตํ อิจฺฉาย ปตฺตพฺพํ, อิทมฺปิ ยมฺปิจฺฉํ น ลภติ ตมฺปิ ทุกฺขํ.}

\prose{กตเม จาวุโส สํขิตฺเตน ปญฺจุปาทานกฺขนฺธา ทุกฺขา.\csromanspacer{} เสยฺยถิทํ, รูปุปาทานกฺขนฺโธ เวทนุปาทานกฺขนฺโธ สญฺญุปาทานกฺขนฺโธ สงฺขารุปาทานกฺขนฺโธ วิญฺญาณุปาทานกฺขนฺโธ.\csromanspacer{} อิเม วุจฺจนฺตาวุโส สํขิตฺเตน ปญฺจุปาทานกฺขนฺธา ทุกฺขา.\csromanspacer{} อิทํ วุจฺจตาวุโส ทุกฺขํ อริยสจฺจํ.}

\proseitem{374}{กตมญฺจาวุโส ทุกฺขสมุทยํ\footnote{ทุกฺขสมุทโย (สฺยา, กํ)} อริยสจฺจํ.\csromanspacer{} ยายํ ตณฺหา โปโนพฺภวิกา\footnote{โปโนภวิกา (สี, อิ)} นนฺทีราสหคตา\footnote{นนฺทิราคสหคตา (สี, สฺยา, กํ, อิ)} ตตฺรตตฺราภินนฺทินี.\csromanspacer{} เสยฺยถิทํ, \csromanfolio{294}กามตณฺหา ภคตณฺหา วิภวตณฺหา.\csromanspacer{} อิทํ วุจฺจตาวุโส ทุกฺขสมุทยํ\footnote{ทุกฺขสมุทโย (สฺยา, กํ)} อริยสจฺจํ.}

\prose{กตมญฺจาวุโส ทุกฺขนิโรธํ\footnote{ทุกฺขนิโรโธ (สฺยา, กํ)} อริยสจฺจํ.\csromanspacer{} โย ตสฺสาเยว ตณฺหาย อเสสวิราคนิโรโธ จาโค ปฏินิสฺสคฺโค มุตฺติ อนาลโย.\csromanspacer{} อิทํ วุจฺจตาวุโส ทุกฺขนิโรธํ อริยสจฺจํ.}

\proseitem{375}{กตมญฺจาวุโส ทุกฺขนิโรธคามินี ปฏิปทา อริยสจฺจํ.\csromanspacer{} อยเมว อริโย อฏฺฐงฺคิโก มคฺโค.\csromanspacer{} เสยฺยถิทํ, สมฺมาทิฏฺฐิ สมฺมาสงฺกปฺโป สมฺมาวาจา สมฺมากมฺมนฺโต สมฺมา-อาชีโว สมฺมาวายาโม สมฺมาสติ สมฺมาสมาธิ.}

\prose{กตมาจาวุโส สมฺมาทิฏฺฐิ.\csromanspacer{} ยํ โข อาวุโส ทุกฺเข ญาณํ ทุกฺขสมุทเย ญาณํ ทุกฺขนิโรเธ ญาณํ ทุกฺขนิโรธคามินิยา ปฏิปทาย ญาณํ.\csromanspacer{} อยํ วุจฺจตาวุโส สมฺมาทิฏฺฐิ.}

\prose{กตโม จาวุโส สมฺมาสงฺกปฺโป.\csromanspacer{} เนกฺขมฺมสงฺกปฺโป อพฺยาปาทสงฺกปฺโป อวิหิํสาสงฺกปฺโป.\csromanspacer{} อยํ วุจฺจตาวุโส สมฺมาสงฺกปฺโป.}

\prose{กตมา จาวุโส สมฺมาวาจา.\csromanspacer{} มุสาวาทา เวรมณี, ปิสุณาย วาจาย เวรมณี, ผรุสาย วาจาย เวรมณี, สมฺผปฺปลาปา เวรมณี.\csromanspacer{} อยํ วุจฺจตาวุโส สมฺมาวาจา.}

\prose{กตโม จาวุโส สมฺมากมฺมนฺโต.\csromanspacer{} ปาณาติปาตา เวรมณี, อทินฺนาทานา เวรมณี, กาเมสุมิจฺฉาจารา เวรมณี.\csromanspacer{} อยํ วุจฺจตาวุโส สมฺมากมฺมนฺโต.}

\prose{กตโม จาวุโส สมฺมา-อาชีโว.\csromanspacer{} อิธาวุโส อริยสาวโก มิจฺฉา-อาชีวํ ปหาย สมฺมา-อาชีเวน ชีวิกํ กปฺเปติ.\csromanspacer{} อยํ วุจฺจตาวุโส สมฺมา-อาชีโว.}

\prose{กตโม จาวุโส สมฺมาวายาโม.\csromanspacer{} อิธาวุโส ภิกฺขุ อนุปฺปนฺนานํ ปาปกานํ อกุสลานํ ธมฺมานํ อนุปฺปาทาย ฉนฺทํ ชเนติ วายมติ วีริยํ อารภติ จิตฺตํ ปคฺคณฺหาติ ปทหติ, อุปฺปนฺนานํ ปาปกานํ อกุสลานํ ธมฺมานํ ปหานาย ฉนฺทํ ชเนติ วายมติ วีริยํ อารภติ จิตฺตํ ปคฺคณฺหาติ ปทหติ, อนุปฺปนฺนานํ กุสลานํ ธมฺมานํ อุปฺปาทาย ฉนฺทํ ชเนติ วายมติ วีริยํ อารภติ จิตฺตํ ปคฺคณฺหาติ ปทหติ, อุปฺปนฺนานํ กุสลานํ ธมฺมานํ}

\csromanheader{2}{12. ทกฺขิณาวิภงฺคสุตฺต (142)}
\prose{\csromanfolio{295}ฐิติยา อสมฺโมสาย ภิยฺโยภาวาย เวปุลฺลาย ภาวนาย ปาริปูริยา ฉนฺทํ ชเนติ วายมติ วีริยํ อารภติ จิตฺตํ ปคฺคณฺหาติ ปทหติ.\csromanspacer{} อยํ วุจฺจตาวุโส สมฺมาวายาโม.}

\prose{กตมา จาวุโส สมฺมาสติ.\csromanspacer{} อิธาวุโส ภิกฺขุ กาเย กายานุปสฺสี วิหรติ อาตาปี สมฺปชาโน สติมา วิเนยฺย โลเก อภิชฺฌาโทมนสฺสํ, เวทนาสุ เวทนานุปสฺสี วิหรติ ฯเปฯ.\csromanspacer{} จิตฺเต จิตฺตานุปสฺสี วิหรติ.\csromanspacer{} ธมฺเมสุ ธมฺมานุปสฺสี วิหรติ อาตาปี สมฺปชาโน สติมา วิเนยฺย โลเก อภิชฺฌาโทมนสฺสํ.\csromanspacer{} อยํ วุจฺจตาวุโส สมฺมาสติ.}

\prose{กตโม จาวุโส สมฺมาสมาธิ.\csromanspacer{} อิธาวุโส ภิกฺขุ วิวิจฺเจว กาเมหิ วิวิจฺจ อกุสเลหิ ธมฺเมหิ สวิตกฺกํ สวิจารํ วิเวกชํ ปีติสุขํ ปฐมํ ฌานํ อุปสมฺปชฺช วิหรติ.\csromanspacer{} วิตกฺกวิจารานํ วูปสมา อชฺฌตฺตํ สมฺปสาทนํ เจตโส เอโกทิภาวํ อวิตกฺกํ อวิจารํ สมาธิชํ ปีติสุขํ ทุติยํ ฌานํ อุปสมฺปชฺช วิหรติ.\csromanspacer{} ปีติยา จ วิราคา อุเปกฺขโก จ วิหรติ ฯเปฯ ตติยํ ฌานํ ฯเปฯ จตุตฺถํ ฌานํ อุปสมฺปชฺช วิหรติ.\csromanspacer{} อยํ วุจฺจตาวุโส สมฺมาสมาธิ.\csromanspacer{} อิทํ วุจฺจตาวุโส ทุกฺขนิโรธคามินี ปฏิปทา อริยสจฺจํ.}

\prose{ตถาคเตนาวุโส อรหตา สมฺมาสมฺพุทฺเธน พาราณสิยํ อิสิปตเน มิคทาเย อนุตฺตรํ ธมฺมจกฺกํ ปวตฺติตํ อปฺปฏิวตฺติยํ สมเณน วา พฺราหฺมเณน วา เทเวน วา มาเรน วา พฺรหฺมุนา วา เกนจิ วาโลกสฺมิํ, ยทิทํ อิเมสํ จตุนฺนํ อริยสจฺจานํ อาจิกฺขนา เทสนา ปญฺญาปนา ปฏฺฐปนา วิวรณา วิภชนา อุตฺตานีกมฺมนฺติ.}

\prose{อิทมโวจ อายสฺมา สาริปุตฺโต.\csromanspacer{} อตฺตมนา เต ภิกฺขู อายสฺมโต สาริปุตฺตสฺส ภาสิตํ อภินนฺทุนฺติ.}

\nitthitamruled{สจฺจวิภงฺคสุตฺตํ นิฏฺฐิตํ เอกาทสมํ.}
\csromanmatikaanchor{mtk.49}
\csromantocmarkref{subsection}{12. ทกฺขิณาวิภงฺคสุตฺต}{mtk.49}
\bookmark[dest={mtk.49},level=2]{12. ทกฺขิณาวิภงฺคสุตฺต}
\csromanheader{2}{12. ทกฺขิณาวิภงฺคสุตฺต}
\proseitem{376}{เอวํ เม สุตํ\csromandash{}เอกํ สมยํ ภควา สกฺเกสุ วิหรติ กปิลวตฺถุสฺมิํ นิโคฺรธาราเม.\csromanspacer{} อถ โข มหาปชาปติ\footnote{มหาปชาปตี (สี, สฺยา, กํ, อิ)} โคตมี นวํ ทุสฺสยุคํ อาทาย เยน ภควา เตนุปสงฺกมิ, อุปสงฺกมิตฺวา ภควนฺตํ \csromanfolio{296}อภิวาเทตฺวา เอกมนฺตํ นิสีทิ, เอกมนฺตํ นิสินฺนา โข มหาปชาปติ โคตมี ภควนฺตํ เอตทโวจ “อิทํ เม ภนฺเต นวํ ทุสฺสยุคํ ภควนฺตํ อุทฺทิสฺส สามํ กนฺตํ สามํ วายิตํ, ตํ เม ภนฺเต ภควา ปฏิคฺคณฺหาตุ อนุกมฺปํ อุปาทายา”ติ.\csromanspacer{} เอวํ วุตฺเต ภควา มหาปชาปติํ โคตมิํ เอตทโวจ “สํเฆ โคตมิ เทหิ, สํเฆ เต ทินฺเน อหญฺเจว ปูชิโต ภวิสฺสามิ สํโฆ จา”ติ.\csromanspacer{} ทุติยมฺปิ โข มหาปชาปติ โคตมี ภควนฺตํ เอตทโวจ “อิทํ เม ภนฺเต นวํ ทุสฺสยุคํ ภควนฺตํ อุทฺทิสฺส สามํ กนฺตํ สามํ วายิตํ, ตํ เม ภนฺเต ภควา ปฏิคฺคณฺหาตุ อนุกมฺปํ อุปาทายา”ติ.\csromanspacer{} ทุติยมฺปิ โข ภควา มหาปชาปติํ โคตมิํ เอตทโวจ “สํเฆ โคตมิ เทหิ, สํเฆ เต ทินฺเน อหญฺเจว ปูชิโต ภวิสฺสามิ สํโฆ จา”ติ.\csromanspacer{} ตติยมฺปิ โข มหาปชาปติ โคตมี ภควนฺตํ เอตทโวจ “อิทํ เม ภนฺเต นวํ ทุสฺสยุคํ ภควนฺตํ อุทฺทิสฺส สามํ กนฺตํ สามํ วายิตํ, ตํ เม ภนฺเต ภควา ปฏิคฺคณฺหาตุ อนุกมฺปํ อุปาทายา”ติ.\csromanspacer{} ตติยมฺปิ โข ภควา มหาปชาปติํ โคตมิํ เอตทโวจ “สํเฆ โคตมิ เทหิ, สํเฆ เต ทินฺเน อหญฺเจว ปูชิโต ภวิสฺสามิ สํโฆ จา”ติ.}

\proseitem{377}{เอวํ วุตฺเต อายสฺมา อานนฺโท ภควนฺตํ เอตทโวจ “ปฏิคฺคณฺหาตุ ภนฺเต ภควา มหาปชาปติยา โคตมิยา นวํ ทุสฺสยุคํ, พหูปการา\footnote{พหุการา (สฺยา, กํ)} ภนฺเต มหาปชาปติ โคตมี ภควโต มาตุจฺฉา อาปาทิกา โปสิกา ขีรสฺส ทายิกา, ภควนฺตํ ชเนตฺติยา กาลงฺกตาย ถญฺญํ ปาเยสิ, ภควาปิ ภนฺเต พหูปกาโร มหาปชาปติยา โคตมิยา, ภควนฺตํ ภนฺเต อาคมฺม มหาปชาปติ โคตมี พุทฺธํ สรณํ คตา, ธมฺมํ สรณํ คตา, สํฆํ สรณํ คตา.\csromanspacer{} ภควนฺตํ ภนฺเต อาคมฺม มหาปชาปติ โคตมี ปาณาติปาตา ปฏิวิรตา, อทินฺนาทานา ปฏิวิรตา, กาเมสุมิจฺฉาจารา ปฏิวิรตา, มุสาวาทา ปฏิวิรตา, สุราเมรยมชฺชปมาทฏฺฐานา ปฏิวิรตา.\csromanspacer{} ภควนฺตํ ภนฺเต อาคมฺม มหาปชาปติ โคตมี พุทฺเธ อเวจฺจปฺปสาเทน สมนฺนาคตา, ธมฺเม อเวจฺจปฺปสาเทน สมนฺนาคตา, สํเฆ อเวจฺจปฺปสาเทน สมนฺนาคตา, อริยกนฺเตหิ สีเลหิ สมนฺนาคตา.\csromanspacer{} ภควนฺตํ ภนฺเต อาคมฺม มหาปชาปติ โคตมี ทุกฺเข นิกฺกงฺขา, ทุกฺขสมุทเย นิกฺกงฺขา, ทุกฺขนิโรเธ นิกฺกงฺขา, ทุกฺขนิโรธคามินิยา ปฏิปทาย นิกฺกงฺขา.\csromanspacer{} ภควาปิ ภนฺเต พหูปกาโร มหาปชาปติยา โคตมิยา”ติ.}

\csromanheader{2}{12. ทกฺขิณาวิภงฺคสุตฺต (142)}
\proseitem{378}{\csromanfolio{297}เอวเมตํ อานนฺท, ยํ หานนฺท ปุคฺคโล ปุคฺคลํ อาคมฺม พุทฺธํ สรณํ คโต โหติ, ธมฺมํ สรณํ คโต โหติ, สํฆํ สรณํ คโต โหติ.\csromanspacer{} อิมสฺสานนฺท ปุคฺคลสฺส อิมินา ปุคฺคเลน น สุปฺปติการํ วทามิ, ยทิทํ อภิวาทนปจฺจุฏฺฐาน-อญฺชลิกมฺม- สามีจิกมฺมจีวรปิณฺฑปาตเสนาสนคิลานปฺปจฺจยเภสชฺช- ปริกฺขารานุปฺปทาเนน.}

\prose{ยํ หานนฺท ปุคฺคโล ปุคฺคลํ อาคมฺม ปาณาติปาตา ปฏิวิรโต โหติ, อทินฺนาทานา ปฏิวิรโต โหติ, กาเมสุมิจฺฉาจารา ปฏิวิรโต โหติ, มุสาวาทา ปฏิวิรโต โหติ, สุราเมรยมชฺชปมาทฏฺฐานา ปฏิวิรโต โหติ.\csromanspacer{} อิมสฺสานนฺท ปุคฺคลสฺส อิมินา ปุคฺคเลน น สุปฺปติการํ วทามิ, ยทิทํ อภิวาทนปจฺจุฏฺฐาน-อญฺชลิกมฺมสามีจิกมฺมจีวรปิณฺฑปาต- เสนาสนคิลานปฺปจฺจยเภสชฺชปริกฺขารานุปฺปทาเนน.}

\prose{ยํ หานนฺท ปุคฺคโล ปุคฺคลํ อาคมฺม พุทฺเธ อเวจฺจปฺปสาเทน สมนฺนาคโต โหติ.\csromanspacer{} ธมฺเม.\csromanspacer{} สํเฆ.\csromanspacer{} อริยกนฺเตหิ สีเลหิ สมนฺนาคโต โหติ.\csromanspacer{} อิมสฺสานนฺท ปุคฺคลสฺส อิมินา ปุคฺคเลน น สุปฺปติการํ วทามิ, ยทิทํ อภิวาทนปจฺจุฏฺฐาน-อญฺชลิกมฺมสามีจิกมฺมจีวรปิณฺฑปาตเสนาสน- คิลานปฺปจฺจยเภสชฺชปริกฺขารานุปฺปทาเนน.}

\prose{ยํ หานนฺท ปุคฺคโล ปุคฺคลํ อาคมฺม ทุกฺเข นิกฺกงฺโข โหติ, ทุกฺขสมุทเย นิกฺกงฺโข โหติ, ทุกฺขนิโรเธ นิกฺกงฺโข โหติ, ทุกฺขนิโรธคามินิยา ปฏิปทาย นิกฺกงฺโข โหติ.\csromanspacer{} อิมสฺสานนฺท ปุคฺคลสฺส อิมินา ปุคฺคเลน น สุปฺปติการํ วทามิ, ยทิทํ อภิวาทน ปจฺจุฏฺฐาน อญฺชลิกมฺม สามีจิกมฺม จีวรปิณฺฑปาต เสนาสน คิลานปฺปจฺจยเภสชฺชปริกฺขารานุปฺปทาเนน.}

\proseitem{379}{จุทฺทส โข ปนิมานนฺท ปาฏิปุคฺคลิกา ทกฺขิณา.\csromanspacer{} กตมา จุทฺทส, ตถาคเต อรหนฺเต สมฺมาสมฺพุทฺเธ ทานํ เทติ, อยํ ปฐมา ปาฏิปุคฺคลิกา ทกฺขิณา.\csromanspacer{} ปจฺเจกสมฺพุทฺเธ\footnote{ปจฺเจกพุทฺเธ (สี, อิ)} ทานํ เทติ, อยํ ทุติยา ปาฏิปุคฺคลิกา ทกฺขิณา.\csromanspacer{} ตถาคตสาวเก อรหนฺเต ทานํ เทติ, อยํ ตติยา ปาฏิปุคฺคลิกา ทกฺขิณา.\csromanspacer{} อรหตฺตผลสจฺฉิ กิริยาย ปฏิปนฺเน ทานํ เทติ, อยํ จตุตฺถี ปาฏิปุคฺคลิกา ทกฺขิณา.\csromanspacer{} อนาคามิสฺส ทานํ เทติ, อยํ ปญฺจมี ปาฏิปุคฺคลิกา ทกฺขิณา.\csromanspacer{} อนาคามิผลสจฺฉิกิริยาย ปฏิปนฺเน ทานํ เทติ, อยํ ฉฏฺฐี ปาฏิปุคฺคลิกา ทกฺขิณา.\csromanspacer{} สกทาคามิสฺส ทานํ เทติ, อยํ สตฺตมี \csromanfolio{298}ปาฏิปุคฺคลิกา ทกฺขิณา.\csromanspacer{} สกทาคามิผลสจฺฉิกิริยาย ปฏิปนฺเน ทานํ เทติ, อยํ อฏฺฐมี ปาฏิปุคฺคลิกา ทกฺขิณา.\csromanspacer{} โสตาปนฺเน ทานํ เทติ, อยํ นวมี ปาฏิปุคฺคลิกา ทกฺขิณา.\csromanspacer{} โสตาปตฺติผลสจฺฉิกิริยาย ปฏิปนฺเน ทานํ เทติ, อยํ ทสมี ปาฏิปุคฺคลิกา ทกฺขิณา.\csromanspacer{} พาหิรเก กาเมสุ วีตราเค ทานํ เทติ, อยํ เอกาทสมี ปาฏิปุคฺคลิกา ทกฺขิณา.\csromanspacer{} ปุถุชฺชนสีลวนฺเต ทานํ เทติ, อยํ ทฺวาทสมี ปาฏิปุคฺคลิกา ทกฺขิณา.\csromanspacer{} ปุถุชฺชนทุสฺสีเล ทานํ เทติ, อยํ เตรสมี ปาฏิปุคฺคลิกา ทกฺขิณา.\csromanspacer{} ติรจฺฉานคเต ทานํ เทติ, อยํ จุทฺทสมี ปาฏิปุคฺคลิกา ทกฺขิณาติ.}

\prose{ตตฺรานนฺท ติรจฺฉานคเต ทานํ ทตฺวา สตคุณา ทกฺขิณา ปาฏิกงฺขิตพฺพา.\csromanspacer{} ปุถุชฺชนทุสฺสีเล ทานํ ทตฺวา สหสฺสคุณา ทกฺขิณา ปาฏิกงฺขิตพฺพา.\csromanspacer{} ปุถุชฺชนสีลวนฺเต ทานํ ทตฺวา สตสหสฺสคุณา ทกฺขิณา ปาฏิกงฺขิตพฺพา.\csromanspacer{} พาหิรเก กาเมสุ วีตราเค ทานํ ทตฺวา โกฏิสตสหสฺส- คุณา ทกฺขิณา ปาฏิกงฺขิตพฺพา.\csromanspacer{} โสตาปตฺติผลสจฺฉิกิริยาย ปฏิปนฺเน ทานํ ทตฺวา อสงฺเขยฺยา อปฺปเมยฺยา ทกฺขิณา ปาฏิกงฺขิตพฺพา.\csromanspacer{} โก ปน วาโท โสตาปนฺเน, โก ปน วาโท สกทาคามิผลสจฺฉิกิริยาย ปฏิปนฺเน, โก ปน วาโท สกทาคามิสฺส, โก ปน วาโท อนาคามิผลสจฺฉิกิริยาย ปฏิปนฺเน, โก ปน วาโท อนาคามิสฺส, โก ปน วาโท อรหตฺตผลสจฺฉิกิริยาย ปฏิปนฺเน, โก ปน วาโท อรหนฺเต, โก ปน วาโท ปจฺเจกสมฺพุทฺเธ, โก ปน วาโท ตถาคเต อรหนฺเต สมฺมาสมฺพุทฺเธ.}

\proseitem{380}{สตฺต โข ปนิมานนฺท สํฆคตา ทกฺขิณา.\csromanspacer{} กตมา สตฺต, พุทฺธปฺปมุเข อุภโตสํเฆ ทานํ เทติ, อยํ ปฐมา สํฆคตา ทกฺขิณา.\csromanspacer{} ตถาคเต ปรินิพฺพุเต อุภโตสํเฆ ทานํ เทติ, อยํ ทุติยา สํฆคตา ทกฺขิณา.\csromanspacer{} ภิกฺขุสํเฆ ทานํ เทติ, อยํ ตติยา สํฆคตา ทกฺขิณา.\csromanspacer{} ภิกฺขุนิสํเฆ ทานํ เทติ, อยํ จตุตฺถี สํฆคตา ทกฺขิณา.\csromanspacer{} เอตฺตกา เม ภิกฺขู จ ภิกฺขุนิโย จ สํฆโต อุทฺทิสฺสถาติ ทานํ เทติ, อยํ ปญฺจมี สํฆคตา ทกฺขิณา.\csromanspacer{} เอตฺตกา เม ภิกฺขู สํฆโต อุทฺทิสฺสถาติ ทานํ เทติ, อยํ ฉฏฺฐี สํฆคตา ทกฺขิณา.\csromanspacer{} เอตฺตกา เม ภิกฺขุนิโย สํฆโต อุทฺทิสฺสถาติ ทานํ เทติ, อยํ สตฺตมี สํฆคตา ทกฺขิณา.}

\csromanheader{2}{12. ทกฺขิณาวิภงฺคสุตฺต (142)}
\prose{\csromanfolio{299}ภวิสฺสนฺติ โข ปนานนฺท อนาคตมทฺธานํ โคตฺรภุโน กาสาวกณฺฐา ทุสฺสีลา ปาปธมฺมา, เตสุ ทุสฺสีเลสุ สํฆํ อุทฺทิสฺส ทานํ ทสฺสนฺติ, ตทาปาหํ อานนฺท สํฆคตํ ทกฺขิณํ อสงฺเขยฺยํ อปฺปเมยฺยํ วทามิ, น เตฺววาหํ อานนฺท เกนจิ ปริยาเยน สํฆคตาย ทกฺขิณาย ปาฏิปุคฺคลิกํ ทานํ มหปฺผลตรํ วทามิ.}

\proseitem{381}{จตสฺโส โข อิมา อานนฺท ทกฺขิณา วิสุทฺธิโย.\csromanspacer{} กตมา จตสฺโส, อตฺถานนฺท ทกฺขิณา ทายกโต วิสุชฺฌติ โน ปฏิคฺคาหกโต, อตฺถานนฺท ทกฺขิณา ปฏิคฺคาหกโต วิสุชฺฌติ โน ทายกโต, อตฺถานนฺท ทกฺขิณา เนว ทายกโต วิสุชฺฌติ โน ปฏิคฺคาหกโต, อตฺถานนฺท ทกฺขิณา ทายกโต เจว วิสุชฺฌติ ปฏิคฺคาหกโต จ.}

\prose{กถญฺจานนฺท ทกฺขิณา ทายกโต วิสุชฺฌติ โน ปฏิคฺคาหกโต.\csromanspacer{} อิธานนฺท ทายโก โหติ สีลวา กลฺยาณธมฺโม, ปฏิคฺคาหกา โหนฺติ ทุสฺสีลา ปาปธมฺมา.\csromanspacer{} เอวํ โข อานนฺท ทกฺขิณา ทายกโต วิสุชฺฌติ โน ปฏิคฺคาหกโต.}

\prose{กถญฺจานนฺท ทกฺขิณา ปฏิคฺคาหกโต วิสุชฺฌติ โน ทายกโต.\csromanspacer{} อิธานนฺท ทายโก โหติ ทุสฺสีโล ปาปธมฺโม, ปฏิคฺคาหกา โหนฺติ สีลวนฺโต\footnote{สีลวนฺตา (สี)} กลฺยาณธมฺมา.\csromanspacer{} เอวํ โข อานนฺท ทกฺขิณา ปฏิคฺคาหกโต วิสุชฺฌติ โน ทายกโต.}

\prose{กถญฺจานนฺท ทกฺขิณา เนว ทายกโต วิสุชฺฌติ โน ปฏิคฺคาหกโต.\csromanspacer{} อิธานนฺท ทายโก จ โหติ ทุสฺสีโล ปาปธมฺโม, ปฏิคฺคาหกา จ โหนฺติ ทุสฺสีลา ปาปธมฺมา.\csromanspacer{} เอวํ โข อานนฺท ทกฺขิณา เนว ทายกโต วิสุชฺฌติ โน ปฏิคฺคาหกโต.}

\prose{กถญฺจานนฺท ทกฺขิณา ทายกโต เจว วิสุชฺฌติ ปฏิคฺคาหกโต จ.\csromanspacer{} อิธานนฺท ทายโก จ โหติ สีลวา กลฺยาณธมฺโม, ปฏิคฺคาหกา จ โหนฺติ สีลวนฺโต กลฺยาณธมฺมา.\csromanspacer{} เอวํ โข อานนฺท ทกฺขิณา ทายกโต เจว วิสุชฺฌติ ปฏิคฺคาหกโต จ.\csromanspacer{} อิมา โข อานนฺท จตสฺโส ทกฺขิณา วิสุทฺธิโยติ.}

\prose{อิทมโวจ ภควา, อิทํ วตฺวาน สุคโต อถาปรํ เอตทโวจ สตฺถา\csromandash{}}

\csromanmatikaanchor{mtk.50}
\csromantocmarkref{subsection}{อุทฺทานคาถา}{mtk.50}
\bookmark[dest={mtk.50},level=2]{อุทฺทานคาถา}
\csromanhangingbegin
\hangingheaditem{382}{\csromanfolio{300}“โย สีลวา ทุสฺสีเลสุ ททาติ ทานํ,}
\hangingbody{ธมฺเมน ลทฺธํ1 สุปสนฺนจิตฺโต.}
\hangingbody{อภิสทฺทหํ กมฺมผลํ อุฬารํ,}
\hangingbody{ทา ทกฺขิณา ทายกโต วิสุชฺฌติ.}
\hangingbody{โย ทุสฺสีโล สีลวนฺเตสุ ททาติ ทานํ,}
\hangingbody{อธมฺเมน ลทฺธํ อปฺปสนฺนจิตฺโต.}
\hangingbody{อนภิสทฺทหํ กมฺมผลํ อุฬารํ,}
\hangingbody{สา ทกฺขิณา ปฏิคฺคาหกโต วิสุชฺฌติ.}
\hangingbody{โย ทุสฺสีโล ทุสฺสีเลสุ ททาติ ทานํ,}
\hangingbody{อธมฺเมน ลทฺธํ อปฺปสนฺนจิตฺโต.}
\hangingbody{อนภิสทฺทหํ กมฺมผลํ อุฬารํ,}
\hangingbody{น ตํ ทานํ วิปุลปฺผลนฺติ พฺรูมิ2.}
\hangingbody{โย สีลวา สีลวนฺเตสุ ททาติ ทานํ,}
\hangingbody{ธมฺเมน ลทฺธํ สุปสนฺนจิตฺโต.}
\hangingbody{อภิสทฺทหํ กมฺมผลํ อุฬารํ,}
\hangingbody{ตํ เว ทานํ วิปุลปฺผลนฺติ พฺรูมิ.}
\hangingbody{โย วีตราโค วีตราเคสุ ททาติ ทานํ,}
\hangingbody{ธมฺเมน ลทฺธํ สุปสนฺนจิตฺโต.}
\hangingbody{อภิสทฺทหํ กมฺมผลํ อุฬารํ,}
\hangingbody{ตํ เว ทานํ อามิสทานานมคฺคนฺ”ติ3.}
\csromanhangingend

\nitthitam{ทกฺขิณาวิภงฺคสุตฺตํ นิฏฺฐิตํ ทฺวาทสมํ.}
\nitthitam{\textbf{วิภงฺควคฺโค นิฏฺฐิโต จตุตฺโถ}}
\titlehead{\textbf{ตสฺสุทฺทานํ}}
\csromangathameasure{ภทฺเทกานนฺทกจฺจาน, โลมสกงฺคิยาสุโภ. \\ มหากมฺมสฬายตนวิภงฺคา, อุทฺเทส-อรณา ธาตุ สจฺจํ. \\ ทกฺขิณาวิภงฺคสุตฺตนฺติ.}
\csromangathagroup{\csromangathastanza{ภทฺเทกานนฺทกจฺจาน, โลมสกงฺคิยาสุโภ. \\ มหากมฺมสฬายตนวิภงฺคา, อุทฺเทส-อรณา ธาตุ สจฺจํ. \\ ทกฺขิณาวิภงฺคสุตฺตนฺติ.}}

\csromanmatikaanchor{mtk.51}
\csromantocmarknopage{chapter}{5. สฬายตนวคฺค}{mtk.51}
\bookmark[dest={mtk.51},level=0]{5. สฬายตนวคฺค}
\chapterheadpageread{5. สฬายตนวคฺค}
\csromanmatikaanchor{mtk.52}
\csromantocmarkref{subsection}{1. อนาถปิณฺฑิโกวาทสุตฺต}{mtk.52}
\bookmark[dest={mtk.52},level=2]{1. อนาถปิณฺฑิโกวาทสุตฺต}
\csromanheader{2}{1. อนาถปิณฺฑิโกวาทสุตฺต}
\proseitem{383}{\csromanfolio{301}เอวํ เม สุตํ\csromandash{}เอกํ สมยํ ภควา สาวตฺถิยํ วิหรติ เชตวเน อนาถปิณฺฑิกสฺส อาราเม.\csromanspacer{} เตน โข ปน สมเยน อนาถปิณฺฑิโก คหปติ อาพาธิโก โหติ ทุกฺขิโต พาฬฺหคิลาโน.\csromanspacer{} อถ โข อนาถปิณฺฑิโก คหปติ อญฺญตรํ ปุริสํ อามนฺเตสิ “เอหิ ตฺวํ อมฺโภ ปุริส เยน ภควา เตนุปสงฺกม, อุปสงฺกมิตฺวา มม วจเนน ภควโต ปาเท สิรสา วนฺทาหิ\footnote{วนฺทาหิ เอวญฺจ วเทหิ (สพฺพตฺถ) อญฺญสุตฺเตสุ ปน นตฺถิ.} ‘อนาถปิณฺฑิโก ภนฺเต คหปติ อาพาธิโก ทุกฺขิโต พาฬฺหคิลาโน, โส ภควโต ปาเท สิรสา วนฺทตี’ติ.\csromanspacer{} เยน จายสฺมา สาริปุตฺโต เตนุปสงฺกม, อุปสงฺกมิตฺวา มม วจเนน อายสฺมโต สาริปุตฺตสฺส ปาเท สิรสา วนฺทาหิ1 ‘อนาถปิณฺฑิโก ภนฺเต คหปติ อาพาธิโก ทุกฺขิโต พาฬฺหคิลาโน, โส อายสฺมโต สาริปุตฺตสฺส ปาเท สิรสา วนฺทตี’ติ.\csromanspacer{} เอวญฺจ วเทหิ ‘สาธุ กิร ภนฺเต อายสฺมา สาริปุตฺโต เยน อนาถปิณฺฑิกสฺส คหปติสฺส นิเวสนํ เตนุปสงฺกมตุ อนุกมฺปํ อุปาทายา’ติ”.}

\prose{“เอวํ ภนฺเต”ติ โข โส ปุริโส อนาถปิณฺฑิกสฺส คหปติสฺส ปฏิสฺสุตฺวา เยน ภควา เตนุปสงฺกมิ, อุปสงฺกมิตฺวา ภควนฺตํ อภิวาเทตฺวา เอกมนฺตํ นิสีทิ, เอกมนฺตํ นิสินฺโน โข โส ปุริโส ภควนฺตํ เอตทโวจ “อนาถปิณฺฑิโก ภนฺเต คหปติ อาพาธิโก ทุกฺขิโต พาฬฺหคิลาโน, โส ภควโต ปาเท สิรสา วนฺทตี”ติ.\csromanspacer{} เยน จายสฺมา สาริปุตฺโต เตนุปสงฺกมิ, อุปสงฺกมิตฺวา อายสฺมนฺตํ สาริปุตฺตํ อภิวาเทตฺวา เอกมนฺตํ นิสีทิ, เอกมนฺตํ นิสินฺโน โข โส ปุริโส อายสฺมนฺตํ สาริปุตฺตํ เอตทโวจ “อนาถปิณฺฑิโก ภนฺเต คหปติ อาพาธิโก ทุกฺขิโต พาฬฺหคิลาโน, โส อายสฺมโต สาริปุตฺตสฺส ปาเท สิรสา วนฺทติ, เอวญฺจ วเทติ ‘สาธุ กิร ภนฺเต อายสฺมา สาริปุตฺโต เยน อนาถปิณฺฑิกสฺส คหปติสฺส นิเวสนํ เตนุปสงฺกมตุ \csromanfolio{302}อนุกมฺปํ อุปาทายา’ติ”.\csromanspacer{} อธิวาเสสิ โข อายสฺมา สาริปุตฺโต ตุณฺหีภาเวน.}

\proseitem{384}{อถ โข อายสฺมา สาริปุตฺโต นิวาเสตฺวา ปตฺตจีวรมาทาย อายสฺมตา อานนฺเทน ปจฺฉาสมเณน เยน อนาถปิณฺฑิกสฺส คหปติสฺส นิเวสนํ เตนุปสงฺกมิ, อุปสงฺกมิตฺวา ปญฺญตฺเต อาสเน นิสีทิ, นิสชฺช โข อายสฺมา สาริปุตฺโต อนาถปิณฺฑิกํ คหปติํ เอตทโวจ “กจฺจิ เต คหปติ ขมนียํ, กจฺจิ ยาปนียํ, กจฺจิ เต ทุกฺขา เวทนา ปฏิกฺกมนฺติ โน อภิกฺกมนฺติ, ปฏิกฺกโมสานํ ปญฺญายติ โน อภิกฺกโม”ติ.}

\prose{น เม ภนฺเต สาริปุตฺต ขมนียํ, น ยาปนียํ, พาฬฺหา เม ทุกฺขา เวทนา อภิกฺกมนฺติ โน ปฏิกฺกมนฺติ, อภิกฺกโมสานํ ปญฺญายติ โน ปฏิกฺกโม.\csromanspacer{} เสยฺยถาปิ ภนฺเต สาริปุตฺต พลวา ปุริโส ติเณฺหน สิขเรน มุทฺธนิ\footnote{มุทฺธานํ (สี, สฺยา, กํ, อิ)} อภิมตฺเถยฺย\footnote{อภิมนฺเถยฺย (สี, อิ) 3. โอหนนฺติ (สฺยา, กํ)}.\csromanspacer{} เอวเมว โข เม ภนฺเต สาริปุตฺต อธิมตฺตา วาตา มุทฺธนิ1 อูหนนฺติ3.\csromanspacer{} น เม ภนฺเต สาริปุตฺต ขมนียํ, น ยาปนียํ, พาฬฺหา เม ทุกฺขา เวทนา อภิกฺกมนฺติ โน ปฏิกฺกมนฺติ, อภิกฺกโมสานํ ปญฺญายติ โน ปฏิกฺกโม.\csromanspacer{} เสยฺยถาปิ ภนฺเต สาริปุตฺต พลวา ปุริโส ทเฬฺหน วรตฺตขณฺเฑน สีเส สีสเวฐํ ทเทยฺย.\csromanspacer{} เอวเมว โข เม ภนฺเต สาริปุตฺต อธิมตฺตา สีเส สีสเวทนา\footnote{อธิมตฺตา วาตา สีสํ ปริกนฺตนฺติ (สี, สฺยา, กํ)}.\csromanspacer{} น เม ภนฺเต สาริปุตฺต ขมนียํ, น ยาปนียํ, พาฬฺหา เม ทุกฺขา เวทนา อภิกฺกมนฺติ โน ปฏิกฺกมนฺติ, อภิกฺกโมสานํ ปญฺญายติ โน ปฏิกฺกโม.\csromanspacer{} เสยฺยถาปิ ภนฺเต สาริปุตฺต ทกฺโข โคฆาตโก วา โคฆาตกนฺเตวาสี วา ติเณฺหน โควิกนฺตเนน กุจฺฉิํ ปริกนฺเตยฺย.\csromanspacer{} เอวเมว โข เม ภนฺเต สาริปุตฺต อธิมตฺตา วาตา กุจฺฉิํ ปริกนฺตนฺติ.\csromanspacer{} น เม ภนฺเต สาริปุตฺต ขมนียํ, น ยาปนียํ, พาฬฺหา เม ทุกฺขา เวทนา อภิกฺกมนฺติ โน ปฏิกฺกมนฺติ, อภิกฺกโมสานํ ปญฺญายติ โน ปฏิกฺกโม.\csromanspacer{} เสยฺยถาปิ ภนฺเต สาริปุตฺต เทฺว พลวนฺโต ปุริสา ทุพฺพลตรํ ปุริสํ นานาพาหาสุ คเหตฺวา องฺคารกาสุยา สนฺตาเปยฺยุํ สมฺปริตาเปยฺยุํ.\csromanspacer{} เอวเมว โข เม ภนฺเต สาริปุตฺต อธิมตฺโต กายสฺมิํ ฑาโห.\csromanspacer{} น เม ภนฺเต สาริปุตฺต ขมนียํ, น ยาปนียํ, พาฬฺหา เม ทุกฺขา}

\csromanheader{2}{1. อนาถปิณฺฑิโกวาทสุตฺต (143)}
\prose{\csromanfolio{303}เวทนา อภิกฺกมนฺติ โน ปฏิกฺกมนฺติ, อภิกฺกโมสานํ ปญฺญายติ โน ปฏิกฺกโมติ.}

\proseitem{385}{ตสฺมาติห เต คหปติ เอวํ สิกฺขิตพฺพํ “น จกฺขุํ อุปาทิยิสฺสามิ, น จ เม จกฺขุนิสฺสิตํ วิญฺญาณํ ภวิสฺสตี”ติ.\csromanspacer{} เอวํ หิ เต คหปติ สิกฺขิตพฺพํ.}

\prose{ตสฺมาติห เต คหปติ เอวํ สิกฺขิตพฺพํ “น โสตํ อุปาทิยิสฺสามิ, น จ เม โสตนิสฺสิตํ วิญฺญาณํ ภวิสฺสตี”ติ, เอวํ หิ เต คหปติ สิกฺขิตพฺพํ.\csromanspacer{} ตสฺมาติห เต คหปติ เอวํ สิกฺขิตพฺพํ “น ฆานํ อุปาทิยิสฺสามิ, น จ เม ฆานนิสฺสิตํ วิญฺญาณํ ภวิสฺสตี”ติ, เอวํ หิ เต คหปติ สิกฺขิตพฺพํ.\csromanspacer{} ตสฺมาติห เต คหปติ เอวํ สิกฺขิตพฺพํ “น ชิวฺหํ อุปาทิยิสฺสามิ, น จ เม ชิวฺหานิสฺสิตํ วิญฺญาณํ ภวิสฺสตี”ติ, เอวํ หิ เต คหปติ สิกฺขิตพฺพํ.\csromanspacer{} ตสฺมาติห เต คหปติ เอวํ สิกฺขิตพฺพํ “น กายํ อุปาทิยิสฺสามิ, น จ เม กายนิสฺสิตํ วิญฺญาณํ ภวิสฺสตี”ติ, เอวํ หิ เต คหปติ สิกฺขิตพฺพํ.\csromanspacer{} ตสฺมาติห เต คหปติ เอวํ สิกฺขิตพฺพํ “น มนํ อุปาทิยิสฺสามิ, น จ เม มโนนิสฺสิตํ วิญฺญาณํ ภวิสฺสตี”ติ, เอวํ หิ เต คหปติ สิกฺขิตพฺพํ. (1)}

\prose{ตสฺมาติห เต คหปติ เอวํ สิกฺขิตพฺพํ “น รูปํ อุปาทิยิสฺสามิ, น จ เม รูปนิสฺสิตํ วิญฺญาณํ ภวิสฺสตี”ติ.\csromanspacer{} เอวํ หิ เต คหปติ สิกฺขิตพฺพํ.\csromanspacer{} ตสฺมาติห เต คหปติ เอวํ สิกฺขิตพฺพํ “น สทฺทํ อุปาทิยิสฺสามิ ฯเปฯ น คนฺธํ อุปาทิยิสฺสามิ.\csromanspacer{} น รสํ อุปาทิยิสฺสามิ.\csromanspacer{} น โผฏฺฐพฺพํ อุปาทิยิสฺสามิ.\csromanspacer{} น ธมฺมํ อุปาทิยิสฺสามิ, น จ เม ธมฺม- นิสฺสิตํ วิญฺญาณํ ภวิสฺสตี”ติ.\csromanspacer{} เอวํ หิ เต คหปติ สิกฺขิตพฺพํ. (2)}

\prose{ตสฺมาติห เต คหปติ เอวํ สิกฺขิตพฺพํ “น จกฺขุวิญฺญาณํ อุปาทิยิสฺสามิ, น จ เม จกฺขุวิญฺญาณนิสฺสิตํ วิญฺญาณํ ภวิสฺสตี”ติ, เอวํ หิ เต คหปติ สิกฺขิตพฺพํ.\csromanspacer{} ตสฺมาติห เต คหปติ เอวํ สิกฺขิตพฺพํ “น โสตวิญฺญาณํ อุปาทิยิสฺสามิ.\csromanspacer{} น ฆานวิญฺญาณํ อุปาทิยิสฺสามิ.\csromanspacer{} น ชิวฺหาวิญฺญาณํ อุปาทิยิสฺสามิ.\csromanspacer{} น กายวิญฺญาณํ อุปาทิยิสฺสามิ.\csromanspacer{} น มโนวิญฺญาณํ อุปาทิยิสฺสามิ, น จ เม มโนวิญฺญาณนิสฺสิตํ วิญฺญาณํ ภวิสฺสตี”ติ, เอวํ หิ เต คหปติ สิกฺขิตพฺพํ. (3)}

\prose{\csromanfolio{304}ตสฺมาติห เต คหปติ เอวํ สิกฺขิตพฺพํ “น จกฺขุสมฺผสฺสํ อุปาทิยิสฺสามิ, น จ เม จกฺขุสมฺผสฺสนิสฺสิตํ วิญฺญาณํ ภวิสฺสตี”ติ, เอวํ หิ เต คหปติ สิกฺขิตพฺพํ.\csromanspacer{} ตสฺมาติห เต คหปติ เอวํ สิกฺขิตพฺพํ “น โสตสมฺผสฺสํ อุปาทิยิสฺสามิ.\csromanspacer{} น ฆานสมฺผสฺสํ อุปาทิยิสฺสามิ.\csromanspacer{} น ชิวฺหาสมฺผสฺสํ อุปาทิยิสฺสามิ.\csromanspacer{} น กายสมฺผสฺส อุปาทิยิสฺสามิ.\csromanspacer{} น มโนสมฺผสฺสํ อุปาทิยิสฺสามิ, น จ เม มโนสมฺผสฺสนิสฺสิตํ วิญฺญาณํ ภวิสฺสตี”ติ, เอวํ หิ เต คหปติ สิกฺขิตพฺพํ.(4)}

\prose{ตสฺมาติห เต คหปติ เอวํ สิกฺขิตพฺพํ “น จกฺขุสมฺผสฺสชํ เวทนํ อุปาทิยิสฺสามิ, น จ เม จกฺขุสมฺผสฺสชาเวทนานิสฺสิตํ วิญฺญาณํ ภวิสฺสตี”ติ, เอวํ หิ เต คหปติ สิกฺขิตพฺพํ.\csromanspacer{} ตสฺมาติห เต คหปติ เอวํ สิกฺขิตพฺพํ “น โสตสมฺผสฺสชํ เวทนํ อุปาทิยิสฺสามิ.\csromanspacer{} น ฆานสมฺผสฺสชํ เวทนํ อุปาทิยิสฺสามิ.\csromanspacer{} น ชิวฺหาสมฺผสฺสชํ เวทนํ อุปาทิยิสฺสามิ.\csromanspacer{} น กายสมฺผสฺสชํ เวทนํ อุปาทิยิสฺสามิ.\csromanspacer{} น มโนสมฺผสมฺผสฺสชํ เวทนํ อุปาทิยิสฺสามิ, น จ เม มโนสมฺผสฺสชาเวทนานิสฺสิตํ วิญฺญาณํ ภวิสฺสตี”ติ, เอวํ หิ เต คหปติ สิกฺขิตพฺพํ. (5)}

\proseitem{386}{ตสฺมาติห เต คหปติ เอวํ สิกฺขิตพฺพํ “น ปถวีธาตุํ อุปาทิยิสฺสามิ, น จ เม ปถวีธาตุนิสฺสิตํ วิญฺญาณํ ภวิสฺสตี”ติ, เอวํ หิ เต คหปติ สิกฺขิตพฺพํ.\csromanspacer{} ตสฺมาติห เต คหปติ เอวํ สิกฺขิตพฺพํ “น อาโปธาตุํ อุปาทิยิสฺสามิ.\csromanspacer{} น เต โชธาตุํ อุปาทิยิสฺสามิ.\csromanspacer{} น วาโยธาตุํ อุปาทิยิสฺสามิ.\csromanspacer{} น อากาสธาตุํ อุปาทิยิสฺสามิ.\csromanspacer{} น วิญฺญาณธาตุํ อุปาทิยิสฺสามิ, น จ เม วิญฺญาณธาตุนิสฺสิตํ วิญฺญาณํ ภวิสฺสตี”ติ, เอวํ หิ เต คหปติ สิกฺขิตพฺพํ. (6)}

\prose{ตสฺมาติห เต คหปติ เอวํ สิกฺขิตพฺพํ “น รูปํ อุปาทิยิสฺสามิ, น จ เม รูปนิสฺสิตํ วิญฺญาณํ ภวิสฺสตี”ติ, เอวํ หิ เต คหปติ สิกฺขิตพฺพํ.\csromanspacer{} ตสฺมาติห เต คหปติ เอวํ สิกฺขิตพฺพํ “น เวทนํ อุปาทิยิสฺสามิ.\csromanspacer{} น สญฺญํ อุปาทิยิสฺสามิ.\csromanspacer{} น สงฺขาเร อุปาทิยิสฺสามิ.\csromanspacer{} น วิญฺญาณํ อุปาทิยิสฺสามิ, น จ เม วิญฺญาณนิสฺสิตํ วิญฺญาณํ ภวิสฺสตี”ติ, เอวํ หิ เต คหปติ สิกฺขิตพฺพํ. (7)}

\prose{ตสฺมาติห เต คหปติ เอวํ สิกฺขิตพฺพํ “น อากาสานญฺจายตนํ อุปาทิยิสฺสามิ, น จ เม อากาสานญฺจายตนนิสฺสิตํ วิญฺญาณํ}

\csromanheader{2}{1. อนาถปิณฺฑิโกวาทสุตฺต (143)}
\prose{\csromanfolio{305}ภวิสฺสตี”ติ, เอวํ หิ เต คหปติ สิกฺขิตพฺพํ.\csromanspacer{} ตสฺมาติห เต คหปติ เอวํ สิกฺขิตพฺพํ “น วิญฺญาณญฺจายตนํ อุปาทิยิสฺสามิ.\csromanspacer{} น อากิญฺจญฺญายตนํ อุปาทิยิสฺสามิ.\csromanspacer{} น เนวสญฺญานาสญฺญายตนํ อุปาทิยิสฺสามิ, น จ เม เนวสญฺญานาสญฺญายตนนิสฺสิตํ วิญฺญาณํ ภวิสฺสตี”ติ, เอวํ หิ เต คหปติ สิกฺขิตพฺพํ. (8)}

\prose{ตสฺมาติห เต คหปติ เอวํ สิกฺขิตพฺพํ “น อิธ โลกํ อุปาทิยิสฺสามิ, น จ เม อิธ โลกนิสฺสิตํ วิญฺญาณํ ภวิสฺสตี”ติ, เอวํ หิ เต คหปติ สิกฺขิตพฺพํ.\csromanspacer{} ตสฺมาติห เต คหปติ เอวํ สิกฺขิตพฺพํ “น ปรโลกํ อุปาทิยิสฺสามิ, น จ เม ปรโลกนิสฺสิตํ วิญฺญาณํ ภวิสฺสตี”ติ, เอวํ หิ เต คหปติ สิกฺขิตพฺพํ.\csromanspacer{} ตสฺมาติห เต คหปติ เอวํ สิกฺขิตพฺพํ “ยมฺปิ เม ทิฏฺฐํ สุตํ มุตํ วิญฺญาตํ ปตฺตํ ปริเยสิตํ อนุปริเยสิตํ อนุจริตํ มนสา, ตมฺปิ น อุปาทิยิสฺสามิ, น จ เม ตํนิสฺสิตํ วิญฺญาณํ ภวิสฺสตี”ติ, เอวํ หิ เต คหปติ สิกฺขิตพฺพนฺติ. (9)}

\proseitem{387}{เอวํ วุตฺเต อนาถปิณฺฑิโก คหปติ ปโรทิ, อสฺสูนิ ปวตฺเตสิ.\csromanspacer{} อถ โข อายสฺมา อานนฺโท อนาถปิณฺฑิกํ คหปติํ เอตทโวจ “โอลียสิ โข ตฺวํ คหปติ, สํสีทสิ โข ตฺวํ คหปตี”ติ.\csromanspacer{} นาหํ ภนฺเต อานนฺท โอลียามิ นปิ สํสีทามิ, อปิ จ เม ทีฆรตฺตํ สตฺถา ปยิรุปาสิโต, มโนภาวนียา จ ภิกฺขู, น จ เม เอวรูปี ธมฺมี กถา สุตปุพฺพาติ.\csromanspacer{} น โข คหปติ คิหีนํ โอทาตวสนานํ เอวรูปี ธมฺมี กถา ปฏิภาติ, ปพฺพชิตานํ โข คหปติ เอวรูปี ธมฺมี กถา ปฏิภาตีติ.\csromanspacer{} เตน หิ ภนฺเต สาริปุตฺต คิหีนมฺปิ โอทาตวสนานํ เอวรูปี ธมฺมี กถา ปฏิภาตุ, สนฺติ หิ ภนฺเต กุลปุตฺตา อปฺปรชกฺขชาติกา, อสฺสวนตา ธมฺมสฺส ปริหายนฺติ, ภวิสฺสนฺติ ธมฺมสฺส อญฺญาตาโรติ.}

\prose{อถ โข อายสฺมา จ สาริปุตฺโต อายสฺมา จ อานนฺโท อนาถปิณฺฑิกํ คหปติํ อิมินา โอวาเทน โอวทิตฺวา อุฏฺฐายาสนา ปกฺกมิํสุ.\csromanspacer{} อถ โข อนาถปิณฺฑิโก คหปติ อจิรปกฺกนฺเต อายสฺมนฺเต จ สาริปุตฺเต อายสฺมนฺเต จ อานนฺเท กาลมกาสิ, ตุสิตํ กายํ อุปปชฺชิ.\csromanspacer{} อถ โข อนาถปิณฺฑิโก เทวปุตฺโต อภิกฺกนฺตาย รตฺติยา อภิกฺกนฺตวณฺโณ เกวลกปฺปํ เชตวนํ โอภาเสตฺวา เยน ภควา เตนุปสงฺกมิ, อุปสงฺกมิตฺวา ภควนฺตํ อภิวาเทตฺวา เอกมนฺตํ อฏฺฐาสิ, \csromanfolio{306}เอกมนฺตํ ฐิโต โข อนาถปิณฺฑิโก เทวปุตฺโต ภควนฺตํ คาถาหิ อชฺฌภาสิ\csromandash{}}

\csromangathameasure{“อิทํ หิ ตํ เชตวนํ, อิสิสํฆนิเสวิตํ. \\ อาวุตฺถํ ธมฺมราเชน, ปีติสญฺชนนํ มม. \\ กมฺมํ วิชฺชา จ ธมฺโม จ, สีลํ ชีวิตมุตฺตมํ. \\ เอเตน มจฺจา สุชฺฌนฺติ, น โคตฺเตน ธเนน วา. \\ ตสฺมา หิ ปณฺฑิโต โปโส, สมฺปสฺสํ อตฺถมตฺตโน. \\ โยนิโส วิจิเน ธมฺมํ, เอวํ ตตฺถ วิสุชฺฌติ. \\ สาริปุตฺโตว ปญฺญาย, สีเลน อุปสเมน จ. \\ โยปิ ปารงฺคโต ภิกฺขุ, เอตาวปรโม สิยา”ติ.}
\csromangathagroup{\csromangathastanza{“อิทํ หิ ตํ เชตวนํ, อิสิสํฆนิเสวิตํ. \\ อาวุตฺถํ ธมฺมราเชน, ปีติสญฺชนนํ มม.}\csromangathastanza{กมฺมํ วิชฺชา จ ธมฺโม จ, สีลํ ชีวิตมุตฺตมํ. \\ เอเตน มจฺจา สุชฺฌนฺติ, น โคตฺเตน ธเนน วา.}\csromangathastanza{ตสฺมา หิ ปณฺฑิโต โปโส, สมฺปสฺสํ อตฺถมตฺตโน. \\ โยนิโส วิจิเน ธมฺมํ, เอวํ ตตฺถ วิสุชฺฌติ.}\csromangathastanza{สาริปุตฺโตว ปญฺญาย, สีเลน อุปสเมน จ. \\ โยปิ ปารงฺคโต ภิกฺขุ, เอตาวปรโม สิยา”ติ.}}

\prose{อิทมโวจ อนาถปิณฺฑิโก เทวปุตฺโต.\csromanspacer{} สมนุญฺโญ สตฺถา อโหสิ.\csromanspacer{} อถ โข อนาถปิณฺฑิโก เทวปุตฺโต “สมนุญฺโญ เม สตฺถา”ติ ภควนฺตํ อภิวาเทตฺวา ปทกฺขิณํ กตฺวา ตตฺเถวนฺตรธายิ.}

\proseitem{388}{อถ โข ภควา ตสฺสา รตฺติยา อจฺจเยน ภิกฺขู อามนฺเตสิ\csromandash{}อิมํ ภิกฺขเว รตฺติํ อญฺญตโร เทวปุตฺโต อภิกฺกนฺตาย รตฺติยา อภิกฺกนฺตวณฺโณ เกวลกปฺปํ เชตวนํ โอภาเสตฺวา เยนาหํ เตนุปสงฺกมิ, อุปสงฺกมิตฺวา มํ อภิวาเทตฺวา เอกมนฺตํ อฏฺฐาสิ, เอกมนฺตํ ฐิโต โข โส เทวปุตฺโต มํ คาถาหิ อชฺฌภาสิ\csromandash{}}

\csromangathameasure{“อิทํ หิ ตํ เชตวนํ, อิสิสํฆนิเสวิตํ. \\ อาวุตฺถํ ธมฺมราเชน, ปีติสญฺชนนํ มม. \\ กมฺมํ วิชฺชา จ ธมฺโม จ, สีลํ ชีวิตมุตฺตมํ. \\ เอเตน มจฺจา สุชฺฌนฺติ, น โคตฺเตน ธเนน วา. \\ ตสฺมา หิ ปณฺฑิโต โปโส, สมฺปสฺสํ อตฺถมตฺตโน. \\ โยนิโส วิจิเน ธมฺมํ, เอวํ ตตฺถ วิสุชฺฌติ. \\ สาริปุตฺโตว ปญฺญาย, สีเลน อุปสเมน จ. \\ โยปิ ปารงฺคโต ภิกฺขุ, เอตาวปรโม สิยา”ติ.}
\csromangathagroup{\csromangathastanza{“อิทํ หิ ตํ เชตวนํ, อิสิสํฆนิเสวิตํ. \\ อาวุตฺถํ ธมฺมราเชน, ปีติสญฺชนนํ มม.}\csromangathastanza{กมฺมํ วิชฺชา จ ธมฺโม จ, สีลํ ชีวิตมุตฺตมํ. \\ เอเตน มจฺจา สุชฺฌนฺติ, น โคตฺเตน ธเนน วา.}\csromangathastanza{ตสฺมา หิ ปณฺฑิโต โปโส, สมฺปสฺสํ อตฺถมตฺตโน. \\ โยนิโส วิจิเน ธมฺมํ, เอวํ ตตฺถ วิสุชฺฌติ.}\csromangathastanza{สาริปุตฺโตว ปญฺญาย, สีเลน อุปสเมน จ. \\ โยปิ ปารงฺคโต ภิกฺขุ, เอตาวปรโม สิยา”ติ.}}

\prose{อิทมโวจ ภิกฺขเว โส เทวปุตฺโต “สมนุญฺโญ เม สตฺถา”ติ มํ อภิวาเทตฺวา ปทกฺขิณํ กตฺวา ตตฺเถวนฺตรธายีติ.}

\csromanheader{2}{2. ฉนฺโนวาทสุตฺต (144)}
\prose{\csromanfolio{307}เอวํ วุตฺเต อายสฺมา อานนฺโท ภควนฺตํ เอตทโวจ “โส หิ นูน โส ภนฺเต อนาถปิณฺฑิโก เทวปุตฺโต ภวิสฺสติ, อนาถปิณฺฑิโก ภนฺเต คหปติ อายสฺมนฺเต สาริปุตฺเต อภิปฺปสนฺโน อโหสี”ติ.\csromanspacer{} สาธุ สาธุ อานนฺท, ยาวตกํ โข อานนฺท ตกฺกาย ปตฺตพฺพํ, อนุปฺปตฺตํ ตํ ตยา, อนาถปิณฺฑิโก โส อานนฺท เทวปุตฺโตติ.}

\prose{อิทมโวจ ภควา.\csromanspacer{} อตฺตมโน อายสฺมา อานนฺโท ภควโต ภาสิตํ อภินนฺทีติ.}

\nitthitamruled{อนาถปิณฺฑิโกวาทสุตฺตํ นิฏฺฐิตํ ปฐมํ.}
\csromanmatikaanchor{mtk.53}
\csromantocmarkref{subsection}{2. ฉนฺโนวาทสุตฺต}{mtk.53}
\bookmark[dest={mtk.53},level=2]{2. ฉนฺโนวาทสุตฺต}
\csromanheader{2}{2. ฉนฺโนวาทสุตฺต}
\proseitem{389}{เอวํ เม สุตํ\csromandash{}เอกํ สมยํ ภควา ราชคเห วิหรติ เวฬุวเน กลนฺทกนิวาเป.\csromanspacer{} เตน โข ปน สมเยน อายสฺมา จ สาริปุตฺโต อายสฺมา จ มหาจุนฺโท อายสฺมา จ ฉนฺโน คิชฺฌกูเฏ ปพฺพเต วิหรนฺติ.\csromanspacer{} เตน โข ปน สมเยน อายสฺมา ฉนฺโน อาพาธิโก โหติ ทุกฺขิโต พาฬฺหคิลาโน.\csromanspacer{} อถ โข อายสฺมา สาริปุตฺโต สายนฺหสมยํ ปฏิสลฺลานา วุฏฺฐิโต เยนายสฺมา มหาจุนฺโท เตนุปสงฺกมิ, อุปสงฺกมิตฺวา อายสฺมนฺตํ มหาจุนฺทํ เอตทโวจ “อายามาวุโส จุนฺท เยนายสฺมา ฉนฺโน เตนุปสงฺกมิสฺสาม คิลานปุจฺฉกา”ติ. “เอวมาวุโส”ติ โข อายสฺมา มหาจุนฺโท อายสฺมโต สาริปุตฺตสฺส ปจฺจสฺโสสิ.}

\prose{อถ โข อายสฺมา จ สาริปุตฺโต อายสฺมา จ มหาจุนฺโท เยนายสฺมา ฉนฺโน เตนุปสงฺกมิํสุ, อุปสงฺกมิตฺวา อายสฺมตา ฉนฺเนน สทฺธิํ สมฺโมทิํสุ, สมฺโมทนียํ กถํ สารณียํ วีติสาเรตฺวา เอกมนฺตํ นิสีทิํสุ, เอกมนฺตํ นิสินฺโน โข อายสฺมา สาริปุตฺโต อายสฺมนฺตํ ฉนฺนํ เอตทโวจ “กจฺจิ เต อาวุโส ฉนฺน ขมนียํ, กจฺจิ ยาปนียํ, กจฺจิ เต ทุกฺขา เวทนา ปฏิกฺกมนฺติ โน อภิกฺกมนฺติ, ปฏิกฺกโมสานํ ปญฺญายติ โน อภิกฺกโม”ติ.}

\prose{น เม อาวุโส สาริปุตฺต ขมนียํ, น ยาปนียํ, พฬฺหา เม ทุกฺขา เวทนา อภิกฺกมนฺติ โน ปฏิกฺกมนฺติ, อภิกฺกโมสานํ ปญฺญายติ โน ปฏิกฺกโม.\csromanspacer{} เสยฺยถาปิ อาวุโส สาริปุตฺต พลวา \csromanfolio{308}ปุริโส ติเณฺหน สิขเรน มุทฺธนิ อภิมตฺเถยฺย.\csromanspacer{} เอวเมว โข เม อาวุโส สาริปุตฺต อธิมตฺตา วาตา มุทฺธนิ อูหนนฺติ.\csromanspacer{} น เม อาวุโส สาริปุตฺต ขมนียํ, น ยาปนียํ, พาฬฺหา เม ทุกฺขา เวทนา อภิกฺกมนฺติ โน ปฏิกฺกมนฺติ, อภิกฺกโมสานํ ปญฺญายติ โน ปฏิกฺกโม.\csromanspacer{} เสยฺยถาปิ อาวุโส สาริปุตฺต พลวา ปุริโส ทเฬฺหน วรตฺตกฺขณฺเฑน สีเส สีสเวฐํ ทเทยฺย.\csromanspacer{} เอวเมว โข เม อาวุโส สาริปุตฺต อธิมตฺตา สีเส สีสเวทนา.\csromanspacer{} น เม อาวุโส สาริปุตฺต ขมนียํ, น ยาปนียํ, พาฬฺหา เม ทุกฺขา เวทนา อภิกฺกมนฺติ โน ปฏิกฺกมนฺติ อภิกฺกโมสานํ ปญฺญายติ โน ปฏิกฺกโม.\csromanspacer{} เสยฺยถาปิ อาวุโส สาริปุตฺต ทกฺโข โคฆาตโก วา โคฆาตกนฺเตวาสี วา ติเณฺหน โควิกนฺตเนน กุจฺฉิํ ปริกนฺเตยฺย.\csromanspacer{} เอวเมว โข เม อาวุโส สาริปุตฺต อธิมตฺตา วาตา กุจฺฉิํ ปริกนฺตนฺติ.\csromanspacer{} น เม อาวุโส สาริปุตฺต ขมนียํ, น ยาปนียํ, พาฬฺหา เม ทุกฺขา เวทนา อภิกฺกมนฺติ โน ปฏิกฺกมนฺติ, อภิกฺกโมสานํ ปญฺญายติ โน ปฏิกฺกโม.\csromanspacer{} เสยฺยถาปิ อาวุโส สาริปุตฺต เทฺว พลวนฺโต ปุริสา ทุพฺพลตรํ ปุริสํ นานาพาหาสุ คเหตฺวา องฺคารกาสุยา สนฺตาเปยฺยุํ สมฺปริตาเปยฺยุํ.\csromanspacer{} เอวเมว โข เม อาวุโส สาริปุตฺต อธิมตฺโต กายสฺมิํ ฑาโห.\csromanspacer{} น เม อาวุโส สาริปุตฺต ขมนียํ, น ยาปนียํ, พาฬา เม ทุกฺขา เวทนา อภิกฺกมนฺติ โน ปฏิกฺกมนฺติ, อภิกฺกโมสานํ ปญฺญายติ โน ปฏิกฺกโม.\csromanspacer{} สตฺถํ อาวุโส สาริปุตฺต อาหริสฺสามิ, นาวกงฺขามิ ชีวิตนฺติ.}

\proseitem{390}{มายสฺมา ฉนฺโน สตฺถํ อาหเรสิ, ยาเปตายสฺมา ฉนฺโน, ยาเปนฺตํ มยํ อายสฺมนฺตํ ฉนฺนํ อิจฺฉาม.\csromanspacer{} สเจ อายสฺมโต ฉนฺนสฺส นตฺถิ สปฺปายานิ โภชนานิ, อหํ อายสฺมโต ฉนฺนสฺส สปฺปายานิ โภชนานิ ปริเยสิสฺสามิ.\csromanspacer{} สเจ อายสฺมโต ฉนฺนสฺส นตฺถิ สปฺปายานิ เภสชฺชานิ, อหํ อายสฺมโต ฉนฺนสฺส สปฺปายานิ เภสชฺชานิ ปริเยสิสฺสามิ.\csromanspacer{} สเจ อายสฺมโต ฉนฺนสฺส นตฺถิ ปติรูปา อุปฏฺฐากา, อหํ อายสฺมนฺตํ ฉนฺนํ อุปฏฺฐหิสฺสามิ.\csromanspacer{} มายสฺมา ฉนฺโน สตฺถํ อาหเรสิ, ยาเปตายสฺมา ฉนฺโน, ยาเปนฺตํ มยํ อายสฺมนฺตํ ฉนฺนํ อิจฺฉามาติ.}

\prose{นปิ เม อาวุโส สาริปุตฺต นตฺถิ สปฺปายานิ โภชนานิ, นปิ เม นตฺถิ สปฺปายานิ เภสชฺชานิ, นปิ เม นตฺถิ ปติรูปา อุปฏฺฐากา.\csromanspacer{} อปิ จาวุโส}

\csromanheader{2}{2. ฉนฺโนวาทสุตฺต (144)}
\prose{\csromanfolio{309}สาริปุตฺต ปริจิณฺโณ เม สตฺถา ทีฆรตฺตํ มนาเปเนว โน อมนาเปน.\csromanspacer{} เอตํ หิ อาวุโส สาริปุตฺต สาวกสฺส ปติรูปํ, ยํ สตฺถารํ ปริจเรยฺย มนาเปเนว โน อมนาเปน.\csromanspacer{} อนุปวชฺชํ ฉนฺโน ภิกฺขุ สตฺถํ อาหริสฺสตีติ เอวเมตํ\footnote{เอวเมว โข ตฺวํ (ก)} อาวุโส สาริปุตฺต ธาเรหีติ.\csromanspacer{} ปุจฺเฉยฺยาม มยํ อายสฺมนฺตํ ฉนฺนํ กิญฺจิเทว เทสํ, สเจ อายสฺมา ฉนฺโน โอกาสํ กโรติ ปญฺหสฺส เวยฺยากรณายาติ.\csromanspacer{} ปุจฺฉาวุโส สาริปุตฺต สุตฺวา เวทิสฺสามีติ.}

\proseitem{391}{จกฺขุํ อาวุโส ฉนฺน จกฺขุ วิญฺญาณํ จกฺขุ วิญฺญาณวิญฺญาตพฺเพ ธมฺเม “เอตํ มม, เอโสหมสฺมิ, เอโส เม อตฺตา”ติ สมนุปสฺสสิ? โสตํ อาวุโส ฉนฺน โสตวิญฺญาณํ ฯเปฯ.\csromanspacer{} ฆานํ อาวุโส ฉนฺน ฆานวิญฺญาณํ.\csromanspacer{} ชิวฺหํ อาวุโส ฉนฺน ชิวฺหาวิญฺญาณํ.\csromanspacer{} กายํ อาวุโส ฉนฺน กายวิญฺญาณํ.\csromanspacer{} มนํ อาวุโส ฉนฺน มโนวิญฺญาณํ มโนวิญฺญาณวิญฺญาตพฺเพ ธมฺเม “เอตํ มม, เอโสหมสฺมิ, เอโส เม อตฺตา”ติ สมนุปสฺสสีติ?}

\prose{จกฺขุํ อาวุโส สาริปุตฺต จกฺขุวิญฺญาณํ จกฺขุวิญฺญาณวิญฺญาตพฺเพ ธมฺเม “เนตํ มม, เนโสหมสฺมิ, น เมโส อตฺตา”ติ สมนุปสฺสามิ.\csromanspacer{} โสตํ อาวุโส สาริปุตฺต ฯเปฯ.\csromanspacer{} ฆานํ อาวุโส สาริปุตฺต.\csromanspacer{} ชิวฺหํ อาวุโส สาริปุตฺต.\csromanspacer{} กายํ อาวุโส สาริปุตฺต.\csromanspacer{} มนํ อาวุโส สาริปุตฺต มโนวิญฺญาณํ มโนวิญฺญาณํวิญฺญาตพฺเพ ธมฺเม “เนตํ มม, เนโส หมสฺมิ, น เมโส อตฺตา”ติ สมนุปสฺสามีติ.}

\proseitem{392}{จกฺขุสฺมิํ อาวุโส ฉนฺน จกฺขุวิญฺญาเณ จกฺขุวิญฺญาณวิญฺญาตพฺเพสุ ธมฺเมสุ กิํ ทิสฺวา กิํ อภิญฺญาย จกฺขุํ จกฺขุวิญฺญาณํ จกฺขุวิญฺญาณวิญฺญาตพฺเพ ธมฺเม “เนตํ มม, เนโสหมสฺมิ, น เมโส อตฺตา”ติ สมนุปสฺสสิ? โสตสฺมิํ อาวุโส ฉนฺน โสตวิญฺญาเณ.\csromanspacer{} ฆานสฺมิํ อาวุโส ฉนฺน ฆานวิญฺญาเณ.\csromanspacer{} ชิวฺหาย อาวุโส ฉนฺน ชิวฺหาวิญฺญาเณ.\csromanspacer{} กายสฺมิํ อาวุโส ฉนฺน กายวิญฺญาเณ.\csromanspacer{} มนสฺมิํ อาวุโส ฉนฺน มโนวิญฺญาเณ มโนวิญฺญาณวิญฺญาตพฺเพสุ ธมฺเมสุ กิํ ทิสฺวา กิํ อภิญฺญาย มนํ มโนวิญฺญาณํ มโนวิญฺญาณวิญฺญาตพฺเพ ธมฺเม “เนตํ มม, เนโสหมสฺมิ, น เมโส อตฺตา”ติ สมนุปสฺสสีติ?}

\prose{\csromanfolio{310}จกฺขุสฺมิํ อาวุโส สาริปุตฺต จกฺขุวิญฺญาเณ จกฺขุวิญฺญาณวิญฺญาตพฺเพสุ ธมฺเมสุ นิโรธํ ทิสฺวา นิโรธํ อภิญฺญาย จกฺขุํ จกฺขุวิญฺญาณํ จกฺขุวิญฺญาณวิญฺญาตพฺเพ ธมฺเม “เนตํ มม, เนโสหมสฺมิ, น เมโส อตฺตา”ติ สมนุปสฺสามิ.\csromanspacer{} โสตสฺมิํ อาวุโส สาริปุตฺต โสตวิญฺญาเณ.\csromanspacer{} ฆานสฺมิํ อาวุโส สาริปุตฺต ฆานวิญฺญาเณ.\csromanspacer{} ชิวฺหาย อาวุโส สาริปุตฺต ชิวฺหาวิญฺญาเณ.\csromanspacer{} กายสฺมิํ อาวุโส สาริปุตฺต กายวิญฺญาเณ.\csromanspacer{} มนสฺมิํ อาวุโส สาริปุตฺต มโนวิญฺญาเณ มโนวิญฺญาณวิญฺญาตพฺเพสุ ธมฺเมสุ นิโรธํ ทิสฺวา นิโรธํ อภิญฺญาย มนํ มโนวิญฺญาณํ มโนวิญฺญาณวิญฺญาตพฺเพ ธมฺเม “เนตํ มม, เนโสหมสฺมิ, น เมโส อตฺตา”ติ สมนุปสฺสามีติ.}

\proseitem{393}{เอวํ วุตฺเต อายสฺมา มหาจุนฺโท อายสฺมนฺตํ ฉนฺนํ เอตทโวจ “ตสฺมาติห อาวุโส ฉนฺน อิทมฺปิ ตสฺส ภควโต สาสนํ\footnote{วจนํ (สี)} นิจฺจกปฺปํ มนสิ กาตพฺพํ ‘นิสฺสิตสฺส จลิตํ, อนิสฺสิตสฺส จลิตํ นตฺถิ, จลิเต อสติ ปสฺสทฺธิ, ปสฺสทฺธิยา สติ นติ น โหติ, นติยา อสติ อาคติคติ น โหติ, อาคติคติยา อสติ จุตูปปาโต น โหติ, จุตูปปาเต อสติ เนวิธ น หุรํ น อุภยมนฺตเรน, เอเสวนฺโต ทุกฺขสฺสา’ติ”.\csromanspacer{} อถ โข อายสฺมา จ สาริปุตฺโต อายสฺมา จ มหาจุนฺโท อายสฺมนฺตํ ฉนฺนํ อิมินา โอวาเทน โอวทิตฺวา อุฏฺฐายาสนา ปกฺกมิํสุ.}

\proseitem{394}{อถ โข อายสฺมา ฉนฺโน อจิรปกฺกนฺเต อายสฺมนฺเต จ สาริปุตฺเต อายสฺมนฺเต จ มหาจุนฺเท สตฺถํ อาหเรสิ.\csromanspacer{} อถ โข อายสฺมา สาริปุตฺโต เยน ภควา เตนุปสงฺกมิ, อุปสงฺกมิตฺวา ภควนฺตํ อภิวาเทตฺวา เอกมนฺตํ นิสีทิ, เอกมนฺตํ นิสินฺโน โข อายสฺมา สาริปุตฺโต ภควนฺตํ เอตทโวจ “อายสฺมตา ภนฺเต ฉนฺเนน สตฺถํ อาหริตํ, ตสฺส กา คติ, โก อภิสมฺปราโย”ติ.\csromanspacer{} นนุ เต สาริปุตฺต ฉนฺเนน ภิกฺขุนา สมฺมุขาเยว อนุปวชฺชตา พฺยากตาติ.\csromanspacer{} อตฺถิ ภนฺเต ปุพฺพชิรํ\footnote{ปปฺปชิตญฺหิตํ (ก), อุปวชฺชิตํ (ก), ปุพฺพวิชฺชนํ, ปุพฺพวิชฺฌนํ, ปุพฺพวิจิรํ (สํยุตฺตเก)} นาม วชฺชิคาโม, ตตฺถายสฺมโต ฉนฺนสฺส มิตฺตกุลานิ สุหชฺชกุลานิ อุปวชฺชกุลานีติ.\csromanspacer{} โหนฺติ\footnote{โปสนฺติ (ก)} เหเต สาริปุตฺต ฉนฺนสฺส ภิกฺขุโน มิตฺตกุลานิ สุหชฺชกุลานิ อุปวชฺชกุลานิ, นาหํ สาริปุตฺต}

\csromanheader{2}{3. ปุณฺโณวาทสุตฺต (145)}
\prose{\csromanfolio{311}เอตฺตาวตา ส-อุปวชฺโชติ วทามิ.\csromanspacer{} โย โข สาริปุตฺต อิมญฺจ กายํ นิกฺขิปติ, อญฺญญฺจ กายํ อุปาทิยติ, ตมหํ ส-อุปวชฺโชติ วทามิ, ตํ ฉนฺนสฺส ภิกฺขุโน นตฺถิ, อนุปวชฺโช ฉนฺโน ภิกฺขุ สตฺถํ อาหเรสีติ เอวเมตํ สาริปุตฺต ธาเรหีติ.}

\prose{อิทมโวจ ภควา.\csromanspacer{} อตฺตมโน อายสฺมา สาริปุตฺโต ภควโต ภาสิตํ อภินนฺทีติ.}

\nitthitamruled{ฉนฺโนวาทสุตฺตํ นิฏฺฐิตํ ทุติยํ.}
\csromanmatikaanchor{mtk.54}
\csromantocmarkref{subsection}{3. ปุณฺโณวาทสุตฺต}{mtk.54}
\bookmark[dest={mtk.54},level=2]{3. ปุณฺโณวาทสุตฺต}
\csromanheader{2}{3. ปุณฺโณวาทสุตฺต}
\proseitem{395}{เอวํ เม สุตํ\csromandash{}เอกํ สมยํ ภควา สาวตฺถิยํ วิหรติ เชตวเน อนาถปิณฺฑิกสฺส อาราเม.\csromanspacer{} อถ โข อายสฺมา ปุณฺโณ สายนฺหสมยํ ปฏิสลฺลานา วุฏฺฐิโต เยน ภควา เตนุปสงฺกมิ, อุปสงฺกมิตฺวา ภควนฺตํ อภิวาเทตฺวา เอกมนฺตํ นิสีทิ, เอกมนฺตํ นิสินฺโน โข อายสฺมา ปุณฺโณ ภควนฺตํ เอตทโวจ “สาธุ มํ ภนฺเต ภควา สํขิตฺเตน โอวาเทน โอวทตุ, ยมหํ ภควโต ธมฺมํ สุตฺวา เอโก วูปกฏฺโฐ อปฺปมตฺโต อาตาปี ปหิตตฺโต วิหเรยฺยนฺ”ติ.\csromanspacer{} เตน หิ ปุณฺณ สุณาหิ สาธุกํ มนสิ กโรหิ, ภาสิสฺสามีติ. “เอวํ ภนฺเต”ติ โข อายสฺมา ปุณฺโณ ภควโต ปจฺจสฺโสสิ.\csromanspacer{} ภควา เอตทโวจ\csromandash{}}

\prose{สนฺติ โข ปุณฺณ จกฺขุวิญฺเญยฺยา รูปา อิฏฺฐา กนฺตา มนาปา ปิยรูปา กามูปสํหิตา รชนียา, ตํ เจ ภิกฺขุ อภินนฺทติ อภิวทติ อชฺโฌสาย ติฏฺฐติ.\csromanspacer{} ตสฺส ตํ อภินนฺทโต อภิวทโต อชฺโฌสาย ติฏฺฐโต อุปฺปชฺชติ นนฺที\footnote{นนฺทิ (สฺยา, กํ)}, นนฺทีสมุทยา ทุกฺขสมุทโย ปุณฺณาติ วทามิ.}

\prose{สนฺติ โข ปุณฺณ โสตวิญฺเญยฺยา สทฺทา.\csromanspacer{} ฆานวิญฺเญยฺยา คนฺธา.\csromanspacer{} ชิวฺหาวิญฺเญยฺยา รสา.\csromanspacer{} กายวิญฺเญยฺยา โผฏฺฐพฺพา.\csromanspacer{} มโนวิญฺเญยฺยา ธมฺมา อิฏฺฐา กนฺตา มนาปา ปิยรูปา กามูปสํหิตา รชนียา, ตํ เจ ภิกฺขุ อภินนฺทติ อภิวทติ อชฺโฌสาย ติฏฺฐติ.\csromanspacer{} ตสฺส ตํ อภินนฺทโต อภิวทโต อชฺโฌสาย ติฏฺฐโต อุปฺปชฺชติ นนฺที, นนฺทีสมุทยา ทุกฺขสมุทโย ปุณฺณาติ วทามิ.}

\prose{\csromanfolio{312}สนฺติ จ โข ปุณฺณ จกฺขุวิญฺเญยฺยา รูปา อิฏฺฐา กนฺตา มนาปา ปิยรูปา กามูปสํหิตา รชนียา, ตํ เจ ภิกฺขุ นาภินนฺทติ นาภิวทติ นาชฺโฌสาย ติฏฺฐติ.\csromanspacer{} ตสฺส ตํ อนภินนฺทโต อนภิวทโต อนชฺโฌสาย ติฏฺฐโต นนฺที นิรุชฺฌติ, นนฺทีนิโรธา ทุกฺขนิโรโธ ปุณฺณาติ วทามิ.}

\prose{สนฺติ จ โข ปุณฺณ โสตวิญฺเญยฺยา สทฺทา.\csromanspacer{} ฆานวิญฺเญยฺยา คนฺธา.\csromanspacer{} ชิวฺหาวิญฺเญยฺยา รสา.\csromanspacer{} กายวิญฺเญยฺยา โผฏฺฐพฺพา.\csromanspacer{} มโนวิญฺเญยฺยา ธมฺมา อิฏฺฐา กนฺตา มนาปา ปิยรูปา กามูปสํหิตา รชนียา, ตํ เจ ภิกฺขุ นาภินนฺทติ นาภิวทติ นาชฺโฌสาย ติฏฺฐติ.\csromanspacer{} ตสฺส ตํ อนภินนฺทโต อนภิวทโต อนชฺโฌสาย ติฏฺฐโต นนฺที นิรุชฺฌติ, นนฺทีนิโรธา ทุกฺขนิโรโธ ปุณฺณาติ วทามิ.}

\prose{อิมินา จ ตฺวํ ปุณฺณ มยา สํขิตฺเตน โอวาเทน โอวทิโต กตรสฺมิํ ชนปเท วิหริสฺสสีติ.\csromanspacer{} อิมินาหํ ภนฺเต ภควตา สํขิตฺเตน โอวาเทน โอวทิโต อตฺถิ สุนาปรนฺโต นาม ชนปโท, ตตฺถาหํ วิหริสฺสามีติ.}

\proseitem{396}{จณฺฑา โข ปุณฺณ สุนาปรนฺตกา มนุสฺสา, ผรุสา โข ปุณฺณ สุนาปรนฺตกา มนุสฺสา.\csromanspacer{} สเจ ตํ ปุณฺณ สุนาปรนฺตกา มนุสฺสา อกฺโกสิสฺสนฺติ ปริภาสิสฺสนฺติ, ตตฺถ เต ปุณฺณ กินฺติ ภวิสฺสตีติ.\csromanspacer{} สเจ มํ ภนฺเต สุนาปรนฺตกา มนุสฺสา อกฺโกสิสฺสนฺติ ปริภาสิสฺสนฺติ, ตตฺถ เม เอวํ ภวิสฺสติ “ภทฺทกา\footnote{ภทฺรกา (ก)} วติเม สุนาปรนฺตกา มนุสฺสา, สุภทฺทกา วติเม สุนาปรนฺตกา มนุสฺสา, ยํ เม นยิเม ปาณินา ปหารํ เทนฺตี”ติ, เอวเมตฺถ\footnote{เอวมฺเมตฺถ (?)} ภควา ภวิสฺสติ, เอวเมตฺถ สุคต ภวิสฺสตีติ.}

\prose{สเจ ปน เต ปุณฺณ สุนาปรนฺตกา มนุสฺสา ปาณินา ปหารํ ทสฺสนฺติ, ตตฺถ ปน เต ปุณฺณ กินฺติ ภวิสฺสตีติ.\csromanspacer{} สเจ เม ภนฺเต สุนาปรนฺตกา มนุสฺสา, ปาณินา ปหารํ ทสฺสนฺติ, ตตฺถ เม เอวํ ภวิสฺสติ “ภทฺทกา วติเม สุนาปรนฺตกา มนุสฺสา, สุภทฺทกา วติเม สุนาปรนฺตกา มนุสฺสา.\csromanspacer{} ยํ เม นยิเม เลฑฺฑุนา ปหารํ เทนฺตี”ติ, เอวเมตฺถ ภควา ภวิสฺสติ, เอวเมตฺถ สุคต ภวิสฺสตีติ.}

\prose{สเจ ปน เต ปุณฺณ สุนาปรนฺตกา มนุสฺสา เลฑฺฑุนา ปหารํ ทสฺสนฺติ, ตตฺถ ปน เต ปุณฺณ กินฺติ ภวิสฺสตีติ.\csromanspacer{} สเจ เม ภนฺเต สุนาปรนฺตกา}

\csromanheader{2}{3. ปุณฺโณวาทสุตฺต (145)}
\prose{\csromanfolio{313}มนุสฺสา เลฑฺฑุนา ปหารํ ทสฺสนฺติ, ตตฺถ เม เอวํ ภวิสฺสติ “ภทฺทกา วติเม สุนาปรนฺตกา มนุสฺสา, สุภทฺทกา วติเม สุนาปรนฺตกา มนุสฺสา, ยํ เม นยิเม ทณฺเฑน ปหารํ เทนฺตี”ติ, เอวเมตฺถ ภควา ภวิสฺสติ, เอวเมตฺถ สุคต ภวิสฺสตีติ.}

\prose{สเจ ปน เต ปุณฺณ สุนาปรนฺตกา มนุสฺสา สณฺเฑน ปหารํ ทสฺสนฺติ, ตตฺถ ปน เต ปุณฺณ กินฺติ ภวิสฺสตีติ.\csromanspacer{} สเจ เม ภนฺเต สุนาปรนฺตกา มนุสฺสา ทณฺเฑน ปหารํ ทสฺสนฺติ, ตตฺถ เม เอวํ ภวิสฺสติ “ภทฺทกา วติเม สุนาปรนฺตกา มนุสฺสา, สุภทฺทกา วติเม สุนาปรนฺตกา มนุสฺสา, ยํ เม นยิเม สตฺเถน ปหารํ เทนฺตี”ติ, เอวเมตฺถ ภควา ภวิสฺสติ, เอวเมตฺถ สุคต ภวิสฺสตีติ.}

\prose{สเจ ปน เต ปุณฺณ สุนาปรนฺตกา มนุสฺสา สตฺเถน ปหารํ ทสฺสนฺติ, ตตฺถ ปน เต ปุณฺณ กินฺติ ภวิสฺสตีติ.\csromanspacer{} สเจ เม ภนฺเต สุนาปรนฺตกา มนุสฺสา สตฺเถน ปหารํ ทสฺสนฺติ, ตตฺถ เม เอวํ ภวิสฺสติ “ภทฺทกา วติเม สุนาปรนฺตกา มนุสฺสา, สุภทฺทกา วติเม สุนาปรนฺตกา มนุสฺสา, ยํ มํ\footnote{ยํ เม (สี, อิ, ก)} นยิเม ติเณฺหน สตฺเถน ชีวิตา โวโรเปนฺติ”ติ, เอวเมตฺถ ภควา ภวิสฺสติ, เอวเมตฺถ สุคต ภวิสฺสตีติ.}

\prose{สเจ ปน ตํ ปุณฺณ สุนาปรนฺตกา มนุสฺสา ติเณฺหน สตฺเถน ชีวิตา โวโรเปสฺสนฺติ, ตตฺถ ปน เต ปุณฺณ กินฺติ ภวิสฺสตีติ.\csromanspacer{} สเจ มํ ภนฺเต สุนาปรนฺตกา มนุสฺสา ติเณฺหน สตฺเถน ชีวิตา โวโรเปสฺสนฺติ, ตตฺถ เม เอวํ ภวิสฺสติ “สนฺติ โข ภควโต สาวกา กาเย จ ชีวิเต จ อฑฺฑียมานา หรายมานา ชิคุจฺฉมานา สตฺถหารกํ ปริเยสนฺติ, ตํ เม อิทํ อปริยิฏฺฐํเยว สตฺถหารกํ ลทฺธนฺ”ติ, เอวเมตฺถ ภควา ภวิสฺสติ, เอวเมตฺถ สุคต ภวิสฺสตีติ.\csromanspacer{} สาธุ สาธุ ปุณฺณ, สกฺขิสฺสสิ โข ตฺวํ ปุณฺณ อิมินา ทมูปสเมน สมนฺนาคโต สุนาปรนฺตสฺมิํ ชนปเท วิหริตุํ, ยสฺสทานิ ตฺวํ ปุณฺณ กาลํ มญฺญสีติ.}

\proseitem{397}{อถ โข อายสฺมา ปุณฺโณ ภควโต ภาสิตํ อภินนฺทิตฺวา อนุโมทิตฺวา อุฏฺฐายาสนา ภควนฺตํ อภิวาเทตฺวา ปทกฺขิณํ กตฺวา เสนาสนํ สํสาเมตฺวา ปตฺตจีวรมาทาย เยน สุนาปรนฺโต ชนปโท \csromanfolio{314}เตน จาริกํ ปกฺกามิ.\csromanspacer{} อนุปพฺเพน จาริกํ จรมาโน เยน สุนาปรนฺโต ชนปโท ตทวสริ.\csromanspacer{} ตตฺร สุทํ อายสฺมา ปุณฺโณ สุนาปรนฺตสฺมิํ ชนปเท วิหรติ.\csromanspacer{} อถ โข อายสฺมา ปุณฺโณ เตเนวนฺตรวสฺเสน ปญฺจมตฺตานิ อุปาสกสตานิ ปฏิเวเทสิ\footnote{ปฏิปาเทสิ (สี, อิ), ปฏิเทเสสิ (สฺยา, กํ)}, เตเนวนฺตรวสฺเสน ปญฺจมตฺตานิ อุปาสิกสตานิ ปฏิเวเทสิ, เตเนวนฺตรวสฺเสน ติสฺโส วิชฺชา สจฺฉากาสิ.\csromanspacer{} อถ โข อายสฺมา ปุณฺโณ อปเรน สมเยน ปรินิพฺพายิ.}

\prose{อถ โข สมฺพหุลา ภิกฺขู เยน ภควา เตนุปสงฺกมิํสุ, อุปสงฺกมิตฺวา ภควนฺตํ อภิวาเทตฺวา เอกมนฺตํ นิสีทิํสุ, เอกมนฺตํ นิสินฺนา โข เต ภิกฺขู ภควนฺตํ เอตทโวจุํ “โย โส ภนฺเต ปุณฺโณ นาม กุลปุตฺโต ภควตา สํขิตฺเตน โอวาเทน โอวทิโต, โส กาลงฺกโต, ตสฺส กา คติ, โก อภิสมฺปราโย”ติ.\csromanspacer{} ปณฺฑิโต ภิกฺขเว ปุณฺโณ กุลปุตฺโต ปจฺจปาทิ\footnote{สจฺจวาที ธมฺมวาที (ก)} ธมฺมสฺสานุธมฺมํ, น จ มํ ธมฺมาธิกรณํ วิเหเสสิ, ปรินิพฺพุโต ภิกฺขเว ปุณฺโณ กุลปุตฺโตติ.}

\prose{อิทมโวจ ภควา.\csromanspacer{} อตฺตมนา เต ภิกฺขู ภควโต ภาสิตํ อภินนฺทุนฺติ.}

\nitthitamruled{ปุณฺโณวาทสุตฺตํ นิฏฺฐิตํ ตติยํ.}
\csromanmatikaanchor{mtk.55}
\csromantocmarkref{subsection}{4. นนฺทโกวาทสุตฺต}{mtk.55}
\bookmark[dest={mtk.55},level=2]{4. นนฺทโกวาทสุตฺต}
\csromanheader{2}{4. นนฺทโกวาทสุตฺต}
\proseitem{398}{เอวํ เม สุตํ\csromandash{}เอกํ สมยํ ภควา สาวตฺถิยํ วิหรติ เชตวเน อนาถปิณฺฑิกสฺส อาราเม.\csromanspacer{} อถ โข มหาปชาปติโคตมี ปญฺจมตฺเตหิ ภิกฺขุนิสเตหิ สทฺธิํ เยน ภควา เตนุปสงฺกมิ, อุปสงฺกมิตฺวา ภควนฺตํ อภิวาเทตฺวา เอกมนฺตํ อฏฺฐาสิ, เอกมนฺตํ ฐิตา โข มหาปชาปติโคตมี ภควนฺตํ เอตทโว “โอวทตุ ภนฺเต ภควา ภิกฺขุนิโย, อนุสาสตุ ภนฺเต ภควา ภิกฺขุนิโย, กโรตุ ภนฺเต ภควา ภิกฺขุนีนํ ธมฺมิํ กถนฺ”ติ\footnote{ธมฺมิกถนฺติ (สฺยา, กํ, ก)}.}

\prose{เตน โข ปน สมเยน เถรา ภิกฺขู ภิกฺขุนิโย โอวทนฺติ ปริยาเยน.\csromanspacer{} อายสฺมา นนฺทโก น อิจฺฉติ ภิกฺขุนิโย โอวทิตุํ ปริยาเยน.}

\csromanheader{2}{4. นนฺทโกวาทสุตฺต (146)}
\prose{\csromanfolio{315}อถ โข ภควา อายสฺมนฺตํ อานนฺทํ อามนฺเตสิ “กสฺส นุ โข อานนฺท อชฺช ปริยาโย ภิกฺขุนิโย โอวทิตุํ ปริยาเยนา”ติ.\csromanspacer{} สพฺเพเหว ภนฺเต กโต\footnote{นนฺทกสฺส ภนฺเต (สี, อิ)} ปริยาโย ภิกฺขุนิโย โอวทิตุํ ปริยาเยน, อยํ ภนฺเต อายสฺมา นนฺทโก น อิจฺฉติ ภิกฺขุนิโย โอวทิตุํ ปริยาเยนาติ.}

\prose{อถ โข ภควา อายสฺมนฺตํ นนฺทกํ อามนฺเตสิ “โอวท นนฺทก ภิกฺขุนิโย, อนุสาส นนฺทก ภิกฺขุนิโย, กโรหิ ตฺวํ พฺราหฺมณ ภิกฺขุนีนํ ธมฺมิํ กถนฺ”ติ. “เอวํ ภนฺเต”ติ โข อายสฺมา นนฺทโก ภควโต ปฏิสฺสุตฺวา ปุพฺพณฺหสมยํ นิวาเสตฺวา ปตฺตจีวรมาทาย สาวตฺถิํ ปิณฺฑาย ปาวิสิ, สาวตฺถิยํ ปิณฺฑาย จริตฺวา ปจฺฉาภตฺตํ ปิณฺฑปาตปฏิกฺกนฺโต อตฺตทุติโย เยน ราชการาโม เตนุปสงฺกมิ.\csromanspacer{} อทฺทสํสุ โข ตา ภิกฺขุนิโย อายสฺมนฺตํ นนฺทกํ ทูรโตว อาคจฺฉนฺตํ, ทิสฺวานํ อาสนํ ปญฺญาเปสุํ อุทกญฺจ ปาทานํ อุปฏฺฐเปสุํ.\csromanspacer{} นิสีทิ โข อายสฺมา นนฺทโก ปญฺญตฺเต อาสเน, นิสชฺช ปาเท ปกฺขาเลสิ.\csromanspacer{} ตาปิ โข ภิกฺขุนิโย อายสฺมนฺตํ นนฺทกํ อภิวาเทตฺวา เอกมนฺตํ นิสีทิํสุ, เอกมนฺตํ นิสินฺนา โข ตา ภิกฺขุนิโย อายสฺมา นนฺทโก เอตทโวจ “ปฏิปุจฺฉกถา โข ภคินิโย ภวิสฺสติ.\csromanspacer{} ตตฺถ อาชานนฺตีหิ ‘อาชานามา’ติสฺส วจนียํ, น อาชานนฺตีหิ ‘น อาชานามา’ติสฺส วจนียํ, ยสฺสา วา ปนสฺส กงฺขา วา วิมติ วา, อหเมว ตตฺถ ปฏิปุจฺฉิตพฺโพ ‘อิทํ ภนฺเต กถํ อิมสฺส กฺวตฺโถ’ติ”.\csromanspacer{} เอตฺตเกนปิ มยํ ภนฺเต อยฺยสฺส นนฺทกสฺส อตฺตมนา อภิรทฺธา\footnote{อภินนฺทาม (สฺยา, กํ)}, ยํ โน อยฺโย นนฺทโก ปวาเรตีติ.}

\proseitem{399}{ตํ กิํ มญฺญถ ภคินิโย, จกฺขุ นิจฺจํ วา อนิจฺจํ วาติ.\csromanspacer{} อนิจฺจํ ภนฺเต.\csromanspacer{} ยํ ปนานิจฺจํ, ทุกฺขํ วา ตํ สุขํ วาติ.\csromanspacer{} ทุกฺขํ ภนฺเต.\csromanspacer{} ยํ ปนานิจฺจํ ทุกฺขํ วิปริณามธมฺมํ, กลฺลํ นุ ตํ สมนุปสฺสิตุํ “เอตํ มม, เอโสหมสฺมิ, เอโส เม อตฺตา”ติ.\csromanspacer{} โน เหตํ ภนฺเต.\csromanspacer{} ตํ กิํ มญฺญถ ภคินิโย, โสตํ นิจฺจํ วา อนิจฺจํ วาติ.\csromanspacer{} อนิจฺจํ ภนฺเต ฯเปฯ.\csromanspacer{} ฆานํ นิจฺจํ วา อนิจฺจํ วาติ.\csromanspacer{} อนิจฺจํ ภนฺเต.\csromanspacer{} ชิวฺหา นิจฺจา วา อนิจฺจา วาติ.\csromanspacer{} อนิจฺจา ภนฺเต.\csromanspacer{} กาโย นิจฺโจ วา อนิจฺโจ วาติ.\csromanspacer{} อนิจฺโจ ภนฺเต.\csromanspacer{} มโน นิจฺโจ วา อนิจฺโจ วาติ.\csromanspacer{} อนิจฺโจ ภนฺเต.\csromanspacer{} ยํ ปนานิจฺจํ, ทุกฺขํ วา ตํ สุขํ วาติ.\csromanspacer{} ทุกฺขํ ภนฺเต.\csromanspacer{} ยํ ปนานิจฺจํ ทุกฺขํ วิปริณามธมฺมํ, \csromanfolio{316}กลฺลํ นุ ตํ สมนุปสฺสิตุํ “เอตํ มม, เอโสหมสฺมิ, เอโส เม อตฺตา”ติ.\csromanspacer{} โน เหตํ ภนฺเต.\csromanspacer{} ตํ กิสฺส เหตุ, ปุพฺเพว โน เอตํ ภนฺเต ยถาภูตํ สมฺมปฺปญฺญาย สุทิฏฺฐํ “อิติปิเม ฉ อชฺฌตฺติกา อายตนา อนิจฺจา”ติ.\csromanspacer{} สาธุ สาธุ ภคินิโย, เอวํ เหตํ ภคินิโย โหติ อริยสาวกสฺส ยถาภูตํ สมฺมปฺปญฺญาย ปสฺสโต.}

\proseitem{400}{ตํ กิํ มญฺญถ ภคินิโย, รูปา นิจฺจา วา อนิจฺจา วาติ.\csromanspacer{} อนิจฺจา ภนฺเต.\csromanspacer{} ยํ ปนานิจฺจํ, ทุกฺขํ วา ตํ สุขํ วาติ.\csromanspacer{} ทุกฺขํ ภนฺเต.\csromanspacer{} ยํ ปนานิจฺจํ ทุกฺขํ วิปริณามธมฺมํ, กลฺลํ นุ ตํ สมนุปสฺสิตุํ “เอตํ มม, เอโสหมสฺมิ, เอโส เม อตฺตา”ติ.\csromanspacer{} โน เหตํ ภนฺเต.\csromanspacer{} ตํ กิํ มญฺญถ ภคินิโย, สทฺทา นิจฺจา วา อนิจฺจา วาติ.\csromanspacer{} อนิจฺจา ภนฺเต ฯเปฯ.\csromanspacer{} คนฺธา นิจฺจา วา อนิจฺจา วาติ.\csromanspacer{} อนิจฺจา ภนฺเต.\csromanspacer{} รสา นิจฺจา วา อนิจฺจา วาติ.\csromanspacer{} อนิจฺจา ภนฺเต.\csromanspacer{} โผฏฺฐพฺพา นิจฺจา วา อนิจฺจา วาติ.\csromanspacer{} อนิจฺจา ภนฺเต.\csromanspacer{} ธมฺมา นิจฺจา วา อนิจฺจา วาติ.\csromanspacer{} อนิจฺจา ภนฺเต.\csromanspacer{} ยํ ปนานิจฺจํ, ทุกฺขํ วา ตํ สุขํ วาติ.\csromanspacer{} ทุกฺขํ ภนฺเต.\csromanspacer{} ยํ ปนานิจฺจํ ทุกฺขํ วิปริณามธมฺมํ, กลฺลํ นุ ตํ สมนุปสฺสิตุํ “เอตํ มม, เอโสหมสฺมิ, เอโส เม อตฺตา”ติ.\csromanspacer{} โน เหตํ ภนฺเต.\csromanspacer{} ตํ กิสฺส เหตุ, ปุพฺเพว โน เอตํ ภนฺเต ยถาภูตํ สมฺมปฺปญฺญาย สุทิฏฺฐํ “อิติปิเม ฉ พาหิรา อายตนา อนิจฺจา”ติ.\csromanspacer{} สาธุ สาธุ ภคินิโย, เอวํ เหตํ ภคินิโย โหติ อริยสาวกสฺส ยถาภูตํ สมฺมปฺปญฺญาย ปสฺสโต.}

\proseitem{401}{ตํ กิํ มญฺญถ ภคินิโย, จกฺขุวิญฺญาณํ นิจฺจํ วา อนิจฺจํ วาติ.\csromanspacer{} อนิจฺจํ ภนฺเต.\csromanspacer{} ยํ ปนานิจฺจํ, ทุกฺขํ วา ตํ สุขํ วาติ.\csromanspacer{} ทุกฺขํ ภนฺเต.\csromanspacer{} ยํ ปนานิจฺจํ ทุกฺขํ วิปริณามธมฺมํ, กลฺลํ นุ ตํ สมนุปสฺสิตุํ “เอตํ มม, เอโสหมสฺมิ, เอโส เม อตฺตา”ติ.\csromanspacer{} โน เหตํ ภนฺเต.\csromanspacer{} ตํ กิํ มญฺญถ ภคินิโย, โสตวิญฺญาณํ นิจฺจํ วา อนิจฺจํ วาติ.\csromanspacer{} อนิจฺจํ ภนฺเต ฯเปฯ.\csromanspacer{} ฆานวิญฺญาณํ นิจฺจํ วา อนิจฺจํ วาติ.\csromanspacer{} อนิจฺจํ ภนฺเต.\csromanspacer{} ชิวฺหาวิญฺญาณํ นิจฺจํ วา อนิจฺจํ วาติ.\csromanspacer{} อนิจฺจํ ภนฺเต.\csromanspacer{} กายวิญฺญาณํ นิจฺจํ วา อนิจฺจํ วาติ.\csromanspacer{} อนิจฺจํ ภนฺเต.\csromanspacer{} มโนวิญฺญาณํ นิจฺจํ วา อนิจฺจํ วาติ.\csromanspacer{} อนิจฺจํ ภนฺเต.\csromanspacer{} ยํ ปนานิจฺจํ, ทุกฺขํ วา ตํ สุขํ วาติ.\csromanspacer{} ทุกฺขํ ภนฺเต.\csromanspacer{} ยํ ปนานิจฺจํ ทุกฺขํ วิปริณามธมฺมํ, กลฺลํ นุ ตํ สมนุปสฺสิตุํ “เอตํ มม, เอโสหมสฺมิ, เอโส เม อตฺตา”ติ.\csromanspacer{} โน เหตํ ภนฺเต.\csromanspacer{} ตํ กิสฺส เหตุ, ปุพฺเพว โน เอตํ ภนฺเต ยถาภูตํ สมฺมปฺปญฺญาย สุทิฏฺฐํ “อิติปิเม ฉ วิญฺญาณกายา อนิจฺจา”ติ.\csromanspacer{} สาธุ สาธุ ภคินิโย, เอวํ เหตํ ภคินิโย โหติ อริยสาวกสฺส ยถาภูตํ สมฺมปฺปญฺญาย ปสฺสโต.}

\csromanheader{2}{4. นนฺทโกวาทสุตฺต (146)}
\proseitem{402}{\csromanfolio{317}เสยฺยถาปิ ภคินิโย เตลปฺปทีปสฺส ฌายโต เตลมฺปิ อนิจฺจํ วิปริณามธมฺมํ, วฏฺฏิปิ อนิจฺจา วิปริณามธมฺมา, อจฺจิปิ อนิจฺจา วิปริณามธมฺมา, อาภาปิ อนิจฺจา วิปริณามธมฺมา.\csromanspacer{} โย นุ โข ภคินิโย เอวํ วเทยฺย “อมุสฺส เตลปฺปทีปสฺส ฌายโต เตลมฺปิ อนิจฺจํ วิปริณามธมฺมํ, วฏฺฏิปิ อนิจฺจา วิปริณามธมฺมํ, อจฺจิปิ อนิจฺจา วิปริณามธมฺมา.\csromanspacer{} ยา จ ขฺวาสฺส อาภา, สา นิจฺจา ธุวา สสฺสตา อวิปริณามธมฺมา”ติ, สมฺมา นุ โข โส ภคินิโย วทมาโน วเทยฺยาติ.\csromanspacer{} โน เหตํ ภนฺเต.\csromanspacer{} ตํ กิสฺส เหตุ, อมุสฺส หิ ภนฺเต เตลปฺปทีปสฺส ฌายโต เตลมฺปิ อนิจฺจํ วิปริณามธมฺมํ, วฏฺฏิปิ อนิจฺจา วิปริณามธมฺมา, อจฺจิปิ อนิจฺจา วิปริณามธมฺมา.\csromanspacer{} ปเควสฺส อาภา อนิจฺจา วิปริณามธมฺมาติ.\csromanspacer{} เอวเมว โข ภคินิโย, โย นุ โข เอวํ วเทยฺย “ฉ โขเม อชฺฌตฺติกา อายตนา อนิจฺจา\footnote{อนิจฺจา วิปริณามธมฺมา (?)}.\csromanspacer{} ยญฺจ โข ฉ อชฺฌตฺติเก อายตเน ปฏิจฺจ ปฏิสเวเทติ สุขํ วา ทุกฺขํ วา อทุกฺขมสุขํ วา, ตํ นิจฺจํ ธุวํ สสฺสตํ อวิปริณามธมฺมนฺ”ติ, สมฺมา นุ โข โส ภคินิโย วทมาโน วเทยฺยาติ.\csromanspacer{} โน เหตํ ภนฺเต.\csromanspacer{} ตํ กิสฺส เหตุ, ตชฺชํ ตชฺชํ ภนฺเต ปจฺจยํ ปฏิจฺจ ตชฺชา ตชฺชา เวทนา อุปฺปชฺชนฺติ, ตชฺชสฺส ตชฺชสฺส ปจฺจยสฺส นิโรธา ตชฺชา ตชฺชา เวทนา นิรุชฺฌนฺตีติ.\csromanspacer{} สาธุ สาธุ ภคินิโย, เอวํ เหตํ ภคินิโย โหติ อริยสาวกสฺส ยถาภูตํ สมฺมปฺปญฺญาย ปสฺสโต.}

\proseitem{403}{เสยฺยถาปิ ภคินิโย มหโต รุกฺขสฺส ติฏฺฐโต สารวโต มูลมฺปิ อนิจฺจํ วิปริณามธมฺมํ, ขนฺโธปิ อนิจฺโจ วิปริณามธมฺโม, สาขาปลาสมฺปิ อนิจฺจํ วิปริณามธมฺมํ, ฉายาปิ อนิจฺจา วิปริณามธมฺมา, โย นุ โข ภคินิโย เอวํ วเทยฺย “อมุสฺส มหโต รุกฺขสฺส ติฏฺฐโต สารวโต มูลมฺปิ อนิจฺจํ วิปริณามธมฺมํ, ขนฺโธปิ อนิจฺโจ วิปริณามธมฺโม, สาขาปลาสมฺปิ อนิจฺจํ วิปริณามธมฺมํ, ยา จ ขฺวาสฺส ฉายา, สา นิจฺจา ธุวา สสฺสตา อวิปริณามธมฺมา”ติ, สมฺมา นุ โข โส ภคินิโย วทมาโน วเทยฺยาติ.\csromanspacer{} โน เหตํ ภนฺเต.\csromanspacer{} ตํ กิสฺส เหตุ, อมุสฺส หิ ภนฺเต มหโต รุกฺขสฺส ติฏฺฐโต สารวโต มูลมฺปิ อนิจฺจํ วิปริณามธมฺมํ, ขนฺโธปิ อนิจฺโจ วิปริณามธมฺโม, สาขาปลาสมฺปิ อนิจฺจํ วิปริณามธมฺมํ.\csromanspacer{} ปเควสฺส ฉายา อนิจฺจา วิปริณามธมฺมาติ.\csromanspacer{} เอวเมว โข ภคินิโย, โย นุ โข เอวํ วเทยฺย “ฉ โขเม พาหิรา \csromanfolio{318}อายตนา อนิจฺจา\footnote{อนิจฺจา วิปริณามธมฺมา (สี, อิ)}.\csromanspacer{} ยญฺจ โข ฉ พาหิเร อายตเน ปฏิจฺจ ปฏิสํเวเทติ สุขํ วา ทุกฺขํ วา อทุกฺขมสุขํ วา, ตํ นิจฺจํ ธุวํ สสฺสตํ อวิปริณามธมฺมนฺ”ติ, สมฺมา นุ โข โส ภคินิโย วทมาโน วเทยฺยาติ.\csromanspacer{} โน เหตํ ภนฺเต.\csromanspacer{} ตํ กิสฺส เหตุ, ตชฺชํ ตชฺชํ ภนฺเต ปจฺจยํ ปฏิจฺจ ตชฺชา ตชฺชา เวทนา อุปฺปชฺชนฺติ, ตชฺชสฺส ตชฺชสฺส ปจฺจยสฺส นิโรธา ตชฺชา ตชฺชา เวทนา นิรุชฺฌนฺตีติ.\csromanspacer{} สาธุ สาธุ ภคินิโย, เอวํ เหตํ ภคินิโย โหติ อริยสาวกสฺส ยถาภูตํ สมฺมปฺปญฺญาย ปสฺสโต.}

\proseitem{404}{เสยฺยถาปิ ภคินิโย ทกฺโข โคฆาตโก วา โคฆาตกนฺเตวาสี วา คาวํ วธิตฺวา ติเณฺหน โควิกนฺตเนน คาวิํ สงฺกนฺเตยฺย อนุปหจฺจ อนฺตรํ มํสกายํ อนุปหจฺจ พาหิรํ จมฺมกายํ.\csromanspacer{} ยํ ยเทว ตตฺถ อนฺตรา วิลิมํสํ\footnote{วิลิมํ (สี, อิ, ก)} อนฺตรา นฺหารุ อนฺตรา พนฺธนํ, ตํ ตเทว ติเณฺหน โควิกนฺตเนน สญฺฉินฺเทยฺย สงฺกนฺเตยฺย สมฺปกนฺเตยฺย สมฺปริกนฺเตยฺย, สญฺฉินฺทิตฺวา สงฺกนฺติตฺวา สมฺปกนฺติตฺวา สมฺปริกนฺติตฺวา วิธุนิตฺวา พาหิรํ จมฺมกายํ เตเนว จมฺเมน ตํ คาวิํ ปฏิจฺฉาเทตฺวา เอวํ วเทยฺย “ตเถวายํ คาวี สํยุตฺตา อิมินาว จมฺเมนา”ติ, สมฺมา นุ โข โส ภคินิโย วทมาโน วเทยฺยาติ.\csromanspacer{} โน เหตํ ภนฺเต.\csromanspacer{} ตํ กิสฺส เหตุ, อมุ หิ ภนฺเต ทกฺโข โคฆาตโก วา โคฆาตกนฺเตวาสี วา คาวิํ วธิตฺวา ติเณฺหน โควิกนฺตเนน คาวิํ สงฺกนฺเตยฺย อนุปหจฺจ อนฺตรํ มํสกายํ อนุปหจฺจ พาหิรํ จมฺมกายํ.\csromanspacer{} ยํ ยเทว ตตฺถ อนฺตรา วิลิมํสํ อนฺตรา นฺหารุ อนฺตรา พนฺธนํ, ตํ ตเทว ติเณฺหน โควิกนฺตเนน สญฺฉินฺเทยฺย สงฺกนฺเตยฺย สมฺปกนฺเตยฺย สมฺปริกนฺเตยฺย, สญฺฉินฺทิตฺวา สงฺกนฺติตฺวา สมฺปกนฺติตฺวา สมฺปริกนฺติตฺวา วิธุนิตฺวา พาหิรํ จมฺมกายํ เตเนว จมฺเมน ตํ คาวิํ ปฏิจฺฉาเทตฺวา กิญฺจาปิ โส เอวํ วเทยฺย “ตเถวายํ คาวี สํยุตฺตา อิมินาว จมฺเมนา”ติ, อถ โข สา คาวี วิสํยุตฺตา เตเนว จมฺเมนาติ.}

\prose{อุปมา โข เม อยํ ภคินิโย กตา อตฺถสฺส วิญฺญาปนาย, อยเมเวตฺถ อตฺโถ. “อนฺตรา มํสกาโย”ติ โข ภคินิโย ฉนฺเนตํ อชฺฌตฺติกานํ อายตนานํ อธิวจนํ. “พาหิโร จมฺมกาโย”ติ โข ภคินิโย ฉนฺเนตํ พาหิรานํ อายตนานํ อธิวจนํ. “อนฺตรา วิลิมํสํ อนฺตรา นฺหารุ อนฺตรา พนฺธนนฺ”ติ โข ภคินิโย นนฺทีราคสฺเสตํ อธิวจนํ.}

\csromanheader{2}{4. นนฺทโกวาทสุตฺต (146)}
\prose{\csromanfolio{319}“ติณฺหํ โควิกนฺตนนฺ”ติ โข ภคินิโย อริยาเยตํ ปญฺญาย อธิวจนํ, ยายํ อริยา ปญฺญา อนฺตรา กิเลสํ อนฺตรา สํโยชนํ อนฺตรา พนฺธนํ สญฺฉินฺทติ สงฺกนฺตติ สมฺปกนฺตติ สมฺปริกนฺตติ.}

\proseitem{405}{สตฺต โข ปนิเม ภคินิโย โพชฺฌงฺคา, เยสํ ภาวิตตฺตา พหุลีกตตฺตา ภิกฺขุ อาสวานํ ขยา อนาสวํ เจโตวิมุตฺติํ ปญฺญาวิมุตฺติํ ทิฏฺเฐว ธมฺเม สยํ อภิญฺญา สจฺฉิกตฺวา อุปสมฺปชฺช วิหรติ.\csromanspacer{} กตเม สตฺต, อิธ ภคินิโย ภิกฺขุ สติสมฺโพชฺฌงฺคํ ภาเวติ วิเวกนิสฺสิตํ วิราคนิสฺสิตํ นิโรธนิสฺสิตํ โวสฺสคฺคปริณามิํ.\csromanspacer{} ธมฺมวิจยสมฺโพชฺฌงฺคํ ภาเวติ ฯเปฯ.\csromanspacer{} วีริยสมฺโพชฺฌงฺคํ ภาเวติ.\csromanspacer{} ปีติสมฺโพชฺฌงฺคํ ภาเวติ.\csromanspacer{} ปสฺสทฺธิสมฺโพชฺฌงฺคํ ภาเวติ.\csromanspacer{} สมาธิสมฺโพชฺฌงฺคํ ภาเวติ.\csromanspacer{} อุเปกฺขาสมฺโพชฺฌงฺคํ ภาเวติ วิเวกนิสฺสิตํ วิราคนิสฺสิตํ นิโรธนิสฺสิตํ โวสฺสคฺคปริณามิํ.\csromanspacer{} อิเม โข ภคินิโย สตฺต โพชฺฌงฺคา.\csromanspacer{} เยสํ ภาวิตตฺตา พหุลีกตตฺตา ภิกฺขุ อาสวานํ ขยา อนาสวํ เจโตวิมุตฺติํ ปญฺญาวิมุตฺติํ ทิฏฺเฐว ธมฺเม สยํ อภิญฺญา สจฺฉิกตฺวา อุปสมฺปชฺช วิหรตีติ.}

\proseitem{406}{อถ โข อายสฺมา นนฺทโก ตา ภิกฺขุนิโย อิมินา โอวาเทน โอวทิตฺวา อุยฺโยเชสิ “คจฺฉถ ภคินิโย กาโล”ติ.\csromanspacer{} อถ โข ตา ภิกฺขุนิโย อายสฺมโต นนฺทกสฺส ภาสิตํ อภินนฺทิตฺวา อนุโมทิตฺวา อุฏฺฐายาสนา อายสฺมนฺตํ นนฺทกํ อภิวาเทตฺวา ปทกฺขิณํ กตฺวา เยน ภควา เตนุปสงฺกมิํสุ, อุปสงฺกมิตฺวา ภควนฺตํ อภิวาเทตฺวา เอกมนฺตํ อฏฺฐํสุ, เอกมนฺตํ ฐิตา โข ตา ภิกฺขุนิโย ภควา เอตทโวจ “คจฺฉถ ภิกฺขุนิโย กาโล”ติ.\csromanspacer{} อถ โข ตา ภิกฺขุนิโย ภควนฺตํ อภิวาเทตฺวา ปทกฺขิณํ กตฺวา ปกฺกมิํสุ.\csromanspacer{} อถ โข ภควา อจิรปกฺกนฺตีสุ ตาสุ ภิกฺขุนีสุ ภิกฺขู อามนฺเตสิ “เสยฺยถาปิ ภิกฺขเว ตทหุโปสเถ จาตุทฺทเส น โหติ พหุโน ชนสฺส กงฺขา วา วิมติ วา ‘อูโน นุ โข จนฺโท, ปุณฺโณ นุ โข จนฺโท’ติ, อถ โข อูโน จนฺโทเตฺวว โหติ.\csromanspacer{} เอวเมว โข ภิกฺขเว ตา ภิกฺขุนิโย นนฺทกสฺส ธมฺมเทสนาย อตฺตมนา โหนฺติ, โน จ โข ปริปุณฺณสงฺกปฺปา”ติ.}

\proseitem{407}{อถ โข ภควา อายสฺมนฺตํ นนฺทกํ อามนฺเตสิ “เตน หิ ตฺวํ นนฺทก เสฺวปิ ตา ภิกฺขุนิโย เตเนโววาเทน โอวเทยฺยาสี”ติ. \csromanfolio{320}“เอวํ ภนฺเต”ติ โข อายสฺมา นนฺทโก ภควโต ปจฺจสฺโสสิ.\csromanspacer{} อถ โข อายสฺมา นนฺทโก ตสฺสา รตฺติยา อจฺจเยน ปุพฺเพณฺหสมยํ นิวาเสตฺวา ปตฺตจีวรมาทาย สาวตฺถิํ ปิณฺฑาย ปาวิสิ, สาวตฺถิยํ ปิณฺฑาย จริตฺวา ปจฺฉาภตฺตํ ปิณฺฑปาตปฏิกฺกนฺโต อตฺตทุติโย เยน ราชการาโม เตนุปสงฺกมิ.\csromanspacer{} อทฺทสํสุ โข ตา ภิกฺขุนิโย อายสฺมนฺตํ นนฺทกํ ทูรโตว อาคจฺฉนฺตํ, ทิสฺวาน อาสนํ ปญฺญาเปสุํ อุทกญฺจ ปาทานํ อุปฏฺฐเปสุํ.\csromanspacer{} นิสีทิ โข อายสฺมา นนฺทโก ปญฺญตฺเต อาสเน, นิสชฺช ปาเท ปกฺขาเลสิ.\csromanspacer{} ตาปิ โข ภิกฺขุนิโย อายสฺมนฺตํ นนฺทกํ อภิวาเทตฺวา เอกมนฺตํ นิสีทิํสุ, เอกมนฺตํ นิสินฺนา โข ตา ภิกฺขุนิโย อายสฺมา นนฺทโก เอตทโวจ “ปฏิปุจฺฉกถา โข ภคินิโย ภวิสฺสติ, ตตฺถ อาชานนฺตีหิ ‘อาชานามา’ติสฺส วจนียํ, น อาชานนฺตีหิ ‘น อาชานามา’ติสฺส วจนียํ, ยสฺสา วา ปนสฺส กงฺขา วา วิมติ วา, อหเมว ตตฺถ ปฏิปุจฺฉิตพฺโพ ‘อิทํ ภนฺเต กถํ อิมสฺส กฺวตฺโถ’ติ”.\csromanspacer{} เอตฺตเกนปิ มยํ ภนฺเต อยฺยสฺส นนฺทกสฺส อตฺตมนา อภิรทฺธา, ยํ โน อยฺโย นนฺทโก ปวาเรตีติ.}

\proseitem{408}{ตํ กิํ มญฺญถ ภคินิโย, จกฺขุ นิจฺจํ วา อนิจฺจํ วาติ.\csromanspacer{} อนิจฺจํ ภนฺเต.\csromanspacer{} ยํ ปนานิจฺจํ, ทุกฺขํ วา ตํ สุขํ วาติ.\csromanspacer{} ทุกฺขํ ภนฺเต.\csromanspacer{} ยํ ปนานิจฺจํ ทุกฺขํ วิปริณามธมฺมํ, กลฺลํ นุ ตํ สมนุปสฺสิตุํ “เอตํ มม, เอโสหมสฺมิ, เอโส เม อตฺตา”ติ.\csromanspacer{} โน เหตํ ภนฺเต.\csromanspacer{} ตํ กิํ มญฺญถ ภคินิโย, โสตํ นิจฺจํ วา อนิจฺจํ วาติ.\csromanspacer{} อนิจฺจํ ภนฺเต ฯเปฯ.\csromanspacer{} ฆานํ นิจฺจํ วา อนิจฺจํ วาติ.\csromanspacer{} อนิจฺจํ ภนฺเต.\csromanspacer{} ชิวฺหา.\csromanspacer{} กาโย.\csromanspacer{} มโน นิจฺโจ วา อนิจฺโจ วาติ.\csromanspacer{} อนิจฺโจ ภนฺเต.\csromanspacer{} ยํ ปนานิจฺจํ, ทุกฺขํ วา ตํ สุขํ วาติ.\csromanspacer{} ทุกฺขํ ภนฺเต.\csromanspacer{} ยํ ปนานิจฺจํ ทุกฺขํ วิปริณามธมฺมํ, กลฺลํ นุ ตํ สมนุปสฺสิตุํ “เอตํ มม, เอโสหมสฺมิ, เอโส เม อตฺตา”ติ.\csromanspacer{} โน เหตํ ภนฺเต.\csromanspacer{} ตํ กิสฺส เหตุ, ปุพฺเพว โน เอตํ ภนฺเต ยถาภูตํ สมฺมปฺปญฺญาย สุทิฏฺฐํ “อิติปิเม ฉ อชฺฌตฺติกา อายตนา อนิจฺจา”ติ.\csromanspacer{} สาธุ สาธุ ภคินิโย, เอวํ เหตํ ภคินิโย โหติ อริยสาวกสฺส ยถาภูตํ สมฺมปฺปญฺญาย ปสฺสโต.\csromanspacer{} ตํ กิํ มญฺญถ ภคินิโย, รูปา นิจฺจา วา อนิจฺจา วาติ.\csromanspacer{} อนิจฺจา ภนฺเต.\csromanspacer{} ยํ ปนานิจฺจํ, ทุกฺขํ วา ตํ สุขํ วาติ.\csromanspacer{} ทุกฺขํ ภนฺเต.\csromanspacer{} ยํ ปนานิจฺจํ ทุกฺขํ วิปริณามธมฺมํ, กลฺลํ นุ ตํ สมนุปสฺสิตุํ “เอตํ มม, เอโสหมสฺมิ,}

\csromanheader{2}{4. นนฺทโกวาทสุตฺต (146)}
\prose{\csromanfolio{321}เอโส เม อตฺตา”ติ.\csromanspacer{} โน เหตํ ภนฺเต.\csromanspacer{} ตํ กิํ มญฺญถ ภคินิโย, สทฺทา นิจฺจา วา อนิจฺจา วาติ.\csromanspacer{} อนิจฺจา ภนฺเต ฯเปฯ.\csromanspacer{} คนฺธา นิจฺจา วา อนิจฺจา วาติ.\csromanspacer{} อนิจฺจา ภนฺเต.\csromanspacer{} รสา นิจฺจา วา อนิจฺจา วาติ.\csromanspacer{} อนิจฺจา ภนฺเต.\csromanspacer{} โผฏฺฐพฺพา นิจฺจา วา อนิจฺจา วาติ.\csromanspacer{} อนิจฺจา ภนฺเต.\csromanspacer{} ธมฺมา นิจฺจา วา อนิจฺจา วาติ.\csromanspacer{} อนิจฺจา ภนฺเต.\csromanspacer{} ยํ ปนานิจฺจํ, ทุกฺขํ วา ตํ สุขํ วาติ.\csromanspacer{} ทุกฺขํ ภนฺเต.\csromanspacer{} ยํ ปนานิจฺจํ ทุกฺขํ วิปริณามธมฺมํ, กลฺลํ นุ ตํ สมนุปสฺสิตุํ “เอตํ มม, เอโสหมสฺมิ, เอโส เม อตฺตา”ติ.\csromanspacer{} โน เหตํ ภนฺเต.\csromanspacer{} ตํ กิสฺส เหตุ, ปุพฺเพว โน เอตํ ภนฺเต ยถาภูตํ สมฺมปฺปญฺญาย สุทิฏฺฐํ “อิติปิเม ฉ พาหิรา อายตนา อนิจฺจา”ติ.\csromanspacer{} สาธุ สาธุ ภคินิโย, เอวํ เหตํ ภคินิโย โหติ อริยสาวกสฺส ยถาภูตํ สมฺมปฺปญฺญาย ปสฺสโต.}

\proseitem{410}{ตํ กิํ มญฺญถ ภคินิโย, จกฺขุวิญฺญาณํ นิจฺจํ วา อนิจฺจํ วาติ.\csromanspacer{} อนิจฺจํ ภนฺเต ฯเปฯ.\csromanspacer{} โสตวิญฺญาณํ นิจฺจํ วา อนิจฺจํ วาติ.\csromanspacer{} อนิจฺจํ ภนฺเต.\csromanspacer{} ฆานวิญฺญาณํ นิจฺจํ วา อนิจฺจํ วาติ.\csromanspacer{} อนิจฺจํ ภนฺเต.\csromanspacer{} ชิวฺหาญฺญาณํ นิจฺจํ วา อนิจฺจํ วาติ.\csromanspacer{} อนิจฺจํ ภนฺเต.\csromanspacer{} กายวิญฺญาณํ นิจฺจํ วา อนิจฺจํ วาติ.\csromanspacer{} อนิจฺจํ ภนฺเต.\csromanspacer{} มโนวิญฺญาณํ นิจฺจํ วา อนิจฺจํ วาติ.\csromanspacer{} อนิจฺจํ ภนฺเต.\csromanspacer{} ยํ ปนานิจฺจํ, ทุกฺขํ วา ตํ สุขํ วาติ.\csromanspacer{} ทุกฺขํ ภนฺเต.\csromanspacer{} ยํ ปนานิจฺจํ ทุกฺขํ วิปริณามธมฺมํ, กลฺลํ นุ ตํ สมนุปสฺสิตุํ “เอตํ มม, เอโสหมสฺมิ, เอโส เม อตฺตา”ติ.\csromanspacer{} โน เหตํ ภนฺเต.\csromanspacer{} ตํ กิสฺส เหตุ, ปุพฺเพว โน เอตํ ภนฺเต ยถาภูตํ สมฺมปฺปญฺญาย สุทิฏฺฐํ “อิติปิเม ฉ วิญฺญาณกายา อนิจฺจา”ติ.\csromanspacer{} สาธุ สาธุ ภคินิโย, เอวํ เหตํ ภคินิโย โหติ อริยสาวกสฺส ยถาภูตํ สมฺมปฺปญฺญาย ปสฺสโต.}

\proseitem{411}{เสยฺยถาปิ ภคินิโย เตลปฺปทีปสฺส ฌายโต เตลมฺปิ อนิจฺจํ วิปริณามธมฺมํ, วฏฺฏิปิ อนิจฺจา วิปริณามธมฺมา, อจฺจิปิ อนิจฺจา วิปริณามธมฺมา, อาภาปิ อนิจฺจา วิปริณามธมฺมา.\csromanspacer{} โย นุ โข ภคินิโย เอวํ วเทยฺย “อมุสฺส เตลปฺปทีปสฺส ฌายโต เตลมฺปิ อนิจฺจํ วิปริณามธมฺมํ, วฏฺฏิปิ อนิจฺจา วิปริณามธมฺมา, อจฺจิปิ อนิจฺจา วิปริณามธมฺมา.\csromanspacer{} ยา จ ขฺวาสฺส อาภา, สา นิจฺจา ธุวา สสฺสตา อวิปริณามธมฺมา”ติ, สมฺมา นุ โข โส ภคินิโย วทมาโน วเทยฺยาติ.\csromanspacer{} โน เหตํ ภนฺเต.\csromanspacer{} ตํ กิสฺส เหตุ, อมุสฺส หิ ภนฺเต เตลปฺปทีปสฺส ฌายโต เตลมฺปิ อนิจฺจํ วิปริณามธมฺมํ, วฏฺฏิปิ อนิจฺจา วิปริณามธมฺมา, อจฺจิปิ อนิจฺจา \csromanfolio{322}วิปริณามธมฺมา.\csromanspacer{} ปเควสฺส อาภา อนิจฺจา วิปริณามธมฺมาติ.\csromanspacer{} เอวเมว โข ภคินิโย, โย นุ โข เอวํ วเทยฺย “ฉ โขเม อชฺฌตฺติกา อายตนา อนิจฺจา.\csromanspacer{} ยญฺจ โข ฉ อชฺฌตฺติเก อายตเน ปฏิจฺจ ปฏิสํเวเทติ สุขํ วา ทุกฺขํ วา อทุกฺขมสุขํ วา, ตํ นิจฺจํ ธุวํ สสฺสตํ อวิปริณามธมฺม นฺ”ติ, สมฺมา นุ โข โส ภคินิโย วทมาโน วเทยฺยาติ.\csromanspacer{} โน เหตํ ภนฺเต.\csromanspacer{} ตํ กิสฺส เหตุ, ตชฺชํ ตชฺชํ ภนฺเต ปจฺจยํ ปฏิจฺจ ตชฺชา ตชฺชา เวทนา อุปฺปชฺชนฺติ, ตชฺชสฺส ตชฺชสฺส ปจฺจยสฺส นิโรธา ตชฺชา ตชฺชา เวทนา นิรุชฺฌนฺตีติ.\csromanspacer{} สาธุ สาธุ ภคินิโย, เอวํ เหตํ ภคินิโย โหติ อริยสาวกสฺส ยถาภูตํ สมฺมปฺปญฺญาย ปสฺสโต.}

\proseitem{412}{เสยฺยถาปิ ภคินิโย มหโต รุกฺขสฺส ติฏฺฐโต สารวโต มูลมฺปิ อนิจฺจํ วิปริณามธมฺมํ, ขนฺโธปิ อนิจฺโจ วิปริณามธมฺโม, สาขาปลาสมฺปิ อนิจฺจํ วิปริณามธมฺมํ, ฉายาปิ อนิจฺจา วิปริณามธมฺมา.\csromanspacer{} โย นุ โข ภคินิโย เอวํ วเทยฺย “อมุสฺส มหโต รุกฺขสฺส ติฏฺฐโต สารวโต มูลมฺปิ อนิจฺจํ วิปริณามธมฺมํ, ขนฺโธปิ อนิจฺโจ วิปริณามธมฺโม, สาขาปลาสมฺปิ อนิจฺจํ วิปริณามธมฺมํ, ยา จ ขฺวาสฺส ฉายา, สา นิจฺจา ธุวา สสฺสตา อวิปริณามธมฺมา”ติ, สมฺมา นุ โข โส ภคินิโย วทมาโน วเทยฺยาติ.\csromanspacer{} โน เหตํ ภนฺเต.\csromanspacer{} ตํ กิสฺส เหตุ, อมุสฺส หิ ภนฺเต มหโต รุกฺขสฺส ติฏฺฐโต สารวโต มูลมฺปิ อนิจฺจํ วิปริณามธมฺมํ, ขนฺโธปิ อนิจฺโจ วิปริณามธมฺโม, สาขาปลาสมฺปิ อนิจฺจํ วิปริณามธมฺมํ.\csromanspacer{} ปเควสฺส ฉายา อนิจฺจา วิปริณามธมฺมาติ.\csromanspacer{} เอวเมว โข ภคินิโย, โย นุ โข เอวํ วเทยฺย “ฉ โขเม พาหิรา อายตนา อนิจฺจา.\csromanspacer{} ยญฺจ โข ฉ พาหิเร อายตเน ปฏิจฺจ ปฏิสํเวเทติ สุขํ วา ทุกฺขํ วา อทุกฺขมสุขํ วา, ตํ นิจฺจํ ธุวํ สสฺสตํ อวิปริณามธมฺมนฺ”ติ, สมฺมา นุ โข โส ภคินิโย วทมาโน วเทยฺยาติ.\csromanspacer{} โน เหตํ ภนฺเต.\csromanspacer{} ตํ กิสฺส เหตุ, ตชฺชํ ตชฺชํ ภนฺเต ปจฺจยํ ปฏิจฺจ ตชฺชา ตชฺชา เวทนา อุปฺปชฺชนฺติ, ตชฺชสฺส ตชฺชสฺส ปจฺจยสฺส นิโรธา ตชฺชา ตชฺชา เวทนา นิรุชฺฌนฺตีติ.\csromanspacer{} สาธุ สาธุ ภคินิโย, เอวํ เหตํ ภคินิโย โหติ อริยสาวกสฺส ยถาภูตํ สมฺมปฺปญฺญาย ปสฺสโต.}

\proseitem{413}{เสยฺยถาปิ ภคินิโย ทกฺโข โคฆาตโก วา โคฆาตกนฺเตวาสี วา คาวิํ วธิตฺวา ติเณฺหน โควิกนฺตเนน คาวิํ สงฺกนฺเตยฺย}

\csromanheader{2}{4. นนฺทโกวาทสุตฺต (146)}
\prose{\csromanfolio{323}อนุปหจฺจ อนฺตรํ มํสกายํ อนุปหจฺจ พาหิรํ จมฺมกายํ.\csromanspacer{} ยํ ยเทว ตตฺถ อนฺตรา วิลิมํสํ อนฺตรา นฺหารุ อนฺตรา พนฺธนํ, ตํ ตเทว ติเณฺหน โควิกนฺตเนน สญฺฉินฺเทยฺย สงฺกนฺเตยฺย สมฺปกนฺเตยฺย สมฺปริกนฺเตยฺย, สญฺฉินฺทิตฺวา สงฺกนฺติตฺวา สมฺปกนฺติตฺวา สมฺปริกนฺติตฺวา วิธุนิตฺวา พาหิรํ จมฺมกาย เตนว จมฺเมน ตํ คาวิํ ปฏิจฺฉาเทตฺวา เอวํ วเทยฺย “ตเถวายํ คาวี สํยุตฺตา อิมินาว จมฺเมนา”ติ, สมฺมา นุ โข โส ภคินิโย วทมาโน วเทยฺยาติ.\csromanspacer{} โน เหตํ ภนฺเต.\csromanspacer{} ตํ กิสฺส เหตุ, อมุ หิ ภนฺเต ทกฺโข โคฆาตโก วา โคฆาตกนฺเตวาสี วา คาวิํ วธิตฺวา ติเณฺหน โควิกนฺตเนน คาวิํ สงฺกนฺเตยฺย อนุปหจฺจ อนฺตรํ มํสกายํ อนุปหจฺจ พาหิรํ จมฺมกายํ.\csromanspacer{} ยํ ยเทว ตตฺถ อนฺตรา วิลิมํสํ อนฺตรา นฺหารุ อนฺตรา พนฺธนํ, ตํ ตเทว ติเณฺหน โควิกนฺตเนน สญฺฉินฺเทยฺย สงฺกนฺเตยฺย สมฺปกนฺเตยฺย สมฺปรินฺเตยฺย, สญฺฉินฺทิตฺวา สงฺกนฺติตฺวา สมฺปกนฺติตฺวา สมฺปริกนฺติตฺวา วิธุนิตฺวา พาหิรํ จมฺมกายํ เตนว จมฺเมน ตํ คาวิํ ปฏิจฺฉาเทตฺวา กิญฺจาปิ โส เอวํ วเทยฺย “ตเถวายํ คาวี สํยุตฺตา อิมินาว จมฺเมนา”ติ, อถ โข สา คาวี วิสํยุตฺตา เตเนว จมฺเมนาติ.}

\prose{อุปมา โข เม อยํ ภคินิโย กตา อตฺถสฺส วิญฺญาปนาย, อยเมเวตฺถ อตฺโถ. “อนฺตรา มํสกาโย”ติ โข ภคินิโย ฉนฺเนตํ อชฺฌตฺติกานํ อายตนานํ อธิวจนํ. “พาหิโร จมฺมกาโย”ติ โข ภคินิโย ฉนฺเนตํ พาหิรานํ อายตนานํ อธิวจนํ. “อนฺตรา วิลิมํสํ อนฺตรา นฺหารุ อนฺตรา พนฺธนนฺ”ติ โข ภคินิโย นนฺทีราคสฺเสตํ อธิวจนํ. “ติณฺหํ โควิกนฺตนนฺ”ติ โข ภคินิโย อริยาเยตํ ปญฺญาย อธิวจนํ, ยายํ อริยา ปญฺญา อนฺตรา กิเลสํ อนฺตรา สํโยชนํ อนฺตรา พนฺธนํ สญฺฉินฺทติ สงฺกนฺตติ สมฺปกนฺตติ สมฺปริกนฺตติ.}

\proseitem{414}{สตฺต โข ปนิเม ภคินิโย โพชฺฌงฺคา, เยสํ ภาวิตตฺตา พหุลีกตตฺตา ภิกฺขุ อาสวานํ ขยา อนาสวํ เจโตวิมุตฺติํ ปญฺญาวิมุตฺติํ ทิฏฺเฐว ธมฺเม สยํ อภิญฺญา สจฺฉิกตฺวา อุปสมฺปชฺช วิหรติ.\csromanspacer{} กตเม สตฺต, อิธ ภคินิโย ภิกฺขุ สติสมฺโพชฺฌงฺคํ ภาเวติ วิเวกนิสฺสิตํ วิราคนิสฺสิตํ นิโรธนิสฺสิตํ โวสฺสคฺคปริณามิํ.\csromanspacer{} ธมฺมวิจยสมฺโพชฺฌงฺคํ ภาเวติ ฯเปฯ.\csromanspacer{} วีริยสมฺโพชฺฌงฺคํ ภาเวติ.\csromanspacer{} ปีติสมฺโพชฺฌงฺคํ ภาเวติ.\csromanspacer{} ปสฺสทฺธิสมฺโพชฺฌงฺคํ ภาเวติ.\csromanspacer{} สมาธิสมฺโพชฺฌงฺคํ ภาเวติ. \csromanfolio{324}อุเปกฺขาสมฺโพชฺฌงฺคํ ภาเวติ วิเวกนิสฺสิตํ วิราคนิสฺสิตํ นิโรธนิสฺสิตํ โวสฺสคฺคปริณามิํ.\csromanspacer{} อิเม โข ภคินิโย สตฺต โพชฺฌงฺคา.\csromanspacer{} เยสํ ภาวิตตฺตา พหุลีกตตฺตา ภิกฺขุ อาสวานํ ขยา อนาสวํ เจโตวิมุตฺติํ ปญฺญาวิมุตฺติํ ทิฏฺเฐว ธมฺเม สยํ อภิญฺญา สจฺฉิกตฺวา อุปสมฺปชฺช วิหรตีติ.}

\proseitem{415}{อถ โข อายสฺมา นนฺทโก ตา ภิกฺขุนิโย อิมินา โอวาเทน โอวทิตฺวา อุยฺโยเชสิ “คจฺฉถ ภคินิโย กาโล”ติ.\csromanspacer{} อถ โข ตา ภิกฺขุนิโย อายสฺมโต นนฺทกสฺส ภาสิตํ อภินนฺทิตฺวา อนุโมทิตฺวา อุฏฺฐายาสนา อายสฺมนฺตํ นนฺทกํ อภิวาเทตฺวา ปทกฺขิณํ กตฺวา เยน ภควา เตนุปสงฺกมิํสุ, อุปสงฺกมิตฺวา ภควนฺตํ อภิวาเทตฺวา เอกมนฺตํ อฏฺฐํสุ, เอกมนฺตํ ฐิตา โข ตา ภิกฺขุนิโย ภควา เอตทโวจ “คจฺฉถ ภิกฺขุนิโย กาโล”ติ.\csromanspacer{} อถ โข ตา ภิกฺขุนิโย ภควนฺตํ อภิวาเทตฺวา ปทกฺขิณํ กตฺวา ปกฺกมิํสุ.\csromanspacer{} อถ โข ภควา อจิรปกฺกนฺตีสุ ตาสุ ภิกฺขุนีสุ ภิกฺขู อามนฺเตสิ “เสยฺยถาปิ ภิกฺขเว ตทหุโปสเถ ปนฺนรเส น โหติ พหุโน ชนสฺส กงฺขา วา วิมติ วา ‘อูโน นุ โข จนฺโท, ปุณฺโณ นุ โข จนฺโท’ติ, อถ โข ปุณฺโณ จนฺโทเตฺวว โหติ.\csromanspacer{} เอวเมว โข ภิกฺขเว ตา ภิกฺขุนิโย นนฺทกสฺส ธมฺมเทสนาย อตฺตมนา เจว ปริปุณฺณาสงฺกปฺปา จ.\csromanspacer{} ตาสํ ภิกฺขเว ปญฺจนฺนํ ภิกฺขุนิสตานํ ยา ปจฺฉิมิกา ภิกฺขุนี, สา\footnote{ยา ปจฺฉิมา ภิกฺขุนี, สา (สี, สฺยา, กํ, อิ), ยา ปจฺฉิมิกา, ตา ภิกฺขุนิโย (ก)} โสตาปนฺนา อวินิปาตธมฺมา นิยตา สมฺโพธิปรายนา”ติ.}

\prose{อิทมโวจ ภควา.\csromanspacer{} อตฺตมนา เต ภิกฺขู ภควโต ภาสิตํ อภินนฺทุนฺติ.}

\nitthitamruled{นนฺทโกวาทสุตฺตํ นิฏฺฐิตํ จตุตฺถํ.}
\csromanmatikaanchor{mtk.56}
\csromantocmarkref{subsection}{5. จูฬราหุโลวาทสุตฺต}{mtk.56}
\bookmark[dest={mtk.56},level=2]{5. จูฬราหุโลวาทสุตฺต}
\csromanheader{2}{5. จูฬราหุโลวาทสุตฺต}
\proseitem{416}{เอวํ เม สุตํ\csromandash{}เอกํ สมยํ ภควา สาวตฺถิยํ วิหรติ เชตวเน อนาถปิณฺฑิกสฺส อาราเม.\csromanspacer{} อถ โข ภควโต รโหคตสฺส ปฏิสลฺลีนสฺส เอวํ เจตโส ปริวิตกฺโก อุทปาทิ “ปริปกฺกา}

\csromanheader{2}{5. จูฬราหุโลวาทสุตฺต (147)}
\prose{\csromanfolio{325}โข ราหุลสฺส วิมุตฺติปริปาจนียา ธมฺมา, ยํนูนาหํ ราหุลํ อุตฺตริํ อาสวานํ ขเย วิเนยฺยนฺ”ติ.\csromanspacer{} อถ โข ภควา ปุพฺพณฺหสมยํ นิวาเสตฺวา ปตฺตจีวรมาทาย สาวตฺถิํ ปิณฺฑาย ปาวิสิ, สาวตฺถิยํ ปิณฺฑาย จริตฺวา ปจฺฉาภตฺตํ ปิณฺฑปาตปฏิกฺกนฺโต อายสฺมนฺตํ ราหุลํ อามนฺเตสิ “คณฺหาหิ ราหุล นิสีทนํ, เยน อนฺธวนํ เตนุปสงฺกมิสฺสาม ทิวาวิหารายา”ติ. “เอวํ ภนฺเต”ติ โข อายสฺมา ราหุโล ภควโต ปฏิสฺสุตฺวา นิสีทนํ อาทาย ภควนฺตํ ปิฏฺฐิโต ปิฏฺฐิโต อนุพนฺธิ.}

\prose{เตน โข ปน สมเยน อเนกานิ เทวตาสหสฺสานิ ภควนฺตํ อนุพนฺธานิ โหนฺติ “อชฺช ภควา อายสฺมนฺตํ ราหุลํ อุตฺตริํ อาสวานํ ขเย วิเนสฺสตี”ติ.\csromanspacer{} อถ โข ภควา อนฺธวนํ อชฺโฌคาเหตฺวา อญฺญตรสฺมิํ รุกฺขมูเล ปญฺญตฺเต อาสเน นิสีทิ, อายสฺมาปิ โข ราหุโล ภควนฺตํ อภิวาเทตฺวา เอกมนฺตํ นิสีทิ, เอกมนฺตํ นิสินฺนํ โข อายสฺมนฺตํ ราหุลํ ภควา เอตทโวจ\csromandash{}}

\proseitem{417}{ตํ กิํ มญฺญสิ ราหุล, จกฺขุ นิจฺจํ วา อนิจฺจํ วาติ.\csromanspacer{} อนิจฺจํ ภนฺเต.\csromanspacer{} ยํ ปนานิจฺจํ, ทุกฺขํ วา ตํ สุขํ วาติ.\csromanspacer{} ทุกฺขํ ภนฺเต.\csromanspacer{} ยํ ปนานิจฺจํ ทุกฺขํ วิปริณามธมฺมํ, กลฺลํ นุ ตํ สมนุปสฺสิตุํ “เอตํ มม, เอโสหมสฺมิ, เอโส เม อตฺตา”ติ.\csromanspacer{} โน เหตํ ภนฺเต.\csromanspacer{} ตํ กิํ มญฺญสิ ราหุล, รูปา นิจฺจา วา อนิจฺจา วาติ.\csromanspacer{} อนิจฺจา ภนฺเต.\csromanspacer{} ยํ ปนานิจฺจํ, ทุกฺขํ วา ตํ สุขํ วาติ.\csromanspacer{} ทุกฺขํ ภนฺเต.\csromanspacer{} ยํ ปนานิจฺจํ ทุกฺขํ วิปริณามธมฺมํ, กลฺลํ นุ ตํ สมนุปสฺสิตุํ “เอตํ มม, เอโสหมสฺมิ, เอโส เม อตฺตา”ติ.\csromanspacer{} โน เหตํ ภนฺเต.\csromanspacer{} ตํ กิํ มญฺญสิ ราหุล, จกฺขุวิญฺญาณํ นิจฺจํ วา อนิจฺจํ วาติ.\csromanspacer{} อนิจฺจํ ภนฺเต.\csromanspacer{} ยํ ปนานิจฺจํ, ทุกฺขํ วา ตํ สุขํ วาติ.\csromanspacer{} ทุกฺขํ ภนฺเต.\csromanspacer{} ยํ ปนานิจฺจํ ทุกฺขํ วิปริณามธมฺมํ, กลฺลํ นุ ตํ สมนุปสฺสิตุํ “เอตํ มม, เอโสหมสฺมิ, เอโส เม อตฺตา”ติ.\csromanspacer{} โน เหตํ ภนฺเต.\csromanspacer{} ตํ กิํ มญฺญสิ ราหุล, จกฺขุสมฺผสฺโส นิจฺโจ วา อนิจฺโจ วาติ.\csromanspacer{} อนิจฺโจ ภนฺเต.\csromanspacer{} ยํ ปนานิจฺจํ, ทุกฺขํ วา ตํ สุขํ วาติ.\csromanspacer{} ทุกฺขํ ภนฺเต.\csromanspacer{} ยํ ปนานิจฺจํ ทุกฺขํ วิปริณามธมฺมํ, กลฺลํ นุ ตํ สมนุปสฺสิตุํ “เอตํ มม, เอโสหมสฺมิ, เอโส เม อตฺตา”ติ.\csromanspacer{} โน เหตํ ภนฺเต.\csromanspacer{} ตํ กิํ มญฺญสิ ราหุล, ยมิทํ\footnote{ยมฺปิทํ (สี, ก)} จกฺขุสมฺผสฺสปจฺจยา อุปฺปชฺชติ เวทนา- คตํ สญฺญาคตํ สงฺขารคตํ วิญฺญาณคตํ, ตมฺปิ นิจฺจํ วา อนิจฺจํ \csromanfolio{326}วาติ.\csromanspacer{} อนิจฺจํ ภนฺเต.\csromanspacer{} ยํ ปนานิจฺจํ, ทุกฺขํ วา ตํ สุขํ วาติ.\csromanspacer{} ทุกฺขํ ภนฺเต.\csromanspacer{} ยํ ปนานิจฺจํ ทุกฺขํ วิปริณามธมฺมํ, กลฺลํ นุ ตํ สมนุปสฺสิตุํ “เอตํ มม, เอโสหมสฺมิ, เอโส เม อตฺตา”ติ.\csromanspacer{} โน เหตํ ภนฺเต.}

\proseitem{418}{ตํ กิํ มญฺญสิ ราหุลา, โสตํ นิจฺจํ วา อนิจฺจํ วาติ.\csromanspacer{} อนิจฺจํ ภนฺเต ฯเปฯ.\csromanspacer{} ฆานํ นิจฺจํ วา อนิจฺจํ วาติ.\csromanspacer{} อนิจฺจํ ภนฺเต.\csromanspacer{} ชิวฺหา นิจฺจา วา อนิจฺจา วาติ.\csromanspacer{} อนิจฺจา ภนฺเต.\csromanspacer{} กาโย นิจฺโจ วา อนิจฺโจ วาติ.\csromanspacer{} อนิจฺโจ ภนฺเต.\csromanspacer{} มโน นิจฺโจ วา อนิจฺโจ วาติ.\csromanspacer{} อนิจฺโจ ภนฺเต.\csromanspacer{} ยํ ปนานิจฺจํ, ทุกฺขํ วา ตํ สุขํ วาติ.\csromanspacer{} ทุกฺขํ ภนฺเต.\csromanspacer{} ยํ ปนานิจฺจํ ทุกฺขํ วิปริณามธมฺมํ, กลฺลํ นุ ตํ สมนุปสฺสิตุํ “เอตํ มม, เอโสหมสฺมิ, เอโส เม อตฺตา”ติ.\csromanspacer{} โน เหตํ ภนฺเต.\csromanspacer{} ตํ กิํ มญฺญสิ ราหุล, ธมฺมา นิจฺจา วา อนิจฺจา วาติ.\csromanspacer{} อนิจฺจา ภนฺเต.\csromanspacer{} ยํ ปนานิจฺจํ, ทุกฺขํ วา ตํ สุขํ วาติ.\csromanspacer{} ทุกฺขํ ภนฺเต.\csromanspacer{} ยํ ปนานิจฺจํ ทุกฺขํ วิปริณามธมฺมํ, กลฺลํ นุ ตํ สมนุปสฺสิตุํ “เอตํ มม, เอโสหมสฺมิ, เอโส เม อตฺตา”ติ.\csromanspacer{} โน เหตํ ภนฺเต.\csromanspacer{} ตํ กิํ มญฺญสิ ราหุล, มโนวิญฺญาณํ นิจฺจํ วา อนิจฺจํ วาติ.\csromanspacer{} อนิจฺจํ ภนฺเต.\csromanspacer{} ยํ ปนานิจฺจํ, ทุกฺขํ วา ตํ สุขํ วาติ.\csromanspacer{} ทุกฺขํ ภนฺเต.\csromanspacer{} ยํ ปนานิจฺจํ ทุกฺขํ วิปริณามธมฺมํ, กลฺลํ นุ ตํ สมนุปสฺสิตุํ “เอตํ มม, เอโสหมสฺมิ, เอโส เม อตฺตา”ติ.\csromanspacer{} โน เหตํ ภนฺเต.\csromanspacer{} ตํ กิํ มญฺญสิ ราหุล, มโนสมฺผสฺโส นิจฺโจ วา อนิจฺโจ วาติ.\csromanspacer{} อนิจฺโจ ภนฺเต.\csromanspacer{} ยํ ปนานิจฺจํ, ทุกฺขํ วา ตํ สุขํ วาติ.\csromanspacer{} ทุกฺขํ ภนฺเต.\csromanspacer{} ยํ ปนานิจฺจํ ทุกฺขํ วิปริณามธมฺมํ, กลฺลํ นุ ตํ สมนุปสฺสิตุํ “เอตํ มม, เอโสหมสฺมิ, เอโส เม อตฺตา”ติ.\csromanspacer{} โน เหตํ ภนฺเต.\csromanspacer{} ตํ กิํ มญฺญสิ ราหุล, ยมิทํ มโนสมฺผสฺสปจฺจยา อุปฺปชฺชติ เวทนาคตํ สญฺญาคตํ สงฺขารคตํ วิญฺญาณคตํ, ตมฺปิ นิจฺจํ วา อนิจฺจํ วาติ.\csromanspacer{} อนิจฺจํ ภนฺเต.\csromanspacer{} ยํ ปนานิจฺจํ, ทุกฺขํ วา ตํ สุขํ วาติ.\csromanspacer{} ทุกฺขํ ภนฺเต.\csromanspacer{} ยํ ปนานิจฺจํ ทุกฺขํ วิปริณามธมฺมํ, กลฺลํ นุ ตํ สมนุปสฺสิตุํ “เอตํ มม, เอโสหมสฺมิ, เอโส เม อตฺตา”ติ.\csromanspacer{} โน เหตํ ภนฺเต.}

\proseitem{419}{เอวํ ปสฺสํ ราหุล สุตวา อริยสาวโก จกฺขุสฺมิํ\footnote{จกฺขุสฺมิมฺปิ (สฺยา, กํ) เอวมิตเรสุปิ.} นิพฺพินฺทติ, รูเปสุ นิพฺพินฺทติ, จกฺขุวิญฺญาเณ นิพฺพินฺทติ, จกฺขุสมฺผสฺเส นิพฺพินฺทติ.\csromanspacer{} ยมิทํ จกฺขุสมฺผสฺสปจฺจยา อุปฺปชฺชติ เวทนาคตํ สญฺญาคตํ สงฺขารคตํ วิญฺญาณคตํ, ตสฺมิมฺปิ นิพฺพินฺทติ.\csromanspacer{} โสตสฺมิํ นิพฺพินฺทติ, สทฺเทสุ}

\csromanheader{2}{6. ฉฉกฺกสุตฺต (148)}
\prose{\csromanfolio{327}นิพฺพินฺทติ ฯเปฯ.\csromanspacer{} ฆานสฺมิํ นิพฺพินฺทติ, คนฺเธสุ นิพฺพินฺทติ.\csromanspacer{} ชิวฺหาย นิพฺพินฺทติ, รเสสุ นิพฺพินฺทติ.\csromanspacer{} กายสฺมิํ นิพฺพินฺทติ, โผฏฺฐพฺเพสุ นิพฺพินฺทติ.\csromanspacer{} มนสฺมิํ นิพฺพินฺทติ, ธมฺเมสุ นิพฺพินฺทติ, มโนวิญฺญาเณ นิพฺพินฺทติ, มโนสมฺผสฺเส นิพฺพินฺทติ.\csromanspacer{} ยมิทํ มโนสมฺผสฺสปจฺจยา อุปฺปชฺชติ เวทนาคตํ สญฺญาคตํ สงฺขารคตํ วิญฺญาณคตํ ตสฺมิมฺปิ นิพฺพินฺทติ.\csromanspacer{} นิพฺพินฺทํ วิรชฺชติ, วิราคา วิมุจฺจติ, วิมุตฺตสฺมิํ “วิมุตฺตมฺ”อิติ ญาณํ โหติ, “ขีณา ชาติ, วุสิตํ พฺรหฺมจริยํ, กตํ กรณียํ, นาปรํ อิตฺถตฺตายา”ติ ปชานาตีติ.}

\prose{อิทมโวจ ภควา.\csromanspacer{} อตฺตมโน อายสฺมา ราหุโล ภควโต ภาสิตํ อภินนฺทีติ.\csromanspacer{} อิมสฺมิํ จ ปน เวยฺยากรณสฺมิํ ภญฺญมาเน อายสฺมโต ราหุลสฺส อนุปาทาย อาสเวหิ จิตฺตํ วิมุจฺจิ, ตาสญฺจ อเนกานํ เทวตาสหสฺสานํ วิรชํ วีตมลํ ธมฺมจกฺขุํ อุทปาทิ “ยํ กิญฺจิ สมุทยธมฺมํ, สพฺพํ ตํ นิโรธธมฺมนฺ”ติ.}

\nitthitamruled{จูฬราหุโลวาทสุตฺตํ นิฏฺฐิตํ ปญฺจมํ.}
\csromanmatikaanchor{mtk.57}
\csromantocmarkref{subsection}{6. ฉฉกฺกสุตฺต}{mtk.57}
\bookmark[dest={mtk.57},level=2]{6. ฉฉกฺกสุตฺต}
\csromanheader{2}{6. ฉฉกฺกสุตฺต}
\proseitem{420}{เอวํ เม สุตํ\csromandash{}เอกํ สมยํ ภควา สาวตฺถิยํ วิหรติ เชตวเน อนาถปิณฺฑิกสฺส อาราเม.\csromanspacer{} ตตฺร โข ภควา ภิกฺขู อามนฺเตสิ “ภิกฺขโว”ติ. “ภทนฺเต”ติ เต ภิกฺขู ภควโต ปจฺจสฺโสสุํ.\csromanspacer{} ภควา เอตทโวจ “ธมฺมํ โว ภิกฺขเว เทเสสฺสามิ อาทิกลฺยาณํ มชฺเฌกลฺยาณํ ปริโยสานกลฺยาณํ สาตฺถํ สพฺยญฺชนํ เกวลปริปุณฺณํ ปริสุทฺธํ พฺรหฺมจริยํ ปกาเสสฺสามิ, ยทิทํ ฉ ฉกฺกานิ, ตํ สุณาถ สาธุกํ มนสิ กโรถ, ภาสิสฺสามี”ติ. “เอวํ ภนฺเต”ติ โข เต ภิกฺขู ภควโต ปจฺจสฺโสสุํ.\csromanspacer{} ภควา เอตทโวจ\csromandash{}}

\prose{ฉ อชฺฌตฺติกานิ อายตนานิ เวทิตพฺพานิ, ฉ พาหิรานิ อายตนานิ เวทิตพฺพานิ, ฉ วิญฺญาณกายา เวทิตพฺพา, ฉ ผสฺสกายา เวทิตพฺพา, ฉ เวทนากายา เวทิตพฺพา, ฉ ตณฺหากายา เวทิตพฺพา.}

\proseitem{411}{“ฉ อชฺฌตฺติกานิ อายตนานิ เวทิตพฺพานี”ติ อิติ โข ปเนตํ วุตฺตํ, กิญฺเจตํ ปฏิจฺจ วุตฺตํ.\csromanspacer{} จกฺขายตนํ โสตายตนํ ฆานายตนํ \csromanfolio{328}ชิวฺหายตนํ กายายตนํ มนายตนํ. “ฉ อชฺฌตฺติกานิ อายตนานิ เวทิตพฺพานี”ติ อิติ ยํ ตํ วุตฺตํ, อิทเมตํ ปฏิจฺจ วุตฺตํ.\csromanspacer{} อิทํ ปฐมํ ฉกฺกํ. (1)}

\prose{“ฉ พาหิรานิ อายตนานิ เวทิตพฺพานี”ติ อิติ โข ปเนตํ วุตฺตํ, กิญฺเจตํ ปฏิจฺจ วุตฺตํ.\csromanspacer{} รูปายตนํ สทฺทายตนํ คนฺธายตนํ รสายตนํ โผฏฺฐพฺพายตนํ ธมฺมายตนํ. “ฉ พาหิรานิ อายตนานิ เวทิตพฺพานี”ติ อิติ ยํ ตํ วุตฺตํ, อิทเมตํ ปฏิจฺจ วุตฺตํ.\csromanspacer{} อิทํ ทุติยํ ฉกฺกํ. (2)}

\prose{“ฉ วิญฺญาณกายา เวทิตพฺพา”ติ อิติ โข ปเนตํ วุตฺตํ, กิญฺเจตํ ปฏิจฺจ วุตฺตํ.\csromanspacer{} จกฺขุญฺจ ปฏิจฺจ รูเป จ อุปฺปชฺชติ จกฺขุวิญฺญาณํ, โสตญฺจ ปฏิจฺจ สทฺเท จ อุปฺปชฺชติ โสตวิญฺญาณํ, ฆานญฺจ ปฏิจฺจ คนฺเธ จ อุปฺปชฺชติ ฆานวิญฺญาณํ, ชิวฺหญฺจ ปฏิจฺจ รเส จ อุปฺปชฺชติ ชิวฺหาวิญฺญาณํ, กายญฺจ ปฏิจฺจ โผฏฺฐพฺเพ จ อุปฺปชฺชติ กายวิญฺญาณํ, มนญฺจ ปฏิจฺจ ธมฺเม จ อุปฺปชฺชติ มโนวิญฺญาณํ. “ฉ วิญฺญาณกายา เวทิตพฺพา”ติ อิติ ยํ ตํ วุตฺตํ, อิทเมตํ ปฏิจฺจ วุตฺตํ.\csromanspacer{} อิทํ ตติย ฉกฺกํ. (3)}

\prose{“ฉ ผสฺสกายา เวทิตพฺพา”ติ อิติ โข ปเนตํ วุตฺตํ, กิญฺเจตํ ปฏิจฺจ วุตฺตํ.\csromanspacer{} จกฺขุญฺจ ปฏิจฺจ รูเป จ อุปฺปชฺชติ จกฺขุวิญฺญาณํ ติณฺณํ สงฺคติ ผสฺโส, โสตญฺจ ปฏิจฺจ สทฺเท จ อุปฺปชฺชติ โสตวิญฺญาณํ ติณฺณํ สงฺคติ ผสฺโส, ฆานญฺจ ปฏิจฺจ คนฺเธ จ อุปฺปชฺชติ ฆานวิญฺญาณํ ติณฺณํ สงฺคติ ผสฺโส, ชิวฺหญฺจ ปฏิจฺจ รเส จ อุปฺปชฺชติ ชิวฺหาวิญฺญาณํ ติณฺณํ สงฺคติ ผสฺโส, กายญฺจ ปฏิจฺจ โผฏฺฐพฺเพ จ อุปฺปชฺชติ กายวิญฺญาณํ ติณฺณํ สงฺคติ ผสฺโส, มนญฺจ ปฏิจฺจ ธมฺเม จ อุปฺปชฺชติ มโนวิญฺญาณํ ติณฺณํ สงฺคติ ผสฺโส. “ฉ ผสฺสกายา เวทิตพฺพา”ติ อิติ ยํ ตํ วุตฺตํ, อิทเมตํ ปฏิจฺจ วุตฺตํ.\csromanspacer{} อิทํ จตุตฺถํ ฉกฺกํ. (4)}

\prose{“ฉ เวทนากายา เวทิตพฺพา”ติ อิติ โข ปเนตํ วุตฺตํ, กิญฺเจตํ ปฏิจฺจ วุตฺตํ.\csromanspacer{} จกฺขุญฺจ ปฏิจฺจ รูเป จ อุปฺปชฺชติ จกฺขุวิญฺญาณํ ติณฺณํ สงฺคติ ผสฺโส ผสฺสปจฺจยา เวทนา.\csromanspacer{} โสตญฺจ ปฏิจฺจ สทฺเท จ อุปฺปชฺชติ โสตวิญฺญาณํ ติณฺณํ สงฺคติ ผสฺโส ผสฺสปจฺจยา เวทนา.\csromanspacer{} ฆานญฺจ ปฏิจฺจ คนฺเธ จ อุปฺปชฺชติ ฆานวิญฺญาณํ ติณฺณํ สงฺคติ ผสฺโส ผสฺสปจฺจยา เวทนา.\csromanspacer{} ชิวฺหญฺจ ปฏิจฺจ รเส จ อุปฺปชฺชติ ชิวฺหาวิญฺญาณํ ติณฺณํ สงฺคติ ผสฺโส ผสฺสปจฺจยา เวทนา.\csromanspacer{} กายญฺจ ปฏิจฺจ โผฏฺฐพฺเพ จ อุปฺปชฺชติ กายวิญฺญาณํ ติณฺณํ สงฺคติ ผสฺโส}

\csromanheader{2}{6. ฉฉกฺกสุตฺต (148)}
\prose{\csromanfolio{329}ผสฺสปจฺจยา เวทนา.\csromanspacer{} มนญฺจ ปฏิจฺจ ธมฺเม จ อุปฺปชฺชติ มโนวิญฺญาณํ ติณฺณํ สงฺคติ ผสฺโส ผสฺสปจฺจยา เวทนา. “ฉ เวทนากายา เวทิตพฺพา”ติ อิติ ยํ ตํ วุตฺตํ, อิทเมตํ ปฏิจฺจ วุตฺตํ.\csromanspacer{} อิทํ ปญฺจมํ ฉกฺกํ. (5)}

\prose{“ฉ ตณฺหากายา เวทิตพฺพา”ติ อิติ โข ปเนตํ วุตฺตํ, กิญฺเจตํ ปฏิจฺจ วุตฺตํ.\csromanspacer{} จกฺขุญฺจ ปฏิจฺจ รูเป จ อุปฺปชฺชติ จกฺขุวิญฺญาณํ ติณฺณํ สงฺคติ ผสฺโส ผสฺสปจฺจยา เวทนา เวทนาปจฺจยา ตณฺหา.\csromanspacer{} โสตญฺจ ปฏิจฺจ สทฺเท จ อุปฺปชฺชติ โสตวิญฺญาณํ ฯเปฯ.\csromanspacer{} ฆานญฺจ ปฏิจฺจ คนฺเธ จ อุปฺปชฺชติ ฆานวิญฺญาณํ ฯเปฯ.\csromanspacer{} ชิวฺหญฺจ ปฏิจฺจ รเส จ อุปฺปชฺชติ ชิวฺหาวิญฺญาณํ ฯเปฯ.\csromanspacer{} กายญฺจ ปฏิจฺจ โผฏฺฐพฺเพ จ อุปฺปชฺชติ กายวิญฺญาณํ ฯเปฯ.\csromanspacer{} มนญฺจ ปฏิจฺจ ธมฺเม จ อุปฺปชฺชติ มโนวิญฺญาณํ ติณฺณํ สงฺคติ ผสฺโส ผสฺสปจฺจยา เวทนา เวทนาปจฺจยา ตณฺหา. “ฉ ตณฺหากายา เวทิตพฺพา”ติ อิติ ยํ ตํ วุตฺตํ, อิทเมตํ ปฏิจฺจ วุตฺตํ.\csromanspacer{} อิทํ ฉฏฺฐํ ฉกฺกํ. (6)}

\proseitem{422}{“จกฺขุ อตฺตา”ติ โย วเทยฺย, ตํ น อุปปชฺชติ จกฺขุสฺส อุปฺปาโทปิ วโยปิ ปญฺญายติ.\csromanspacer{} ยสฺส โข ปน อุปฺปาโทปิ วโยปิ ปญฺญายติ “อตฺตา เม อุปฺปชฺชติ จ เวติ จา”ติ อิจฺจสฺส เอวมาคตํ โหติ, ตสฺมา ตํ น อุปปชติ, “จกฺขุ อตฺตา”ติ โย วเทยฺย.\csromanspacer{} อิติ จกฺขุ อนตฺตา.}

\prose{“รูปา อตฺตา”ติ โย วเทยฺย, ตํ น อุปปชฺชติ รูปานํ อุปฺปาโทปิ วโยปิ ปญฺญายติ.\csromanspacer{} ยสฺส โข ปน อุปฺปาโทปิ วโยปิ ปญฺญายติ “อตฺตา เม อุปฺปชฺชติ จ เวติ จา”ติ อิจฺจสฺส เอวมาคตํ โหติ, ตสฺมา ตํ น อุปปชฺชติ, “รูปา อตฺตา”ติ โย วเทยฺย.\csromanspacer{} อิติ จกฺขุ อนตฺตา, รูปา อนตฺตา.}

\prose{“จกฺขุวิญฺญาณํ อตฺตา”ติ โย วเทยฺย, ตํ น อุปปชฺชติ จกฺขุวิญฺญาณสฺส อุปฺปาโทปิ วโยปิ ปญฺญายติ.\csromanspacer{} ยสฺส โข ปน อุปฺปาโทปิ วโยปิ ปญฺญายติ “อตฺตา เม อุปฺปชฺชติ จ เวติ จา”ติ อิจฺจสฺส เอวมาคตํ โหติ, ตสฺมา ตํ น อุปปชฺชติ, “จกฺขุวิญฺญาณํ อตฺตา”ติ โย วเทยฺย.\csromanspacer{} อิติ จกฺขุ อนตฺตา, รูปา อนตฺตา, จกฺขุวิญฺญาณํ อนตฺตา.}

\prose{“จกฺขุสมฺผสฺโส อตฺตา”ติ โย วเทยฺย, ตํ น อุปปชฺชติ จกฺขุสมฺผสฺสสฺส อุปฺปาโทปิ วโยปิ ปญฺญายติ.\csromanspacer{} ยสฺส โข ปน อุปฺปาโทปิ วโยปิ ปญฺญายติ “อตฺตา เม อุปฺปชฺชติ จ เวติ จา”ติ อิจฺจสฺส เอวมาคตํ โหติ, ตสฺมา ตํ น อุปปชฺชติ, “จกฺขุสมฺผสฺโส อตฺตา”ติ \csromanfolio{330}โย วเทยฺย.\csromanspacer{} อิติ จกฺขุ อนตฺตา, รูปา อนตฺตา, จกฺขุวิญฺญาณํ อนตฺตา, จกฺขุสมฺผสฺโส อนตฺตา.}

\prose{“เวทนา อตฺตา”ติ โย วเทยฺย, ตํ น อุปปชฺชติ เวทนาย อุปฺปาโทปิ วโยปิ ปญฺญายติ.\csromanspacer{} ยสฺส โข ปน อุปฺปาโทปิ วโยปิ ปญฺญายติ “อตฺตา เม อุปฺปชฺชติ จ เวติ จา”ติ อิจฺจสฺส เอวมาคตํ โหติ, ตสฺมา ตํ น อุปปชฺชติ, “เวทนา อตฺตา”ติ โย วเทยฺย.\csromanspacer{} อิติ จกฺขุ อนตฺตา, รูปา อนตฺตา, จกฺขุวิญฺญาณํ อนตฺตา, จกฺขุสมฺผสฺโส อนตฺตา, เวทนา อนตฺตา.}

\prose{“ตณฺหา อตฺตา”ติ โย วเทยฺย, ตํ น อุปปชฺชติ ตณฺหาย อุปฺปาโทปิ วโยปิ ปญฺญายติ.\csromanspacer{} ยสฺส โข ปน อุปฺปาโทปิ วโยปิ ปญฺญายติ “อตฺตา เม อุปฺปชฺชติ จ เวติ จา”ติ อิจฺจสฺส เอวมาคตํ โหติ, ตสฺมา ตํ น อุปปชฺชติ, “ตณฺหา อตฺตา”ติ โย วเทยฺย.\csromanspacer{} อิติ จกฺขุ อนตฺตา, รูปา อนตฺตา, จกฺขุวิญฺญาณํ อนตฺตา, จกฺขุสมฺผสฺโส อนตฺตา, เวทนา อนตฺตา, ตณฺหา อนตฺตา.}

\proseitem{423}{“โสตํ อตฺตา”ติ โย วเทยฺย ฯเปฯ. “ฆานํ อตฺตา”ติ โย วเทยฺย ฯเปฯ. “ชิวฺหา อตฺตา”ติ โย วเทยฺย ฯเปฯ.\csromanspacer{} กาโย “อตฺตา”ติ โย วเทยฺย ฯเปฯ. “มโน อตฺตา”ติ โย วเทยฺย, ตํ น อุปปชฺชติ มนสฺส อุปฺปาโทปิ วโยปิ ปญฺญายติ.\csromanspacer{} ยสฺส โข ปน อุปฺปาโทปิ วโยปิ ปญฺญายติ “อตฺตา เม อุปฺปชฺชติ จ เวติ จา”ติ อิจฺจสฺส เอวมาคตํ โหติ, ตสฺมา ตํ น อุปปชฺชติ, “มโน อตฺตา”ติ โย วเทยฺย.\csromanspacer{} อิติ มโน อนตฺตา.}

\prose{“ธมฺมา อตฺตา”ติ โย วเทยฺย, ตํ น อุปปชฺชติ ธมฺมานํ อุปฺปาโทปิ วโยปิ ปญฺญายติ.\csromanspacer{} ยสฺส โข ปน อุปฺปาโทปิ วโยปิ ปญฺญายติ “อตฺตา เม อุปฺปชฺชติ จ เวติ จา”ติ อิจฺจสฺส เอวมาคตํ โหติ, ตสฺมา ตํ น อุปปชฺชติ, “ธมฺมา อตฺตา”ติ โย วเทยฺย.\csromanspacer{} อิติ มโน อนตฺตา, ธมฺมา อนตฺตา.}

\prose{“มโนวิญฺญาณํ อตฺตา”ติ โย วเทยฺย, ตํ น อุปปชฺชติ มโนวิญฺญาณสฺส อุปฺปาโทปิ วโยปิ ปญฺญายติ.\csromanspacer{} ยสฺส โข ปน อุปฺปาโทปิ วโยปิ ปญฺญายติ “อตฺตา เม อุปฺปชฺชติ จ เวติ จา”ติ อิจฺจสฺส เอวมาคตํ โหติ, ตสฺมา ตํ น อุปปชฺชติ, “มโนวิญฺญาณํ}

\csromanheader{2}{6. ฉฉกฺกสุตฺต (148)}
\prose{\csromanfolio{331}อตฺตา”ติ โย วเทยฺย.\csromanspacer{} อิติ มโน อนตฺตา, ธมฺมา อนตฺตา, มโนวิญฺญาณํ อนตฺตา.}

\prose{“มโนสมฺผสฺโส อตฺตา”ติ โย วเทยฺย, ตํ น อุปฺปชฺชติ มโนสมฺผสฺสสฺส อุปฺปาโทปิ วโยปิ ปญฺญายติ.\csromanspacer{} ยสฺส โข ปน อุปฺปาโทปิ วโยปิ ปญฺญายติ “อตฺตา เม อุปฺปชฺชติ จ เวติ จา”ติ อิจฺจสฺส เอวมาคตํ โหติ, ตสฺมา ตํ น อุปปชฺชติ, “มโนสมฺผสฺโส อตฺตา”ติ โย วเทยฺย.\csromanspacer{} อิติ มโน อนตฺตา, ธมฺมา อนตฺตา, มโนวิญฺญาณํ อนตฺตา, มโนสมฺผสฺโส อนตฺตา.}

\prose{“เวทนา อตฺตา”ติ โย วเทยฺย, ตํ น อุปปชฺชติ เวทนาย อุปฺปาโทปิ วโยปิ ปญฺญายติ.\csromanspacer{} ยสฺส โข ปน อุปฺปาโทปิ วโยปิ ปญฺญายติ “อตฺตา เม อุปฺปชฺชติ จ เวติ จา”ติ อิจฺจสฺส เอวมาคตํ โหติ, ตสฺมา ตํ น อุปปชฺชติ, “เวทนา อตฺตา”ติ โย วเทยฺย.\csromanspacer{} อิติ มโน อนตฺตา, ธมฺมา อนตฺตา, มโนวิญฺญาณํ อนตฺตา, มโนสมฺผสฺโส อนตฺตา, เวทนา อนตฺตา.}

\prose{“ตณฺหา อตฺตา”ติ โย วเทยฺย, ตํ น อุปปชฺชติ ตณฺหาย อุปฺปาโทปิ วโยปิ ปญฺญายติ.\csromanspacer{} ยสฺส โข ปน อุปฺปาโทปิ วโยปิ ปญฺญายติ “อตฺตา เม อุปฺปชฺชติ จ เวติ จา”ติ อิจฺจสฺส เอวมาคตํ โหติ, ตสฺมา ตํ น อุปปชฺชติ, “ตณฺหา อตฺตา”ติ โย วเทยฺย.\csromanspacer{} อิติ มโน อนตฺตา, ธมฺมา อนตฺตา, มโนวิญฺญาณํ อนตฺตา, มโนสมฺผสฺโส อนตฺตา, เวทนา อนตฺตา, ตณฺหา อนตฺตา.}

\proseitem{424}{อยํ โข ปน ภิกฺขเว สกฺกายสมุทยคามินี ปฏิปทา\csromandash{} จกฺขุํ “เอตํ มม, เอโสหมสฺมิ, เอโส เม อตฺตา”ติ สมนุปสฺสติ.\csromanspacer{} รูเป “เอตํ มม, เอโสหมสฺมิ, เอโส เม อตฺตา”ติ สมนุปสฺสติ.\csromanspacer{} จกฺขุวิญฺญาณํ “เอตํ มม, เอโสหมสฺมิ, เอโส เม อตฺตา”ติ สมนุปสฺสติ.\csromanspacer{} จกฺขุสมฺผสฺสํ “เอตํ มม, เอโสหมสฺมิ, เอโส เม อตฺตา”ติ สมนุปสฺสติ.\csromanspacer{} เวทนํ “เอตํ มม, เอโสหมสฺมิ, เอโส เม อตฺตา”ติ สมนุปสฺสติ.\csromanspacer{} ตณฺหํ “เอตํ มม, เอโสหมสฺมิ, เอโส เม อตฺตา”ติ สมนุปสฺสติ.\csromanspacer{} โสตํ “เอตํ มม, เอโสหมสฺมิ, เอโส เม อตฺตา”ติ สมนุปสฺสติ ฯเปฯ.\csromanspacer{} ฆานํ “เอตํ มม, เอโสหมสฺมิ, เอโส เม อตฺตา”ติ สมนุปสฺสติ ฯเปฯ.\csromanspacer{} ชิวฺหํ “เอตํ มม, เอโสหมสฺมิ, เอโส เม \csromanfolio{332}อตฺตา”ติ สมนุปสฺสติ ฯเปฯ.\csromanspacer{} กายํ “เอตํ มม, เอโสหมสฺมิ, เอโส เม อตฺตา”ติ สมนุปสฺสติ ฯเปฯ.\csromanspacer{} มนํ “เอตํ มม, เอโสหมสฺมิ, เอโส เม อตฺตา”ติ สมนุปสฺสติ.\csromanspacer{} ธมฺเม “เอตํ มม, เอโสหมสฺมิ, เอโส เม อตฺตา”ติ สมนุปสฺสติ.\csromanspacer{} มโนวิญฺญาณํ “เอตํ มม, เอโสหมสฺมิ, เอโส เม อตฺตา”ติ สมนุปสฺสติ.\csromanspacer{} มโนสมฺผสฺสํ “เอตํ มม, เอโสหมสฺมิ, เอโส เม อตฺตา”ติ สมนุปสฺสติ.\csromanspacer{} เวทนํ “เอตํ มม, เอโสหมสฺมิ, เอโส เม อตฺตา”ติ สมนุปสฺสติ.\csromanspacer{} ตณฺหํ “เอตํ มม, เอโสหมสฺมิ, เอโส เม อตฺตา”ติ สมนุปสฺสติ.}

\prose{อยํ โข ปน ภิกฺขเว สกฺกายนิโรธคามินี ปฏิปทา\csromandash{}จกฺขุํ “เนตํ มม, เนโสหมสฺมิ, น เมโส อตฺตา”ติ สมนุปสฺสติ.\csromanspacer{} รูเป “เนตํ มม, เนโสหมสฺมิ, น เมโส อตฺตา”ติ สมนุปสฺสติ.\csromanspacer{} จกฺขุวิญฺญาณํ “เนตํ มม, เนโสหมสฺมิ, น เมโส อตฺตา’ติ สมนุปสฺสติ.\csromanspacer{} จกฺขุสมฺผสฺสํ “เนตํ มม, เนโสหมสฺมิ, น เมโส อตฺตา”ติ สมนุปสฺสติ.\csromanspacer{} เวทนํ “เนตํ มม, เนโสหมสฺมิ, น เมโส อตฺตา”ติ สมนุปสฺสติ.\csromanspacer{} ตณฺหํ “เนตํ มม, เนโสหมสฺมิ, น เมโส อตฺตา”ติ สมนุปสฺสติ.\csromanspacer{} โสตํ “เนตํ มม, เนโสหมสฺมิ, น เมโส อตฺตา”ติ สมนุปสฺสติ ฯเปฯ.\csromanspacer{} ฆานํ “เนตํ มม, เนโสหมสฺมิ, น เมโส อตฺตา”ติ สมนุปสฺสติ ฯเปฯ.\csromanspacer{} ชิวฺหํ “เนตํ มม, เนโสหมสฺมิ, น เมโส อตฺตา”ติ สมนุปสฺสติ ฯเปฯ.\csromanspacer{} กายํ “เนตํ มม, เนโสหมสฺมิ, น เมโส อตฺตา”ติ สมนุปสฺสติ ฯเปฯ.\csromanspacer{} มนํ “เนตํ มม, เนโสหมสฺมิ, น เมโส อตฺตา”ติ สมนุปสฺสติ.\csromanspacer{} ธมฺเม “เนตํ มม, เนโสหมสฺมิ, น เมโส อตฺตา”ติ สมนุปสฺสติ.\csromanspacer{} มโนวิญฺญาณํ “เนตํ มม, เนโสหมสฺมิ, น เมโส อตฺตา”ติ สมนุปสฺสติ.\csromanspacer{} มโนสมฺผสฺสํ “เนตํ มม, เนโสหมสฺมิ, น เมโส อตฺตา”ติ สมนุปสฺสติ.\csromanspacer{} เวทนํ “เนตํ มม, เนโสหมสฺมิ, น เมโส อตฺตา”ติ สมนุปสฺสติ.\csromanspacer{} ตณฺหํ “เนตํ มม, เนโสหมสฺมิ, น เมโส อตฺตา”ติ สมนุปสฺสติ.}

\proseitem{425}{จกฺขุญฺจ ภิกฺขเว ปฏิจฺจ รูเป จ อุปฺปชฺชติ จกฺขุวิญฺญาณํ, ติณฺณํ สงฺคติ ผสฺโส, ผสฺสปจฺจยา อุปฺปชฺชติ เวทยิตํ สุขํ วา ทุกฺขํ วา อทุกฺขมสุขํ วา.\csromanspacer{} โส สุขาย เวทนาย ผุฏฺโฐ สมาโน อภินนฺทติ อภิวทติ อชฺโฌสาย ติฏฺฐติ, ตสฺส ราคานุสโย อนุเสติ.\csromanspacer{} ทุกฺขาย เวทนาย ผุฏฺโฐ สมาโน โสจติ กิลมติ ปริเทวติ}

\csromanheader{2}{6. ฉฉกฺกสุตฺต (148)}
\prose{\csromanfolio{333}อุรตฺตาฬิํ กนฺทติ สมฺโมหํ อาปชฺชติ, ตสฺส ปฏิฆานุสโย อนุเสติ.\csromanspacer{} อทุกฺขมสุขาย เวทนาย ผุฏฺโฐ สมาโน ตสฺสา เวทนาย สมุทยญฺจ อตฺถงฺคมญฺจ อสฺสาทญฺจ อาทีนวญฺจ นิสฺสรณญฺจ ยถาภูตํ นปฺปชานาติ, ตสฺส อวิชฺชานุสโย อนุเสติ.\csromanspacer{} โส วต ภิกฺขเว สุขาย เวทนาย ราคานุสยํ อปฺปหาย ทุกฺขาย เวทนาย ปฏิฆานุสยํ อปฺปฏิวิโนเทตฺวา อทุกฺขมสุขาย เวทนาย อวิชฺชานุสยํ อสมูหนิตฺวา อวิชฺชํ อปฺปหาย วิชฺชํ อนุปฺปาเทตฺวา ทิฏฺเฐว ธมฺเม ทุกฺขสฺสนฺตกโร ภวิสฺสตีติ เนตํ ฐานํ วิชฺชติ.}

\prose{โสตญฺจ ภิกฺขเว ปฏิจฺจ สทฺเท จ อุปฺปชฺชติ โสตวิญฺญาณํ ฯเปฯ.\csromanspacer{} ฆานญฺจ ภิกฺขเว ปฏิจฺจ คนฺเธ จ อุปฺปชฺชติ ฆานวิญฺญาณํ ฯเปฯ.\csromanspacer{} ชิวฺหญฺจ ภิกฺขเว ปฏิจฺจ รเส จ อุปฺปชฺชติ ชิวฺหาวิญฺญาณํ ฯเปฯ.\csromanspacer{} กายญฺจ ภิกฺขเว ปฏิจฺจ โผฏฺฐพฺเพ จ อุปฺปชฺชติ กายวิญฺญาณํ ฯเปฯ.\csromanspacer{} มนญฺจ ภิกฺขเว ปฏิจฺจ ธมฺเม จ อุปฺปชฺชติ มโนวิญฺญาณํ, ติณฺณํ สงฺคติ ผสฺโส, ผสฺสปจฺจยา อุปฺปชฺชติ เวทยิตํ สุขํ วา ทุกฺขํ วา อทุกฺขมสุขํ วา.\csromanspacer{} โส สุขาย เวทนาย ผุฏฺโฐ สมาโน อภินนฺทติ อภิวทติ อชฺโฌสาย ติฏฺฐติ, ตสฺส ราคานุสโย อนุเสติ.\csromanspacer{} ทุกฺขาย เวทนาย ผุฏฺโฐ สมาโน โสจติ กิลมติ ปริเทวติ อุรตฺตาลิํ กนฺทติ สมฺโมหํ อาปชฺชติ, ตสฺส ปฏิฆานุสโย อนุเสติ.\csromanspacer{} อทุกฺขมสุขาย เวทนาย ผุฏฺโฐ สมาโน ตสฺสา เวทนาย สมุทยญฺจ อตฺถงฺคมญฺจ อสฺสาทญฺจ อาทีนวญฺจ นิสฺสรณญฺจ ยถาภูตํ นปฺปชานาติ, ตสฺส อวิชฺชานุสโย อนุเสติ.\csromanspacer{} โส วต ภิกฺขเว สุขาย เวทนาย ราคานุสยํ อปฺปหาย ทุกฺขาย เวทนาย ปฏิฆานุสยํ อปฺปฏิวิโนเทตฺวา อทุกฺขมสุขาย เวทนาย อวิชฺชานุสยํ อสมูหนิตฺวา อวิชฺชํ อปฺปหาย วิชฺชํ อนุปฺปาเทตฺวา ทิฏฺเฐว ธมฺเม ทุกฺขสฺสนฺตกโร ภวิสฺสตีติ เนตํ ฐานํ วิชฺชติ.}

\proseitem{426}{จกฺขุญฺจ ภิกฺขเว ปฏิจฺจ รูเป จ อุปฺปชฺชติ จกฺขุวิญฺญาณํ, ติณฺณํ สงฺคติ ผสฺโส, ผสฺสปจฺจยา อุปฺปชฺชติ เวทยิตํ สุขํ วา ทุกฺขํ วา อทุกฺขมสุขํ วา.\csromanspacer{} โส สุขาย เวทนาย ผุฏฺโฐ สมาโน นาภินนฺทติ นาภิวทติ นาชฺโฌสาย ติฏฺฐติ, ตสฺส ราคานุสโย นานุเสติ.\csromanspacer{} ทุกฺขาย เวทนาย ผุฏฺโฐ สมาโน น โสจติ น กิลมติ น ปริเทวติ น อุรตฺตาฬิํ กนฺทติ น สมฺโมหํ อาปชฺชติ, ตสฺส ปฏิฆานุสโย นานุเสติ.\csromanspacer{} อทุกฺขมสุขาย เวทนาย ผุฏฺโฐ สมาโน \csromanfolio{334}ตสฺสา เวทนาย สมุทยญฺจ อตฺถงฺคมญฺจ อสฺสาทญฺจ อาทีนวญฺจ นิสฺสรณญฺจ ยถาภูตํ ปชานาติ, ตสฺส อวิชฺชานุสโย นานุเสติ.\csromanspacer{} โส วต ภิกฺขเว สุขาย เวทนาย ราคานุสยํ ปหาย ทุกฺขาย เวทนาย ปฏิฆานุสยํ ปฏิวิโนเทตฺวา อทุกฺขมสุขาย เวทนาย อวิชฺชานุสยํ สมูหนิตฺวา อวิชฺชํ ปหาย วิชฺชํ อุปฺปาเทตฺวา ทิฏฺเฐว ธมฺเม ทุกฺขสฺสนฺตกโร ภวิสฺสตีติ ฐานเมตํ วิชฺชติ.}

\prose{โสตญฺจ ภิกฺขเว ปฏิจฺจ สทฺเท จ อุปฺปชฺชติ โสตวิญฺญาณํ ฯเปฯ.\csromanspacer{} ฆานญฺจ ภิกฺขเว ปฏิจฺจ คนฺเธ จ อุปฺปชฺชติ ฆานวิญฺญาณํ ฯเปฯ.\csromanspacer{} ชิวฺหญฺจ ภิกฺขเว ปฏิจฺจ รเส จ อุปฺปชฺชติ ชิวฺหาวิญฺญาณํ ฯเปฯ.\csromanspacer{} กายญฺจ ภิกฺขเว ปฏิจฺจ โผฏฺฐพฺเพ จ อุปฺปชฺชติ กายวิญฺญาณํ ฯเปฯ.\csromanspacer{} มนญฺจ ภิกฺขเว ปฏิจฺจ ธมฺเม จ อุปฺปชฺชติ มโนวิญฺญาณํ, ติณฺณํ สงฺคติ ผสฺโส, ผสฺสปจฺจยา อุปฺปชฺชติ เวทยิตํ สุขํ วา ทุกฺขํ วา อทุกฺขมสุขํ วา.\csromanspacer{} โส สุขาย เวทนาย ผุฏฺโฐ สมาโน นาภินนฺทติ นาภิวทติ นาชฺโฌสาย ติฏฺฐติ, ตสฺส ราคานุสโย นานุเสติ.\csromanspacer{} ทุกฺขาย เวทนาย ผุฏฺโฐ สมาโน น โสจติ น กิลมติ น ปริเทวติ น อุรตฺตาฬิํ กนฺทติ น สมฺโมหํ อาปชฺชติ, ตสฺส ปฏิฆานุสโย นานุเสติ.\csromanspacer{} อทุกฺขมสุขาย เวทนาย ผุฏฺโฐ สมาโน ตสฺสา เวทนาย สมุทยญฺจ อตฺถงฺคมญฺจ อสฺสาทญฺจ อาทีนวญฺจ นิสฺสรณญฺจ ยถาภูตํ ปชานาติ, ตสฺส อวิชฺชานุสโย นานุเสติ.\csromanspacer{} โส วต ภิกฺขเว สุขาย เวทนาย ราคานุสยํ ปหาย ทุกฺขาย เวทนาย ปฏิฆานุสยํ ปฏิวิโนเทตฺวา อทุกฺขมสุขาย เวทนาย อวิชฺชานุสยํ สมูหนิตฺวา อวิชฺชํ ปหาย วิชฺชํ อุปฺปาเทตฺวา ทิฏฺเฐว ธมฺเม ทุกฺขสฺสนฺตกโร ภวิสฺสตีติ ฐานเมตํ วิชฺชติ.}

\proseitem{427}{เอวํ ปสฺสํ ภิกฺขเว สุตวา อริยสาวโก จกฺขุสฺมิํ\footnote{จกฺขุสฺมิมฺปิ (สฺยา, กํ) เอวมิตเรสุปิ.} นิพฺพินฺทติ, รูเปสุ นิพฺพินฺทติ, จกฺขุวิญฺญาเณ นิพฺพินฺทติ, จกฺขุสมฺผสฺเส นิพฺพินฺทติ, เวทนาย นิพฺพินฺทติ, ตณฺหาย นิพฺพินฺทติ, โสตสฺมิํ นิพฺพินฺทติ, สทฺเทสุ นิพฺพินฺทติ ฯเปฯ.\csromanspacer{} ฆานสฺมิํ นิพฺพินฺทติ, คนฺเธสุ นิพฺพินฺทติ.\csromanspacer{} ชิวฺหาย นิพฺพินฺทติ, รเสสุ นิพฺพินฺทติ.\csromanspacer{} กายสฺมิํ นิพฺพินฺทติ, โผฏฺฐพฺเพสุ นิพฺพินฺทติ.\csromanspacer{} มนสฺมิํ นิพฺพินฺทติ, ธมฺเมสุ นิพฺพินฺทติ, มโนวิญฺญาเณ นิพฺพินฺทติ, มโนสมฺผสฺเส นิพฺพินฺทติ, เวทนาย นิพฺพินฺทติ, ตณฺหาย นิพฺพินฺทติ.\csromanspacer{} นิพฺพินฺทํ วิรชฺชติ, วิราคา วิมุจฺจติ, วิมุตฺตสฺมิํ “วิมุตฺตมฺ”อิติ}

\csromanheader{2}{7. มหาสฬายตนิกสุตฺต (149)}
\prose{\csromanfolio{335}ญาณํ โหติ, “ขีณา ชาติ, วุสิตํ พฺรหฺมจริยํ, กตํ กรณียํ, นาปรํ อิตฺถตฺตายา”ติ ปชานาตีติ.}

\prose{อิทมโวจ ภควา.\csromanspacer{} อตฺตมนา เต ภิกฺขู ภควโต ภาสิตํ อภินนฺทุนฺติ.\csromanspacer{} อิมสฺมิํ โข ปน เวยฺยากรณสฺมิํ ภญฺญมาเน สฏฺฐิมตฺตานํ ภิกฺขูนํ อนุปาทาย อาสเวหิ จิตฺตานิ วิมุจฺจิํสูติ.}

\nitthitamruled{ฉฉกฺกสุตฺตํ นิฏฺฐิตํ ฉฏฺฐํ.}
\csromanmatikaanchor{mtk.58}
\csromantocmarkref{subsection}{7. มหาสฬายตนิกสุตฺต}{mtk.58}
\bookmark[dest={mtk.58},level=2]{7. มหาสฬายตนิกสุตฺต}
\csromanheader{2}{7. มหาสฬายตนิกสุตฺต}
\proseitem{428}{เอวํ เม สุตํ\csromandash{}เอกํ สมยํ ภควา สาวตฺถิยํ วิหรติ เชตวเน อนาถปิณฺฑิกสฺส อาราเม.\csromanspacer{} ตตฺร โข ภควา ภิกฺขู อามนฺเตสิ “ภิกฺขโว”ติ. “ภทนฺเต”ติ เต ภิกฺขู ภควโต ปจฺจสฺโสสุํ.\csromanspacer{} ภควา เอตทโวจ “มหาสฬายตนิกํ โว ภิกฺขเว เทเสสฺสามิ, ตํ สุณาถ สาธุกํ มนสิ กโรถ, ภาสิสฺสามี”ติ. “เอวํ ภนฺเต”ติ โข เต ภิกฺขู ภควโต ปจฺจสฺโสสุํ.\csromanspacer{} ภควา เอตทโวจ\csromandash{}}

\proseitem{429}{จกฺขุํ ภิกฺขเว อชานํ อปสฺสํ ยถาภูตํ, รูเป อชานํ อปสฺสํ ยถาภูตํ, จกฺขุวิญฺญาณํ อชานํ อปสฺสํ ยถาภูตํ, จกฺขุสมฺผสฺสํ อชานํ อปสฺสํ ยถาภูตํ, ยมิทํ จกฺขุสมฺผสฺสปจฺจยา อุปฺปชฺชติ เวทยิตํ สุขํ วา ทุกฺขํ วา อทุกฺขมสุขํ วา, ตมฺปิ อชานํ อปสฺสํ ยถาภูตํ จกฺขุสฺมิํ สารชฺชติ, รูเปสุ สารชฺชติ, จกฺขุวิญฺญาเณ สารชฺชติ, จกฺขุสมฺผสฺเส สารชฺชติ, ยมิทํ จกฺขุสมฺผสฺสปจฺจยา อุปฺปชฺชติ เวทยิตํ สุขํ วา ทุกฺขํ วา อทุกฺขมสุขํ วา, ตสฺมิมฺปิ สารชฺชติ.}

\prose{ตสฺส สารตฺตสฺส สํยุตฺตสฺส สมฺมูฬฺหสฺส อสฺสาทานุปสฺสิโน วิหรโต อายติํ ปญฺจุปาทานกฺขนฺธา อุปจยํ คจฺฉนฺติ.\csromanspacer{} ตณฺหา จสฺส โปโนพฺภวิกา นนฺทีราคสหคตา ตตฺรตตฺราภินนฺทินี, สา จสฺส ปวฑฺฒติ.\csromanspacer{} ตสฺส กายิกาปิ ทรถา ปวฑฺฒนฺติ, เจตสิกาปิ ทรถา ปวฑฺฒนฺติ.\csromanspacer{} กายิกาปิ สนฺตาปา ปวฑฺฒนฺติ, เจตสิกาปิ สนฺตาปา ปวฑฺฒนฺติ.\csromanspacer{} กายิกาปิ ปริฬาหา ปวฑฺฒนฺติ, เจตสิกาปิ ปริฬาหา ปวฑฺฒนฺติ.\csromanspacer{} โส กายทุกฺขมฺปิ\footnote{กายิกทุกฺขมฺปิ (สฺยา, กํ), กายิกํ ทุกฺขมฺปิ (ก)} เจโตทุกฺขมฺปิ ปฏิสํเวเทติ.}

\prose{\csromanfolio{336}โสตํ ภิกฺขเว อชานํ อปสฺสํ ยถาภูตํ ฯเปฯ.\csromanspacer{} ฆานํ ภิกฺขเว อชานํ อปสฺสํ ยถาภูตํ ฯเปฯ.\csromanspacer{} ชิวฺหํ ภิกฺขเว อชานํ อปสฺสํ ยถาภูตํ ฯเปฯ.\csromanspacer{} กายํ ภิกฺขเว อชานํ อปสฺสํ ยถาภูตํ ฯเปฯ.\csromanspacer{} มนํ ภิกฺขเว อชานํ อปสฺสํ ยถาภูตํ, ธมฺเม ภิกฺขเว อชานํ อปสฺสํ ยถาภูตํ, มโนวิญฺญาณํ ภิกฺขเว อชานํ อปสฺสํ ยถาภูตํ, มโนสมฺผสฺสํ ภิกฺขเว อชานํ อปสฺสํ ยถาภูตํ, ยมิทํ มโนสมฺผสฺสปจฺจยา อุปฺปชฺชติ เวทยิตํ สุขํ วา ทุกฺขํ วา อทุกฺขมสุขํ วา, ตมฺปิ อชานํ อปสฺสํ ยถาภูตํ มนสฺมิํ สารชฺชติ, ธมฺเมสุ สารชฺชติ, มโนวิญฺญาเณ สารชฺชติ, มโนสมฺผสฺเส สารชฺชติ, ยมิทํ มโนสมฺผสฺสปจฺจยา อุปฺปชฺชติ เวทยิตํ สุขํ วา ทุกฺขํ วา อทุกฺขมสุขํ วา, ตสฺมิมฺปิ สารชฺชติ.}

\prose{ตสฺส สารตฺตสฺส สํยุตฺตสฺส สมฺมูฬฺหสฺส อสฺสาทานุปสฺสิโน วิหรโต อายติํ ปญฺจุปาทานกฺขนฺธา อุปจยํ คจฺฉนฺติ.\csromanspacer{} ตณฺหา จสฺส โปโนพฺภวิกา นนฺทีราคสหคตา ตตฺรตตฺราภินนฺทินี, สา จสฺส ปวฑฺฒติ.\csromanspacer{} ตสฺส กายิกาปิ ทรถา ปวฑฺฒนฺติ, เจตสิกาปิ ทรถา ปวฑฺฒนฺติ.\csromanspacer{} กายิกาปิ สนฺตาปา ปวฑฺฒนฺติ, เจตสิกาปิ สนฺตาปา ปวฑฺฒนฺติ.\csromanspacer{} กายิกาปิ ปริฬาหา ปวฑฺฒนฺติ, เจตสิกาปิ ปริฬาหา ปวฑฺฒนฺติ.\csromanspacer{} โส กายทุกฺขมฺปิ เจโตทุกฺขมฺปิ ปฏิสํเวเทติ.}

\proseitem{430}{จกฺขุญฺจ โข ภิกฺขเว ชานํ ปสฺสํ ยถาภูตํ, รูเป ชานํ ปสฺสํ ยถาภูตํ, จกฺขุวิญฺญาณํ ชานํ ปสฺสํ ยถาภูตํ, จกฺขุสมฺผสฺสํ ชานํ ปสฺสํ ยถาภูตํ, ยมิทํ จกฺขุสมฺผสฺสปจฺจยา อุปฺปชฺชติ เวทยิตํ สุขํ วา ทุกฺขํ วา อทุกฺขมสุขํ วา, ตมฺปิ ชานํ ปสฺสํ ยถาภูตํ จกฺขุสฺมิํ น สารชฺชติ, รูเปสุ น สารชฺชติ, จกฺขุวิญฺญาเณ น สารชฺชติ, จกฺขุสมฺผสฺเส น สารชฺชติ, ยมิทํ จกฺขุสมฺผสฺสปจฺจยา อุปฺปชฺชติ เวทยิตํ สุขํ วา ทุกฺขํ วา อทุกฺขมสุขํ วา, ตสฺมิมฺปิ น สารชฺชติ.}

\prose{ตสฺส อสารตฺตสฺส อสํยุตฺตสฺส อสมฺมูฬฺหสฺส อาทีนวานุปสฺสิโน วิหรโต อายติํ ปญฺจุปาทานกฺขนฺธา อปจยํ คจฺฉนฺติ.\csromanspacer{} ตณฺหา จสฺส โปโนพฺภวิกา นนฺทีราคสหคตา ตตฺรตตฺราภินนฺทินี, สา จสฺส ปหียติ.\csromanspacer{} ตสฺส กายิกาปิ ทรถา ปหียนฺติ, เจตสิกาปิ ทรถา ปหียนฺติ.\csromanspacer{} กายิกาปิ สนฺตาปา ปหียนฺติ, เจตสิกาปิ สนฺตาปา ปหียนฺติ.\csromanspacer{} กายิกาปิ ปริฬาหา ปหียนฺติ, เจตสิกาปิ ปริฬาหา ปหียนฺติ.\csromanspacer{} โส กายสุขมฺปิ เจโตสุขมฺปิ ปฏิสํเวเทติ.}

\csromanheader{2}{7. มหาสฬายตนิกสุตฺต (149)}
\proseitem{431}{\csromanfolio{337}ยา ตถาภูตสฺส\footnote{ยถาภูตสฺส (สี, อิ)} ทิฏฺฐิ, สาสฺส โหติ สมฺมาทิฏฺฐิ.\csromanspacer{} โย ตถาภูตสฺส1 สงฺกปฺโป, สฺวาสฺส โหติ สมฺมาสงฺกปฺโป.\csromanspacer{} โย ตถาภูตสฺส1 วายาโม, สฺวาสฺส โหติ สมฺมาวายาโม.\csromanspacer{} ยา ตถาภูตสฺส1 สติ, สาสฺส โหติ สมฺมาสติ.\csromanspacer{} โย ตถาภูตสฺส1 สมาธิ, สฺวาสฺส โหติ สมฺมาสมาธิ.\csromanspacer{} ปุพฺเพว โข ปนสฺส กายกมฺมํ วจีกมฺมํ อาชีโว สุปริสุทฺโธ โหติ.\csromanspacer{} เอวมสฺสายํ อริโย อฏฺฐงฺคิโก มคฺโค ภาวนาปาริปูริํ คจฺฉติ.}

\prose{ตสฺส เอวํ อิมํ อริยํ อฏฺฐงฺคิกํ มคฺคํ ภาวยโต จตฺตาโรปิ สติปฏฺฐานา ภาวนาปาริปูริํ คจฺฉนฺติ, จตฺตาโรปิ สมฺมปฺปธานา ภาวนาปาริปูริํ คจฺฉนฺติ, จตฺตาโรปิ อิทฺธิปาทา ภาวนาปาริปูริํ คจฺฉนฺติ, ปญฺจปิ อินฺทฺริยานิ ภาวนาปาริปูริํ คจฺฉนฺติ, ปญฺจปิ พลานิ ภาวนาปาริปูริํ คจฺฉนฺติ, สตฺตปิ โพชฺฌงฺคา ภาวนาปาริปูริํ คจฺฉนฺติ.}

\prose{ตสฺสิเม เทฺว ธมฺมา ยุคนนฺธา\footnote{ยุคนทฺธา (สี, สฺยา, กํ)} วตฺตนฺติ, สมโถ จ วิปสฺสนา จ.\csromanspacer{} โส เย ธมฺมา อภิญฺญา ปริญฺเญยฺยา, เต ธมฺเม อภิญฺญา ปริชานาติ.\csromanspacer{} เย ธมฺมา อภิญฺญา ปหาตพฺพา, เต ธมฺเม อภิญฺญา ปชหติ.\csromanspacer{} เย ธมฺมา อภิญฺญา ภาเวตพฺพา, เต ธมฺเม อภิญฺญา ภาเวติ.\csromanspacer{} เย ธมฺมา อภิญฺญา สจฺฉิกาตพฺพา, เต ธมฺเม อภิญฺญา สจฺฉิกโรติ.}

\prose{กตเม จ ภิกฺขเว ธมฺมา อภิญฺญา ปริญฺเญยฺยา, ปญฺจุปาทานกฺขนฺธาติสฺส วจนียํ.\csromanspacer{} เสยฺยถิทํ, รูปุปาทานกฺขนฺโธ เวทนุปาทานกฺขนฺโธ สญฺญุปาทานกฺขนฺโธ สงฺขารุปาทานกฺขนฺโธ วิญฺญาณุปาทานกฺขนฺโธ.\csromanspacer{} อิเม ธมฺมา อภิญฺญา ปริญฺเญยฺยา.}

\prose{กตเม จ ภิกฺขเว ธมฺมา อภิญฺญา ปหาตพฺพา, อวิชฺชา จ ภวตณฺหา จ.\csromanspacer{} อิเม ธมฺมา อภิญฺญา ปหาตพฺพา.}

\prose{กตเม จ ภิกฺขเว ธมฺมา อภิญฺญา ภาเวตพฺพา, สมโถ จ วิปสฺสนา จ.\csromanspacer{} อิเม ธมฺมา อภิญฺญา ภาเวตพฺพา.}

\prose{กตเม จ ภิกฺขเว ธมฺมา อภิญฺญา สจฺฉิกาตพฺพา, วิชฺชา จ วิมุตฺติ จ.\csromanspacer{} อิเม ธมฺมา อภิญฺญา สจฺฉิกาตพฺพา.}

\proseitem{432}{\csromanfolio{338}โสตํ ภิกฺขเว ชานํ ปสฺสํ ยถาภูตํ ฯเปฯ.\csromanspacer{} ฆานํ ภิกฺขเว ชานํ ปสฺสํ ยถาภูตํ ฯเปฯ.\csromanspacer{} ชิวฺหํ ภิกฺขเว ชานํ ปสฺสํ ยถาภูตํ ฯเปฯ.\csromanspacer{} กายํ ภิกฺขเว ชานํ ปสฺสํ ยถาภูตํ ฯเปฯ.\csromanspacer{} มนํ ภิกฺขเว ชานํ ปสฺสํ ยถาภูตํ, ธมฺเม ชานํ ปสฺสํ ยถาภูตํ, มโนวิญฺญาณํ ชานํ ปสฺสํ ยถาภูตํ, มโนสมฺผสฺสํ ชานํ ปสฺสํ ยถาภูตํ, ยมิทํ มโนสมฺผสฺสปจฺจยา อุปฺปชฺชติ เวทยิตํ สุขํ วา ทุกฺขํ วา อทุกฺขมสุขํ วา, ตมฺปิ ชานํ ปสฺสํ ยถาภูตํ มนสฺมิํ น สารชฺชติ, ธมฺเมสุ น สารชฺชติ, มโนวิญฺญาเณ น สารชฺชติ, มโนสมฺผสฺเส น สารชฺชติ, ยมิทํ มโนสมฺผสฺสปจฺจยา อุปฺปชฺชติ เวทยิตํ สุขํ วา ทุกฺขํ วา อทุกฺขมสุขํ วา, ตสฺมิมฺปิ น สารชฺชติ.}

\prose{ตสฺส อสารตฺตสฺส อสํยุตฺตสฺส อสมฺมูฬฺหสฺส อาทีนวานุปสฺสิโน วิหรโต อายติํ ปาญฺจุปาทานกฺขนฺธา อปจยํ คจฺฉนฺติ.\csromanspacer{} ตณฺหา จสฺส โปโนพฺภวิกา นนฺทีราคสหคตา ตตฺรตตฺราภินนฺทินี, สา จสฺส ปหียติ.\csromanspacer{} ตสฺส กายิกาปิ ทรถา ปหียนฺติ, เจตสิกาปิ ทรถา ปหียนฺติ.\csromanspacer{} กายิกาปิ สนฺตาปา ปหียนฺติ, เจตสิกาปิ สนฺตาปา ปหียนฺติ.\csromanspacer{} กายิกาปิ ปริฬาหา ปหียนฺติ, เจตสิกาปิ ปริฬาหา ปหียนฺติ.\csromanspacer{} โส กายสุขมฺปิ เจโตสุขมฺปิ ปฏิสํเวเทติ.}

\proseitem{433}{ยา ตถาภูตสฺส ทิฏฺฐิ, สาสฺส โหติ สมฺมาทิฏฺฐิ.\csromanspacer{} โย ตถาภูตสฺส สงฺกปฺโป, สฺวาสฺส โหติ สมฺมาสงฺกปฺโป.\csromanspacer{} โย ตถาภูตสฺส วายาโม, สฺวาสฺส โหติ สมฺมาวายาโม.\csromanspacer{} ยา ตถาภูตสฺส สติ, สาสฺส โหติ สมฺมาสติ.\csromanspacer{} โย ตถาภูตสฺส สมาธิ, สฺวาสฺส โหติ สมฺมาสมาธิ.\csromanspacer{} ปุพฺเพว โข ปนสฺส กายกมฺมํ วจีกมฺมํ อาชีโว สุปริสุทฺโธ โหติ.\csromanspacer{} เอวมสฺสายํ อริโย อฏฺฐงฺคิโก มคฺโค ภาวนาปาริปูริํ คจฺฉติ.}

\prose{ตสฺส เอวํ อิมํ อริยํ อฏฺฐงฺคิกํ มคฺคํ ภาวยโต จตฺตาโรปิ สติปฏฺฐานา ภาวนาปาริปูริํ คจฺฉนฺติ, จตฺตาโรปิ สมฺมปฺปธานา ภาวนาปาริปูริํ คจฺฉนฺติ, จตฺตาโรปิ อิทฺธิปาทา ภาวนาปาริปูริํ คจฺฉนฺติ, ปญฺจปิ อินฺทฺริยานิ ภาวนาปาริปูริํ คจฺฉนฺติ, ปญฺจปิ พาลานิ ภาวนาปาริปูริํ คจฺฉนฺติ, สตฺตปิ โพชฺฌงฺคา ภาวนาปาริปูริํ คจฺฉนฺติ.}

\prose{ตสฺสิเม เทฺว ธมฺมา ยุคนนฺธา วตฺตนฺติ, สมโถ จ วิปสฺสนา จ.\csromanspacer{} โส เย ธมฺมา อภิญฺญา ปริญฺเญยฺยา, เต ธมฺเม อภิญฺญา ปริชานาติ.}

\csromanheader{2}{8. นครวินฺเทยฺยสุตฺต (150)}
\prose{\csromanfolio{339}เย ธมฺมา อภิญฺญา ปหาตพฺพา, เต ธมฺเม อภิญฺญา ปชหติ.\csromanspacer{} เย ธมฺมา อภิญฺญา ภาเวตพฺพา, เต ธมฺเม อภิญฺญา ภาเวติ.\csromanspacer{} เย ธมฺมา อภิญฺญา สจฺฉิกาตพฺพา, เต ธมฺเม อภิญฺญา สจฺฉิกโรติ.}

\prose{กตเม จ ภิกฺขเว ธมฺมา อภิญฺญา ปริญฺเญยฺยา, ปญฺจุปาทานกฺขนฺธาติสฺส วจนียํ.\csromanspacer{} เสยฺยถิทํ, รูปุปาทานกฺขนฺโธ เวทนุปาทานกฺขนฺโธ สญฺญุปาทานกฺขนฺโธ สงฺขารุปาทานกฺขนฺโธ วิญฺญาณุปาทานกฺขนฺโธ.\csromanspacer{} อิเม ธมฺมา อภิญฺญา ปริญฺเญยฺยา.}

\prose{กตเม ภิกฺขเว ธมฺมา อภิญฺญา ปหาตพฺพา, อวิชฺชา จ ภวตณฺหา จ.\csromanspacer{} อิเม ธมฺมา อภิญฺญา ปหาตพฺพา.}

\prose{กตเม จ ภิกฺขเว ธมฺมา อภิญฺญา ภาเวตพฺพา, สมโถ จ วิปสฺสนา จ.\csromanspacer{} อิเม ธมฺมา อภิญฺญา ภาเวตพฺพา.}

\prose{กตเม จ ภิกฺขเว ธมฺมา อภิญฺญา สจฺฉิกาตพฺพา, วิชฺชา จ วิมุตฺติ จ.\csromanspacer{} อิเม ธมฺมา อภิญฺญา สจฺฉิกาตพฺพาติ.}

\prose{อิทมโวจ ภควา.\csromanspacer{} อตฺตมนา เต ภิกฺขู ภควโต ภาสิตํ อภินฺทุนฺติ.}

\nitthitamruled{มหาสฬายตนิกสุตฺตํ นิฏฺฐิตํ สตฺตมํ.}
\csromanmatikaanchor{mtk.59}
\csromantocmarkref{subsection}{8. นครวินฺเทยฺยสุตฺต}{mtk.59}
\bookmark[dest={mtk.59},level=2]{8. นครวินฺเทยฺยสุตฺต}
\csromanheader{2}{8. นครวินฺเทยฺยสุตฺต}
\proseitem{434}{เอวํ เม สุตํ\csromandash{}เอกํ สมยํ ภควา โกสเลสุ จาริกํ จรมาโน มหตา ภิกฺขุสํเฆน สทฺธิํ เยน นครวินฺทํ นาม โกสลานํ พฺราหฺมณานํ คาโม ตทวสริ.\csromanspacer{} อสฺโสสุํ โข นครวินฺเทยฺยกา\footnote{นครวินฺเทยฺยา (ก)} พฺราหฺมณคหปติกา “สมโณ ขลุ โภ โคตโม สกฺยปุตฺโต สกฺยกุลา ปพฺพชิโต โกสเลสุ จาริกํ จรมาโน มหตา ภิกฺขุสํเฆน สทฺธิํ นครวินฺทํ อนุปฺปตฺโต.\csromanspacer{} ตํ โข ปน ภวนฺตํ โคตมํ เอวํ กลฺยาโณ กิตฺติสทฺโท อพฺภุคฺคโต ‘อิติปิ โส ภควา อรหํ สมฺมาสมฺพุทฺโธ วิชฺชาจรณสมฺปนฺโน สุคโต โลกวิทู อนุตฺตโร ปุริสทมฺมสารถิ สตฺถา เทวมนุสฺสานํ พุทฺโธ ภควา’ติ. \csromanfolio{340}โส อิมํ โลกํ สเทวกํ สมารกํ สพฺรหฺมกํ สสฺสมณพฺราหฺมณิํ ปชํ สเทวมนุสฺสํ สยํ อภิญฺญา สจฺฉิกตฺวา ปเวเทติ, โส ธมฺมํ เทเสติ อาทิกลฺยาณํ มชฺเฌกลฺยาณํ ปริโยสานกลฺยาณํ สาตฺถํ สพฺยญฺชนํ เกวลปริปุณฺณํ ปริสุทฺธํ พฺรหฺมจริยํ ปกาเสติ, สาธุ โข ปน ตถารูปานํ อรหตํ ทสฺสนํ โหตี”ติ.}

\prose{อถ โข นครวินฺเทยฺยกา พฺราหฺมณคหปติกา เยน ภควา เตนุปสงฺกมิํสุ, อุปสงฺกมิตฺวา อปฺเปกจฺเจ ภควนฺตํ อภิวาเทตฺวา เอกมนฺตํ นิสีทิํสุ, อปฺเปกจฺเจ ภควตา สทฺธิํ สมฺโมทิํสุ, สมฺโมทนียํ กถํ สารณียํ วีติสาเรตฺวา เอกมนฺตํ นิสีทิํสุ, อปฺเปกจฺเจ เยน ภควา เตนญฺชลิํ ปณาเมตฺวา เอกมนฺตํ นิสีทิํสุ, อปฺเปกจฺเจ ภควโต สนฺติเก นามโคตฺตํ สาเวตฺวา เอกมนฺตํ นิสีทิํสุ, อปฺเปกจฺเจ ตุณฺหีภูตา เอกมนฺตํ นิสีทิํสุ, เอกมนฺตํ นิสินฺเน โข นครวินฺเทยฺยเก พฺราหฺมณคหปติเก ภควา เอตทโวจ\csromandash{}}

\proseitem{435}{สเจ โว คหปตโย อญฺญติตฺถิยา ปริพฺพาชกา เอวํ ปุจฺเฉยฺยุํ “กถํภูตา คหปตโย สมณพฺราหฺมณา น สกฺกาตพฺพา น ครุกาตพฺพา น มาเนตพฺพา น ปูเชตพฺพา”ติ, เอวํ ปุฏฺฐา ตุเมฺห คหปตโย เตสํ อญฺญติตฺถิยานํ ปริพฺพาชกานํ เอวํ พฺยากเรยฺยาถ “เย เต สมณพฺราหฺมณา จกฺขุวิญฺเญยฺเยสุ รูเปสุ อวีตราคา อวีตโทสา อวีตโมหา อชฺฌตฺตํ อวูปสนฺตจิตฺตา สมวิสมํ จรนฺติ กาเยน วาจาย มนสา, เอวรูปา สมณพฺราหฺมณา น สกฺกาตพฺพา น ครุกาตพฺพา น มาเนตพฺพา น ปูเชตพฺพา.\csromanspacer{} ตํ กิสฺส เหตุ, มยมฺปิ หิ จกฺขุวิญฺเญยฺเยสุ รูเปสุ อวีตราคา อวีตโทสา อวีตโมหา อชฺฌตฺตํ อวูปสนฺตจิตฺตา สมวิสมํ จราม กาเยน วาจาย มนสา.\csromanspacer{} เตสํ โน สมจริยมฺปิ เหตํ อุตฺตริ อปสฺสตํ ตสฺมา เต โภนฺโต สมณพฺราหฺมณา น สกฺกาตพฺพา น ครุกาตพฺพา น มาเนตพฺพา น ปูเชตพฺพา.\csromanspacer{} เย เต สมณพฺราหฺมณา โสตวิญฺเญยฺเยสุ สทฺเทสุ.\csromanspacer{} ฆานวิญฺเญยฺเยสุ คนฺเธสุ.\csromanspacer{} ชิวฺหาวิญฺเญยฺเยสุ รเสสุ.\csromanspacer{} กายวิญฺเญยฺเยสุ โผฏฺฐพฺเพสุ.\csromanspacer{} มโนวิญฺเญยฺเยสุ ธมฺเมสุ อวีตราคา อวีตโทสา อวีตโมหา อชฺฌตฺตํ อวูปสนฺตจิตฺตา สมวิสมํ จรนฺติ กาเยน วาจาย มนสา, เอวรูปา สมณพฺราหฺมณา น สกฺกาตพฺพา น ครุกาตพฺพา น มาเนตพฺพา น ปูเชตพฺพา.\csromanspacer{} ตํ กิสฺส เหตุ, มยมฺปิ หิ มโนวิญฺเญยฺเยสุ ธมฺเมสุ อวีตราคา อวีตโทสา}

\csromanheader{2}{8. นครวินฺเทยฺยสุตฺต (150)}
\prose{\csromanfolio{341}อวีตโมหา อชฺฌตฺตํ อวูปสนฺตจิตฺตา สมวิสมํ จราม กาเยน วาจาย มนสา, เตสํ โน สมจริยมฺปิ เหตํ อุตฺตริ อปสฺสตํ.\csromanspacer{} ตสฺมา เต โภนฺโต สมณพฺราหฺมณา น สกฺกาตพฺพา น ครุกาตพฺพา น มาเนตพฺพา น ปูเชตพฺพา”ติ.\csromanspacer{} เอวํ ปุฏฺฐา ตุเมฺห คหปตโย เตสํ อญฺญติตฺถิยานํ ปริพฺพาชกานํ เอวํ พฺยากเรยฺยาถ.}

\proseitem{436}{สเจ ปน โว คหปตโย อญฺญติตฺถิยา ปริพฺพาชกา เอวํ ปุจฺเฉยฺยุํ “กถํภูตา คหปตโย สมณพฺราหฺมณา สกฺกาตพฺพา ครุกาตพฺพา มาเนตพฺพา ปูเชตพฺพา”ติ, เอวํ ปุฏฺฐา ตุเมฺห คหปตโย เตสํ อญฺญติตฺถิยานํ ปริพฺพาชกานํ เอวํ พฺยากเรยฺยาถ “เย เต สมณพฺราหฺมณา จกฺขุวิญฺเญยฺเยสุ รูเปสุ วีตราคา วีตโทสา วีตโมหา อชฺฌตฺตํ วูปสนฺตจิตฺตา สมจริยํ จรนฺติ กาเยน วาจาย มนสา, เอวรูปา สมณพฺราหฺมณา สกฺกาตพฺพา ครุกาตพฺพา มาเนตพฺพา ปูเชตพฺพา.\csromanspacer{} ตํ กิสฺส เหตุ, มยมฺปิ หิ\footnote{มยํ หิ (?)} จกฺขุวิญฺเญยฺเยสุ รูเปสุ อวีตราคา อวีตโทสา อวีตโมหา อชฺฌตฺตํ อวูปสนฺตจิตฺตา สมวิสมํ จราม กาเยน วาจาย มนสา.\csromanspacer{} เตสํ โน สมจริยมฺปิ เหตํ อุตฺตริ ปสฺสตํ ตสฺมา เต โภนฺโต สมณพฺราหฺมณา สกฺกาตพฺพา ครุกาตพฺพา มาเนตพฺพา ปูเชตพฺพา.\csromanspacer{} เย เต สมณพฺราหฺมณา โสตวิญฺเญยฺเยสุ สทฺเทสุ.\csromanspacer{} ฆานวิญฺเญยฺเยสุ คนฺเธสุ.\csromanspacer{} ชิวฺหาวิญฺเญยฺเยสุ รเสสุ.\csromanspacer{} กายวิญฺเญยฺเยสุ โผฏฺฐพฺเพสุ.\csromanspacer{} มโนวิญฺเญยฺเยสุ ธมฺเมสุ วีตราคา วีตโทสา วีตโมหา อชฺฌตฺตํ วูปสนฺตจิตฺตา สมจริยํ จรนฺติ กาเยน วาจาย มนสา, เอวรูปา สมณพฺราหฺมณา สกฺกาตพฺพา ครุกาตพฺพา มาเนตพฺพา ปูเชตพฺพา.\csromanspacer{} ตํ กิสฺส เหตุ, มยมฺปิ หิ มโนวิญฺเญยฺเยสุ ธมฺเมสุ อวีตราคา อวีตโทสา อวีตโมหา อชฺฌตฺตํ อวูปสนฺตจิตฺตา สมวิสมํ จราม กาเยน วาจาย มนสา, เตสํ โน สมจริยมฺปิ เหตํ อุตฺตริ ปสฺสตํ.\csromanspacer{} ตสฺมา เต โภนฺโต สมณาพฺราหฺมณา สกฺกาตพฺพา ครุกาตพฺพา มาเนตพฺพา ปูเชตพฺพา”ติ.\csromanspacer{} เอวํ ปุฏฺฐา ตุเมฺห คหปตโย เตสํ อญฺญติตฺถิยานํ ปริพฺพาชกานํ เอวํ พฺยากเรยฺยาถ.}

\proseitem{437}{สเจ ปน โว\footnote{สเจ เต (สฺยา, กํ, อิ, ก)} คหปตโย อญฺญติตฺถิยา ปริพฺพาชกา เอวํ ปุจฺเฉยฺยุํ “เก ปนายสฺมนฺตานํ อาการา เก อนฺวยา, เยน ตุเมฺห \csromanfolio{342}อายสฺมนฺโต เอวํ วเทถ ‘อทฺธา เต อายสฺมนฺโต วีตราคา วา ราควินยาย วา ปฏิปนฺนา, วีตโทสา วา โทสวินยาย วา ปฏิปนฺนา, วีตโมหา วา โมหวินยาย วา ปฏิปนฺนา’ติ”, เอวํ ปุฏฺฐา ตุเมฺห คหปตโย เตสํ อญฺญติตฺถิยานํ ปริพฺพาชกานํ เอวํ พฺยากเรยฺยาถ “ตถา หิ เต อายสฺมนฺโต อรญฺญวนปตฺถานิ ปนฺตานิ เสนาสนานิ ปฏิเสวนฺติ, นตฺถิ โข ปน ตตฺถ ตถารูปา จกฺขุวิญฺเญยฺยา รูปา, เย ทิสฺวา ทิสฺวา อภิรเมยฺยุํ.\csromanspacer{} นตฺถิ โข ปน ตตฺถ ตถารูปา โสตวิญฺเญยฺยา สทฺทา, เย สุตฺวา สุตฺวา อภิรเมยฺยุํ.\csromanspacer{} นตฺถิ โข ปน ตตฺถ ตถารูปา ฆานวิญฺเญยฺยา คนฺธา, เย ฆายิตฺวา ฆายิตฺวา อภิรเมยฺยุํ.\csromanspacer{} นตฺถิ โข ปน ตตฺถ ตถารูปา ชิวฺหาวิญฺเญยฺยา รสา, เย สายิตฺวา สายิตฺวา อภิรเมยฺยุํ.\csromanspacer{} นตฺถิ โข ปน ตตฺถ ตถารูปา กายวิญฺเญยฺยา โผฏฺฐพฺพา, เย ผุสิตฺวา ผุสิตฺวา อภิรเมยฺยุํ.\csromanspacer{} อิเม โข โน อาวุโส อาการา อิเม อนฺวยา, เยน มยํ\footnote{เยน มยํ อายสฺมนฺโต (สี, อิ), เยน มยํ อายสฺมนฺเต (สฺยา, กํ)} เอวํ วเทม ‘อทฺธา เต อายสฺมนฺโต วีตราคา วา ราควินยาย วา ปฏิปนฺนา, วีตโทสา วา โทสวินยาย วา ปฏิปนฺนา, วีตโมหา วา โมหวินยาย วา ปฏิปนฺนา’ติ”, เอวํ ปุฏฺฐา ตุเมฺห คหปตโย เตสํ อญฺญติตฺถิยานํ ปริพฺพาชกานํ เอวํ พฺยากเรยฺยาถาติ.}

\prose{เอวํ วุตฺเต นครวินฺเทยฺยกา พฺราหฺมณคหปติกา ภควนฺตํ เอตทโวจุํ “อภิกฺกนฺตํ โภ โคตม, อภิกฺกนฺตํ โภ โคตม, เสยฺยถาปิ โภ โคตม นิกฺกุชฺชิตํ วา อุกฺกุชฺเชยฺย, ปฏิจฺฉนฺนํ วา วิวเรยฺย, มูฬฺหสฺส วา มคฺคํ อาจิกฺเขยฺย, อนฺธกาเร วา เตลปชฺโชตํ ธาเรยฺย ‘จกฺขุมนฺโต รูปานิ ทกฺขนฺตี’ติ.\csromanspacer{} เอวเมวํ โภตา โคตเมน อเนกปริยาเยน ธมฺโม ปกาสิโต, เอเต มยํ ภวนฺตํ โคตมํ สรณํ คจฺฉาม ธมฺมญฺจ ภิกฺขุสํฆญฺจ, อุปาสเก โน ภวํ โคตโม ธาเรตุ อชฺชตคฺเค ปาณุเปเต สรณํ คเต”ติ.}

\nitthitamruled{นครวินฺเทยฺยสุตฺตํ นิฏฺฐิตํ อฏฺฐมํ.}
\csromanmatikaanchor{mtk.60}
\csromantocmarkref{subsection}{9. ปิณฺฑปาตปาริสุทฺธิสุตฺต}{mtk.60}
\bookmark[dest={mtk.60},level=2]{9. ปิณฺฑปาตปาริสุทฺธิสุตฺต}
\csromanheader{2}{9. ปิณฺฑปาตปาริสุทฺธิสุตฺต}
\proseitem{438}{เอวํ เม สุตํ\csromandash{}เอกํ สมยํ ภควา ราชคเห วิหรติ เวฬุวเน กลนฺทกนิวาเป.\csromanspacer{} อถ โข อายสฺมา สาริปุตฺโต สายนฺหสมยํ}

\csromanheader{2}{9. ปิณฺฑปาตปาริสุทฺธิสุตฺต (151)}
\prose{\csromanfolio{343}ปฏิสลฺลานา วุฏฺฐิโต เยน ภควา เตนุปสงฺกมิ, อุปสงฺกมิตฺวา ภควนฺตํ อภิวาเทตฺวา เอกมนฺตํ นิสีทิ.\csromanspacer{} เอกมนฺตํ นิสินฺนํ โข อายสฺมนฺตํ สาริปุตฺตํ ภควา เอตทโวจ\csromandash{}}

\prose{วิปฺปสนฺนานิ โข เต สาริปุตฺต อินฺทฺริยานิ, ปริสุทฺโธ ฉวิวณฺโณ ปริโยทาโต, กตเมน โข ตฺวํ สาริปุตฺต วิหาเรน เอตรหิ พหุลํ วิหรสีติ.\csromanspacer{} สุญฺญตาวิหาเรน โข อหํ ภนฺเต เอตรหิ พหุลํ วิหรามีติ.\csromanspacer{} สาธุ สาธุ สาริปุตฺต, มหาปุริสวิหาเรน กิร ตฺวํ สาริปุตฺต เอตรหิ พหุลํ วิหรสิ.\csromanspacer{} มหาปุริสวิหาโร เอโส\footnote{เหส (สี, สฺยา, กํ, อิ)} สาริปุตฺต ยทิทํ สุญฺญตา.\csromanspacer{} ตสฺมาติห สาริปุตฺต ภิกฺขุ สเจ อากงฺเขยฺย “สุญฺญตาวิหาเรน พหุลํ\footnote{เอตรหิ พหุลํ (สี, อิ)}, วิหเรยฺยนฺ”ติ.\csromanspacer{} เตน สาริปุตฺต ภิกฺขุนา อิติ ปฏิสญฺจิกฺขิตพฺพํ “เยน จาหํ มคฺเคน คามํ ปิณฺฑาย ปาวิสิํ, ยสฺมิํ จ ปเทเส ปิณฺฑาย อจริํ, เยน จ มคฺเคน คามโต ปิณฺฑาย ปฏิกฺกมิํ.\csromanspacer{} อตฺถิ นุ โข เม ตตฺถ จกฺขุวิญฺเญยฺเยสุ รูเปสุ ฉนฺโท วา ราโค วา โทโส วา โมโห วา ปฏิฆํ วาปิ เจตโส”ติ.\csromanspacer{} สเจ สาริปุตฺต ภิกฺขุ ปจฺจเวกฺขมาโน เอวํ ชานาติ “เยน จาหํ มคฺเคน คามํ ปิณฺฑาย ปาวิสิํ, ยสฺมิํ จ ปเทเส ปิณฺฑาย อจริํ, เยน จ มคฺเคน คามโต ปิณฺฑาย ปฏิกฺกมิํ.\csromanspacer{} อตฺถิ เม ตตฺถ จกฺขุวิญฺเญยฺเยสุ รูเปสุ ฉนฺโท วา ราโค วา โทโส วา โมโห วา ปฏิฆํ วาติ เจตโส”ติ, เตน สาริปุตฺต ภิกฺขุนา เตสํเยว ปาปกานํ อกุสลานํ ธมฺมานํ ปหานาย วายมิตพฺพํ.\csromanspacer{} สเจ ปน สาริปุตฺต ภิกฺขุ ปจฺจเวกฺขมาโน เอวํ ชานาติ “เยน จายํ มคฺเคน คามํ ปิณฺฑาย ปาวิสิํ, ยสฺมิํ จ ปเทเส ปิณฺฑาย อจริํ, เยน จ มคฺเคน คามโต ปิณฺฑาย ปฏิกฺกมิํ.\csromanspacer{} นตฺถิ เม ตตฺถ จกฺขุวิญฺเญยฺเยสุ รูเปสุ ฉนฺโท วา ราโค วา โทโส วา โมโห วา ปฏิฆํ วาปิ เจตโส”ติ, เตน สาริปุตฺต ภิกฺขุนา เตเนว ปีติปาโมชฺเชน วิหาตพฺพํ อโหรตฺตานุสิกฺขินา กุสเลสุ ธมฺเมสุ.}

\proseitem{439}{ปุน จปรํ สาริปุตฺต ภิกฺขุนา อิติ ปฏิสญฺจิกฺขิตพฺพํ “เยน จาหํ มคฺเคน คามํ ปิณฺฑาย ปาวิสิํ, ยสฺมิํ จ ปเทเส ปิณฺฑาย อจริํ, เยน จ มคฺเคน คามโต ปิณฺฑาย ปฏิกฺกมิํ.\csromanspacer{} อตฺถิ นุ โข เม ตตฺถ โสตวิญฺเญยฺเยสุ สทฺเทสุ ฯเปฯ.\csromanspacer{} ฆานวิญฺเญยฺเยสุ คนฺเธสุ.\csromanspacer{} ชิวฺหาวิญฺเญยฺเยสุ \csromanfolio{344}รเสสุ.\csromanspacer{} กายวิญฺเญยฺเยสุ โผฏฺฐพฺเพสุ.\csromanspacer{} มโนวิญฺเญยฺเยสุ ธมฺเมสุ ฉนฺโท วา ราโค วา โทโส วา โมโห วา ปฏิฆํ วาปิ เจตโส”ติ.\csromanspacer{} สเจ สาริปุตฺต ภิกฺขุ ปจฺจเวกฺขมาโน เอวํ ชานาติ “เยน จาหํ มคฺเคน คามํ ปิณฺฑาย ปาวิสิํ, ยสฺมิํ จ ปเทเส ปิณฺฑาย อจริํ, เยน จ มคฺเคน คามโต ปิณฺฑาย ปฏิกฺกมิํ.\csromanspacer{} อตฺถิ เม ตตฺถ มโนวิญฺเญยฺเยสุ ธมฺเมสุ ฉนฺโท วา ราโค วา โทโส วา โมโห วา ปฏิฆํ วาปิ เจตโส”ติ, เตน สาริปุตฺต ภิกฺขุนา เตสํเยว ปาปกานํ อกุสลานํ ธมฺมานํ ปหานาย วายมิตพฺพํ.\csromanspacer{} สเจ ปน สาริปุตฺต ภิกฺขุ ปจฺจเวกฺขมาโน เอวํ ชานาติ “เยน จาหํ มคฺเคน คามํ ปิณฺฑาย ปาวิสิํ, ยสฺมิํ จ ปเทเส ปิณฺฑาย อจริํ, เยน จ มคฺเคน คามโต ปิณฺฑาย ปฏิกฺกมิํ.\csromanspacer{} นตฺถิ เม ตตฺถ มโนวิญฺเญยฺเยสุ ธมฺเมสุ ฉนฺโท วา ราโค วา โทโส วา โมโห วา ปฏิฆํ วาปิ เจตโส”ติ, เตน สาริปุตฺต ภิกฺขุนา เตนว ปีติปาโมชฺเชน วิหาตพฺพํ อโหรตฺตานุสิกฺขินา กุสเลสุ ธมฺเมสุ.}

\proseitem{440}{ปุน จปรํ สาริปุตฺต ภิกฺขุนา อิติ ปฏิสญฺจิกฺขิตพฺพํ “ปหีนา นุ โข เม ปญฺจ กามคุณา”ติ.\csromanspacer{} สเจ สาริปุตฺต ภิกฺขุ ปจฺจเวกฺขมาโน เอวํ ชานาติ “อปฺปหีนา โข เม ปญฺจ กามคุณา”ติ, เตน สาริปุตฺต ภิกฺขุนา ปญฺจนฺนํ กามคุณานํ ปหานาย วายมิตพฺพํ.\csromanspacer{} สเจ ปน สาริปุตฺต ภิกฺขุ ปจฺจเวกฺขมาโน เอวํ ชานาติ “ปหีนา โข เม ปญฺจ กามคุณา”ติ, เตน สาริปุตฺต ภิกฺขุนา เตเนว ปีติปาโมชฺเชน วิหาตพฺพํ อโหรตฺตานุสิกฺขินา กุสเลสุ ธมฺเมสุ.}

\proseitem{441}{ปุน จปรํ สาริปุตฺต ภิกฺขุนา อิติ ปฏิสญฺจิกฺขิตพฺพํ “ปหีนา นุ โข เม ปญฺจ นีวรณา”ติ.\csromanspacer{} สเจ สาริปุตฺต ภิกฺขุ ปจฺจเวกฺขมาโน เอวํ ชานาติ “อปฺปหีนา โข เม ปญฺจ นีวรณา”ติ, เตน สาริปุตฺต ภิกฺขุนา ปญฺจนฺนํ นีวรณานํ ปหานาย วายมิตพฺพํ.\csromanspacer{} สเจ ปน สาริปุตฺต ภิกฺขุนา ปจฺจเวกฺขมาโน เอวํ ชานาติ “ปหีนา โข เม ปญฺจ นีวรณา”ติ, เตน สาริปุตฺต ภิกฺขุนา เตเนว ปีติปาโมชฺเชน วิหาตพฺพํ อโหรตฺตานุสิกฺขินา กุสเลสุ ธมฺเมสุ.}

\proseitem{442}{ปุน จปรํ สาริปุตฺต ภิกฺขุนา อิติ ปฏิสญฺจิกฺขิตพฺพํ “ปริญฺญาตา นุ โข เม ปญฺจุปาทานกฺขนฺธา”ติ.\csromanspacer{} สเจ สาริปุตฺต ภิกฺขุ ปจฺจเวกฺขมาโน เอวํ ชานาติ “อปริญฺญาตา โข เม ปญฺจุปาทานกฺขนฺธา”ติ, เตน สาริปุตฺต}

\csromanheader{2}{9. ปิณฺฑปาตปาริสุทฺธิสุตฺต (151)}
\prose{\csromanfolio{345}ภิกฺขุนา ปญฺจนฺนํ อุปาทานกฺขนฺธานํ ปริญฺญาย วายมิตพฺพํ.\csromanspacer{} สเจ ปน สาริปุตฺต ภิกฺขุ ปจฺจเวกฺขมาโน เอวํ ชานาติ “ปริญฺญาตา โข เม ปญฺจุปาทานกฺขนฺธา”ติ, เตน สาริปุตฺต ภิกฺขุนา เตเนว ปีติปาโมชฺเชน วิหาตพฺพํ อโหรตฺตานุสิกฺขินา กุสเลสุ ธมฺเมสุ.}

\proseitem{443}{ปุน จปรํ สาริปุตฺต ภิกฺขุนา อิติ ปฏิสญฺจิกฺขิตพฺพํ “ภาวิตา นุ โข เม จตฺตาโร สติปฏฺฐานา”ติ.\csromanspacer{} สเจ สาริปุตฺต ภิกฺขุ ปจฺจเวกฺขมาโน เอวํ ชานาติ “อภาวิตา โข เม จตฺตาโร สติปฏฺฐานา”ติ, เตน สาริปุตฺต ภิกฺขุนา จตุนฺนํ สติปฏฺฐานานํ ภาวนาย วายมิตพฺพํ.\csromanspacer{} สเจ ปน สาริปุตฺต ภิกฺขุ ปจฺจเวกฺขมาโน เอวํ ชานาติ “ภาวิตา โข เม จตฺตาโร สติปฏฺฐานา”ติ, เตน สาริปุตฺต ภิกฺขุนา เตเนว ปีติปาโมชฺเชน วิหาตพฺพํ อโหรตฺตานุสิกฺขินา กุสเลสุ ธมฺเมสุ.}

\proseitem{444}{ปุน จปรํ สาริปุตฺต ภิกฺขุนา อิติ ปฏิสญฺจิกฺขิตพฺพํ “ภาวิตา นุ โข เม จตฺตาโร สมฺมปฺปธานา”ติ.\csromanspacer{} สเจ สาริปุตฺต ภิกฺขุ ปจฺจเวกฺขมาโน เอวํ ชานาติ “อภาวิตา โข เม จตฺตาโร สมฺมปฺปธานา”ติ, เตน สาริปุตฺต ภิกฺขุนา จตุนฺนํ สมฺมปฺปธานานํ ภาวนาย วายมิตพฺพํ.\csromanspacer{} สเจ ปน สาริปุตฺต ภิกฺขุ ปจฺจเวกฺขมาโน เอวํ ชานาติ “ภาวิตา โข เม จตฺตาโร สมฺมปฺปธานา”ติ, เตน สาริปุตฺต ภิกฺขุนา เตเนว ปีติปาโมชฺเชน วิหาตพฺพํ อโหรตฺตานุสิกฺขินา กุสเลสุ ธมฺเมสุ.}

\proseitem{445}{ปุน จปรํ สาริปุตฺต ภิกฺขุนา อิติ ปฏิสญฺจิกฺขิตพฺพํ “ภาวิตา นุ โข เม จตฺตาโร อิทฺธิปาทา”ติ.\csromanspacer{} สเจ สาริปุตฺต ภิกฺขุ ปจฺจเวกฺขมาโน เอวํ ชานาติ “อภาวิตา โข เม จตฺตาโร อิทฺธิปาทา”ติ, เตน สาริปุตฺต ภิกฺขุนา จตุนฺนํ อิทฺธิปาทานํ ภาวนาย วายมิตพฺพํ.\csromanspacer{} สเจ ปน สาริปุตฺต ภิกฺขุ ปจฺจเวกฺขมาโน เอวํ ชานาติ “ภาวิตา โข เม จตฺตาโร อิทฺธิปาทา”ติ, เตน สาริปุตฺต ภิกฺขุนา เตเนว ปีติปาโมชฺเชน วิหาตพฺพํ อโหรตฺตานุสิกฺขินา กุสเลสุ ธมฺเมสุ.}

\proseitem{446}{ปุน จปรํ สาริปุตฺต ภิกฺขุนา อิติ ปฏิสญฺจิกฺขิตพฺพํ “ภาวิตานิ นุ โข เม ปญฺจินฺทฺริยานี”ติ.\csromanspacer{} สเจ สาริปุตฺต ภิกฺขุ ปจฺจเวกฺขมาโน เอวํ ชานาติ “อภาวิตานิ โข เม ปญฺจินฺทฺริยานี”ติ, เตน สาริปุตฺต ภิกฺขุนา ปญฺจนฺนํ อินฺทฺริยานํ ภาวนาย วายมิตพฺพํ.\csromanspacer{} สเจ ปน สาริปุตฺต ภิกฺขุ \csromanfolio{346}ปจฺจเวกฺขมาโน เอวํ ชานาติ “ภาวิตานิ โข เม ปญฺจินฺทฺริยานี”ติ, เตน สาริปุตฺต ภิกฺขุนา เตเนว ปีติปาโมชฺเชน วิหาตพฺพํ อโหรตฺตานุสิกฺขินา กุสเลสุ ธมฺเมสุ.}

\proseitem{447}{ปุน จปรํ สาริปุตฺต ภิกฺขุนา อิติ ปฏิสญฺจิกฺขิตพฺพํ “ภาวิตานิ นุ โข เม ปญฺจ พลานี”ติ.\csromanspacer{} สเจ สาริปุตฺต ภิกฺขุ ปจฺจเวกฺขมาโน เอวํ ชานาติ “อภาวิตานิ โข เม ปญฺจ พลานี”ติ, เตน สาริปุตฺต ภิกฺขุนา ปญฺจนฺนํ พลานํ ภาวนาย วายมิตพฺพํ.\csromanspacer{} สเจ ปน สาริปุตฺต ภิกฺขุ ปจฺจเวกฺขมาโน เอวํ ชานาติ “ภาวิตานิ โข เม ปญฺจ พลานี”ติ, เตน สาริปุตฺต ภิกฺขุนา เตเนน ปีติปาโมชฺเชน วิหาตพฺพํ อโหรตฺตานุสิกฺขินา กุสเลสุ ธมฺเมสุ.}

\proseitem{448}{ปุน จปรํ สาริปุตฺต ภิกฺขุนา อิติ ปฏิสญฺจิกฺขิตพฺพํ “ภาวิตา นุ โข เม สตฺต โพชฺฌงฺคา”ติ.\csromanspacer{} สเจ สาริปุตฺต ภิกฺขุ ปจฺจเวกฺขมาโน เอวํ ชานาติ “อภาวิตา โข เม สตฺต โพชฺฌงฺคา”ติ, เตน สาริปุตฺต ภิกฺขุนา สตฺตนฺนํ โพชฺฌงฺคานํ ภาวนาย วายมิตพฺพํ.\csromanspacer{} สเจ ปน สาริปุตฺต ภิกฺขุ ปจฺจเวกฺขมาโน เอวํ ชานาติ “ภาวิตา โข เม สตฺต โพชฺฌงฺคา”ติ, เตน สาริปุตฺต ภิกฺขุนา เตเนว ปีติปาโมชฺเชน วิหาตพฺพํ อโหรตฺตานุสิกฺขินา กุสเลสุ ธมฺเมสุ.}

\proseitem{449}{ปุน จปรํ สาริปุตฺต ภิกฺขุนา อิติ ปฏิสญฺจิกฺขิตพฺพํ “ภาวิโต นุ โข เม อริโย อฏฺฐงฺคิโก มคฺโค”ติ.\csromanspacer{} สเจ สาริปุตฺต ภิกฺขุ ปจฺจเวกฺขมาโน เอวํ ชานาติ “อภาวิโต โข เม อริโย อฏฺฐงฺคิโก มคฺโค”ติ, เตน สาริปุตฺต ภิกฺขุนา อริยสฺส อฏฺฐงฺคิกสฺส มคฺคสฺส ภาวนาย วายมิตพฺพํ.\csromanspacer{} สเจ ปน สาริปุตฺต ภิกฺขุ ปจฺจเวกฺขมาโน เอวํ ชานาติ “ภาวิโต โข เม อริโย อฏฺฐงฺคิโก มคฺโค”ติ, เตน สาริปุตฺต ภิกฺขุนา เตเนว ปีติปาโมชฺเชน วิหาตพฺพํ อโหรตฺตานุสิกฺขินา กุสเลสุ ธมฺเมสุ.}

\proseitem{450}{ปุน จปรํ สาริปุตฺต ภิกฺขุนา อิติ ปฏิสญฺจิกฺขิตพฺพํ “ภาวิตา นุ โข เม สมโถ จ วิปสนา จา”ติ.\csromanspacer{} สเจ สาริปุตฺต ภิกฺขุ ปจฺจเวกฺขมาโน เอวํ ชานาติ “อภาวิตา โข เม สมโถ จ วิปสฺสนา จา”ติ, เตน สาริปุตฺต ภิกฺขุนา สมถวิปสฺสนานํ ภาวนาย วายมิตพฺพํ.\csromanspacer{} สเจ ปน สาริปุตฺต ภิกฺขุ ปจฺจเวกฺขมาโน เอวํ ชานาติ “ภาวิตา}

\csromanheader{2}{10. อินฺทฺริยภาวนาสุตฺต (152)}
\prose{\csromanfolio{347}โข เม สมโถ จ วิปสฺสนา จา”ติ, เตน สาริปุตฺต ภิกฺขุนา เตเนว ปีติปาโมชฺเชน วิหาตพฺพํ อโหรตฺตานุสิกฺขินา กุสเลสุ ธมฺเมสุ.}

\proseitem{451}{ปุน จปรํ สาริปุตฺต ภิกฺขุนา อิติ ปฏิสญฺจิกฺขิตพฺพํ “สจฺฉิกตา นุ โข เม วิชฺชา จ วิมุตฺติ จา”ติ.\csromanspacer{} สเจ สาริปุตฺต ภิกฺขุ ปจฺจเวกฺขมาโน เอวํ ชานาติ “อสจฺฉิกตา โข เม วิชฺชา จ วิมุตฺติ จา”ติ, เตน สาริปุตฺต ภิกฺขุนา วิชฺชาย วิมุตฺติยา สจฺฉิกิริยาย วายมิตพฺพํ.\csromanspacer{} สเจ ปน สาริปุตฺต ภิกฺขุ ปจฺจเวกฺขมาโน เอวํ ชานาติ “สจฺฉิกตา โข เม วิชฺชา จ วิมุตฺติ จา”ติ, เตน สาริปุตฺต ภิกฺขุนา เตเนว ปีติปาโมชฺเชน วิหาตพฺพํ อโหรตฺตานุสิกฺขินา กุสเลสุ ธมฺเมสุ.}

\proseitem{452}{เย หิ เกจิ สาริปุตฺต อตีตมทฺธานํ สมณา วา พฺราหฺมณา วา ปิณฺฑปาตํ ปริโสเธสุํ, สพฺเพ เต เอวเมว ปจฺจเวกฺขิตฺวา ปจฺจเวกฺขิตฺวา ปิณฺฑปาตํ ปริโสเธสุํ.\csromanspacer{} เยปิ หิ เกจิ สาริปุตฺต อนาคตมทฺธานํ สมณา วา พฺราหฺมณา วา ปิณฺฑปาตํ ปริโสเธสฺสนฺติ, สพฺเพ เต เอวเมว ปจฺจเวกฺขิตฺวา ปจฺจเวกฺขิตฺวา ปิณฺฑปาตํ ปริโสเธสฺสนฺติ.\csromanspacer{} เยปิ หิ เกจิ สาริปุตฺต เอตรหิ สมณา วา พฺราหฺมณา วา ปิณฺฑปาตํ ปริโสเธนฺติ, สพฺเพ เต เอวเมว ปจฺจเวกฺขิตฺวา ปจฺจเวกฺขิตฺวา ปิณฺฑปาตํ ปริโสเธนฺติ.\csromanspacer{} ตสฺมาติห สาริปุตฺต\footnote{โว สาริปุตฺต เอวํ สิกฺขิตพฺพํ (สี, อิ)} “ปจฺจเวกฺขิตฺวา ปจฺจเวกฺขิตฺวา ปิณฺฑปาตํ ปริโสเธสฺสามา”ติ เอวํ หิ โว สาริปุตฺต สิกฺขิตพฺพนฺติ.}

\prose{อิทมโวจ ภควา.\csromanspacer{} อตฺตมโน อายสฺมา สาริปุตฺตา ภควโต ภาสิตํ อภินนฺทีติ.}

\nitthitamruled{ปิณฺฑปาตปาริสุทฺธิสุตฺตํ นิฏฺฐิตํ นวมํ.}
\csromanmatikaanchor{mtk.61}
\csromantocmarkref{subsection}{10. อินฺทฺริยภาวนาสุตฺต}{mtk.61}
\bookmark[dest={mtk.61},level=2]{10. อินฺทฺริยภาวนาสุตฺต}
\csromanheader{2}{10. อินฺทฺริยภาวนาสุตฺต}
\proseitem{453}{เอวํ เม สุตํ\csromandash{}เอกํ สมยํ ภควา คชงฺคลายํ\footnote{กชงฺคลายํ (สี, อิ), กชฺชงฺคลายํ (สฺยา, กํ)} วิหรติ สุเวฬุวเน\footnote{เวฬุวเน (สฺยา, กํ), มุเขลุวเน (สี, อิ)}.\csromanspacer{} อถ โข อุตฺตโร มาณโว ปาราสิวิยนฺเตวาสี\footnote{ปาราสริยนฺเตวาสี (สี, อิ), ปาราสิริยนฺเตวาสี (สฺยา, กํ)} เยน \csromanfolio{348}ภควา เตนุปสงฺกมิ, อุปสงฺกมิตฺวา ภควตา สทฺธิํ สมฺโมทิ, สมฺโมทนียํ กถํ สารณียํ วีติสาเรตฺวา เอกมนฺตํ นิสีทิ.\csromanspacer{} เอกมนฺตํ นิสินฺนํ โข อุตฺตรํ มาณวํ ปาราสิวิยนฺเตวาสิํ ภควา เอตทโวจ “เทเสติ อุตฺตร ปาราสิวิโย พฺราหฺมโณ สาวกานํ อินฺทฺริยภาวนนฺ”ติ.\csromanspacer{} เทเสติ โภ โคตม ปาราสิวิโย พฺราหฺมโณ สาวกานํ อินฺทฺริยภาวนนฺติ.\csromanspacer{} ยถา กถํ ปน อุตฺตร เทเสติ ปาราสิวิโย พฺราหฺมโณ สาวกานํ อินฺทฺริยภาวนนฺติ.\csromanspacer{} อิธ โภ โคตม จกฺขุนา รูปํ น ปสฺสติ, โสเตน สทฺทํ น สุณาติ.\csromanspacer{} เอวํ โข โภ โคตม เทเสติ ปาราสิวิโย พฺราหฺมโณ สาวกานํ อินฺทฺริยภาวนนฺติ.\csromanspacer{} เอวํ สนฺเต โข อุตฺตร อนฺโธ ภาวิตินฺทฺริโย ภวิสฺสติ, พธิโร ภาวิตินฺทฺริโย ภวิสฺสติ, ยถา ปาราสิวิยสฺส พฺราหฺมณสฺส วจนํ.\csromanspacer{} อนฺโธ หิ อุตฺตร จกฺขุนา รูปํ น ปสฺสติ, พธิโร โสเตน สทฺทํ น สุณาตีติ.\csromanspacer{} เอวํ วุตฺเต อุตฺตโร มาณโว ปาราสิวิยนฺเตสี ตุณฺหีภูโต มงฺกุภูโต ปตฺตกฺขนฺโธ อโธมุโข ปชฺฌายนฺโต อปฺปฏิภาโน นิสีทิ.}

\prose{อถ โข ภควา อุตฺตรํ มาณวํ ปาราสิวิยนฺเตวาสิํ ตุณฺหีภูตํ มงฺกุภูตํ ปตฺตกฺขนฺธํ อโธมุขํ ปชฺโฌยนฺตํ อปฺปฏิภานํ วิทิตฺวา อายสฺมนฺตํ อานนฺทํ อามนฺเตสิ “อญฺญถา โข อานนฺท เทเสติ ปาราสิวิโย พฺราหฺมโณ สาวกานํ อินฺทฺริยภาวนํ, อญฺญถา จ ปนานนฺท อริยสฺส วินเย อนุตฺตรา อินฺทฺริยภาวนา โหตี”ติ.\csromanspacer{} เอตสฺส ภควา กาโล, เอตสฺส สุคต กาโล, ยํ ภควา อริยสฺส วินเย อนุตฺตรํ อินฺทฺริยภาวนํ เทเสยฺย, ภควโต สุตฺวา ภิกฺขู ธาเรสฺสนฺตีติ.\csromanspacer{} เตนหานนฺท สุณาหิ สาธุกํ มนสิ กโรหิ, ภาสิสฺสามีติ. “เอวํ ภนฺเต”ติ โข อายสฺมา อานนฺโท ภควโต ปจฺจสฺโสสิ.\csromanspacer{} ภควา เอตทโวจ\csromandash{}}

\proseitem{454}{กถญฺจานนฺท อริยสฺส วินเย อนุตฺตรา อินฺทฺริยภาวนา โหติ, อิธานนฺท ภิกฺขุโน จกฺขุนา รูปํ ทิสฺวา อุปฺปชฺชติ มนาปํ อุปฺปชฺชติ อมนาปํ อุปฺปชฺชติ มนาปามนาปํ, โส เอวํ ปชานาติ “อุปฺปนฺนํ โข เม อิทํ มนาปํ อุปฺปนฺนํ อมนาปํ อุปฺปนฺนํ มนาปามนาปํ, ตญฺจ โข สงฺขตํ โอฬาริกํ ปฏิจฺจสมุปฺปนฺนํ, เอตํ สนฺตํ เอตํ ปณีตํ ยทิทํ อุเปกฺขา”ติ.\csromanspacer{} ตสฺส ตํ อุปฺปนฺนํ มนาปํ อุปฺปนฺนํ อมนาปํ อุปฺปนฺนํ มนาปามนาปํ นิรุชฺฌติ, อุเปกฺขา สณฺฐาติ.\csromanspacer{} เสยฺยถาปิ อานนฺท จกฺขุมา ปุริโส อุมฺมีเลตฺวา วา นิมีเลยฺย, นิมีเลตฺวา วา อุมฺมีเลยฺย.\csromanspacer{} เอวเมว โข อานนฺท ยสฺส กสฺสจิ เอวํสีฆํ}

\csromanheader{2}{10. อินฺทฺริยภาวนาสุตฺต (152)}
\prose{\csromanfolio{349}เอวํตุวฏํ เอวํอปฺปกสิเรน อุปฺปนฺนํ มนาปํ อุปฺปนฺนํ อมนาปํ อุปฺปนฺนํ มนาปามนาปํ นิรุชฺฌติ, อุเปกฺขา สณฺฐาติ.\csromanspacer{} อยํ วุจฺจตานนฺท อริยสฺส วินเย อนุตฺตรา อินฺทฺริยภาวนา จกฺขุวิญฺเญยฺเยสุ รูเปสุ.}

\proseitem{455}{ปุน จปรํ อานนฺท ภิกฺขุโน โสเตน สทฺทํ สุตฺวา อุปฺปชฺชติ มนาปํ อุปฺปชฺชติ อมนาปํ อุปฺปชฺชติ มนาปามนาปํ, โส เอวํ ปชานาติ “อุปฺปนฺนํ โข เม อิทํ มนาปํ อุปฺปนฺนํ อมนาปํ อุปฺปนฺนํ มนาปามนาปํ, ตญฺจ โข สงฺขตํ โอฬาริกํ ปฏิจฺจสมุปฺปนฺนํ, เอตํ สนฺตํ เอตํ ปณีตํ ยทิทํ อุเปกฺขา”ติ.\csromanspacer{} ตสฺส ตํ อุปฺปนฺนํ มนาปํ อุปฺปนฺนํ อมนาปํ อุปฺปนฺนํ มนาปามนาปํ นิรุชฺฌติ, อุเปกฺขา สณฺฐาติ.\csromanspacer{} เสยฺยถาปิ อานนฺท พลวา ปุริโส อปฺปกสิเรเนว อจฺฉรํ\footnote{อจฺฉริกํ (สฺยา, กํ, อิ, ก)} ปหเรยฺย.\csromanspacer{} เอวเมว โข อานนฺท ยสฺส กสฺสจิ เอวํสีฆํ เอวํตุวฏํ เอวํอปฺปกสิเรน อุปฺปนฺนํ มนาปํ อุปฺปนฺนํ อมนาปํ อุปฺปนฺนํ มนาปามนาปํ นิรุชฺฌติ, อุเปกฺขา สณฺฐาติ.\csromanspacer{} อยํ วุจฺจตานนฺท อริยสฺส วินเย อนุตฺตรา อินฺทฺริยภาวนา โสตวิญฺเญยฺเยสุ สทฺเทสุ.}

\proseitem{456}{ปุน จปรํ อานนฺท ภิกฺขุโน ฆาเนน คนฺธํ ฆายิตฺวา อุปฺปชฺชติ มนาปํ อุปฺปชฺชติ อมนาปํ อุปฺปชฺชติ มนาปามนาปํ, โส เอวํ ปชานาติ “อุปฺปนฺนํ โข เม อิทํ มนาปํ อุปฺปนฺนํ อมนาปํ อุปฺปนฺนํ มนาปามนาปํ, ตญฺจ โข สงฺขตํ โอฬาริกํ ปฏิจฺจสมุปฺปนฺนํ, เอตํ สนฺตํ เอตํ ปณีตํ ยทิทํ อุเปกฺขา”ติ.\csromanspacer{} ตสฺส ตํ อุปฺปนฺนํ มนาปํ อุปฺปนฺนํ อมนาปํ อุปฺปนฺนํ มนาปามนาปํ นิรุชฺฌติ, อุเปกฺขา สณฺฐาติ.\csromanspacer{} เสยฺยถาปิ อานนฺท อีสกํโปเณ\footnote{¢สกโปเณ (สี, สฺยา, กํ, อิ), อีสกผเณ (สี-ฏฺฐ), “มชฺเฌ อุจฺจํ หุตฺวา”ติ ฏีกาย สํสนฺทิตพฺพา.} ปทุมปลาเส\footnote{ปทุมินิปตฺเต (สี, สฺยา, กํ, อิ)} อุทกผุสิตานิ ปวตฺตนฺติ น สณฺฐนฺติ.\csromanspacer{} เอวเมว โข อานนฺท ยสฺส กสฺสจิ เอวํสีฆํ เอวํตุวฏํ เอวํอปฺปกสิเรน อุปฺปนฺนํ มนาปํ อุปฺปนฺนํ อมนาปํ อุปฺปนฺนํ มนาปามนาปํ นิรุชฺฌติ, อุเปกฺขา สณฺฐาติ.\csromanspacer{} อยํ วุจฺจตานนฺท อริยสฺส วินเย อนุตฺตรา อินฺทฺริยภาวนา ฆานวิญฺเญยฺเยสุ คนฺเธสุ.}

\proseitem{457}{ปุน จปรํ อานนฺท ภิกฺขุโน ชิวฺหาย รสํ สายิตฺวา อุปฺปชฺชติ มนาปํ อุปฺปชฺชติ อมนาปํ อุปฺปชฺชติ มนาปามนาปํ, โส เอวํ ปชานาติ “อุปฺปนฺนํ โข เม อิทํ มนาปํ อุปฺปนฺนํ อมนาปํ อุปฺปนฺนํ มนาปามนาปํ, ตญฺจ โข สงฺขตํ โอฬาริกํ ปฏิจฺจสมุปฺปนฺนํ, เอตํ สนฺตํ เอตํ ปณีตํ ยทิทํ อุเปกฺขา”ติ.\csromanspacer{} ตสฺส ตํ อุปฺปนฺนํ มนาปํ อุปฺปนฺนํ อมนาปํ อุปฺปนฺนํ มนาปามนาปํ นิรุชฺฌติ, อุเปกฺขา \csromanfolio{350}สณฺฐาติ.\csromanspacer{} เสยฺยถาปิ อานนฺท พลวา ปุริโส ชิวฺหคฺเค เขฬปิณฺฑํ สํยูหิตฺวา อปฺปกสิเรน วเมยฺย\footnote{สนฺธเมยฺย (ก)}.\csromanspacer{} เอวเมว โข อานนฺท ยสฺส กสฺสจิ เอวํสีฆํ เอวํตุวฏํ เอวํอปฺปกสิเรน อุปฺปนฺนํ มนาปํ อุปฺปนฺนํ อมนาปํ อุปฺปนฺนํ มนาปามนาปํ นิรุชฺฌติ, อุเปกฺขา สณฺฐาติ.\csromanspacer{} อยํ วุจฺจตานนฺท อริยสฺส วินเย อนุตฺตรา อินฺทฺริยภาวนา ชิวฺหาวิญฺเญยฺเยสุ รเสสุ.}

\proseitem{458}{ปุน จปรํ อานนฺท ภิกฺขุโน กาเยน โผฏฺฐพฺพํ ผุสิตฺวา อุปฺปชฺชติ มนาปํ อุปฺปชฺชติ อมนาปํ อุปฺปชฺชติ มนาปามนาปํ, โส เอวํ ปชานาติ “อุปฺปนฺนํ โข เม อิทํ มนาปํ อุปฺปนฺนํ อมนาปํ อุปฺปนฺนํ มนาปามนาปํ, ตญฺจ โข สงฺขตํ โอฬาริกํ ปฏิจฺจสมุปฺปนฺนํ, เอตํ สนฺตํ เอตํ สนฺตํ ปณีตํ ยทิทํ อุเปกฺขา”ติ.\csromanspacer{} ตสฺส ตํ อุปฺปนฺนํ มนาปํ อุปฺปนฺนํ อมนาปํ อุปฺปนฺนํ มนาปามนาปํ นิรุชฺฌติ, อุเปกฺขา สณฺฐาติ.\csromanspacer{} เสยฺยถาปิ อานนฺท พลวา ปุริโส สมิญฺชิตํ วา พาหํ ปสาเรยฺย, ปสาริตํ วา พาหํ สมิญฺเชยฺย.\csromanspacer{} เอวเมว โข อานนฺท ยสฺส กสฺสจิ เอวํสีฆํ เอวํตุวฏํ เอวํ-อปฺปกสิเรน อุปฺปนฺนํ มนาปํ อุปฺปนฺนํ อมนาปํ อุปฺปนฺนํ มนาปามนาปํ นิรุชฺฌติ, อุเปกฺขา สณฺฐาติ.\csromanspacer{} อยํ วุจฺจตานนฺท อริยสฺส วินเย อนุตฺตรา อินฺทฺริยภาวนา กายวิญฺเญยฺเยสุ โผฏฺฐพฺเพสุ.}

\proseitem{459}{ปุน จปรํ อานนฺท ภิกฺขุโน มนสา ธมฺมํ วิญฺญาย อุปฺปชฺชติ มนาปํ อุปฺปชฺชติ อมนาปํ อุปฺปชฺชติ มนาปามนาปํ, โส เอวํ ปชานาติ “อุปฺปนฺนํ โข เม อิทํ มนาปํ อุปฺปนฺนํ อมนาปํ อุปฺปนฺนํ มนาปามนาปํ, ตญฺจ โข สงฺขตํ โอฬาริกํ ปฏิจฺจสมุปฺปนฺนํ, เอตํ สนฺตํ เอตํ ปณีตํ ยทิทํ อุเปกฺขา”ติ.\csromanspacer{} ตสฺส ตํ อุปฺปนฺนํ มนาปํ อุปฺปนฺนํ อมนาปํ อุปฺปนฺนํ มนาปามนาปํ นิรุชฺฌติ, อุเปกฺขา สณฺฐาติ.\csromanspacer{} เสยฺยถาปิ อานนฺท พลวา ปุริโส ทิวสํสนฺตตฺเต\footnote{ทิวสสนฺตตฺเต (สี)} อโยกฏาเห เทฺว วา ตีณิ วา อุทกผุสิตานิ นิปาเตยฺย, ทนฺโธ อานนฺท อุทกผุสิตานํ นิปาโต, อถ โข นํ ขิปฺปเมว ปริกฺขยํ ปริยาทานํ คจฺเฉยฺย.\csromanspacer{} เอวเมว โข อานนฺท ยสฺส กสฺสจิ เอวํสีฆํ เอวํตุวฏํ เอวํอปฺปกสิเรน อุปฺปนฺนํ มนาปํ อุปฺปนฺนํ อมนาปํ อุปฺปนฺนํ มนาปามนาปํ นิรุชฺฌติ, อุเปกฺขา สณฺฐาติ.\csromanspacer{} อยํ วุจฺจตานนฺท อริยสฺส วินเย อนุตฺตรา อินฺทฺริยภาวนา มโนวิญฺเญยฺเยสุ ธมฺเมสุ.\csromanspacer{} เอวํ โข อานนฺท อริยสฺส วินเย อนุตฺตรา อินฺทฺริยภาวนา โหติ.}

\proseitem{460}{กถญฺจานนฺท เสโข โหติ ปาฏิปโท, อิธานนฺท ภิกฺขุโน จกฺขุนา รูปํ ทิสฺวา อุปฺปชฺชติ มนาปํ อุปฺปชฺชติ อมนาปํ อุปฺปชฺชติ มนาปามนาปํ,}

\csromanheader{2}{10. อินฺทฺริยภาวนาสุตฺต (152)}
\prose{\csromanfolio{351}โส เตน อุปฺปนฺเนน มนาเปน อุปฺปนฺเนน อมนาเปน อุปฺปนฺเนน มนาปามนาเปน อฏฺฏียติ หรายติ ชิคุจฺฉติ.\csromanspacer{} โสเตน สทฺทํ สุตฺวา ฯเปฯ.\csromanspacer{} ฆาเนน คนฺธํ ฆายิตฺวา.\csromanspacer{} ชิวฺหาย รสํ สายิตฺวา.\csromanspacer{} กาเยน โผฏฺฐพฺพํ ผุสิตฺวา.\csromanspacer{} มนสา ธมฺมํ วิญฺญาย อุปฺปชฺชติ มนาปํ อุปฺปชฺชติ อมนาปํ อุปฺปชฺชติ มนาปามนาปํ, โส เตน อุปฺปนฺเนน มนาเปน อุปฺปนฺเนน อมนาเปน อุปฺปนฺเนน มนาปามนาเปน อฏฺฏียติ หรายติ ชิคุจฺฉติ.\csromanspacer{} เอวํ โข อานนฺท เสโข โหติ ปาฏิปโท.}

\proseitem{461}{กถญฺจานนฺท อริโย โหติ ภาวิตินฺทฺริโย, อิธานนฺท ภิกฺขุโน จกฺขุนา รูปํ ทิสฺวา อุปฺปชฺชติ มนาปํ อุปฺปชฺชติ อมนาปํ อุปฺปชฺชติ มนาปามนาปํ, โส สเจ อากงฺขติ “ปฏิกูเล\footnote{ปฏิกฺกูเล (สพฺพตฺถ)} อปฺปฏิกูลสญฺญี วิหเรยฺยนฺ”ติ, อปฺปฏิกูลสญฺญี ตตฺถ วิหรติ.\csromanspacer{} สเจ อากงฺขติ “อปฺปฏิกูเล ปฏิกูลสญฺญี วิหเรยฺยนฺ”ติ, ปฏิกูลสญฺญี ตตฺถ วิหรติ.\csromanspacer{} สเจ อากงฺขติ “ปฏิกูเล จ อปฺปฏิกูเล จ อปฺปฏิกูลสญฺญี วิหเรยฺยนฺ”ติ, อปฺปฏิกูลสญฺญี ตตฺถ วิหรติ.\csromanspacer{} สเจ อากงฺขติ “อปฺปฏิกูเล จ ปฏิกูเล จ ปฏิกูเล จ ปฏิกูลสญฺญี วิหเรยฺยนฺ”ติ, ปฏิกูลสญฺญี ตตฺถ วิหรติ.\csromanspacer{} สเจ อากงฺขติ “ปฏิกูลญฺจ อปฺปฏิกูลญฺจ ตทุภยํ อภินิวชฺเชตฺวา อุเปกฺขโก วิหเรยฺยํ สโต สมฺปชาโน”ติ, อุเปกฺขโก ตตฺถ วิหรติ สโต สมฺปชาโน.}

\proseitem{462}{ปุน จปรํ อานนฺท ภิกฺขุโน โสเตน สทฺทํ สุตฺวา ฯเปฯ.\csromanspacer{} ฆาเนน คนฺธํ ฆายิตฺวา.\csromanspacer{} ชิวฺหาย รสํ สายิตฺวา.\csromanspacer{} กาเยน โผฏฺฐพฺพํ ผุสิตฺวา.\csromanspacer{} มนสา ธมฺมํ วิญฺญาย อุปฺปชฺชติ มนาปํ อุปฺปชฺชติ อมนาปํ อุปฺปชฺชติ มนาปามนาปํ.\csromanspacer{} โส สเจ อากงฺขติ “ปฏิกูเล อปฺปฏิกูลสญฺญี วิหเรยฺยนฺ”ติ, อปฺปฏิกูลสญฺญี ตตฺถ วิหรติ.\csromanspacer{} สเจ อากงฺขติ ‘อปฺปฏิกูเล ปฏิกูลสญฺญี วิหเรยฺยนฺ”ติ, ปฏิกูลสญฺญี ตตฺถ วิหรติ.\csromanspacer{} สเจ อากงฺขติ “ปฏิกูเล จ อปฺปฏิกูเล จ อปฺปฏิกูลสญฺญี วิหเรยฺยนฺ”ติ, อปฺปฏิกูลสญฺญี ตตฺถ วิหรติ.\csromanspacer{} สเจ อากงฺขติ “อปฺปฏิกูเล จ ปฏิกูเล จ ปฏิกูลสญฺญี วิหเรยฺยนฺ”ติ, ปฏิกูลสญฺญี ตตฺถ วิหรติ.\csromanspacer{} สเจ อากงฺขติ “ปฏิกูลญฺจ อปฺปฏิกูลญฺจ ตทุภยมฺปิ อภินิวชฺเชตฺวา อุเปกฺขโก วิหเรยฺยํ สโต สมฺปชาโน”ติ, อุเปกฺขโก ตตฺถ วิหรติ สโต สมฺปชาโน.\csromanspacer{} เอวํ โข อานนฺท อริโย โหติ ภาวิตินฺทฺริโย.}

\csromanmatikaanchor{mtk.62}
\csromantocmarkref{subsection}{อุทฺทานคาถา}{mtk.62}
\bookmark[dest={mtk.62},level=2]{อุทฺทานคาถา}
\proseitem{463}{\csromanfolio{352}อิติ โข อานนฺท เทสิตา มยา อริยสฺส วินเย อนุตฺตรา อินฺทฺริยภาวนา, เทสิโต เสโข ปาฏิปโท, เทสิโต อริโย ภาวิตินฺทฺริโย.\csromanspacer{} ยํ โข อานนฺท สตฺถารา กรณียํ สาวกานํ หิเตสินา อนุกมฺปเกน อนุกมฺปํ อุปาทาย, กตํ โว ตํ มยา.\csromanspacer{} เอตานิ อานนฺท รุกฺขมูลานิ เอตานิ สุญฺญาคารานิ.\csromanspacer{} ฌายถานนฺท, มา ปมาทตฺถ, มา ปจฺฉา วิปฺปฏิสาริโน อหุวตฺถ, อยํ โว อมฺหากํ อนุสาสนีติ.}

\prose{อิทมโวจ ภควา.\csromanspacer{} อตฺตมโน อายสฺมา อานนฺโท ภควโต ภาสิตํ อภินนฺทีติ.}

\nitthitam{อินฺทฺริยภาวนาสุตฺตํ นิฏฺฐิตํ ทสมํ.}
\nitthitam{\textbf{สฬายตนวคฺโค นิฏฺฐิโต ปญฺจโม.}}
\titlehead{\textbf{ตสฺสุทฺทานํ}}
\csromangathameasure{อนาถปิณฺฑิโก ฉนฺโน, ปุณฺโณ นนฺทกราหุลา. \\ ฉฉกฺกํ สฬายตนิกํ, นครวินฺเทยฺยสุทฺธิกา. \\ อินฺทฺริยภาวนา จาปิ, วคฺโค โอวาทปญฺจโมติ.}
\csromangathagroup{\csromangathastanza{อนาถปิณฺฑิโก ฉนฺโน, ปุณฺโณ นนฺทกราหุลา. \\ ฉฉกฺกํ สฬายตนิกํ, นครวินฺเทยฺยสุทฺธิกา. \\ อินฺทฺริยภาวนา จาปิ, วคฺโค โอวาทปญฺจโมติ.}}

\csromanheader{2}{\textbf{อิทํ วคฺคานมุทฺทานํ}}
\csromangathameasure{เทวทโหนุปโท จ, สุญฺญโต จ วิภงฺคโก. \\ สฬายตโนติ วคฺคา, อุปริปณฺณาสเก ฐิตาติ. \\ \textbf{อุปริปณฺณาสกํ สมตฺตํ.} \textbf{ตีหิ ปณฺณาสเกหิ ปฏิมณฺฑิโต สกโล มชฺฌิมนิกาโย สมตฺโต.}}
\csromangathagroup{\csromangathastanza{เทวทโหนุปโท จ, สุญฺญโต จ วิภงฺคโก. \\ สฬายตโนติ วคฺคา, อุปริปณฺณาสเก ฐิตาติ. \\ \textbf{อุปริปณฺณาสกํ สมตฺตํ.} \textbf{ตีหิ ปณฺณาสเกหิ ปฏิมณฺฑิโต สกโล มชฺฌิมนิกาโย สมตฺโต.}}}

\orphannote{อนุปาทาวิโมกฺโข. ตยิทํ ภิกฺขเว ตถาคเตน อนุตฺตรํ สนฺติวรปทํ อภิสมฺพุทฺธํ, ยทิทํ ฉนฺนํ ผสฺสายตนานํ สมุทยญฺจ อตฺถงฺคมญฺจ อสฺสาทญฺจ อาทีนวญฺจ นิสฺสรณญฺจ ยถาภูตํ วิทิตฺวา อนุปาทาวิโมกฺโขติ (สี, สฺยา, กํ, อิ)}

\orphannote{( ) นตฺถิ สี-อิ-โปตฺถเกสุ}

